\documentclass[11pt,a4paper]{report}
\usepackage[T1]{fontenc}
\usepackage[utf8]{inputenc}
\usepackage[polish]{babel}
\usepackage{lmodern}
\usepackage{microtype}
\usepackage{amsmath,amssymb}
\usepackage[a4paper,margin=2.3cm]{geometry}
\usepackage{booktabs,tabularx,array}
\usepackage[dvipsnames]{xcolor}
\usepackage{listings}
\usepackage{hyperref}
\usepackage{enumitem}
\usepackage{fancyhdr}

\hypersetup{colorlinks=true,linkcolor=blue!45!black, citecolor=blue!45!black,urlcolor=MidnightBlue,
            pdftitle={C++: Wspolbieznosc i asynchronicznosc},
            pdfauthor={Material dydaktyczny}}
\setlength{\parindent}{0pt}
\setlength{\parskip}{5pt}
\setlist{itemsep=2pt,topsep=3pt}
\setcounter{tocdepth}{1}
\setcounter{secnumdepth}{3}
\pagestyle{fancy}
\fancyhf{}
\fancyhead[L]{\small Programowanie współbieżne w C++}
\fancyhead[R]{\small Rozdział 1}
\fancyfoot[C]{\thepage}
\renewcommand{\headrulewidth}{0.35pt}
\setlength{\headheight}{14pt}
\definecolor{codemid}{RGB}{245,247,250}
\definecolor{codecomment}{RGB}{67,111,82}
\definecolor{codekeyword}{RGB}{25,69,139}
\lstdefinestyle{cpp}{
    language=C++,
    basicstyle=\ttfamily\footnotesize,
    keywordstyle=\color{codekeyword}\bfseries,
    commentstyle=\color{codecomment},
    stringstyle=\color{BrickRed},
    backgroundcolor=\color{codemid},
    frame=single,
    rulecolor=\color{gray!30},
    numbers=left,
    numberstyle=\tiny\color{gray},
    numbersep=9pt,
    tabsize=4,
    breaklines=true,
    breakatwhitespace=true,
    showstringspaces=false,
    columns=fullflexible,
    keepspaces=true,
    upquote=true,
    xleftmargin=1.3em,
    framexleftmargin=1.3em,
    aboveskip=8pt,
    belowskip=7pt,
    captionpos=b
}
\lstset{style=cpp}
\newcommand{\term}[1]{\textbf{#1}}
\newcommand{\code}[1]{\lstinline[columns=fixed]!#1!}
\newcommand{\nota}[1]{\par\smallskip\noindent\textbf{Uwaga. }#1\par}


\title{\textbf{Programowanie w C++}}
\author{Materiał dydaktyczny}
\date{Wersja dydaktyczna --- C++11--C++23, wybrane nowości C++26}

\begin{document}

\maketitle
\tableofcontents
% Rozdzial 19. Wymaga makr \code, \term, \nota oraz pakietow:
% listings, tabularx, booktabs, enumitem, hyperref, amsmath.
% Fragment rozdzialu (bez preambuly). C++20; nowsze API zaznaczono.
\chapter{Wejście, wyjście, system plików i przetwarzanie tekstu}
\label{chap:cpp19-io}

\section{Wprowadzenie: dane, reprezentacja i granice odpowiedzialności}

Operacje wejścia i wyjścia (\emph{I/O}) przenoszą dane między programem
oraz jego otoczeniem: terminalem, plikiem, potokiem czy urządzeniem.
W odróżnieniu od zwykłych obliczeń mogą zakończyć się częściowym sukcesem,
błędem systemu albo blokowaniem. Dobrze zaprojektowany kod oddziela trzy
warstwy: \emph{pozyskanie bajtów}, \emph{parsowanie reprezentacji} oraz
\emph{przetworzenie obiektów}. Dzięki temu test parsera nie wymaga dysku,
a zmiana formatu pliku nie wymusza zmiany algorytmu numerycznego.

\begin{center}\small
\begin{tabularx}{\textwidth}{@{}p{0.28\textwidth}X@{}}
\toprule
Mechanizm & Typowe zastosowanie \\
\midrule
\code{iostream}, \code{fstream} & Formatowany tekst, prosta obsługa plików, strumienie. \\
\code{filesystem} & Ścieżki, katalogi, metadane, operacje na plikach. \\
\code{string}, \code{string_view} & Własność tekstu i nieposiadający go widok. \\
\code{charconv} & Szybkie parsowanie i formatowanie liczb bez ustawień regionalnych. \\
\code{format} (C++20) & Składanie tekstu według opisu formatowania. \\
\code{print} (C++23) & Wygodne formatowane wypisywanie; zależy od implementacji biblioteki. \\
\bottomrule
\end{tabularx}
\end{center}

\nota{Przykłady są pisane dla C++20. W tym rozdziale \code{std::filesystem}
nie oznacza, że każda operacja jest atomowa ani że istnieją uprawnienia
do każdego wskazanego pliku. Kod obsługujący niezaufane ścieżki wymaga
osobnej analizy bezpieczeństwa.}

\section{Strumienie: stan, formatowanie i buforowanie}
\label{sec:cpp19-streams}

\code{std::istream} jest interfejsem odczytu, a \code{std::ostream}
interfejsem zapisu. \code{std::cin}, \code{std::cout}, \code{std::cerr}
i \code{std::clog} to gotowe strumienie standardowe. Operatory
\code{>>} i \code{<<} wykonują zwykle operacje \emph{formatowane}.
Strumień przechowuje także bity stanu: \code{goodbit}, \code{eofbit},
\code{failbit} i \code{badbit}. Stan jest obserwowalny przez
\code{good()}, \code{eof()}, \code{fail()}, \code{bad()} i \code{rdstate()}.

\subsection{Poprawna pętla odczytu}

Nie należy pisać \code{while (!in.eof())}: wskaźnik końca pliku jest
ustawiany \emph{po nieudanej próbie odczytu}, a nie przed nią.
Warunkiem pętli powinien być sam odczyt.

\begin{lstlisting}[caption={Odczyt liczb i kontrola bledu formatu (kompletny program, C++20).}]
#include <iostream>
#include <sstream>

int main() {
    std::istringstream in("10 20 30");
    int value = 0;
    int sum = 0;
    while (in >> value)
        sum += value;

    if (!in.eof()) {
        std::cerr << "Invalid input\n";
        return 1;
    }
    std::cout << sum << '\n'; // 60
}
\end{lstlisting}

Odczyt \code{int} ignoruje początkowe białe znaki. Gdy napotka
\code{abc}, ustawia \code{failbit}; naprawienie \emph{stanu strumienia}
przez \code{clear()} nie usuwa automatycznie błędnego tekstu z bufora.
Połączenie \code{>>} z \code{getline} wymaga uwzględnienia pozostawionego
znaku końca linii.

\begin{lstlisting}[caption={Liczba, a potem pelna linia (kompletny program).}]
#include <iostream>
#include <limits>
#include <sstream>
#include <string>

int main() {
    std::istringstream in("7\nhello world\n");
    int id{};
    in >> id;
    in.ignore(std::numeric_limits<std::streamsize>::max(), '\n');
    std::string line;
    std::getline(in, line);
    std::cout << id << ":" << line << '\n';
}
\end{lstlisting}

\subsection{Flagi formatowania i wyjątki strumieni}

Manipulatory \code{std::fixed}, \code{std::scientific},
\code{std::setprecision}, \code{std::setw}, \code{std::setfill} i
\code{std::boolalpha} zmieniają sposób prezentacji. Część ustawień
pozostaje w strumieniu na kolejne operacje; \code{setw} dotyczy tylko
następnego pola. Ustawienie \code{exceptions(mask)} sprawia, że zmiana
odpowiednich bitów stanu może zgłosić \code{std::ios_base::failure}.

\begin{lstlisting}[caption={Lokalne formatowanie bez zmieniania std::cout.}]
#include <iomanip>
#include <iostream>
#include <sstream>

int main() {
    std::ostringstream out;
    out << std::fixed << std::setprecision(2) << 3.14159;
    std::cout << out.str() << '\n'; // 3.14
}
\end{lstlisting}

\nota{\code{std::endl} nie jest tylko nową linią: dodatkowo opróżnia
bufor. W pętli wypisującej wiele rekordów używamy zwykle \code{'\\n'},
a \code{flush} wywołujemy, gdy rzeczywiście wymagamy wypchnięcia bufora.
\code{flush} nie oznacza trwałego zapisania danych na fizycznym nośniku.}

\subsection{Formatowanie a ustawienia regionalne}

Operatory strumieni zależą od przypisanej im lokalizacji
(\code{std::locale}); separator dziesiętny i grupowanie cyfr mogą się
różnić. Do stabilnego formatu wymiany danych należy jawnie wybrać
konwencję, najlepiej niezależną od regionalnych ustawień użytkownika.
\code{std::format} (C++20) zapewnia formatowanie według tekstu formatu,
lecz dostępność nagłówka \code{<format>} zależy od biblioteki.

\begin{lstlisting}[caption={Formatowanie C++20 -- wymaga obslugi <format>.}]
#include <format>
#include <string>

std::string describe(int n, double x) {
    return std::format("n={}, x={:.3f}", n, x);
}
\end{lstlisting}

\section{Pliki tekstowe i binarne}
\label{sec:cpp19-files}

\code{std::ifstream} czyta plik, \code{std::ofstream} zapisuje,
a \code{std::fstream} obsługuje oba kierunki. Otwarcie odbywa się
w konstruktorze lub przez \code{open}; destruktor zamyka uchwyt
zgodnie z RAII. Tryby obejmują \code{std::ios::in}, \code{out},
\code{app}, \code{trunc}, \code{binary} i \code{ate}.

\begin{lstlisting}[caption={Zapis i odczyt pliku tekstowego (kompletny program).}]
#include <fstream>
#include <iostream>
#include <stdexcept>
#include <string>

int main() {
    const std::string path = "cpp19_example.txt";
    {
        std::ofstream out(path, std::ios::trunc);
        if (!out) throw std::runtime_error("open for write");
        out << "alpha\nbeta\n";
        out.close();
        if (!out) throw std::runtime_error("write/close failed");
    }
    std::ifstream in(path);
    if (!in) throw std::runtime_error("open for read");
    std::string line;
    while (std::getline(in, line))
        std::cout << line << '\n';
    if (in.bad()) throw std::runtime_error("read failed");
}
\end{lstlisting}

Sprawdzanie sukcesu otwarcia nie wystarcza: zapis może zakończyć się
błędem także po częściowym wykonaniu. Jawne \code{close()} umożliwia
sprawdzenie stanu przed destrukcją. Sam fakt \code{close()} nie daje
gwarancji trwałości danych na dysku (\emph{durability}).

\subsection{Binarna reprezentacja obiektu nie jest formatem pliku}

\code{std::ios::binary} usuwa pewne transformacje trybu tekstowego;
nie wprowadza przenośnego sposobu kodowania struktur. Zapis
\code{reinterpret_cast<const char*>(&obj)} pomija kwestie
\emph{endianness}, wyrównania, paddingu, rozmiaru typów i wersji
reprezentacji. W plikach wymiany zalecamy własny jawny format albo
przetestowaną bibliotekę serializacji.

\begin{lstlisting}[caption={Jawny zapis liczby 32-bitowej w kolejnosci little-endian (C++20).}]
#include <array>
#include <cstdint>
#include <fstream>
#include <stdexcept>

void write_u32_le(std::ostream& out, std::uint32_t v) {
    std::array<char, 4> bytes{};
    for (unsigned i = 0; i < 4; ++i)
        bytes[i] = static_cast<char>((v >> (8 * i)) & 0xffu);
    out.write(bytes.data(), bytes.size());
    if (!out) throw std::runtime_error("write failed");
}
\end{lstlisting}

Odbiorca powinien sprawdzać długość danych, wersję protokołu,
dozwolone zakresy i limity alokacji, zanim zacznie tworzyć obiekty.
\code{std::bit_cast} ułatwia analizę reprezentacji typów o odpowiednich
własnościach, ale nie rozwiązuje samoistnie problemu przenośnej serializacji.

\section{System plików: \texttt{std::filesystem}}
\label{sec:cpp19-filesystem}

Nagłówek \code{<filesystem>} (C++17) udostępnia typ \code{path}
oraz operacje na ścieżkach i katalogach. \code{path} reprezentuje
ścieżkę w notacji systemu, nie otwiera jeszcze pliku. Operator
\code{/} poprawnie dołącza elementy ścieżki.

\begin{lstlisting}[caption={Sciezki i katalogi (kompletny program).}]
#include <filesystem>
#include <iostream>

namespace fs = std::filesystem;

int main() {
    const fs::path base = fs::temp_directory_path();
    const fs::path path = base / "results" / "data.csv";
    std::cout << path.filename().string() << '\n';
    std::cout << path.extension().string() << '\n';
    std::cout << path.parent_path().filename().string() << '\n';
}
\end{lstlisting}

\subsection{Wyjątki i \texttt{std::error\_code}}

Wiele funkcji ma przeciążenia zgłaszające
\code{std::filesystem::filesystem_error} oraz warianty przyjmujące
\code{std::error_code&}. Drugie przydają się w kodzie, w którym brak
pliku jest normalnym wynikiem, a nie wyjątkiem.

\begin{lstlisting}[caption={Kontrolowany odczyt informacji o pliku.}]
#include <filesystem>
#include <iostream>
#include <system_error>

int main() {
    std::error_code ec;
    const bool exists = std::filesystem::exists("config.txt", ec);
    if (ec) {
        std::cerr << "filesystem error: " << ec.message() << '\n';
        return 1;
    }
    std::cout << std::boolalpha << exists << '\n';
}
\end{lstlisting}

\subsection{Iterowanie i symbole łączy}

\code{directory_iterator} przechodzi po jednym poziomie, natomiast
\code{recursive_directory_iterator} schodzi do podkatalogów.
\code{exists}, \code{is_regular_file}, \code{is_symlink}, \code{status},
\code{symlink_status}, \code{canonical} i \code{weakly_canonical}
mają różne znaczenie. Domyślnie nie należy utożsamiać ścieżki
tekstowej z tożsamością rzeczywistego pliku: łącza symboliczne,
uprawnienia i zmiany katalogów mogą zmienić wynik.

\begin{lstlisting}[caption={Wypisanie zwyklych plikow z katalogu (C++17).}]
#include <filesystem>
#include <iostream>
#include <system_error>

int main() {
    namespace fs = std::filesystem;
    std::error_code ec;
    for (fs::directory_iterator it(fs::current_path(), ec), end;
         !ec && it != end; it.increment(ec)) {
        std::error_code item_ec;
        if (it->is_regular_file(item_ec) && !item_ec)
            std::cout << it->path().filename().string() << '\n';
    }
    if (ec) {
        std::cerr << ec.message() << '\n';
        return 1;
    }
}
\end{lstlisting}

\nota{Wzorzec \emph{sprawdź-then-użyj} nie zapobiega wyścigom
\emph{TOCTOU}: między \code{exists(path)} i \code{open(path)} ktoś może
zmienić obiekt wskazywany przez ścieżkę. Przy wymaganiach
bezpieczeństwa potrzebne są mechanizmy systemowe pracujące na
uchwytach do otwartych plików i właściwe polityki dostępu.}

\subsection{Zapis przez plik tymczasowy i podmianę}

Częsty wzorzec to zapis nowej zawartości do osobnego pliku w tym
samym katalogu, zamknięcie go i zastąpienie pliku docelowego przez
\code{rename}. Dokładne gwarancje atomowości podmiany, obsługę
istniejącego celu oraz trwałości trzeba ustalić dla danego systemu
plików i systemu operacyjnego. Sama biblioteka C++ nie tworzy
transakcyjnego systemu plików. Dla plików poufnych konieczne jest
również bezpieczne tworzenie unikalnego pliku tymczasowego.

\section{Własność tekstu i \texttt{std::string\_view}}
\label{sec:cpp19-text}

\code{std::string} zarządza własną pamięcią. \code{std::string_view}
przechowuje wskaźnik i długość fragmentu znaków, nie przedłużając
czasu życia danych. Widok bywa idealnym parametrem funkcji, jeżeli
funkcja nie zachowuje go na później.

\begin{lstlisting}[caption={Widok do bufora -- bez kopiowania (kompletny program).}]
#include <iostream>
#include <string>
#include <string_view>

void print_prefix(std::string_view text, std::size_t n) {
    std::cout << text.substr(0, n) << '\n';
}

int main() {
    std::string text = "programming";
    print_prefix(text, 7); // program
    print_prefix("hello", 3); // hel
}
\end{lstlisting}

\begin{lstlisting}[caption={Blad czasu zycia -- tylko ilustracja, nie uruchamiac.}]
std::string_view bad() {
    std::string local = "temporary";
    return local; // dangling view po zakonczeniu funkcji
}
\end{lstlisting}

Zmiana rozmiaru \code{std::string} może unieważnić wcześniejsze widoki.
Dodatkowo \code{string_view::substr} nie kopiuje znaków, natomiast
\code{string::substr} zwraca nowy \code{string}. W API akceptującym
tekst wejściowy dobrze jest dokumentować, kto posiada pamięć i jak
długo wywoływana funkcja może jej używać.

\subsection{UTF-8, Unicode i kodowanie}

\code{std::string} jest ciągiem jednostek typu \code{char}; sam z siebie
nie jest zbiorem znaków Unicode. W UTF-8 jeden punkt kodowy może
zajmować 1--4 bajty, a widoczny dla użytkownika znak (\emph{grapheme
cluster}) może składać się z kilku punktów kodowych. Dlatego
\code{text.size()} liczy bajty, a nie litery. \code{char8_t} (C++20)
reprezentuje jednostkę kodową UTF-8, lecz \code{u8"..."} ma typ
\code{const char8_t[N]}, nie \code{const char[N]}.

\begin{lstlisting}[caption={Rozmiar w bajtach i liczba punktow kodowych UTF-8.}]
#include <cstddef>
#include <iostream>
#include <string_view>

std::size_t codepoints_utf8(std::string_view bytes) {
    std::size_t n = 0;
    for (unsigned char c : bytes)
        if ((c & 0xc0u) != 0x80u) ++n;
    return n;
}

int main() {
    std::string_view text = "A\xC4\x85"; // A oraz 'a' z ogonkiem
    std::cout << text.size() << ' ' << codepoints_utf8(text) << '\n';
}
\end{lstlisting}

Funkcja z przykładu \emph{zakłada poprawne UTF-8} i liczy jedynie
jednostki niebędące bajtami kontynuacji: nie jest walidatorem kodowania.
Dla danych niezaufanych wymagane jest pełne sprawdzenie składni UTF-8,
a przy operacjach językowych --- biblioteka z obsługą normalizacji,
segmentacji i porównywania Unicode. \code{std::codecvt} nie jest zalecanym
nowym API (część jego komponentów została wycofana lub uznana za
przestarzałą).

\section{Parsowanie liczb przez \texttt{from\_chars}}
\label{sec:cpp19-charconv}

Funkcje \code{std::from_chars} i \code{std::to_chars} pracują na
przedziałach bufora, nie wymagają tworzenia strumienia i nie zależą
od lokalizacji \code{std::locale}. Wynik \code{from_chars} zawiera
wskaźnik na pierwszy nieprzetworzony znak oraz \code{std::errc}.
W szczególności udany odczyt prefiksu \code{123} z \code{123abc}
nie oznacza, że całe wejście jest poprawną liczbą.

\begin{lstlisting}[caption={Poprawny parser pelnej liczby calkowitej (C++17).}]
#include <charconv>
#include <optional>
#include <string_view>
#include <system_error>

std::optional<int> parse_int(std::string_view text) {
    int value{};
    auto result = std::from_chars(text.data(),
                                  text.data() + text.size(), value);
    if (result.ec != std::errc{} ||
        result.ptr != text.data() + text.size())
        return std::nullopt;
    return value;
}
\end{lstlisting}

Inaczej niż odczyt formatowany, \code{from_chars} nie pomija
początkowych białych znaków. Dla typów zmiennoprzecinkowych należy
uwzględnić dostępność odpowiedniej implementacji standardowej biblioteki
oraz wymagania dotyczące wartości nieskończonych, NaN i zaokrągleń.

\section{Parser pliku CSV: przykład kontraktu formatu}

CSV nie jest po prostu \code{split(',')}: poprawne pola mogą zawierać
przecinek i znaki cudzysłowu, a cytowane pola także nowe linie.
Najpierw należy określić dialekt, kodowanie, maksymalny rozmiar pola
i zachowanie wobec niepoprawnych danych. Poniższy prosty parser
obsługuje jeden rekord przechowywany już jako \code{string_view},
cudzysłowy i podwojone cudzysłowy, ale \emph{nie obsługuje wieloliniowych
rekordów} (ich odczyt wymaga osobnej warstwy).

\begin{lstlisting}[caption={Parser pojedynczego rekordu CSV (C++20).}]
#include <stdexcept>
#include <string>
#include <string_view>
#include <vector>
#include <utility>

std::vector<std::string> parse_csv_record(std::string_view s) {
    std::vector<std::string> fields;
    std::string field;
    bool quoted = false;
    bool just_closed = false;
    bool at_field_start = true;
    for (std::size_t i = 0; i < s.size(); ++i) {
        const char c = s[i];
        if (quoted) {
            if (c == '"') {
                if (i + 1 < s.size() && s[i + 1] == '"') {
                    field.push_back('"'); ++i;
                } else {
                    quoted = false;
                    just_closed = true;
                }
            } else {
                field.push_back(c);
            }
        } else if (c == ',' ) {
            fields.push_back(std::move(field));
            field.clear();
            at_field_start = true;
            just_closed = false;
        } else if (c == '"' && at_field_start) {
            quoted = true;
            at_field_start = false;
        } else {
            if (just_closed || c == '"')
                throw std::invalid_argument("invalid CSV quote");
            field.push_back(c);
            at_field_start = false;
        }
    }
    if (quoted) throw std::invalid_argument("unclosed quote");
    fields.push_back(std::move(field));
    return fields;
}
\end{lstlisting}

Istotne jest ograniczenie pamięci: parser przyjmujący dowolnie długi
rekord może przy dużym wejściu spowodować wyczerpanie zasobów.
Zachowanie \code{std::regex} również nie zastępuje pełnego parsera CSV;
wyrażenia regularne nie są ponadto narzędziem do ogólnej segmentacji Unicode.

\section{Wydajność i projektowanie interfejsów I/O}

Koszt operacji na nośniku i koszty systemowe często dominują nad
samymi obliczeniami. Nie należy jednak bez pomiaru utożsamiać
\code{mmap} z szybkością, a \code{fstream} z wolnym rozwiązaniem.
Przed optymalizacją warto mierzyć liczbę wywołań systemowych,
wielkość bufora, przepustowość, koszt parsowania i alokacji.

\begin{lstlisting}[caption={Wczytanie calego pliku binarnego z kontrola bledu (C++20).}]
#include <fstream>
#include <iterator>
#include <stdexcept>
#include <string>

std::string read_all(const std::string& path) {
    std::ifstream input(path, std::ios::binary);
    if (!input) throw std::runtime_error("open failed");
    std::string result((std::istreambuf_iterator<char>(input)),
                        std::istreambuf_iterator<char>{});
    if (input.bad()) throw std::runtime_error("read failed");
    return result;
}
\end{lstlisting}

Funkcja jest prosta i przydatna dla plików o kontrolowanym rozmiarze.
Jeśli plik może być wielogigabajtowy lub pochodzi od niezaufanego
nadawcy, lepiej przetwarzać go strumieniowo, z limitami i obsługą
częściowych rekordów. Dla operacji formatowania warto rozważyć
\code{std::format} (C++20) i \code{std::print} (C++23), ale przy
\code{std::print} należy sprawdzić wspierającą implementację.

\section{Konwersja liczb w obu kierunkach: \texttt{to\_chars}}

Dla protokołów tekstowych często potrzeba zamieniać liczby na tekst
w już przydzielonym buforze. \code{std::to_chars} nie tworzy
\code{std::string}, nie używa lokalizacji i informuje, gdzie skończył
pisać. Użytkownik odpowiada za wielkość bufora.

\begin{lstlisting}[caption={Formatowanie liczby calkowitej bez alokacji (C++17).}]
#include <array>
#include <charconv>
#include <stdexcept>
#include <string>
#include <system_error>

std::string format_integer(int value) {
    std::array<char, 32> buffer{};
    auto [ptr, ec] = std::to_chars(buffer.data(),
                                   buffer.data() + buffer.size(),
                                   value);
    if (ec != std::errc{})
        throw std::runtime_error("buffer too small");
    return std::string(buffer.data(), ptr);
}
\end{lstlisting}

Warto odróżnić brak alokacji w samej konwersji od alokacji przy
zbudowaniu wynikowego \code{std::string}. Można zamiast tego pisać
do bufora zarządzanego przez wywołującego; wtedy jego czas życia i
pojemność stają się elementem kontraktu.

\section{Małe protokoły binarne: odczyt kontrolowany}

Zapis wersjonowany nie musi oznaczać złożonej biblioteki. Jeżeli
format przewiduje czterobajtową liczbę little-endian, parsowanie
powinno sprawdzić, czy rzeczywiście otrzymaliśmy dokładnie tyle
bajtów. Należy jasno ustalić, czy nadmiarowe dane są dozwolone.

\begin{lstlisting}[caption={Odczyt 32-bitowej liczby little-endian (C++20).}]
#include <array>
#include <cstdint>
#include <istream>
#include <stdexcept>

std::uint32_t read_u32_le(std::istream& input) {
    std::array<unsigned char, 4> b{};
    for (auto& byte : b) {
        char c{};
        if (!input.get(c))
            throw std::runtime_error("truncated input");
        byte = static_cast<unsigned char>(c);
    }
    std::uint32_t result = 0;
    for (unsigned i = 0; i < 4; ++i)
        result |= std::uint32_t{b[i]} << (8 * i);
    return result;
}
\end{lstlisting}

W plikach o zmiennym formacie warto stosować pola: sygnaturę
(\emph{magic}), wersję, długość rekordu i jego zawartość. Jeśli
długość pochodzi z pliku, konieczny jest limit, kontrola przepełnienia
przy obliczaniu rozmiaru alokacji i obsługa przerwania odczytu.

\section{Wyjście wielowątkowe i \texttt{std::osyncstream}}

Wielu producentów logów może wywoływać \code{std::cout} jednocześnie.
Nawet jeśli biblioteka zapobiega wewnętrznemu wyścigowi danych,
poszczególne znaki wielu wiadomości mogą się przeplatać.
C++20 dostarcza \code{std::osyncstream}, który grupuje wypisywanie
i emituje zebrany fragment przy synchronizacji bufora lub destrukcji.

\begin{lstlisting}[caption={Niezachodzace na siebie linie logu (C++20).}]
#include <iostream>
#include <syncstream>
#include <thread>
#include <vector>

int main() {
    std::vector<std::jthread> workers;
    for (int id = 0; id < 4; ++id) {
        workers.emplace_back([id] {
            std::osyncstream(std::cout) << "worker " << id << '\n';
        });
    }
}
\end{lstlisting}

Kolejność wyświetlenia całych komunikatów nadal zależy od wykonania
wątków. \code{osyncstream} nie służy do synchronizacji obliczeń ani
do wymuszania kolejności zdarzeń w algorytmie.

\section{Szczegóły przenośności: ścieżki i kodowanie nazw plików}

Reprezentacja \code{filesystem::path} jest zależna od platformy.
Na Windows ścieżka może mieć natywną reprezentację opartą na szerokich
znakach, podczas gdy systemy POSIX posługują się sekwencjami bajtów.
Nie należy utożsamiać \code{path.string()} z konwersją Unicode
z identyczną semantyką na wszystkich platformach.

Również separator katalogów, wielkość liter, dozwolone nazwy,
obsługa ścieżek sieciowych i normalizacja są platformowe.
\code{lexically_normal()} upraszcza składnię ścieżki, lecz
nie wykonuje operacji I/O i nie rozstrzyga, dokąd prowadzą łącza
symboliczne. \code{canonical()} odnosi się do faktycznego systemu
plików, a zatem może zgłosić błąd.

\section{Projektowanie parserów i diagnostyka błędów}

Dobre API parsera zwraca nie tylko rezultat, ale -- gdy jest to
potrzebne -- także pozycję błędu. Format powinien określać,
czy dopuszcza nadmiarowe białe znaki, puste pola, separator na końcu
i dodatkowe dane. W warstwie parsera warto rozróżnić błędny zapis
liczby od jej przepełnienia oraz zbyt długi rekord od przedwczesnego
końca pliku. Takie rozróżnienie ułatwia generowanie precyzyjnych
komunikatów i testów regresyjnych.

W systemach obsługujących dane niezaufane kod przetwarzający tekst
powinien obowiązkowo mieć testy brzegowe: pusty plik, zerowy bajt,
wielobajtowy znak na granicy bufora, cytowanie CSV, rekord przekraczający
limit oraz zmianę pliku podczas odczytu. Dla parserów złożonych
pomocne są fuzzery i sanitizatory opisane w rozdziale 18.


\section{Podsumowanie: kiedy używać którego narzędzia?}
\begin{center}\small
\begin{tabularx}{\textwidth}{@{}p{0.32\textwidth}X@{}}
\toprule
Potrzeba & Narzędzie i uwagi \\
\midrule
Niewielkie pliki konfiguracyjne & \code{ifstream}, \code{getline}, oddzielny parser, walidacja. \\
Formatowanie komunikatów & \code{ostringstream} albo \code{format} przy dostępnej implementacji. \\
Duże przetwarzanie liczb tekstowych & \code{from_chars}, buforowanie, pomiar narzutów. \\
Ścieżki i metadane & \code{filesystem}, obsługa \code{error_code} i wyjątków. \\
API odczytujące tekst & \code{string_view}, pod warunkiem poprawnego czasu życia. \\
Trwałe dane binarne & Wersjonowany format; nie bezpośredni zapis struktur. \\
Unicode i operacje językowe & Walidacja UTF-8 i wyspecjalizowana biblioteka. \\
\bottomrule
\end{tabularx}
\end{center}

\nota{Trzy częste błędy to: \code{while(!eof())}, ignorowanie błędu
zapisu oraz zachowywanie \code{string_view} do zniszczonego bufora.
Z kolei w kodzie wrażliwym na bezpieczeństwo samo sprawdzenie, czy
ścieżka istnieje, nie zapewnia bezpiecznego otwarcia.}

\section*{Literatura i dokumentacja}
\addcontentsline{toc}{section}{Literatura i dokumentacja}
\begingroup\sloppy
\begin{itemize}
\item \url{https://en.cppreference.com/w/cpp/io} --- strumienie i ich stany.
\item \url{https://en.cppreference.com/w/cpp/filesystem} --- \code{filesystem}.
\item \url{https://en.cppreference.com/w/cpp/utility/from_chars} --- parsowanie liczb.
\item \url{https://en.cppreference.com/w/cpp/string/basic_string_view} --- widoki tekstowe.
\item \url{https://en.cppreference.com/w/cpp/utility/format/format} --- formatowanie.
\item The Unicode Consortium, \emph{The Unicode Standard}, \url{https://unicode.org/standard/standard.html}.
\end{itemize}
\endgroup

% ROZDZIAL 13 — Kontenery STL i algorytmy w C++20.
% Fragment do zbiorczego skryptu (bez \documentclass i \begin{document}).
% Wymagania: listings, amsmath, amssymb, booktabs, tabularx, array, hyperref.
% Kompilator kodu: C++20; nowe elementy standardu sa oznaczane w tekscie.
% Dolaczenie: \input{rozdzial_13_kontenery_stl_algorytmy_cpp.tex}

\chapter{Kontenery STL i algorytmy}
\label{chap:cpp-stl-containers}

\section{Cel i zakres rozdziału}

Kontener przechowuje grupę elementów i pozwala je odnajdywać lub
zmieniać. Najprostsza mapa wyboru wygląda tak: \texttt{vector} jest
listą z dostępem przez indeks, \texttt{map} słownikiem klucz--wartość,
a \texttt{set} zbiorem wartości bez powtórzeń. Jeśli kolejność kluczy
nie ma znaczenia, istnieją także haszujące warianty
\texttt{unordered\_map} i \texttt{unordered\_set}.

Biblioteka standardowa C++ udostępnia gotowe struktury danych, których wybór ma
istotne konsekwencje dla złożoności operacji, lokalności pamięci, bezpieczeństwa
wyjątkowego i poprawności iteratorów. Samo rozróżnienie między
\texttt{std::vector}, \texttt{std::list} i \texttt{std::map} nie wystarcza:
trzeba rozumieć również \emph{własność elementów}, \emph{unieważnianie
iteratorów}, wymagania dotyczące porządku i haszowania oraz relację między
kontenerami a algorytmami. W dalszej części domyślnie używamy C++20.

\paragraph{Główna zasada projektowa.}
Najpierw ustalamy rzeczywiste operacje: czy potrzebujemy dostępu przez
indeks, wyszukiwania po kluczu, zachowania kolejności, stabilnych adresów,
czy częstych wstawek i usunięć. Dopiero potem wybieramy kontener. Prosty
\texttt{std::vector} pozostaje często najlepszym rozwiązaniem, nawet gdy
jego teoretyczna złożoność niektórych operacji jest gorsza, ponieważ dane
leżą w sąsiadujących komórkach pamięci.

\section{Klasyfikacja kontenerów i podstawowe pojęcia}
\label{sec:cpp13-taxonomy}

Biblioteka standardowa rozróżnia m.in. kontenery sekwencyjne
(\texttt{array}, \texttt{vector}, \texttt{deque}, \texttt{list},
\texttt{forward\_list}), uporządkowane asocjacyjne (\texttt{map},
\texttt{set} i ich warianty \texttt{multi}) oraz nieuporządkowane
asocjacyjne (\texttt{unordered\_map}, \texttt{unordered\_set} i warianty
\texttt{multi}). \texttt{stack}, \texttt{queue} i \texttt{priority\_queue}
są \emph{adapterami kontenerów}: oferują inny interfejs dostępu do
wewnętrznej struktury danych. \texttt{std::string} jest wyspecjalizowaną
strukturą sekwencyjną, ale semantykę tekstu omawiamy oddzielnie.

\begin{center}
\small
\begin{tabularx}{\textwidth}{@{}lXXX@{}}
\toprule
\textbf{Kontener} & \textbf{Dostęp} & \textbf{Wstawienie / usunięcie} & \textbf{Typowe zastosowanie} \\
\midrule
\texttt{vector} & Indeks $O(1)$ & Koniec: zamortyzowane $O(1)$; środek: $O(n)$ & Dane zwarte, iteracja, obliczenia \\
\texttt{deque} & Indeks $O(1)$ & Oba końce: $O(1)$; środek: $O(n)$ & Kolejka dwustronna \\
\texttt{list} & Iteracja $O(n)$ & $O(1)$ przy znanym iteratorze & Stabilne węzły, \texttt{splice} \\
\texttt{forward\_list} & Iteracja $O(n)$ & $O(1)$ za poprzednikiem & Lista jednokierunkowa \\
\texttt{map} / \texttt{set} & Szukanie $O(\log n)$ & $O(\log n)$ & Klucze w ustalonej kolejności \\
\texttt{unordered\_map} & Średnio $O(1)$ & Średnio $O(1)$ & Słownik, szybkie wyszukiwanie \\
\bottomrule
\end{tabularx}
\end{center}

W tabeli $n$ oznacza liczbę elementów, a koszty zależą od wymagań
poszczególnych operacji. \texttt{unordered\_map} może w najgorszym
przypadku wykonywać operacje w czasie $O(n)$ (np. przy wielu kolizjach).
Sformułowanie ,,\emph{zamortyzowane} $O(1)$'' oznacza, że pojedyncza
operacja może być kosztowna, ale średni koszt w długiej sekwencji operacji
jest ograniczony stałą.

\paragraph{Wartości, wskaźniki, referencje.}
Kontenery zwykle przechowują \emph{własne obiekty elementów}.
\texttt{std::vector<T>} nie jest kontenerem referencji do obiektów
\texttt{T}; może za to przechowywać np. \texttt{std::reference\_wrapper<T>},
\texttt{T*} lub \texttt{std::unique\_ptr<T>}. Wybór wpływa na własność
obiektów, kopiowanie, czas życia i poprawność dostępu.

\section{\texttt{std::vector}: domyślny kontener sekwencyjny}

\begin{center}\small
\begin{tabularx}{\textwidth}{@{}p{0.34\textwidth}X@{}}
\toprule
Operacja & Za co odpowiada \\
\midrule
\code{size()}, \code{empty()} & Podają liczbę elementów i informację, czy kontener jest pusty. \\
\code{capacity()} & Podaje liczbę elementów, które zmieszczą się bez ponownej alokacji. \\
\code{reserve(n)} & Rezerwuje miejsce na co najmniej \code{n} elementów; nie zmienia \code{size()}. \\
\code{resize(n)} & Zmienia liczbę elementów: tworzy nowe albo niszczy nadmiarowe. \\
\code{push_back(x)} & Dodaje istniejącą wartość na koniec przez kopiowanie lub przeniesienie. \\
\code{emplace_back(args...)} & Konstruuje nowy element na końcu z podanych argumentów. \\
\code{erase(pos)} & Usuwa element; kolejne elementy przesuwają się, a iteratory od miejsca usunięcia tracą ważność. \\
\bottomrule
\end{tabularx}
\end{center}
\label{sec:cpp13-vector}

\texttt{std::vector<T>} przechowuje elementy \emph{ciągle} w pamięci,
co pozwala przekazywać je do interfejsów oczekujących tablicy (np. przez
\texttt{data()} lub \texttt{std::span}). Obiekt wektora zarządza własną
pamięcią i niszczy elementy po zakończeniu swojego czasu życia.

\begin{lstlisting}[caption={Podstawowe operacje na vector (C++11); kompletny program.}]
#include <cassert>
#include <vector>

int main() {
    std::vector<int> v{4, 8, 15};
    v.push_back(16);           // kopiowanie lub przenoszenie argumentu
    v.emplace_back(23);        // konstruowanie nowego elementu

    assert(v.size() == 5);
    assert(v[0] == 4);
    assert(v.at(4) == 23);
    assert(v.front() == 4 && v.back() == 23);

    v.pop_back();
    assert(v.size() == 4);
    v.clear();
    assert(v.empty());
}
\end{lstlisting}

\texttt{operator[]} \emph{nie sprawdza} granic: użycie indeksu spoza zakresu
prowadzi do niezdefiniowanego zachowania. Metoda \texttt{at()} sprawdza
zakres i zgłasza \texttt{std::out\_of\_range}. Metody \texttt{front()},
\texttt{back()} i \texttt{pop\_back()} wymagają niepustego kontenera.

\subsection{\texttt{size}, \texttt{capacity}, \texttt{reserve} i \texttt{resize}}

\texttt{size()} oznacza liczbę \emph{istniejących} elementów, a
\texttt{capacity()} liczbę elementów, dla których przydzielono miejsce
bez konieczności realokacji. \texttt{reserve(k)} zmienia rezerwację
pamięci, nie tworzy elementów; \texttt{resize(k)} zmienia ich liczbę.

\begin{lstlisting}[caption={Rezerwacja pamięci nie zwiększa rozmiaru (C++11); kompletny program.}]
#include <cassert>
#include <vector>

int main() {
    std::vector<int> data;
    data.reserve(100);
    assert(data.size() == 0);
    assert(data.capacity() >= 100);

    data.push_back(7);
    data.resize(4);          // dopisuje trzy elementy o wartosci 0
    assert(data.size() == 4);
    assert(data[0] == 7 && data[3] == 0);
}
\end{lstlisting}

Nie należy wywoływać \texttt{reserve(size()+1)} przed każdym
\texttt{push\_back()}: wymuszenie niemal liniowego wzrostu pojemności
może zniweczyć politykę geometrycznego powiększania bufora i pogorszyć
wydajność. \texttt{shrink\_to\_fit()} jest tylko żądaniem, a nie gwarancją
zmniejszenia pojemności.

\subsection{Realokacja i unieważnianie iteratorów}

Po realokacji \texttt{vector} zmienia adresy wszystkich elementów:
tracą ważność iteratory, referencje i wskaźniki do elementów, również
poprzednia wartość \texttt{data()}. Bez realokacji \texttt{push\_back()}
unieważnia iterator \texttt{end()}, lecz nie unieważnia iteratorów i
referencji do wcześniej istniejących elementów. Wstawienie w środku
bez realokacji unieważnia iteratory i referencje do pozycji wstawienia
oraz wszystkich późniejszych pozycji (włącznie z dawnym \texttt{end()}).

\begin{lstlisting}[caption={Bezpieczne ponowne pobranie iteratora po wstawieniu; kompletny program.}]
#include <cassert>
#include <vector>

int main() {
    std::vector<int> v{1, 3, 4};
    auto pos = v.begin() + 1;
    pos = v.insert(pos, 2);   // przyjmuje zwrocony iterator
    assert(*pos == 2);
    assert((v == std::vector<int>{1, 2, 3, 4}));
}
\end{lstlisting}

\paragraph{Przenoszenie elementów i bezpieczeństwo wyjątkowe.}
Podczas realokacji \texttt{vector} może przenosić lub kopiować elementy.
Dla pewnych typów i operacji warunki silnej gwarancji wyjątkowej zależą
od tego, czy konstruktor przenoszący jest \texttt{noexcept} oraz czy
dostępne jest kopiowanie. Zagadnienie łączy się bezpośrednio z semantyką
przenoszenia omówioną w rozdziale 6.

\subsection{\texttt{push\_back} kontra \texttt{emplace\_back}}

\texttt{emplace\_back(args...)} przekazuje argumenty bezpośrednio do
konstruktora nowego elementu. Nie oznacza to jednak automatycznie
,,zawsze szybciej'': jeśli mamy już obiekt, \texttt{push\_back(std::move(x))}
bywa równie naturalne; przy realokacji nadal konieczne może być
przeniesienie innych elementów. Obie funkcje mogą zgłaszać wyjątki.

\begin{lstlisting}[caption={Konstruowanie elementu z argumentow (C++17); kompletny program.}]
#include <cassert>
#include <string>
#include <utility>
#include <vector>

struct Person {
    std::string name;
    int age;
    Person(std::string n, int a) : name(std::move(n)), age(a) {}
};

int main() {
    std::vector<Person> people;
    people.emplace_back("Ada", 31);
    people.push_back(Person{"Alan", 42});
    assert(people[1].age == 42);
}
\end{lstlisting}

\section{\texttt{std::array}, \texttt{std::deque} i listy}
\label{sec:cpp13-sequential}

\subsection{\texttt{std::array<T,N>}: rozmiar znany w kompilacji}

\texttt{std::array} przechowuje dokładnie $N$ elementów. Nie dokonuje
osobnej dynamicznej alokacji dla swojej zawartości; pamięć elementów
stanowi część samego obiektu \texttt{array}. Można używać iteratorów,
\texttt{size()}, \texttt{data()} i algorytmów, analogicznie do wektora.

\begin{lstlisting}[caption={std::array i algorytmy standardowe (C++17); kompletny program.}]
#include <algorithm>
#include <array>
#include <cassert>
#include <tuple>

int main() {
    std::array<int, 4> x{5, 1, 4, 2};
    std::sort(x.begin(), x.end());
    assert((x == std::array<int, 4>{1, 2, 4, 5}));
    static_assert(std::tuple_size_v<decltype(x)> == 4);
}
\end{lstlisting}

\subsection{\texttt{std::deque}: dwa tanie końce}

\texttt{deque} udostępnia dostęp indeksowany w czasie $O(1)$, lecz
zwykle nie jest jednym ciągłym blokiem pamięci. Wstawianie i usuwanie
na początku i końcu ma stałą złożoność; w przeciwieństwie do
\texttt{vector} ma operacje \texttt{push\_front()} i \texttt{pop\_front()}.
\texttt{deque} nie udostępnia \texttt{data()} dla wszystkich elementów.

\begin{lstlisting}[caption={Kolejka dwustronna (C++11); kompletny program.}]
#include <cassert>
#include <deque>

int main() {
    std::deque<int> q{2, 3};
    q.push_front(1);
    q.push_back(4);
    assert(q[0] == 1 && q[3] == 4);
    q.pop_front();
    assert(q.front() == 2);
}
\end{lstlisting}

Wstawienie na końcu \texttt{deque} może unieważnić \emph{iteratory},
ale zachowuje referencje do istniejących elementów. Szczegółowe reguły
różnią się dla wstawień i usunięć w środku oraz na końcach; nie należy
mechanicznie przenosić reguł \texttt{vector} na \texttt{deque}.

\subsection{\texttt{std::list} i \texttt{std::forward\_list}}

\texttt{list} jest listą dwukierunkową. Wstawianie i usuwanie przy
znanym iteratorze odbywa się w stałym czasie; iteratory i referencje
do pozostałych elementów pozostają ważne. Wyszukanie elementu przez
wartość wymaga jednak przejścia listy $O(n)$, podobnie jak uzyskanie
elementu o zadanym indeksie. W przeciwieństwie do \texttt{vector},
\texttt{list} nie oferuje operatora \texttt{[]} ani iteratorów dostępu
swobodnego. Do sortowania używamy \texttt{list::sort()}, a nie
\texttt{std::sort()}.

\begin{lstlisting}[caption={Przeniesienie wezla przez list::splice (C++11); kompletny program.}]
#include <cassert>
#include <iterator>
#include <list>

int main() {
    std::list<int> source{1, 2, 3};
    std::list<int> target{10};
    auto it = std::next(source.begin()); // wskazuje 2

    target.splice(target.end(), source, it);
    assert(*it == 2); // iterator pozostaje wazny
    assert(source.size() == 2 && target.size() == 2);
    target.sort();
    assert(target.front() == 2);
}
\end{lstlisting}

\texttt{forward\_list} przechowuje węzły listy jednokierunkowej;
wstawianie i usuwanie wykonuje \emph{za} wskazanym poprzednikiem
(\texttt{insert\_after}, \texttt{erase\_after}), a \texttt{before\_begin()}
udostępnia pozycję przed pierwszym elementem. Kontener nie posiada
metody \texttt{size()} (rozmiar można policzyć liniowo przez
\texttt{std::distance}).

\paragraph{Kiedy lista jest rzeczywiście potrzebna?}
Lista bywa przydatna, gdy musimy często \emph{przenosić węzły bez
kopiowania obiektów} i wymagamy stabilnych iteratorów. Dla częstej
iteracji i obliczeń numerycznych \texttt{vector} z reguły zapewnia
lepszą lokalność pamięci. Wybór należy jednak potwierdzić pomiarem.

\section{Kontenery uporządkowane: \texttt{set} i \texttt{map}}
\label{sec:cpp13-ordered}

\texttt{std::set<Key>} przechowuje unikalne klucze, a
\texttt{std::map<Key,T>} unikalne pary klucz--wartość. Elementy są
iterowane w kolejności zdefiniowanej przez komparator; wyszukiwanie,
wstawianie i usuwanie są logarytmiczne. Warianty
\texttt{multiset}/\texttt{multimap} dopuszczają powtarzające się klucze.

\subsection{\texttt{map::operator[]}, \texttt{at} i \texttt{find}}

Odwołanie \texttt{m[key]} do nieistniejącego klucza \emph{wstawia} nowy
element z wartością inicjalizowaną domyślnie; może zatem modyfikować
strukturę danych. \texttt{at()} zgłasza wyjątek, jeżeli klucz nie
istnieje. \texttt{find()} zwraca iterator, a \texttt{contains()} (C++20)
zwraca wartość logiczną bez wstawiania.

\begin{lstlisting}[caption={Liczenie czestosci slow z map (C++20); kompletny program.}]
#include <cassert>
#include <map>
#include <string>
#include <vector>

int main() {
    std::vector<std::string> words{"kot", "pies", "kot"};
    std::map<std::string, int> count;
    for (const auto& word : words) ++count[word];

    assert(count.at("kot") == 2);
    assert(count.contains("pies"));
    assert(!count.contains("lis"));
    assert(count.find("lis") == count.end());
}
\end{lstlisting}

\subsection{\texttt{insert}, \texttt{try\_emplace}, \texttt{insert\_or\_assign}}

\texttt{insert} podejmuje próbę dodania gotowej pary.
\texttt{try\_emplace(key,args...)} (C++17) konstruuje wartość
w miejscu z przekazanych argumentów tylko wtedy, gdy klucz nie występuje.
Ważne zastrzeżenie: \emph{wyrażenia będące argumentami wywołania}
są obliczane jeszcze przed wejściem do \texttt{try\_emplace}; nie
należy więc podawać niepotrzebnie kosztownego wyrażenia tworzącego
tymczasowy obiekt. \texttt{insert\_or\_assign} dodaje lub zastępuje wartość.

\begin{lstlisting}[caption={Rozne operacje modyfikacji map (C++17); kompletny program.}]
#include <cassert>
#include <map>
#include <memory>
#include <string>

struct Heavy {
    static inline int constructions = 0; // C++17
    int value;
    explicit Heavy(int x) : value(x) { ++constructions; }
};

int main() {
    std::map<std::string, Heavy> m;
    auto [it, added] = m.try_emplace("a", 3);
    assert(added && it->second.value == 3);

    auto [again, added_again] = m.try_emplace("a", 4);
    assert(!added_again && again->second.value == 3);
    assert(Heavy::constructions == 1);

    m.insert_or_assign("a", Heavy{9});
    assert(m.at("a").value == 9);

    std::map<int, std::unique_ptr<int>> owned;
    owned.try_emplace(1, std::make_unique<int>(10));
    assert(*owned.at(1) == 10);
}
\end{lstlisting}

\subsection{Własny komparator i ścisły słaby porządek}

Porównanie musi spełniać warunek \emph{strict weak ordering}; w
szczególności komparator powinien być irrefleksyjny
($\neg\operatorname{comp}(x,x)$) i przechodni. Jeżeli dwie wartości
$ a,b $ nie spełniają ani \texttt{comp(a,b)}, ani \texttt{comp(b,a)},
kontener uznaje je za równoważne jako klucze --- nawet gdy
\texttt{operator==} zwraca \texttt{false}. Wymagania tego porządku
są istotne również dla \texttt{sort} i \texttt{lower\_bound}.
Na przykład komparator ignorujący wielkość liter uzna \texttt{"ABC"}
i \texttt{"abc"} za ten sam klucz: druga próba wstawienia do
\texttt{map} nie utworzy nowego elementu.

\begin{lstlisting}[caption={Mapowanie bez rozrozniania wielkosci liter: prosty komparator; kompletny program.}]
#include <cassert>
#include <cctype>
#include <map>
#include <string>

struct CaseInsensitiveLess {
    bool operator()(const std::string& a,
                    const std::string& b) const {
        auto ai = a.begin();
        auto bi = b.begin();
        while (ai != a.end() && bi != b.end()) {
            auto x = static_cast<unsigned char>(*ai++);
            auto y = static_cast<unsigned char>(*bi++);
            int lx = std::tolower(x);
            int ly = std::tolower(y);
            if (lx != ly) return lx < ly;
        }
        return a.size() < b.size();
    }
};

int main() {
    std::map<std::string, int, CaseInsensitiveLess> m;
    m.emplace("ABC", 1);
    m.emplace("abc", 2); // rownowazny klucz
    assert(m.size() == 1 && m.begin()->second == 1);
}
\end{lstlisting}

Przykład porównuje bajty w podstawowym środowisku znakowym;
\texttt{std::tolower} z bieżącą lokalizacją C nie realizuje ogólnego
Unicode case folding. Do pełnego porównywania napisów w Unicode trzeba
zastosować osobne, świadomie wybrane reguły normalizacji i kolacji.

\subsection{Wyszukiwanie heterogeniczne}

Jeżeli komparator mapy jest \emph{transparentny} (np. \texttt{std::less<>}),
możemy szukać w mapie o kluczach \texttt{std::string} za pomocą
\texttt{std::string\_view}, bez tworzenia tymczasowego ciągu znaków.
Heterogeniczne przeciążenia \texttt{find} są dostępne
od C++14, a \texttt{contains} od C++20.

\begin{lstlisting}[caption={Heterogeniczne wyszukiwanie w map (C++20); kompletny program.}]
#include <cassert>
#include <functional>
#include <map>
#include <string>
#include <string_view>

int main() {
    std::map<std::string, int, std::less<>> m{{"alpha", 3}};
    std::string_view key = "alpha";
    auto it = m.find(key);
    assert(it != m.end() && it->second == 3);
    assert(m.contains(key));
}
\end{lstlisting}

\subsection{Zmiana klucza przez uchwyt węzła (C++17)}

Kluczy w \texttt{map} nie wolno modyfikować bezpośrednio przez iterator.
Uchwyt \texttt{node\_handle} uzyskany przez \texttt{extract()} pozwala
wyjąć węzeł, zmienić jego klucz i wstawić z powrotem. Trzeba pamiętać,
że ponowne wstawienie może się nie udać z powodu kolizji kluczy.

\begin{lstlisting}[caption={Zmiana klucza przez extract (C++17); kompletny program.}]
#include <cassert>
#include <map>
#include <string>
#include <utility>

int main() {
    std::map<std::string, int> m{{"old", 7}};
    auto node = m.extract("old");
    assert(!node.empty());
    node.key() = "new";
    auto result = m.insert(std::move(node));
    assert(result.inserted && m.at("new") == 7);
}
\end{lstlisting}

\section{Kontenery haszujące: \texttt{unordered\_map} i \texttt{unordered\_set}}
\label{sec:cpp13-unordered}

Kontener \texttt{unordered\_map<Key,T,Hash,Eq>} rozmieszcza elementy
w kubełkach. Funkcja \texttt{Hash} przekształca klucz w wartość typu
\texttt{std::size\_t}; \texttt{Eq} rozstrzyga równość kluczy.
Wymagany jest warunek zgodności:
\[
\operatorname{Eq}(a,b) \Longrightarrow
\operatorname{Hash}(a)=\operatorname{Hash}(b).
\]
Warunek odwrotny nie zachodzi: kolizje haszy są dozwolone i
rozstrzygane przez porównanie kluczy. Iteracja po kontenerze
nie ma stabilnego porządku i nie należy z niego korzystać jako z
interfejsu logicznego programu.
Można wyobrazić sobie hasz jako wskazówkę, w którym kubełku szukać.
Jeśli dwa różne klucze trafią do tego samego kubełka, \texttt{Eq}
sprawdza, czy są naprawdę równe. Stała w poniższym \texttt{PointHash}
służy jedynie do wymieszania haszy obu współrzędnych.

\begin{lstlisting}[caption={Wlasny hash i porownanie kluczy (C++20); kompletny program.}]
#include <cassert>
#include <cstddef>
#include <functional>
#include <unordered_map>

struct Point {
    int x, y;
    friend bool operator==(const Point&, const Point&) = default;
};

struct PointHash {
    std::size_t operator()(const Point& p) const noexcept {
        std::size_t h1 = std::hash<int>{}(p.x);
        std::size_t h2 = std::hash<int>{}(p.y);
        return h1 ^ (h2 + std::size_t{0x9e3779b9}
                        + (h1 << 6) + (h1 >> 2));
    }
};

int main() {
    std::unordered_map<Point, int, PointHash> m;
    m[{2, 5}] = 42;
    assert(m.at(Point{2, 5}) == 42);
}
\end{lstlisting}

Ta funkcja mieszająca służy demonstracji składni, a nie zastosowaniom
kryptograficznym. Dla publicznie sterowanych kluczy przy projektowaniu
systemów wystawionych na ataki kolizyjne należy rozważyć również
odporność na złośliwe dane wejściowe.

\subsection{\texttt{load\_factor}, \texttt{reserve} i \texttt{rehash}}

\texttt{load\_factor()} to stosunek liczby elementów do liczby kubełków.
\texttt{max\_load\_factor()} podaje lub ustawia pożądany maksymalny
współczynnik. \texttt{reserve(n)} żąda przygotowania na $n$ elementów,
a \texttt{rehash(k)} na przynajmniej $k$ kubełków, z uwzględnieniem
limitów wynikających z współczynnika wypełnienia.

\begin{lstlisting}[caption={Rezerwowanie kubelek przed licznymi wstawkami (C++20); kompletny program.}]
#include <cassert>
#include <unordered_set>

int main() {
    std::unordered_set<int> seen;
    seen.reserve(1000);
    for (int i = 0; i < 1000; ++i) seen.insert(i);
    assert(seen.size() == 1000);
    assert(seen.contains(17));
}
\end{lstlisting}

Rehashing unieważnia iteratory, lecz \emph{nie unieważnia} wskaźników
i referencji do elementów. Wstawienie może unieważnić iteratory, jeśli
wywoła rehash. Usunięcie unieważnia iteratory i referencje do usuniętych
elementów. Nie wolno modyfikować kluczy w sposób zaburzający ich hash
lub równość podczas przebywania elementów w kontenerze.

\section{Adaptery: \texttt{stack}, \texttt{queue}, \texttt{priority\_queue}}
\label{sec:cpp13-adaptors}

\texttt{stack} udostępnia schemat LIFO, \texttt{queue} --- FIFO,
a \texttt{priority\_queue} --- dostęp do elementu o najwyższym
priorytecie według komparatora (domyślnie do największego).
Adaptery zwykle nie zapewniają bezpośredniej iteracji po elementach.

\begin{lstlisting}[caption={Trzy adaptery standardowe (C++11); kompletny program.}]
#include <cassert>
#include <functional>
#include <queue>
#include <stack>
#include <vector>

int main() {
    std::stack<int> s;
    s.push(1); s.push(2);
    assert(s.top() == 2);
    s.pop();

    std::queue<int> q;
    q.push(1); q.push(2);
    assert(q.front() == 1);
    q.pop();

    std::priority_queue<int, std::vector<int>, std::greater<int>> pq;
    pq.push(7); pq.push(2); pq.push(5);
    assert(pq.top() == 2); // kopiec minimum
}
\end{lstlisting}

\texttt{priority\_queue} zwykle opiera się na kopcu: \texttt{top()} ma
koszt $O(1)$, a \texttt{push()} i \texttt{pop()} $O(\log n)$. Nie jest
to kolejka FIFO ani struktura z aktualizacją priorytetu dowolnego elementu
przez gotowy publiczny interfejs.

\section{Iteratory, stabilność adresów i unieważnianie}
\label{sec:cpp13-invalidation}

Iterator nie jest ogólnie tożsamy ze wskaźnikiem. Może przechowywać
stan niezbędny do poruszania się po drzewie lub kubełkach. Po operacji
modyfikującej trzeba wiedzieć, czy iterator, referencja i wskaźnik do
elementu pozostają ważne. Typowe reguły podsumowano w tabeli.

\begin{center}
\footnotesize
\begin{tabularx}{\textwidth}{@{}lXX@{}}
\toprule
\textbf{Kontener} & \textbf{Wstawienie} & \textbf{Usunięcie} \\
\midrule
\texttt{vector} & Przy realokacji: wszystkie iteratory i referencje; bez realokacji: pozycja wstawienia i dalej & Pozycja usunięcia i dalej (także dawny \texttt{end()}) \\
\texttt{deque} & Na końcach: iteratory tracą ważność, referencje do elementów zwykle pozostają; w środku: także referencje & Zależy od pozycji; usunięcie w środku unieważnia wszystkie iteratory i referencje \\
\texttt{list} & Istniejące iteratory i referencje zachowane & Tylko usunięte elementy \\
\texttt{map}, \texttt{set} & Istniejące iteratory i referencje zachowane & Tylko usunięte elementy \\
\texttt{unordered\_*} & Iteratory mogą utracić ważność przy rehash; referencje zachowane & Tylko usunięte elementy \\
\bottomrule
\end{tabularx}
\end{center}

Dla \texttt{deque} szczególnie ostrożnie trzeba traktować
\texttt{end()}, którego ważność zależy od rodzaju usunięcia.
W tabeli podano reguły najważniejsze dla praktyki, a nie kompletną
specyfikację każdej przeciążonej operacji. Po operacjach zmieniających
strukturę bezpieczną praktyką jest korzystanie z iteratora zwróconego
przez metodę (np. \texttt{erase}, \texttt{insert}).

\subsection{Bezpieczne usuwanie podczas iteracji}

\begin{lstlisting}[caption={Usuwanie elementow map podczas iteracji (C++20); kompletny program.}]
#include <cassert>
#include <map>

int main() {
    std::map<int, int> m{{1, 2}, {2, 3}, {3, 4}};
    for (auto it = m.begin(); it != m.end();) {
        if (it->second % 2 == 0) it = m.erase(it);
        else ++it;
    }
    assert(m.size() == 1 && m.contains(2));
}
\end{lstlisting}

\paragraph{Szczególny przypadek \texttt{vector<bool>}.}
\texttt{std::vector<bool>} jest specjalizacją oszczędzającą pamięć,
która może reprezentować wartości w postaci bitów. Jej \texttt{operator[]}
nie musi zwracać \texttt{bool\&}, tylko specjalny obiekt proxy. Nie
należy oczekiwać normalnej tablicy obiektów \texttt{bool}, ciągłego
przechowywania pojedynczych \texttt{bool} ani bezpiecznej modyfikacji
odrębnych bitów z różnych wątków bez synchronizacji.

\section{Algorytmy standardowe i wymagania wobec iteratorów}

\begin{center}\small
\begin{tabularx}{\textwidth}{@{}p{0.34\textwidth}X@{}}
\toprule
Algorytm & Działanie i warunki \\
\midrule
\code{std::find} & Szuka pierwszej wartości równej wzorcowi; przy braku zwraca iterator końcowy. \\
\code{std::transform} & Stosuje funkcję do zakresu i zapisuje wyniki przez iterator wyjściowy. \\
\code{std::sort} & Sortuje zakres z iteratorami dostępu swobodnego; nie zachowuje kolejności elementów równoważnych. \\
\code{std::stable_sort} & Sortuje, zachowując wzajemną kolejność elementów równoważnych. \\
\code{std::lower_bound} & Szuka pierwszego elementu nie mniejszego od wzorca; wymaga odpowiednio uporządkowanego zakresu. \\
\code{std::erase_if} & Usuwa z kontenera elementy spełniające predykat; dostępne od C++20 dla wspieranych kontenerów. \\
\bottomrule
\end{tabularx}
\end{center}
\label{sec:cpp13-algorithms}

Algorytmy oddzielają \emph{operację} od \emph{reprezentacji danych}.
Wersje klasyczne przyjmują iteratory, a w C++20 wiele algorytmów
w \texttt{std::ranges} można wywołać bezpośrednio na zakresie. Szczegóły
kategorii iteratorów i widoków omawiamy w rozdziale~\ref{chap:ranges-views-cpp}; tutaj skupiamy
się na doborze algorytmu i jego kontrakcie.

\subsection{Wyszukiwanie, transformacje i agregacje}

Do częstych operacji należą \texttt{find}, \texttt{find\_if},
\texttt{count\_if}, \texttt{transform}, \texttt{accumulate},
\texttt{all\_of}, \texttt{any\_of}, \texttt{none\_of} i
\texttt{for\_each}. Algorytmy numeryczne znajdują się w nagłówku
\texttt{<numeric>}, pozostałe głównie w \texttt{<algorithm>}.

\begin{lstlisting}[caption={Algorytmy klasyczne i Ranges (C++20); kompletny program.}]
#include <algorithm>
#include <cassert>
#include <numeric>
#include <vector>

int main() {
    std::vector<int> values{1, 2, 3, 4, 5};
    auto even = [](int x) { return x % 2 == 0; };

    assert(std::ranges::count_if(values, even) == 2);
    assert(std::ranges::all_of(values,
        [](int x) { return x > 0; }));
    assert(std::accumulate(values.begin(), values.end(), 0) == 15);

    std::vector<int> squares(values.size());
    std::ranges::transform(values, squares.begin(),
                           [](int x) { return x * x; });
    assert(squares[2] == 9);
}
\end{lstlisting}

\subsection{Sortowanie i stabilność kolejności}

\texttt{std::sort} wymaga iteratorów dostępu swobodnego i ma
złożoność $O(n\log n)$ w najgorszym przypadku (w liczbie porównań).
\texttt{stable\_sort} zachowuje kolejność elementów, które komparator
uznaje za równoważne, zwykle kosztem dodatkowych zasobów; jego
szczegółowe gwarancje złożoności zależą od dostępnej pamięci.
\texttt{nth\_element} umieszcza wybrany element na pozycji, którą
zajmowałby po sortowaniu, bez sortowania całej sekwencji; przeciętna
liczba porównań jest liniowa.

\begin{lstlisting}[caption={Stabilne sortowanie rekordow po kluczu (C++20); kompletny program.}]
#include <algorithm>
#include <cassert>
#include <string>
#include <vector>

struct Record {
    std::string name;
    int group;
};

int main() {
    std::vector<Record> v{{"A", 2}, {"B", 1}, {"C", 2}};
    std::ranges::stable_sort(v, {}, &Record::group);
    assert(v[0].name == "B");
    assert(v[1].name == "A" && v[2].name == "C");
}
\end{lstlisting}

\subsection{Wyszukiwanie binarne i porządek wejścia}

\texttt{binary\_search}, \texttt{lower\_bound}, \texttt{upper\_bound}
i \texttt{equal\_range} wymagają odpowiedniego uporządkowania
(lub partycjonowania względem danego klucza) zakresu. Zastosowanie
do nieuporządkowanych danych daje wynik niespełniający kontraktu
algorytmu. W wersjach na iteratorach nie-losowego dostępu liczba
porównań może pozostać logarytmiczna, lecz przesuwanie iteratorów
może być liniowe.

\begin{lstlisting}[caption={Wyszukiwanie miejsca wstawienia w posortowanym vector (C++11); kompletny program.}]
#include <algorithm>
#include <cassert>
#include <vector>

int main() {
    std::vector<int> sorted{1, 3, 3, 7};
    auto it = std::lower_bound(sorted.begin(), sorted.end(), 3);
    assert(it == sorted.begin() + 1);

    auto upper = std::upper_bound(sorted.begin(), sorted.end(), 3);
    assert(upper == sorted.begin() + 3);
    assert(std::binary_search(sorted.begin(), sorted.end(), 7));
}
\end{lstlisting}

\subsection{Idiom \emph{erase--remove} i \texttt{std::erase\_if}}

\texttt{std::remove} nie usuwa elementów \emph{z kontenera}:
przestawia elementy pozostające i zwraca logiczny nowy koniec.
Rozmiar kontenera zmniejsza dopiero wywołanie \texttt{erase}.
Od C++20 wygodniej używać \texttt{std::erase} i
\texttt{std::erase\_if} dla wspieranych kontenerów.

\begin{lstlisting}[caption={Dwa sposoby usuwania z vector (C++20); kompletny program.}]
#include <algorithm>
#include <cassert>
#include <vector>

int main() {
    std::vector<int> a{1, 2, 3, 4, 5};
    auto last = std::remove_if(a.begin(), a.end(),
                               [](int x) { return x % 2 == 0; });
    a.erase(last, a.end());
    assert((a == std::vector<int>{1, 3, 5}));

    std::vector<int> b{1, 2, 3, 4, 5};
    auto removed = std::erase_if(b,
                                 [](int x) { return x % 2 == 0; });
    assert(removed == 2 && b == a);
}
\end{lstlisting}

Dla kontenerów asocjacyjnych nie wolno stosować klasycznego
\texttt{remove\_if} tak jak dla wektora: ich elementy mają klucze,
których pozycja jest regulowana przez strukturę danych. Stosujemy
\texttt{erase} z iteratorem lub \texttt{std::erase\_if} (C++20).

\subsection{Algorytmy kopca i \texttt{std::priority\_queue}}

\texttt{make\_heap}, \texttt{push\_heap}, \texttt{pop\_heap} i
\texttt{sort\_heap} pozwalają użyć własnego wektora jako kopca.
\texttt{make\_heap} działa w czasie liniowym, a operacje dopisania i
usunięcia maksimum w $O(\log n)$. Są to operacje inne niż pełne sortowanie.

\section{Przypadki projektowe: dobór kontenera do problemu}
\label{sec:cpp13-use-cases}

\subsection{Histogram, deduplikacja i lista wyników}

Załóżmy, że chcemy zebrać wartości, usunąć duplikaty i otrzymać
wynik w kolejności rosnącej. Możemy użyć \texttt{std::set}, albo
\texttt{vector} z \texttt{sort} i \texttt{unique}. Drugi sposób bywa
korzystniejszy w przetwarzaniu wsadowym, ponieważ \texttt{vector}
przechowuje dane ciągle i ogranicza liczbę alokacji węzłów.

\begin{lstlisting}[caption={Deduplikacja przez sort i unique (C++11); kompletny program.}]
#include <algorithm>
#include <cassert>
#include <vector>

int main() {
    std::vector<int> keys{5, 2, 5, 1, 2, 8};
    std::sort(keys.begin(), keys.end());
    keys.erase(std::unique(keys.begin(), keys.end()), keys.end());
    assert((keys == std::vector<int>{1, 2, 5, 8}));
}
\end{lstlisting}

\subsection{Własność elementów w kontenerach polimorficznych}

Przechowywanie obiektów klas pochodnych przez wartość w
\texttt{vector<Base>} może powodować \emph{object slicing}. Dla
polimorfizmu dynamicznego używamy np. wektora wskaźników
\texttt{unique\_ptr<Base>}. Przy zamkniętej rodzinie typów można użyć
wektora z elementami \texttt{variant} (zob. rozdział~\ref{chap:cpp14-type-erasure}).
Różnicę między polimorfizmem a obcięciem obiektu omawia
rozdział~\ref{chap:cpp-object-model}.

\begin{lstlisting}[caption={Kontener wlasciwych obiektow polimorficznych (C++14); kompletny program.}]
#include <cassert>
#include <memory>
#include <vector>

struct Shape {
    virtual ~Shape() = default;
    virtual int area() const = 0;
};
struct Square : Shape {
    int side;
    explicit Square(int x) : side(x) {}
    int area() const override { return side * side; }
};

int main() {
    std::vector<std::unique_ptr<Shape>> shapes;
    shapes.push_back(std::make_unique<Square>(3));
    assert(shapes.front()->area() == 9);
}
\end{lstlisting}

\subsection{Kontenery a współbieżność}

Samo wybranie kontenera standardowego nie czyni jego operacji
modyfikujących bezpiecznymi dla współbieżnego wykonania. Równoczesne
modyfikowanie wspólnego \texttt{vector} przez niezależne
\texttt{push\_back()} bez synchronizacji jest wyścigiem danych.
Można natomiast równolegle odczytywać elementy kontenera, jeśli żaden
wątek ich nie modyfikuje, a czas życia kontenera jest zachowany.
Dla wielu kontenerów można także modyfikować \emph{różne istniejące}
elementy równolegle, jeśli nie dochodzi do konfliktu pamięci; wyjątek
stanowi m.in. specjalizacja \texttt{vector<bool>}. Zmiany strukturalne
wymagają osobnej analizy i synchronizacji (zob. rozdział~\ref{chap:concurrency}).

\section{Złożoność, lokalność pamięci i pomiar wydajności}
\label{sec:cpp13-performance}

Teoretyczna złożoność to jeden z elementów oceny. Istotne są również:
liczba alokacji, liczba odwołań do rozproszonych adresów, rozmiar
elementu, przewidywalność gałęzi i zachowanie cache procesora.
Dla wielu zastosowań \texttt{vector} z wyszukiwaniem liniowym
może być szybszy od struktury drzewiastej przy małej liczbie kluczy.
To nie przeczy złożoności asymptotycznej: stałe i charakter obciążenia
są częścią kosztu rzeczywistego programu.

\subsection{Koszt sumowania ciągłych danych}

\begin{lstlisting}[caption={Prosty szkielet pomiaru vector i list (C++20); kompletny program.}]
#include <chrono>
#include <cstdint>
#include <iostream>
#include <list>
#include <numeric>
#include <utility>
#include <vector>

int main() {
    constexpr int N = 200000;
    std::vector<int> v(N, 1);
    std::list<int> l(v.begin(), v.end());

    auto measure = [](const auto& xs) {
        auto begin = std::chrono::steady_clock::now();
        std::int64_t sum = std::accumulate(xs.begin(), xs.end(),
                                           std::int64_t{0});
        auto end = std::chrono::steady_clock::now();
        auto ns = std::chrono::duration_cast<
            std::chrono::nanoseconds>(end - begin).count();
        return std::pair{sum, ns};
    };

    auto [sum_v, time_v] = measure(v);
    auto [sum_l, time_l] = measure(l);
    std::cout << "vector: " << sum_v << ", " << time_v << " ns\n";
    std::cout << "list:   " << sum_l << ", " << time_l << " ns\n";
}
\end{lstlisting}

To \emph{ilustracja sposobu pomiaru}, a nie miarodajny benchmark.
Jednorazowy wynik bywa zdominowany przez rozgrzanie cache, optymalizacje
kompilatora, scheduler i zakłócenia systemowe. Poprawniejszy eksperyment
powtarza pomiary, rozdziela przygotowanie danych od mierzonej operacji,
uwzględnia pamięć, różne rozmiary danych, medianę i rozrzut pomiarów.
Dla wiarygodnych porównań warto użyć np. Google Benchmark i profilera.

\subsection{Przybliżone reguły wyboru}

\begin{center}
\small
\begin{tabularx}{\textwidth}{@{}lX@{}}
\toprule
\textbf{Wymaganie} & \textbf{Rozsądny wybór początkowy} \\
\midrule
Dużo iteracji, obliczeń, dostęp indeksowany & \texttt{vector}; przy stałym rozmiarze \texttt{array} \\
Częste dodawanie i usuwanie z obu końców & \texttt{deque} \\
Utrzymanie kolejności kluczy, wyszukiwanie przedziałowe & \texttt{map} / \texttt{set} \\
Wyszukiwanie po kluczu bez wymogu kolejności & \texttt{unordered\_map} / \texttt{unordered\_set} \\
Przenoszenie istniejących węzłów i stabilne iteratory & \texttt{list} \\
FIFO / LIFO & \texttt{queue} / \texttt{stack} \\
Wielokrotne pobieranie największego lub najmniejszego & \texttt{priority\_queue} \\
\bottomrule
\end{tabularx}
\end{center}

\section{Najczęstsze błędy i dobre praktyki}
\label{sec:cpp13-pitfalls}

\begin{enumerate}
\item Używanie \texttt{vector::operator[]} bez pewności, że indeks jest
  poprawny; w kodzie granicznym warto stosować \texttt{at()}.
\item Traktowanie \texttt{reserve} jako tworzenia elementów (mylenie
  \texttt{size} z \texttt{capacity}).
\item Korzystanie z iteratorów lub referencji po unieważnieniu przez
  realokację, \texttt{erase} lub \texttt{rehash}.
\item Zakładanie, że \texttt{emplace\_back} zawsze jest szybsze od
  \texttt{push\_back}.
\item Używanie \texttt{map::operator[]} do samego sprawdzania istnienia
  klucza, z niezamierzonym wstawieniem nowej wartości.
\item Pisanie komparatorów naruszających \emph{strict weak ordering};
  komparator \texttt{a <= b} nie jest poprawną relacją porządkującą.
\item Używanie \texttt{std::sort} na \texttt{std::list} lub
  \texttt{std::lower\_bound} na danych bez odpowiedniego porządku.
\item Wnioskowanie o kolejności elementów \texttt{unordered\_map}
  na podstawie przypadkowego przebiegu programu.
\item Niepotrzebne stosowanie \texttt{list} z przekonaniem, że
  wstawienie $O(1)$ zawsze wygrywa z \texttt{vector}.
\item Wprowadzanie równoczesnych modyfikacji kontenera z wielu wątków
  bez uzasadnienia braku konfliktu i odpowiedniej synchronizacji.
\end{enumerate}

\section{Pytania kontrolne}

\begin{enumerate}
\item Jaka jest różnica między \texttt{reserve} a \texttt{resize}, i
  które z nich tworzy nowe elementy?
\item Kiedy \texttt{push\_back} na \texttt{vector} unieważnia wszystkie
  referencje, a kiedy zachowuje referencje do wcześniejszych elementów?
\item Dlaczego \texttt{map::operator[]} może być niepożądany podczas
  sprawdzania istnienia klucza?
\item Jakie wymagania muszą spełniać \texttt{Hash} i \texttt{Eq} w
  \texttt{unordered\_map}?
\item Dlaczego \texttt{std::remove} nie zmienia rozmiaru \texttt{vector}?
\item Jak różnią się koszty i kontrakty \texttt{sort}, \texttt{stable\_sort},
  \texttt{nth\_element} i \texttt{lower\_bound}?
\item W jakich okolicznościach \texttt{list::splice} zapewnia korzyść
  niedostępną dla zwykłego \texttt{vector}?
\end{enumerate}

\section{Podsumowanie rozdziału}

Dobry kod C++ wybiera kontener zgodnie z potrzebnymi operacjami i
własnościami pamięci, a nie jedynie na podstawie pojedynczego rzędu
złożoności. \texttt{vector} to często najlepszy punkt wyjścia;
\texttt{deque}, kontenery drzewiaste, haszujące i listy rozwiązują
bardziej szczegółowe problemy. Praktycznie równie ważne jak złożoności
są reguły unieważniania iteratorów, wymagania komparatorów,
bezpieczeństwo czasu życia i świadome wykorzystywanie algorytmów
standardowych zamiast ręcznie pisanych pętli.

\section*{Literatura i dokumentacja}
\addcontentsline{toc}{section}{Literatura i dokumentacja}
\begin{itemize}
\item \url{https://en.cppreference.com/w/cpp/container} --- kontenery i
  zestawienie reguł unieważniania iteratorów.
\item \url{https://en.cppreference.com/w/cpp/algorithm} --- algorytmy
  biblioteki standardowej i ich kontrakty.
\item \url{https://en.cppreference.com/w/cpp/container/vector} ---
  \texttt{std::vector} i zarządzanie pojemnością.
\item \url{https://en.cppreference.com/w/cpp/container/unordered_map} ---
  struktury haszujące.
\item Scott Meyers, \emph{Effective Modern C++}, O'Reilly, 2014.
\item Nicolai M. Josuttis, \emph{The C++ Standard Library}, wydanie drugie,
  Addison-Wesley, 2012.
\end{itemize}

% ROZDZIAL 12 — Model obiektowy C++ i polimorfizm
% Fragment do zbiorczego skryptu; wczytac przez \input{...} PO rozdziale 11.
% Standard podstawowy: C++20. Wymagane pakiety: listings, amsmath, amssymb,
% booktabs, tabularx, array, hyperref, xcolor (jak we wczesniejszych rozdzialach).
% Ten plik celowo nie zawiera \documentclass ani \begin{document}.

\chapter{Model obiektowy C++ i polimorfizm}
\label{chap:cpp-object-model}

\section{Po co poznawać model obiektowy?}

Klasy w C++ są jednocześnie narzędziem modelowania danych, organizowania
interfejsów i budowania różnych rodzajów polimorfizmu. Dwie implementacje
tego samego interfejsu mogą jednak mieć bardzo różne własności:
\begin{itemize}
\item mogą rozstrzygać wywołanie już w trakcie kompilacji albo dopiero w
  czasie wykonania;
\item mogą wymagać alokacji dynamicznej lub działać bez niej;
\item mogą przechowywać obiekty bezpośrednio, przez referencję albo przez
  wskaźnik zarządzający czasem życia;
\item mogą zachowywać tożsamość dynamicznego typu obiektu lub ją utracić
  wskutek \emph{object slicing}.
\end{itemize}

W tym rozdziale rozdzielamy \emph{gwarancje standardu C++} od typowych
technik używanych przez kompilatory, takich jak tablice funkcji wirtualnych
(\emph{vtable}). Pozwala to czytać kod zarówno od strony języka, jak i
możliwych kosztów wykonania.

\section{Obiekty, podobiekty i tożsamość typu}
\label{sec:cpp12-subobjects}

\emph{Obiekt} zajmuje pewien obszar pamięci, ma określony typ i czas życia.
Obiekt klasy może zawierać \emph{podobiekty składowych} oraz
\emph{podobiekty klas bazowych}. Dziedziczenie nie oznacza zatem
,,zamiany'' klasy pochodnej na bazową: obiekt pochodny zawiera w sobie
odpowiedni podobiekt bazowy.

\begin{lstlisting}[caption={Klasa pochodna zawiera podobiekt bazowy (C++11); kompletny program.}]
#include <cassert>
#include <string>

struct Entity {
    int id = 0;
};

struct Employee : Entity {
    std::string name;
};

int main() {
    Employee e;
    e.id = 7;
    e.name = "Ada";
    Entity& base = e;  // odniesienie do podobiektu Entity
    assert(&base == static_cast<Entity*>(&e));
    assert(base.id == 7);
}
\end{lstlisting}

\paragraph{Typ statyczny a typ dynamiczny.}
\emph{Typ statyczny} wyrażenia wynika z reguł języka i jest znany
kompilatorowi. Gdy wyrażenie wskazuje na obiekt klasy polimorficznej
(czyli klasy posiadającej lub dziedziczącej funkcję wirtualną),
\emph{typ dynamiczny} obiektu może być klasą pochodną.

\begin{lstlisting}[caption={Statyczny i dynamiczny typ obiektu; fragment.}]
struct Animal { virtual ~Animal() = default; };
struct Dog : Animal {};

Dog dog;
Animal& ref = dog;  // statyczny typ: Animal&; obiekt: Dog
Animal* ptr = &dog; // statyczny typ: Animal*; obiekt: Dog
\end{lstlisting}

\noindent W C++ nie ma wymogu, aby podobiekt klasy bazowej zajmował
\emph{pierwsze} bajty obiektu pochodnego. Nie należy również zakładać,
że wskaźnik do całego obiektu i wskaźnik do dowolnej jego bazy mają
tę samą reprezentację liczbową. Odpowiednie konwersje wykonuje język.

\subsection{Składnia dziedziczenia i kontrola dostępu}

\begin{lstlisting}[caption={Dziedziczenie publiczne, chronione i prywatne; fragment.}]
class Base {
public:
    int a = 1;
protected:
    int b = 2;
private:
    int c = 3;
};

class PublicDerived    : public Base {};
class ProtectedDerived : protected Base {};
class PrivateDerived   : private Base {};
\end{lstlisting}

Przy dziedziczeniu publicznym publiczne elementy interfejsu bazowego
pozostają publiczne, a chronione pozostają chronione. Przy dziedziczeniu
chronionym oba rodzaje stają się chronione, natomiast przy prywatnym
oba stają się prywatne dla użytkowników klasy pochodnej.
Prywatne elementy bazy pozostają niedostępne \emph{bezpośrednio}
wewnątrz klasy pochodnej niezależnie od rodzaju dziedziczenia.

Sama dostępność elementu nie zmienia jego faktycznego istnienia
w podobiekcie bazowym. Publiczne dziedziczenie najczęściej wyraża
relację \emph{is-a}: obiekt pochodny może być używany tam, gdzie
oczekiwany jest interfejs bazowy.

\section{Przeciążanie, przesłanianie i funkcje wirtualne}
\label{sec:cpp12-virtual}

\emph{Przeciążanie} (\emph{overloading}) polega na deklarowaniu kilku
funkcji o tej samej nazwie, różniących się listą parametrów.
\emph{Przesłanianie nazwy} (\emph{name hiding}) zachodzi, gdy deklaracja
w klasie pochodnej zasłania nazwę z klasy bazowej.
\emph{Nadpisywanie} (\emph{overriding}) oznacza dostarczenie implementacji
funkcji wirtualnej z bazy przez klasę pochodną.

\begin{lstlisting}[caption={Dynamiczne wiązanie wywołania przez virtual (C++11); kompletny program.}]
#include <iostream>

struct Shape {
    virtual ~Shape() = default;
    virtual double area() const = 0;
    virtual void print() const {
        std::cout << "Shape\n";
    }
};

struct Square final : Shape {
    explicit Square(double side) : side_(side) {}
    double area() const override { return side_ * side_; }
    void print() const override {
        std::cout << "Square: " << area() << '\n';
    }
private:
    double side_;
};

int main() {
    Square square{3.0};
    Shape& shape = square;
    shape.print(); // Square: 9
    std::cout << shape.area() << '\n';
}
\end{lstlisting}

Klasa \texttt{Shape} jest \emph{abstrakcyjna}, ponieważ jej funkcja
\texttt{area()} jest czysto wirtualna, co oznaczamy \texttt{= 0}.
Nie można skonstruować obiektu \texttt{Shape}, lecz można utworzyć
referencję lub wskaźnik do jej podobiektu w obiekcie konkretnym.

\paragraph{Reguła wyboru funkcji.}
Dla wywołania \texttt{base.f()} najpierw rozstrzygane są dostępność
nazwy, przeciążenia i poprawność wywołania według \emph{typu statycznego}.
Jeśli wybrana funkcja jest wirtualna, a wywołanie nie wyłącza mechanizmu
wirtualnego przez kwalifikację nazwy, uruchamiany jest odpowiedni
\emph{final overrider} dla dynamicznego typu obiektu.

\subsection{Dlaczego stosować \texttt{override} i \texttt{final}?}

\texttt{override} oznacza, że dana metoda ma nadpisywać metodę wirtualną.
Jeśli sygnatura nie pasuje, kompilator zgłosi błąd. \texttt{final}
uniemożliwia dalsze nadpisywanie metody albo dziedziczenie po klasie.

\begin{lstlisting}[caption={override wykrywa niezgodność sygnatur; fragment.}]
struct Base {
    virtual void draw() const {}
    virtual ~Base() = default;
};

struct Good : Base {
    void draw() const override {} // poprawnie
};

// BLEDNE -- celowo niekompilowalne:
// struct Bad : Base {
//     void draw() override {} // brak const
// };
\end{lstlisting}

\noindent Przy nadpisywaniu znaczenie mają między innymi kwalifikatory
\texttt{const}/\texttt{volatile} metody i kwalifikator referencyjny
\texttt{\&}/\texttt{\&\&}. Specyfikacja wyjątków też podlega regułom
zgodności: nadpisanie metody \texttt{noexcept} metodą potencjalnie
zgłaszającą wyjątek jest niedozwolone. Zwykłego typu wyniku nie można
zmienić, z wyjątkiem określonych \emph{kowariantnych} typów zwracanych.

\begin{lstlisting}[caption={noexcept i kwalifikatory referencyjne w nadpisywaniu; fragment.}]
struct Service {
    virtual int read() const & noexcept { return 1; }
    virtual ~Service() = default;
};

struct FastService : Service {
    int read() const & noexcept override { return 2; }
};
\end{lstlisting}

\subsection{Przesłanianie przeciążeń i \texttt{using}}

\begin{lstlisting}[caption={Odsłanianie przeciążeń klasy bazowej; kompletny program.}]
#include <cassert>

struct Printer {
    int print(int) const { return 1; }
    int print(double) const { return 2; }
};

struct FancyPrinter : Printer {
    using Printer::print;
    int print(const char*) const { return 3; }
};

int main() {
    FancyPrinter p;
    assert(p.print(10) == 1);
    assert(p.print(2.5) == 2);
    assert(p.print("text") == 3);
}
\end{lstlisting}

Bez deklaracji \texttt{using Printer::print;} nazwa \texttt{print}
z klasy pochodnej przesłaniałaby zestaw przeciążeń o tej nazwie
z klasy bazowej. Mechanizm przesłaniania występuje również dla
funkcji \emph{niewirtualnych}.

\subsection{Argumenty domyślne i wiązanie statyczne}

Wartości domyślnych argumentów nie wybiera się polimorficznie:
są związane ze statycznym typem wyrażenia, mimo że sama implementacja
funkcji może być wybrana dynamicznie.

\begin{lstlisting}[caption={Domyślny argument pochodzi z deklaracji bazy (C++11); kompletny program.}]
#include <cassert>

struct Base {
    virtual int calculate(int x = 10) const { return x; }
    virtual ~Base() = default;
};

struct Derived : Base {
    int calculate(int x = 20) const override { return 2 * x; }
};

int main() {
    Derived object;
    Base& base = object;
    assert(base.calculate() == 20); // 2 * 10
    assert(object.calculate() == 40); // 2 * 20
}
\end{lstlisting}

Z tego powodu w praktycznych interfejsach polimorficznych warto unikać
różnych domyślnych argumentów dla nadpisywanych funkcji.

\section{Destruktor wirtualny i reguły czasu życia}
\label{sec:cpp12-lifetime}

Jeżeli obiekt klasy pochodnej ma być niszczony przez wskaźnik do
klasy bazowej, destruktor bazy powinien być \texttt{public virtual}
(alternatywnie projekt może zabraniać usuwania przez bazę,
np. za pomocą \texttt{protected} destruktora niewirtualnego).

\begin{lstlisting}[caption={Bezpieczne usuwanie obiektu przez wskaźnik bazowy (C++14); kompletny program.}]
#include <iostream>
#include <memory>

struct Resource {
    virtual ~Resource() { std::cout << "~Resource\n"; }
};

struct FileResource : Resource {
    ~FileResource() override {
        std::cout << "~FileResource\n";
    }
};

int main() {
    std::unique_ptr<Resource> object =
        std::make_unique<FileResource>();
    // Najpierw ~FileResource, potem ~Resource.
}
\end{lstlisting}

Jeżeli bazowy destruktor nie jest wirtualny, standardowe
\texttt{delete base\_ptr;} przy dynamicznym typie różnym od statycznego
może prowadzić do niezdefiniowanego zachowania; istnieją szczególne
reguły związane z tzw. \emph{destroying delete}, ale nie należy
traktować ich jako zwykłego substytutu destruktora wirtualnego.

\subsection{Kolejność konstrukcji i destrukcji}

Dla obiektu klasy pochodnej kolejność konstrukcji jest następująca:
\begin{enumerate}
\item wirtualne podobiekty bazowe (przez najbardziej pochodny obiekt);
\item niewirtualne podobiekty bazowe w kolejności ich deklaracji;
\item składowe klasy w kolejności \emph{deklaracji pól}, niezależnie od
      kolejności w liście inicjalizacyjnej;
\item ciało konstruktora klasy.
\end{enumerate}
Destrukcja przebiega w odpowiedniej kolejności odwrotnej.
W prostym pojedynczym dziedziczeniu można to zilustrować następująco.

\begin{lstlisting}[caption={Kolejność konstruowania podobiektów (C++11); kompletny program.}]
#include <iostream>

struct A {
    A()  { std::cout << "A+ "; }
    ~A() { std::cout << "A- "; }
};
struct B {
    B()  { std::cout << "B+ "; }
    ~B() { std::cout << "B- "; }
};
struct C : A {
    B member;
    C()  { std::cout << "C+ "; }
    ~C() { std::cout << "C- "; }
};

int main() {
    C object;
    // A+ B+ C+ C- B- A-
}
\end{lstlisting}

\subsection{Funkcje wirtualne w konstruktorach i destruktorach}

Podczas wykonywania konstruktora \texttt{Base} wywołanie wirtualne
nie trafia do implementacji z \texttt{Derived}, której część nie
została jeszcze skonstruowana. Analogicznie podczas wykonywania
destruktora \texttt{Base} część pochodna została już zniszczona.

\begin{lstlisting}[caption={Dynamiczne wiązanie podczas konstrukcji (C++11); kompletny program.}]
#include <cassert>

struct Base {
    Base() { assert(code() == 1); }
    virtual ~Base() = default;
    virtual int code() const { return 1; }
};
struct Derived : Base {
    int code() const override { return 2; }
};
int main() {
    Derived object;
    Base& b = object;
    assert(b.code() == 2);
}
\end{lstlisting}

\noindent Nie należy wywoływać funkcji \emph{czysto wirtualnej}
przez mechanizm wirtualny z konstruktora albo destruktora klasy,
na której poziomie funkcja jest czysto wirtualna: takie wywołanie
jest niezdefiniowanym zachowaniem. Preferuj dwufazową konstrukcję
kontrolowaną przez fabrykę tylko wtedy, gdy naprawdę jest potrzebna,
a w typowych sytuacjach inicjalizację kompletnych podobiektów
przez konstruktory i RAII.

\section{Object slicing: utrata części pochodnej}
\label{sec:cpp12-slicing}

\emph{Object slicing} nie jest dynamicznym rzutowaniem. Gdy obiekt
pochodny kopiujemy \emph{przez wartość} do obiektu bazy,
powstaje osobny obiekt bazy, zawierający tylko część bazową.

\begin{lstlisting}[caption={Kopiowanie przez wartość powoduje slicing (C++11); kompletny program.}]
#include <cassert>

struct Base {
    virtual ~Base() = default;
    virtual int value() const { return 1; }
};
struct Derived : Base {
    int value() const override { return 2; }
};

int by_value(Base object) { return object.value(); }
int by_reference(const Base& object) {
    return object.value();
}

int main() {
    Derived d;
    assert(by_value(d) == 1);      // nowy obiekt Base
    assert(by_reference(d) == 2);  // oryginalny Derived
}
\end{lstlisting}

Zwykle przechowujemy heterogeniczne obiekty polimorficzne
jako \texttt{std::unique\_ptr<Base>} (własność wyłączna),
\texttt{std::shared\_ptr<Base>} (współdzielona własność) albo
nieposiadające obiektu \texttt{Base\&}/\texttt{Base*} (widoki).
Bezpośredni \texttt{std::vector<Base>} przechowuje wyłącznie obiekty
\texttt{Base} i nie zachowuje typów pochodnych przy kopiowaniu.

\subsection{Bezpieczny kontener polimorficzny}

\begin{lstlisting}[caption={Własność obiektów pochodnych w kontenerze unikalnych wskaźników (C++14); kompletny program.}]
#include <cassert>
#include <memory>
#include <vector>

struct Operation {
    virtual ~Operation() = default;
    virtual int apply(int x) const = 0;
};
struct Add : Operation {
    explicit Add(int delta) : delta_(delta) {}
    int apply(int x) const override { return x + delta_; }
private:
    int delta_;
};
struct Multiply : Operation {
    explicit Multiply(int factor) : factor_(factor) {}
    int apply(int x) const override { return x * factor_; }
private:
    int factor_;
};

int main() {
    std::vector<std::unique_ptr<Operation>> operations;
    operations.push_back(std::make_unique<Add>(2));
    operations.push_back(std::make_unique<Multiply>(3));

    int result = 4;
    for (const auto& operation : operations)
        result = operation->apply(result);
    assert(result == 18); // (4 + 2) * 3
}
\end{lstlisting}

Warto pamiętać, że unikalny wskaźnik na typ pochodny można
przenieść do unikalnego wskaźnika na bazę tylko wtedy, gdy
spełnione są wymagania dotyczące konwersji wskaźników i deletera. Domyślny deleter w typowym
przypadku usuwa obiekt przez \texttt{Base*}, dlatego destruktor
polimorficznej bazy musi być poprawnie zaprojektowany.

\section{RTTI: \texttt{dynamic\_cast} oraz \texttt{typeid}}
\label{sec:cpp12-rtti}

RTTI (\emph{run-time type information}) pozwala w pewnych sytuacjach
sprawdzić typ dynamiczny i wykonać sprawdzoną konwersję w hierarchii
klas polimorficznych.

\subsection{\texttt{dynamic\_cast} — sprawdzony downcast}

\begin{lstlisting}[caption={Sprawdzona konwersja w dół hierarchii (C++11); kompletny program.}]
#include <cassert>
#include <typeinfo>

struct Animal { virtual ~Animal() = default; };
struct Dog : Animal { int bark() const { return 10; } };
struct Cat : Animal {};

int main() {
    Dog dog;
    Animal* animal = &dog;

    if (auto* ptr = dynamic_cast<Dog*>(animal)) {
        assert(ptr->bark() == 10);
    }
    assert(dynamic_cast<Cat*>(animal) == nullptr);

    Animal& reference = dog;
    assert(&dynamic_cast<Dog&>(reference) == &dog);
    try {
        (void)dynamic_cast<Cat&>(reference);
        assert(false); // nie powinno zostac wykonane
    } catch (const std::bad_cast&) {
        // Nieudany dynamic_cast do referencji rzuca bad_cast.
    }
}
\end{lstlisting}

W typowym \emph{downcast} źródłowy typ klasy musi być polimorficzny.
Nieudane rzutowanie do wskaźnika daje \texttt{nullptr}, a do referencji
zgłasza \texttt{std::bad\_cast}. Konwersja zwykłego wskaźnika
\texttt{Derived*} do dostępnej, jednoznacznej bazy \texttt{Base*}
nie potrzebuje RTTI ani \texttt{dynamic\_cast}.

\paragraph{\texttt{static\_cast} nie sprawdza typu dynamicznego.}

\begin{lstlisting}[caption={Niebezpieczny downcast; fragment celowo błędny.}]
struct Base { virtual ~Base() = default; };
struct Derived : Base { int only_derived = 7; };

Base base;
Base* ptr = &base;
// Kompilator moze zaakceptowac ponizsze rzutowanie,
// ale wskazywany obiekt nie jest obiektem Derived:
// Derived* wrong = static_cast<Derived*>(ptr);
// int x = wrong->only_derived; // niezdefiniowane zachowanie
\end{lstlisting}

\noindent Stosuj \texttt{static\_cast} przy downcast tylko wtedy,
gdy poprawny typ dynamiczny wynika z udowodnionego niezmiennika.
Jeżeli architektura wymaga bardzo wielu wywołań \texttt{dynamic\_cast},
rozważ lepszy interfejs wirtualny, visitor albo \texttt{std::variant}.

\subsection{\texttt{typeid} i typ dynamiczny}

\begin{lstlisting}[caption={Odczyt informacji o typie dynamicznym (C++11); kompletny program.}]
#include <cassert>
#include <typeinfo>

struct Base { virtual ~Base() = default; };
struct Derived : Base {};

int main() {
    Derived object;
    Base& base = object;
    assert(typeid(base) == typeid(Derived));
    assert(typeid(object) == typeid(Derived));
}
\end{lstlisting}

\noindent Dla wyrażenia typu klasy polimorficznej, będącego \emph{glvalue},
\texttt{typeid(expression)} zwraca informację o typie dynamicznym.
W innych sytuacjach zazwyczaj uwzględnia typ statyczny wyrażenia.
Nazwy zwracane przez \texttt{type\_info::name()} są zależne
od implementacji; nie należy używać ich jako stabilnego identyfikatora
w zapisywanych plikach lub w protokołach sieciowych.
Dla \texttt{typeid(*ptr)} z zerowym wskaźnikiem do typu polimorficznego
może zostać zgłoszony \texttt{std::bad\_typeid}.

\section{Wielodziedziczenie i wirtualne klasy bazowe}
\label{sec:cpp12-multiple}

C++ pozwala dziedziczyć po wielu klasach. Jest to przydatne
np. dla kilku niezależnych interfejsów, ale wymaga zrozumienia
niejednoznaczności oraz struktury podobiektów bazowych.

\begin{lstlisting}[caption={Dwie niezależne bazy-interfejsy (C++11); kompletny program.}]
#include <cassert>

struct Readable {
    virtual ~Readable() = default;
    virtual int read() const = 0;
};
struct Writable {
    virtual ~Writable() = default;
    virtual void write(int x) = 0;
};
struct Buffer : Readable, Writable {
    int value = 0;
    int read() const override { return value; }
    void write(int x) override { value = x; }
};

int main() {
    Buffer buffer;
    Writable& out = buffer;
    Readable& in = buffer;
    out.write(17);
    assert(in.read() == 17);
}
\end{lstlisting}

\subsection{Problem diamentu}

Załóżmy, że \texttt{Left} i \texttt{Right} dziedziczą po
\texttt{Base}, a \texttt{Diamond} po obu. Przy zwykłym
dziedziczeniu istnieją \emph{dwa różne} podobiekty \texttt{Base}.
Przy dziedziczeniu wirtualnym najbardziej pochodny obiekt zawiera
jedną współdzieloną wirtualną bazę \texttt{Base}.

\begin{lstlisting}[caption={Jedna wirtualna baza we wzorcu diamentu (C++11); kompletny program.}]
#include <cassert>

struct Root {
    explicit Root(int n) : value(n) {}
    int value;
};
struct Left : virtual Root {
    Left() : Root(1) {} // ignorowane jako inicjalizacja
};                          // bazy w obiekcie Diamond
struct Right : virtual Root {
    Right() : Root(2) {}
};
struct Diamond : Left, Right {
    Diamond() : Root(42), Left(), Right() {}
};

int main() {
    Diamond d;
    Root* from_left = static_cast<Left*>(&d);
    Root* from_right = static_cast<Right*>(&d);
    assert(from_left == from_right);
    assert(d.value == 42);
}
\end{lstlisting}

\noindent Najbardziej pochodna klasa \texttt{Diamond} odpowiada
za konstrukcję wirtualnej bazy \texttt{Root}. W implementacjach
wielodziedziczenie i wirtualne bazy mogą wymagać korekty wskaźników
oraz dodatkowych informacji o położeniu podobiektów.
Nie należy jednak utożsamiać modelu języka z konkretną reprezentacją
pamięci w danym ABI.

\section{Kowariantne typy zwracane}

W szczególnym przypadku nadpisująca funkcja wirtualna może zwracać
wskaźnik lub referencję do typu bardziej pochodnego, zamiast
wskaźnika lub referencji do klasy bazowej.

\begin{lstlisting}[caption={Kowariantny typ wyniku metody wirtualnej (C++11); kompletny program.}]
#include <cassert>

struct Base {
    virtual ~Base() = default;
    virtual Base* self() { return this; }
};
struct Derived : Base {
    Derived* self() override { return this; }
    int payload = 5;
};

int main() {
    Derived d;
    Base& b = d;
    Base* a = b.self(); // typ statyczny Base*
    Derived* x = d.self(); // typ statyczny Derived*
    assert(a == x);
    assert(x->payload == 5);
}
\end{lstlisting}

\noindent Kowariancja nie obejmuje dowolnych szablonów: na przykład
\texttt{std::unique\_ptr<Derived>} nie jest kowariantnym typem zwracanym
nadpisującym metodę zwracającą \texttt{std::unique\_ptr<Base>}.
Dlatego metody \texttt{clone()} zwracające \texttt{unique\_ptr}
w klasach pochodnych zwykle zachowują wspólny typ wyniku.

\section{Klasy abstrakcyjne jako kontrakty i wzorzec NVI}
\label{sec:cpp12-nvi}

Dobrą praktyką może być rozdzielenie publicznego, niewirtualnego
interfejsu kontrolującego niezmienniki od prywatnej implementacji
wirtualnej (\emph{non-virtual interface}, NVI).

\begin{lstlisting}[caption={NVI: publiczny kontrakt i prywatna metoda wirtualna (C++11); kompletny program.}]
#include <cassert>
#include <stdexcept>

class Encoder {
public:
    virtual ~Encoder() = default;

    int encode(int value) const {
        if (value < 0)
            throw std::invalid_argument("negative input");
        return do_encode(value);
    }

private:
    virtual int do_encode(int value) const = 0;
};

class DoubleEncoder final : public Encoder {
private:
    int do_encode(int value) const override {
        return 2 * value;
    }
};

int main() {
    DoubleEncoder encoder;
    Encoder& api = encoder;
    assert(api.encode(21) == 42);
}
\end{lstlisting}

\noindent Wersja publiczna może centralnie realizować walidację,
logowanie, kontrolę błędów i niezmienników. Prywatna funkcja
wirtualna w klasie bazowej \emph{może} zostać nadpisana przez klasę
pochodną, mimo że nie jest z niej bezpośrednio wywoływalna
przez zwykłą nazwę; w przykładzie nadpisanie jest prywatne.

\subsection{Polimorficzne kopiowanie przez \texttt{clone}}

Gdy chcemy uzyskać niezależną kopię obiektu o zachowanym typie
dynamicznym, możemy zastosować wirtualną fabrykę \texttt{clone()}.

\begin{lstlisting}[caption={Kopiowanie polimorficzne (C++14); kompletny program.}]
#include <cassert>
#include <memory>

struct Node {
    virtual ~Node() = default;
    virtual std::unique_ptr<Node> clone() const = 0;
    virtual int value() const = 0;
};
struct Number final : Node {
    explicit Number(int x) : x_(x) {}
    std::unique_ptr<Node> clone() const override {
        return std::make_unique<Number>(*this);
    }
    int value() const override { return x_; }
private:
    int x_;
};

int main() {
    std::unique_ptr<Node> a = std::make_unique<Number>(42);
    std::unique_ptr<Node> b = a->clone();
    assert(a.get() != b.get());
    assert(a->value() == b->value());
}
\end{lstlisting}

Takie podejście przywraca semantykę kopiowania do hierarchii,
której obiekty zwykle są trzymane przez wskaźniki.
Trzeba jednak jawnie zdefiniować, co oznacza \emph{głęboka kopia},
zwłaszcza jeśli obiekt współdzieli inne zasoby.

\section{Model implementacyjny: \emph{vptr}, \emph{vtable} i ABI}
\label{sec:cpp12-vtable}

Standard C++ określa \emph{zachowanie}, ale nie nakazuje użycia
żadnej konkretnej struktury wewnętrznej do realizacji funkcji
wirtualnych. Popularne ABI używają mniej więcej następującego modelu:
\begin{itemize}
\item obiekt polimorficzny przechowuje ukryty wskaźnik \emph{vptr};
\item wskaźnik ten prowadzi do tablicy \emph{vtable}, która zawiera
  informacje potrzebne do wywołań wirtualnych i często do RTTI;
\item wywołanie wirtualne może wymagać odczytania \texttt{vptr},
  odpowiedniego wpisu i wykonania wywołania pośredniego.
\end{itemize}

To \emph{opis typowej implementacji}, nie kontrakt języka.
W szczególności liczba i położenie wskaźników \texttt{vptr},
układ tablic oraz narzut pamięci są zależne od ABI i kompilatora.
Wielodziedziczenie może komplikować ten schemat, ponieważ niektóre
wywołania wymagają korekty wskaźnika \texttt{this}.

\begin{center}
\fbox{\begin{minipage}{0.91\textwidth}
\textbf{Schemat poglądowy (nie jest układem wymaganym przez standard)}
\medskip

\begin{tabular}{@{}p{0.43\textwidth}cp{0.36\textwidth}@{}}
\textbf{Obiekt \texttt{Derived}} & & \textbf{Tablica wirtualna} \\
\hline
\texttt{vptr} & $\longrightarrow$ & adres implementacji \texttt{draw()} \\
\texttt{Base::id} & & adres odpowiedniego destruktora \\
\texttt{Derived::payload} & & metadane zależne od ABI \\
\end{tabular}
\end{minipage}}
\end{center}

\subsection{\texttt{sizeof}, wyrównanie i brak gwarancji binarnego układu}

\begin{lstlisting}[caption={Empiryczny rozmiar klasy polimorficznej (C++11); kompletny program.}]
#include <iostream>

struct Plain { int x = 0; };
struct Poly {
    int x = 0;
    virtual ~Poly() = default;
};

int main() {
    std::cout << sizeof(Plain) << '\n';
    std::cout << sizeof(Poly) << '\n';
}
\end{lstlisting}

Nie wpisujemy w skrypcie stałych wartości tych rozmiarów:
na wielu platformach typ \texttt{Poly} będzie większy, ale
rozmiary wynikają z implementacji i obowiązujących reguł
wyrównania. Sama obecność \texttt{virtual} nie gwarantuje określonego
liczbowo narzutu ani ustalonej reprezentacji binarnej.
Nie należy serializować surowych bajtów obiektu polimorficznego
jako trwałego formatu danych.

\subsection{Koszty i devirtualizacja}

Wywołanie wirtualne może być droższe od zwykłego z powodu
pośrednictwa oraz potencjalnie trudniejszego \emph{inliningu}.
W praktyce jednak koszt zależy od rozkładu typów dynamicznych,
lokalności pamięci, optymalizacji międzyjednostkowej i architektury CPU.
Kompilator może zastosować \emph{devirtualizację}: jeśli potrafi
ustalić docelową implementację, wywołanie może stać się bezpośrednie,
a następnie zostać rozwinięte inline.

\begin{lstlisting}[caption={final może ułatwiać ustalenie celu wywołania; fragment.}]
struct Interface {
    virtual int eval(int) const = 0;
    virtual ~Interface() = default;
};
struct Square final : Interface {
    int eval(int x) const override { return x * x; }
};

int evaluate(const Square& x, int value) {
    return x.eval(value); // cel wywolania jest znany
}
\end{lstlisting}

\noindent Nie oznacza to, że \texttt{final} zawsze poprawia wydajność.
Porównania należy prowadzić na przykładach reprezentatywnych
i przy realistycznych opcjach optymalizacji. Wąskim gardłem może
być w istocie rozproszony układ obiektów w pamięci, a nie sam
koszt wywołania funkcji wirtualnej.

\section{Alternatywa: polimorfizm statyczny przez szablony i CRTP}
\label{sec:cpp12-static-poly}

Polimorfizm \emph{statyczny} wykorzystuje dedukcję typów i
instancjowanie szablonów, dzięki czemu kompilator zna konkretny typ
w miejscu użycia. Pozwala to generować kod wyspecjalizowany
bez konieczności utrzymywania wspólnej bazy wirtualnej.

\subsection{Szablon funkcji i concepts (C++20)}

\begin{lstlisting}[caption={Polimorfizm statyczny przez concept (C++20); kompletny program.}]
#include <cassert>
#include <concepts>

struct Circle {
    double radius;
    double area() const { return 3.141592653589793 * radius * radius; }
};
struct Rectangle {
    double a, b;
    double area() const { return a * b; }
};

template<class T>
concept HasArea = requires(const T& x) {
    { x.area() } -> std::convertible_to<double>;
};

template<HasArea T>
double area_of(const T& shape) {
    return shape.area();
}

int main() {
    assert(area_of(Rectangle{2, 3}) == 6.0);
    assert(area_of(Circle{2}) > 12.5);
}
\end{lstlisting}

Powyższy kod nie wymaga wspólnej klasy bazowej. Ograniczeniem jest to,
że w typowym kodzie szablonowym zestaw typów musi być znany przy
kompilacji miejsca instancjowania. Nie otrzymujemy automatycznie
heterogenicznego \texttt{std::vector<T>} dla wielu różnych typów.

\subsection{CRTP (\emph{Curiously Recurring Template Pattern})}

\begin{lstlisting}[caption={CRTP z metodą dostarczaną przez bazę (C++11); kompletny program.}]
#include <cassert>

template<class Derived>
struct Printable {
    int twice() const {
        const auto& self = static_cast<const Derived&>(*this);
        return 2 * self.value();
    }
};

struct Number : Printable<Number> {
    int data;
    explicit Number(int x) : data(x) {}
    int value() const { return data; }
};

int main() {
    Number n{21};
    assert(n.twice() == 42);
}
\end{lstlisting}

CRTP nie jest bezpośrednim zamiennikiem wszystkich interfejsów
wirtualnych: wyrażenie \texttt{Printable<A>*} nie jest uniwersalnym
wskaźnikiem do dowolnego \texttt{Printable<B>}. Może natomiast
umożliwiać \emph{static dispatch}, implementację domieszek
(\emph{mixins}) i wykrywanie błędów już przy instancjowaniu szablonów.

\section{Polimorfizm bez dziedziczenia: \texttt{std::variant}}
\label{sec:cpp12-variant}

Jeżeli lista możliwych typów jest zamknięta i znana przy kompilacji,
\texttt{std::variant} (C++17) pozwala przechowywać jedną z kilku
alternatyw \emph{przez wartość}. Zamiast wirtualnego wywołania
używa się zwykle \texttt{std::visit}.

\begin{lstlisting}[caption={Zamknięty zbiór typów i std::visit (C++17); kompletny program.}]
#include <cassert>
#include <variant>

struct Add { int x; };
struct Multiply { int x; };
using Instruction = std::variant<Add, Multiply>;

int execute(const Instruction& instruction, int input) {
    return std::visit([input](const auto& action) -> int {
        using T = std::decay_t<decltype(action)>;
        if constexpr (std::is_same_v<T, Add>)
            return input + action.x;
        else
            return input * action.x;
    }, instruction);
}

int main() {
    assert(execute(Add{2}, 4) == 6);
    assert(execute(Multiply{3}, 4) == 12);
}
\end{lstlisting}

\noindent To podejście ułatwia przechowywanie wartości bez osobnych
alokacji i ma inną charakterystykę wydajnościową. Ceną jest
konieczność jawnego wyliczenia alternatywnych typów w definicji
\texttt{variant} oraz aktualizowania odwiedzających, gdy pojawiają
się nowe alternatywy. W przykładzie użyto też
\texttt{std::decay\_t} oraz \texttt{std::is\_same\_v}
z nagłówka \texttt{<type\_traits>}.

\section{Polimorfizm dynamiczny i statyczny — porównanie}
\label{sec:cpp12-comparison}

\begin{center}
\small
\begin{tabularx}{\textwidth}{@{}p{0.23\textwidth}XXX@{}}
\toprule
\textbf{Cecha} & \textbf{\texttt{virtual}} & \textbf{Szablon/CRTP} & \textbf{\texttt{variant}} \\
\midrule
Zbiór implementacji & Otwarty w granicach hierarchii & Ustalany przy instancjowaniu & Zamknięty zestaw alternatyw \\
Wspólny uchwyt & \texttt{Base\&}, \texttt{Base*} & Zależny od typu & Jeden typ sumaryczny \\
Mechanizm wywołania & Zwykle dynamiczny & Zwykle statyczny & Wybór aktywnej alternatywy \\
Przechowywanie & Często wskaźniki do obiektów & Zazwyczaj konkretne obiekty & Wartość wewnątrz \texttt{variant} \\
Rozszerzanie o nowy typ & Zwykle bez zmiany wywołującego & Nowa instancjacja & Zmiana listy typów \\
Typowe koszty & Dyspozycja i lokalność danych & Rozmiar kodu po instancjacjach & Rozmiar największej alternatywy i odwiedzanie \\
\bottomrule
\end{tabularx}
\end{center}

Wybór zależy przede wszystkim od architektury programu:
\begin{itemize}
\item \textbf{Interfejsy pluginów i otwarte hierarchie}: zwykle
  polimorfizm dynamiczny, z uwzględnieniem stabilności ABI.
\item \textbf{Algorytmy generyczne na znanych typach}: zwykle
  szablony i concepts.
\item \textbf{Zamknięta rodzina wariantów wartości}: często
  \texttt{std::variant}.
\item \textbf{Wiele małych obiektów w intensywnej pętli}: przy wyborze
  należy zmierzyć także koszty alokacji, lokalność danych i wynik
  optymalizacji kompilatora.
\end{itemize}

\section{Reguła zera/pięciu w klasach polimorficznych}

Wprowadzenie destruktora wirtualnego powoduje, że automatyczne
generowanie operacji przenoszących nie działa tak samo jak w prostej
klasie agregującej pola. Dodatkowo publiczne kopiowanie samego
podobiektu bazowego może umożliwiać przypadkowe obcięcie obiektu.

W wielu hierarchiach interfejsowych rozsądne jest zabronienie
kopiowania i świadome zarządzanie obiektami przez smart pointery:

\begin{lstlisting}[caption={Niekopiowalna klasa bazowa interfejsu (C++11); fragment.}]
class Interface {
public:
    Interface() = default;
    virtual ~Interface() = default;
    Interface(const Interface&) = delete;
    Interface& operator=(const Interface&) = delete;
    Interface(Interface&&) = delete;
    Interface& operator=(Interface&&) = delete;

    virtual void run() = 0;
};
\end{lstlisting}

\noindent Nie jest to reguła uniwersalna: klasy reprezentujące wartości
mogą być kopiowalne, a wybrane hierarchie wspierają jawne
\texttt{clone()}. Istotne, aby polityka kopiowania, przenoszenia
i własności była świadomym elementem kontraktu.

\section{Typowe błędy i pytania kontrolne}

\begin{enumerate}
\item \textbf{Brak wirtualnego destruktora bazy}, mimo usuwania
  obiektów pochodnych przez wskaźnik do bazy.
\item \textbf{Brak \texttt{override}} i przypadkowe utworzenie innej metody
  przez pominięcie \texttt{const}, zmianę parametrów albo
  kwalifikatora \texttt{\&}/\texttt{\&\&}.
\item \textbf{Slicing} podczas przekazania przez wartość lub
  wstawiania do kontenera bazowych wartości.
\item \textbf{Bezpodstawny \texttt{static\_cast} w dół hierarchii}
  bez gwarancji rzeczywistego typu obiektu.
\item \textbf{Wywołanie funkcji czysto wirtualnej} podczas konstrukcji
  lub destrukcji na etapie, gdy implementacja pochodna nie jest aktywna.
\item \textbf{Mylenie implementacyjnej vtable z gwarancją standardu}.
\item \textbf{Zależność od domyślnych argumentów metod wirtualnych}
  różniących się między bazą a klasą pochodną.
\item \textbf{Nadmierna hierarchia klas}, gdy prostsza jest kompozycja,
  szablon lub \texttt{std::variant}.
\end{enumerate}

\noindent Pytania, które warto umieć rozstrzygnąć samodzielnie:
\begin{enumerate}
\item Czy \texttt{virtual} oznacza, że \emph{każde} wywołanie metody
  danej klasy musi być realizowane pośrednio?
\item Dlaczego \texttt{Base value = derived;} może dać inny wynik
  niż \texttt{Base\& ref = derived;}?
\item Co stanie się z funkcją wirtualną nadpisaną w \texttt{Derived},
  jeśli wywoła ją konstruktor \texttt{Base}?
\item W jakiej sytuacji \texttt{dynamic\_cast<Derived\&>} zgłasza
  \texttt{std::bad\_cast}?
\item Kiedy użyć \texttt{std::variant}, a kiedy interfejsu wirtualnego?
\end{enumerate}

\section{Wskazówki dotyczące diagnostyki i pomiarów}

Przy analizie klas polimorficznych przydatne są ostrzeżenia
\texttt{-Wall -Wextra -Wpedantic -Woverloaded-virtual} i, zależnie
od kompilatora, \texttt{-Wsuggest-override}. Błędy czasu życia
mogą być wykrywane przez AddressSanitizer i UBSan (w szczególności
odpowiednie kontrole vptr, jeśli są wspierane).

\begin{lstlisting}[language=bash,caption={Kompilacja przykładowego programu w GCC lub Clangu.}]
g++ -std=c++20 -O2 -Wall -Wextra -Wpedantic \
    -Woverloaded-virtual example.cpp -o example
./example

# Testy diagnostyczne (zalezne od platformy):
g++ -std=c++20 -O1 -g -fsanitize=address,undefined \
    -fno-omit-frame-pointer example.cpp -o example_san
./example_san
\end{lstlisting}

Należy porównywać z \emph{realnym} obciążeniem, a nie zakładać,
że jeden mechanizm jest zawsze szybszy od drugiego. Przy pomiarach
uwzględniamy optymalizację, lokalność danych, rozmiar próby,
rozrzut pomiarów i to, czy kompilator nie wyeliminował pracy.

\section{Podsumowanie rozdziału}

Kluczowe jest rozdzielenie trzech kwestii: \emph{jakie typy obiektów}
występują, \emph{kto odpowiada za ich czas życia} oraz
\emph{kiedy wybierana jest implementacja operacji}.
\texttt{virtual} umożliwia polimorfizm dynamiczny, lecz wymaga
rozważenia destruktorów i własności obiektów. Szablony i CRTP
oferują polimorfizm statyczny, a \texttt{std::variant} jest
wygodną alternatywą dla zamkniętej rodziny typów.
Praktyczna wydajność zależy nie tylko od kosztu wywołania,
ale również od reprezentacji i rozmieszczenia danych.

\section*{Literatura i dalsze źródła}
\addcontentsline{toc}{section}{Literatura do rozdziału 12}
\begin{itemize}
\item ISO C++ Standard: model obiektowy, klasy, dziedziczenie,
      funkcje wirtualne i RTTI (normatywne źródło reguł języka).
\item \url{https://en.cppreference.com/w/cpp/language/virtual}
      --- funkcje wirtualne i reguły nadpisywania.
\item \url{https://en.cppreference.com/w/cpp/language/dynamic_cast}
      --- sprawdzane konwersje w hierarchiach.
\item \url{https://en.cppreference.com/w/cpp/language/derived_class}
      --- dziedziczenie i wirtualne klasy bazowe.
\item Bjarne Stroustrup, \emph{The C++ Programming Language},
      wydanie czwarte, Addison-Wesley, 2013.
\end{itemize}

% ROZDZIAŁ 2 -- fragment do dołączenia do głównego skryptu.
% W preambule dokumentu umieść:
% \usepackage{listings}
% \lstset{
%   language=C++,
%   basicstyle=\ttfamily\small,
%   columns=fullflexible,
%   keepspaces=true,
%   breaklines=true,
%   showstringspaces=false,
%   frame=single,
%   tabsize=4
% }
% Przykładowe programy zakładają C++20, o ile nie podano inaczej.

\chapter{Specyfikatory i kwalifikatory deklaracji w C++}
\label{chap:specyfikatory-cpp}

\section{Wprowadzenie: czym właściwie są \texttt{static}, \texttt{const} i \texttt{inline}?}

Zanim zapamiętamy słowa kluczowe, warto zadać trzy pytania:
czy wartość wolno zmienić, jak długo istnieje obiekt i czy inne pliki
programu widzą ten sam obiekt. \texttt{const} odpowiada głównie na
pierwsze pytanie. \texttt{static} może wpływać na czas życia albo
widoczność nazwy, zależnie od miejsca deklaracji. \texttt{inline}
pozwala umieścić odpowiednią definicję w nagłówku dołączanym przez
kilka plików \texttt{.cpp}.

Słowa kluczowe \texttt{static}, \texttt{extern}, \texttt{inline}, \texttt{const},
\texttt{constexpr} i podobne nie są jednym rodzajem ``modyfikatorów''.
Mogą określać różne, w znacznej mierze niezależne własności deklaracji:
\begin{itemize}
    \item \emph{czas przechowywania} obiektu (\emph{storage duration}), czyli
    jak długo istnieje jego pamięć;
    \item \emph{łączność} nazwy (\emph{linkage}), czyli czy ta sama nazwa
    może oznaczać ten sam byt w różnych jednostkach translacji;
    \item \emph{modyfikowalność} obiektu lub dostępu przez dany typ;
    \item \emph{możliwość obliczenia wyrażenia podczas kompilacji};
    \item \emph{reguły wielokrotnego definiowania} funkcji lub zmiennej;
    \item \emph{własności funkcji i metod}, np. zakaz niejawnej konwersji
    albo gwarancje dotyczące wyjątków.
\end{itemize}

\paragraph{Jednostka translacji.}
W uproszczeniu plik \texttt{.cpp} po rozwinięciu dyrektyw
\texttt{\#include} tworzy jednostkę translacji
(\emph{translation unit}, TU). Plik nagłówkowy bywa włączany do wielu
jednostek translacji, dlatego umieszczanie w nim definicji wymaga uwagi.
Praktycznie: jedna TU to plik \texttt{.cpp} razem z treścią nagłówków,
którą kompilator widzi podczas jego kompilowania.

\paragraph{Czas przechowywania.}
Rozróżniamy w szczególności:
\begin{itemize}
    \item \emph{automatyczny} --- zwykle zmienne lokalne, tworzone przy
    wejściu w odpowiedni zakres i niszczone przy jego opuszczaniu;
    \item \emph{statyczny} --- obiekt istnieje przez czas działania programu;
    \item \emph{wątkowy} --- oddzielny obiekt istnieje dla każdego wątku;
    \item \emph{dynamiczny} --- czas życia kontrolowany jawnie lub przez
    obiekty zarządzające zasobem (np. \texttt{std::unique\_ptr}).
\end{itemize}

\paragraph{Łączność.}
Nazwa może nie mieć łączności (typowo zmienna lokalna), mieć
\emph{łączność wewnętrzną} (\emph{internal linkage}: dostępna jako ten
sam byt tylko w obrębie TU) albo \emph{łączność zewnętrzną}
(\emph{external linkage}: może odnosić się do jednego bytu w wielu TU).
C++20 wprowadza również łączność modułową; w tym rozdziale zajmujemy się
głównie klasycznym modelem plików nagłówkowych i źródłowych.

\section{\texttt{const}: zakaz modyfikacji przez dany dostęp}

\begin{center}\small
\begin{tabularx}{\textwidth}{@{}p{0.30\textwidth}X@{}}
\toprule
Specyfikator & Znaczenie \\
\midrule
\code{const} & Zabrania modyfikacji przez dany dostęp; wartość nie musi być znana podczas kompilacji. \\
\code{constexpr} & Pozwala użyć wartości lub funkcji w obliczeniu stałym, gdy spełnione są jego warunki. \\
\code{consteval} & Wymaga obliczenia wywołania funkcji w czasie kompilacji. \\
\code{constinit} & Wymaga statycznej inicjalizacji zmiennej o statycznym lub wątkowym czasie przechowywania; nie czyni jej stałą. \\
\code{mutable} & Pozwala zmieniać wskazaną składową także przez stały obiekt. \\
\bottomrule
\end{tabularx}
\end{center}

Deklaracja \texttt{const} kwalifikuje typ. Dla obiektu oznacza, że
nie wolno zmieniać jego wartości przez dostęp, w którym typ jest
\texttt{const} (z uwzględnieniem szczególnych zasad np. \texttt{mutable}).
Samo \texttt{const} \emph{nie gwarantuje} obliczenia wartości w czasie kompilacji.

\begin{lstlisting}
int read_value();

void example() {
    const int a = 10;
    // a = 11;              // error: a is const

    const int b = read_value();  // OK: runtime initialization
    int c = read_value();
    c += b;                      // OK: c is mutable
}
\end{lstlisting}

\subsection{Wskaźniki, referencje i pozycja \texttt{const}}

Różnica między kwalifikacją wskaźnika i wskazywanego obiektu jest
fundamentalna:

\begin{lstlisting}
int a = 1;
int b = 2;

const int* p = &a;  // pointer to const int
p = &b;             // OK: pointer can change
// *p = 3;          // error: pointee is const through p

int* const q = &a;  // const pointer to int
*q = 3;             // OK: a can change
// q = &b;          // error: pointer is const

const int* const r = &a;  // both are const through r
const int& ref = a;      // read-only access through ref
\end{lstlisting}

\paragraph{Argumenty i metody.}
Dla dużych obiektów często używamy \texttt{const T\&}, aby uniknąć
kopiowania i zadeklarować brak modyfikacji przez referencję.
Kwalifikator \texttt{const} za listą parametrów metody odnosi się
do obiektu \texttt{*this}:

\begin{lstlisting}
#include <string>

class Person {
    std::string name_;
public:
    explicit Person(std::string name) : name_(name) {}

    const std::string& name() const { return name_; }
    void rename(std::string s) { name_ = s; }
};
\end{lstlisting}

\paragraph{Kiedy stosować?}
Używaj \texttt{const} dla wartości, wskaźników, referencji i metod,
jeśli modyfikacja przez dany interfejs nie powinna być dozwolona.
Nie używaj go jako zamiennika \texttt{constexpr}: zwykła stała
zainicjalizowana wynikiem funkcji może być wyznaczona dopiero podczas
wykonywania programu.

\section{\texttt{constexpr}: stałe i funkcje dostępne dla obliczeń stałych}

\texttt{constexpr} dla zmiennej wymaga inicjalizacji nadającej się
do obliczenia jako wyrażenie stałe; taka zmienna jest kwalifikowana
\texttt{const}. Dla funkcji \texttt{constexpr} oznacza \emph{możliwość}
wywołania w kontekście wymagającym wyniku obliczanego w czasie kompilacji,
a nie obowiązek obliczenia każdego wywołania podczas kompilacji.

\begin{lstlisting}
#include <array>

constexpr int square(int x) {
    return x * x;
}

int read_value();

void example() {
    constexpr int n = square(5);    // constant expression
    std::array<int, n> data{};     // valid array extent

    int x = read_value();
    int y = square(x);             // runtime call is valid

    // constexpr int z = square(x);  // error: x not constant
    (void)data;
    (void)y;
}
\end{lstlisting}

\paragraph{Ważne rozróżnienie.}
Zwykle \texttt{const int n = 5;} także może uczestniczyć w wyrażeniach
stałych, bo spełnia odpowiednie reguły języka; nie oznacza to jednak,
że \emph{każda} zmienna \texttt{const} jest stała czasu kompilacji.
\texttt{constexpr} wyraża ten zamiar jednoznacznie i jest sprawdzany
przez kompilator.

\paragraph{Kiedy stosować?}
Używaj \texttt{constexpr}, gdy wartość ma być wykorzystywana jako stała
czasu kompilacji, np. do rozmiarów tablic \texttt{std::array} i argumentów
szablonów. Nadaje się także do \texttt{static\_assert}. Oznaczaj nim również
funkcje, które naturalnie nadają się zarówno do obliczeń stałych,
jak i zwykłych wywołań.

\subsection{\texttt{if constexpr} (od C++17)}

To osobna konstrukcja sterująca, typowo stosowana w szablonach.
W odróżnieniu od zwykłego \texttt{if} może odrzucić gałąź podczas
instancjacji szablonu:

\begin{lstlisting}
#include <type_traits>

template<class T>
const char* category() {
    if constexpr (std::is_integral_v<T>) {
        return "integral";
    } else {
        return "non-integral";
    }
}
\end{lstlisting}

\section{\texttt{consteval} i \texttt{constinit} (od C++20)}

\subsection{\texttt{consteval}: wywołanie wymagające stałej ewaluacji}

Funkcja \texttt{consteval} wymaga wyniku obliczanego podczas kompilacji.
Dlatego można przekazać jej stałą, ale zwykle nie wartość wczytaną
dopiero podczas działania programu. Dokładne reguły języka obejmują
również wywołania wewnątrz innych funkcji natychmiastowych.

\begin{lstlisting}
consteval int cube(int x) {
    return x * x * x;
}

int read_value();

void example() {
    constexpr int a = cube(3);  // OK
    int b = cube(4);            // OK: call evaluated as constant
    int x = read_value();
    // int c = cube(x);         // error: x not constant
    (void)b;
    (void)x;
}
\end{lstlisting}

\texttt{consteval} wybieraj wtedy, kiedy wywołanie \emph{musi}
być obliczone w kontekście stałym; nie jest to ogólny przepis
na szybsze wykonywanie programu.

\subsection{\texttt{constinit}: gwarancja inicjalizacji statycznej}

\texttt{constinit} dotyczy zmiennych o statycznym lub wątkowym czasie
przechowywania. Wymusza inicjalizację statyczną, a więc wyklucza
inicjalizację dynamiczną. \emph{Nie czyni zmiennej stałą}:

\begin{lstlisting}
constexpr int twice(int x) { return 2 * x; }

constinit int global_count = twice(10);

void update() {
    global_count += 1;        // OK: mutable variable
}

int read_value();
// constinit int bad = read_value();  // error: dynamic init
\end{lstlisting}

Nie wolno łączyć \texttt{constinit} i \texttt{constexpr} w tej samej
deklaracji. \texttt{constinit} ma sens np. dla mutowalnego stanu
globalnego, który musi być inicjalizowany bez inicjalizacji dynamicznej.
Nie rozwiązuje natomiast problemów synchronizacji dostępu wielu wątków.

\section{\texttt{static}: trzy różne zastosowania}

\begin{center}\small
\begin{tabularx}{\textwidth}{@{}p{0.30\textwidth}X@{}}
\toprule
Zapis & Skutek \\
\midrule
\code{static} lokalne & Daje zmiennej lokalnej czas życia do końca programu; inicjalizacja następuje przy pierwszym przejściu przez deklarację. \\
\code{static} składowa & Tworzy składową wspólną dla obiektów klasy. \\
\code{static} w pliku & Nadaje nazwie na poziomie przestrzeni nazw łączność wewnętrzną. \\
\code{extern} & Deklaruje nazwę o łączności zewnętrznej; definicja może być w innym pliku. \\
\code{inline} & Umożliwia zgodne z ODR definicje funkcji lub zmiennej w wielu jednostkach translacji; nie nakazuje optymalizacji wywołania. \\
\bottomrule
\end{tabularx}
\end{center}

Znaczenie \texttt{static} zależy od miejsca deklaracji. Jest to jedno
z najczęstszych źródeł nieporozumień.

\subsection{Nazwa w przestrzeni nazw: łączność wewnętrzna}

W przestrzeni nazw (w tym globalnej) \texttt{static} nadaje nazwie
łączność wewnętrzną. Definicja umieszczona w dwóch plikach
\texttt{.cpp} opisuje wtedy dwa różne byty.

\begin{lstlisting}
// file: first.cpp
static int cache_size = 100;
static int helper(int x) { return x + 1; }

// file: second.cpp
static int cache_size = 200;  // a different variable
\end{lstlisting}

We współczesnym C++ często preferuje się \emph{anonimową przestrzeń nazw}
dla lokalnych szczegółów implementacji:

\begin{lstlisting}
// implementation.cpp
namespace {
    int cache_size = 100;
    int helper(int x) { return x + 1; }
}
\end{lstlisting}

\paragraph{Czego unikać?}
Nie definiuj \texttt{static} zmiennej w nagłówku z zamiarem stworzenia
\emph{jednego} wspólnego obiektu. Każda TU otrzyma swoją instancję.
Jeżeli celem jest wspólny obiekt w nagłówku, rozważ \texttt{inline}
(od C++17), opisane dalej.

Jeśli \texttt{a.cpp} i \texttt{b.cpp} dołączają ten sam nagłówek,
\texttt{static int x = 0;} tworzy dwa oddzielne \texttt{x}.
\texttt{inline int x = 0;} tworzy jeden wspólny obiekt, a
\texttt{extern int x;} jedynie ogłasza wspólny obiekt, którego
jedna definicja musi znaleźć się w pliku \texttt{.cpp}.

\subsection{Zmienna lokalna: czas przechowywania statyczny}

Wewnątrz funkcji \texttt{static} nie nadaje nazwie łączności zewnętrznej
ani wewnętrznej: lokalna nazwa nadal nie ma łączności. Zmienia natomiast
czas przechowywania obiektu. Inicjalizacja lokalnej zmiennej statycznej,
jeśli nie została wykonana wcześniej, następuje przy pierwszym przejściu
sterowania przez jej deklarację.

\begin{lstlisting}
int next_id() {
    static int counter = 0;
    return ++counter;
}

// next_id() returns 1, then 2, then 3, ...
\end{lstlisting}

Od C++11 jednorazowa inicjalizacja lokalnej zmiennej statycznej jest
bezpieczna wobec współbieżności: tylko jeden wątek przeprowadza
inicjalizację, a pozostałe odpowiednio czekają. \emph{Nie wynika z tego},
że późniejsze modyfikacje są bezpieczne. Przy wywołaniach równoległych
powyższe \texttt{++counter} może powodować wyścig danych.

\begin{lstlisting}
#include <atomic>

int next_id_thread_safe() {
    static std::atomic<int> counter{0};
    return counter.fetch_add(1, std::memory_order_relaxed) + 1;
}
\end{lstlisting}

\paragraph{Kiedy stosować?}
Lokalna \texttt{static} jest dobra dla stanu należącego do funkcji
lub obiektu tworzonego leniwie przy pierwszym użyciu. Unikaj jej,
jeśli miałaby ukrywać zależności, utrudniać testy lub tworzyć
współdzielony mutowalny stan bez synchronizacji.

\subsection{Składowa klasy: wspólna dla wszystkich obiektów}

Statyczne pola klasy należą do klasy, a nie do każdej instancji
oddzielnie. Statyczna metoda nie ma wskaźnika \texttt{this}; dlatego
nie może bezpośrednio odwoływać się do niestatycznych pól.

\begin{lstlisting}
#include <cstddef>

struct Config {
    inline static int attempts = 0;          // C++17
    static constexpr std::size_t limit = 256; // inline since C++17

    static void record_attempt() {
        ++attempts;
    }
};

void example() {
    Config::record_attempt();
    auto n = Config::limit;
    (void)n;
}
\end{lstlisting}

Przed C++17 zwykłe statyczne pole niestałe deklarowano w klasie,
a definiowano raz poza nią:

\begin{lstlisting}
// counter.hpp
struct Counter {
    static int value;
};

// counter.cpp
#include "counter.hpp"
int Counter::value = 0;
\end{lstlisting}

Od C++17 \texttt{inline static} pozwala na definicję z inicjalizacją
bezpośrednio w klasie. Pole \texttt{static constexpr} jest od C++17
niejawnie \texttt{inline} i nie wymaga osobnej definicji poza klasą.
W starszych standardach dodatkowa definicja bywała potrzebna przy
\emph{odr-use}, np. przy pobieraniu adresu.

\section{\texttt{extern}: deklaracja bytu zdefiniowanego gdzie indziej}

Typowy wzorzec to deklaracja w pliku nagłówkowym oraz \emph{jedna}
definicja w pliku źródłowym. \texttt{extern} w takiej deklaracji
wskazuje, że nazwa może odnosić się do bytu z łącznością zewnętrzną;
nie alokuje nowego obiektu.

\begin{lstlisting}
// config.hpp
#pragma once
extern int connection_count;
extern const int max_connections;

// config.cpp
#include "config.hpp"
int connection_count = 0;
const int max_connections = 64;

// main.cpp
#include "config.hpp"
void connect() {
    if (connection_count < max_connections)
        ++connection_count;
}
\end{lstlisting}

Ponieważ deklaracja w \texttt{config.cpp} zawiera wcześniejsze
\texttt{extern}, definicja tej zmiennej ma łączność zewnętrzną.
Natomiast samodzielna definicja zmiennej typu \texttt{const int}
na poziomie przestrzeni nazw ma domyślnie łączność wewnętrzną
(w klasycznym kodzie bez modułów).

\paragraph{Uwaga na definicje.}
Zapis \texttt{extern int x = 5;} \emph{jest definicją}, gdy występuje
na poziomie przestrzeni nazw: inicjalizator zmienia znaczenie deklaracji.
Dlatego nie umieszczamy takiego zapisu w zwykłym nagłówku dołączanym
do wielu TU.

\subsection{Dlaczego \texttt{extern const} nie zastąpi \texttt{constexpr}?}

\begin{lstlisting}
// constants.hpp
extern const int array_size;

// constants.cpp
#include "constants.hpp"
const int array_size = 16;

// client.cpp
#include "constants.hpp"
// std::array<int, array_size> a;  // invalid: not constant here
\end{lstlisting}

Inna TU nie widzi inicjalizatora, więc nie może wykorzystywać tej
zmiennej jako znanej stałej szablonowej. Jeśli wartość musi być
stała czasu kompilacji i wspólna w wielu TU, od C++17 zwykle
wybieramy \texttt{inline constexpr} w nagłówku.

\subsection{\texttt{extern "C"}: łączność językowa}

Jest to inne zastosowanie słowa \texttt{extern}: określa
\emph{łączność językową}, np. przy współpracy z biblioteką C.

\begin{lstlisting}
extern "C" int c_library_function(int);
\end{lstlisting}

Nie oznacza ono, że wywołanie funkcji ma semantykę języka C;
przede wszystkim dotyczy zgodności łączenia i typu funkcji.

\section{\texttt{inline}: dozwolone definicje w wielu TU}

Historycznie \texttt{inline} sugerowało wstawienie ciała funkcji
w miejsce wywołania, ale \emph{nie wymusza} takiej optymalizacji.
Najważniejsza współczesna rola polega na odpowiednich regułach
jednej definicji (\emph{One Definition Rule}, ODR).

\subsection{Funkcje \texttt{inline}}

Definicja funkcji \texttt{inline} może znaleźć się w wielu TU,
np. przez dołączenie nagłówka, jeśli spełnione są wymagania ODR,
w szczególności zgodność definicji. Dla funkcji \texttt{inline}
o łączności zewnętrznej wszystkie takie definicje odnoszą się
do jednej funkcji (w tym do tego samego adresu).

\begin{lstlisting}
// math.hpp
#pragma once
inline int add(int a, int b) {
    return a + b;
}
\end{lstlisting}

Gdybyśmy usunęli \texttt{inline} i dołączyli nagłówek do dwóch
\texttt{.cpp}, mielibyśmy niedozwolone wielokrotne definicje zwykłej
funkcji o łączności zewnętrznej.

Funkcje zdefiniowane w obrębie definicji klasy są zazwyczaj niejawnie
\texttt{inline} (z zastrzeżeniem szczególnych reguł modułów C++20).
Funkcje \texttt{constexpr} i \texttt{consteval} też są niejawnie
\texttt{inline} przy pierwszej deklaracji.

\subsection{Zmienne \texttt{inline} (od C++17)}

Zmienna \texttt{inline} zdefiniowana w nagłówku może oznaczać
\emph{jeden wspólny obiekt} w wielu TU:

\begin{lstlisting}
// state.hpp
#pragma once
#include <atomic>

inline std::atomic<int> requests{0};
inline constexpr int buffer_size = 4096;

inline int register_request() {
    return requests.fetch_add(1, std::memory_order_relaxed) + 1;
}
\end{lstlisting}

To praktyczna podstawa bibliotek \emph{header-only}. Na poziomie
przestrzeni nazw zmienna \texttt{inline constexpr} ma zazwyczaj
łączność zewnętrzną. Sama zmienna \texttt{constexpr} ma zaś
domyślnie łączność wewnętrzną.

\paragraph{Ważne: \texttt{static inline} to nie to samo.}
Na poziomie przestrzeni nazw \texttt{static inline} \texttt{int x = 0;}
ma łączność wewnętrzną: każda TU posiada własne \texttt{x}.
Dodanie \texttt{inline} nie zmienia tego w jeden globalny obiekt.

\subsection{Reguła ODR w praktyce}

\begin{itemize}
    \item Zwykła funkcja lub zmienna o łączności zewnętrznej:
    zasadniczo jedna definicja w całym programie.
    \item Funkcja \texttt{inline}, zmienna \texttt{inline}, a także
    odpowiednio zdefiniowane szablony: dozwolone definicje w wielu TU,
    z dodatkowymi wymaganiami co do ich identyczności.
    \item Sama deklaracja w nagłówku jest zwykle bezpieczna;
    problem dotyczy definicji.
\end{itemize}

\section{Łączenie specyfikatorów: najważniejsze przypadki}

Poniższe zapisy nie są wzajemnie zamienne:

\begin{lstlisting}
const int a = 7;                // const; namespace: internal linkage
constexpr int b = 7;            // compile-time const; internal linkage
inline constexpr int c = 7;     // C++17: shared across TUs
static int d = 7;               // namespace: internal linkage
static constexpr int e = 7;     // compile-time; internal linkage
inline int f = 7;               // C++17: one shared object
extern int g;                   // declaration, definition elsewhere
extern const int h;             // declaration, definition elsewhere
constinit int i = 7;            // C++20: static init, mutable
static constinit int j = 7;     // C++20: internal linkage
\end{lstlisting}

\noindent W powyższym przykładzie zakładamy, że wszystkie deklaracje
znajdują się na poziomie przestrzeni nazw, poza modułami nazwanymi.
W rzeczywistym projekcie \texttt{g} i \texttt{h} muszą mieć definicje,
jeśli są odpowiednio używane; nie należy też kopiować wielu definicji
\texttt{i} do nagłówków.

\begin{center}
\small
\renewcommand{\arraystretch}{1.22}
\begin{tabular}{|p{0.24\textwidth}|p{0.31\textwidth}|p{0.34\textwidth}|}
\hline
\textbf{Konstrukcja} & \textbf{Najważniejsza cecha} & \textbf{Typowe zastosowanie} \\
\hline
\texttt{const} & Brak modyfikacji przez dany dostęp & Parametry, referencje, stałe wartości \\
\hline
\texttt{constexpr} & Możliwość użycia wyniku jako wyrażenia stałego & Parametry szablonów, rozmiary, funkcje obliczeniowe \\
\hline
\texttt{consteval} & Wymagana stała ewaluacja wywołania & Walidacja i generowanie stałych przy kompilacji \\
\hline
\texttt{constinit} & Zakaz inicjalizacji dynamicznej & Mutowalny stan o inicjalizacji statycznej \\
\hline
\texttt{static} lokalne & Jeden obiekt o statycznym czasie przechowywania & Leniwy stan funkcji \\
\hline
\texttt{static} globalne & Łączność wewnętrzna & Szczegóły implementacji w \texttt{.cpp} \\
\hline
\texttt{static} w klasie & Składowa wspólna dla instancji & Liczniki, konfiguracja klasy \\
\hline
\texttt{extern} & Deklaracja nazwy odnoszącej się do innej definicji & Podział \texttt{.hpp}/\texttt{.cpp} \\
\hline
\texttt{inline} & Szczególne reguły wielokrotnych definicji & Funkcje i zmienne w nagłówkach \\
\hline
\end{tabular}
\end{center}

\section{\texttt{thread\_local}: osobny obiekt dla każdego wątku}

\texttt{thread\_local} oznacza wątkowy czas przechowywania. Każdy
wątek posiada własną instancję obiektu, więc modyfikacje tej instancji
nie są współdzielone z innymi wątkami.

\begin{lstlisting}
int next_local_id() {
    thread_local int id = 0;
    return ++id;
}
\end{lstlisting}

Różnica względem \texttt{static int id = 0;} jest istotna:
\texttt{static} daje jedną instancję współdzieloną przez wątki,
\texttt{thread\_local} --- po jednej na wątek.

Na poziomie przestrzeni nazw można użyć \texttt{static thread\_local}
(dla łączności wewnętrznej) albo \texttt{extern thread\_local}
(w deklaracji obiektu zdefiniowanego w innym TU).

\paragraph{Kiedy stosować?}
Dla stanu rzeczywiście związanego z wątkiem, np. lokalnego generatora
lub bufora roboczego. Nie jest to ogólny zamiennik blokad i nie nadaje
się do danych, które muszą być wspólne między wątkami.

\section{\texttt{mutable} i \texttt{volatile}: dwa inne rodzaje kwalifikacji}

\subsection{\texttt{mutable}: kontrolowana modyfikacja wewnątrz obiektu stałego}

\texttt{mutable} dotyczy niestatycznych pól klasy. Pozwala modyfikować
takie pole nawet w metodzie \texttt{const}. Typowy przypadek to blokada,
która jest szczegółem technicznym i nie stanowi logicznego stanu obiektu:

\begin{lstlisting}
#include <mutex>

class SafeCounter {
    mutable std::mutex mutex_;
    int value_ = 0;

public:
    void increment() {
        std::lock_guard<std::mutex> lock(mutex_);
        ++value_;
    }

    int value() const {
        std::lock_guard<std::mutex> lock(mutex_);
        return value_;
    }
};
\end{lstlisting}

Nie stosuj \texttt{mutable} tylko po to, aby obchodzić prawidłowo
zaprojektowany interfejs \texttt{const}. Sam \texttt{mutable}
nie zapewnia bezpieczeństwa wątkowego.

\subsection{\texttt{volatile}: nie jest mechanizmem synchronizacji}

\texttt{volatile} służy do określonych sytuacji, w których dostępy
do obiektu mają być obserwowalnymi operacjami (np. w pewnych
zastosowaniach niskopoziomowych). Nie zapewnia atomowości,
porządkowania pamięci ani braku wyścigów danych w C++.

\begin{lstlisting}
// Not a thread-safe counter:
volatile int wrong_counter = 0;

// For concurrent increments, use e.g.:
#include <atomic>
std::atomic<int> correct_counter{0};
\end{lstlisting}

Dla synchronizacji wątków stosuj \texttt{std::atomic}, muteksy i inne
narzędzia współbieżności, a nie \texttt{volatile}.

\section{\texttt{auto} i \texttt{decltype(auto)}: dedukcja typu}

We współczesnym C++ \texttt{auto} w deklaracji zmiennej z inicjalizatorem
włącza dedukcję typu. Dedukcja przez wartość odrzuca
najwyższy poziom \texttt{const} oraz referencyjność wyrażenia.

Natomiast \texttt{decltype(auto)} zachowuje semantykę \texttt{decltype}.

\begin{lstlisting}
int& shared_value() {
    static int n = 10;
    return n;
}

void example() {
    const int k = 4;
    auto a = k;             // int (copy)
    const auto b = k;       // const int

    auto x = shared_value();           // int: kopia
    auto& original = shared_value();   // int&: ten sam obiekt
    x = 15;                            // nie zmienia shared_value()
    original = 20;                     // zmienia shared_value()
    decltype(auto) y = shared_value(); // int&: także referencja
    y = 21;
    (void)a;
    (void)b;
    (void)x;
}
\end{lstlisting}

Dedukcja typu z \texttt{auto} poprawia czytelność, zwłaszcza
przy iteratorach i rozbudowanych typach. Nie należy jednak
niezauważenie tracić referencji, gdy kod wymaga operowania
na oryginalnym obiekcie. Wtedy używaj np. \texttt{auto\&} lub
\texttt{const auto\&}.

\section{Powiązane specyfikatory metod i funkcji}

Nie są to specyfikatory czasu przechowywania, lecz często występują
obok poprzednich słów kluczowych:

\begin{lstlisting}
#include <cstddef>

struct Base {
    virtual ~Base() = default;
    virtual int evaluate() const noexcept { return 1; }
};

struct Derived final : Base {
    int evaluate() const noexcept override { return 2; }
};

struct Index {
    explicit Index(std::size_t x) : value(x) {}
    std::size_t value;
};

[[nodiscard]] int compute_result() { return 42; }
\end{lstlisting}

\begin{itemize}
    \item \texttt{explicit} kontroluje niejawne konwersje
    przez konstruktory i operatory konwersji (szczegóły poniżej);
    \item \texttt{virtual} umożliwia dynamiczny polimorfizm,
    \texttt{override} pozwala sprawdzić, czy metoda rzeczywiście
    nadpisuje metodę bazową, a \texttt{final} zabrania dalszego
    nadpisania metody lub dziedziczenia po klasie;
    \item \texttt{noexcept} deklaruje funkcję nierzucającą wyjątków
    (naruszenie tej gwarancji prowadzi do \texttt{std::terminate});
    \item \texttt{[[nodiscard]]} pozwala kompilatorowi ostrzegać,
    gdy wynik został zignorowany.

    Atrybut \texttt{[[maybe\_unused]]} pozwala ograniczyć
    ostrzeżenia o nieużywanych bytach.
\end{itemize}

\subsection{\texttt{explicit}: kontrola niejawnych konwersji}

Konstruktor bez \texttt{explicit} może zostać użyty przez kompilator do
niejawnej konwersji, jeżeli pozostałe reguły konwersji na to pozwalają.
\texttt{explicit} wymaga, by w takim przypadku zamiar utworzenia obiektu
był zapisany jawnie. Dotyczy to nie tylko konstruktorów z jednym
parametrem: także konstruktor z kilkoma parametrami może uczestniczyć
w inicjalizacji przez listę.

\begin{center}\small
\begin{tabularx}{\textwidth}{@{}p{0.37\textwidth}X@{}}
\toprule
Zapis & Skutek dla konstruktora \code{explicit Meter(int)} \\
\midrule
\code{Meter a(3)} & Inicjalizacja bezpośrednia: dozwolona. \\
\code{Meter a{3}} & Inicjalizacja bezpośrednia przez listę: dozwolona. \\
\code{Meter a = 3} & Inicjalizacja kopiująca wymagałaby niejawnej konwersji: błąd kompilacji. \\
\code{Meter a = {3}} & Inicjalizacja kopiująca przez listę nie może zakończyć się wyborem konstruktora \code{explicit}: błąd kompilacji. \\
\code{use(Meter{3})} & Jawne utworzenie argumentu: dozwolone. Samo \code{use(3)} byłoby niedozwolone. \\
\bottomrule
\end{tabularx}
\end{center}

\begin{lstlisting}[caption={Konstruktor i operator konwersji explicit (kompletny program, C++20).}]
#include <iostream>

struct Meter {
    explicit Meter(int n) : value(n) {}
    int value;
};

void use(Meter m) {
    std::cout << m.value << '\n';
}

struct Ready {
    bool value;
    explicit operator bool() const { return value; }
};

int main() {
    Meter first(3);
    Meter second{4};
    use(first);
    use(Meter{5});
    // Meter bad = 3;   // blad: niejawna konwersja
    // Meter bad2 = {3}; // blad: copy-list-initialization
    // use(6);           // blad: niejawna konwersja argumentu

    Ready ready{true};
    if (ready) std::cout << "ready\n"; // kontekst wymaga bool
    bool flag = static_cast<bool>(ready); // jawna konwersja
    // bool bad = ready; // blad: operator jest explicit
    (void)second;
    (void)flag;
}
\end{lstlisting}

\texttt{explicit operator bool()} nadal działa w warunkach takich jak
\texttt{if}, \texttt{while} czy operand \texttt{!}: język dopuszcza tam
\emph{kontekstową konwersję do bool}. Nie pozwala natomiast na zwykłe
\texttt{bool flag = ready}. Słowo \texttt{explicit} nie zabrania
inicjalizacji bezpośredniej ani jawnego \texttt{static\_cast}.

Od C++20 można uzależnić tę własność od warunku obliczanego podczas
kompilacji: \texttt{explicit(condition)}. Prawdziwy warunek działa jak
\texttt{explicit}, a fałszywy jak jego brak:

\begin{lstlisting}[caption={Warunkowe explicit od C++20 (kompletny program).}]
template<bool AllowImplicit>
struct Token {
    explicit(!AllowImplicit) Token(int n) : value(n) {}
    int value;
};

int main() {
    Token<true> a = 3;  // dozwolona konwersja
    Token<false> b{3};  // jawna inicjalizacja
    // Token<false> c = 3; // blad: wymaga jawnej inicjalizacji
    return a.value + b.value == 6 ? 0 : 1;
}
\end{lstlisting}

\texttt{explicit} jest szczególnie przydatne dla typów jednostek,
identyfikatorów i uchwytów, w których przypadkowe utworzenie obiektu
z liczby lub napisu mogłoby ukryć błąd wywołania.

\section{Praktyczna organizacja plików: przykład całościowy}

Rozważmy licznik wspólny dla programu, stała kompilacyjna
oraz funkcję dostępną z kilku TU:

\begin{lstlisting}
// statistics.hpp
#pragma once
#include <cstddef>

inline constexpr std::size_t block_size = 1024;
extern unsigned total_jobs;

inline unsigned completed_jobs() {
    return total_jobs;
}
\end{lstlisting}

\begin{lstlisting}
// statistics.cpp
#include "statistics.hpp"

unsigned total_jobs = 0;
\end{lstlisting}

\begin{lstlisting}
// worker.cpp
#include "statistics.hpp"

void do_one_job() {
    ++total_jobs;
}
\end{lstlisting}

\begin{lstlisting}
// main.cpp
#include "statistics.hpp"
#include <iostream>

void do_one_job();

int main() {
    do_one_job();
    std::cout << block_size << " " << completed_jobs() << '\n';
}
\end{lstlisting}

Kompilacja przykładowa:
\begin{lstlisting}
g++ -std=c++20 -Wall -Wextra -pedantic \
    main.cpp worker.cpp statistics.cpp -o program
\end{lstlisting}

Mamy tu jedną definicję \texttt{total\_jobs} w \texttt{statistics.cpp},
jedną stałą \texttt{block\_size} udostępnioną przez
\texttt{inline constexpr}, a także legalną definicję funkcji
\texttt{completed\_jobs} w nagłówku. Przykład zakłada jeden wątek;
przy równoczesnych zapisach do licznika trzeba użyć synchronizacji,
którą omawia rozdział~\ref{chap:concurrency}.

\section{Najczęstsze błędy i reguły wyboru}

\subsection{Błędy}

\begin{enumerate}
    \item \textbf{Definicja zwykłej zmiennej w nagłówku:}
    zapis \texttt{int counter = 0;} powoduje wielokrotne definicje,
    jeśli nagłówek włączymy do kilku TU. Użyj \texttt{extern}
    i jednej definicji w \texttt{.cpp} albo \texttt{inline} (C++17).
    \item \textbf{\texttt{static} w nagłówku jako wspólny stan:}
    w rzeczywistości tworzy osobny stan w każdej TU.
    \item \textbf{Utożsamianie \texttt{const} z \texttt{constexpr}:}
    obiekt \texttt{const} może być inicjalizowany w runtime.
    \item \textbf{Utożsamianie \texttt{constexpr} z \texttt{consteval}:}
    zwykłą funkcję \texttt{constexpr} wolno wywoływać w runtime.
    \item \textbf{Utożsamianie \texttt{inline} z optymalizacją:}
    kompilator może wstawić lub nie wstawić ciała dowolnej funkcji,
    zgodnie z dozwolonymi optymalizacjami.
    \item \textbf{Utożsamianie \texttt{static} z synchronizacją:}
    bezpieczna inicjalizacja lokalnego statyka nie zabezpiecza
    późniejszych zapisów.
    \item \textbf{Używanie \texttt{volatile} zamiast \texttt{atomic}:}
    \texttt{volatile} nie eliminuje wyścigów danych.
    \item \textbf{Nadmierne korzystanie z globalnego mutowalnego stanu:}
    nawet poprawny pod względem ODR program może być trudny
    do testowania oraz podatny na problemy kolejności inicjalizacji.
\end{enumerate}

\subsection{Kolejność inicjalizacji obiektów globalnych}

Inicjalizacja dynamiczna obiektów o statycznym czasie przechowywania
w różnych TU może mieć nieoczywistą kolejność. Jest to źródło
\emph{static initialization order fiasco}. Jedną z praktycznych
technik jest przeniesienie zależności do lokalnego obiektu
statycznego, inicjalizowanego przy pierwszym wywołaniu:

\begin{lstlisting}
#include <string>

const std::string& application_name() {
    static const std::string name = "MyApplication";
    return name;
}
\end{lstlisting}

Gdy możliwe jest \emph{stałe inicjalizowanie}, preferuj proste
\texttt{constexpr} lub \texttt{constinit}. Pamiętaj jednak, że
\texttt{constinit} wyklucza inicjalizację dynamiczną, lecz
nie narzuca niemutowalności obiektu.

\subsection{Zalecenia w skrócie}

\begin{itemize}
    \item Potrzebujesz niemutowalnego dostępu? Użyj \texttt{const}.
    \item Potrzebujesz wyrażenia stałego? Użyj \texttt{constexpr}.
    \item Wywołanie musi być stałe? Rozważ \texttt{consteval}.
    \item Mutowalny obiekt ma wymagać inicjalizacji statycznej?
    Rozważ \texttt{constinit}.
    \item Chcesz ukryć szczegół implementacji w \texttt{.cpp}?
    Rozważ anonimową przestrzeń nazw.
    \item Potrzebujesz stanu zachowywanego między wywołaniami?
    Rozważ lokalną \texttt{static}, pamiętając o wątkach.
    \item Potrzebujesz jednego symbolu zdefiniowanego w \texttt{.cpp}?
    Deklaruj go w nagłówku przez \texttt{extern}.
    \item Potrzebujesz definicji w nagłówku współdzielonej przez TU?
    Użyj \texttt{inline}, jeśli jest odpowiednie; dla stałych
    kompilacyjnych \texttt{inline constexpr}.
    \item Potrzebujesz oddzielnego stanu wątku? Użyj \texttt{thread\_local}.
    \item Potrzebujesz komunikacji i synchronizacji między wątkami?
    Użyj atomików lub blokad, a nie \texttt{volatile}.
    \item Konstruktor lub operator konwersji mógłby uruchamiać
    niezamierzoną konwersję? Rozważ \texttt{explicit}.
\end{itemize}

\paragraph{Konkluzja.}
Najważniejsze pytanie nie brzmi ``które słowo kluczowe jest szybsze?'',
lecz: czy deklaracja ma określać \emph{czas życia}, \emph{łączność},
\emph{niemutowalność}, \emph{stałą ewaluację}, czy
\emph{dopuszczalność wielu definicji}? Dopiero wtedy wybór
specyfikatora jest jednoznaczny.

% ROZDZIAL 3 -- fragment do dolaczenia do glownego skryptu.
% Przykladowe programy zakladaja C++20, jesli nie podano inaczej.
% W preambule dokumentu potrzebny jest pakiet listings:
% \usepackage{listings}
% \lstset{language=C++,basicstyle=\ttfamily\small,
%   columns=fullflexible,keepspaces=true,breaklines=true,
%   showstringspaces=false,frame=single,tabsize=4}
% Fragment mozna wlaczyc przez:
% \input{rozdzial_3_wskazniki_smart_pointery_cpp.tex}

\chapter[Wskaźniki i inteligentne wskaźniki w C++]{Wskaźniki, zarządzanie pamięcią i inteligentne wskaźniki w C++}
\label{chap:wskazniki-cpp}

\section{Po co istnieją wskaźniki?}

Wskaźnik (ang. \emph{pointer}) przechowuje wartość pozwalającą
odwołać się do obiektu albo funkcji. Typowy wskaźnik do obiektu ma
postać \texttt{T*}. W odróżnieniu od obiektu typu \texttt{T},
wskaźnik może wskazywać na inny obiekt, może być pusty, a samo
wskazywanie \emph{nie oznacza posiadania wskazywanego obiektu}.

Wskaźniki wykorzystuje się między innymi do:
\begin{itemize}
    \item przekazywania i modyfikowania istniejących obiektów bez ich kopiowania;
    \item opisywania opcjonalnej relacji ``obiekt istnieje albo go nie ma'';
    \item implementowania dynamicznych struktur danych, takich jak listy,
    drzewa i grafy;
    \item realizowania polimorfizmu i dynamicznego wiązania;
    \item współpracy z interfejsami C, sterownikami i bibliotekami systemowymi;
    \item kontrolowania czasu życia zasobów za pomocą
    \texttt{std::unique\_ptr}, \texttt{std::shared\_ptr} i
    \texttt{std::weak\_ptr}.
\end{itemize}

\paragraph{Najważniejsza zasada.}
W nowoczesnym C++ wskaźnik \texttt{T*} traktujemy przede wszystkim
jako \emph{nieposiadający dostęp} do obiektu (\emph{observer}).
Odpowiedzialność za czas życia obiektu powinna być reprezentowana
osobno i jednoznacznie. Najczęściej wystarcza obiekt przechowywany
bezpośrednio, kontener standardowy albo \texttt{std::unique\_ptr}.

\section{Podstawy: adres, dereferencja i \texttt{nullptr}}

\begin{center}\small
\begin{tabularx}{\textwidth}{@{}p{0.32\textwidth}X@{}}
\toprule
Zapis & Za co odpowiada \\
\midrule
\texttt{\&x} & Pobiera adres obiektu \code{x}. \\
\code{*p} & Uzyskuje obiekt wskazywany przez \code{p}; wymaga ważnego, niepustego wskaźnika. \\
\code{p->field} & Udostępnia składową przez wskaźnik; odpowiada \code{(*p).field}. \\
\code{nullptr} & Oznacza brak wskazywanego obiektu. \\
\code{std::span<T>} & Łączy dostęp do kolejnych elementów z ich liczbą, ale nie posiada pamięci. \\
\code{std::addressof(x)} & Pobiera rzeczywisty adres obiektu nawet przy przeciążonym operatorze \texttt{\&}. \\
\bottomrule
\end{tabularx}
\end{center}

\subsection{Deklarowanie wskaźników}

Operator \texttt{\&} zastosowany do obiektu daje jego adres, natomiast
operator jednoargumentowy \texttt{*} służy do dereferencji (dostępu do
wskazywanego obiektu):

\begin{lstlisting}
#include <iostream>

int main() {
    int x = 10;
    int* p = &x;             // p points to x

    std::cout << x << '\n';  // 10
    std::cout << *p << '\n'; // 10

    *p = 25;                 // modifies x
    std::cout << x << '\n';  // 25
}
\end{lstlisting}

W deklaracji \texttt{int* p} gwiazdka należy do deklaratora, nie jest
operacją dereferencji. Nie należy też utożsamiać wskaźnika z samym
obiektem: \texttt{p} jest osobnym obiektem o własnej wartości i adresie.

\begin{lstlisting}
int x = 7;
int* p = &x;
int** pp = &p;  // pointer to a pointer

**pp = 9;       // changes x
\end{lstlisting}

Zapis \texttt{int* a, b;} deklaruje wskaźnik \texttt{a} oraz
zwykłą zmienną \texttt{b}. Dla czytelności lepiej deklarować je
oddzielnie.

\subsection{Wskaźnik pusty i poprawna dereferencja}

Pusty wskaźnik zapisujemy za pomocą \texttt{nullptr} (od C++11).
Nie wolno go dereferencjonować. Samo porównanie z
\texttt{nullptr} jest bezpieczne:

\begin{lstlisting}
int* p = nullptr;

if (p != nullptr) {
    *p = 5;
}

if (p) {                 // equivalent non-null test
    *p = 6;
}
\end{lstlisting}

\texttt{nullptr} ma osobny typ \texttt{std::nullptr\_t}; w odróżnieniu od
historycznych zapisów \texttt{0} i \texttt{NULL} nie jest zwykłą
liczbą całkowitą. Pomaga to przy przeciążaniu funkcji:

\begin{lstlisting}
void f(int);
void f(int*);

f(0);         // calls f(int)
f(nullptr);   // calls f(int*)
\end{lstlisting}

\paragraph{Uwaga o wskaźnikach niezainicjalizowanych.}
Lokalne \texttt{int* p;} bez inicjalizacji ma nieokreśloną wartość;
korzystanie z niej może prowadzić do niezdefiniowanego zachowania.
Zawsze inicjalizuj wskaźnik poprawnym adresem lub \texttt{nullptr}.

\subsection{Dostęp do składowych: \texttt{.} i \texttt{->}}

Operator \texttt{.} wybiera składową obiektu, natomiast
\texttt{->} wybiera składową przez wskaźnik:

\begin{lstlisting}
#include <string>

struct User {
    std::string name;
    int score = 0;
};

User u{"Anna", 12};
User* p = &u;

u.score = 13;
p->score = 14;       // same as (*p).score = 14
\end{lstlisting}

Nawiasy w \texttt{(*p).score} są potrzebne ze względu na priorytety
operatorów. Zapis \texttt{*p.score} nie oznacza tego samego.

\section{Wskaźniki, referencje i kwalifikator \texttt{const}}

\subsection{Kiedy wskaźnik, a kiedy referencja?}

Referencja \texttt{T\&} jest zwykle odpowiednia, gdy obiekt musi
istnieć i parametr nie jest opcjonalny. Wskaźnik \texttt{T*} nadaje
się, gdy brak obiektu jest dopuszczalny lub gdy chcemy zmieniać
obiekt, na który wskazujemy.

\begin{lstlisting}
void increase(int& x) { ++x; }

void increase_if_present(int* x) {
    if (x) { ++*x; }
}

int n = 3;
increase(n);             // n == 4
increase_if_present(&n); // n == 5
increase_if_present(nullptr); // no operation
\end{lstlisting}

Referencja również może stać się \emph{wisząca}, gdy obiekt,
do którego się odnosi, przestanie istnieć. Typ \texttt{T\&}
nie gwarantuje samodzielnie, że czas życia obiektu jest wystarczająco długi.

\subsection{Co jest stałe: wskaźnik czy obiekt?}

\begin{lstlisting}
int a = 1;
int b = 2;

const int* p = &a;       // read-only access through p
p = &b;                 // OK: p can point elsewhere
// *p = 3;              // error

int* const q = &a;       // q cannot be reassigned
*q = 3;                 // OK: a is mutable
// q = &b;              // error

const int* const r = &a; // neither reassignment nor write
int const* s = &a;       // same type as const int*
\end{lstlisting}

Reguła praktyczna: \texttt{const} po \texttt{*} dotyczy samego
wskaźnika, a przed \texttt{*} --- wskazywanego typu.

\paragraph{Pozorna niemutowalność.}
Wskaźnik \texttt{const T*} nie pozwala modyfikować obiektu
\emph{przez ten wskaźnik}. Nie oznacza, że obiekt nie może ulec zmianie
inną drogą:

\begin{lstlisting}
int x = 10;
const int* read_only = &x;
x = 20;                 // allowed
// *read_only = 30;     // not allowed
\end{lstlisting}

\section{Czas życia obiektu i wskaźniki wiszące}

\subsection{Adres nie oznacza prawa dostępu}

Wskaźnik może przechowywać adres miejsca, w którym obiekt już
nie istnieje. Jest wtedy wskaźnikiem \emph{wiszącym}
(ang. \emph{dangling pointer}). Dereferencja takiego wskaźnika
prowadzi do niezdefiniowanego zachowania (\emph{undefined behavior}, UB).

\begin{lstlisting}
int* bad() {
    int local = 42;
    return &local;       // WRONG: local dies at return
}

// int x = *bad();       // UB: dangling pointer
\end{lstlisting}

Prawidłowo można zwrócić samą wartość:

\begin{lstlisting}
int good() {
    int local = 42;
    return local;        // safe value return
}
\end{lstlisting}

\subsection{Unieważnianie wskaźników przez kontenery}

Nie każdy wskaźnik staje się wiszący przez wyjście z funkcji.
Relokacja elementów \texttt{std::vector} także może unieważnić
wskaźniki, referencje i iteratory do tych elementów:

\begin{lstlisting}
#include <vector>

std::vector<int> values{10, 20};
int* p = &values[0];

values.reserve(100);     // may reallocate storage
// *p = 50;              // unsafe if reallocation occurred
p = &values[0];          // reacquire a valid pointer
\end{lstlisting}

Po zmianie pojemności wymagającej realokacji adresy elementów
się zmieniają. W ogólności po operacji mogącej unieważnić wskaźnik
należy pobrać go ponownie, a nie próbować odgadnąć, czy nadal działa.

\subsection{Trzy niezależne pytania}

Przy każdym wskaźniku rozstrzygaj:
\begin{enumerate}
    \item Czy wartość wskaźnika jest poprawna (np. nie jest
    niezainicjalizowana)?
    \item Czy obiekt, na który wskazujemy, nadal istnieje?
    \item Czy operacja jest dozwolona dla danego typu i zakresu pamięci
    (np. czy nie wychodzi poza tablicę)?
\end{enumerate}

Sam test \texttt{p != nullptr} rozwiązuje co najwyżej fragment
pierwszego problemu, ale \emph{nie} wykrywa wiszących wskaźników.

\section{Tablice, arytmetyka wskaźników i \texttt{std::span}}

\subsection{Arytmetyka wskaźników}

Dla wskaźnika do elementu tablicy operacje \texttt{p + k} i
\texttt{p[k]} odnoszą się do jej elementów, pod warunkiem że
poruszamy się w dozwolonym zakresie. Można utworzyć wskaźnik
\emph{jeden za końcem} tablicy, ale nie wolno go dereferencjonować.

\begin{lstlisting}
int a[]{4, 8, 12};
int* p = a;          // points to a[0]

int x = *(p + 1);    // 8
int y = p[2];       // 12
int* end = a + 3;   // valid one-past-the-end pointer
// int z = *end;     // UB
\end{lstlisting}

Dodawanie i odejmowanie wskaźników ma sens w obrębie \emph{tej samej}
tablicy (z uwzględnieniem pozycji jeden za końcem). Nie należy
wykorzystywać takich działań do przemieszczania się między dowolnymi
obiektami w pamięci.

\subsection{Tablica to nie wskaźnik}

Typ \texttt{int[3]} i typ \texttt{int*} są różne, choć w wielu
wyrażeniach tablica jest automatycznie konwertowana na wskaźnik
do pierwszego elementu (tzw. \emph{array-to-pointer decay}).

\begin{lstlisting}
int a[3] = {1, 2, 3};
int* p = a;

static_assert(sizeof(a) == 3 * sizeof(int));
// sizeof(p) is the size of a pointer, not of the array
\end{lstlisting}

Funkcja \texttt{void f(int* p)} nie zna długości tablicy.
Jeśli trzeba pracować z ciągłym fragmentem pamięci o znanym
rozmiarze, C++20 udostępnia \texttt{std::span}:

\begin{lstlisting}
#include <span>
#include <vector>

int sum(std::span<const int> xs) {
    int result = 0;
    for (int x : xs) {
        result += x;
    }
    return result;
}

int a[]{1, 2, 3};
std::vector<int> v{4, 5, 6};

int s1 = sum(a);  // 6
int s2 = sum(v);  // 15
\end{lstlisting}

\texttt{std::span<T>} przechowuje widok ciągłych danych
(z informacją o liczbie elementów), ale \emph{nie jest właścicielem}
pamięci. Gdy dane znikną albo przeniosą się w inne miejsce,
\texttt{span} również może stać się niepoprawny. Dla większości
dynamicznych tablic własnościowych preferuj \texttt{std::vector<T>}.

\section{\texttt{void*}, konwersje i rzutowania wskaźników}

\subsection{\texttt{void*}: adres bez informacji o typie}

Wskaźnik \texttt{void*} może przechowywać adres obiektu bez
informacji o jego konkretnym typie. Bywa używany w interfejsach C
i bibliotekach niskopoziomowych, ale nie pozwala na bezpośrednią
dereferencję:

\begin{lstlisting}
int x = 42;
void* opaque = &x;              // type erased
// *opaque = 7;                 // error: cannot dereference void*

int* p = static_cast<int*>(opaque);
*p = 7;                        // valid: original object is int
\end{lstlisting}

Przywrócenie błędnego typu i późniejsza próba użycia obiektu
mogą prowadzić do niezdefiniowanego zachowania. Jeśli typ jest
znany, zwykle lepiej przechowywać od razu \texttt{T*}.

\subsection{\texttt{static\_cast}, \texttt{dynamic\_cast} i
\texttt{reinterpret\_cast}}

\begin{itemize}
    \item \texttt{static\_cast<T*>} nadaje się m.in. do poprawnych
    konwersji pomiędzy \texttt{void*} a wskaźnikiem obiektowym,
    jeśli rzeczywisty typ i wyrównanie na to pozwalają.
    \item \texttt{dynamic\_cast<T*>} służy do sprawdzanych w czasie
    wykonania konwersji w hierarchiach klas polimorficznych.
    Niepowodzenie konwersji wskaźnika zwraca \texttt{nullptr}.
    \item \texttt{reinterpret\_cast<T*>} umożliwia niskopoziomową
    reinterpretację reprezentacji wskaźnika, ale sama konwersja
    nie uprawnia do dowolnej dereferencji ani nie tworzy obiektu
    odpowiedniego typu.
\end{itemize}

\begin{lstlisting}
struct Base {
    virtual ~Base() = default;
};

struct Derived : Base {
    int value = 12;
};

Derived d;
Base* base = &d;

if (Derived* p = dynamic_cast<Derived*>(base)) {
    p->value = 30;          // checked downcast succeeded
}
\end{lstlisting}

Przy przekazywaniu parametrów funkcji nie używaj rzutowań
jako sposobu na obejście \texttt{const}, sprawdzania czasu życia
ani zasad własności. Dobrze zaprojektowany interfejs zazwyczaj
w ogóle nie wymaga rzutowania wskaźników.

\section{Pamięć dynamiczna: \texttt{new}, \texttt{delete} i ich ograniczenia}

\subsection{Klasyczny mechanizm}

Operator \texttt{new} tworzy obiekt w pamięci przydzielonej
dynamicznie i zwraca wskaźnik. Odpowiadający mu
\texttt{delete} niszczy obiekt i zwalnia pamięć:

\begin{lstlisting}
int* p = new int(42);
// use *p

delete p;
p = nullptr;           // avoids this one stale alias
\end{lstlisting}

Dla tablic obowiązuje para \texttt{new[]} i \texttt{delete[]}:

\begin{lstlisting}
int* a = new int[5]{}; // zero-initialized elements
// use a[0], ..., a[4]

delete[] a;
\end{lstlisting}

\textbf{Nie wolno} łączyć \texttt{new[]} z \texttt{delete} ani
zwykłego \texttt{new} z \texttt{delete[]}. Ponowne
\texttt{delete} tego samego żywego przydziału również jest błędem.
\texttt{delete nullptr;} jest natomiast bezpieczne.

\subsection{Dlaczego ręczne zarządzanie pamięcią jest niebezpieczne?}

\begin{lstlisting}
void fragile() {
    int* p = new int(42);
    operation_that_may_throw();
    delete p;   // skipped if the operation throws
}
\end{lstlisting}

Jeśli funkcja zgłosi wyjątek, pamięć zostaje utracona
(\emph{memory leak}). Analogiczny problem powstaje przy
wielu instrukcjach \texttt{return}, kopiowaniu wskaźników i
niespójnym przekazywaniu odpowiedzialności za \texttt{delete}.

\paragraph{Zalecenie.}
W typowym kodzie aplikacyjnym nie pisz bezpośrednio
\texttt{new}/\texttt{delete}. Użyj obiektu lokalnego,
\texttt{std::vector<T>} albo inteligentnego wskaźnika.
Ręczne operacje alokowania pozostają potrzebne w implementacjach
niskopoziomowych, alokatorach i przy integracji z określonym API.

\section{RAII i semantyka własności}

\subsection{RAII: zasób należy do obiektu}

\emph{Resource Acquisition Is Initialization} (RAII) oznacza
powiązanie przejęcia zasobu z utworzeniem obiektu zarządzającego,
a zwolnienia zasobu z jego destruktorem. Stosuje się je nie tylko
do pamięci, ale też do plików, blokad, gniazd sieciowych i uchwytów
systemowych.

\begin{lstlisting}
#include <fstream>
#include <string>

void save(const std::string& path) {
    std::ofstream file(path);
    file << "example\n";
} // file is closed by its destructor
\end{lstlisting}

Destruktor uruchamia się podczas normalnego opuszczania zakresu
oraz podczas odwijania stosu po wyjątku. Stąd RAII upraszcza
bezpieczeństwo wyjątkowe.

\subsection{Trzy rodzaje relacji z obiektem}

\begin{description}
    \item[Własność wyłączna (\emph{unique ownership}).]
    Dokładnie jeden obiekt zarządzający odpowiada za zniszczenie
    zasobu: najczęściej \texttt{std::unique\_ptr<T>}.
    \item[Własność współdzielona (\emph{shared ownership}).]
    Wielu właścicieli może utrzymywać jeden zasób przy życiu:
    \texttt{std::shared\_ptr<T>}.
    \item[Dostęp bez własności (\emph{non-owning access}).]
    Użytkownik widzi obiekt, ale nie przedłuża jego życia:
    zwykły wskaźnik, referencja, \texttt{std::span<T>} lub
    \texttt{std::weak\_ptr<T>} (ten ostatni obserwuje stan
    zarządzany przez \texttt{shared\_ptr} i pozwala warunkowo
    uzyskać właściciela).
\end{description}

Dobór reprezentacji wynika z \emph{kontraktu własności}, a nie
z samego faktu, że obiekt znajduje się w pamięci dynamicznej.

\section{\texttt{std::unique\_ptr}: jeden właściciel}

\begin{center}\small
\begin{tabularx}{\textwidth}{@{}p{0.36\textwidth}X@{}}
\toprule
Mechanizm & Własność i czas życia \\
\midrule
\code{std::make_unique<T>(...)} & Tworzy obiekt i zwraca jego jedynego właściciela. \\
\code{unique_ptr::get()} & Zwraca surowy wskaźnik bez przekazania własności. \\
\code{unique_ptr::reset(p)} & Niszczy obecnie posiadany obiekt i przejmuje wskaźnik \code{p}. \\
\code{unique_ptr::release()} & Oddaje surowy wskaźnik bez zniszczenia obiektu; odtąd trzeba osobno zadbać o jego zwolnienie. \\
\code{std::shared_ptr<T>} & Współdzieli własność; obiekt znika po odejściu ostatniego właściciela. \\
\code{weak_ptr::lock()} & Próbuje uzyskać \code{shared_ptr}; zwraca pusty wynik, jeśli obiekt wygasł. \\
\bottomrule
\end{tabularx}
\end{center}

\subsection{Tworzenie obiektu i automatyczne niszczenie}

\texttt{std::unique\_ptr<T>} (nagłówek \texttt{<memory>}, C++11)
zarządza pojedynczym obiektem przez wyłączną własność. Od C++14
najczęściej tworzymy go przez \texttt{std::make\_unique<T>}:

\begin{lstlisting}
#include <memory>
#include <string>

#include <utility>

struct Person {
    explicit Person(std::string n) : name(std::move(n)) {}
    std::string name;
};

void example() {
    auto p = std::make_unique<Person>("Ala");
    p->name = "Ola";
    // no delete needed
} // destructor destroys Person
\end{lstlisting}

\texttt{unique\_ptr} można traktować jako obiekt zarządzający:
jego destruktor wywołuje odpowiednią funkcję usuwającą zasób
(\emph{deleter}). Domyślnie jest nią \texttt{std::default\_delete<T>}.

\subsection{Brak kopiowania i przenoszenie własności}

Wskaźnik o własności wyłącznej nie może być kopiowany.
Może natomiast zostać \emph{przeniesiony} za pomocą
\texttt{std::move}:

\begin{lstlisting}
#include <memory>
#include <utility>

auto first = std::make_unique<int>(10);
// auto bad = first;                  // error: copy

auto second = std::move(first);      // transfers ownership
// first == nullptr
*second = 20;
\end{lstlisting}

\texttt{std::move} samo niczego nie zwalnia i nie przenosi fizycznie
wskazywanego obiektu. Zwraca wyrażenie pozwalające wywołać konstruktor
lub przypisanie przenoszące. W tym przypadku przenoszona jest
\emph{własność wskaźnika}, a obiekt typu \texttt{int} nie zmienia adresu.

\subsection{Przekazywanie do funkcji: kto jest właścicielem?}

Rozróżnijmy trzy interfejsy:

\begin{lstlisting}
#include <memory>
#include <utility>

struct Widget { void run() {} };

void inspect(const Widget& w);             // borrows, non-null
void maybe_inspect(const Widget* w);       // borrows, nullable
void take(std::unique_ptr<Widget> w);      // takes ownership

void demo() {
    auto p = std::make_unique<Widget>();
    inspect(*p);                // p still owns the object
    maybe_inspect(p.get());     // p still owns the object
    take(std::move(p));         // ownership transferred
    // p is now empty
}
\end{lstlisting}

Jeżeli funkcja \emph{jedynie używa} obiektu, przekazywanie przez
\texttt{T\&}, \texttt{const T\&} lub \texttt{T*} jest zwykle
czytelniejsze niż przekazywanie \texttt{unique\_ptr} przez referencję.
Parametr \texttt{unique\_ptr<T>} przyjmowany \emph{przez wartość}
sygnalizuje przejęcie własności. Dla funkcji, która ma zastąpić
obiekt będący własnością wywołującego, można użyć
\texttt{std::unique\_ptr<T>\&}.

\subsection{Metody \texttt{get}, \texttt{reset}, \texttt{release}, \texttt{swap}}

\begin{lstlisting}
#include <memory>
#include <utility>

auto p = std::make_unique<int>(7);
int* observer = p.get();       // borrow; do not delete

p.reset(new int(9));           // destroys old int
// observer is now dangling

auto q = std::make_unique<int>(11);
p.swap(q);                    // exchanges managed objects

int* raw = q.release();        // q becomes empty; raw owns now
// responsibility for delete raw is now explicit
delete raw;
\end{lstlisting}

\texttt{get()} zwraca dostęp bez własności. \texttt{reset()} zastępuje
zarządzany obiekt, zwalniając dotychczasowy. \texttt{release()} oddaje
surowy wskaźnik \emph{bez zwolnienia obiektu}; dlatego jest potrzebne
głównie przy przekazywaniu własności do interfejsu, który tego wymaga.
Nie należy mylić \texttt{release()} z \texttt{reset()}.

\subsection{Tablice: \texttt{unique\_ptr<T[]>} a \texttt{vector<T>}}

\begin{lstlisting}
#include <memory>
#include <vector>

auto a = std::make_unique<int[]>(5); // array of 5 zeros
a[2] = 17;

std::vector<int> b(5);         // usually more convenient
b.at(2) = 17;                  // checked indexing
\end{lstlisting}

\texttt{unique\_ptr<T[]>} dobiera \texttt{delete[]} i udostępnia
\texttt{operator[]}, lecz sam nie przechowuje rozmiaru.
Zwykle preferuj \texttt{std::vector<T>}, bo posiada rozmiar,
iteratory i bogatszy interfejs. Dla tablicy o stałym rozmiarze
określonym w kodzie rozważ \texttt{std::array<T, N>}.

\subsection{Własny deleter: zasób inny niż pamięć \texttt{new}}

Inteligentny wskaźnik może wywoływać dowolną funkcję zwalniającą
zasób. Typowym przykładem jest plik w interfejsie C:

\begin{lstlisting}
#include <cstdio>
#include <memory>

struct FileCloser {
    void operator()(std::FILE* file) const noexcept {
        if (file) { std::fclose(file); }
    }
};

using FilePtr = std::unique_ptr<std::FILE, FileCloser>;

void example() {
    FilePtr file(std::fopen("log.txt", "w"));
    if (!file) { return; }

    std::fputs("hello\n", file.get());
} // fclose is called automatically
\end{lstlisting}

W praktyce deleter bywa osobną strukturą, a jego typ wchodzi
w skład typu \texttt{unique\_ptr<T, Deleter>}. Mechanizm pozwala
owinąć RAII również zasoby zwalniane przez \texttt{free()},
\texttt{fclose()} lub bibliotekę zewnętrzną.

\subsection{Polimorfizm i wirtualny destruktor}

\begin{lstlisting}
#include <memory>

struct Shape {
    virtual ~Shape() = default;
    virtual double area() const = 0;
};

struct Circle : Shape {
    explicit Circle(double r) : radius(r) {}
    double area() const override { return 3.14159 * radius * radius; }
    double radius;
};

std::unique_ptr<Shape> make_shape() {
    return std::make_unique<Circle>(2.0);
}
\end{lstlisting}

Przy niszczeniu obiektu pochodnego przez wskaźnik do bazy
zwykle potrzebny jest wirtualny destruktor bazy.
Dotyczy to również przenoszenia własności z klasy pochodnej
do bazowej przy użyciu \texttt{unique\_ptr} i domyślnego deletera.

\section{\texttt{std::shared\_ptr}: współdzielona własność}

\subsection{Jak działa licznik właścicieli?}

\texttt{std::shared\_ptr<T>} (C++11) udostępnia współdzieloną
własność zasobu. Zwykle tworzymy go za pomocą
\texttt{std::make\_shared<T>}:

\begin{lstlisting}
#include <memory>

auto a = std::make_shared<int>(10);
auto b = a;             // shared ownership

// a.use_count() == 2 (provided there are no other owners)
b.reset();              // releases b's ownership
// a.use_count() == 1
\end{lstlisting}

Gdy liczba \emph{silnych właścicieli} spadnie do zera, zarządzany
obiekt zostaje zniszczony. Implementacja utrzymuje tzw.
\emph{blok kontrolny} (\emph{control block}) zawierający między
innymi liczniki właścicieli, informacje potrzebne do zniszczenia
zasobu i ewentualny deleter.

\paragraph{Schemat ideowy.}
\begin{lstlisting}
shared_ptr a ---\
                +--> [ control block: strong = 2 ] --> [ T ]
shared_ptr b ---/           weak = 0
\end{lstlisting}

\texttt{use\_count()} ma znaczenie przede wszystkim diagnostyczne.
Nie powinno się na podstawie odczytanego licznika podejmować decyzji
o bezpieczeństwie wątkowym ani o wyłącznej własności, która może
zmienić się równolegle.

\subsection{Tworzenie i przekazywanie \texttt{shared\_ptr}}

\begin{lstlisting}
#include <memory>

struct Data { int x = 0; };

std::shared_ptr<Data> create() {
    return std::make_shared<Data>();
}

void keep(std::shared_ptr<Data> p) {
    // stores a copy of p somewhere, if needed
}

void read_only(const Data& data) {
    // does not require shared ownership
}

void demo() {
    auto p = create();
    read_only(*p); // borrow
    keep(p);       // creates another temporary shared owner
}
\end{lstlisting}

Parametr \texttt{shared\_ptr<T>} przyjmowany przez wartość oznacza,
że funkcja otrzymuje własny udział we własności na czas wywołania
i ewentualnie może go przechować. Jeśli funkcja wyłącznie
czyta obiekt, zwykle wystarcza \texttt{const T\&} lub \texttt{T*}.
Samo przekazanie \texttt{const shared\_ptr<T>\&} unika inkrementacji
licznika, ale bardziej wiąże interfejs z modelem własności.

\subsection{\texttt{make\_shared} a konstruktor ze wskaźnika}

W typowych implementacjach \texttt{make\_shared} umieszcza obiekt
i blok kontrolny w jednym przydziale pamięci, co może ograniczać
narzut alokacji. Natomiast konstrukcja \texttt{shared\_ptr<T>(new T)}
zwykle wymaga osobnych przydziałów.

\begin{lstlisting}
auto preferred = std::make_shared<int>(42);
std::shared_ptr<int> also_valid(new int(42));
\end{lstlisting}

Nie wolno jednak tworzyć dwóch \texttt{shared\_ptr} niezależnie
z \emph{tego samego} surowego wskaźnika:

\begin{lstlisting}
int* raw = new int(42);
std::shared_ptr<int> a(raw);
// std::shared_ptr<int> b(raw); // WRONG: second control block
\end{lstlisting}

Dwa niezależne bloki kontrolne próbowałyby zwolnić ten sam obiekt.
Poprawne współdzielenie powstaje przez \emph{kopiowanie istniejącego}
\texttt{shared\_ptr}: \texttt{auto b = a;}.

\paragraph{Kiedy nie stosować \texttt{make\_shared}?}
Zintegrowany przydział może zatrzymywać pamięć zajmowaną przez
obiekt w przydzielonym bloku do momentu zniszczenia także ostatniego
\texttt{weak\_ptr}, nawet gdy destruktor obiektu już się wykonał.
Przy dużych obiektach i długo żyjących słabych referencjach
warto rozważyć osobny przydział. Innymi powodami mogą być
specjalne wymagania alokacji lub konstruktory niedostępne
dla \texttt{make\_shared}.

\subsection{Współdzielenie własności a współdzielenie mutacji}

Sam \texttt{shared\_ptr} nie synchronizuje dostępu do pól obiektu:

\begin{lstlisting}
#include <memory>
#include <mutex>

struct Counter {
    std::mutex mutex;
    int value = 0;

    void increment() {
        std::lock_guard<std::mutex> lock(mutex);
        ++value;
    }
};

auto p = std::make_shared<Counter>();
// many threads may own copies of p, but must synchronize
// concurrent modifications of Counter::value
\end{lstlisting}

Blok kontrolny pozwala różnym egzemplarzom \texttt{shared\_ptr}
bezpiecznie współdzielić własność. Równoczesna modyfikacja
\emph{tego samego} egzemplarza wymaga jednak synchronizacji.
Od C++20 można użyć \texttt{std::atomic<std::shared\_ptr<T>>}.

\section{\texttt{std::weak\_ptr}: obserwacja bez utrzymywania przy życiu}

\subsection{Problem cykli własności}

Gdy obiekty wzajemnie posiadają się przez \texttt{shared\_ptr},
liczniki właścicieli mogą nigdy nie spaść do zera:

\begin{lstlisting}
#include <memory>

struct Node {
    std::shared_ptr<Node> next;
};

void leak_cycle() {
    auto a = std::make_shared<Node>();
    auto b = std::make_shared<Node>();
    a->next = b;
    b->next = a;          // cycle of strong ownership
} // nodes remain alive after both locals are destroyed
\end{lstlisting}

To nie jest problem \emph{niezerowego wskaźnika}, lecz błędnego
modelu własności: każdy węzeł utrzymuje przy życiu drugi.

\subsection{Przerywanie cyklu i \texttt{lock()}}

\texttt{std::weak\_ptr<T>} obserwuje ten sam blok kontrolny,
ale nie zwiększa liczby silnych właścicieli:

\begin{lstlisting}
#include <memory>

struct Node {
    std::weak_ptr<Node> parent;         // non-owning
    std::shared_ptr<Node> child;        // owning
};

void example() {
    auto root = std::make_shared<Node>();
    auto leaf = std::make_shared<Node>();

    root->child = leaf;
    leaf->parent = root;

    if (auto parent = leaf->parent.lock()) {
        // parent is a temporary shared owner
        (void)parent;
    }
} // no ownership cycle
\end{lstlisting}

Metoda \texttt{lock()} atomowo próbuje uzyskać
\texttt{shared\_ptr<T>}. Jeżeli obiekt jeszcze istnieje,
zwraca właściciela utrzymującego go przy życiu przez czas
używania. W przeciwnym razie zwraca pusty \texttt{shared\_ptr}.

\begin{lstlisting}
std::weak_ptr<int> observer;
{
    auto p = std::make_shared<int>(8);
    observer = p;
    if (auto alive = observer.lock()) {
        *alive += 1;           // safe while alive owns it
    }
}

// observer.expired() is now true
// observer.lock() now returns an empty shared_ptr
\end{lstlisting}

W kodzie współbieżnym nie stosuj schematu
\texttt{if (!w.expired()) { ... }} z późniejszym odwołaniem:
między sprawdzeniem a użyciem obiekt może zniknąć.
Preferuj jedną operację \texttt{lock()} i posługuj się
otrzymanym \texttt{shared\_ptr}.

\section{\texttt{enable\_shared\_from\_this} i wskaźnik \texttt{this}}

Gdy metoda obiektu zarządzanego przez \texttt{shared\_ptr}
chce przekazać dalej \emph{współdzieloną własność do siebie},
nie może wykonać \texttt{std::shared\_ptr<T>(this)}:
powstałby drugi, niezależny blok kontrolny.

Można zamiast tego dziedziczyć publicznie po
\texttt{std::enable\_shared\_from\_this<T>}:

\begin{lstlisting}
#include <memory>

struct Session : std::enable_shared_from_this<Session> {
    std::shared_ptr<Session> self() {
        return shared_from_this();
    }
};

void example() {
    auto session = std::make_shared<Session>();
    auto second_owner = session->self();
    // both share the same control block
}
\end{lstlisting}

Metodę \texttt{shared\_from\_this()} wolno wywołać dopiero wtedy,
gdy obiektem zarządza już \texttt{shared\_ptr}. Jej wywołanie
na obiekcie lokalnym albo w konstruktorze, zanim powstanie
właściciel, zgłasza wyjątek \texttt{std::bad\_weak\_ptr}.

Od C++17 dostępna jest również metoda \texttt{weak\_from\_this()}.

\section{Zaawansowany przykład: wskaźnik aliasujący \texttt{shared\_ptr}}

\texttt{shared\_ptr} rozróżnia \emph{obiekt zarządzany} (którego
czas życia kontroluje blok kontrolny) i \emph{przechowywany
wskaźnik} (zwracany przez \texttt{get()}). Specjalny konstruktor
aliasujący pozwala udostępnić adres składowej, jednocześnie
utrzymując przy życiu cały obiekt:

\begin{lstlisting}
#include <memory>

struct Pair {
    int first = 1;
    int second = 2;
};

void example() {
    auto whole = std::make_shared<Pair>();
    std::shared_ptr<int> member(whole, &whole->second);

    whole.reset();       // Pair is still alive
    *member = 20;        // accesses Pair::second safely
} // member destroys the last strong owner
\end{lstlisting}

Konstruktor aliasujący jest przydatny np. przy udostępnianiu
podobiektów, ale wymaga prawidłowej relacji czasów życia.
Nie należy go traktować jako sposobu na naprawienie
wcześniej wiszącego wskaźnika.

\section{Wskaźniki w interfejsach funkcji i klas}

\subsection{Katalog typowych kontraktów}

\begin{center}
\small
\begin{tabular}{|p{0.31\textwidth}|p{0.56\textwidth}|}
\hline
\textbf{Interfejs} & \textbf{Znaczenie} \\
\hline
\texttt{void f(const T\&)} & Obiekt wymagany, tylko odczyt; brak przekazania własności. \\
\hline
\texttt{void f(T\&)} & Obiekt wymagany, możliwa modyfikacja; brak przekazania własności. \\
\hline
\texttt{void f(const T*)} & Obiekt opcjonalny, tylko odczyt; brak własności. \\
\hline
\texttt{void f(T*)} & Obiekt opcjonalny, możliwa modyfikacja; brak własności. \\
\hline
\texttt{void f(unique\_ptr<T>)} & Przejęcie wyłącznej własności. \\
\hline
\texttt{unique\_ptr<T> f()} & Przekazanie wyłącznej własności wywołującemu. \\
\hline
\texttt{void f(shared\_ptr<T>)} & Uzyskanie udziału we współdzielonej własności. \\
\hline
\texttt{void f(weak\_ptr<T>)} & Możliwość obserwacji obiektu, którego życie kontrolują inni. \\
\hline
\texttt{void f(span<const T>)} & Widok na ciągłą sekwencję, bez własności. \\
\hline
\end{tabular}
\end{center}

W tabeli pominięto przestrzeń nazw \texttt{std::} dla czytelności.
To \emph{konwencje projektowania API}, a nie automatyczne
gwarancje długowieczności surowych wskaźników i referencji.

\subsection{Dlaczego nie zawsze używać \texttt{shared\_ptr}?}

Rozważmy funkcję, która wyświetla nazwę:

\begin{lstlisting}
#include <iostream>
#include <memory>
#include <string>

struct Product { std::string name; };

void print(const Product& p) {
    std::cout << p.name << '\n';
}

void example() {
    Product local{"book"};
    print(local);                  // no allocation

    auto owned = std::make_unique<Product>(Product{"pen"});
    print(*owned);                 // still no extra owner
}
\end{lstlisting}

Funkcja \texttt{print} nie potrzebuje znać strategii
przechowywania obiektu. Przekazanie referencji oddziela
\emph{sposób dostępu} od \emph{sposobu zarządzania czasem życia}.

\section{PImpl: ukrywanie implementacji przez \texttt{unique\_ptr}}

Technika \emph{pointer to implementation} pozwala utrzymać
szczegóły implementacji poza publicznym nagłówkiem.
Typ implementacyjny może być niekompletny w nagłówku, ale
miejsce wykonania domyślnego deletera musi widzieć jego definicję.
Dlatego destruktor klasy definiujemy w pliku \texttt{.cpp}.

\begin{lstlisting}
// widget.hpp
#pragma once
#include <memory>

class Widget {
public:
    Widget();
    ~Widget();                 // defined in .cpp

    Widget(Widget&&) noexcept;
    Widget& operator=(Widget&&) noexcept;

    Widget(const Widget&) = delete;
    Widget& operator=(const Widget&) = delete;

    void work();

private:
    struct Impl;
    std::unique_ptr<Impl> impl_;
};
\end{lstlisting}

\begin{lstlisting}
// widget.cpp
#include "widget.hpp"
#include <memory>

struct Widget::Impl {
    int value = 0;
    void run() { ++value; }
};

Widget::Widget() : impl_(std::make_unique<Impl>()) {}
Widget::~Widget() = default;
Widget::Widget(Widget&&) noexcept = default;
Widget& Widget::operator=(Widget&&) noexcept = default;

void Widget::work() { impl_->run(); }
\end{lstlisting}

Ten przykład prezentuje semantykę przenoszenia. Po przeniesieniu
źródłowy obiekt \texttt{Widget} ma pusty \texttt{impl\_};
nie wolno wywoływać na nim \texttt{work()} bez odpowiedniego
sprawdzenia. W produkcyjnym API można ten stan jawnie obsłużyć
lub określić w kontrakcie.

\section{Wskaźniki do funkcji i składowych klas}

Wskaźniki nie muszą odnosić się do danych. Można wskazywać
funkcje, a także składowe klas, choć te ostatnie wymagają
innej składni.

\subsection{Wskaźnik do funkcji}

\begin{lstlisting}
int add(int a, int b) { return a + b; }
int mul(int a, int b) { return a * b; }

using Operation = int(*)(int, int);

int apply(int x, int y, Operation op) {
    return op(x, y);
}

int a = apply(2, 3, &add); // 5
int b = apply(2, 3, &mul); // 6
\end{lstlisting}

Obiekty wywoływalne różnych typów można przechowywać za pomocą
\texttt{std::function}. Prosty szablon lub lambda bywają jednak
bardziej efektywne.

\subsection{Wskaźnik do składowej}

\begin{lstlisting}
struct Record {
    int value = 0;
    int twice() const { return 2 * value; }
};

int Record::* field = &Record::value;
int (Record::* method)() const = &Record::twice;

Record r{7};
int x = r.*field;          // 7
int y = (r.*method)();     // 14

Record* p = &r;
int z = p->*field;         // 7
\end{lstlisting}

Wskaźnik do składowej nie jest zwykłym adresem gotowego obiektu.
Dopiero w połączeniu z konkretną instancją klasy pozwala
odwołać się do odpowiedniej składowej.

\section{Wskaźniki i wielowątkowość}

\subsection{\texttt{shared\_ptr} nie eliminuje wyścigów danych}

Możliwość bezpiecznego utrzymania czasu życia obiektu nie oznacza,
że jego pola można dowolnie zmieniać z wielu wątków.
Jeśli dwa wątki zapisują do \texttt{*p}, potrzebna jest osobna
synchronizacja (np. \texttt{std::mutex} albo atomiki).

Od C++20 dostępna jest specjalizacja
\texttt{std::atomic<std::shared\_ptr<T>>}, przydatna do atomowego
publikowania nowej wersji współdzielonego obiektu:

\begin{lstlisting}
#include <atomic>
#include <memory>

struct Config {
    int workers = 4;
};

std::atomic<std::shared_ptr<const Config>> current{
    std::make_shared<const Config>(Config{4})
};

void publish(int count) {
    current.store(std::make_shared<const Config>(Config{count}));
}

int read_workers() {
    auto snapshot = current.load(); // keeps snapshot alive
    return snapshot->workers;
}
\end{lstlisting}

Tu czytelnicy otrzymują stabilny, niemutowalny \emph{snapshot}.
Atomowość dotyczy odczytu i zapisu \texttt{shared\_ptr},
a nie arbitralnych operacji na jego zawartości.

\section{Najczęstsze błędy zebrane w jednym miejscu}

\begin{enumerate}
    \item \textbf{Dereferencja \texttt{nullptr} lub wskaźnika
    niezainicjalizowanego.} Zawsze sprawdzaj kontrakt i inicjalizuj
    wskaźniki.
    \item \textbf{Wskaźnik wiszący.} \texttt{get()}, \texttt{span}
    i surowa referencja nie przedłużają życia zasobu.
    \item \textbf{Podwójne \texttt{delete}.} Jeden przydział musi mieć
    jedną, jednoznaczną strategię zwalniania.
    \item \textbf{Mieszanie \texttt{new[]} i \texttt{delete}.}
    Do tablic służy \texttt{delete[]}; zwykle lepiej wybrać kontener.
    \item \textbf{Zwracanie adresu lokalnej zmiennej.}
    Jej czas życia kończy się przy wyjściu z zakresu.
    \item \textbf{Przekazywanie surowego wskaźnika z \texttt{get()}
    do nowego \texttt{shared\_ptr}.} Tworzy niezależny blok kontrolny.
    \item \textbf{Nieuzasadnione \texttt{shared\_ptr} wszędzie.}
    Współdzielenie własności kosztuje i zaciera odpowiedzialność.
    \item \textbf{Cykl silnych właścicieli.} Przerywaj go przez
    \texttt{weak\_ptr} lub przeprojektowanie relacji.
    \item \textbf{Nadużywanie \texttt{release()}.} Po wywołaniu
    ktoś musi jawnie przejąć odpowiedzialność za zasób.
    \item \textbf{Brak wirtualnego destruktora w bazie.}
    Usuwanie obiektu pochodnego przez wskaźnik bazowy z domyślnym
    deleterem może wówczas prowadzić do UB.
    \item \textbf{Założenie, że \texttt{shared\_ptr} chroni obiekt
    przed wyścigami.} Chroni zarządzanie własnością, nie samą mutację.
    \item \textbf{Przechowywanie wskaźników do elementów \texttt{vector}
    po realokacji.} Przy operacjach zmieniających bufor pobieraj
    adresy ponownie.
\end{enumerate}

\section{Jak wybrać właściwy typ?}

\begin{center}
\small
\begin{tabular}{|p{0.43\textwidth}|p{0.43\textwidth}|}
\hline
\textbf{Sytuacja} & \textbf{Najczęstszy wybór} \\
\hline
Pojedynczy obiekt o prostym czasie życia & Zwykła zmienna \texttt{T} \\
\hline
Dynamiczna lista elementów & \texttt{std::vector<T>} \\
\hline
Obiekt wymagany tylko do odczytu & \texttt{const T\&} \\
\hline
Obiekt wymagany i modyfikowany & \texttt{T\&} \\
\hline
Obiekt opcjonalny bez przekazania własności & \texttt{T*} lub \texttt{const T*} \\
\hline
Pojedynczy właściciel dynamicznego obiektu & \texttt{std::unique\_ptr<T>} \\
\hline
Kilku niezależnych właścicieli & \texttt{std::shared\_ptr<T>} \\
\hline
Obserwacja obiektu współdzielonego & \texttt{std::weak\_ptr<T>} \\
\hline
Fragment istniejącej tablicy, bez własności & \texttt{std::span<T>} \\
\hline
Polimorficzna kolekcja obiektów & \texttt{std::vector} zawierający \texttt{unique\_ptr} \\
\hline
\end{tabular}
\end{center}

\subsection{Praktyczny przykład: kolekcja obiektów polimorficznych}

\begin{lstlisting}
#include <iostream>
#include <memory>
#include <vector>

struct Task {
    virtual ~Task() = default;
    virtual void run() const = 0;
};

struct PrintTask : Task {
    void run() const override { std::cout << "hello\n"; }
};

int main() {
    std::vector<std::unique_ptr<Task>> tasks;
    tasks.push_back(std::make_unique<PrintTask>());

    for (const auto& task : tasks) {
        task->run();
    }
} // vector and tasks destroy everything automatically
\end{lstlisting}

\section{Podsumowanie i praktyczne reguły}

\begin{itemize}
    \item Preferuj \emph{przechowywanie przez wartość}, jeśli
    obiekt nie wymaga odrębnej, dynamicznej tożsamości.
    \item Gdy obiekt ma jednego właściciela, używaj
    \texttt{std::unique\_ptr}; własność przenoś przez
    \texttt{std::move}.
    \item Sięgaj po \texttt{std::shared\_ptr} dopiero wtedy,
    gdy kilka części programu \emph{rzeczywiście} musi niezależnie
    utrzymywać ten sam obiekt przy życiu.
    \item Jeśli obserwator nie powinien utrzymywać obiektu przy życiu,
    użyj \texttt{std::weak\_ptr} i metody \texttt{lock()}.
    \item Do niewłasnościowego dostępu używaj
    \texttt{T\&}, \texttt{const T\&}, \texttt{T*} lub
    \texttt{std::span<T>} w zależności od kontraktu.
    \item Nie dereferencjonuj wskaźnika tylko dlatego, że jest
    różny od \texttt{nullptr}: obiekt może już nie istnieć.
    \item Preferuj \texttt{std::vector} i \texttt{std::array}
    zamiast ręcznie alokowanych tablic.
    \item Zasoby systemowe obejmuj RAII i odpowiednim deleterem;
    osobno analizuj bezpieczeństwo wątkowe wskazywanego obiektu.
\end{itemize}

\paragraph{Konkluzja.}
Najważniejsze pytanie brzmi nie ``jak zdobyć adres?'', ale
``\emph{kto jest właścicielem obiektu, jak długo on istnieje
i kto ma prawo z niego korzystać?}''. Właśnie te trzy kwestie
rozstrzygają, czy w danym miejscu potrzebujemy referencji,
surowego wskaźnika, \texttt{unique\_ptr}, \texttt{shared\_ptr}
czy \texttt{weak\_ptr}.

% Dalsza lektura:
% https://en.cppreference.com/cpp/memory/unique_ptr
% https://en.cppreference.com/cpp/memory/shared_ptr
% https://en.cppreference.com/cpp/memory/weak_ptr
% https://en.cppreference.com/cpp/memory/enable_shared_from_this
% https://en.cppreference.com/cpp/container/span

% ROZDZIAL 6 -- Semantyka przenoszenia, kategorie wartosci, perfect forwarding (C++20)
% Fragment do wstawienia za rozdzialem 5.
% Wymagana preambula: pakiet listings i konfiguracja lstlisting.
% Uzycie: \input{rozdzial_6_move_semantics_cpp.tex}

\chapter[Semantyka przenoszenia w C++]{Semantyka przenoszenia, kategorie wartości i przekazywanie doskonałe}
\label{chap:move-semantics-cpp}

\section{Dlaczego wprowadzono semantykę przenoszenia?}

Rozważmy kontener przechowujący milion liczb. Skopiowanie go oznacza
zwykle przydzielenie nowej pamięci i skopiowanie każdego elementu.
Jeżeli dotychczasowy obiekt nie będzie już potrzebny, bardziej naturalna
jest operacja \emph{przekazania własności} jego zasobu zamiast kopiowania.
Od C++11 język udostępnia \emph{semantykę przenoszenia} (\textit{move
semantics}): programista może rozróżnić sytuację, w której obiekt źródłowy
ma pozostać niezmieniony, od sytuacji, w której jego zasoby wolno wykorzystać
ponownie.

\begin{lstlisting}
#include <string>
#include <utility>
#include <vector>

std::vector<std::string> a(100'000, "data");
auto b = a;                // copy: a remains intact

auto c = std::move(a);     // move may transfer a's allocation
// c now owns its elements; a is still valid, but its state
// after this move should not be assumed without a guarantee.
\end{lstlisting}

Nie oznacza to, że każde przeniesienie jest szybsze niż kopiowanie.
Dla typu \texttt{int} oba działania są w praktyce tożsame, a niektóre typy
kopiują część małych danych również podczas przenoszenia.
Kluczowa idea jest semantyczna: \emph{przeniesienie pozwala zmienić stan
źródła, a kopiowanie nie powinno go zmieniać}.

\paragraph{Zakres rozdziału.}
Zakładamy \textbf{C++20}. Kategorie wartości istnieją od początku języka,
referencje do r-wartości i \texttt{std::move} wprowadzono w C++11,
\emph{guaranteed copy elision} w C++17, a wiele przykładów poniżej
wykorzystuje koncepcje omówione w rozdziale~5.
Wszystkie przykłady zawierają komentarze angielskie, aby łatwo
przenosić je do projektów.

\section{Typ wyrażenia a kategoria wartości}
\label{sec:move-categories}

\textbf{Typ} i \textbf{kategoria wartości} są dwiema różnymi własnościami
wyrażenia. Typ mówi, jaki rodzaj danych otrzymujemy; kategoria mówi
m.in., czy wyrażenie identyfikuje istniejący obiekt i jak można
wiązać je z referencjami. Dokładnie trzy podstawowe kategorie to:
\emph{lvalue}, \emph{xvalue} i \emph{prvalue}.

\begin{center}
\begin{tabular}{|l|p{10.5cm}|}
\hline
\textbf{Kategoria} & \textbf{Intuicja i przykłady} \\
\hline
\texttt{lvalue} & Wyrażenie identyfikujące obiekt lub funkcję; np. \texttt{x}, \texttt{*ptr}, \texttt{++x} dla wbudowanego \texttt{int}. \\
\hline
\texttt{xvalue} & \emph{Expiring value}: wyrażenie identyfikujące obiekt, którego zasoby mogą być ponownie wykorzystane; np. \texttt{std::move(x)}. \\
\hline
\texttt{prvalue} & \emph{Pure rvalue}: wyrażenie obliczające wartość lub inicjalizujące obiekt; np. \texttt{42}, \texttt{x + 1}, \texttt{std::string("hi")}. \\
\hline
\end{tabular}
\end{center}

Pozostałe dwie nazwy oznaczają sumy kategorii:
\[
\underbrace{\mathrm{glvalue}}_{\mathrm{lvalue}\,\cup\,\mathrm{xvalue}},
\qquad
\underbrace{\mathrm{rvalue}}_{\mathrm{prvalue}\,\cup\,\mathrm{xvalue}}.
\]
Nie należy utożsamiać \texttt{rvalue} z tymczasowym obiektem ani
\texttt{lvalue} z wyrażeniem znajdującym się po lewej stronie przypisania.
Są to techniczne kategorie zdefiniowane przez język.

\subsection{Kilka wyrażeń o zaskakujących kategoriach}

\begin{lstlisting}
#include <string>
#include <utility>

int x = 7;
int* p = &x;

// x                 : lvalue
// *p                : lvalue
// (x)               : lvalue
// 7                 : prvalue
// x + 1             : prvalue
// std::move(x)      : xvalue
// std::string{"abc"}: prvalue

int& lref = x;
int&& rref = 9;

// lref              : lvalue
// rref              : lvalue (!) - it has a name
// std::move(rref)   : xvalue
\end{lstlisting}

To fundamentalna zasada: \textbf{nazwa zmiennej jest wyrażeniem lvalue,
nawet jeśli typ tej zmiennej to referencja rvalue \texttt{T\&\&}}.
Samo zadeklarowanie \texttt{T\&\&} nie oznacza, że każde użycie takiej
zmiennej będzie r-wartością.

\subsection{Badanie kategorii za pomocą \texttt{decltype((expr))}}

Dla prostego, nieujętego w nawiasy identyfikatora składnia
\texttt{decltype(x)} zwraca deklarowany typ. Dla ogólnego wyrażenia
\texttt{decltype((expr))} oddaje jego kategorię poprzez typ:
\[
\mathrm{decltype}((e))=
\left\{
\begin{array}{ll}
T\&, & \mbox{gdy } e \mbox{ jest lvalue typu } T,\\
T\&\&, & \mbox{gdy } e \mbox{ jest xvalue typu } T,\\
T, & \mbox{gdy } e \mbox{ jest prvalue typu } T.
\end{array}
\right.
\]

\begin{lstlisting}
#include <type_traits>
#include <utility>

int x = 10;
static_assert(std::is_same_v<decltype(x), int>);
static_assert(std::is_same_v<decltype((x)), int&>);
static_assert(std::is_same_v<decltype((std::move(x))), int&&>);
static_assert(std::is_same_v<decltype((x + 1)), int>);
\end{lstlisting}

\paragraph{Uwaga.} \texttt{decltype} stosuje szczególne reguły dla
nieujętych w nawiasy nazw i dostępu do składowych. Dlatego
\texttt{decltype(x)} oraz \texttt{decltype((x))} mogą zwrócić
różne typy. Nie jest to sprzeczność, lecz zamierzona cecha języka.

\subsection{Różnica między \texttt{auto} i \texttt{decltype(auto)}}

Zwykłe \texttt{auto} w deklaracji zmiennej najczęściej dedukuje typ
przez wartość i pomija referencję wynikającą z inicjalizatora.
Natomiast \texttt{decltype(auto)} stosuje reguły \texttt{decltype}
i może zachować kategorię wartości wyrażenia.

\begin{lstlisting}
#include <type_traits>

int n = 5;
auto copy = (n);           // int: independent value
decltype(auto) alias = (n); // int&: refers to n

static_assert(std::is_same_v<decltype(copy), int>);
static_assert(std::is_same_v<decltype(alias), int&>);

auto&& first = n;          // deduces int&
auto&& second = 5;         // deduces int&&
static_assert(std::is_same_v<decltype(first), int&>);
static_assert(std::is_same_v<decltype(second), int&&>);
\end{lstlisting}

W typie wyniku funkcji zapis \texttt{decltype(auto)} wymaga
ostrożności. Instrukcja \texttt{return (x);} może zwrócić
referencję do zmiennej lokalnej, która przestanie istnieć.
Zwykłe \texttt{auto} jest często bezpieczniejszym wyborem,
gdy funkcja ma zwrócić własny obiekt przez wartość.

\section{Referencje lvalue i rvalue}
\label{sec:move-references}

\begin{lstlisting}
#include <string>

std::string name = "Ada";
std::string& a = name;            // binds to mutable lvalue
const std::string& b = name;      // binds to lvalue
const std::string& c = "temporary"; // can bind to temporary
std::string&& d = std::string{"x"}; // binds to rvalue

// std::string& e = std::string{"x"}; // error
// std::string&& f = name;           // error
\end{lstlisting}

Referencja \texttt{T\&} umożliwia pracę z istniejącym, zwykle
nazwanym obiektem; \texttt{T\&\&} może wiązać się z r-wartością.
\texttt{const T\&} wiąże się zarówno z lvalue, jak i z rvalue,
nie dając prawa do modyfikowania obiektu poprzez tę referencję.

\subsection{Wybór przeciążenia}

\begin{lstlisting}
#include <iostream>
#include <string>
#include <utility>

void inspect(const std::string&) {
    std::cout << "const lvalue reference\n";
}
void inspect(std::string&&) {
    std::cout << "rvalue reference\n";
}

int main() {
    std::string text = "hello";
    inspect(text);             // const lvalue reference
    inspect(std::string{"x"}); // rvalue reference
    inspect(std::move(text));  // rvalue reference
}
\end{lstlisting}

Istnienie przeciążenia \texttt{f(T\&\&)} samo w sobie nie przesądza
o rzeczywistym przeniesieniu czegokolwiek. Funkcja może tylko
odczytać argument; \emph{move} jest operacją wykonywaną przez
konstruktor lub przypisanie, a nie przez mechanizm wiązania referencji.

\subsection{Wydłużenie życia obiektu tymczasowego}

W typowych deklaracjach lokalnych związanie obiektu tymczasowego
bezpośrednio z referencją może wydłużyć jego życie do końca życia
tej referencji:
\begin{lstlisting}
#include <string>

const std::string& a = std::string{"safe until a ends"};
std::string&& b = std::string{"safe until b ends"};
\end{lstlisting}
Nie wolno jednak uogólniać tej reguły na zwracanie referencji ani
na dowolne przekazywanie referencji dalej. Przykładowo:
\begin{lstlisting}
#include <string>

const std::string& broken() {
    return std::string{"dangling"}; // dangling reference
}
\end{lstlisting}
Zwrócona referencja nie utrzymuje tego obiektu przy życiu.
Jeżeli funkcja ma \emph{wytworzyć} wynik, zwykle powinna zwracać go
przez wartość.

\section{Co naprawdę robi \texttt{std::move}?}
\label{sec:move-std-move}

Nazwa jest nieco myląca: \texttt{std::move(x)} \textbf{nie przenosi
zasobów samodzielnie}. W przybliżeniu wykonuje rzutowanie do
referencji rvalue odpowiedniego typu. Dopiero dalszy konstruktor,
operator przypisania lub inna funkcja może potraktować wynik
jako obiekt nadający się do przeniesienia.

\begin{lstlisting}
#include <type_traits>
#include <utility>

template<class T>
constexpr std::remove_reference_t<T>&&
my_move(T&& value) noexcept {
    return static_cast<std::remove_reference_t<T>&&>(value);
}
\end{lstlisting}

Powyższa implementacja ilustruje istotę działania
\texttt{std::move}. W rzeczywistym kodzie korzystamy z funkcji
standardowej z nagłówka \texttt{<utility>}.

\begin{lstlisting}
#include <string>
#include <utility>

std::string first = "long text";
std::string second = std::move(first);
// A move constructor is selected for second.
// first is valid; do not assume its exact contents.

first = "reused";             // valid: assignment
first.clear();                // valid: clear()
\end{lstlisting}

\paragraph{Stan obiektu po przeniesieniu.}
Typy biblioteczne zwykle gwarantują stan \emph{poprawny, lecz
nieokreślony co do zawartości} (\textit{valid but unspecified}), o ile
kontrakt konkretnego typu nie obiecuje więcej. Takiego obiektu
można bezpiecznie zniszczyć, przypisać mu nową wartość albo
wykonać operację, której wymagania wstępne na to pozwalają.
Nie należy jednak \emph{zakładać}, że każdy przeniesiony
\texttt{std::string} jest pusty. Dla niektórych typów, np.
\texttt{std::unique\_ptr}, stan po przeniesieniu określono dokładniej.

\subsection{Problem obiektu \texttt{const}}

\begin{lstlisting}
#include <string>
#include <utility>

const std::string original = "cannot be emptied here";
std::string other = std::move(original);
// Typically: copy, not move (const string&& cannot bind to
// string's non-const move constructor).
\end{lstlisting}

\texttt{std::move(original)} ma w tym przypadku typ
\texttt{const std::string\&\&}. Typowy konstruktor przenoszący
przyjmuje \texttt{std::string\&\&}, więc nie może otrzymać
argumentu \texttt{const}. Kompilator wybiera konstruktor kopiujący
przyjmujący \texttt{const std::string\&}. \textbf{Samo wystąpienie
\texttt{std::move} w kodzie nie gwarantuje przeniesienia}.\footnote{Dla
typów niekopiowalnych ten sam zapis może skończyć się błędem
kompilacji.}

\section{Konstruktor przenoszący i przenoszący operator przypisania}
\label{sec:move-special-members}

W klasie \texttt{T} konstruktor przenoszący ma zwykle sygnaturę
\texttt{T(T\&\&)}, natomiast przenoszący operator przypisania:
\texttt{T\& operator=(T\&\&)}. Nie należy mylić ich z kopiowaniem:
\begin{lstlisting}
struct Item {
    Item() = default;
    Item(const Item&);            // copy constructor
    Item& operator=(const Item&); // copy assignment
    Item(Item&&) noexcept;        // move constructor
    Item& operator=(Item&&) noexcept; // move assignment
    ~Item();                      // destructor
};
\end{lstlisting}

\subsection{Przykład z obserwacją operacji}
\label{sec:move-trace}

\begin{lstlisting}
#include <iostream>
#include <utility>

struct Trace {
    Trace() { std::cout << "default\n"; }
    Trace(const Trace&) { std::cout << "copy ctor\n"; }
    Trace(Trace&&) noexcept { std::cout << "move ctor\n"; }
    Trace& operator=(const Trace&) {
        std::cout << "copy assign\n";
        return *this;
    }
    Trace& operator=(Trace&&) noexcept {
        std::cout << "move assign\n";
        return *this;
    }
};

int main() {
    Trace a;                    // default
    Trace b = a;                // copy ctor
    Trace c = std::move(a);     // move ctor
    b = c;                      // copy assign
    b = std::move(c);           // move assign
}
\end{lstlisting}

Konstruktor tworzy \emph{nowy} obiekt, a operator przypisania
zmienia już istniejący. W tym przykładzie \texttt{Trace} jedynie
wypisuje komunikaty. W klasie zarządzającej zasobem konstruktor
przenoszący powinien przekazać własność zasobu, a operator
przypisania musi dodatkowo poprawnie pozbyć się zasobu
posiadanego wcześniej przez cel.

\subsection{Pełny przykład: klasa samodzielnie zarządzająca tablicą}
\label{sec:move-raw-buffer}

To przykład \emph{edukacyjny}. W rzeczywistym projekcie zwykle
lepiej użyć standardowego kontenera, np. \texttt{std::vector<int>}.

\begin{lstlisting}
#include <algorithm>
#include <cstddef>
#include <utility>

class RawBuffer {
    std::size_t size_ = 0;
    int* data_ = nullptr;

public:
    explicit RawBuffer(std::size_t n)
        : size_(n), data_(n ? new int[n]{} : nullptr) {}

    ~RawBuffer() { delete[] data_; }

    RawBuffer(const RawBuffer& other)
        : RawBuffer(other.size_) {
        if (size_ != 0)
            std::copy_n(other.data_, size_, data_);
    }

    RawBuffer& operator=(const RawBuffer& other) {
        if (this != &other) {
            RawBuffer copy(other); // may throw; *this unchanged
            swap(copy);
        }
        return *this;
    }

    RawBuffer(RawBuffer&& other) noexcept
        : size_(std::exchange(other.size_, 0)),
          data_(std::exchange(other.data_, nullptr)) {}

    RawBuffer& operator=(RawBuffer&& other) noexcept {
        if (this != &other) {
            delete[] data_;
            size_ = std::exchange(other.size_, 0);
            data_ = std::exchange(other.data_, nullptr);
        }
        return *this;
    }

    void swap(RawBuffer& other) noexcept {
        using std::swap;
        swap(size_, other.size_);
        swap(data_, other.data_);
    }

    std::size_t size() const noexcept { return size_; }
    int& operator[](std::size_t i) { return data_[i]; }
    const int& operator[](std::size_t i) const { return data_[i]; }
};
\end{lstlisting}

\texttt{std::exchange(obj, value)} zwraca starą wartość \texttt{obj}
i zastępuje ją wartością \texttt{value}. Konstruktor przenoszący
przejmuje więc wskaźnik i rozmiar, a źródło ustawia w jawnie
określony stan pusty. Operator przenoszącego przypisania
usuwa dotychczasową tablicę celu, przejmuje nową i zabezpiecza
przypadek \texttt{b = std::move(b)}. W implementacji kopiowania
widać z kolei \emph{silną gwarancję wyjątkową}: jeśli utworzenie
kopii nie powiedzie się, pierwotny obiekt pozostaje bez zmian.

\section{Reguła pięciu, reguła trzech i reguła zera}
\label{sec:move-rule-of-five}

Klasy zarządzające zasobem często muszą kontrolować pięć operacji:
(1) destruktor, (2) konstruktor kopiujący, (3) kopiujący operator
przypisania, (4) konstruktor przenoszący oraz (5) przenoszący
operator przypisania. To \emph{reguła pięciu}.

Historyczna \emph{reguła trzech} odnosiła się do destruktora,
konstruktora kopiującego i kopiującego operatora przypisania,
zanim język udostępnił move semantics.

\paragraph{Ważne: deklarowanie destruktora może blokować generowanie
konstruktora przenoszącego.}

\begin{lstlisting}
#include <memory>
#include <type_traits>

struct WithDestructor {
    std::unique_ptr<int> p;
    ~WithDestructor() = default;
};

static_assert(!std::is_move_constructible_v<WithDestructor>);

struct WithMove {
    std::unique_ptr<int> p;
    ~WithMove() = default;
    WithMove() = default;
    WithMove(WithMove&&) noexcept = default;
    WithMove& operator=(WithMove&&) noexcept = default;
};
static_assert(std::is_move_constructible_v<WithMove>);
\end{lstlisting}

W \texttt{WithDestructor} użytkownik zadeklarował destruktor,
więc kompilator nie deklaruje automatycznie operacji przenoszących.
Kopiowanie jest niedostępne z powodu \texttt{unique\_ptr}.
Samo \texttt{= default} przy destruktorze nie usuwa tej reguły.

\subsection{Najlepsza praktyka: reguła zera}

\emph{Reguła zera} mówi, że klasy, które nie są bezpośrednio
odpowiedzialne za posiadanie surowych zasobów, zwykle nie powinny
same definiować żadnej z pięciu operacji.

\begin{lstlisting}
#include <string>
#include <vector>

struct Dataset {
    std::string title;
    std::vector<double> values;
};

Dataset original{"results", {1.0, 2.0, 3.0}};
Dataset copy = original;           // deep copy of vector data
Dataset moved = std::move(original); // member-wise move
\end{lstlisting}

To kontenery standardowe odpowiadają za przydział i zwalnianie
pamięci. Automatycznie generowane operacje \texttt{Dataset}
wykonują odpowiednie operacje na jego składowych. \textbf{Zalecenie:}
zacznij od reguły zera; własny konstruktor przenoszący pisz
jedynie wtedy, gdy zachowanie domyślne nie jest właściwe.

\section{Przenoszenie a wyjątki: \texttt{noexcept}}
\label{sec:move-noexcept}

Kontenery takie jak \texttt{std::vector<T>} przy realokacji muszą
przenieść lub skopiować dotychczasowe elementy do nowego obszaru.
Jeśli przenoszenie \texttt{T} może rzucać wyjątki, a \texttt{T}
można kopiować, biblioteka może wybrać kopiowanie, by zachować
silniejsze gwarancje wyjątkowe. Właśnie dlatego poprawne operacje
przenoszące, które nie mogą rzucać wyjątków, warto deklarować
\texttt{noexcept}.

\begin{lstlisting}
#include <type_traits>
#include <utility>

struct Resource {
    Resource() = default;
    Resource(const Resource&) = default;
    Resource(Resource&&) noexcept = default;
};

static_assert(std::is_nothrow_move_constructible_v<Resource>);

// move_if_noexcept can provide an lvalue reference if
// moving might throw and a copy is available.
\end{lstlisting}

\texttt{std::move\_if\_noexcept(x)} zwraca odpowiednio
\texttt{T\&\&} albo \texttt{const T\&}, zależnie od własności
konstruktorów danego typu. To narzędzie używane przy implementowaniu
kontenerów, a nie uniwersalny zamiennik \texttt{std::move}.

\paragraph{Ostrożnie z \texttt{noexcept}.}
Wyjście wyjątku poza funkcję zadeklarowaną jako
\texttt{noexcept} wywołuje \texttt{std::terminate}.
Dlatego bezwyjątkowość musi wynikać z rzeczywistych operacji składowych.

\section{Zwijanie referencji (\textit{reference collapsing})}
\label{sec:move-ref-collapse}

Gdy referencje występują pośrednio poprzez aliasy lub
podstawienie argumentów szablonu, stosuje się \emph{zwijanie
referencji}. Wszystkie przypadki przedstawia tabela:

\begin{center}
\begin{tabular}{|c|c|c|}
\hline
\textbf{Typ pośredni} & \textbf{Zapis} & \textbf{Wynik} \\
\hline
\texttt{T\&} i \texttt{\&} & \texttt{T\& \&} & \texttt{T\&} \\
\texttt{T\&} i \texttt{\&\&} & \texttt{T\& \&\&} & \texttt{T\&} \\
\texttt{T\&\&} i \texttt{\&} & \texttt{T\&\& \&} & \texttt{T\&} \\
\texttt{T\&\&} i \texttt{\&\&} & \texttt{T\&\& \&\&} & \texttt{T\&\&} \\
\hline
\end{tabular}
\end{center}

Krótko: jeżeli \emph{choć jeden} składnik jest \texttt{\&},
wynik jest \texttt{\&}; tylko \texttt{\&\&} połączone
z \texttt{\&\&} daje \texttt{\&\&}.

\begin{lstlisting}
#include <type_traits>

using L = int&;
using R = int&&;

static_assert(std::is_same_v<L&, int&>);
static_assert(std::is_same_v<L&&, int&>);
static_assert(std::is_same_v<R&, int&>);
static_assert(std::is_same_v<R&&, int&&>);
\end{lstlisting}

Zwijanie referencji jest podstawą \emph{perfect forwarding}:
dzięki niemu jeden szablon może przyjmować zarówno lvalue,
jak i rvalue i przekazywać tę różnicę dalej.

\section{Referencje przekazujące i dedukcja \texttt{T\&\&}}
\label{sec:move-forwarding-ref}

Zapis \texttt{T\&\&} ma dwa znaczenia w zależności od kontekstu.
Jeśli \texttt{T} jest dedukowanym, niekwalifikowanym parametrem
typu \emph{tego samego} szablonu funkcji, \texttt{T\&\&}
to \emph{referencja przekazująca} (\textit{forwarding reference}).
W pozostałych typowych kontekstach jest to po prostu
referencja rvalue.

\begin{lstlisting}
#include <string>

template<class T>
void forwarding(T&& x);         // forwarding reference

template<class T>
void not_forwarding(const T&& x); // rvalue reference

void fixed(std::string&& x);    // rvalue reference

template<class T>
struct Holder {
    void receive(T&& x);        // NOT forwarding: T comes
                               // from the class template

    template<class U>
    void relay(U&& x);          // forwarding: U is deduced
};
\end{lstlisting}

Dla wywołania \texttt{forwarding(x)} z lvalue typu \texttt{int}
kompilator dedukuje \texttt{T = int\&}, zatem parametr ma po
zwinięciu typ \texttt{int\&}. Dla \texttt{forwarding(5)}
dedukuje \texttt{T = int}, a parametr ma typ \texttt{int\&\&}.

\begin{lstlisting}
#include <type_traits>

template<class T>
void check(T&& x) {
    if constexpr (std::is_lvalue_reference_v<T>) {
        // Called with an lvalue.
    } else {
        // Called with an rvalue.
    }
}
\end{lstlisting}

\paragraph{Drobne wyjątki.} \texttt{auto\&\&} zwykle również
pozwala na przekazywanie kategorii, lecz dedukcja \texttt{auto}
ma szczególne reguły dla list inicjalizacyjnych.
\texttt{const T\&\&} nie jest referencją przekazującą.

\section{Perfect forwarding i \texttt{std::forward}}
\label{sec:move-perfect-forwarding}

Przez \emph{perfect forwarding} rozumiemy przekazanie argumentu
do następnej funkcji przy zachowaniu kategorii wartości,
a także odpowiednich kwalifikatorów typu.

\begin{lstlisting}
#include <iostream>
#include <utility>

void target(int&)  { std::cout << "lvalue\n"; }
void target(int&&) { std::cout << "rvalue\n"; }

template<class T>
void incorrect(T&& value) {
    target(value); // value is a named variable: always lvalue
}

template<class T>
void correct(T&& value) {
    target(std::forward<T>(value));
}

int main() {
    int x = 1;
    incorrect(2); // lvalue (!)
    correct(x);   // lvalue
    correct(2);   // rvalue
}
\end{lstlisting}

W przybliżeniu \texttt{std::forward<T>(value)} zwraca
lvalue, jeżeli \texttt{T} wydedukowano jako referencję lvalue,
i xvalue w przeciwnym razie. \texttt{std::move(value)} zawsze
wymusza kategorię xvalue: w funkcji przekazującej byłoby to
błędem, ponieważ niepotrzebnie pozwalałoby na przeniesienie
zasobu z argumentu przekazanego jako lvalue.

\subsection{Przykład: fabryka obiektów}

\begin{lstlisting}
#include <memory>
#include <string>
#include <utility>

struct Person {
    std::string name;
    int age;

    Person(std::string n, int a)
        : name(std::move(n)), age(a) {}
};

template<class T, class... Args>
std::unique_ptr<T> make_object(Args&&... args) {
    return std::unique_ptr<T>(
        new T(std::forward<Args>(args)...)
    );
}

int main() {
    std::string first = "Ada";
    auto p = make_object<Person>(first, 37); // first copied
    auto q = make_object<Person>(std::move(first), 38);
}
\end{lstlisting}

To edukacyjna wersja fabryki. W produkcyjnym kodzie
\texttt{std::make\_unique<T>(...)} robi to samo w prostszy
sposób. W drugim wywołaniu użytkownik świadomie pozwolił na
przeniesienie \texttt{first}; w pierwszym pozostaje on
nienaruszony. \texttt{std::forward<Args>(args)...} rozwija pakiet
argumentów omówiony w rozdziale~4.

\subsection{Nie przekazuj tego samego zasobu dwa razy bez zastanowienia}

\begin{lstlisting}
#include <utility>

template<class T>
void bad_double_forward(T&& x) {
    consume(std::forward<T>(x));
    consume(std::forward<T>(x)); // may use moved-from object
}
\end{lstlisting}

Ten fragment jest \emph{schematem błędu}: zakładamy, że istnieje
funkcja \texttt{consume}. Jeśli pierwsze wywołanie wykorzysta
argument rvalue przez przeniesienie, drugie otrzyma już obiekt
po przeniesieniu. Referencję przekazującą traktuj zatem jako
przeznaczoną zwykle do \emph{jednego} przekazania dalej.

\section{Perfect forwarding a przeciążenia i szablony klas}

Konstruktor z referencją przekazującą potrafi być bardzo zachłanny
podczas rozstrzygania przeciążeń. Może np. wygrać z konstruktorem
kopiującym przy argumentach nie-\texttt{const}.

\begin{lstlisting}
#include <string>
#include <utility>

struct Named {
    std::string name;

    Named(const Named&) = default;

    template<class U>
    explicit Named(U&& u)
        : name(std::forward<U>(u)) {}
};

// Named a{"text"};
// Named b(a);  // may select the constructor template,
//              // then fail to initialize std::string from Named&
\end{lstlisting}

Zwykle lepiej zacząć od zwykłych przeciążeń albo konstruktora
przyjmującego konkretny typ. Jeżeli szablon jest potrzebny,
można ograniczyć go koncepcją:

\begin{lstlisting}
#include <concepts>
#include <string>
#include <type_traits>
#include <utility>

struct SafeNamed {
    std::string name;

    SafeNamed(const SafeNamed&) = default;
    SafeNamed(SafeNamed&&) = default;

    template<class U>
      requires (!std::same_as<std::remove_cvref_t<U>, SafeNamed>
                && std::constructible_from<std::string, U&&>)
    explicit SafeNamed(U&& u)
        : name(std::forward<U>(u)) {}
};
\end{lstlisting}

Takie powiązanie rozdziału o perfect forwarding z koncepcjami
z rozdziału~5 jest bardzo przydatne
w bibliotekach generycznych.

\section{Copy elision: kiedy nie ma ani kopii, ani przeniesienia?}
\label{sec:move-elision}

Prosty model \emph{zwrócenie przez wartość = kopia lub move}
jest niepoprawny w nowoczesnym C++. Od C++17 pewne wyrażenia
\texttt{prvalue} inicjalizują wynik \emph{bezpośrednio w miejscu
docelowym}. Nie istnieje dodatkowy obiekt, który trzeba by
skopiować lub przenieść.

\begin{lstlisting}
struct NonMovable {
    NonMovable() = default;
    NonMovable(const NonMovable&) = delete;
    NonMovable(NonMovable&&) = delete;
};

NonMovable create() {
    return NonMovable{}; // guaranteed elision, C++17+
}

int main() {
    NonMovable x = create(); // OK in C++17+
}
\end{lstlisting}

Nie jest to zwykła optymalizacja, którą kompilator \emph{może}
wyłączyć: w pokazanym przypadku wynika to z semantyki języka.

\subsection{NRVO a gwarantowana elizja}

\begin{lstlisting}
#include <string>

std::string first() {
    return std::string{"hello"}; // mandatory direct construction
}

std::string second() {
    std::string s = "hello";
    return s;                    // NRVO may occur
}
\end{lstlisting}

\textit{Named Return Value Optimization} (NRVO) dotyczy
zwracania nazwanej lokalnej zmiennej. W C++20 kompilator \emph{może}
pominąć utworzenie dodatkowego obiektu. Jeśli NRVO nie nastąpi,
wyrażenie \texttt{s} w instrukcji \texttt{return} może
uczestniczyć w niejawnej próbie przeniesienia.

\paragraph{Nie pisz automatycznie \texttt{return std::move(s);}.}

\begin{lstlisting}
std::string make_text() {
    std::string s = "result";
    return s;             // allows NRVO
}

std::string less_good() {
    std::string s = "result";
    return std::move(s);  // inhibits NRVO; move may be required
}
\end{lstlisting}

\texttt{std::move(s)} zmienia wynik w xvalue i tym samym
nie spełnia warunków NRVO dotyczących nazwy obiektu.
Jeśli funkcja zwraca własną lokalną zmienną przez wartość,
zazwyczaj najlepsze jest \texttt{return s;}.

\section{Argumenty funkcji: kiedy kopiować, kiedy przenosić?}
\label{sec:move-api}

Wybór sposobu przyjmowania argumentu jest częścią kontraktu API.
Nie chodzi jedynie o wydajność, ale przede wszystkim o własność
i intencję.

\begin{center}
\begin{tabular}{|p{3.4cm}|p{4.3cm}|p{4.7cm}|}
\hline
\textbf{Parametr} & \textbf{Znaczenie} & \textbf{Kiedy używać} \\
\hline
\texttt{const T\&} & Odczyt bez przejmowania własności. & Duże obiekty, gdy wystarczy tylko odczyt. \\
\hline
\texttt{T\&} & Modyfikacja obiektu należącego do wywołującego. & Parametry wyjściowe i operacje in-place. \\
\hline
\texttt{T} & Własna wartość; może być skopiowana albo przeniesiona. & Małe typy i funkcje mające zachować wartość. \\
\hline
\texttt{T\&\&} & Dostęp do argumentu będącego r-wartością. & Przeciążenia zależne od kategorii oraz operacje konsumujące. \\
\hline
\texttt{T\&\&} z dedukcją \texttt{T} & Referencja przekazująca. & Opakowania i fabryki zachowujące kategorię argumentów. \\
\hline
\end{tabular}
\end{center}

\subsection{Idiom \emph{pass by value, then move}}

\begin{lstlisting}
#include <string>
#include <utility>

class User {
    std::string name_;
public:
    explicit User(std::string name)
        : name_(std::move(name)) {}
};

std::string a = "Alice";
User first(a);                    // copy into parameter, then move
User second(std::string{"Bob"});  // parameter from prvalue, then move
\end{lstlisting}

To dobre rozwiązanie, gdy funkcja \emph{zawsze} potrzebuje własnej
kopii/przeniesienia argumentu. Nie jest jednak magicznie
najszybsze w każdym przypadku: dla lvalue może wymagać kopiowania
i jeszcze jednego przeniesienia, a duże obiekty niemożliwe do
przeniesienia mogą wymagać osobnego projektu API.

\subsection{Przekazywanie własności \texttt{unique\_ptr}}

\begin{lstlisting}
#include <memory>
#include <utility>

void adopt(std::unique_ptr<int> p) {
    // Ownership transferred into this function.
}

int main() {
    auto p = std::make_unique<int>(42);
    adopt(std::move(p));
    // p == nullptr after moving this unique_ptr
}
\end{lstlisting}

Przekazanie \texttt{unique\_ptr} przez wartość oznacza przejęcie
własności. Jeżeli funkcja ma jedynie korzystać z obiektu,
zwykle lepiej przyjąć \texttt{T\&} lub \texttt{const T\&}.
Temat własności zasobów omówiono w rozdziale~3.

\section{Rvalue-qualified methods i świadome udostępnianie zasobów}
\label{sec:move-ref-qualified}

C++ pozwala ograniczać metodę do lvalue albo rvalue odbiorcy,
dopisując \texttt{\&} lub \texttt{\&\&} po liście parametrów.
Można to wykorzystać np. do udostępnienia zawartości obiektu
bez kopiowania, gdy całe opakowanie jest tymczasowe.

\begin{lstlisting}
#include <string>
#include <utility>

class Message {
    std::string data_;
public:
    explicit Message(std::string s) : data_(std::move(s)) {}

    const std::string& text() const & noexcept {
        return data_;              // lvalue owner: borrow
    }
    std::string text() && noexcept {
        return std::move(data_);   // rvalue owner: consume
    }
};

Message m{"hello"};
const std::string& borrowed = m.text();
std::string owned = Message{"world"}.text();
\end{lstlisting}

Wariant \texttt{\&\&} może bezpiecznie oddać wartość przez
przeniesienie. Gdybyśmy zwracali w nim referencję do składowej
tymczasowego obiektu, powstałaby referencja wisząca.
Projektując takie API, warto rozważyć również wywołania na
\texttt{const} r-wartościach i ich wymagane zachowanie.

\section{Kontenery, iteratory przenoszące i wzorce użytkowe}

\subsection{Dodawanie elementów do \texttt{std::vector}}

\begin{lstlisting}
#include <string>
#include <utility>
#include <vector>

std::vector<std::string> words;
std::string word = "one";

words.push_back(word);             // copy from lvalue
words.push_back(std::move(word));  // move from lvalue cast to xvalue
words.emplace_back(3, 'x');       // construct "xxx" in place
\end{lstlisting}

Metoda \texttt{emplace\_back} konstruuje nowy element z podanych
argumentów. Nie musi być szybsza od \texttt{push\_back}.
Gotowy obiekt można przekazać z pomocą \texttt{std::move}.
Podczas realokacji kontener nadal może kopiować lub przenosić
wcześniej dodane elementy.

\subsection{Przenoszenie całego zakresu}

\begin{lstlisting}
#include <algorithm>
#include <iterator>
#include <string>
#include <vector>

std::vector<std::string> source{"a", "b"};
std::vector<std::string> dest;

std::move(source.begin(), source.end(),
          std::back_inserter(dest)); // algorithm std::move

// source still contains two valid strings, but their
// contents are not assumed.
\end{lstlisting}

Nazwa \texttt{std::move} występuje również dla algorytmu z nagłówka
\texttt{<algorithm>}. W przeciwieństwie do pojedynczego
\texttt{std::move(x)} z nagłówka \texttt{<utility>}, algorytm
przypisuje kolejno elementy z zakresu docelowego, korzystając
z możliwości ich przenoszenia.

\section{Typowe błędy i pułapki}
\label{sec:move-errors}

\subsection{Błąd 1: przekonanie, że \texttt{std::move} niszczy obiekt}

\texttt{std::move} jest rzutowaniem, a nie destruktorem. Obiekt
źródłowy nadal istnieje aż do końca własnego czasu życia.
Może zmienić stan, jeśli na jego podstawie wykonano
rzeczywistą operację przeniesienia.

\subsection{Błąd 2: \texttt{std::move} z argumentu przekazanego jako \texttt{const T\&}}

\begin{lstlisting}
#include <string>
#include <utility>

std::string ineffective(const std::string& x) {
    return std::move(x); // copies for std::string
}
\end{lstlisting}

Jeżeli funkcja potrzebuje zwrócić niezależną kopię, czytelniej
napisać \texttt{return x;}. Jeśli ma konsumować argument,
powinna otrzymać wartość lub odpowiednią modyfikowalną
referencję. Nie wolno usuwać \texttt{const} przez rzutowanie,
by wymusić przeniesienie z oryginalnie stałego obiektu.

\subsection{Błąd 3: użycie nazwanego \texttt{T\&\&} bez \texttt{std::move}}

\begin{lstlisting}
#include <memory>
#include <utility>

void take(std::unique_ptr<int>&& p) {
    // auto q = p;            // error: p is an lvalue
    auto q = std::move(p);    // explicit ownership transfer
}
\end{lstlisting}

\subsection{Błąd 4: przechowywanie referencji do argumentu tymczasowego}

\begin{lstlisting}
#include <string>
#include <string_view>

std::string_view dangerous() {
    return std::string{"temporary"}; // dangling view
}
\end{lstlisting}

\texttt{std::string\_view} jest widokiem, a nie właścicielem
pamięci. Semantyka przenoszenia nie rozwiązuje problemów
związanych z czasem życia. Gdy wynik funkcji ma być własnością
wywołującego, zwracaj \texttt{std::string} przez wartość.

\subsection{Błąd 5: przenoszenie nie gwarantuje zachowania wskaźników}

Przeniesienie kontenera może spowodować unieważnienie wcześniej
zachowanych iteratorów, referencji lub wskaźników zależnie od
typu kontenera, operacji i alokatora. Nie zakładaj ogólnej
zasady, że po \texttt{std::move} adresy danych zawsze pozostają
takie same. Sprawdź dokumentację odpowiedniego typu.

\subsection{Błąd 6: nadmierne stosowanie \texttt{std::forward}}

\texttt{std::forward<T>} ma sens głównie tam, gdzie \texttt{T}
pochodzi z dedukcji dla referencji przekazującej. Przy zwykłej
lokalnej zmiennej przeznaczonej do przekazania własności
należy używać \texttt{std::move}; przy zwracaniu lokalnej
zmiennej przez wartość zazwyczaj żadnego z nich.

\section{Przykład przekrojowy: kolejka zadań}
\label{sec:move-project}

Rozważmy prosty bufor zadań przechowujący funkcje do
późniejszego wykonania. Użyjemy \texttt{std::unique\_ptr},
aby zadanie jednoznacznie posiadało zasób. Poniższy przykład
jest niezależny od wielowątkowości; \emph{kolejka zadań}
nie uruchamia niczego w osobnym wątku.

\begin{lstlisting}
#include <functional>
#include <iostream>
#include <memory>
#include <utility>
#include <vector>

class Tasks {
    std::vector<std::function<void()>> items_;
public:
    template<class F>
    void add(F&& f) {
        items_.emplace_back(std::forward<F>(f));
    }

    void run() {
        for (auto& f : items_) f();
        items_.clear();
    }
};

int main() {
    Tasks tasks;
    auto value = std::make_shared<int>(10);
    tasks.add([value] { std::cout << *value << '\n'; });
    tasks.run();
}
\end{lstlisting}

\texttt{std::function<void()>} w C++20 wymaga, aby przechowywany
obiekt wywoływalny był kopiowalny. Dlatego w przykładzie lambda
przechowuje \texttt{shared\_ptr}. Gdybyśmy zamiast niego
przechwycili przez przeniesienie \texttt{unique\_ptr}, lambda
stałaby się niekopiowalna i taki obiekt nie spełniłby
wymagań \texttt{std::function} w C++20. C++23 wprowadza
\texttt{std::move\_only\_function}, która rozwiązuje ten problem.
To pokazuje, że prawidłowy wybór sposobu przenoszenia zależy
od kontraktów pozostałych typów.

\section{Praktyczna ściąga i zalecenia projektowe}

\begin{center}
\begin{tabular}{|p{4.0cm}|p{8.4cm}|}
\hline
\textbf{Sytuacja} & \textbf{Zalecenie} \\
\hline
Zwracam utworzony obiekt & Zwróć go przez wartość; korzystaj z elizji lub NRVO. \\
\hline
Przenoszę zasób z lokalnego obiektu & Użyj \texttt{std::move(x)} w miejscu przekazania własności. \\
\hline
Piszę wrapper o argumentach \texttt{Args\&\&...} & Wykorzystaj \texttt{std::forward<Args>(args)...}. \\
\hline
Klasa ma \texttt{vector}/\texttt{string} jako składowe & Preferuj regułę zera i generowane operacje. \\
\hline
Klasa posiada goły wskaźnik do alokacji & Preferuj zmianę projektu na RAII; jeśli zarządzasz nim ręcznie, rozważ wszystkie operacje specjalne. \\
\hline
Kontener przenosi elementy przy realokacji & Właściwy, bezpieczny \texttt{noexcept} w konstruktorze move pomaga w gwarancjach wyjątkowych. \\
\hline
Argument będzie tylko odczytywany & Zwykle \texttt{const T\&} dla dużego \texttt{T}, a małe typy przez wartość. \\
\hline
Funkcja przejmuje \texttt{unique\_ptr} & Przyjmij go przez wartość; wywołujący świadomie używa \texttt{std::move}. \\
\hline
\end{tabular}
\end{center}

\paragraph{Najważniejsze rozróżnienie.}
\begin{enumerate}
\item \texttt{std::move} wyraża zgodę na potraktowanie obiektu jako źródła przeniesienia, ale samo zasobów nie przenosi.
\item \texttt{std::forward<T>} zachowuje kategorię argumentu dedukowanego w szablonie.
\item Konstruktor przenoszący to zwykle operacja \texttt{T(T\&\&)}, ale może wykonywać dowolną operację zgodną z kontraktem typu.
\item \texttt{T\&\&} jako typ parametru nie zawsze jest referencją przekazującą.
\item Nazwany parametr typu \texttt{T\&\&} jest wyrażeniem \texttt{lvalue}.
\item Zwracanie przez wartość jest często najlepszym sposobem przekazywania wyniku.
\item Kiedy tylko się da, stosuj regułę zera zamiast własnych pięciu funkcji specjalnych.
\end{enumerate}

\paragraph{Pytania kontrolne.}
\begin{enumerate}
\item Dlaczego wywołanie \texttt{g(x)} w funkcji \texttt{f(T\&\& x)} traktuje \texttt{x} jako lvalue?
\item Kiedy \texttt{std::move(x)} kończy się kopiowaniem zamiast przeniesienia?
\item Czym różnią się \texttt{return T\{\};} i \texttt{return local;} w C++20?
\item Dlaczego przy powiększaniu \texttt{vector} konstruktor move z \texttt{noexcept} ma znaczenie?
\item Kiedy użyć \texttt{std::forward<T>(x)}, a kiedy \texttt{std::move(x)}?
\item Dlaczego metoda zwracająca \texttt{const std::string\&} nie powinna ujawniać referencji do składowej obiektu tymczasowego?
\end{enumerate}

\paragraph{Dalsza lektura.} \emph{cppreference}: \textit{Value categories},
\textit{Move constructors}, \textit{Reference declaration},
\textit{Copy elision}, \textit{std::move}, \textit{std::forward}.

% ROZDZIAL 11 — Wyjatki i bezpieczenstwo wyjatkowe w C++20
% Fragment rozdzialu do skryptu. Wymagane pakiety: listings, amsmath, array, hyperref.
% Uzycie: \input{rozdzial_11_wyjatki_bezpieczenstwo_cpp.tex}

\chapter{Wyjątki i bezpieczeństwo wyjątkowe w C++}
\label{chap:cpp-exceptions}

\section{Po co C++ ma wyjątki?}

Mechanizm wyjątków oddziela kod realizujący zwykły przebieg algorytmu
od kodu reagującego na niepowodzenia. Jeżeli funkcja nie może spełnić
swojego kontraktu, może zgłosić wyjątek poleceniem \texttt{throw}.
Wówczas wykonanie jest przekazywane do odpowiedniego bloku
\texttt{catch}; po drodze niszczone są obiekty automatyczne, których
czas życia dobiegł końca. To ostatnie zjawisko nazywamy
\emph{odwijaniem stosu} (\emph{stack unwinding}).

Najprostszy schemat wygląda tak: \texttt{throw} zgłasza problem,
a pasujący \texttt{catch} go odbiera. Kod po \texttt{throw} w tym samym
bloku nie zostanie wykonany:
\begin{lstlisting}
try {
    throw std::runtime_error("bad input");
} catch (const std::runtime_error& error) {
    std::cerr << error.what() << '\n';
}
\end{lstlisting}
W samodzielnym programie ten fragment wymaga nagłówków
\texttt{<iostream>} i \texttt{<stdexcept>}.

Z perspektywy projektowania API warto odróżniać:
\begin{itemize}
\item \emph{zwykły wynik} (w tym informację, że elementu nie znaleziono),
      dla którego często lepszy jest \texttt{std::optional};
\item \emph{błąd operacji, który wywołujący powinien świadomie obsłużyć},
      reprezentowany np. przez kod błędu lub, od C++23, \texttt{std::expected};
\item \emph{niepowodzenie uniemożliwiające dotrzymanie kontraktu},
      dla którego wyjątek bywa naturalnym sposobem propagacji.
\end{itemize}
Nie ma uniwersalnej reguły nakazującej używanie wyjątków wszędzie.
Znaczenie mają kontrakty funkcji, konwencje projektu, ograniczenia
środowiska oraz wymagania dotyczące opóźnień wykonania.

\section{Składnia \texttt{try}, \texttt{throw} i \texttt{catch}}

\begin{center}\small
\begin{tabularx}{\textwidth}{@{}p{0.32\textwidth}X@{}}
\toprule
Konstrukcja & Za co odpowiada \\
\midrule
\code{try} & Obejmuje blok, z którego wyjątki mogą zostać przechwycone przez następujące bloki \code{catch}. \\
\code{throw value} & Tworzy obiekt wyjątku i rozpoczyna poszukiwanie pasującej obsługi. \\
\texttt{catch(const E\& e)} & Przechwytuje wyjątek zgodnego typu bez kopiowania i bez obcinania części pochodnej. \\
\code{catch(...)} & Przechwytuje dowolny wyjątek, także nieznanego typu. \\
\code{throw;} & Wewnątrz obsługi ponownie zgłasza bieżący wyjątek, zachowując jego typ dynamiczny. \\
\code{std::exception::what()} & Zwraca tekstowy opis błędu; służy do diagnostyki, nie do sterowania logiką. \\
\bottomrule
\end{tabularx}
\end{center}
\label{sec:cpp11-basics}

\begin{lstlisting}
#include <iostream>
#include <stdexcept>

double divide(double a, double b) {
    if (b == 0.0)
        throw std::invalid_argument("zero denominator");
    return a / b;
}

int main() {
    try {
        std::cout << divide(10.0, 2.0) << '\n';
        std::cout << divide(10.0, 0.0) << '\n';
        std::cout << "not executed\n";
    } catch (const std::invalid_argument& e) {
        std::cerr << "invalid input: " << e.what() << '\n';
    }
}
\end{lstlisting}

Instrukcja \texttt{throw E\{...\}} tworzy obiekt wyjątku. Obsługa
szuka pasującej klauzuli \texttt{catch}; jeśli wyjątek nie zostanie
obsłużony, program kończy działanie poprzez \texttt{std::terminate}.
Nie należy zakładać, że nieobsłużony wyjątek zawsze spowoduje pełne
odwinięcie stosu: szczegóły zakończenia programu mogą zależeć od
implementacji.

\subsection{Przechwytywanie przez referencję i kolejność klauzul}

Przyjmujemy zwykle \texttt{const E\&}: pozwala to uniknąć niepotrzebnego
kopiowania i \emph{object slicing} przy przekazywaniu obiektu typu
pochodnego jako typu bazowego przez wartość.

\begin{lstlisting}
#include <exception>
#include <iostream>
#include <stdexcept>

void handle() {
    try {
        throw std::out_of_range("bad index");
    } catch (const std::out_of_range& e) {
        std::cout << "out_of_range: " << e.what() << '\n';
    } catch (const std::exception& e) {
        std::cout << "other std::exception: " << e.what() << '\n';
    } catch (...) {
        std::cout << "unknown exception\n";
    }
}
\end{lstlisting}

Klauzule są sprawdzane w kolejności występowania: najpierw umieszcza
się typy bardziej szczegółowe, a na końcu bazowe. Wyrażenie
\texttt{catch (...)} przechwytuje wyjątek dowolnego typu, ale nie
udostępnia bezpośrednio jego obiektu.

\paragraph{Ponowne zgłoszenie.}
\texttt{throw;} wewnątrz aktywnej obsługi ponownie zgłasza bieżący
wyjątek z zachowaniem jego dynamicznego typu. Natomiast
\texttt{throw e;} zgłasza nowe wyrażenie; w zależności od statycznego
typu \texttt{e} może prowadzić do kopiowania i obcięcia do klasy
bazowej.

\begin{lstlisting}
void log_and_rethrow() {
    try {
        throw std::runtime_error("disk error");
    } catch (const std::exception& e) {
        std::cerr << e.what() << '\n';
        throw;                  // keeps original exception object
        // throw e;              // may slice to std::exception
    }
}
\end{lstlisting}

\section{Odwijanie stosu i RAII}
\label{sec:cpp11-unwinding}

Najważniejsza reguła praktyczna brzmi: \emph{zwalnianie zasobu powinno
wynikać z destruktora obiektu zarządzającego zasobem}. Stanowi to
rozwinięcie idei RAII poznanej w rozdziale~\ref{chap:wskazniki-cpp}.

\begin{lstlisting}
#include <iostream>
#include <stdexcept>
#include <string>
#include <utility>

struct Trace {
    std::string name;
    explicit Trace(std::string n) : name(std::move(n)) {
        std::cout << "create " << name << '\n';
    }
    ~Trace() noexcept {
        std::cout << "destroy " << name << '\n';
    }
};

void third() {
    Trace c{"C"};
    throw std::runtime_error("failure");
}
void second() {
    Trace b{"B"};
    third();
}
void first() {
    Trace a{"A"};
    second();
}
// try { first(); } catch (const std::exception&) {}
// Destruction order: C, B, A.
\end{lstlisting}

Wyjątek powoduje zakończenie odpowiednich zakresów, a obiekty
z lokalnym automatycznym czasem przechowywania są niszczone w odwrotnej
kolejności ich skonstruowania. Jeżeli konstruktor zgłosi wyjątek,
destruktor \emph{całego, nieutworzonego obiektu} nie jest wywoływany,
ale niszczone są już w pełni skonstruowane podobiekty bazowe i składowe.

\subsection{Zarządzanie plikiem przy pomocy RAII}

\begin{lstlisting}
#include <cstdio>
#include <memory>
#include <stdexcept>

struct FileCloser {
    void operator()(std::FILE* file) const noexcept {
        if (file) std::fclose(file);
    }
};
using File = std::unique_ptr<std::FILE, FileCloser>;

File open_file(const char* name) {
    File file{std::fopen(name, "rb")};
    if (!file) throw std::runtime_error("cannot open file");
    return file;
}
\end{lstlisting}

Nawet gdy po otwarciu pliku dalsze operacje zgłoszą wyjątek,
\texttt{unique\_ptr} wykona deleter. Zauważmy, że
\texttt{std::fclose} może sygnalizować błąd przy zamykaniu; prosty
noexcept-deleter ignoruje ten kod błędu. Jeżeli informacja o błędzie
zamknięcia jest istotna, należy zapewnić osobną, jawną operację
finalizacji, której wynik można obsłużyć \emph{przed} destrukcją.

\section{Gwarancje bezpieczeństwa wyjątkowego}

\begin{center}\small
\begin{tabularx}{\textwidth}{@{}p{0.32\textwidth}X@{}}
\toprule
Gwarancja lub narzędzie & Skutek w razie wyjątku \\
\midrule
Brak wycieku & Zasoby pozostają poprawnie zwalniane, zwykle przez RAII. \\
Gwarancja podstawowa & Obiekty pozostają poprawne, choć ich stan może się zmienić. \\
Gwarancja silna & Operacja kończy się sukcesem albo pozostawia obserwowalny stan bez zmian. \\
Gwarancja braku wyjątków & Operacja nie pozwala wyjątkowi opuścić funkcji. \\
\code{noexcept} & Deklaruje gwarancję braku wyjątku; wyjątek opuszczający funkcję kończy program przez \code{std::terminate}. \\
\bottomrule
\end{tabularx}
\end{center}
\label{sec:cpp11-guarantees}

Bezpieczeństwo wyjątkowe nie oznacza tylko braku wycieków pamięci.
Załóżmy, że lista zawiera \texttt{[1, 2]} i dopisanie \texttt{3}
zgłosi wyjątek. Gwarancja podstawowa oznacza, że lista nadal nadaje
się do użycia, choć może mieć zmienioną zawartość. Gwarancja silna
oznacza, że nadal zawiera dokładnie \texttt{[1, 2]}. Gwarancja braku
wyjątków mówi tylko, że wyjątek nie opuści funkcji; jeśli operacja
ma dodatkowo zawsze osiągnąć cel, mówimy o gwarancji \emph{no-fail}.

Silna gwarancja nie oznacza automatycznie cofnięcia efektów
zewnętrznych, np. wysłanego pakietu czy wpisu w bazie danych.
Zakres transakcyjności trzeba sprecyzować w kontrakcie.

\subsection{Silna gwarancja dzięki operacji „przygotuj, a potem zatwierdź”}

\begin{lstlisting}
#include <utility>
#include <vector>

class Table {
    std::vector<int> data_;
public:
    explicit Table(std::vector<int> values)
        : data_(std::move(values)) {}

    void append_many(const std::vector<int>& more) {
        auto candidate = data_;          // may throw
        candidate.insert(candidate.end(),
                         more.begin(), more.end()); // may throw
        data_.swap(candidate);           // noexcept for std::vector<int>
    }

    const std::vector<int>& data() const noexcept { return data_; }
};
\end{lstlisting}

Jeżeli kopiowanie lub \texttt{insert} zawiedzie, \texttt{data\_}
pozostaje nietknięte. Dopiero bezwyjątkowe \texttt{swap} zatwierdza
przygotowaną wersję. Kosztem jest dodatkowa pamięć i kopiowanie;
silna gwarancja bywa więc świadomym wyborem projektowym, a nie darmową
optymalizacją.

\paragraph{Uwaga na aliasing.}
W przykładzie argument \texttt{more} może nawet odnosić się do
\texttt{data\_} obiektu \texttt{Table}: kopia \texttt{candidate}
powstaje przed modyfikacją oryginału. W kodzie operującym bezpośrednio
na tym samym kontenerze trzeba dodatkowo analizować ważność iteratorów
i referencji.

\subsection{Konstruktor i niezmienniki klasy}

Dobrze napisany konstruktor tworzy w pełni poprawny obiekt albo
zgłasza wyjątek. Bezpiecznie zarządzane składowe automatycznie
posprzątają po częściowo wykonanej inicjalizacji.

\begin{lstlisting}
#include <memory>
#include <stdexcept>

class Session {
    std::unique_ptr<int[]> buffer_;
    std::size_t size_;
public:
    explicit Session(std::size_t n)
        : buffer_(std::make_unique<int[]>(n)), size_(n) {
        if (n == 0) throw std::invalid_argument("empty session");
    }
    std::size_t size() const noexcept { return size_; }
};
\end{lstlisting}

Nie należy zastępować RAII ręcznie pisanymi \texttt{catch}, w których
sprząta się kolejne zasoby. Takie procedury są podatne na błędy,
zwłaszcza gdy konstruktor zdobywa wiele zasobów po kolei.

\section{\texttt{noexcept}: deklarowanie i sprawdzanie}
\label{sec:cpp11-noexcept}

\texttt{noexcept} w specyfikacji funkcji oznacza, że wyjątek nie może
opuścić jej ciała. Naruszenie tego wymagania kończy się wywołaniem
\texttt{std::terminate} (nie jest to „ciche zignorowanie” błędu).

\begin{lstlisting}
void close_handle() noexcept {
    // cleanup that must not let exceptions escape
}

void may_fail();
static_assert(noexcept(close_handle()));
static_assert(!noexcept(may_fail()));
\end{lstlisting}

W drugim użyciu \texttt{noexcept(expr)} jest \emph{operatorem}
zwracającym wartość logiczną w czasie kompilacji. Nie uruchamia on
samego wyrażenia.

\subsection{Warunkowe \texttt{noexcept}}

\begin{lstlisting}
#include <type_traits>
#include <utility>

template<class T>
void swap_values(T& a, T& b)
    noexcept(std::is_nothrow_move_constructible_v<T> &&
             std::is_nothrow_move_assignable_v<T>) {
    T tmp(std::move(a));
    a = std::move(b);
    b = std::move(tmp);
}
\end{lstlisting}

Jest to przykład \emph{warunkowej specyfikacji}. Nie oznacza,
że przedstawiony algorytm daje silną gwarancję, jeżeli przenoszenie
rzuca wyjątki; dla takich typów stan \texttt{a} lub \texttt{b} może
się zmienić przed niepowodzeniem. W kodzie generycznym trzeba
analizować nie tylko składnię, ale też gwarancje operacji typu \texttt{T}.

\subsection{Dlaczego \texttt{noexcept} jest ważne dla \texttt{vector}?}

Przy powiększaniu pojemności kontener może potrzebować przenieść
lub skopiować stare elementy do nowej pamięci. Dla typów kopiowalnych
konstrukcja kopiująca bywa bezpieczniejsza, jeśli konstruktor
przenoszący może zgłosić wyjątek.

\begin{lstlisting}
#include <type_traits>
#include <utility>

struct Item {
    Item() = default;
    Item(const Item&) {}              // copying available
    Item(Item&&) noexcept(false) {}   // move may throw
};

static_assert(std::is_copy_constructible_v<Item>);
static_assert(!std::is_nothrow_move_constructible_v<Item>);

// std::move_if_noexcept(x) usually gives const Item& here.
\end{lstlisting}

\texttt{std::move\_if\_noexcept} zwraca \texttt{T\&\&}, jeżeli
przenoszenie jest zadeklarowane jako bezwyjątkowe albo typ nie jest
kopiowalny; w przeciwnym razie zwraca \texttt{const T\&}.
Zasada ta pomaga kontenerom utrzymywać silną gwarancję w wielu,
ale nie we wszystkich przypadkach. Szczegóły są zależne od
konkretnej operacji i własności typu elementów.

\subsection{Destruktory i \texttt{std::terminate}}

Domyślnie destruktor jest niewyrzucający, o ile destruktory jego
składowych i klas bazowych nie wymagają innej specyfikacji.
Szczególnie podczas odwijania stosu wypuszczenie kolejnego wyjątku
z destruktora może spowodować \texttt{std::terminate}.

\begin{lstlisting}
struct SafeCleanup {
    ~SafeCleanup() noexcept {
        try {
            risky_cleanup();
        } catch (...) {
            // Log, record an error, or use another no-throw fallback.
            // Never let it escape this destructor.
        }
    }
    static void risky_cleanup();
};
\end{lstlisting}

Nie oznacza to, że każdą awarię należy ukryć. Jeśli operacja
finalizacji może się nie udać i użytkownik musi znać wynik, projektujmy
jawną metodę \texttt{close()} lub \texttt{commit()}, a destruktor
niech pełni rolę bezpiecznej ścieżki awaryjnej.

\section{Własne klasy wyjątków i hierarchia standardowa}
\label{sec:cpp11-types}

W nagłówkach \texttt{<stdexcept>}, \texttt{<exception>},
\texttt{<new>} i \texttt{<system\_error>} znajdziemy wiele gotowych
typów. W typowych zastosowaniach najczęściej spotkamy
\texttt{std::invalid\_argument}, \texttt{std::out\_of\_range},
\texttt{std::runtime\_error}, \texttt{std::bad\_alloc} oraz
\texttt{std::system\_error}.

\begin{lstlisting}
#include <cstddef>
#include <stdexcept>
#include <string>
#include <utility>

class ParseError : public std::runtime_error {
    std::size_t position_;
public:
    ParseError(std::size_t pos, std::string message)
        : std::runtime_error(std::move(message)),
          position_(pos) {}

    std::size_t position() const noexcept { return position_; }
};

int parse_digit(char c, std::size_t pos) {
    if (c < '0' || c > '9')
        throw ParseError(pos, "expected decimal digit");
    return c - '0';
}
\end{lstlisting}

Własna klasa może przechowywać dodatkowe informacje, np. identyfikator
operacji, pozycję w tekście lub kod błędu. Nie należy umieszczać tam
niepotrzebnie ogromnych danych. Metoda \texttt{what()} klasy
\texttt{std::exception} jest wirtualna i \texttt{noexcept}; jeśli
ją przesłaniamy, musimy zachować zgodny kontrakt i poprawny czas życia
zwracanego łańcucha.

\paragraph{Błędy systemowe.}
\texttt{std::system\_error} umożliwia przenoszenie kodu błędu
\texttt{std::error\_code} wraz z kategorią, np. błędu systemowego,
bez uzależniania całego API od przypadkowego tekstu \texttt{what()}.

\section{Wyjątki w funkcjach szablonowych i neutralność wyjątkowa}

Komponent generyczny powinien zwykle zachować niezmienniki i pozwolić
wyjątkom operacji wywołanych na typie \texttt{T} przejść do klienta,
zamiast pochłaniać je przez \texttt{catch (...)}.

\begin{lstlisting}
#include <algorithm>
#include <vector>

template<class T, class Compare>
void sort_elements(std::vector<T>& xs, Compare comp) {
    std::sort(xs.begin(), xs.end(), comp);
    // An exception from comp may propagate to the caller.
}
\end{lstlisting}

Należy jednak odróżnić \emph{neutralność wyjątkową} od silnej
gwarancji. Jeśli komparator zgłosi wyjątek w trakcie sortowania,
kontener pozostaje poprawnym obiektem, ale jego kolejność nie musi
odpowiadać stanowi sprzed wywołania.

\section{Przenoszenie wyjątków między funkcjami i wątkami}
\label{sec:cpp11-exception-ptr}

Ta i dwie kolejne sekcje pokazują zastosowania bardziej zaawansowane.
Do podstaw wystarczą \texttt{try}/\texttt{catch}, RAII oraz gwarancje
stanu po niepowodzeniu.

Typ \texttt{std::exception\_ptr} pozwala zachować wyjątek poza blokiem
\texttt{catch}, a później zgłosić go ponownie. Jest przydatny
szczególnie w programowaniu współbieżnym.

\begin{lstlisting}
#include <exception>
#include <iostream>
#include <stdexcept>
#include <string>
#include <thread>

std::exception_ptr error;
std::thread worker([&] {
    try {
        throw std::runtime_error("worker failed");
    } catch (...) {
        error = std::current_exception();
    }
});
worker.join();

if (error) {
    std::rethrow_exception(error);
}
\end{lstlisting}

Zmienna \texttt{error} nie jest tutaj współdzielona jednocześnie przez
dwa aktywne dostępy: zapis zachodzi w wątku roboczym, a odczyt
dopiero po \texttt{join()}, które zapewnia wymaganą synchronizację.
Nie wolno dopuścić, aby wyjątek opuścił funkcję startową zwykłego
\texttt{std::thread}: typowym skutkiem jest \texttt{std::terminate}.

\subsection{\texttt{std::future} przechowuje wyjątek}

Wspomniane w rozdziale o współbieżności \texttt{std::async} i
\texttt{std::future} ułatwiają propagowanie niepowodzenia do
odbiorcy wyniku.

\begin{lstlisting}
#include <future>
#include <stdexcept>

auto fut = std::async(std::launch::async, []() -> int {
    throw std::runtime_error("calculation failed");
});
try {
    (void)fut.get();  // rethrows the stored exception
} catch (const std::runtime_error&) {
    // report the failure in the receiving thread
}
\end{lstlisting}

Przez \texttt{get()} wynik (lub wyjątek) trafia do wątku odbiorcy.
W \texttt{std::future} wynik można pobrać tylko raz; inne własności
tego mechanizmu omówiliśmy w rozdziale o asynchroniczności.

\section{Wyjątki zagnieżdżone i dodatkowy kontekst}
\label{sec:cpp11-nested}

Gdy funkcja wyższego poziomu rozumie więcej niż warstwa niższa, może
wzbogacić informację o błędzie i zachować oryginalną przyczynę.

\begin{lstlisting}
#include <exception>
#include <stdexcept>

void load_config() {
    try {
        throw std::runtime_error("invalid token");
    } catch (...) {
        std::throw_with_nested(
            std::runtime_error("cannot load configuration"));
    }
}

void print_nested(const std::exception& e, int depth = 0) {
    std::cerr << std::string(depth * 2, ' ') << e.what() << '\n';
    try {
        std::rethrow_if_nested(e);
    } catch (const std::exception& inner) {
        print_nested(inner, depth + 1);
    } catch (...) {
        std::cerr << "nested non-standard exception\n";
    }
}
\end{lstlisting}

Takie opakowywanie wyjątków nie jest tym
samym co zgłoszenie \texttt{throw e;}: nowy wyjątek przechowuje
kontekst, a przyczynę można odzyskać przez \texttt{rethrow\_if\_nested}.

\section{\texttt{std::uncaught\_exceptions} i strażnicy zakresu}

Funkcja \texttt{std::uncaught\_exceptions()} (C++17) zwraca liczbę
aktualnie nieobsługiwanych obiektów wyjątków w bieżącym wątku.
Można jej użyć do wykrycia, czy wyjście z zakresu następuje wskutek
odwijania stosu.

\begin{lstlisting}
#include <exception>

class FailureGuard {
    int initial_ = std::uncaught_exceptions();
    int& failure_count_;
public:
    explicit FailureGuard(int& count) noexcept : failure_count_(count) {}
    FailureGuard(const FailureGuard&) = delete;
    FailureGuard& operator=(const FailureGuard&) = delete;

    ~FailureGuard() noexcept {
        if (std::uncaught_exceptions() > initial_) {
            ++failure_count_; // no-throw bookkeeping
        }
    }
};
\end{lstlisting}

Porównanie z \emph{początkową liczbą}, a nie tylko z zerem, jest
ważne przy zagnieżdżonych wyjątkach. Strażnik nie daje jednak
pełnego, transakcyjnego mechanizmu cofania zmian. W projektach
korzystających z C++20 można napisać własny scope guard, ale trzeba
uważnie określić, czy przechowuje funkcję przez kopię, przez
przeniesienie czy przez referencję i czy jej wywołanie może rzucać.

\section{Wyjątki a alternatywne modele obsługi błędów}
\label{sec:cpp11-alternatives}

\subsection{\texttt{std::optional}: brak wyniku nie musi być błędem}

\begin{lstlisting}
#include <optional>
#include <string_view>

std::optional<std::size_t> find_colon(std::string_view s) {
    const auto i = s.find(':');
    if (i == std::string_view::npos) return std::nullopt;
    return i;
}
\end{lstlisting}

Brak dwukropka może być normalnym scenariuszem. Zgłaszanie wyjątku
przy każdym nieudanym wyszukiwaniu bywa niepotrzebne.

\subsection{\texttt{std::expected<T,E>} — C++23, nie C++20}

\texttt{std::expected} reprezentuje \emph{albo wartość typu T, albo
błąd typu E}. Został ustandaryzowany w C++23 i nie należy zakładać,
że nagłówek \texttt{<expected>} jest dostępny w kompilacji C++20.

\begin{lstlisting}
// This example requires C++23 and <expected>.
#include <expected>
#include <string>
#include <string_view>

std::expected<int, std::string>
parse_positive(std::string_view text) {
    if (text == "1") return 1;
    if (text == "2") return 2;
    return std::unexpected("expected 1 or 2");
}

void use_expected() {
    auto result = parse_positive("3");
    if (!result) {
        const std::string& error = result.error();
        (void)error;
    }
}
\end{lstlisting}

Dzięki temu błędy oczekiwane przez API są jawne w sygnaturze.
\texttt{std::expected} nie eliminuje wszystkich wyjątków: konstrukcja
zawartego obiektu czy inne operacje nadal mogą się nie udać.
W C++20 podobny kontrakt można zrealizować własnym typem
\texttt{Result}, \texttt{std::variant<T,E>} lub odpowiednią biblioteką.

\subsection{Kody błędów, wyjątki i wydajność}

W konwencjonalnych implementacjach C++ obsługa sukcesu na ścieżce,
na której żaden wyjątek nie wystąpił, może być bardzo tania, ale samo
zgłoszenie i obsłużenie wyjątku zwykle kosztuje więcej niż zwrócenie
prostego kodu błędu. Nie należy jednak zamieniać każdego \texttt{throw}
na kod zwrotny tylko na podstawie intuicji: warto mierzyć rzeczywistą
częstotliwość błędów, profilować aplikację i uwzględniać prostotę
kontraktów.

\section{Typowe błędy i dobre praktyki}
\label{sec:cpp11-guidelines}

\begin{enumerate}
\item \textbf{Nie rzucaj wskaźnikami do lokalnych obiektów.}
      Preferuj wyjątki będące obiektami, zwykle pochodnymi od
      \texttt{std::exception}.
\item \textbf{Przechwytuj wyjątki polimorficzne przez
      \texttt{const\&}.} Unikaj obcinania obiektu (\emph{slicing}).
\item \textbf{Do ponownego zgłoszenia używaj \texttt{throw;}.}
      Nie zastępuj go bez namysłu \texttt{throw e;}.
\item \textbf{Stosuj RAII.} Surowe wywołania \texttt{new}/\texttt{delete}
      rozdzielone operacjami rzucającymi wyjątki to częsta przyczyna
      wycieków pamięci.
\item \textbf{Nie wypuszczaj wyjątków z destruktorów.}
      Błędy finalizacji zgłaszaj jawnie w odpowiednio zaprojektowanym API.
\item \textbf{Nie oznaczaj funkcji \texttt{noexcept} bez analizy.}
      To obietnica z konsekwencją w postaci \texttt{std::terminate}.
\item \textbf{Nie ukrywaj błędów w pustym \texttt{catch (...)}.}
      Obsłuż wyjątek, zapisz informację diagnostyczną albo przekaż go wyżej.
\item \textbf{Określ gwarancję niepowodzenia.} Dla operacji mutujących
      dane rozważ silną gwarancję i strategię „przygotuj–zatwierdź”.
\item \textbf{Złap wyjątek w funkcji wątku.}
      Nie pozwól mu opuścić funkcji startowej \texttt{std::thread}.
      Do przekazania błędu użyj \texttt{std::exception\_ptr},
      \texttt{promise/future} lub API wyższego poziomu.
\item \textbf{Wybierz właściwy model błędu.} Brak zwykłego wyniku,
      przewidywalna porażka operacji i naruszenie kontraktu to różne
      sytuacje.
\end{enumerate}

\section{Przykład końcowy: transakcyjna aktualizacja danych}

Poniższy program ilustruje współpracę wyjątku, RAII i silnej gwarancji.
Załóżmy, że aktualizacja dodaje wartości do rejestru, ale w razie błędu
rejestr ma pozostać bez zmian.

\begin{lstlisting}
#include <iostream>
#include <stdexcept>
#include <utility>
#include <vector>

class Register {
    std::vector<int> values_;
public:
    Register() : values_{1, 2, 3} {}

    void add_checked(const std::vector<int>& numbers) {
        std::vector<int> proposed = values_;
        for (int x : numbers) {
            if (x < 0) {
                throw std::invalid_argument("negative input");
            }
            proposed.push_back(x);
        }
        values_.swap(proposed);
    }

    const std::vector<int>& values() const noexcept {
        return values_;
    }
};

int main() {
    Register r;
    try {
        r.add_checked({4, -5, 6});
    } catch (const std::invalid_argument& e) {
        std::cerr << e.what() << '\n';
    }
    // Still {1, 2, 3}; no partial update.
    for (int x : r.values()) std::cout << x << ' ';
}
\end{lstlisting}

Najpierw powstaje kopia danych. Sprawdzanie i alokowanie nowej pamięci
mogą zgłaszać wyjątki, lecz działają na kopii. Po pomyślnym wykonaniu
następuje zatwierdzenie przez \texttt{swap}. Jest to wzorzec, który
można stosować również w rozbudowanych klasach, przy świadomym
uwzględnieniu kosztów i własności ich składowych.

\section{Podsumowanie rozdziału}

W praktyce mechanizm wyjątków jest najbardziej użyteczny wtedy,
gdy łączy się go z dobrym modelem własności zasobów i z precyzyjnie
określonymi gwarancjami operacji. \texttt{throw}, \texttt{catch}
i \texttt{noexcept} to tylko składnia; o poprawności projektu
decydują RAII, niezmienniki obiektów, semantyka transakcyjna oraz
świadome oddzielanie błędów wyjątkowych od zwykłych wariantów wyniku.


% ROZDZIAL 10 — Organizacja pamieci i wydajnosc w C++20
% Fragment do wlaczenia po rozdziale 9.
% Wymagane pakiety: listings, array, amsmath, hyperref.
% Uzycie: \input{rozdzial_10_pamiec_wydajnosc_cpp.tex}

\chapter[Organizacja pamięci i wydajność]{Organizacja pamięci, wyrównanie i wydajność niskopoziomowa w C++}
\label{chap:cpp-memory-performance}

\section{Po co analizować organizację pamięci?}

W rozdziałach o wskaźnikach (rozdział~3), semantyce przenoszenia
(rozdział~6) i iteratorach (rozdział~8) interesowało nas głównie
\emph{poprawne zarządzanie obiektami}. Teraz przyjrzymy się temu, jak
obiekty są reprezentowane w pamięci oraz w jaki sposób wybory
programisty wpływają na liczbę alokacji i lokalność dostępu do danych.

Rozróżniajmy trzy poziomy opisu:
\begin{enumerate}
\item \textbf{Gwarancje języka C++:} rozmiar obiektu, wymagane wyrównanie,
      czas życia obiektu, legalność dostępu i zasady alokowania pamięci.
\item \textbf{Szczegóły implementacji:} rozmieszczenie pól, rozmiary
      wskaźników, zachowanie konkretnego alokatora czy układ binarny ABI.
\item \textbf{Właściwości sprzętu:} linie cache, hierarchia pamięci,
      prefetching, przepustowość i zjawisko \emph{false sharing}.
\end{enumerate}

Kod poprawny według standardu może być wolny; szybki kod działający
na jednym komputerze może być powolny na innym. Dlatego optymalizacja
wymaga najpierw \emph{pomiaru}, a następnie weryfikacji hipotezy.

\section{Obiekt, bajt, adres i czas życia}
\label{sec:cpp10-object}

\subsection{Pamięć nie jest tym samym co obiekt}

W modelu C++ \emph{obszar pamięci} może istnieć bez żywego obiektu typu
\texttt{T}. Obiekt charakteryzują między innymi typ, rozmiar, czas życia
i, dla pewnych klas, odpowiednio wywołane konstruktory i destruktory.
Samo uzyskanie adresu wystarczająco dużego bloku nie oznacza, że
wolno traktować ten blok jak dowolny obiekt C++.

\begin{lstlisting}
#include <cstddef>
#include <type_traits>

static_assert(sizeof(char) == 1);
static_assert(sizeof(std::byte) == 1);
static_assert(sizeof(int) >= 1);

struct Point { double x, y; };
static_assert(sizeof(Point) >= 2 * sizeof(double));
\end{lstlisting}

\texttt{sizeof(T)} podaje rozmiar obiektu w \emph{bajtach C++}.
Jeden bajt nie musi mieć ośmiu bitów; ich liczbę określa
\texttt{CHAR\_BIT} z nagłówka \texttt{<climits>}.
Rozmiar typu obiektowego jest dodatni. Każdy obiekt ma też
wymagane \emph{wyrównanie}, które omówimy dalej.

\paragraph{Czas przechowywania a czas życia.}
Należy odróżnić \emph{storage duration} od \emph{object lifetime}:
\begin{itemize}
\item obiekt lokalny ma zazwyczaj automatyczny czas przechowywania;
\item obiekt z \texttt{static} ma statyczny czas przechowywania;
\item obiekt \texttt{thread\_local} ma czas przechowywania wątku;
\item pamięć otrzymana przez \texttt{new} lub alokator jest pamięcią
      o dynamicznym czasie przechowywania.
\end{itemize}
Nie należy utożsamiać tych kategorii z gwarantowanymi fizycznymi
obszarami „stos” i „sterta”: te nazwy opisują zwyczajowe implementacje,
a nie pełny model abstrakcyjny języka.

\begin{lstlisting}
#include <memory>

int global = 17;                 // static storage duration
thread_local int per_thread = 0; // thread storage duration

void demo() {
    int local = 5;               // automatic storage duration
    auto p = std::make_unique<int>(10);
    // Obiekt *p zyje w pamieci zarzadzanej dynamicznie.
}
\end{lstlisting}

\subsection{Reprezentacja obiektu i bezpieczne oglądanie bajtów}

C++ pozwala badać reprezentację bajtową obiektu przez
\texttt{char}, \texttt{unsigned char} lub \texttt{std::byte}.
Nie wolno jednak wyciągać z tego wniosku, że dowolny wskaźnik
\texttt{T*} można przekonwertować na \texttt{U*} i odczytać dane jako
\texttt{U}: takie postępowanie może naruszyć reguły wyrównania,
czasu życia i dostępu przez glvalue odpowiedniego typu.

\begin{lstlisting}
#include <array>
#include <cstddef>
#include <cstring>
#include <cstdint>

std::uint32_t x = 0x12345678;
std::array<std::byte, sizeof(x)> bytes{};
std::memcpy(bytes.data(), &x, sizeof(x));

std::uint32_t y = 0;
std::memcpy(&y, bytes.data(), sizeof(y));
// y == x
\end{lstlisting}

Kolejność bajtów zależy od \emph{endianness}. Można ją rozpoznać przez
\texttt{std::endian} (C++20, nagłówek \texttt{<bit>}), jednak do
serializacji najlepiej jawnie ustalać porządek bajtów, format pól
i wersję protokołu. Nie należy wysyłać całej struktury przez
\texttt{write(fd, \&object, sizeof(object))} jako przenośnego formatu:
padding, reprezentacja liczb i ABI mogą się różnić.

\subsection{\texttt{std::bit\_cast} i ograniczenia rzutowań}

\texttt{std::bit\_cast<To>(from)} (C++20, \texttt{<bit>}) tworzy wartość
\texttt{To} przez skopiowanie reprezentacji bitowej, o ile oba typy
są trywialnie kopiowalne i mają równy rozmiar. Dodatkowo reprezentacja
musi oznaczać poprawną wartość docelowego typu.

\begin{lstlisting}
#include <bit>
#include <cstdint>

float f = 1.0f;
// Przyklad zaklada sizeof(float) == sizeof(uint32_t).
// Na typowej platformie wynik to kod bitowy IEEE 754.
if constexpr (sizeof(float) == sizeof(std::uint32_t)) {
    auto bits = std::bit_cast<std::uint32_t>(f);
    (void)bits;
}
\end{lstlisting}

Nie mylmy \texttt{bit\_cast} z \texttt{reinterpret\_cast}:
pierwszy wytwarza \emph{wartość}, drugi może tylko zmienić sposób
interpretacji wskaźnika; sam nie tworzy obiektu docelowego.

\section{Layout struktury: \texttt{sizeof}, \texttt{alignof}, padding}

\begin{center}\small
\begin{tabularx}{\textwidth}{@{}p{0.33\textwidth}X@{}}
\toprule
Narzędzie & Co pokazuje lub zmienia \\
\midrule
\code{sizeof(T)} & Podaje rozmiar obiektu typu \code{T} w bajtach, wraz z ewentualnym paddingiem. \\
\code{alignof(T)} & Podaje wymagane wyrównanie adresu obiektu typu \code{T}. \\
\code{offsetof(T, field)} & Podaje przesunięcie składowej w typie o standardowym układzie. \\
\code{alignas(n)} & Żąda wyrównania co najmniej do granicy \code{n} dla obiektu lub typu. \\
\code{std::byte} & Reprezentuje jednostkę surowej pamięci bez znaczenia znaku tekstowego. \\
\bottomrule
\end{tabularx}
\end{center}
\label{sec:cpp10-layout}

\subsection{Skąd bierze się padding?}

Każdy typ obiektowy ma wymaganie wyrównania. Kompilator może wstawić
\emph{padding}, czyli nieużywane bajty, aby kolejne pola miały
odpowiednie adresy, a elementy tablicy struktur były poprawnie
wyrównane. Rozmieszczenie pól zależy od typu klasy i reguł języka;
konkretna liczba bajtów paddingu zależy od implementacji.

\begin{lstlisting}
#include <cstddef>
#include <iostream>

struct Unordered {
    char flag;
    double value;
    int count;
};

struct Ordered {
    double value;
    int count;
    char flag;
};

std::cout << sizeof(Unordered) << '\n';
std::cout << sizeof(Ordered) << '\n';
std::cout << alignof(Ordered) << '\n';
\end{lstlisting}

Na częstym ABI, przy \texttt{sizeof(double)==8},
\texttt{alignof(double)==8} i \texttt{sizeof(int)==4},
pierwszy typ zajmuje zwykle 24 bajty, a drugi 16. \textbf{Nie są to
wartości wymagane przez standard}. Samo przestawienie pól może
zmniejszyć zapotrzebowanie na pamięć milionów rekordów, ale może też
pogorszyć czytelność, stabilność ABI lub lokalność dostępu.

\begin{center}
\small
\begin{tabular}{|l|p{4.4cm}|p{4.4cm}|}
\hline
 & \textbf{Unordered (przykładowo)} & \textbf{Ordered (przykładowo)} \\
\hline
Początek & \texttt{flag}: 1 B & \texttt{value}: 8 B \\
\hline
Dalej & padding: 7 B & \texttt{count}: 4 B \\
\hline
Dalej & \texttt{value}: 8 B & \texttt{flag}: 1 B \\
\hline
Koniec & \texttt{count}: 4 B + padding: 4 B & padding: 3 B \\
\hline
Razem & 24 B & 16 B \\
\hline
\end{tabular}
\end{center}

\subsection{\texttt{offsetof}, standard-layout i ABI}

Dla typów \emph{standard-layout} makro \texttt{offsetof} pozwala
poznać przesunięcie pola od początku obiektu. Używanie go z innymi
typami klas od C++17 jest warunkowo wspierane, dlatego do analiz
przenośnych przyjmujemy tutaj standard-layout.

\begin{lstlisting}
#include <cstddef>
#include <cstdint>
#include <type_traits>

struct Header {
    std::uint32_t id;
    std::uint16_t version;
};
static_assert(std::is_standard_layout_v<Header>);
static_assert(offsetof(Header, version) >=
              sizeof(std::uint32_t));
\end{lstlisting}

W analizie ABI trzeba pamiętać o polach wirtualnych, dziedziczeniu,
przestrzeniach adresowych i konwencjach wywołania. Zmiana kolejności
pól publicznej struktury używanej w bibliotekach binarnych może
zmienić jej ABI nawet wtedy, gdy program źródłowy nadal się kompiluje.

\subsection{\texttt{alignof} i \texttt{alignas}}

\texttt{alignof(T)} zwraca wyrównanie typu \texttt{T};
\texttt{alignas(N)} pozwala narzucić silniejsze wyrównanie.
Należy zachować ograniczenia języka dotyczące dopuszczalnych wartości
\texttt{N}. Wyrównanie może zwiększyć rozmiar obiektu, co jest
istotne przy tablicach.

\begin{lstlisting}
#include <atomic>
#include <cstdint>

struct alignas(64) AlignedCounter {
    std::atomic<std::uint64_t> value{0};
};
static_assert(alignof(AlignedCounter) >= 64);
\end{lstlisting}

Wyrównanie do 64 bajtów ma sens tylko wtedy, gdy wynika z wymagań
sprzętu lub pomiarów. Nie jest to uniwersalna wielkość linii cache.
Ponadto \texttt{alignas} nie czyni dostępu do obiektu atomowym,
nie zastępuje \texttt{std::atomic} i samo nie zapewnia synchronizacji.

\subsection{Czy warto „pakować” struktury?}

Atrybuty kompilatora w rodzaju \texttt{\_\_attribute\_\_((packed))}
lub \texttt{\#pragma pack} nie stanowią przenośnych mechanizmów
standardowego C++20. Zmniejszenie rozmiaru może wymusić niewyrównany
dostęp, obniżyć wydajność, a na niektórych architekturach być
problemem poprawności. W kodzie protokołów lepiej często dekodować
pola jawnie z bufora bajtów i nie wiązać formatu z layoutem klasy.

\section{Wyrównane alokacje i czas życia obiektów}
\label{sec:cpp10-lifetime}

\subsection{Różnica między alokacją a konstruowaniem}

Klasyczne \texttt{new T(args)} wykonuje logicznie dwa zadania:
pozyskuje odpowiednio wyrównaną pamięć i inicjalizuje w niej obiekt.
Analogicznie \texttt{delete} niszczy obiekt i zwalnia pamięć.
Dla narzędzi niskopoziomowych warto rozdzielić te operacje.

\begin{lstlisting}
#include <new>

struct Widget { explicit Widget(int) {} };

void* mem = ::operator new(sizeof(Widget));
Widget* w = nullptr;
try {
    w = ::new (mem) Widget(5); // placement new
} catch (...) {
    ::operator delete(mem);
    throw;
}
w->~Widget();                 // koniec zycia obiektu
::operator delete(mem);      // zwolnienie samej pamieci
\end{lstlisting}

Taka konstrukcja jest poprawna dla \texttt{Widget} o wyrównaniu
obsługiwanym przez zwykłe \texttt{operator new}. Dla typów
\emph{over-aligned} należy korzystać z odpowiedniej alokacji
obsługującej \texttt{std::align\_val\_t} lub ze standardowego
alokatora, który uwzględnia wymagania wyrównania.

\paragraph{Dlaczego to nie jest zalecany idiom aplikacyjny?}
Obsługa wyjątków, wielokrotnych destruktorów, rozmiaru tablicy
i zgodności alokacji z dealokacją łatwo prowadzi do błędów.
W większości kodu wystarczają \texttt{std::vector},
\texttt{std::make\_unique} i obiekty automatyczne. Rozdzielanie
alokacji od konstruowania ma sens głównie przy implementowaniu
kontenerów, pul pamięci i innych abstrakcji zasobów.

\subsection{\texttt{std::construct\_at} i \texttt{std::destroy\_at} (C++20)}

Od C++20 można użyć funkcji z nagłówka \texttt{<memory>} do
jawnego rozpoczęcia i zakończenia czasu życia obiektu w odpowiednio
przygotowanej pamięci.

\begin{lstlisting}
#include <cstddef>
#include <memory>

struct Tracked {
    int value;
    explicit Tracked(int n) : value(n) {}
};

alignas(Tracked) std::byte buffer[sizeof(Tracked)];
Tracked* obj = std::construct_at(
    reinterpret_cast<Tracked*>(buffer), 42);
// ... korzystamy z *obj ...
std::destroy_at(obj);
\end{lstlisting}

\textbf{Ważne:} \texttt{std::destroy\_at} kończy czas życia obiektu, ale
nie zwalnia jego pamięci. Nie wolno wykonywać dwóch destrukcji tego
samego żywego obiektu ani odwoływać się do obiektu po destrukcji.
Ponadto rozmiar i wyrównanie bufora muszą odpowiadać konstruowanemu
typowi; sam \texttt{reinterpret\_cast} nie naprawia niewłaściwego
wyrównania.

\subsection{Wyrównany \texttt{operator new}}

Dla typów o szczególnie silnym wyrównaniu istnieją przeciążenia
alokacji i dealokacji wykorzystujące \texttt{std::align\_val\_t}
(C++17). Zwykłe \texttt{new Overaligned{}} wybiera odpowiedni
mechanizm automatycznie.

\begin{lstlisting}
#include <new>

struct alignas(128) Block { char data[128]; };
Block* b = new Block{};
delete b;
\end{lstlisting}

Jeśli ręcznie używamy \texttt{::operator new} z argumentem
\texttt{std::align\_val\_t\{128\}}, musimy dobrać zgodną operację
\texttt{::operator delete} z parametrem wyrównania. Mieszanie tych
rodzin funkcji jest błędem.

\subsection{Ponowne użycie tej samej pamięci i \texttt{std::launder}}

W tej samej pamięci można zakończyć życie jednego obiektu, a potem
rozpocząć życie kolejnego, ale dostęp przez wcześniej utworzone
wskaźniki podlega szczegółowym regułom języka. \texttt{std::launder}
(C++17) ma zastosowanie w określonych sytuacjach, gdy nowy obiekt
nie jest przezroczystym zamiennikiem poprzedniego; nie jest
uniwersalnym remedium na aliasing, nieprawidłowy czas życia
czy źle wyrównaną pamięć.

\paragraph{Reguła praktyczna.}
W kodzie aplikacyjnym nie sięgaj po \texttt{std::launder} bez
konkretnego problemu z czasem życia i bez zrozumienia odpowiednich
reguł. Preferuj własność zasobu wyrażoną typem RAII.

\section{Alokatory: pamięć, obiekty i bezpieczeństwo wyjątkowe}
\label{sec:cpp10-allocators}

\subsection{Standardowy model alokatora}

Kontenery standardowe używają \emph{alokatorów} do uzyskiwania
pamięci na swoje elementy. Alokator nie jest tym samym co funkcja
\texttt{new} i nie musi budować obiektu w chwili rezerwacji pamięci.
\texttt{std::allocator\_traits} zapewnia jednolity interfejs do
alokowania, konstruowania i niszczenia obiektów.

\begin{lstlisting}
#include <memory>

struct Item { explicit Item(int v) : value(v) {} int value; };
std::allocator<Item> alloc;
Item* p = alloc.allocate(2); // pamiec na dwa obiekty

std::allocator_traits<decltype(alloc)>::construct(alloc, p, 5);
std::allocator_traits<decltype(alloc)>::construct(alloc, p + 1, 7);

// teraz dwa obiekty sa zywe
std::allocator_traits<decltype(alloc)>::destroy(alloc, p + 1);
std::allocator_traits<decltype(alloc)>::destroy(alloc, p);
alloc.deallocate(p, 2);
\end{lstlisting}

Przykład dla przejrzystości zakłada, że oba konstruktory nie rzucają
wyjątków. \textbf{Pełna implementacja kontenera} musi zliczać
pomyślnie skonstruowane obiekty i w razie wyjątku zniszczyć tylko je,
a potem zwolnić pamięć. Właśnie z tego powodu lepiej opierać się na
gotowych kontenerach niż ręcznie pisać alokację elementów.

\subsection{Rezerwacja pojemności i realokacja \texttt{vector}}

\texttt{std::vector} przechowuje elementy w pamięci ciągłej, dzięki
czemu zapewnia dobrą lokalność dostępu. \texttt{reserve} może
ograniczyć liczbę realokacji, jeśli znamy przewidywaną wielkość
kontenera.

\begin{lstlisting}
#include <vector>

std::vector<int> values;
values.reserve(10'000);
for (int i = 0; i < 10'000; ++i)
    values.push_back(i);
\end{lstlisting}

\texttt{reserve(n)} nie zmienia \texttt{size()}, lecz zapewnia
\texttt{capacity() >= n}, chyba że zgłoszony zostanie wyjątek.
Realokacja unieważnia wskaźniki, referencje i iteratory do elementów.
Nie jest dobrym pomysłem wywoływanie \texttt{reserve(size()+1)}
przy każdym dodaniu elementu: można w ten sposób utracić
korzystną politykę geometrycznego wzrostu pojemności.

\subsection{Dlaczego własny alokator to nie pierwszy krok?}

Własny alokator w szablonie kontenera wiąże się z wymaganiami
dotyczącymi propagowania przy kopiowaniu, przenoszeniu i zamianie
kontenerów oraz z kwestią porównywania alokatorów.
Dla większości zastosowań wygodniejsza jest warstwa
\texttt{std::pmr}, którą poznamy teraz.

\section{\texttt{std::pmr}: polimorficzne zasoby pamięci}

\begin{center}\small
\begin{tabularx}{\textwidth}{@{}p{0.37\textwidth}X@{}}
\toprule
Mechanizm & Zastosowanie i ograniczenie \\
\midrule
\code{pmr::memory_resource} & Interfejs strategii przydziału i zwalniania pamięci używany przez alokatory PMR. \\
\code{pmr::vector<T>} & Kontener, którego alokator korzysta z wybranego zasobu pamięci. \\
\code{monotonic_buffer_resource} & Przydziela kolejne bloki szybko; pojedyncze zwolnienia nie odzyskują miejsca, zasób zwalnia je zbiorczo. \\
\code{unsynchronized_pool_resource} & Pula małych bloków bez synchronizacji między wątkami. \\
\code{synchronized_pool_resource} & Pula z synchronizacją, odpowiednia do współdzielenia między wątkami. \\
\bottomrule
\end{tabularx}
\end{center}
\label{sec:cpp10-pmr}

\subsection{Alokator jako decyzja w czasie działania}

\texttt{std::pmr} (C++17, \texttt{<memory\_resource>}) oddziela
\emph{rodzaj kontenera} od \emph{sposobu przydzielania pamięci}.
\texttt{std::pmr::vector<T>} jest aliasem \texttt{std::vector<T, ...>}
z alokatorem polimorficznym. Zasób \texttt{memory\_resource}
można wybierać w czasie działania programu.

\begin{lstlisting}
#include <memory_resource>
#include <vector>

std::pmr::monotonic_buffer_resource arena;
std::pmr::vector<int> numbers{&arena};
for (int i = 0; i < 100; ++i)
    numbers.push_back(i);
\end{lstlisting}

Kontener przechowuje \emph{wskaźnik do zasobu}, ale nie przejmuje
własności zasobu pamięci. Zasób musi
pozostać żywy wystarczająco długo, aby obsłużyć wszystkie powiązane
alokacje i dealokacje. \textbf{Nie wolno niszczyć lub resetować
areny, gdy żywe obiekty wciąż korzystają z jej pamięci}.

\subsection{Arena monotoniczna: wiele alokacji, jedno zwolnienie}

\texttt{std::pmr::monotonic\_buffer\_resource} wydaje kolejne
fragmenty pamięci. Jej \texttt{deallocate} nie zwalnia pojedynczych
bloków; odzyskanie następuje przy \texttt{release()} lub destrukcji
zasobu. To pasuje do wielu obiektów budowanych razem i usuwanych
w jednej fazie, np. drzewa parsowania pojedynczego dokumentu.

\begin{lstlisting}
#include <cstddef>
#include <memory_resource>
#include <vector>

alignas(std::max_align_t) std::byte storage[4096];
std::pmr::monotonic_buffer_resource arena(
    storage, sizeof(storage),
    std::pmr::null_memory_resource());

{
    std::pmr::vector<int> data{&arena};
    for (int i = 0; i < 32; ++i) data.push_back(i);
} // kontener znika; jego deallocate nie uwalnia blokow
// arena jest teraz bezpieczna do zniszczenia lub release()
\end{lstlisting}

\texttt{null\_memory\_resource()} wymusza niepowodzenie alokacji,
gdy lokalny bufor się wyczerpie (rzucane jest \texttt{std::bad\_alloc}).
Jeżeli zamiast niego użyjemy zwykłego zasobu nadrzędnego
(\emph{upstream}), arena może dobierać kolejne bufory.

\paragraph{Ważna konsekwencja.}
Gdy \texttt{vector} zwiększa pojemność, porzuca stary blok.
W przypadku areny monotonicznej ten blok zwykle pozostaje zajęty
aż do końca całej fazy. Arena nie oznacza zatem automatycznie
najmniejszego zużycia pamięci; warto stosować \texttt{reserve}.

\subsection{Zasoby pulowe: \texttt{unsynchronized\_pool\_resource}}

Gdy obiekty są często alokowane i usuwane w różnych momentach,
przydatne bywają pule przechowujące bloki o podobnych rozmiarach.

\begin{lstlisting}
#include <memory_resource>
#include <vector>

std::pmr::unsynchronized_pool_resource pool;
std::pmr::vector<int> items{&pool};
items.resize(500);
\end{lstlisting}

\texttt{unsynchronized\_pool\_resource} nie zapewnia bezpieczeństwa
jednoczesnego użycia z wielu wątków. Do współdzielonego użycia
istnieje \texttt{synchronized\_pool\_resource}. Wybór zależy od
profilu alokacji i pomiarów: pule nie są uniwersalnie szybsze.

\subsection{Własny zasób: liczymy rzeczywiste alokacje}

Wariant diagnostyczny może zliczać żądania pamięci i delegować je
do \texttt{new\_delete\_resource()}. Trzy wirtualne funkcje
\texttt{do\_allocate}, \texttt{do\_deallocate}, \texttt{do\_is\_equal}
tworzą podstawę własnej implementacji.

\begin{lstlisting}
#include <cstddef>
#include <memory_resource>

struct CountingResource : std::pmr::memory_resource {
    std::pmr::memory_resource* upstream =
        std::pmr::new_delete_resource();
    std::size_t allocations = 0;
    std::size_t bytes = 0;
private:
    void* do_allocate(std::size_t n,
                      std::size_t alignment) override {
        void* p = upstream->allocate(n, alignment);
        ++allocations;
        bytes += n;
        return p;
    }
    void do_deallocate(void* p, std::size_t n,
                       std::size_t alignment) override {
        upstream->deallocate(p, n, alignment);
    }
    bool do_is_equal(
        const std::pmr::memory_resource& other
    ) const noexcept override {
        return this == &other;
    }
};
\end{lstlisting}

\begin{lstlisting}
CountingResource stats;
{
    std::pmr::vector<int> data{&stats};
    for (int i = 0; i < 100; ++i) data.push_back(i);
}
// stats.allocations i stats.bytes opisuja zadania
// przeslane do tego konkretnego zasobu.
\end{lstlisting}

\textbf{Uwaga:} nie są to wszystkie alokacje programu, a
\texttt{bytes} oznacza sumę żądanych rozmiarów, nie największe
jednoczesne zużycie pamięci. Warto liczyć także dealokacje,
\emph{peak usage} i rozkład rozmiarów bloków.

\subsection{Zagnieżdżone kontenery i propagacja zasobu}

\begin{lstlisting}
#include <memory_resource>
#include <string>
#include <vector>

CountingResource stats;
std::pmr::vector<std::pmr::string> names{&stats};
names.emplace_back("alpha");
names.emplace_back("a rather long text that needs allocation");
\end{lstlisting}

\texttt{pmr::string} również korzysta z polimorficznego alokatora,
ale dokładne zachowanie małych napisów zależy od implementacji
(\emph{small string optimization}). Nie należy zakładać, że każdy
\texttt{emplace\_back} przydziela pamięć ani że krótki napis wymaga
alokacji dynamicznej.

\section{Lokalność danych i hierarchia cache}
\label{sec:cpp10-cache}

\subsection{Dlaczego asymptotyka nie wystarcza?}

Dwa algorytmy o złożoności \(O(n)\) mogą mieć bardzo różną wydajność,
jeżeli pierwszy czyta sąsiadujące elementy tablicy, a drugi podąża
za wskaźnikami do rozproszonych obiektów. Procesor pobiera dane
porcjami (\emph{cache lines}); dostęp sekwencyjny sprzyja
prefetchingowi, zaś losowy dostęp bywa ograniczony opóźnieniem pamięci.

Nie ma jednego rozmiaru linii cache dla wszystkich procesorów ani
jednej wartości kosztu trafienia i chybienia. Zależą one od sprzętu,
aktualnego obciążenia i rodzaju dostępu.

\begin{lstlisting}
#include <numeric>
#include <vector>

std::vector<int> v(100000, 1);
long long sum = std::accumulate(v.begin(), v.end(), 0LL);
// Dostep sekwencyjny do pamieci ciaglej.
\end{lstlisting}

Wskaźniki do osobno alokowanych małych węzłów mogą pogorszyć
lokalność i powiększyć narzut pamięciowy. Nie oznacza to,
że \texttt{std::vector} jest zawsze lepszy od \texttt{std::list}:
wybór zależy od wymaganych operacji, unieważniania iteratorów
i rzeczywistego profilu obciążenia.

\subsection{AoS kontra SoA}

W układzie \emph{Array of Structures} (AoS) przechowujemy obok siebie
wszystkie składowe kolejnych rekordów. W \emph{Structure of Arrays}
(SoA) przechowujemy osobne tablice poszczególnych składowych.

\begin{lstlisting}
#include <vector>

struct PointAoS { float x, y, z, intensity; };
std::vector<PointAoS> aos;

struct PointSoA {
    std::vector<float> x, y, z, intensity;
};
PointSoA soa;
\end{lstlisting}

Gdy obliczamy wyłącznie sumę współrzędnych \texttt{x}, układ SoA
umożliwia odczytywanie samych potrzebnych danych, a nie całych
rekordów. Ułatwia to też wektoryzację. Jeśli natomiast w każdej
operacji potrzebujemy wszystkich pól jednego rekordu, AoS może być
korzystniejszy. Wydajność zależy od rozmiaru danych, kodu i sprzętu.

\begin{center}
\small
\begin{tabular}{|p{6.0cm}|p{6.0cm}|}
\hline
\textbf{AoS: rekordy obok siebie} & \textbf{SoA: kolumny obok siebie} \\
\hline
\texttt{[x1 y1 z1 i1] [x2 y2 z2 i2]} &
\texttt{[x1 x2 ...] [y1 y2 ...]} \\
\hline
Wygodny dostęp do całego punktu & Wygodny dostęp do jednej cechy wszystkich punktów \\
\hline
Często mniej kontenerów i prostszy model & Często lepsze wykorzystanie cache w obliczeniach kolumnowych \\
\hline
\end{tabular}
\end{center}

\subsection{Stride, skoki i przewidywanie dostępu}

Przechodzenie po pamięci ze skokiem \(k\) oznacza, że kolejne
operacje mogą dotykać odległych linii cache. Dla macierzy w układzie
wierszowym zwykle korzystniej jest przechodzić po kolumnach
\emph{wewnątrz pojedynczego wiersza}, a następnie po wierszach.
Zanim zmienimy kolejność pętli, należy uwzględnić zależności
obliczeniowe, autowektoryzację i to, czy kompilator już dokonał
danej transformacji.

\begin{lstlisting}
#include <vector>

constexpr std::size_t rows = 100;
constexpr std::size_t cols = 200;
std::vector<double> matrix(rows * cols, 1.0);
double sum = 0.0;
for (std::size_t i = 0; i < rows; ++i)
    for (std::size_t j = 0; j < cols; ++j)
        sum += matrix[i * cols + j];
\end{lstlisting}

\section{False sharing i współbieżne liczniki}
\label{sec:cpp10-false-sharing}

\subsection{Na czym polega false sharing?}

Jeżeli dwa wątki wielokrotnie zapisują \emph{różne} zmienne, ale
zmienne znajdują się w tej samej linii cache, protokół spójności
pamięci podręcznej może wymuszać kosztowną wymianę własności tej
linii między rdzeniami. Nazywamy to \emph{false sharing}.
Nie oznacza to \emph{data race}, jeżeli poszczególne zmienne są
atomowe lub inaczej poprawnie synchronizowane.

\begin{lstlisting}
#include <atomic>
#include <cstdint>

struct NaiveCounters {
    std::atomic<std::uint64_t> first{0};
    std::atomic<std::uint64_t> second{0};
};
\end{lstlisting}

Dwa liczniki mogą być blisko siebie, jednak dokładny layout obiektu
i mapowanie do linii cache nie są gwarantowane. Na niektórych
platformach ich równoczesna modyfikacja będzie kosztowna.

\subsection{Rozdzielenie liczników w pamięci}

\begin{lstlisting}
#include <atomic>
#include <cstdint>

struct alignas(64) IsolatedAtomic {
    std::atomic<std::uint64_t> value{0};
};
struct SeparatedCounters {
    IsolatedAtomic first, second;
};
\end{lstlisting}

Tu użyliśmy 64 bajtów wyłącznie jako \emph{przykładowej} wartości.
Od C++17 w nagłówku \texttt{<new>} istnieją dwie stałe
opisujące zalecane odstępy dla danych często modyfikowanych przez
różne wątki (\emph{destructive interference}) oraz dla danych,
które powinny pozostawać blisko siebie (\emph{constructive
interference}). Pełną nazwę pierwszej z nich pokazuje kod poniżej.
Wartości stałych zależą od implementacji; nie stanowią pomiaru
rozmiaru linii cache w uruchomionym procesorze.

\begin{lstlisting}
#include <atomic>
#include <new>

struct alignas(
    std::hardware_destructive_interference_size
) Counter {
    std::atomic<unsigned> value{0};
};
\end{lstlisting}

Rozdzielenie liczników zwiększa rozmiar struktury. Przy milionie
liczników koszt pamięciowy może przewyższyć korzyść. Technika ma
największy sens, gdy oddzielne wątki często zapisują niewielką
liczbę niezależnych pól.

\subsection{False sharing a semantyka atomowa}

\texttt{std::memory\_order\_relaxed} wystarczy do niektórych
niezależnych liczników, gdy zależy nam wyłącznie na atomowym
zliczaniu; nie gwarantuje jednak synchronizacji innych danych.
Usunięcie false sharing nie zastępuje poprawnego modelu pamięci
współbieżnej. Szczegółowe zasady operacji atomowych były omawiane
w rozdziale~1 o współbieżności.

\section{Jak rzetelnie mierzyć wydajność?}
\label{sec:cpp10-benchmark}

\subsection{Zasady pomiarów}

Najważniejsze pytanie brzmi nie „czy ten kod wygląda szybciej?”,
lecz „czy w rzeczywistym scenariuszu uzyskujemy powtarzalną poprawę?”.
Doświadczenie powinno uwzględniać:
\begin{enumerate}
\item kompilację zoptymalizowaną (np. \texttt{-O2} lub \texttt{-O3});
\item \emph{warm-up} i wielokrotne uruchomienia;
\item jednakowe dane oraz taki sam rezultat dla obu wariantów;
\item ochronę przed usunięciem testowanej pracy przez optymalizator;
\item pomiar rozkładu wyników, a nie pojedynczego odczytu;
\item ograniczenie wpływu alokacji, I/O, schedulerów i skalowania CPU;
\item kontrolę zużycia pamięci, a nie tylko czasu;
\item narzędzia profilujące (np. Linux \texttt{perf}, odpowiedniki
      sprzętowe lub profilery danego środowiska).
\end{enumerate}

\subsection{Minimalna funkcja pomiarowa}

\begin{lstlisting}
#include <chrono>
#include <utility>

template<class F>
auto measure(F&& f) {
    auto start = std::chrono::steady_clock::now();
    auto result = std::forward<F>(f)();
    auto stop = std::chrono::steady_clock::now();
    return std::pair{result, stop - start};
}
\end{lstlisting}

Taka funkcja przydaje się do nauki, ale \emph{nie jest kompletna}
biblioteka benchmarkowa. Dla zbyt krótkich fragmentów narzut
pomiaru może przeważać nad czasem pracy. Dla liczb zmiennoprzecinkowych
porównywanie różnych kolejności sumowania wymaga ostrożności,
ponieważ zmiana kolejności działań może zmienić zaokrąglenia.

\subsection{Sanitizatory i wydajność}

Sanitizatory, np. AddressSanitizer i UndefinedBehaviorSanitizer,
są narzędziami do badania poprawności pamięci i innych błędów.
Wprowadzają narzut, dlatego \emph{pomiarów porównawczych wydajności}
nie należy bezrefleksyjnie wykonywać na buildach sanitizowanych.
Bezpieczny cykl pracy wygląda następująco: test poprawności,
profilowanie zwykłego buildu, zmiana jednej rzeczy, ponowny test
poprawności i ponowny pomiar.

\section{Przykładowy projekt: szybkie przetwarzanie wielu punktów}
\label{sec:cpp10-project}

Przyjmijmy, że program przechowuje milion punktów \((x,y,z,I)\)
i dla każdego oblicza tylko prosty filtr zależny od \(x\) i \(I\).
Możemy rozważyć dwa modele:
\begin{enumerate}
\item \textbf{AoS:} \texttt{vector<Point>} --- prostszy interfejs,
      wygodny dostęp do pełnego punktu.
\item \textbf{SoA:} osobne \texttt{vector<float>} na \(x,y,z,I\) ---
      możliwy mniejszy transfer danych potrzebnych do filtru.
\end{enumerate}

Następnie oceniamy, czy warto użyć \texttt{reserve} przed zapełnieniem
kontenerów; jeśli rekordy są tworzone na potrzeby jednego zapytania,
możemy sprawdzić arenę monotoniczną. Trzeba jednak rozróżnić
\emph{koszt wypełnienia struktur} od \emph{kosztu przetwarzania}: arena
zmienia zachowanie alokacji, nie gwarantuje poprawy pracy procesora
nad danymi.

Dla kodu numerycznego duży wpływ może mieć też wybór kolejności
wymiarów macierzy, reprezentacji sąsiedztwa w siatce czy sposobu
obliczania operatora na konkretnym punkcie zamiast konstruowania
dużych obiektów pośrednich. Właściwa reprezentacja danych potrafi
dać większy zysk niż drobne optymalizacje pojedynczych instrukcji.

\section{Jak dobrać technikę?}

\begin{center}
\small
\begin{tabular}{|p{5.7cm}|p{6.7cm}|}
\hline
\textbf{Problem} & \textbf{Pierwszy sensowny krok} \\
\hline
Za dużo dynamicznych alokacji & \texttt{reserve}, ograniczenie obiektów pośrednich, pomiar alokacji \\
\hline
Wiele obiektów znika razem & \texttt{pmr::monotonic\_buffer\_resource} po kontroli czasu życia \\
\hline
Dużo małych alokacji o różnej żywotności & \texttt{pmr::...pool\_resource} i pomiar \\
\hline
Praca nad jedną kolumną wielu rekordów & Porównanie AoS z SoA \\
\hline
Losowe skoki po pamięci & Sprawdzenie lokalności, zmian formatu danych i blokowania pętli \\
\hline
Niezależne atomiki spowalniają się nawzajem & Analiza false sharing i separacja intensywnie zapisywanych pól \\
\hline
Trzeba uzyskać silne wyrównanie & \texttt{alignas}, odpowiednia alokacja i kontrola \texttt{alignof} \\
\hline
Trzeba napisać kontener lub pulę & \texttt{allocator\_traits}, zarządzanie czasem życia, RAII \\
\hline
Potrzeba przenośnej serializacji & Jawny format bajtowy, bez kopiowania całych struktur \\
\hline
\end{tabular}
\end{center}

\section{Typowe błędy i pułapki}

\begin{enumerate}
\item \textbf{Wniosek „\texttt{sizeof(T)} to suma pól”.}
      Padding może zwiększyć rozmiar, a reguły klas mogą być złożone.
\item \textbf{Uznanie layoutu z jednej kompilacji za standard.}
      ABI, wyrównanie i rozmiary typów różnią się między platformami.
\item \textbf{Nieprawidłowe \texttt{reinterpret\_cast}.}
      Sama zmiana typu wskaźnika nie legalizuje dostępu do danych.
\item \textbf{Mylenie pamięci z czasem życia obiektu.}
      \texttt{allocate} nie zastępuje konstruowania obiektów klasowych.
\item \textbf{Zwalnianie pamięci bez destruktora.}
      Dla nietrywialnych typów może wyciekać zasób zarządzany przez obiekt.
\item \textbf{Niezgodna para alokacja--dealokacja.}
      Szczególnie ważne przy over-alignment i ręcznej alokacji.
\item \textbf{Pozostawienie kontenera po zniszczeniu areny \texttt{pmr}.}
      Jego pamięć staje się niepoprawna; wynik to niezdefiniowane zachowanie.
\item \textbf{Wywołanie \texttt{release()} przy żywych obiektach.}
      Arena może unieważnić ich pamięć bez wywoływania destruktorów.
\item \textbf{Założenie, że \texttt{monotonic\_buffer\_resource} oszczędza
      pamięć.} Zwolnienie pojedynczych bloków zwykle nie następuje.
\item \textbf{Współdzielenie \texttt{unsynchronized\_pool\_resource}
      między wątkami bez ochrony.} Zasób nie jest thread-safe.
\item \textbf{Optymalizacja poprzez mniejszą strukturę bez pomiaru.}
      Mniejszy \texttt{sizeof} nie gwarantuje krótszego czasu wykonania.
\item \textbf{Rzekoma optymalizacja przez \texttt{volatile}.}
      \texttt{volatile} nie zapewnia synchronizacji między wątkami.
\item \textbf{Pomiar jednego przebiegu.} Scheduler, częstotliwości
      zegara i cache mogą silnie zniekształcić wynik.
\item \textbf{Użycie sanitizatora jako miary szybkości.}
      Narzut instrumentacji zmienia charakterystykę programu.
\end{enumerate}

\paragraph{Pytania kontrolne.}
\begin{enumerate}
\item Czym różnią się \texttt{sizeof(T)} i \texttt{alignof(T)}?
\item Dlaczego zmiana kolejności pól może zmienić rozmiar struktury?
\item Czy placement \texttt{new} przydziela pamięć?
\item Jaka jest różnica między \texttt{destroy\_at} a \texttt{deallocate}?
\item Kiedy arena monotoniczna może zużywać więcej pamięci niż oczekujemy?
\item Dlaczego \texttt{pmr::vector<T>} nie jest właścicielem zasobu?
\item Kiedy SoA może przyspieszyć obliczenia na rekordach?
\item Czy false sharing i data race są tym samym problemem?
\item Dlaczego \texttt{std::memory\_order\_relaxed} nie rozwiązuje
      wszystkich problemów synchronizacji?
\item Jak odróżnić realną poprawę wydajności od szumu pomiarowego?
\end{enumerate}

\paragraph{Dalsza lektura.}
Dokumentacja \emph{cppreference} obejmuje hasła \emph{Object},
\emph{Object lifetime}, \texttt{alignof}, \texttt{alignas},
\texttt{std::construct\_at} oraz \texttt{allocator\_traits}.
Osobno warto przeczytać opis zasobów \texttt{std::pmr}
i własności \emph{hardware interference size}.
W praktyce polecam analizę śladów alokacji, profilery
wydajności i \emph{Compiler Explorer} do porównywania
wygenerowanego kodu maszynowego.

% ROZDZIAL 7 -- Lambdy, obiekty funkcyjne i wywolywalne obiekty (C++20)
% Fragment do wlaczenia po rozdziale 6.
% Wymagana preambula: listings, array, amsmath (opcjonalny), hyperref.
% Uzycie: \input{rozdzial_7_lambdy_obiekty_funkcyjne_cpp.tex}

\chapter[Lambdy i obiekty funkcyjne]{Wyrażenia lambda, obiekty funkcyjne i abstrakcje wywoływalne}
\label{chap:lambdas-cpp}

\section{Po co istnieją obiekty funkcyjne?}

W C++ funkcję można reprezentować nie tylko przez jej nazwę lub wskaźnik.
Można również utworzyć \emph{obiekt wywoływalny} (\textit{callable object}):
obiekt, który daje się zastosować w wyrażeniu przypominającym wywołanie
funkcji. Pozwala to przekazywać zachowanie jako parametr algorytmu,
zapamiętywać stan pomiędzy wywołaniami oraz tworzyć elastyczne API
(np. sortowanie z własnym porządkiem, filtrowanie, obsługa zdarzeń).

\begin{lstlisting}
#include <algorithm>
#include <vector>

std::vector<int> numbers{5, 2, 8, 1};
std::sort(numbers.begin(), numbers.end(),
          [](int a, int b) { return a > b; });
// numbers == {8, 5, 2, 1}
\end{lstlisting}

Trzeci argument \texttt{std::sort} to reguła porównywania. Zamiast pisać
osobną funkcję o globalnej nazwie, definiujemy ją dokładnie w miejscu
użycia. Wyrażenie \texttt{[](int a, int b) \{ return a > b; \}}
jest \emph{lambdą}: tworzy obiekt unikatowego typu domknięcia
(\textit{closure type}).

\paragraph{Zakres i wymagania.}
Przyjmujemy \textbf{C++20}. Lambdy wprowadzono w C++11, przechwytywanie
inicjalizujące i generyczne lambdy w C++14, przechwytywanie kopii
\texttt{*this} w C++17, a jawne listy parametrów szablonu lambdy
w C++20. Wykorzystamy semantykę przenoszenia z rozdziału~\ref{chap:move-semantics-cpp},
szablony wariadyczne z rozdziału~\ref{chap:variadic-templates-cpp}
i koncepcje z rozdziału~\ref{chap:concepts-cpp}.

\section{Trzy podstawowe postacie kodu wywoływalnego}
\label{sec:callables-forms}

\subsection{Funkcja i wskaźnik do funkcji}

Zwykła funkcja nie przechowuje kontekstu wywołania. Jej adres można
przypisać wskaźnikowi o odpowiednim typie:

\begin{lstlisting}
int twice(int x) { return 2 * x; }

int (*fp)(int) = &twice;
int a = fp(7);       // 14
int b = (*fp)(7);    // also 14
\end{lstlisting}

W deklaracji \texttt{int (*fp)(int)} nawiasy są istotne. Sam wskaźnik
przechowuje adres funkcji; nie ma miejsca na dodatkowy stan, taki jak
współczynnik mnożenia ustalany w czasie działania programu.

\subsection{Funktor: klasa z \texttt{operator()}}

\begin{lstlisting}
struct Multiply {
    int factor;

    int operator()(int x) const {
        return factor * x;
    }
};

Multiply times_three{3};
int result = times_three(7); // 21
\end{lstlisting}

\texttt{Multiply} jest klasą, a \texttt{times\_three} jej egzemplarzem.
Dzięki polu \texttt{factor} każdy egzemplarz może reprezentować inne
zachowanie. Operator wywołania może być przeciążany, szablonowy,
opatrzony kwalifikatorami \texttt{const} lub \texttt{\&\&} i mieć
złożoną implementację.

\subsection{Lambda jako wygodny zapis funktora}

\begin{lstlisting}
int factor = 3;
auto times_three = [factor](int x) {
    return factor * x;
};
int result = times_three(7); // 21
\end{lstlisting}

Kompilator tworzy odrębny, nienazwany typ klasy odpowiadający lambdzie.
Uproszczony model jej działania jest następujący:

\begin{lstlisting}
struct ClosureModel {
    int factor;

    int operator()(int x) const {
        return factor * x;
    }
};
\end{lstlisting}

Jest to \emph{model pojęciowy}, nie obiecana przez standard reprezentacja
pamięci. Każde wystąpienie wyrażenia lambda ma własny typ, nawet jeżeli
zapis kodu dwóch lambd wygląda identycznie.

\section{Budowa wyrażenia lambda}
\label{sec:lambda-syntax}

Najważniejsze elementy zapisu przedstawia schemat:

\begin{lstlisting}
[captures](parameters) mutable noexcept -> ResultType {
    // body
}
\end{lstlisting}

W C++20 można również jawnie określić parametry szablonu, koncepcje
oraz klauzulę \texttt{requires}. Typ wyniku zwykle jest wywnioskowany
z instrukcji \texttt{return}; podanie \texttt{-> ResultType} bywa
przydatne, gdy chcemy wymusić konkretny typ lub zwracać referencję.

\begin{lstlisting}
auto square = [](int x) { return x * x; };
auto half = [](double x) -> double { return x / 2; };
auto truth = [] { return true; }; // () may be omitted

int a = square(4);  // 16
double b = half(5); // 2.5
\end{lstlisting}

\paragraph{Najważniejsze zastrzeżenie.}
Lambda nie jest wskaźnikiem do funkcji. Zwykle tworzy obiekt z
\texttt{operator()}; jedynie pewne lambdy bez przechwyceń mają
konwersję na wskaźnik do funkcji o pasującej sygnaturze.

\section{Przechwytywanie zmiennych (\textit{captures})}

\begin{center}\small
\begin{tabularx}{\textwidth}{@{}p{0.30\textwidth}X@{}}
\toprule
Zapis lambdy & Znaczenie \\
\midrule
\code{[x]} & Przechwytuje wartość \code{x} przez kopię w obiekcie domknięcia. \\
\texttt{[\&x]} & Przechwytuje referencję do \code{x}; wywołanie wymaga, by \code{x} nadal żyło. \\
\code{[=]} & Domyślnie przechwytuje użyte zmienne lokalne przez wartość. \\
\texttt{[\&]} & Domyślnie przechwytuje użyte zmienne lokalne przez referencję. \\
\code{[p = std::move(p)]} & Przenosi wartość do składowej domknięcia przez inicjalizację przechwycenia. \\
\code{mutable} & Pozwala zmieniać kopie przechwycone przez wartość wewnątrz lambdy. \\
\bottomrule
\end{tabularx}
\end{center}
\label{sec:lambda-captures}

\subsection{Przechwycenie przez wartość}

\begin{lstlisting}
int threshold = 10;
auto above = [threshold](int x) {
    return x > threshold;
};
auto above_current = [&threshold](int x) {
    return x > threshold;
};
threshold = 100;

bool first = above(20);         // true: zapamietane 10
bool second = above_current(20); // false: aktualne 100
\end{lstlisting}

\texttt{[threshold]} inicjalizuje składową obiektu domknięcia wartością
\texttt{threshold} w chwili tworzenia lambdy. Późniejsza zmiana zmiennej
zewnętrznej nie modyfikuje tej kopii. Jeśli kopiowany obiekt ma kosztowne
kopiowanie, samo utworzenie lub skopiowanie lambdy może także mieć koszt.

\subsection{Przechwycenie przez referencję}

\begin{lstlisting}
int total = 0;
auto add = [&total](int value) {
    total += value;
};
add(4);
add(8); // total == 12
\end{lstlisting}

\texttt{[\&total]} uzyskuje dostęp do tej samej zmiennej, nie do jej
kopii. Lambda \textbf{nie wydłuża czasu życia} obiektu przechwyconego
przez referencję. Dlatego jest wygodna, gdy callback wykonywany jest
synchronicznie, ale potencjalnie niebezpieczna, jeśli jest zapamiętywany
albo wykonywany później.

\subsection{Przechwycenia domyślne i mieszane}

\begin{lstlisting}
int a = 2;
int b = 3;
int count = 0;

auto first  = [=]  { return a + b; }; // by copy
auto second = [&]  { ++count; };     // by reference
auto third  = [=, &count] {
    ++count;
    return a + b;
};
\end{lstlisting}

\texttt{[=]} i \texttt{[\&]} obejmują te użyte w treści zmienne lokalne,
które podlegają niejawnemu przechwyceniu; nie oznaczają automatycznego
skopiowania \emph{wszystkich} zmiennych z otoczenia. W większym kodzie
zwykle lepiej wypisać przechwycenia jawnie: dzięki temu widać zależności
i łatwiej ocenić bezpieczeństwo czasu życia.

\subsection{Przechwytywanie inicjalizujące (C++14)}

W nawiasach kwadratowych można zdefiniować nową nazwę i sposób jej
inicjalizacji. Dzięki temu lambdy mogą przechowywać wartości pochodne,
a także zasoby, których nie wolno kopiować:

\begin{lstlisting}
#include <memory>
#include <utility>

auto p = std::make_unique<int>(42);
auto read = [q = std::move(p)] {
    return *q;
};
// p == nullptr; ownership is now in the closure
int answer = read(); // 42
\end{lstlisting}

Taka lambda jest \emph{move-only}, ponieważ zawiera
\texttt{std::unique\_ptr}. Nie można jej skopiować. To istotne przy
wyborze miejsca przechowywania callbacków: w C++20
\texttt{std::function} wymaga kopiowalnego obiektu docelowego.

\subsection{Niebezpieczne przechwytywanie: wisząca referencja}

\begin{lstlisting}
// WRONG: do not call the returned lambda!
auto broken() {
    int local = 5;
    return [&local] { return local; };
}

// Correct: capture the value.
auto safe() {
    int local = 5;
    return [local] { return local; };
}
\end{lstlisting}

\texttt{broken()} zwraca obiekt z dostępem do zmiennej, która przestaje
istnieć wraz z wyjściem z funkcji. Późniejsze wywołanie prowadziłoby
do niezdefiniowanego zachowania (\textit{undefined behavior}).
Przechwycenie przez kopię rozwiązuje ten konkretny problem.

\section{Przechwytywanie \texttt{this} i \texttt{*this}}
\label{sec:lambda-this}

W metodzie klasy można korzystać z jej pól wewnątrz lambdy. Trzeba
jednak rozróżnić przechwycenie wskaźnika \texttt{this} od skopiowania
bieżącego obiektu.

\begin{lstlisting}
#include <string>

struct Label {
    std::string text;

    auto by_pointer() const {
        return [this] { return text; };
    }

    auto by_copy() const {
        return [*this] { return text; }; // C++17
    }
};
\end{lstlisting}

Lambda zwrócona przez \texttt{by\_pointer()} odwołuje się do oryginalnego
obiektu. Nie może być bezpiecznie wywoływana po jego zniszczeniu.
Lambda z \texttt{[*this]} przechowuje kopię obiektu (jeśli da się go
skopiować); dla składowych zarządzanych przez RAII jest to zwykle
bardziej przewidywalne, ale może być kosztowne.

\textbf{Uwaga:} zapis \texttt{[=]} użyty w metodzie nie kopiuje automatycznie
całego obiektu. Niejawne przechwytywanie \texttt{this} przez \texttt{[=]}
jest w C++20 przestarzałe (\textit{deprecated}); zapisujmy
\texttt{[this]} jawnie albo \texttt{[*this]}, gdy potrzebujemy kopii.

\section{Domyślnie stały operator wywołania i \texttt{mutable}}

Dla typowej lambdy bez słowa \texttt{mutable} operator wywołania ma
kwalifikator \texttt{const}. Oznacza to, że przechwyconych przez wartość
składowych nie wolno w nim zwyczajnie modyfikować:

\begin{lstlisting}
int counter = 0;

// ERROR: captured counter is not writable in this closure.
// auto bad = [counter] { ++counter; };

auto next = [counter]() mutable {
    return ++counter;
};

int a = next(); // 1
int b = next(); // 2
// outer counter is still 0
\end{lstlisting}

\texttt{mutable} modyfikuje zachowanie \emph{kopii wewnątrz lambdy},
nie powoduje przechwycenia oryginalnej zmiennej przez referencję.
Różne kopie obiektu domknięcia mogą przechowywać różny stan:

\begin{lstlisting}
auto a = [n = 0]() mutable { return ++n; };
auto b = a; // independent copy of state

a(); // 1
a(); // 2
b(); // 1
\end{lstlisting}

Przy przechwyceniu przez referencję do modyfikacji zewnętrznego
obiektu \texttt{mutable} nie jest potrzebne: stałość metody
\texttt{operator()} nie zmienia stałości obiektu osiąganego przez
referencję.

\section{Typ lambdy, kopiowanie i konwersja na wskaźnik}

\begin{lstlisting}
#include <type_traits>

inline constexpr auto f = [](int x) { return x + 1; };

auto a = [](int x) { return x; };
auto b = [](int x) { return x; };
static_assert(!std::is_same_v<decltype(a), decltype(b)>);

int (*pf)(int) = +[](int x) { return x + 1; };
\end{lstlisting}

Dwie osobno zapisane lambdy mają różne typy, nawet jeśli robią to samo.
W C++20 typ domknięcia lambdy bez przechwyceń jest domyślnie
konstruowalny; dostępność różnych operacji kopiowania/przypisywania
zależy od tego, czy lambda przechwytuje zmienne. Dla lambd z
przechwyceniami obowiązują też ograniczenia narzucone przez typy pól
(na przykład niekopiowalność \texttt{unique\_ptr}).

\texttt{+[]} jest znanym idiomem wymuszającym konwersję
nieprzechwytującej żadnej zmiennej lambdy na wskaźnik do funkcji.
Nie działa dla lambd przechwytujących:

\begin{lstlisting}
int offset = 10;
auto with_state = [offset](int x) { return x + offset; };
// int (*bad)(int) = with_state;  // ERROR
\end{lstlisting}

\section{Lambdy generyczne i własne parametry szablonu}
\label{sec:lambda-generic}

\subsection{Parametry \texttt{auto} (C++14)}

\begin{lstlisting}
#include <string>
auto print_twice = [](const auto& value) {
    return value + value;
};

auto x = print_twice(7);                 // 14
auto s = print_twice(std::string("ab")); // "abab"
\end{lstlisting}

Parametr \texttt{auto} sprawia, że \texttt{operator()} obiektu domknięcia
jest szablonem funkcji. Przy wywołaniu kompilator dedukuje odpowiednie
typy tak jak przy zwykłej funkcji szablonowej.

\subsection{Jawne parametry szablonu w lambdzie (C++20)}

\begin{lstlisting}
#include <concepts>

auto equal_types = []<class T>(const T& a, const T& b) {
    return a == b;
};

auto sum_integrals = []<std::integral... Ts>(Ts... values) {
    return (0 + ... + values);
};

bool yes = equal_types(3, 3);
int total = sum_integrals(1, 2, 3, 4); // 10
\end{lstlisting}

W przypadku \texttt{equal\_types} oba argumenty muszą dać ten sam
dedukowany typ \texttt{T}. W drugim przykładzie łączymy \emph{concepts}
z pakietem parametrów i wyrażeniem fold znanym z
rozdziałów~\ref{chap:variadic-templates-cpp} i~\ref{chap:concepts-cpp}.

\subsection{Uniwersalny wrapper z perfect forwarding}

\begin{lstlisting}
#include <functional>
#include <utility>

auto invoke_once = []<class F, class... Args>(
    F&& f, Args&&... args) -> decltype(auto) {
    return std::invoke(std::forward<F>(f),
                       std::forward<Args>(args)...);
};
\end{lstlisting}

\texttt{F\&\&} i \texttt{Args\&\&...} są tu referencjami przekazującymi,
a \texttt{decltype(auto)} zachowuje typ wyniku, w tym jego
referencyjność. W odróżnieniu od \texttt{std::move},
\texttt{std::forward} nie wymusza kategorii rvalue dla każdego argumentu.
Dzięki \texttt{std::invoke} wrapper działa nie tylko dla zwykłych
funkcji i lambd, lecz także dla wskaźników do składowych.

\section{Wymuszanie sygnatury: \texttt{noexcept}, \texttt{requires}, zwracanie referencji}

\begin{lstlisting}
#include <concepts>
#include <utility>

auto positive = [](int x) noexcept {
    return x > 0;
};

auto increment = []<std::integral T>(T x) -> T {
    return x + 1;
};

auto first = []<class T>(T& x) -> T& {
    return x;
};

int value = 3;
first(value) = 17; // modifies value
\end{lstlisting}

Ważne są dwie subtelności. Po pierwsze, \texttt{noexcept} jest obietnicą
programisty; jeśli wyjątek wydostanie się z funkcji oznaczonej
\texttt{noexcept}, wywołane zostanie \texttt{std::terminate}.
Po drugie, domyślna dedukcja wyniku lambdy deklarowanej bez
\texttt{decltype(auto)} lub jawnego \texttt{-> T\&} zwykle usuwa
referencję. Właściwe zwracanie referencji wymaga zatem świadomego
określenia typu wyniku i zapewnienia jego poprawnego czasu życia.

\section{\texttt{std::invoke}: jeden mechanizm wywołania}

\begin{center}\small
\begin{tabularx}{\textwidth}{@{}p{0.34\textwidth}X@{}}
\toprule
Narzędzie & Za co odpowiada \\
\midrule
\code{std::invoke(f, args...)} & Jednolicie wywołuje funkcję, lambdę lub wskaźnik do składowej. \\
\code{std::function<R(Args...)>} & Przechowuje kopiowalny obiekt wywoływalny o ustalonej sygnaturze przez wymazywanie typu. \\
\code{std::ref(x)} & Opakowuje referencję, gdy API normalnie kopiuje argument. \\
\code{std::bind_front(f, args...)} & Wiąże pierwsze argumenty funkcji i zwraca nowy obiekt wywoływalny. \\
\code{noexcept} przy lambdzie & Określa, że operator wywołania nie powinien pozwolić wyjątkowi opuścić funkcji. \\
\bottomrule
\end{tabularx}
\end{center}
\label{sec:std-invoke}

\texttt{std::invoke} z nagłówka \texttt{<functional>} wywołuje funkcję,
funktor lub lambdę, ale obsługuje także wskaźnik do metody i pola.

\begin{lstlisting}
#include <functional>

struct Box {
    int value = 10;
    int add(int x) const { return value + x; }
};

Box b;
int a = std::invoke([](int x) { return x * 2; }, 4);
int c = std::invoke(&Box::add, b, 7);      // 17
int& r = std::invoke(&Box::value, b);
r = 25;                                   // b.value == 25
\end{lstlisting}

To podstawowa cegiełka generycznych algorytmów. W nagłówku
\texttt{<type\_traits>} zdefiniowano m.in. \texttt{std::invoke\_result\_t},
\texttt{std::is\_invocable\_v} i \texttt{std::is\_nothrow\_invocable\_v},
które pozwalają sprawdzić w czasie kompilacji poprawność wywołania.

\begin{lstlisting}
#include <type_traits>

static_assert(std::is_invocable_v<decltype(positive), int>);
static_assert(!std::is_invocable_v<decltype(positive),
                                   std::string>);
static_assert(std::is_nothrow_invocable_v<decltype(positive),
                                          int>);
\end{lstlisting}

\section{\texttt{std::function}: jeden typ dla wielu zachowań}
\label{sec:std-function}

Lambdy o różnych typach nie mogą bezpośrednio zastępować się w jednej
zmiennej zadeklarowanej przez \texttt{auto}. Jeżeli chcemy przechowywać
\emph{dowolny} kopiowalny obiekt wywoływalny o zadanej sygnaturze,
możemy użyć \texttt{std::function<R(Args...)>}.

\begin{lstlisting}
#include <functional>
#include <string>

std::function<int(int)> operation =
    [](int x) { return 2 * x; };

int factor = 5;
operation = [factor](int x) { return factor * x; };
int result = operation(3); // 15

operation = nullptr;
if (!operation) {
    // no callable stored
}
\end{lstlisting}

\texttt{std::function} stosuje \emph{type erasure}: ukrywa rzeczywisty
typ obiektu docelowego za wspólnym interfejsem wywołania. Może
przechowywać funkcje, funktory i lambdy z przechwyceniami. Nie jest
odpowiednikiem \texttt{auto}: deklaracja \texttt{auto} zachowuje
konkretny typ i pozwala zwykle kompilatorowi łatwiej zoptymalizować
wywołanie.

\subsection{Pusty \texttt{std::function} i wyjątek}

\begin{lstlisting}
#include <functional>

std::function<void()> callback;
if (callback) {
    callback();
}
// Calling callback() while empty throws std::bad_function_call.
\end{lstlisting}

Sprawdzenie \texttt{if (callback)} jest zatem standardową częścią
obsługi opcjonalnych callbacków. Nie należy przechowywać
\texttt{std::function} bez potrzeby, jeśli konkretny typ wywoływalny
jest znany w czasie kompilacji.

\subsection{Koszty i ograniczenie kopiowalności}

Zwykle \texttt{std::function} wymaga pośredniego wywołania. Przechowywanie
większego domknięcia może też wymagać dynamicznej alokacji (zależy to
od implementacji i rozmiaru obiektu). To koszt elastyczności API, nie
argument przeciw używaniu \texttt{std::function} w każdym przypadku.

W C++20 przechowywany obiekt docelowy musi być kopiowalny, dlatego
następujący kod nie jest poprawny:

\begin{lstlisting}
#include <functional>
#include <memory>

// ERROR in C++20: the closure is not copy-constructible.
// std::function<int()> f =
//     [p = std::make_unique<int>(9)] { return *p; };
\end{lstlisting}

Od \textbf{C++23} biblioteka udostępnia
\texttt{std::move\_only\_function}, która służy do przechowywania
nienadających się do kopiowania obiektów wywoływalnych. W C++20
można natomiast przechowywać taki obiekt bezpośrednio jako \texttt{auto}
lub przekazać go do szablonu, który nie wymaga kopiowania.

\section{Co wybrać: \texttt{auto}, szablon, wskaźnik czy \texttt{std::function}?}
\label{sec:lambda-choice}

\begin{center}
\begin{tabular}{|p{3.0cm}|p{4.4cm}|p{5.0cm}|}
\hline
\textbf{Forma} & \textbf{Zalety} & \textbf{Ograniczenia} \\
\hline
\texttt{auto f = lambda} & Konkretny typ, bez konieczności wymazywania typu, naturalna obsługa stanu & Zmienna ma ustalony typ domknięcia \\
\hline
\texttt{template<class F>} & Działa z wieloma typami, umożliwia inline i forwarding & Definicja szablonu zwykle dostępna w nagłówku; możliwe powiększanie kodu \\
\hline
Wskaźnik do funkcji & Prosty nieposiadający stanu interfejs & Nie przyjmuje lambd z przechwyceniami \\
\hline
\texttt{std::function} & Jedna sygnatura, zmienne zachowanie w czasie wykonania & Type erasure; koszt pośredniego wywołania; wymóg kopiowalności celu w C++20 \\
\hline
\end{tabular}
\end{center}

\paragraph{Praktyczna reguła.}
Jeżeli callback wykorzystujemy natychmiast, wybierajmy zwykle szablon
lub \texttt{auto}. Jeżeli musimy przechowywać i później podmieniać różne
callbacki o jednej sygnaturze, \texttt{std::function} jest często
czytelniejsza niż ręczne modelowanie polimorfizmu.

\section{Lambdy ze standardowymi algorytmami}
\label{sec:lambda-algorithms}

\subsection{\texttt{find\_if}, \texttt{count\_if}, \texttt{transform}}

\begin{lstlisting}
#include <algorithm>
#include <vector>

std::vector<int> xs{1, 2, 3, 4, 5, 6};

int limit = 3;
auto it = std::find_if(xs.begin(), xs.end(),
                       [limit](int x) { return x > limit; });
// *it == 4

auto even_count = std::count_if(xs.begin(), xs.end(),
                                [](int x) { return x % 2 == 0; });
// even_count == 3

std::vector<int> squares(xs.size());
std::transform(xs.begin(), xs.end(), squares.begin(),
               [](int x) { return x * x; });
\end{lstlisting}

Dla \texttt{std::find\_if} trzeba przed dereferencją sprawdzić, czy
wynik jest różny od \texttt{xs.end()}. Dla \texttt{std::transform} zbiór
docelowy musi już zawierać wystarczającą liczbę elementów: samo
\texttt{reserve()} zwiększa \texttt{capacity()}, lecz nie \texttt{size()}.
Alternatywą jest pusty wektor i iterator \texttt{std::back\_inserter}.

\subsection{\texttt{std::ranges} i projekcje (C++20)}

\begin{lstlisting}
#include <algorithm>
#include <ranges>
#include <string>
#include <vector>

struct Person {
    std::string name;
    int age;
};

std::vector<Person> people{{"Ada", 31}, {"Bob", 18}, {"Eve", 25}};

std::ranges::sort(people, {}, &Person::age);
// sorted by age ascending

auto adults = people | std::views::filter(
    [](const Person& p) { return p.age >= 18; });
\end{lstlisting}

Trzeci argument \texttt{ranges::sort} to \emph{projekcja}:
\texttt{\&Person::age} pobiera klucz sortowania, a domyślny komparator
porównuje klucze. \texttt{views::filter} nie buduje automatycznie nowego
wektora; jest leniwym widokiem, który odwołuje się do danych źródłowych.
Przy widokach szczególnie uważnie trzeba kontrolować czas życia
kontenera i przechwyceń predykatu.

\section{\texttt{std::ref}, \texttt{std::bind\_front} i częściowe wiązanie argumentów}
\label{sec:bind-front}

\begin{lstlisting}
#include <functional>

int add(int x, int y) { return x + y; }

auto add_ten = std::bind_front(add, 10); // C++20
int x = add_ten(5);                      // 15
\end{lstlisting}

\texttt{std::bind\_front} przechowuje funkcję i ustalone argumenty,
a pozostałe przyjmuje podczas wywołania. W wielu zastosowaniach
równoważna lambda jest czytelniejsza:

\begin{lstlisting}
auto add_ten_lambda = [](int x) { return add(10, x); };
\end{lstlisting}

Domyślnie argumenty przechowywane przez \texttt{bind\_front} są
przechowywane przez wartość, z uwzględnieniem przenoszenia. Jeśli
potrzebujemy odwołania do oryginalnego obiektu, stosujemy
\texttt{std::ref} (a dla dostępu tylko do odczytu
\texttt{std::cref}), pamiętając o czasie jego życia:

\begin{lstlisting}
#include <functional>

struct Accumulator {
    int total = 0;
    void add(int n) { total += n; }
};

Accumulator acc;
auto update = std::bind_front(&Accumulator::add, std::ref(acc));
update(7); // acc.total == 7
\end{lstlisting}

Starsze \texttt{std::bind} jest dostępne, ale ma bardziej skomplikowane
reguły dotyczące argumentów i podstawiania. W nowym kodzie zwykle
preferujemy lambdy lub \texttt{std::bind\_front}.

\section{Praktyczny projekt: rejestr callbacków}
\label{sec:callbacks}

Rozważmy klasę powiadamiającą o zmianie wartości. Do przechowywania
procedury obsługi zdarzenia użyjemy \texttt{std::function}. Dzięki
temu można ją zastąpić inną nawet po utworzeniu obiektu.

\begin{lstlisting}
#include <functional>
#include <utility>

class Counter {
public:
    using Callback = std::function<void(int)>;

    void on_change(Callback callback) {
        callback_ = std::move(callback);
    }

    void increment() {
        ++value_;
        if (callback_) callback_(value_);
    }

    int value() const noexcept { return value_; }

private:
    int value_ = 0;
    Callback callback_;
};

Counter counter;
counter.on_change([](int value) {
    // report the new value
});
counter.increment();
\end{lstlisting}

Kod w pojedynczym wątku jest prosty i wystarczający dla wielu
interfejsów zdarzeniowych. W programie wielowątkowym równoczesna
zmiana \texttt{callback\_} oraz wykonywanie \texttt{increment()}
wymaga dodatkowej synchronizacji. \texttt{std::function} nie zapewnia
automatycznie ochrony wewnętrznego stanu przed wyścigami danych.

\subsection{Callback a czas życia właściciela}

Typowy przebieg to: rejestrujemy callback, właściciel danych zostaje
zniszczony, a ktoś później wywołuje callback. Jeśli lambda zachowała
samą referencję lub surowy wskaźnik, dostęp byłby wtedy błędny.
Słaby wskaźnik w poniższym przykładzie pozwala sprawdzić czas życia.

\begin{lstlisting}
#include <memory>

struct Worker : std::enable_shared_from_this<Worker> {
    int total = 0;

    auto callback() {
        std::weak_ptr<Worker> weak = shared_from_this();
        return [weak](int n) {
            if (auto self = weak.lock()) {
                self->total += n;
            }
        };
    }
};
\end{lstlisting}

Przykład zakłada, że obiekt \texttt{Worker} jest już zarządzany przez
\texttt{std::shared\_ptr} w chwili wywołania \texttt{callback()}.
\texttt{weak\_ptr} pozwala sprawdzić, czy obiekt wciąż istnieje,
bez przedłużania jego życia przez przechwycenie silnego wskaźnika.
To ważne również dla unikania cykli własności, np. gdy sam obiekt
przechowuje callback odnoszący się do niego.

\section{Rekurencyjne lambdy bez \texttt{std::function}}

Lambda zwykle nie zna bezpośrednio własnej nazwy podczas definiowania.
Można jednak przekazywać obiekt wywoływalny jako argument:

\begin{lstlisting}
auto factorial = [](auto&& self, unsigned n) -> unsigned {
    return n <= 1 ? 1 : n * self(self, n - 1);
};

unsigned x = factorial(factorial, 5); // 120
\end{lstlisting}

To odmiana wzorca \emph{self-parameter}. Dla realnych zastosowań
rekurencyjnych warto ponadto zwrócić uwagę na przepełnienie liczb
całkowitych i maksymalną głębokość stosu. W C++20 metoda ta pozwala
uniknąć pośredniej alokacji i wymazywania typu, których mogłaby wymagać
rekurencja oparta na \texttt{std::function}.

\section{Lambdy w metaprogramowaniu: \texttt{std::visit} i \texttt{std::apply}}
\label{sec:lambda-visit-apply}

Dwa przydatne zastosowania łączą lambdy z mechanizmami z
rozdziałów~\ref{chap:cpp14-type-erasure} i~\ref{chap:variadic-templates-cpp}.
\texttt{std::visit} wybiera wywołanie na podstawie aktywnego typu
\texttt{std::variant}, natomiast \texttt{std::apply} rozwija argumenty
z krotki lub podobnej struktury.

\subsection{Zestaw przeciążeń oparty na dziedziczeniu}

\begin{lstlisting}
#include <string>
#include <variant>

template <class... Fs>
struct overloaded : Fs... {
    using Fs::operator()...;
};

std::variant<int, std::string> value = std::string{"Hello"};

auto stringify = overloaded{
    [](int x) { return std::to_string(x); },
    [](const std::string& x) { return x; }
};

std::string result = std::visit(stringify, value); // "Hello"
\end{lstlisting}

Każda lambda dostarcza jeden \texttt{operator()}. Dziedziczenie
wariadyczne i deklaracja \texttt{using Fs::operator()...} tworzą zbiór
przeciążeń, spośród których kompilator wybiera odpowiednią funkcję.
W tym przykładzie wszystkie warianty wywołania zwracają ten sam typ,
\texttt{std::string}, czego wymaga tu pojedyncze wywołanie
\texttt{std::visit}. Konstrukcja jest przydatna przy przetwarzaniu
stanów programu i heterogenicznych danych.

\subsection{Rozwijanie krotki przez \texttt{std::apply}}

\begin{lstlisting}
#include <tuple>

std::tuple<int, int, int> data{2, 3, 5};
int sum = std::apply(
    [](auto... values) { return (0 + ... + values); },
    data);
// sum == 10
\end{lstlisting}

Tutaj \texttt{std::apply} bierze elementy krotki i wywołuje lambdę
jak zwykłą funkcję trójargumentową. Sama lambda wykorzystuje
\emph{fold expression}, a więc nie trzeba pisać rekurencyjnego kodu.

\section{Pakiety przechwyceń i obliczenia odroczone (C++20)}

C++20 pozwala stosować \emph{pack expansion} również w przechwytywaniu
inicjalizującym. Możemy w ten sposób zbudować jeden obiekt domknięcia,
który przechowuje cały pakiet wartości do późniejszego przetworzenia:

\begin{lstlisting}
#include <concepts>

template <std::integral... Ts>
auto delayed_sum(Ts... xs) {
    return [...stored = xs] {
        return (0 + ... + stored);
    };
}

auto job = delayed_sum(2, 3, 5);
int result = job(); // 10
\end{lstlisting}

Lambdę można zwrócić z funkcji, ponieważ wartości są przechowywane
wewnątrz domknięcia. Przechwycenie przez referencję miałoby inne
wymagania dotyczące czasu życia. W rzeczywistych fabrykach callbacków
pakiet może zawierać obiekty niekopiowalne; wtedy należy zadbać
o przekazanie własności (np. przez \texttt{std::move} i
\texttt{std::forward}) oraz o wymagania interfejsu, który callback
przechowuje.

\section{Lambdy w obliczeniach stałych}

Jeśli ciało lambdy dopuszcza obliczanie stałe, jej wywołanie może
być częścią wyrażenia stałego. Od C++17 odpowiedni operator wywołania
jest domyślnie \texttt{constexpr}, gdy spełnia wymagania takiej funkcji:

\begin{lstlisting}
constexpr auto cube = [](int x) { return x * x * x; };
static_assert(cube(3) == 27);

// C++20: each call must be evaluated at compile time.
auto fourth_power = [](int x) consteval {
    return x * x * x * x;
};
static_assert(fourth_power(2) == 16);
\end{lstlisting}

\texttt{consteval} jest mocniejszą deklaracją niż \texttt{constexpr}:
wymaga, aby właściwe wywołania funkcji natychmiastowej spełniały
reguły obliczania stałego. Nie należy traktować tego jako gwarancji,
że każda zwykła lambda zostanie wykonana podczas kompilacji.

\section{Koncepcje dla obiektów wywoływalnych}
\label{sec:callable-concepts}

W C++20 można ograniczać typy wywoływalne m.in. następującymi
koncepcjami:
\begin{itemize}
\item \texttt{std::invocable} -- poprawność wywołania z danymi argumentami;
\item \texttt{std::regular\_invocable} -- dodatkowe wymagania semantyczne;
\item \texttt{std::predicate} -- poprawnie wywoływalny predykat z
wynikiem dającym się traktować jako wartość logiczną.
\end{itemize}
Koncepcje nie dowodzą wszystkich własności semantycznych. Przykładowo
kompilator nie zweryfikuje matematycznie stabilności wyników predykatu.

\begin{lstlisting}
#include <concepts>
#include <functional>
#include <utility>

template <class F, class T>
    requires std::predicate<F&, const T&>
bool test(const T& value, F&& predicate) {
    return std::invoke(std::forward<F>(predicate), value);
}

bool ok = test(8, [](int x) { return x % 2 == 0; });
\end{lstlisting}

\textbf{Istotne rozróżnienie:} koncepcja weryfikuje przede wszystkim
własności typów i poprawność wyrażeń. Nie dowodzi matematycznie
własności predykatu, np. że comparator tworzy ścisły słaby porządek.
Programista nadal odpowiada za kontrakt semantyczny.

\section{Typowe błędy i trudniejsze przypadki}
\label{sec:lambda-pitfalls}

\subsection{Referencja przechwycona do operacji asynchronicznej}

\begin{lstlisting}
#include <future>

std::future<int> start() {
    int x = 7;
    // WRONG: the task could outlive x.
    // return std::async(std::launch::async, [&x] { return x; });
    return std::async(std::launch::async, [x] { return x; });
}
\end{lstlisting}

W zadaniach asynchronicznych preferuj przechwycenia przez wartość
lub świadomie zarządzaną własność, dopóki nie mamy ścisłej gwarancji,
że obiekt referowany przetrwa wykonanie zadania. Sama biblioteka
\texttt{std::async} nie wydłuża życia lokalnych zmiennych funkcji
uruchamiającej.

\subsection{Predykat sortowania nie może naruszać porządku}

\begin{lstlisting}
// WRONG: <= is not a strict ordering relation.
// std::sort(xs.begin(), xs.end(),
//           [](int a, int b) { return a <= b; });

std::sort(xs.begin(), xs.end(),
          [](int a, int b) { return a < b; });
\end{lstlisting}

Komparator do \texttt{sort} musi definiować \emph{ścisły słaby
porządek}. Przykładowo \texttt{comp(x,x)} ma dawać \texttt{false}.
Naruszenie wymagań algorytmu może prowadzić do niezdefiniowanego
zachowania, mimo że typy argumentów się zgadzają.

\subsection{Pomylenie kopii z oryginalnym stanem}

\begin{lstlisting}
int sum = 0;
std::for_each(xs.begin(), xs.end(), [sum](int x) mutable {
    sum += x; // modifies only the closure's copy
});
// outer sum remains 0

std::for_each(xs.begin(), xs.end(), [&sum](int x) {
    sum += x;
});
\end{lstlisting}

Nie należy zakładać, że \texttt{mutable} zmienia zewnętrzną
zmienną. Należy też pamiętać, że algorytmy
mogą kopiować obiekt funkcyjny; nie opierajmy semantyki algorytmu
na nieudokumentowanych szczegółach kopiowania predykatu.

\subsection{Przechwycenie wskaźnika nie przedłuża życia obiektu}

\begin{lstlisting}
int* p = /* pointer to some object */ nullptr;
auto f = [p] {
    if (p) return *p;
    return 0;
};
\end{lstlisting}

Wartość wskaźnika może zostać skopiowana do lambdy, ale wskazywany
obiekt nadal musi istnieć w chwili dereferencji. Przechwycenie
\texttt{shared\_ptr} przenosi własność współdzieloną, a
\texttt{weak\_ptr} pozwala ją warunkowo odzyskać. To nie jest to
samo co skopiowanie wskaźnika surowego.

\section{Zestawienie: kiedy stosować którą konstrukcję?}

\begin{center}
\begin{tabular}{|p{5.7cm}|p{6.6cm}|}
\hline
\textbf{Potrzeba} & \textbf{Zalecany wybór} \\
\hline
Krótki predykat do algorytmu & Lambda z jawnymi przechwyceniami \\
\hline
Wywołanie bez stanu w API C & Zwykła funkcja albo zgodna lambda bez przechwyceń jako wskaźnik do funkcji \\
\hline
Wydajny callback o typie znanym w czasie kompilacji & Parametr szablonowy \texttt{F}, ewentualnie \texttt{std::invocable} \\
\hline
Callback z możliwością późniejszej podmiany & \texttt{std::function} (C++20, kopiowalny cel) \\
\hline
Przechowanie \texttt{unique\_ptr} wewnątrz lambdy & Init-capture i przeniesienie; nie przechowuj takiej lambdy w \texttt{std::function} w C++20 \\
\hline
Callback żyjący dłużej niż zmienne lokalne & Kopiuj potrzebny stan lub zarządzaj życiem przez RAII i smart pointery \\
\hline
Uniwersalne wywołanie dowolnego callable & \texttt{std::invoke} i \texttt{std::forward} \\
\hline
Częściowe podstawienie kilku argumentów & Lambda lub \texttt{std::bind\_front} \\
\hline
Bezpieczny predykat dla szablonu & \texttt{std::predicate} plus kontrola jego semantyki \\
\hline
\end{tabular}
\end{center}

\paragraph{Pytania kontrolne.}
\begin{enumerate}
\item Dlaczego dwie lambdy o identycznym kodzie mają różne typy?
\item Jaka jest różnica między \texttt{[x]}, \texttt{[\&x]} i \texttt{[x]() mutable}?
\item Kiedy bezpiecznie zwrócić lambdę przechwytującą zmienną lokalną?
\item Dlaczego lambda posiadająca \texttt{unique\_ptr} nie może być celem \texttt{std::function} w C++20?
\item Jak \texttt{std::invoke} wywołuje wskaźnik do metody klasy?
\item Czym różnią się \texttt{[this]} i \texttt{[*this]}?
\item Co może pójść źle, gdy \texttt{std::views::filter} korzysta z lokalnej zmiennej przez referencję?
\item Kiedy interfejs przyjmujący szablon \texttt{F} jest lepszy od \texttt{std::function}?
\item Dlaczego komparator \texttt{a <= b} nie jest poprawny dla \texttt{std::sort}?
\item W jakiej sytuacji przechwycenie \texttt{weak\_ptr} jest lepsze od \texttt{shared\_ptr}?
\end{enumerate}

\paragraph{Dalsza lektura.}
W dokumentacji \emph{cppreference} warto przeczytać hasła
\emph{Lambda expressions}, \emph{Function objects} i \emph{Ranges}.
Dokładne opisy narzędzi znajdują się również pod hasłami:
\begin{itemize}
\item \texttt{std::invoke} oraz \texttt{std::function};
\item \texttt{std::bind\_front} oraz \texttt{std::predicate}.
\end{itemize}

% ROZDZIAL 14 — Type erasure, typy sumaryczne i interfejsy generyczne (C++20).
% Fragment do zbiorczego skryptu, BEZ preambuly i \begin{document}.
% Wymagane: listings, amsmath, amssymb, booktabs, tabularx, array, hyperref.
% Dolaczenie: \input{rozdzial_14_type_erasure_typy_sumaryczne_cpp.tex}
% Kompletne programy sa oznaczone w podpisach listingow.

\chapter{Wymazywanie typów i typy sumaryczne}
\label{chap:cpp14-type-erasure}

\section{Problem: jak reprezentować różne rodzaje wartości?}

W programowaniu generycznym często chcemy posługiwać się różnymi typami
przez wspólny interfejs. C++ oferuje kilka mechanizmów realizacji tego celu,
ale mechanizmy te nie są zamienne.
Wyobraźmy sobie aplikację z trzema poleceniami: zapisaniem pliku,
dodaniem dwóch liczb i zakończeniem pracy. Gdy znamy wszystkie trzy
rodzaje poleceń, możemy je wymienić w \texttt{std::variant}. Jeśli
zewnętrzne moduły mają dodawać nowe polecenia bez zmiany tej listy,
potrzebujemy otwartego interfejsu, np. klasy bazowej z funkcją wirtualną.
Wybór mechanizmu zaczynamy więc od pytań:
\begin{enumerate}
\item Czy zestaw dozwolonych typów jest \emph{zamknięty} (ustalony z góry),
      czy \emph{otwarty} (może zostać rozszerzony o nowy typ)?
\item Czy chcemy przechować \emph{wartość}, czy jedynie referencję/uchwyt
      do istniejącego obiektu?
\item Czy wywoływane operacje mają być znane już podczas kompilacji,
      czy wybierane dynamicznie w czasie działania programu?
\item Czy dopuszczamy alokację dynamiczną, kopiowanie oraz koszt
      pośredniego wywołania funkcji?
\end{enumerate}

\emph{Typ sumaryczny} (ang. \emph{sum type}) reprezentuje alternatywę:
wartość należy do jednego z kilku przypadków. W C++17 taką rolę pełni
\texttt{std::variant<Ts...>}; \texttt{std::optional<T>} opisuje wariant
,,brak wartości albo wartość \texttt{T}''. Natomiast \emph{type erasure}
(\emph{wymazywanie typu}) pozwala ukryć konkretny typ obiektu za
interfejsem nieparametryzowanym przez ten typ. Standardowe przykłady
to \texttt{std::any} i \texttt{std::function}.

\begin{center}
\small
\begin{tabularx}{\textwidth}{@{}lXX@{}}
\toprule
\textbf{Mechanizm} & \textbf{Informacja o typie} & \textbf{Zasadnicze zastosowanie} \\
\midrule
\texttt{optional<T>} & Jedna możliwość \texttt{T} lub brak & Wynik opcjonalny, wartość nieobecna \\
\texttt{variant<Ts...>} & Zamknięty zbiór alternatyw & Zdarzenia, komendy, model stanu \\
\texttt{any} & Typ ustalany w czasie wykonania & Heterogeniczne wartości bez wspólnego interfejsu operacji \\
\texttt{function<R(Args...)>} & Dowolny zgodny obiekt wywoływalny & Przechowywany callback, funkcja jako wartość \\
\texttt{unique\_ptr<Base>} & Otwarta hierarchia klas & Polimorfizm dynamiczny i tożsamość obiektu \\
Szablon/\texttt{concept} & Typ znany przy instancjowaniu & Polimorfizm statyczny, specjalizacja kodu \\
\bottomrule
\end{tabularx}
\end{center}

Najprostsza ściąga to: \texttt{optional<int>} oznacza liczbę albo jej
brak, \texttt{variant<int, std::string>} jedną z dwóch znanych
postaci wartości, a \texttt{any} wartość o typie, którego autor
pojemnika nie musi znać. \texttt{function<void()>} przechowuje
różne działania, które można wywołać w taki sam sposób. W kolejnych
sekcjach zobaczymy, co z tych wyborów wynika dla sprawdzania typów,
własności danych i kosztu wykonania.

\section{\texttt{std::optional<T>}: wartość, której może nie być}

\begin{center}\small
\begin{tabularx}{\textwidth}{@{}p{0.33\textwidth}X@{}}
\toprule
Typ lub operacja & Co reprezentuje \\
\midrule
\code{std::optional<T>} & Wartość typu \code{T} albo jej brak; nie opisuje przyczyny błędu. \\
\code{has_value()} & Sprawdza, czy \code{optional} przechowuje wartość. \\
\code{value()} & Zwraca wartość lub zgłasza \code{std::bad_optional_access}, gdy jej brak. \\
\code{value_or(fallback)} & Zwraca wartość albo podany zamiennik, konwertowany do \code{T}. \\
\code{std::variant<Ts...>} & Dokładnie jedną z ustalonych alternatyw typów. \\
\code{std::any} & Wartość typu ustalonego dopiero w czasie wykonania; przy odczycie trzeba sprawdzić typ. \\
\bottomrule
\end{tabularx}
\end{center}
\label{sec:cpp14-optional}

Typ \texttt{std::optional<T>} (C++17, nagłówek \texttt{<optional>})
przechowuje bezpośrednio obiekt typu \texttt{T}, jeżeli jest on obecny,
lub reprezentuje stan pusty. Jest to pełnoprawna \emph{wartość}, a nie
wskaźnik. W szczególności pusty \texttt{optional} nie oznacza wskaźnika
zerowego, a niepusty nie wymaga automatycznie osobnej alokacji dynamicznej.

\begin{lstlisting}[caption={Najważniejsze operacje optional (C++17); kompletny program.}]
#include <cassert>
#include <optional>
#include <string>

int main() {
    std::optional<int> x;         // stan pusty
    assert(!x.has_value());
    assert(x == std::nullopt);

    x = 42;
    assert(x.has_value());
    assert(*x == 42);
    assert(x.value() == 42);
    assert(x.value_or(-1) == 42);

    x.reset();
    assert(x.value_or(-1) == -1);
    x.emplace(7);
    assert(*x == 7);

    std::optional<std::string> name{std::in_place, 3, 'a'};
    assert(*name == "aaa");
}
\end{lstlisting}

Dostęp przez \texttt{*x} i \texttt{x->...} wymaga obecności wartości:
\emph{użycie pustego obiektu w taki sposób jest niezdefiniowanym
zachowaniem}. Metoda \texttt{value()} wykonuje sprawdzenie i dla stanu
pustego zgłasza \texttt{std::bad\_optional\_access}. Warto pamiętać,
że \texttt{value\_or(default)} przyjmuje wartość domyślną jako argument
funkcji: wyrażenie tworzące argument jest obliczane również wtedy,
gdy \texttt{optional} zawiera wartość. Może ponadto następować kopiowanie
lub przeniesienie przechowywanego obiektu.

\subsection{Praktyczny przykład: parsowanie liczby}

Funkcja parsująca dane wejściowe może zwrócić \texttt{optional}, jeśli
interesuje nas tylko informacja ,,udało się/nie udało się''.
\texttt{std::from\_chars} unika wyjątków i nie wymaga tworzenia dodatkowego
łańcucha zakończonego zerem. W przykładzie wymagamy, aby \emph{cały}
łańcuch był prawidłową liczbą.

\begin{lstlisting}[caption={Parsowanie int do optional (C++17); kompletny program.}]
#include <cassert>
#include <charconv>
#include <optional>
#include <string_view>
#include <system_error>

std::optional<int> parse_int(std::string_view text) {
    int result = 0;
    auto [ptr, ec] = std::from_chars(
        text.data(), text.data() + text.size(), result);
    if (ec != std::errc{} || ptr != text.data() + text.size())
        return std::nullopt;
    return result;
}

int main() {
    assert(parse_int("123") == 123);
    assert(parse_int("123x") == std::nullopt);
    assert(parse_int("") == std::nullopt);
    assert(parse_int("99999999999999999999") == std::nullopt);
}
\end{lstlisting}

\subsection{Własność, referencje i koszt}

\texttt{optional<T>} posiada swoją wartość i zwykle zajmuje więcej pamięci
niż sam \texttt{T}, ponieważ musi przechowywać informację, czy wartość
istnieje. Rozmiar zależy od typu i implementacji. Standard C++20
nie udostępnia \texttt{std::optional<T\&>}. Jeśli wynik ma wskazywać na
istniejący obiekt, można użyć wskaźnika nieposiadającego własności albo
\texttt{std::optional<std::reference\_wrapper<T>>}. W obu przypadkach
programista odpowiada za czas życia wskazywanego obiektu.

\begin{lstlisting}[caption={Opcjonalna referencja: brak przenoszenia własności; fragment.}]
#include <functional>
#include <optional>
#include <string>

std::optional<std::reference_wrapper<const std::string>>
find_name(bool exists, const std::string& name) {
    if (exists) return std::cref(name);
    return std::nullopt;
}
// Zwracana referencja jest poprawna tylko gdy `name` nadal istnieje.
\end{lstlisting}

Gdy brak wartości jest sytuacją normalną, \texttt{optional} najczęściej
czytelniej wyraża intencję niż wyjątkowe zakończenie funkcji. Jeżeli
konieczna jest informacja \emph{dlaczego} nie uzyskano wyniku, należy
zaprojektować osobny typ błędu, użyć \texttt{variant} lub (w C++23)
\texttt{std::expected<T,E>}.

\section{\texttt{std::variant<Ts...>}: zamknięta alternatywa typów}
\label{sec:cpp14-variant}

\texttt{std::variant<Ts...>} (C++17, nagłówek \texttt{<variant>})
przechowuje \emph{jedną aktywną alternatywę} z listy typów \texttt{Ts...}.
Nie jest odpowiednikiem \texttt{std::tuple<Ts...>}: krotka zawiera
\emph{wszystkie} składowe, a wariant dokładnie jeden z dopuszczonych
przypadków (poza szczególnym stanem \texttt{valueless\_by\_exception()}).

\begin{lstlisting}[caption={Tworzenie i badanie aktywnej alternatywy (C++17); kompletny program.}]
#include <cassert>
#include <string>
#include <variant>

int main() {
    using Value = std::variant<int, std::string>;
    Value x = 7;
    assert(x.index() == 0);
    assert(std::holds_alternative<int>(x));
    assert(std::get<int>(x) == 7);

    x = std::string("hello");
    assert(x.index() == 1);
    assert(std::get_if<int>(&x) == nullptr);
    assert(std::get_if<std::string>(&x) != nullptr);

    x.emplace<0>(19);
    assert(std::get<0>(x) == 19);
}
\end{lstlisting}

Niepoprawne \texttt{std::get<T>(x)} (gdy aktywny jest inny typ)
zgłasza \texttt{std::bad\_variant\_access}. Alternatywa
\texttt{std::get\_if<T>(\&x)} zwraca \texttt{nullptr}, a zatem nadaje się
do sprawdzania bez wyjątku. \texttt{std::get<T>} i
\texttt{std::holds\_alternative<T>} wymagają, aby \texttt{T} występował
\emph{dokładnie raz} w liście alternatyw; dla powtarzających się typów
używamy indeksów. Nie należy także zapominać o możliwym niepowodzeniu
operacji konwersji przy konstruowaniu wariantu.

\subsection{Domyślna alternatywa i \texttt{std::monostate}}

Domyślnie konstruowany wariant próbuje skonstruować pierwszą alternatywę.
Jeżeli chcemy jawnie reprezentować stan pusty, można umieścić na początku
\texttt{std::monostate}; zapewnia on także domyślną konstruowalność
wariantu, gdy inny pierwszy typ nie jest domyślnie konstruowalny.

\begin{lstlisting}[caption={Jawny stan początkowy monostate; kompletny program.}]
#include <cassert>
#include <string>
#include <variant>

int main() {
    using State = std::variant<std::monostate, int, std::string>;
    State s;
    assert(std::holds_alternative<std::monostate>(s));
    s = std::string("ready");
    assert(std::get<std::string>(s) == "ready");
}
\end{lstlisting}

\subsection{Odwiedzanie wariantu: \texttt{std::visit}}

Samo wielokrotne sprawdzanie \texttt{holds\_alternative} szybko
prowadzi do rozbudowanych instrukcji \texttt{if}. Mechanizm
\texttt{std::visit(visitor, variant)} wywołuje odwiedzającego z aktywną
wartością odpowiedniego typu. Gdy odwiedzający jest lambdą generyczną,
kompilator instancjuje jej operator wywołania dla różnych alternatyw.

\begin{lstlisting}[caption={Jeden odwiedzający generyczny; kompletny program.}]
#include <cassert>
#include <cstddef>
#include <string>
#include <type_traits>
#include <variant>

int main() {
    std::variant<int, double, std::string> v = "hi";

    const std::size_t size = std::visit([](const auto& x) {
        using T = std::decay_t<decltype(x)>;
        if constexpr (std::is_same_v<T, std::string>)
            return x.size();
        else
            return sizeof(x);
    }, v);

    assert(size == 2);
}
\end{lstlisting}

Do osobnej obsługi poszczególnych typów wygodny jest wzorzec
\texttt{overloaded}, znany również z rozdziału o lambdach.
Wykorzystuje on dziedziczenie po kilku obiektach funkcyjnych oraz
ujawnienie wszystkich operatorów wywołania za pomocą \texttt{using}.

\begin{lstlisting}[caption={overloaded oraz std::visit (C++17); kompletny program.}]
#include <cassert>
#include <cstddef>
#include <string>
#include <variant>

template<class... Fs>
struct overloaded : Fs... {
    using Fs::operator()...;
};
template<class... Fs> overloaded(Fs...) -> overloaded<Fs...>;

int main() {
    using Data = std::variant<int, double, std::string>;
    Data data = std::string("abc");

    auto length = std::visit(overloaded{
        [](int) { return std::size_t{1}; },
        [](double) { return std::size_t{1}; },
        [](const std::string& s) { return s.size(); }
    }, data);
    assert(length == 3);
}
\end{lstlisting}

W przypadku odwiedzania zwracającego wartość wymagany jest zgodny typ
wyniku wszystkich możliwych wywołań (w praktyce należy zapewnić
identyczny typ i kategorię wartości zwracanych wyrażeń, z wyjątkiem
szczególnych reguł dla \texttt{void}). Przy wielu wariantach
\texttt{std::visit} uwzględnia \emph{kombinacje} alternatyw. To daje
czytelny kod przy małej liczbie przypadków, ale może zwiększyć liczbę
instancjacji szablonów.

\begin{lstlisting}[caption={Odwiedzenie dwóch wariantów naraz; kompletny program.}]
#include <cassert>
#include <variant>

int main() {
    std::variant<int, double> a = 3;
    std::variant<int, double> b = 2.5;
    const double sum = std::visit(
        [](const auto& x, const auto& y) -> double {
            return static_cast<double>(x) + static_cast<double>(y);
        }, a, b);
    assert(sum == 5.5);
}
\end{lstlisting}

\subsection{Modyfikacja stanu i kwalifikacja \texttt{const}}

\texttt{std::visit} respektuje typ i kategorię wartości przekazanego
wariantu. Odwiedzając obiekt przez \texttt{variant\&}, możemy modyfikować
aktywną alternatywę; przez \texttt{const variant\&} uzyskujemy
referencję stałą. Przy przekazaniu \texttt{std::move(v)} odwiedzający
może otrzymać referencję rvalue i przenieść przechowywaną wartość.
Przeniesienie nie musi zmieniać aktywnego indeksu wariantu, ale zawartość
przeniesionego obiektu może zostać zmieniona zgodnie z semantyką typu.

\begin{lstlisting}[caption={Modyfikacja aktywnej alternatywy; kompletny program.}]
#include <cassert>
#include <variant>

int main() {
    std::variant<int, double> v = 10;
    std::visit([](auto& x) { x += 1; }, v);
    assert(std::get<int>(v) == 11);
}
\end{lstlisting}

\subsection{Stan \texttt{valueless\_by\_exception}}

Wyjątek podczas przypisania lub tworzenia nowej wartości może
pozostawić \texttt{variant} bez aktywnej alternatywy. Ten stan
wykrywa metoda \texttt{valueless\_by\_exception()}.
Wtedy \texttt{index()} zwraca \texttt{std::variant\_npos}.
Wywołanie \texttt{std::visit} zgłasza \texttt{std::bad\_variant\_access}.
Nie jest to zwykła ,,pusta'' alternatywa. Jeśli jest potrzebna,
można użyć \texttt{std::monostate}.

\paragraph{Ważne ograniczenia.}
\texttt{variant} nie przechowuje typów referencyjnych ani \texttt{void};
nie można także dodać dowolnego nowego typu bez zmiany listy alternatyw.
Duży wariant może zajmować pamięć odpowiadającą największej alternatywie
plus informację o aktywnym przypadku i ewentualne wyrównanie; dokładny
rozmiar nie jest stałą wielkością określoną przez standard.

\section{Przykład projektowy: interpreter prostych poleceń}
\label{sec:cpp14-commands}

Rozpatrzmy miniaturowy model poleceń. Każdy rodzaj polecenia ma inne pola,
ale wszystkie należą do zamkniętej rodziny. \texttt{variant} oddziela
\emph{dane} od \emph{operacji}, a \texttt{visit} pozwala jednoznacznie
obsłużyć każdy przypadek.

\begin{lstlisting}[caption={Modelowanie poleceń za pomocą variant; kompletny program.}]
#include <cassert>
#include <string>
#include <variant>
#include <vector>

struct Add  { int a, b; };
struct Save { std::string filename; };
struct Quit {};
using Command = std::variant<Add, Save, Quit>;

struct Execute {
    int& total;
    bool& stopped;

    void operator()(const Add& c) const {
        total += c.a + c.b;
    }
    void operator()(const Save& c) const {
        if (c.filename.empty()) stopped = true;
    }
    void operator()(const Quit&) const {
        stopped = true;
    }
};

int main() {
    std::vector<Command> commands{
        Add{2, 3}, Save{"result.txt"}, Add{4, 5}, Quit{}
    };
    int total = 0;
    bool stopped = false;
    for (const auto& command : commands) {
        if (stopped) break;
        std::visit(Execute{total, stopped}, command);
    }
    assert(total == 14);
    assert(stopped);
}
\end{lstlisting}

Dodanie nowego typu polecenia, np. \texttt{Delete}, wymaga rozszerzenia
\texttt{Command} i obsługi nowego przypadku. W dobrze skonstruowanym
odwiedzającym brakująca obsługa będzie błędem kompilacji. To zaleta
typów sumarycznych przy modelowaniu \emph{zamkniętego} katalogu zdarzeń.

\section{\texttt{std::any}: wartość o nieznanym typie}
\label{sec:cpp14-any}

\texttt{std::any} (C++17, nagłówek \texttt{<any>}) może przechowywać
obiekt dowolnego typu spełniającego wymaganie kopiowalności i pozwala
później odzyskać dokładny typ przez \texttt{std::any\_cast}. To klasyczna
implementacja \emph{type erasure}: użytkownik obiektu \texttt{any}
nie zna szablonowego typu przechowywanej wartości.

\begin{lstlisting}[caption={std::any i bezpieczne wydobycie typu (C++17); kompletny program.}]
#include <any>
#include <cassert>
#include <string>

int main() {
    std::any value = 17;
    assert(value.has_value());
    assert(std::any_cast<int>(value) == 17);

    if (auto ptr = std::any_cast<int>(&value)) {
        *ptr += 1;
    }
    assert(std::any_cast<int>(value) == 18);

    value.emplace<std::string>("text");
    assert(std::any_cast<int>(&value) == nullptr);
    assert(std::any_cast<const std::string&>(value) == "text");

    value.reset();
    assert(!value.has_value());
}
\end{lstlisting}

\texttt{std::any\_cast<T>(value)} zgłasza \texttt{std::bad\_any\_cast},
jeśli aktywny typ nie odpowiada żądanemu. Wariant wskaźnikowy
\texttt{std::any\_cast<T>(\&value)} zwraca \texttt{nullptr}; jest zatem
wygodny do ostrożnego sprawdzania typu. \texttt{type()} udostępnia
\texttt{std::type\_info} opisujący aktualnie przechowywany typ,
a dla pustego obiektu zwraca \texttt{typeid(void)}.

\begin{lstlisting}[caption={\texttt{any\_cast} nie jest konwersją numeryczną; kompletny program.}]
#include <any>
#include <cassert>
#include <typeinfo>

int main() {
    std::any value = 42;
    assert(value.type() == typeid(int));
    assert(std::any_cast<double>(&value) == nullptr);
    assert(std::any_cast<int>(&value) != nullptr);
}
\end{lstlisting}

\texttt{any} przechowuje \emph{wartość}, nie udostępnia jednak uniwersalnego
operatora \texttt{+}, porównania czy drukowania tej wartości. Aby wykonać
operację, trzeba znać typ albo dostarczyć odrębny mechanizm dyspozycji.
Ponadto przechowywany typ musi być kopiowalny: \texttt{std::any}
nie nadaje się do bezpośredniego przechowywania \texttt{std::unique\_ptr<T>}.
Niektóre implementacje wykorzystują \emph{small-object optimization},
ale standard nie obiecuje uniknięcia alokacji dla konkretnego małego typu.

\paragraph{Kiedy używać?}
\texttt{std::any} ma sens m.in. w luźno typowanych zasobnikach kontekstu,
gdy autor zasobnika nie zna listy możliwych typów, a odbiorca umie
jawnie odtworzyć typ właściwego wpisu. Jeśli rodzaje wartości są znane,
\texttt{variant} daje z reguły silniejszą kontrolę statyczną i czytelniejszą
obsługę przypadków. Unikaj \texttt{any} jako sposobu na ukrycie
źle określonych kontraktów interfejsu.

\section{\texttt{std::function}: wymazywanie typu obiektu wywoływalnego}
\label{sec:cpp14-function}

\texttt{std::function<R(Args...)>} (C++11, nagłówek
\texttt{<functional>}) przechowuje kopiowalny obiekt wywoływalny zgodny
z podaną sygnaturą. Może to być wskaźnik do funkcji, funktor,
wyrażenie lambda, a także obiekt związany z metodą składową. Konkretny
typ obiektu zostaje ukryty, ale interfejs jego wywołania jest znany:
\texttt{R(Args...)}.

\begin{lstlisting}[caption={Jeden typ std::function dla różnych obiektów wywoływalnych; kompletny program.}]
#include <cassert>
#include <functional>

int times_two(int x) { return 2 * x; }
struct PlusThree {
    int operator()(int x) const { return x + 3; }
};

int main() {
    std::function<int(int)> f = times_two;
    assert(f(4) == 8);

    f = PlusThree{};
    assert(f(4) == 7);

    int delta = 5;
    f = [delta](int x) { return x + delta; };
    assert(f(4) == 9);
}
\end{lstlisting}

\subsection{Pusty obiekt i koszt abstrakcji}

Domyślnie skonstruowany \texttt{std::function} jest pusty. Jego
\texttt{operator bool()} informuje, czy istnieje obiekt wywoływalny;
próba wywołania pustego obiektu zgłasza \texttt{std::bad\_function\_call}.
W porównaniu z bezpośrednim wywołaniem szablonu \texttt{std::function}
może wymagać pośredniego wywołania i dodatkowego przechowywania stanu,
a czasami także alokacji dynamicznej. Nie jest jednak poprawne twierdzenie,
że każdy \texttt{std::function} \emph{zawsze} wykonuje alokację.

\begin{lstlisting}[caption={Pusty std::function i jawne sprawdzanie; kompletny program.}]
#include <cassert>
#include <functional>

int main() {
    std::function<void()> callback;
    assert(!callback);
    callback = [] {};
    assert(static_cast<bool>(callback));
    callback();
    callback = nullptr;
    assert(!callback);
}
\end{lstlisting}

\subsection{Ograniczenie kopiowalności i czas życia przechwyceń}

W C++20 \texttt{std::function} wymaga kopiowalności typu przechowywanego
obiektu wywoływalnego. Dlatego lambda zawierająca \emph{move-only}
\texttt{std::unique\_ptr} nie może być bezpośrednio umieszczona w
\texttt{std::function}. W C++23 pojawił się
\texttt{std::move\_only\_function} dla takiego zastosowania.

Przechwytywanie przez referencję wymaga ponadto, aby obiekt odniesienia
pozostawał żywy podczas każdego późniejszego wywołania. Sam
\texttt{std::function} kontroluje czas życia \emph{przechowywanej lambdy},
ale nie przedłuża automatycznie czasu życia obiektów, do których lambda
odwołuje się przez referencje lub nieposiadające własności wskaźniki.

\begin{lstlisting}[caption={Callback utrzymujący obiekt przy życiu przez \texttt{shared\_ptr}; kompletny program.}]
#include <cassert>
#include <functional>
#include <memory>

int main() {
    auto state = std::make_shared<int>(10);
    std::function<int()> get_value = [state] {
        return *state;
    };
    state.reset();
    assert(get_value() == 10);
}
\end{lstlisting}

Trwałe przechwycenie \texttt{shared\_ptr} może z kolei utworzyć cykl
własności, np. gdy obiekt przechowuje callback, a callback przechwytuje
\texttt{shared\_ptr} do tego samego obiektu. Wówczas często lepiej
przechwycić \texttt{weak\_ptr} i uzyskać tymczasową własność przez
\texttt{lock()}.

\begin{lstlisting}[caption={Callback niewymuszający przedłużania życia obiektu; kompletny program.}]
#include <cassert>
#include <functional>
#include <memory>
#include <optional>

int main() {
    auto state = std::make_shared<int>(42);
    std::weak_ptr<int> weak = state;

    std::function<std::optional<int>()> read = [weak]()
        -> std::optional<int> {
        if (auto shared = weak.lock()) return *shared;
        return std::nullopt;
    };

    assert(read() == 42);
    state.reset();
    assert(!read().has_value());
}
\end{lstlisting}

\subsection{Szablon, wskaźnik do funkcji czy \texttt{std::function}?}

Jeżeli funkcja korzysta z obiektu wywoływalnego tylko podczas bieżącego
wywołania, często najprościej użyć szablonu. Pozwala to kompilatorowi
znać konkretny typ i daje okazję do \emph{inliningu}. Wskaźnik do funkcji
jest lekki, ale nie przechowa samodzielnie stanu przechwycenia lambdy.
\texttt{std::function} ułatwia natomiast przechowywanie i zastępowanie
heterogenicznych callbacków o wspólnej sygnaturze.

\begin{lstlisting}[caption={Wywołanie generyczne bez type erasure; kompletny program.}]
#include <cassert>
#include <functional>
#include <utility>

template<class F, class... Args>
decltype(auto) call_now(F&& f, Args&&... args) {
    return std::invoke(std::forward<F>(f),
                       std::forward<Args>(args)...);
}

int main() {
    int factor = 3;
    auto f = [factor](int x) { return factor * x; };
    assert(call_now(f, 7) == 21);
}
\end{lstlisting}

\section{Własny mechanizm type erasure: \texttt{Printable}}
\label{sec:cpp14-custom-erasure}

To rozszerzenie pokazuje, co może znajdować się pod gotowym
opakowaniem wymazującym typ. Przy pierwszej lekturze wystarczy
rozumieć, kiedy użyć \texttt{variant}, \texttt{any} lub
\texttt{function}; własną implementację można przeczytać później.

Pokażemy teraz, jak samodzielnie zbudować typ o stałym interfejsie,
który może przechowywać różne wartości drukowalne do strumienia. Celowo
nie wykorzystamy dziedziczenia typów użytkownika po wspólnej klasie
bazowej. Użytkownik powinien móc przekazać \texttt{int},
\texttt{std::string} albo własny typ, dla którego działa
\texttt{operator<<}. Wymagamy przy tym \emph{semantyki wartości}:
kopiowanie \texttt{Printable} ma kopiować przechowywany obiekt.

\paragraph{Idea implementacji.}
Wewnątrz opakowania znajdują się:
\begin{enumerate}
\item abstrakcyjny interfejs \texttt{Concept} z operacjami \texttt{print}
      i \texttt{clone};
\item szablonowy model \texttt{Model<T>}, który przechowuje konkretną
      wartość i realizuje interfejs;
\item wskaźnik \texttt{unique\_ptr<Concept>}, ukrywający konkretny
      typ w miejscu użycia \texttt{Printable}.
\end{enumerate}

\begin{lstlisting}[caption={Własne kopiowalne type erasure (C++20); kompletny program.}]
#include <cassert>
#include <concepts>
#include <memory>
#include <ostream>
#include <sstream>
#include <string>
#include <type_traits>
#include <utility>

class Printable {
    struct Concept {
        virtual ~Concept() = default;
        virtual void print(std::ostream&) const = 0;
        virtual std::unique_ptr<Concept> clone() const = 0;
    };

    template<class T>
    struct Model final : Concept {
        T value;
        template<class U>
        explicit Model(U&& x) : value(std::forward<U>(x)) {}
        void print(std::ostream& os) const override {
            os << value;
        }
        std::unique_ptr<Concept> clone() const override {
            return std::make_unique<Model<T>>(value);
        }
    };

    std::unique_ptr<Concept> self_;

public:
    Printable() = default;

    template<class T>
      requires (!std::same_as<std::remove_cvref_t<T>, Printable> &&
                std::copy_constructible<std::decay_t<T>> &&
                requires(std::ostream& os,
                         const std::decay_t<T>& value) {
                    { os << value } -> std::same_as<std::ostream&>;
                })
    Printable(T&& value)
        : self_(std::make_unique<Model<std::decay_t<T>>>(
              std::forward<T>(value))) {}

    Printable(const Printable& other)
        : self_(other.self_ ? other.self_->clone() : nullptr) {}

    Printable(Printable&&) noexcept = default;

    Printable& operator=(const Printable& other) {
        if (this != &other) {
            Printable tmp(other); // najpierw bezpieczna kopia
            swap(tmp);            // potem nie-rzucajaca zamiana
        }
        return *this;
    }

    Printable& operator=(Printable&&) noexcept = default;

    void swap(Printable& other) noexcept {
        self_.swap(other.self_);
    }

    void print(std::ostream& os) const {
        if (self_) self_->print(os);
        else os << "<empty>";
    }
};

int main() {
    Printable a = 17;
    Printable b = std::string("hello");
    Printable c = b; // gleboka kopia wartosci
    b = 7.5;

    std::ostringstream stream;
    a.print(stream);
    c.print(stream);
    assert(stream.str() == "17hello");
}
\end{lstlisting}

\subsection{Analiza własności i kosztów}

Ten interfejs jest \emph{otwarty}: każdy nowy typ spełniający wymagania
może zostać opakowany bez modyfikowania klasy \texttt{Printable}.
Użytkownik nie musi dziedziczyć po klasie bazowej, co jest ważną zaletą
w porównaniu ze zwykłym polimorfizmem opartym na \texttt{Base}.
\texttt{Model<T>} jest instancjonowany w czasie kompilacji, natomiast
\texttt{print()} jest wywoływane dynamicznie przez ukryty interfejs.
Kopiowanie przez \texttt{clone()} zachowuje rzeczywisty typ wartości.

Wersja edukacyjna wykorzystuje osobną alokację dynamiczną przy
przechowywaniu niepustej wartości i kolejną alokację przy głębokim
kopiowaniu. Nie wykorzystuje \emph{small-buffer optimization}, które
mogłoby przechowywać niewielkie obiekty bezpośrednio wewnątrz opakowania.
W produkcyjnym type erasure trzeba dodatkowo przemyśleć \emph{wyrównanie},
przenoszenie, bezpieczne niszczenie, usuwanie typów niekopiowalnych,
kwalifikatory \texttt{const}, obsługę wyjątków i stabilność ABI.

\paragraph{Gwarancja silna przypisania kopiującego.}
Tworzenie obiektu \texttt{tmp(other)} może zakończyć się wyjątkiem.
Wówczas obiekt docelowy pozostaje niezmieniony. Jeśli kopia się uda,
\texttt{swap(tmp)} jest operacją \texttt{noexcept}; stara zawartość
zostanie zniszczona razem z \texttt{tmp}. Jest to klasyczny idiom
\emph{copy-and-swap}. Alternatywą jest niekopiowalne opakowanie,
które usuwa operacje kopiowania i przechowuje tylko obiekty przenoszalne.

\section{Type erasure a dziedziczenie i polimorfizm statyczny}
\label{sec:cpp14-comparison}

\subsection{Wspólna baza klas}

Gdy typy zostały zaprojektowane jako część wspólnej hierarchii,
\texttt{std::unique\_ptr<Base>} pozwala przechowywać obiekty o różnych
typach dynamicznych i wywoływać ich operacje wirtualne. Implementacje
muszą dziedziczyć po \texttt{Base}, a tworzenie kopii polimorficznej
wymaga własnego protokołu (np. \texttt{clone()}). Z drugiej strony
\texttt{unique\_ptr} jasno wyraża własność i nie kopiuje obiektów
wskutek zwykłego kopiowania uchwytu (ponieważ kopiowanie jest zabronione).

\begin{lstlisting}[caption={Polimorfizm przez bazę i \texttt{unique\_ptr}; kompletny program.}]
#include <cassert>
#include <memory>
#include <string>
#include <vector>

struct Shape {
    virtual ~Shape() = default;
    virtual int area() const = 0;
};
struct Square final : Shape {
    int side;
    explicit Square(int s) : side(s) {}
    int area() const override { return side * side; }
};
struct Rectangle final : Shape {
    int width, height;
    Rectangle(int w, int h) : width(w), height(h) {}
    int area() const override { return width * height; }
};

int main() {
    std::vector<std::unique_ptr<Shape>> shapes;
    shapes.push_back(std::make_unique<Square>(3));
    shapes.push_back(std::make_unique<Rectangle>(2, 5));
    assert(shapes[0]->area() + shapes[1]->area() == 19);
}
\end{lstlisting}

\subsection{Kiedy nie wymazywać typu?}

Jeżeli algorytm potrzebuje konkretnego typu wyłącznie podczas jednego
wywołania, parametr szablonowy z \texttt{concept} jest często prostszy
i pozwala unikać alokacji. Type erasure ma sens przede wszystkim na
granicy interfejsów, gdy wartości o różnych typach muszą być
\emph{przechowywane w jednej strukturze}, przekazywane przez stabilny
interfejs albo zastępowane w trakcie działania programu.

\begin{center}
\small
\begin{tabularx}{\textwidth}{@{}lXX@{}}
\toprule
\textbf{Rozwiązanie} & \textbf{Zalety} & \textbf{Ograniczenia} \\
\midrule
\texttt{variant<Ts...>} & Typy znane statycznie, sprawdzanie przypadków, lokalne przechowywanie wartości & Zamknięty katalog typów, złożoność odwiedzania i rozmiar największej alternatywy \\
\texttt{any} & Dowolna kopiowalna wartość, wspólny typ pojemnika & Trzeba odzyskać dokładny typ; brak wspólnego interfejsu operacji \\
\texttt{function} & Jedna sygnatura dla różnych callbacków & Kopiowalność (C++20), pośrednictwo, możliwe alokacje \\
Własne type erasure & Otwarty katalog modeli bez wspólnej bazy typów klienta & Złożona implementacja, zarządzanie pamięcią i wyjątkami \\
\texttt{unique\_ptr<Base>} & Jawna hierarchia, dziedziczone operacje, kontrola własności & Wymaga wspólnej bazy; operacje wirtualne \\
Szablon z \texttt{concept} & Typ znany kompilatorowi, możliwość optymalizacji & Wiele instancjacji; zwykle nie przechowa różnych typów w jednej kolekcji \\
\bottomrule
\end{tabularx}
\end{center}

\section{Wydajność: układ pamięci i koszty dynamicznej dyspozycji}
\label{sec:cpp14-performance}

Nie ma uniwersalnej odpowiedzi na pytanie, który mechanizm jest
najszybszy. Poniżej wymieniamy dominujące koszty; faktyczne zachowanie
zależy od rozmiaru typów, optymalizacji kompilatora, struktury programu
i rozkładu typów w danych.

\begin{itemize}
\item \texttt{optional<T>} dodaje znacznik obecności; dostęp do zawartości
      może wymagać sprawdzenia warunku.
\item \texttt{variant<Ts...>} przechowuje wystarczającą ilość miejsca
      na największy typ i indeks aktywnej alternatywy; odwiedzający
      może być implementowany przez instrukcję wyboru lub tablicę
      funkcji, a szczegóły zależą od kompilatora.
\item \texttt{any} przechowuje metadane typu oraz wartość, potencjalnie
      w osobno alokowanym buforze.
\item \texttt{function} może wymagać pośredniego wywołania i przechowania
      stanu lambdy; czasami da się uniknąć alokacji.
\item Własna implementacja z \texttt{unique\_ptr<Concept>} wykonuje
      alokację dynamiczną i pośrednie wywołanie wirtualne.
\end{itemize}

Przede wszystkim trzeba mierzyć \emph{rzeczywisty} przypadek użycia.
Jednorazowy koszt konstruowania callbacku może być nieistotny, gdy jego
wykonanie trwa długo, ale może dominować w pętli składającej się z
milionów krótkich operacji.

\subsection{Test bez mylenia czasu kompilacji z wykonaniem}

Rzetelne porównanie wymaga rozgrzania kodu, wykorzystania odpowiedniego
zestawu danych, uwzględnienia optymalizacji i sprawdzenia, czy kompilator
nie usunął mierzonej pracy. Warto rozdzielać koszt \emph{tworzenia}
opakowania od kosztu \emph{wywołania} i analizować nie tylko średni
czas, ale też rozkład wyników oraz liczbę alokacji. Benchmarku
porównującego dwa mechanizmy nie traktujemy jako dowodu ich przewagi
we wszystkich zastosowaniach.

\section{Granice bezpieczeństwa: wyjątki, czas życia i współbieżność}

\paragraph{Wyjątki.}
\texttt{optional::value()}, \texttt{std::get} dla wariantu, \texttt{any\_cast}
i wywołanie pustego \texttt{std::function} mogą zgłaszać odpowiednie
wyjątki. Konstrukcja zawartości i przypisanie także mogą rzucać,
ponieważ wykonują operacje typu przechowywanego. Należy rozróżniać
\emph{brak wartości} (normalny stan \texttt{optional}) od
\emph{niepowodzenia operacji}, któremu warto nadać osobną reprezentację.

\paragraph{Czas życia.}
\texttt{variant}, \texttt{optional} i \texttt{any} posiadają własne
przechowywane wartości. Można jednak umieścić w nich wskaźniki,
\texttt{reference\_wrapper} albo inne nieposiadające własności uchwyty;
wtedy odpowiedzialność za czas życia wskazywanych danych pozostaje po
stronie programisty. To samo dotyczy referencji przechwyconych przez
lambdę umieszczoną w \texttt{std::function}.

\paragraph{Wątki.}
Żaden z opisanych typów nie czyni dowolnych współbieżnych modyfikacji
tego samego obiektu bezpiecznymi. Przykładowo równoczesny zapis do
jednego \texttt{variant} i odczyt jego aktywnej zawartości bez
synchronizacji mogą powodować wyścig danych. Gdy callback jest
uruchamiany równolegle, trzeba osobno określić zasady synchronizacji
stanu przechwyconego przez lambdę.

\section{Nowości po C++20 i świadomy wybór standardu}

W całym rozdziale przykłady są zgodne z C++20. W nowszych standardach
pojawiają się udogodnienia, które warto znać, ale których \emph{nie wolno}
zakładać przy kompilacji z \texttt{-std=c++20}:
\begin{itemize}
\item \textbf{C++23:} \texttt{std::expected<T,E>} --- wynik albo
      jawna wartość błędu;
\item \textbf{C++23:} \texttt{std::move\_only\_function} --- opakowanie
      dla obiektów wywoływalnych, które nie muszą być kopiowalne;
\item \textbf{C++23:} operacje monadyczne dla \texttt{std::optional}:
      \texttt{and\_then}, \texttt{transform} i \texttt{or\_else};
\item \textbf{C++26:} \texttt{std::function\_ref} --- niewłaścicielski
      widok na obiekt wywoływalny, którego czas życia musi zostać
      zapewniony przez wywołującego.
\end{itemize}

Dostępność funkcji bibliotecznych w konkretnym kompilatorze i jego
bibliotece standardowej może być opóźniona w stosunku do daty
opublikowania standardu.

\section{Zasady projektowania i typowe błędy}
\label{sec:cpp14-guidelines}

\begin{enumerate}
\item Wybieraj \texttt{optional<T>}, gdy brak wyniku jest normalny i
      nie potrzebujesz szczegółowego typu błędu.
\item Wybieraj \texttt{variant}, gdy zestaw przypadków jest znany
      i zamknięty; obsługuj je przez \texttt{visit}, jeśli pozwala to
      uniknąć łańcucha ręcznych kontroli typu.
\item Stosuj \texttt{any} tylko tam, gdzie świadomie rezygnujesz
      ze statycznie określonego zestawu typów i masz jasny protokół
      odzyskiwania ich wartości.
\item Używaj \texttt{std::function} do \emph{przechowywanych},
      kopiowalnych callbacków; dla funkcji wykonywanej natychmiast
      rozważ parametr szablonowy.
\item Nie zakładaj, że \texttt{std::function} zawsze alokuje ani
      że \texttt{std::variant} zawsze jest szybszy. Zmierz przypadek
      zastosowania.
\item Nie pozostawiaj referencji lub wskaźników do obiektów, których
      czas życia zakończy się przed użyciem przechowywanego callbacku.
\item Unikaj \texttt{std::get} i \texttt{any\_cast} w sytuacjach,
      w których nie masz gwarancji typu; do sprawdzania wykorzystuj
      \texttt{get\_if} i wskaźnikowy \texttt{any\_cast}.
\item Nie używaj \texttt{std::any} do bezpośredniego opakowania
      typów przenoszalnych, ale niekopiowalnych, np. \texttt{unique\_ptr}.
\item Rozdziel decyzję o \emph{otwartym zbiorze implementacji} od
      decyzji o \emph{dynamicznym czasie życia i alokacji}.
\end{enumerate}

\section*{Podsumowanie rozdziału}
\addcontentsline{toc}{section}{Podsumowanie rozdziału 10}

\texttt{optional} modeluje brak lub obecność wartości, a
\texttt{variant} modeluje zamkniętą sumę przypadków. \texttt{any}
ukrywa typ kopiowalnej wartości, nie narzucając jej operacji.
\texttt{function} ukrywa typ obiektu wywoływalnego przy zachowaniu
ustalonej sygnatury. Własne type erasure pozwala opakować otwartą
rodzinę obcych typów bez wymagania wspólnej klasy bazowej, ale wymaga
starannego określenia semantyki własności, kopiowania i wyjątków.
Dobry projekt zaczyna się więc od określenia kontraktu: \emph{które
przypadki istnieją, kto posiada obiekty, jakie operacje udostępniamy
oraz czy zbiór typów ma pozostać otwarty}.

\section*{Dokumentacja i literatura}
\addcontentsline{toc}{section}{Dokumentacja i literatura rozdziału 10}
\begin{itemize}
\item \url{https://en.cppreference.com/w/cpp/utility/optional} --- \texttt{std::optional}.
\item \url{https://en.cppreference.com/w/cpp/utility/variant} --- \texttt{std::variant}.
\item \url{https://en.cppreference.com/w/cpp/utility/any} --- \texttt{std::any}.
\item \href{https://en.cppreference.com/w/cpp/utility/functional/function}{cppreference: \texttt{std::function}} --- dokumentacja opakowania wywołań.
\item Bjarne Stroustrup, \emph{A Tour of C++}, wyd. 3, Addison--Wesley, 2022.
\item Klaus Iglberger, \emph{C++ Software Design}, O'Reilly, 2022.
\end{itemize}

% ROZDZIAL 4 -- Variadic templates w C++ (C++20)
% Fragment dolaczany do glownego dokumentu po rozdziale 3.
% Wymagane w preambule: \usepackage{listings}, konfiguracja lstlisting
% (taka sama jak w rozdzialach 2 i 3).
% Uzycie: \input{rozdzial_4_variadic_templates_cpp.tex}

\chapter[Variadic templates w C++]{Szablony o zmiennej liczbie argumentów (\textit{variadic templates})}
\label{chap:variadic-templates-cpp}

\section{Idea i zastosowania}

Zwykły szablon funkcji przyjmuje z góry ustaloną liczbę parametrów
szablonu. \emph{Variadic template} (szablon o zmiennej liczbie parametrów)
może natomiast przyjmować \emph{dowolną liczbę} argumentów szablonu,
również zero. Mechanizmem, który to umożliwia, jest
\emph{pakiet parametrów} (ang. \emph{parameter pack}).

\begin{lstlisting}
template<class T>
void single(const T&) {}        // one type

template<class... Ts>
void many(const Ts&...) {}      // zero or more types

void demo() {
    many();                     // zero arguments
    many(10);                   // one argument
    many(10, 3.5, "abc");      // three arguments
}
\end{lstlisting}

Pakiety parametrów są ważną częścią nowoczesnego C++.
Wykorzystują je m.in. klasy \texttt{std::tuple} i \texttt{std::variant}
oraz fabryki \texttt{std::make\_unique} i \texttt{std::make\_shared}.
Nie trzeba dzięki nim pisać przeciążeń dla każdej
liczby argumentów.

\paragraph{Wymagane wersje języka.}
Pakiety parametrów i perfect forwarding pojawiły się w C++11.
\emph{Fold expressions} dodano w C++17, a koncepcje
(\texttt{concepts}) i klauzulę \texttt{requires} w C++20.
Dalej zakładamy C++20.

\section{Trzy kropki: deklaracja pakietu a rozwinięcie}

\begin{center}\small
\begin{tabularx}{\textwidth}{@{}p{0.35\textwidth}X@{}}
\toprule
Zapis & Działanie \\
\midrule
\code{typename... Ts} & Deklaruje pakiet parametrów typowych szablonu. \\
\code{Ts...} & Rozwija pakiet typów w miejscu, które dopuszcza taki wzorzec. \\
\code{args...} & Rozwija pakiet argumentów, tworząc listę wyrażeń. \\
\code{sizeof...(Ts)} & Podaje liczbę elementów pakietu w czasie kompilacji. \\
\code{std::forward<Ts>(args)...} & Przekazuje każdy argument z zachowaniem jego kategorii wartości. \\
\bottomrule
\end{tabularx}
\end{center}
\label{sec:variadic-trzy-kropki}

W różnych miejscach token \texttt{...} oznacza różne operacje.
Rozróżnienie to stanowi podstawę poprawnego używania szablonów
wariadycznych.

\begin{lstlisting}
template<class... Ts>             // declare type pack Ts
void f(Ts... xs) {                // declare function parameter pack xs
    // ...
}
\end{lstlisting}

Deklaracja \texttt{class... Ts} wprowadza pakiet \emph{typów}.
Zapis \texttt{Ts... xs} rozwija pakiet typów i wprowadza odpowiadający
mu pakiet parametrów funkcji. Wywołanie
\texttt{f(1, 2.0, 'x')} tworzy w uproszczeniu funkcję o parametrach
\texttt{(int, double, char)}.

\subsection{Rozwinięcie pakietu (\textit{pack expansion})}

Rozwinięcie składa się ze \emph{wzorca} (\emph{pattern}) oraz następującego
po nim \texttt{...}. Dla każdego elementu pakietu kompilator podstawia
kolejny element do całego wzorca:

\begin{lstlisting}
template<class... Ts>
void demo(Ts... xs) {
    g(xs...);                // g(x1, x2, ..., xn)
    h((xs * 2)...);          // h(x1*2, x2*2, ..., xn*2)
}
\end{lstlisting}

W drugim wywołaniu wzorcem jest \texttt{(xs * 2)}, a \emph{nie}
pojedyncza zmienna \texttt{xs}. Ważne: \texttt{xs...} nie jest
kontenerem ani tablicą; nie można po nim bezpośrednio iterować zwykłą
pętlą \texttt{for}. To składnia kompilatora służąca do generowania
listy wyrażeń.

\paragraph{Pakiet typów a pakiet wartości.}
W przykładzie \texttt{Ts} reprezentuje typy, natomiast
\texttt{xs} to konkretne parametry funkcji. Jeden typ może wystąpić
w pakiecie wielokrotnie, a każdy element pakietu wartości zachowuje
swój właściwy typ.

\subsection{Liczba elementów: \texttt{sizeof...}}

Operator \texttt{sizeof...} daje liczbę elementów pakietu jako
stałą czasu kompilacji. Nie należy go mylić z operatorem
\texttt{sizeof} mierzącym rozmiar typu lub obiektu.

\begin{lstlisting}
#include <cstddef>

template<class... Ts>
constexpr std::size_t count(Ts&&... xs) {
    static_assert(sizeof...(Ts) == sizeof...(xs));
    return sizeof...(xs);
}

static_assert(count() == 0);
static_assert(count(2, 3.5, 'a') == 3);
\end{lstlisting}

Nazwy \texttt{xs} nie używamy w ciele funkcji do odczytywania
wartości; operacja liczenia jest czysto składniowa.

\section{Pierwsze praktyczne funkcje}

\subsection{Wypisywanie dowolnej liczby argumentów}

Najprostszy sposób wykorzystania pakietu to jego rozwinięcie
w wywołaniu innej funkcji:

\begin{lstlisting}
#include <iostream>
#include <string>

void output(int x, double y, const std::string& z) {
    std::cout << x << ' ' << y << ' ' << z << '\n';
}

template<class... Args>
void forward_to_output(Args... args) {
    output(args...);
}

int main() {
    forward_to_output(3, 2.5, std::string{"hello"});
}
\end{lstlisting}

Funkcja jest poprawna tylko dla takich wywołań, z których można
utworzyć poprawne wywołanie \texttt{output}. Sama wariadyczność
nie oznacza, że każda liczba i każdy rodzaj argumentów będą
akceptowane przez funkcję docelową.

\subsection{Przekazywanie przez wartość i przez referencję}

\begin{lstlisting}
template<class... Ts>
void by_value(Ts... xs);            // may copy/move

template<class... Ts>
void by_const_ref(const Ts&... xs); // read-only; no copies

template<class... Ts>
void by_ref(Ts&... xs);             // needs lvalues
\end{lstlisting}

Wybór postaci parametrów ma te same skutki co w zwykłych funkcjach.
W szczególności \texttt{Ts... xs} tworzy osobne parametry przez wartość,
a \texttt{const Ts\&... xs} zwykle pozwala uniknąć niepotrzebnego
kopiowania. Przypadek \texttt{Ts\&\&... xs} omówimy dokładniej
w sekcji o \emph{perfect forwarding}.

\section{Rekurencyjne przetwarzanie pakietu: metoda C++11}
\label{sec:variadic-rekurencja}

Przed C++17 często rozbijano pakiet na \emph{pierwszy argument}
i \emph{pozostałe argumenty}, a następnie wywoływano funkcję
rekurencyjnie. Potrzebny był osobny przypadek kończący rekurencję.

\begin{lstlisting}
#include <iostream>

void print_recursive() {
    std::cout << '\n';        // base case
}

template<class First, class... Rest>
void print_recursive(const First& first,
                     const Rest&... rest) {
    std::cout << first << ' ';
    print_recursive(rest...); // one argument fewer
}

int main() {
    print_recursive(1, "two", 3.0);
}
\end{lstlisting}

Przy trzech argumentach instancjonowane są kolejno wersje
z pakietami o długościach $3$, $2$, $1$ oraz przypadek końcowy.
Jest to \emph{rekurencja szablonowa}; nie należy automatycznie
utożsamiać jej z dynamiczną rekurencją wykonywaną w programie.
Kompilator może wygenerować zwykłą sekwencję operacji.

\paragraph{Kiedy stosować?}
Rekurencyjne rozkładanie pakietu nadal bywa przydatne, jeśli kolejne
kroki mają odrębne warunki lub należy budować typy rekurencyjne.
Jeżeli wykonujemy tę samą operację dla wszystkich argumentów,
w C++17 i nowszych zazwyczaj lepszym wyborem jest \emph{fold expression}.

\section{Fold expressions (C++17)}

\begin{center}\small
\begin{tabularx}{\textwidth}{@{}p{0.35\textwidth}X@{}}
\toprule
Forma złożenia & Sposób rozwinięcia \\
\midrule
\code{(... + args)} & Złożenie lewostronne: kolejne dodawania łączone od lewej. \\
\code{(args + ...)} & Złożenie prawostronne: kolejne dodawania łączone od prawej. \\
\code{(init + ... + args)} & Złożenie z wartością początkową; działa także dla pustego pakietu. \\
\code{(f(args), ...)} & Wywołuje \code{f} dla każdego argumentu od lewej do prawej dzięki operatorowi przecinka. \\
\texttt{(args \&\& ...)} & Łączy warunki z krótkim spięciem; dla pustego pakietu daje \code{true}. \\
\bottomrule
\end{tabularx}
\end{center}
\label{sec:variadic-fold}

\emph{Fold expression} redukuje pakiet wyrażeń z użyciem jednego
operatora binarnego, np. \texttt{+}, \texttt{\&\&}, \texttt{||}
lub operatora przecinka. Nie wymaga jawnego przypadku
rekurencyjnego.

\subsection{Cztery postacie składni}

Niech \texttt{xs} oznacza pakiet, $\circ$ ustalony operator,
a $z$ wartość początkową. Istnieją cztery postacie składni:

\begin{center}
\small
\begin{tabular}{|p{0.39\textwidth}|p{0.47\textwidth}|}
\hline
\textbf{Postać C++} & \textbf{Sposób łączenia dla trzech argumentów} \\
\hline
\texttt{(xs op ...)} & $x_1\circ(x_2\circ x_3)$ (prawostronne, bez inicjalizatora) \\
\hline
\texttt{(... op xs)} & $(x_1\circ x_2)\circ x_3$ (lewostronne, bez inicjalizatora) \\
\hline
\texttt{(xs op ... op z)} & $x_1\circ(x_2\circ(x_3\circ z))$ (prawostronne, z inicjalizatorem) \\
\hline
\texttt{(z op ... op xs)} & $((z\circ x_1)\circ x_2)\circ x_3$ (lewostronne, z inicjalizatorem) \\
\hline
\end{tabular}
\end{center}

W tabeli \texttt{op} jest jedynie oznaczeniem miejsca na rzeczywisty
operator, np. \texttt{+}; nie jest słowem kluczowym języka.
\emph{Nawiasy całego fold expression są obowiązkowe}.
Dwa operatory występujące w fold z inicjalizatorem muszą być takie same.

\subsection{Suma, iloczyn i warunki logiczne}

\begin{lstlisting}
#include <concepts>

template<class... Ts>
    requires (std::integral<Ts> && ...)
constexpr auto sum(Ts... xs) {
    return (0 + ... + xs);       // left fold with identity 0
}

template<class... Ts>
constexpr bool all_true(Ts... xs) {
    return (true && ... && static_cast<bool>(xs));
}

template<class... Ts>
constexpr bool any_true(Ts... xs) {
    return (false || ... || static_cast<bool>(xs));
}

static_assert(sum() == 0);
static_assert(sum(2, 3, 4) == 9);
static_assert(all_true());
static_assert(!any_true());
static_assert(all_true(1, 2, 3));
\end{lstlisting}

Wyrażenie \texttt{(0 + ... + xs)} może być używane także dla
pustego pakietu. Początkowe \texttt{0} wpływa jednak na typ wyniku
i ewentualne promocje arytmetyczne; nie jest to uniwersalna
implementacja sumy dowolnych typów liczbowych.

\subsection{Lewostronne i prawostronne składanie nie zawsze są równe}

Dla operatora odejmowania kolejność nawiasowania zmienia wynik:

\begin{lstlisting}
template<class... Ts>
constexpr auto left_sub(Ts... xs) {
    return (100 - ... - xs);
}

template<class... Ts>
constexpr auto right_sub(Ts... xs) {
    return (xs - ... - 100);
}

static_assert(left_sub(10, 3) == 87);  // (100 - 10) - 3
static_assert(right_sub(10, 3) == 107); // 10 - (3 - 100)
\end{lstlisting}

Fold określa \emph{strukturę nawiasów}, ale nie należy mylić tej
własności z ogólną gwarancją kolejności wykonywania wszystkich
podwyrażeń dla dowolnego operatora. Reguły ewaluacji wynikają
również z semantyki danego operatora.

\subsection{Pusty pakiet i element neutralny}

Dla \emph{jednoargumentowego} fold z pustym pakietem C++ dopuszcza
zaledwie trzy operatory: \texttt{\&\&} zwraca \texttt{true},
\texttt{||} zwraca \texttt{false}, a \texttt{,} daje
wyrażenie typu \texttt{void}. Pusty pakiet w
\texttt{(... + xs)} nie jest poprawny. Stosując wersję
z wartością początkową, możemy jawnie obsłużyć pusty pakiet:

\begin{lstlisting}
template<class... Ts>
constexpr auto product(Ts... xs) {
    return (1 * ... * xs);
}

static_assert(product() == 1);
static_assert(product(2, 3, 4) == 24);
\end{lstlisting}

\section{Operator przecinka i kolejność wykonywania}
\label{sec:variadic-przecinek}

Operator przecinka jest wygodny, gdy dla każdego elementu pakietu
chcemy wykonać instrukcję o efektach ubocznych. W odróżnieniu
od kolejności ewaluacji zwykłych argumentów funkcji, jego
wbudowana postać gwarantuje wykonanie lewej strony przed prawą.

\begin{lstlisting}
#include <iostream>

template<class... Ts>
void print_all(const Ts&... xs) {
    ((std::cout << xs << ' '), ...);
    std::cout << '\n';
}

int main() {
    print_all(10, 2.5, "hello");
    print_all();                // only a newline
}
\end{lstlisting}

Zewnętrzne nawiasy tworzą fold, a wewnętrzne oddzielają
wyrażenie wypisujące od operatora \texttt{,}.

\paragraph{Kolejność argumentów zwykłej funkcji.}

\begin{lstlisting}
template<class... Ts>
void sink(Ts...);

// sink(change_a(), change_b(), change_c());
\end{lstlisting}

Samo rozwinięcie \texttt{...} w liście argumentów funkcji
\emph{nie} gwarantuje, że \texttt{change\_a()} zostanie obliczone
przed \texttt{change\_b()}. W C++17 i nowszych wywołania
te nie przeplatają swoich efektów, ale ich wzajemna kolejność
nie jest określona. Gdy kolejność jest istotna, można użyć
operatora przecinka lub innej konstrukcji dającej odpowiednie
gwarancje sekwencjonowania.

\section{Pakiety typów i metaprogramowanie}
\label{sec:variadic-type-list}

Wariadyczność może dotyczyć nie tylko funkcji, lecz także
\emph{szablonów klas}. W takim przypadku pakiet przechowuje
listę typów, a instancjonowanie klasy odbywa się podczas kompilacji.

\subsection{Własna lista typów}

\begin{lstlisting}
#include <cstddef>
#include <type_traits>

template<class... Ts>
struct type_list {};

template<class List>
struct list_size;

template<class... Ts>
struct list_size<type_list<Ts...>>
    : std::integral_constant<std::size_t, sizeof...(Ts)> {};

using Types = type_list<int, double, char>;
static_assert(list_size<Types>::value == 3);
static_assert(list_size<type_list<>>::value == 0);
\end{lstlisting}

Jest to przykład \emph{częściowej specjalizacji} szablonu klasy.
Wzorzec \texttt{type\_list<Ts...>} dopasowuje dowolną listę
typów, a \texttt{sizeof...(Ts)} oblicza jej długość.

\subsection{Wydobywanie pierwszego typu}

\begin{lstlisting}
template<class List>
struct first_type;   // intentionally undefined for empty list

template<class T, class... Rest>
struct first_type<type_list<T, Rest...>> {
    using type = T;
};

using X = first_type<type_list<long, float>>::type;
static_assert(std::is_same_v<X, long>);
\end{lstlisting}

Rozkład na \texttt{T} oraz \texttt{Rest...} jest odpowiednikiem
rekurencyjnego rozdzielenia argumentów funkcji na pierwszy
i pozostałe. Dla pustej listy \texttt{type\_list<>} nazwa
\texttt{first\_type<...>::type} celowo nie istnieje.

\subsection{Czy wszystkie typy spełniają własność?}

\begin{lstlisting}
#include <type_traits>

template<class... Ts>
inline constexpr bool all_integral_v =
    (std::is_integral_v<Ts> && ...);

template<class... Ts>
inline constexpr std::size_t count_integral_v =
    (std::size_t{0} + ... +
        (std::is_integral_v<Ts> ? 1u : 0u));

static_assert(all_integral_v<int, long, char>);
static_assert(!all_integral_v<int, double>);
static_assert(count_integral_v<int, double, char> == 2);
\end{lstlisting}

Jest to częsty wzorzec: pakiet typów łączymy z cechami
typów (\emph{type traits}) i operatorami logicznymi fold.

\section{Pakiety wartości i parametry nietypowe}

Pakiet może także zawierać wartości znane w czasie kompilacji,
a nie typy. W szablonie klasy służą do tego
\emph{nietypowe parametry szablonu} (\emph{non-type template parameters}).

\begin{lstlisting}
#include <array>
#include <cstddef>

template<int... Ns>
struct int_pack {
    static constexpr std::size_t size = sizeof...(Ns);
    static constexpr int sum = (0 + ... + Ns);
};

static_assert(int_pack<>::sum == 0);
static_assert(int_pack<2, 4, 6>::sum == 12);
\end{lstlisting}

Od C++17 można zadeklarować pakiet \texttt{auto...}:

\begin{lstlisting}
template<auto... Values>
struct value_pack {
    static constexpr std::size_t size = sizeof...(Values);
};

static_assert(value_pack<42, 'x', true>::size == 3);
\end{lstlisting}

Dla \texttt{auto...} dopuszczalność konkretnej wartości jako
argumentu szablonu nadal podlega regułom parametrów nietypowych;
nie można przekazać do niego dowolnej wartości z czasu wykonania.

\section{Przekazywanie doskonałe (\textit{perfect forwarding})}
\label{sec:variadic-forwarding}

Jednym z najważniejszych zastosowań pakietów jest budowanie
\emph{funkcji pośredniczących}, które przekazują argumenty do
innej funkcji lub konstruktora, możliwie wiernie zachowując
ich kategorię wartości (\emph{value category}).

\subsection{Różnica między lvalue a rvalue}

\begin{lstlisting}
#include <iostream>

void use(int&)  { std::cout << "lvalue\n"; }
void use(int&&) { std::cout << "rvalue\n"; }

int main() {
    int x = 10;
    use(x);       // lvalue
    use(20);      // rvalue
}
\end{lstlisting}

Parametr funkcji, choćby zadeklarowany jako \texttt{T\&\&}, jest
\emph{nazwanym obiektem}; użycie jego nazwy w wyrażeniu ma
kategorię \emph{lvalue}. Dlatego samo \texttt{use(arg)}
nie zachowa kategorii pierwotnego argumentu.

\subsection{Referencje przekazujące i \texttt{std::forward}}

\begin{lstlisting}
#include <utility>

template<class T>
void wrapper(T&& arg) {
    use(std::forward<T>(arg));
}

void forwarding_example() {
    int y = 0;
    wrapper(y);  // T = int&, forwarded as lvalue
    wrapper(30); // T = int, forwarded as rvalue
}
\end{lstlisting}

\texttt{T\&\&} jest \emph{referencją przekazującą}
(ang. \emph{forwarding reference}), gdy \texttt{T} jest
odpowiednio dedukowanym, niekwalifikowanym parametrem typu.
Działa tu \emph{składanie referencji}:

\[
T\&\;\& \longrightarrow T\&,
\qquad
T\&\;\&\& \longrightarrow T\&,
\qquad
T\&\&\;\&\& \longrightarrow T\&\&.
\]

W zapisie \texttt{std::forward<T>(arg)} jawny argument
\texttt{T} jest konieczny, by odtworzyć pierwotną kategorię.
\texttt{std::forward} nie wykonuje bezwarunkowego przeniesienia:
przepuszcza lvalue jako lvalue, a rvalue jako rvalue.

\subsection{Pakiet referencji przekazujących}

\begin{lstlisting}
#include <functional>
#include <utility>

template<class F, class... Args>
decltype(auto) call(F&& f, Args&&... args) {
    return std::invoke(
        std::forward<F>(f),
        std::forward<Args>(args)...
    );
}
\end{lstlisting}

Tutaj rozwijany jest \emph{cały wzorzec}
\texttt{std::forward<Args>(args)}, a nie tylko nazwa
\texttt{args}. Każdy element używa właściwego typu z pakietu
\texttt{Args}.
\texttt{std::invoke} potrafi wywoływać zwykłe funkcje,
obiekty funkcyjne i odpowiednio obsłużone wskaźniki do składowych.
\texttt{decltype(auto)} zachowuje typ zwracany przez docelowe
wywołanie, w tym referencję, jeżeli docelowa funkcja ją zwraca.

\begin{lstlisting}
#include <string>

std::string join(std::string a, std::string b) {
    return a + b;
}

auto result = call(join, std::string{"ab"}, "cd");
// result == "abcd"
\end{lstlisting}

\subsection{Czego nie robić: \texttt{std::move} i nazwy parametrów}

\begin{lstlisting}
template<class T>
void wrong_1(T&& arg) {
    use(arg);              // arg is always an lvalue here
}

template<class T>
void wrong_2(T&& arg) {
    use(std::move(arg));   // forces an rvalue even for lvalues
}
\end{lstlisting}

\texttt{std::move} to rzutowanie umożliwiające potraktowanie
obiektu jak rvalue, a nie polecenie ``zachowaj kategorię''.
Stosuj je, gdy rzeczywiście przekazujesz obiekt do operacji,
która może pobrać z niego zasób. W uniwersalnym wrapperze
używaj \texttt{std::forward}.

\paragraph{Ograniczenia żywotności.}
Referencje przekazujące \emph{nie przedłużają automatycznie}
żywotności obiektu do czasu dowolnego późniejszego wykonania.
Przechowanie referencji do argumentu tymczasowego w strukturze,
która przeżyje wywołanie, może prowadzić do wskaźników lub
referencji wiszących. Ten sam problem dotyczy między innymi
\texttt{std::forward\_as\_tuple}.

\section{Tworzenie obiektów z dowolnych argumentów}
\label{sec:variadic-emplace}

\subsection{Własna fabryka obiektów}

\begin{lstlisting}
#include <memory>
#include <utility>

template<class T, class... Args>
std::unique_ptr<T> create(Args&&... args) {
    return std::make_unique<T>(
        std::forward<Args>(args)...
    );
}

struct Point {
    int x, y;
    Point(int a, int b) : x(a), y(b) {}
};

auto point = create<Point>(4, 7);
\end{lstlisting}

\texttt{std::make\_unique} już dostarcza taką funkcjonalność,
więc \texttt{create} ma tu charakter dydaktyczny. Pokazuje
jednak wzorzec używany w fabrykach oraz funkcjach
\texttt{emplace} i \texttt{make}.

\subsection{Tworzenie kontenera bez wymogu kopiowalności}

\begin{lstlisting}
#include <concepts>
#include <string>
#include <utility>
#include <vector>

template<class T, class... Args>
    requires (std::constructible_from<T, Args&&> && ...)
std::vector<T> make_vector(Args&&... args) {
    std::vector<T> result;
    result.reserve(sizeof...(Args));
    (result.emplace_back(std::forward<Args>(args)), ...);
    return result;
}

auto words = make_vector<std::string>("A", "BB", "CCC");
\end{lstlisting}

Operator przecinka daje tu jednoznaczną kolejność wywołań
\texttt{emplace\_back}. Funkcja przyjmuje także pusty pakiet
argumentów. W przeciwieństwie do budowania wektora z listy
inicjalizacyjnej pozwala łatwo przekazywać do kolejnych
\texttt{emplace\_back} argumenty przenoszalne, ale niekopiowalne.
Dla elementów konstruowanych z więcej niż jednego parametru
wymagałaby innego interfejsu grupowania argumentów.

\section{\texttt{std::tuple}, \texttt{std::apply} i \texttt{std::index\_sequence}}
\label{sec:variadic-tuple}

\subsection{Heterogeniczna krotka \texttt{std::tuple}}

\texttt{std::tuple<Ts...>} to kontener o strukturze opisanej
listą typów znaną w czasie kompilacji. Każda pozycja może
mieć inny typ:

\begin{lstlisting}
#include <tuple>
#include <string>

std::tuple<int, std::string, double> t{7, "Ada", 3.5};

int n = std::get<0>(t);
const std::string& name = std::get<1>(t);
\end{lstlisting}

W przeciwieństwie do \texttt{std::vector<T>}, krotka nie
przechowuje dowolnej liczby elementów jednego typu, lecz
\emph{ustaloną na etapie kompilacji} listę typów.

\subsection{Rozwinięcie krotki do argumentów funkcji}

Funkcja \texttt{std::apply} (C++17) wywołuje funkcję,
przekazując elementy obiektu \emph{tuple-like} jako osobne
argumenty:

\begin{lstlisting}
#include <iostream>
#include <tuple>

void apply_example() {
    std::apply([](const auto&... xs) {
        ((std::cout << xs << ' '), ...);
        std::cout << '\n';
    }, t);
}
\end{lstlisting}

\texttt{std::apply} oszczędza ręcznego konstruowania
sekwencji indeksów, ale warto wiedzieć, jak taki mechanizm
można zaimplementować.

\subsection{Pakiety indeksów: \texttt{std::index\_sequence}}

\begin{lstlisting}
#include <cstddef>
#include <iostream>
#include <tuple>
#include <utility>

template<class Tuple, std::size_t... I>
void print_tuple_impl(const Tuple& t,
                      std::index_sequence<I...>) {
    ((std::cout << std::get<I>(t) << ' '), ...);
    std::cout << '\n';
}

template<class... Ts>
void print_tuple(const std::tuple<Ts...>& t) {
    print_tuple_impl(t, std::index_sequence_for<Ts...>{});
}

void tuple_example() {
    print_tuple(std::make_tuple(1, "abc", 3.0));
}
\end{lstlisting}

\texttt{std::index\_sequence<0,1,2>} jest specjalizacją
szablonu przyjmującego pakiet indeksów całkowitych.
\texttt{std::index\_sequence\_for<Ts...>} generuje indeksy
$0,1,\ldots,n-1$ dla pakietu długości $n$.
Wzorzec \texttt{std::get<I>(t)...} rozwija wyrażenie
osobno dla każdego indeksu.

\subsection{Równoległe rozwijanie pakietów}

W jednym wzorcu można użyć kilku pakietów. Muszą mieć
\emph{tę samą długość}, jeśli są rozwijane równolegle:

\begin{lstlisting}
#include <tuple>
#include <type_traits>
#include <utility>

template<class A, class B>
struct zipped_types;

template<class... Ts, class... Us>
struct zipped_types<std::tuple<Ts...>, std::tuple<Us...>> {
    static_assert(sizeof...(Ts) == sizeof...(Us));
    using type = std::tuple<std::pair<Ts, Us>...>;
};

using Z = zipped_types<
    std::tuple<int, char>,
    std::tuple<double, bool>
>::type;

static_assert(std::is_same_v<
    Z, std::tuple<std::pair<int, double>,
                  std::pair<char, bool>>>);
\end{lstlisting}

Elementy \texttt{Ts} i \texttt{Us} są łączone pozycjami:
$\bigl(T_1,U_1\bigr),\ldots,\bigl(T_n,U_n\bigr)$.
Nie jest to iloczyn kartezjański pakietów.

\section{Dziedziczenie wariadyczne i \texttt{std::visit}}
\label{sec:variadic-overloaded}

Pakiet \texttt{Bases...} można rozwinąć na liście baz.
Z kolei zapis \texttt{using Bases::operator()...}
wprowadza ich operatory wywołania do jednej klasy.
W ten sposób otrzymujemy zestaw przeciążeń.

\begin{lstlisting}
#include <iostream>
#include <string>
#include <variant>

template<class... Fs>
struct overloaded : Fs... {
    using Fs::operator()...;
};

template<class... Fs>
overloaded(Fs...) -> overloaded<Fs...>; // C++17 deduction guide

int main() {
    std::variant<int, std::string> value = "hello";

    std::visit(overloaded{
        [](int n) { std::cout << "int: " << n << '\n'; },
        [](const std::string& s) {
            std::cout << "string: " << s << '\n';
        }
    }, value);
}
\end{lstlisting}

Jest to standardowy idiom C++17: każda lambda definiuje
osobny \texttt{operator()}, a dziedziczenie wariadyczne
łączy je w jeden obiekt wywoływalny. Gdy stosujemy ten wzorzec
w wielu miejscach, warto umieścić \texttt{overloaded} w
niewielkim nagłówku narzędziowym.

\section{Ograniczanie pakietów za pomocą \texttt{concepts} (C++20)}
\label{sec:variadic-concepts}

Nie każda operacja ma sens dla każdej możliwej kombinacji
typów. W C++20 można czytelnie określać wymagania dla
\emph{wszystkich} elementów pakietu.

\begin{lstlisting}
#include <concepts>
#include <type_traits>

template<class... Ts>
    requires (std::integral<Ts> && ...)
constexpr bool all_nonnegative(Ts... xs) {
    return ((xs >= 0) && ...);
}

static_assert(all_nonnegative(0, 2, 7));
static_assert(!all_nonnegative(0, -1, 7));

// all_nonnegative(3.14); // rejected by constraints
\end{lstlisting}

Możemy również ograniczać pakiet w odniesieniu do
konkretnego typu docelowego:

\begin{lstlisting}
#include <concepts>

template<class T, class... Args>
concept ConstructibleWith =
    std::constructible_from<T, Args...>;

static_assert(ConstructibleWith<std::string, const char*>);
static_assert(!ConstructibleWith<int, std::string>);
\end{lstlisting}

Pierwszy przykład stosuje fold nad \emph{warunkami dla pojedynczych
typów}; drugi sprawdza \emph{jedną własność całej sekwencji
typów}, tj. czy odpowiedni konstruktor przyjmuje argumenty
\texttt{Args...} w tej kolejności. Tych dwóch mechanizmów
nie należy utożsamiać.

\subsection{\texttt{if constexpr} i struktura pakietu}

\begin{lstlisting}
#include <iostream>

template<class... Ts>
void report(const Ts&... xs) {
    if constexpr (sizeof...(Ts) == 0) {
        std::cout << "empty\n";
    } else {
        ((std::cout << xs << ' '), ...);
        std::cout << '\n';
    }
}
\end{lstlisting}

\texttt{if constexpr} pozwala w zależności od długości pakietu
wybierać gałąź w czasie kompilacji. W dobrze sformułowanym
szablonie odrzucona gałąź nie jest instancjonowana dla danej
specjalizacji. To użyteczne przy przypadkach brzegowych,
które nie mają sensu dla pustego pakietu.

\section{Ograniczenia, pułapki i diagnostyka}
\label{sec:variadic-pulapki}

\subsection{Pakiet typów nie jest pakietem argumentów funkcji}

\begin{lstlisting}
template<class... Ts>
void f(Ts... args) {
    // sizeof...(Ts): number of types
    // sizeof...(args): number of function parameters
    // args[0];  // ERROR: args is not an array
}
\end{lstlisting}

Do wybierania argumentu na ustalonej pozycji można użyć
\texttt{std::tuple} i \texttt{std::get} albo wywoływać funkcję
przez \texttt{std::apply}. Sam pakiet \texttt{args...}
nie jest indeksowanym kontenerem.

\subsection{Brak zgodności typów po rozwinięciu}

\begin{lstlisting}
template<class... Ts>
auto add(Ts... xs) {
    return (0 + ... + xs);
}

// add(1, std::string{"x"}); // ERROR: unsuitable operator+
\end{lstlisting}

Pakiet nie omija systemu typów. Każde wyrażenie powstałe po
jego rozwinięciu musi być poprawne, podobnie jak zwykły kod
C++. Jeżeli szablon ma określone wymagania, najlepiej
zapisać je za pomocą \texttt{requires}.

\subsection{Niewłaściwe przekazywanie kategorii wartości}

\begin{lstlisting}
template<class F, class... Args>
void bad_call(F&& f, Args&&... args) {
    std::forward<F>(f)(args...); // loses rvalue categories
}

// Correct:
// std::forward<F>(f)(std::forward<Args>(args)...);
\end{lstlisting}

Dla typów niekopiowalnych taki błąd często ujawnia się
natychmiast podczas kompilacji. Dla typów kopiowalnych
może spowodować jedynie dodatkowe kopie albo zmianę
wybranego przeciążenia.

\subsection{Wnioskowanie typów i \texttt{\{...\}}}

\begin{lstlisting}
#include <initializer_list>

template<class... Args>
void foo(Args&&...) {}

void braced_example() {
    // foo({1, 2, 3}); // ERROR: cannot deduce Args
    foo(std::initializer_list<int>{1, 2, 3}); // OK
}
\end{lstlisting}

Sama lista ujęta w nawiasy klamrowe \texttt{\{...\}} nie ma
typu, który dowolny pakiet \texttt{Args\&\&...} potrafiłby
wydedukować. Potrzebny jest określony typ parametru,
np. \texttt{std::initializer\_list}\allowbreak\texttt{<int>}, lub jawnie
utworzony obiekt.

\subsection{Ograniczenia położenia pakietu}

W \emph{głównej definicji szablonu klasy} pakiet parametrów
szablonu musi być ostatni. Dla szablonów funkcji istnieją
bardziej elastyczne reguły, ale kolejne parametry muszą
być dedukowalne lub mieć argumenty domyślne, stosownie do
kontekstu:

\begin{lstlisting}
// ERROR for a primary class template:
// template<class... Ts, class U> struct invalid;

template<class U, class... Ts>
struct valid {};
\end{lstlisting}

\subsection{Skomplikowany kod zamiast prostego rozwiązania}

Szablon wariadyczny to narzędzie, a nie obowiązkowy styl.
Jeżeli argumenty mają \emph{jednolity typ i ich liczba
zmienia się w czasie wykonania}, lepszy bywa
\texttt{std::span<T>}, \texttt{std::vector<T>} lub para
iteratorów. Pakiet ma liczbę elementów i strukturę typów
ustalaną na etapie kompilacji; nie zastępuje kontenera.

\section{Variadic templates a funkcje wariadyczne w stylu C}

Historyczne funkcje z deklaracją \texttt{void f(int, ...)}
korzystają z mechanizmu argumentów wariadycznych języka C,
a nie z pakietu parametrów szablonu C++:

\begin{lstlisting}
#include <cstdio>

int main() {
    std::printf("%d %f\n", 7, 2.5); // C-style variadic call
}
\end{lstlisting}

W formacie C zgodność argumentów ze specyfikatorami
formatowania jest obowiązkiem programisty (niektóre
kompilatory potrafią wykrywać część pomyłek). W szablonach
wariadycznych typy argumentów pozostają dostępne
kompilatorowi, dzięki czemu można wymuszać ograniczenia
i stosować przeciążanie oraz \texttt{std::forward}.
Nie należy jednak wyciągać z tego wniosku, że dowolny
szablon automatycznie bezbłędnie formatuje argumenty.

\section{Pełny przykład: prosty rejestrator zdarzeń}
\label{sec:variadic-projekt}

Zastosujmy omówione mechanizmy do funkcji tworzącej tekstowy
komunikat z dowolnej liczby argumentów możliwych do
wypisania do strumienia. Zamiast ręcznie doklejać spacje
w wielu przeciążeniach, użyjemy fold nad operatorem przecinka.

\begin{lstlisting}
#include <iostream>
#include <sstream>
#include <string>

// Each argument must support insertion into std::ostream.
template<class... Args>
std::string make_message(const Args&... args) {
    std::ostringstream out;
    ((out << args), ...);
    return out.str();
}

template<class... Args>
void log_line(const Args&... args) {
    std::cout << "[LOG] " << make_message(args...) << '\n';
}

int main() {
    log_line("count=", 3, ", ratio=", 0.75);
    log_line("finished");
    log_line();
}
\end{lstlisting}

\paragraph{Analiza.}
Dla argumentów \texttt{("count=",3,", ratio=",0.75)}
pakiet zawiera cztery elementy. Wyrażenie
\texttt{((out << args), ...)} wstawia je kolejno do strumienia,
a \texttt{make\_message} zwraca powstały napis. Dla
pustego pakietu fold z operatorem przecinka jest poprawny,
a wynikowy napis jest pusty. W wersji produkcyjnej należałoby
rozważyć wymagania dotyczące formatowania, zachowanie w
obliczu wyjątków, współbieżność oraz ewentualne skorzystanie
z narzędzia formatującego wspieranego przez używany kompilator.

\section{Zestawienie: kiedy używać poszczególnych mechanizmów?}

\begin{center}
\small
\begin{tabular}{|p{0.40\textwidth}|p{0.46\textwidth}|}
\hline
\textbf{Cel} & \textbf{Typowe rozwiązanie} \\
\hline
Dowolna liczba argumentów o różnych typach & \texttt{template<class... Ts>} i \texttt{Ts...} \\
\hline
Wykonanie tej samej operacji dla każdego argumentu & Fold z operatorem \texttt{,} (C++17) \\
\hline
Suma, iloczyn lub wspólny warunek & Fold z \texttt{+}, \texttt{*}, \texttt{\&\&} lub \texttt{||} \\
\hline
Przekazanie argumentów dalej z zachowaniem kategorii & \texttt{Args\&\&...}, \texttt{std::forward<Args>(args)...} \\
\hline
Zapis heterogenicznych wartości & \texttt{std::tuple<Ts...>} \\
\hline
Wywołanie funkcji elementami krotki & \texttt{std::apply} (C++17) \\
\hline
Rozwinięcie krotki przy pomocy indeksów & \texttt{std::index\_sequence} \\
\hline
Lista typów do obliczeń podczas kompilacji & Wariadyczny szablon klasy i specjalizacje \\
\hline
Ograniczenie do odpowiednich typów & Fold z koncepcjami i \texttt{requires} (C++20) \\
\hline
Zmienna długość listy jednolitych danych w czasie wykonania & \texttt{std::vector} lub \texttt{std::span}, nie pakiet \\
\hline
\end{tabular}
\end{center}

\section{Podsumowanie i zadania do samodzielnego sprawdzenia}

\begin{enumerate}
    \item \texttt{class... Ts} deklaruje \emph{pakiet typów},
    a \texttt{Ts... xs} wprowadza \emph{pakiet parametrów funkcji}.
    \item \texttt{sizeof...(Ts)} oblicza liczbę elementów pakietu
    podczas kompilacji.
    \item \texttt{pattern...} rozwija cały \emph{wzorzec},
    nie tylko pojedynczy identyfikator.
    \item Fold expressions (od C++17) zastępują wiele prostych
    rekurencyjnych implementacji z C++11.
    \item Dla pustego pakietu istotne są element neutralny i
    reguły operatora; nie każda forma fold jest poprawna.
    \item W wrapperach zachowuj kategorię wartości argumentów,
    używając referencji przekazujących i \texttt{std::forward}.
    \item \texttt{std::tuple}, \texttt{std::apply} i
    \texttt{std::index\_sequence} współpracują z pakietami,
    umożliwiając operacje na heterogenicznych danych.
    \item W C++20 używaj \texttt{concepts}/\texttt{requires},
    aby błędne wywołania odrzucać na granicy interfejsu.
\end{enumerate}

\paragraph{Pytania kontrolne.}
Spróbuj przewidzieć wyniki poniższych wyrażeń bez uruchamiania
programu. Następnie zweryfikuj je za pomocą \texttt{static\_assert}:

\begin{lstlisting}
// static_assert(sizeof...(int) == 1); // ERROR: int is not a pack

// For the functions defined above:
static_assert(product() == 1);
static_assert(left_sub(10, 3) == 87);
static_assert(right_sub(10, 3) == 107);
static_assert(count_integral_v<int, double, char> == 2);
\end{lstlisting}

Pierwsza linia jest \emph{celowo niepoprawna}: operator
\texttt{sizeof...} wymaga nazwy pakietu, a nie zwykłego typu.

% Wykorzystane punkty odniesienia (bez kopiowania tresci):
% https://en.cppreference.com/w/cpp/language/parameter_pack
% https://en.cppreference.com/w/cpp/language/fold
% https://en.cppreference.com/w/cpp/utility/forward
% https://en.cppreference.com/w/cpp/utility/apply
% https://en.cppreference.com/w/cpp/utility/integer_sequence

% Rozdzial 12 w kolejnosci cpp.tex — metaprogramowanie w C++20.
% Fragment do wlaczenia po rozdziale 11.
% Wymagane pakiety: listings, array, amsmath, hyperref.
% Uzycie: \input{rozdzial_9_metaprogramowanie_cpp.tex}

\chapter[Metaprogramowanie w C++]{Zaawansowane metaprogramowanie: szablony, cechy typów i obliczenia podczas kompilacji}
\label{chap:meta-cpp}

\section{Czym jest metaprogramowanie?}

\emph{Metaprogramowanie} oznacza pisanie kodu, który na etapie kompilacji
wyznacza typy, własności typów, wartości lub strukturę później
kompilowanego programu. Szablon nie jest gotową funkcją czy klasą:
stanowi przepis, z którego kompilator tworzy konkretne wersje
dla podanych typów. Nazywamy to \emph{instancjowaniem szablonu}.
Na przykład wywołanie
\texttt{twice(21)} w programie poniżej każe kompilatorowi przygotować
wersję funkcji dla \texttt{int}; program podczas działania wywołuje
już tę wersję. W C++20 do dyspozycji mamy kilka narzędzi:
\begin{itemize}
\item \textbf{szablony klas i funkcji} — generowanie kodu dla różnych typów;
\item \textbf{\texttt{type\_traits}} — obliczanie własności i przekształcanie typów;
\item \textbf{specjalizacje i SFINAE} — wybór poprawnych wersji szablonów;
\item \textbf{\texttt{if constexpr}, \texttt{constexpr}, \texttt{consteval}} —
      wykonywanie obliczeń i selekcja kodu w czasie kompilacji;
\item \textbf{\texttt{concepts}} — czytelne formułowanie wymagań wobec typów;
\item \textbf{CRTP} — wykorzystanie dziedziczenia do polimorfizmu statycznego.
\end{itemize}

Metaprogramowanie ma praktyczne cele: eliminowanie powtarzalnego kodu,
uzyskiwanie lepszych błędów kompilacji, wybór wydajnej implementacji oraz
zapewnienie poprawności określonych operacji \emph{przed uruchomieniem} programu.
Nie oznacza automatycznie zwiększenia wydajności — kosztem mogą być dłuższa
kompilacja, powiększenie kodu maszynowego i trudniejsze komunikaty diagnostyczne.

\begin{lstlisting}
#include <type_traits>

template<class T>
T twice(T x) {
    static_assert(std::is_arithmetic_v<T>);
    return x + x;
}

int a = twice(21);       // T = int
// twice("text");       // static_assert fails
\end{lstlisting}

\paragraph{Związki z innymi rozdziałami.}
Wykorzystamy szablony wariadyczne (rozdział~11),
przenoszenie i referencje (rozdział~6) oraz lambdy (rozdział~9).
Koncepcje szerzej poznamy w rozdziale~13; tutaj pokazujemy
przede wszystkim mechanizmy, na których się opierają.
Wszystkie oznaczone jako poprawne przykłady odnoszą się do \textbf{C++20}.

\section{Instancjowanie szablonów, parametry i zależne nazwy}
\label{sec:cpp9-instantiation}

\subsection{Parametry typowe, nietypowe i szablonowe}

Szablon może przyjmować argumenty będące typami, wartościami stałymi
lub innymi szablonami. Parametr wartości (\textit{non-type template parameter},
NTTP) w C++20 może również mieć odpowiedni typ klasowy spełniający
wymagania tzw. typu strukturalnego.

\begin{lstlisting}
#include <array>
#include <cstddef>

template<class T, std::size_t N>
struct Buffer {
    std::array<T, N> data{};
};

Buffer<double, 4> samples; // T=double, N=4
\end{lstlisting}

\texttt{Buffer<double,4>} i \texttt{Buffer<double,8>} są różnymi typami.
Parametr \texttt{N} jest znany podczas kompilacji i może wystąpić tam,
gdzie język wymaga wyrażenia stałego, np. jako rozmiar \texttt{std::array}.

\subsection{Parametr będący obiektem klasy (C++20)}

Od C++20 odpowiednie \emph{typy strukturalne} mogą występować jako
parametry szablonu będące wartościami. Umożliwia to np. identyfikowanie
szablonu przez napis znany podczas kompilacji:

\begin{lstlisting}
#include <cstddef>

template<std::size_t N>
struct fixed_string {
    char data[N]{};
    constexpr fixed_string(const char (&s)[N]) {
        for (std::size_t i = 0; i < N; ++i)
            data[i] = s[i];
    }
};

template<fixed_string Name>
struct NamedCommand {
    static constexpr auto name = Name;
};

static_assert(NamedCommand<"run">::name.data[0] == 'r');
\end{lstlisting}

To przykład zaawansowanej techniki; zwykły obiekt \texttt{std::string}
nie jest w C++20 strukturalnym argumentem szablonu. Rozmiar i
zawartość napisu wpływają tutaj na tożsamość specjalizacji szablonu.

\subsection{Nazwy zależne: \texttt{typename} i \texttt{template}}

W definicji szablonu kompilator nie zawsze wie, czy nazwa zależna od
parametru jest typem, czy wartością. W wielu kontekstach musimy to
jednoznacznie zaznaczyć:

\begin{lstlisting}
#include <vector>

template<class Container>
auto first_copy(const Container& c) {
    typename Container::value_type x = c.front();
    return x;
}

template<class T>
struct Wrapper {
    template<class U>
    static constexpr bool same = false;
};

template<class T>
constexpr bool use_member_template() {
    return Wrapper<T>::template same<int>;
}
\end{lstlisting}

Słowo \texttt{typename} wskazuje, że zależna nazwa członkowska
jest typem. Zapis \texttt{::template same<int>} informuje,
że po operatorze \texttt{::} występuje szablon składowy, a znaki
\texttt{<} i \texttt{>} są delimitatorami argumentów szablonu.
Dokładne reguły zależą od kontekstu; obu słów nie należy dopisywać
mechanicznie wszędzie.

\section{Metafunkcje i \texttt{std::integral\_constant}}
\label{sec:cpp9-integral-constant}

Poniższa silnia obliczana przez szablon pokazuje starszy styl
metaprogramowania. Gdy wynikiem ma być zwykła liczba, funkcja
\texttt{constexpr} jest dziś zazwyczaj prostsza. Szablony klas
pozostają potrzebne zwłaszcza wtedy, gdy wynikiem obliczenia ma być
\emph{typ}.

Klasyczne metaprogramowanie szablonowe opierało się na klasach, które
\emph{przenosiły wyniki obliczenia w swoich składowych statycznych}.
Dla przykładu silnię można wyznaczyć przez rekurencyjne instancjowanie
szablonów:

\begin{lstlisting}
template<unsigned N>
struct Factorial {
    static constexpr unsigned value =
        N * Factorial<N - 1>::value;
};

template<>
struct Factorial<0> {
    static constexpr unsigned value = 1;
};

static_assert(Factorial<5>::value == 120);
\end{lstlisting}

Specjalizacja dla zera pełni rolę warunku stopu. Bez niej kompilator
próbowałby tworzyć kolejne instancje, aż osiągnąłby limit głębokości.
Obliczenie ma charakter rekursji \emph{w strukturze typów}, a nie
rekursji wywołań funkcji wykonywanych w czasie działania programu.

Standardowym fundamentem takich metafunkcji jest
\texttt{std::integral\_constant<T,v>}:

\begin{lstlisting}
#include <type_traits>

using TrueType = std::integral_constant<bool, true>;
static_assert(TrueType::value);
static_assert(std::bool_constant<true>::value);
static_assert(std::true_type::value);
static_assert(!std::false_type::value);

template<unsigned N>
struct PowerOfTwo :
    std::integral_constant<unsigned, 2 * PowerOfTwo<N-1>::value> {};

template<>
struct PowerOfTwo<0> : std::integral_constant<unsigned, 1> {};

static_assert(PowerOfTwo<5>::value == 32);
\end{lstlisting}

Predykaty takie jak \texttt{std::is\_pointer} udostępniają wartość
\texttt{::value}. Od C++17 można używać krótszej wersji
z przyrostkiem \texttt{\_v}, na przykład:
\begin{center}
\texttt{std::is\_pointer\_v<T>}.
\end{center}

\subsection{Metafunkcja szablonowa czy funkcja \texttt{constexpr}?}

\begin{lstlisting}
constexpr unsigned factorial(unsigned n) {
    unsigned result = 1;
    for (unsigned i = 2; i <= n; ++i)
        result *= i;
    return result;
}

static_assert(factorial(5) == 120);
unsigned x = 6;
unsigned y = factorial(x); // allowed at run time
\end{lstlisting}

Jeśli wynik jest zwykłą wartością, funkcja \texttt{constexpr} jest
najczęściej prostsza niż rekurencyjna metafunkcja. Klasyczne metafunkcje
pozostają użyteczne, gdy \emph{wynikiem jest typ}, gdy korzystamy ze
specjalizacji lub interoperujemy ze starszym API.

\section{Specjalizacje: pełne i częściowe}
\label{sec:cpp9-specialization}

Szablon podstawowy definiuje przypadek ogólny. \emph{Specjalizacja pełna}
dotyczy konkretnego zestawu argumentów, a \emph{specjalizacja częściowa}
obsługuje pewną ich rodzinę:

\begin{lstlisting}
#include <type_traits>

template<class T>
struct is_raw_pointer : std::false_type {};

template<class T>
struct is_raw_pointer<T*> : std::true_type {};

static_assert(!is_raw_pointer<int>::value);
static_assert(is_raw_pointer<double*>::value);
static_assert(!is_raw_pointer<int* const>::value);
\end{lstlisting}

Ostatni wynik nie jest pomyłką: \texttt{int* const} jest
\emph{stałym wskaźnikiem}, nie typem postaci \texttt{T*}.
Jeśli chcemy ignorować kwalifikatory zewnętrzne, możemy najpierw
użyć \texttt{std::remove\_cv\_t<T>}.

\begin{lstlisting}
template<class T>
inline constexpr bool is_any_raw_pointer_v =
    is_raw_pointer<std::remove_cv_t<T>>::value;

static_assert(is_any_raw_pointer_v<int* const>);
\end{lstlisting}

\paragraph{Ważne ograniczenie.}
Dla \textbf{szablonów funkcji nie ma specjalizacji częściowej}.
W takich przypadkach wykorzystujemy przeciążanie, koncepcje,
\texttt{if constexpr} lub specjalizację pomocniczego szablonu klasy.
Pełna specjalizacja funkcji istnieje, ale jej interakcja z przeciążeniami
bywa zaskakująca; w nowym kodzie zwykle preferujemy przeciążenia.

\section{\texttt{type\_traits}: pytania o typy}
\label{sec:cpp9-traits}

Biblioteka \texttt{<type\_traits>} jest zbiorem standardowych metafunkcji.
Dzielimy je na \emph{predykaty} (wynik logiczny) i \emph{transformacje}
(wynik będący typem). Ich składnia opiera się na ustalonych konwencjach:
\begin{itemize}
\item \texttt{Trait<T>::value} lub \texttt{Trait\_v<T>} dla predykatów;
\item \texttt{typename Trait<T>::type} lub \texttt{Trait\_t<T>}
      dla transformacji.
\end{itemize}

\begin{lstlisting}
#include <type_traits>
#include <string>

static_assert(std::is_integral_v<int>);
static_assert(std::is_floating_point_v<double>);
static_assert(std::is_same_v<int, int>);
static_assert(std::is_pointer_v<const char*>);
static_assert(std::is_constructible_v<std::string, const char*>);
static_assert(std::is_move_constructible_v<std::string>);

using A = std::remove_const_t<const int>; // int
using B = std::remove_reference_t<int&>; // int
using C = std::remove_cvref_t<const int&>; // int (C++20)
using D = std::add_pointer_t<int>;        // int*
static_assert(std::is_same_v<A, B>);
static_assert(std::is_same_v<C, int>);
static_assert(std::is_same_v<D, int*>);
\end{lstlisting}

\begin{center}
\small
\begin{tabular}{|>{\raggedright\arraybackslash}p{5.5cm}|>{\raggedright\arraybackslash}p{7.0cm}|}
\hline
\textbf{Cecha} & \textbf{Co sprawdza lub wyznacza?} \\
\hline
\texttt{is\_same\_v<A,B>} & Identyczność typów (wraz z istotnymi kwalifikatorami). \\
\hline
\texttt{is\_convertible\_v<A,B>} & Dostępność konwersji niejawnej \texttt{A} do \texttt{B}. \\
\hline
\texttt{is\_constructible\_v<T,Args...>} & Możliwość skonstruowania \texttt{T} z argumentów zadanych typów. \\
\hline
\texttt{is\_nothrow\_move\_}\newline\texttt{constructible\_v<T>} & Deklarowana bezwyjątkowość konstrukcji przenoszącej. \\
\hline
\texttt{is\_trivially\_copyable\_v<T>} & Własność pozwalająca na określone operacje z reprezentacją bajtową. \\
\hline
\texttt{remove\_cvref\_t<T>} & Usunięcie referencji i kwalifikatorów \texttt{const}/\texttt{volatile} z poziomu wierzchniego. \\
\hline
\texttt{common\_type\_t<A,B>} & Wspólny typ wynikowy dla danych typów (o ile istnieje). \\
\hline
\texttt{invoke\_result\_t<F,Args...>} & Typ wyniku wywołania poprzez \texttt{std::invoke}. \\
\hline
\end{tabular}
\end{center}

\paragraph{\texttt{remove\_cvref} a \texttt{decay}.}
\texttt{std::decay\_t<T>} oprócz usunięcia referencji i zewnętrznych
kwalifikatorów może przekształcić tablicę w wskaźnik oraz typ funkcji
w wskaźnik do funkcji. Nie są to zatem synonimy.

\begin{lstlisting}
using X = std::remove_cvref_t<int(&)[3]>; // int[3]
using Y = std::decay_t<int(&)[3]>;        // int*
static_assert(std::is_same_v<Y, int*>);
\end{lstlisting}

\paragraph{Czego predykat nie gwarantuje?}
\texttt{std::is\_move\_constructible\_v<T>} może być prawdziwe również
wtedy, gdy wywołanie z argumentem rvalue używa konstruktora kopiującego.
Podobnie \texttt{is\_trivially\_copyable} nie oznacza, że każdy odczyt
bajtów, rzutowanie lub serializacja do pliku będzie poprawna bądź
przenośna. Cechy typów opisują wybrane własności języka, nie
pełny kontrakt semantyczny.

\subsection{Składanie predykatów z zachowaniem krótkiego obliczania}

\texttt{std::conjunction}, \texttt{std::disjunction} i
\texttt{std::negation} łączą predykaty, które mają składową
\texttt{::value}. Dwa pierwsze narzędzia zatrzymują odczytywanie \texttt{::value}
z kolejnych predykatów, kiedy wynik jest już znany. Trzeba jednak
poprawnie sformować same argumenty typowe tych metafunkcji.

\begin{lstlisting}
#include <type_traits>

template<class T>
using SmallSignedInteger = std::conjunction<
    std::is_integral<T>,
    std::is_signed<T>,
    std::bool_constant<(sizeof(T) <= sizeof(int))>>;

static_assert(SmallSignedInteger<short>::value);
static_assert(!SmallSignedInteger<unsigned>::value);
\end{lstlisting}

Dla prostych stałych logicznych zapis
\texttt{(cond1 \&\& cond2)} bywa czytelniejszy.
\texttt{conjunction} jest użyteczne przede wszystkim wtedy, gdy
poszczególne warunki są odrębnymi metafunkcjami, a ich instancjowanie
może być kosztowne lub niepoprawne bez wcześniejszego warunku.

\section{Metafunkcje zwracające typy i leniwy wybór}
\label{sec:cpp9-transforms}

Tworząc własną transformację typu, zwyczajowo zapisujemy wynik jako
\texttt{using type = ...} i udostępniamy wygodny alias:

\begin{lstlisting}
template<class T>
struct storage_type {
    using type = std::remove_cvref_t<T>;
};

template<class T>
using storage_type_t = typename storage_type<T>::type;

static_assert(std::is_same_v<storage_type_t<const long&>, long>);
\end{lstlisting}

Ważnym narzędziem jest \texttt{std::conditional\_t<B,T,F>}, wybierające
jeden z dwóch typów w zależności od znanej w kompilacji wartości logicznej:

\begin{lstlisting}
template<class T>
using accumulator_t = std::conditional_t<
    std::is_integral_v<T>, long long, long double>;

static_assert(std::is_same_v<accumulator_t<int>, long long>);
static_assert(std::is_same_v<accumulator_t<float>, long double>);
\end{lstlisting}

\paragraph{Pułapka eager instantiation.}
\texttt{conditional\_t} wybiera typ spośród \emph{już utworzonych
argumentów typowych}. Nie chroni przed błędami powstającymi podczas
formowania któregoś z nich:

\begin{lstlisting}
// Incorrect for int: typename T::value_type is invalid.
// template<class T>
// using bad = std::conditional_t<
//     std::is_integral_v<T>, T, typename T::value_type>;
\end{lstlisting}

Aby uzyskać prawdziwie warunkowe tworzenie typu, można zastosować
specjalizację pomocniczej klasy:

\begin{lstlisting}
template<class T, bool = std::is_integral_v<T>>
struct value_or_self {
    using type = typename T::value_type;
};

template<class T>
struct value_or_self<T, true> {
    using type = T;
};

template<class T>
using value_or_self_t = typename value_or_self<T>::type;

static_assert(std::is_same_v<value_or_self_t<int>, int>);
\end{lstlisting}

Podstawowa definicja zawiera nazwę zależną, która jest sprawdzana przy
instancjowaniu odpowiedniego przypadku; nie jest potrzebna dla \texttt{int}.

\section{\texttt{decltype}, \texttt{declval} i konteksty bez obliczania}

\begin{center}\small
\begin{tabularx}{\textwidth}{@{}p{0.32\textwidth}X@{}}
\toprule
Narzędzie & Za co odpowiada \\
\midrule
\code{decltype(expr)} & Ustala typ wyrażenia bez wykonywania go; szczególna reguła dotyczy nieujętej w nawiasy nazwy. \\
\code{decltype((x))} & Dla lvalue \code{x} zwraca typ referencyjny, np. \texttt{T\&}. \\
\code{std::declval<T>()} & Udostępnia fikcyjne wyrażenie typu \code{T} wyłącznie w kontekście nieobliczanym. \\
\code{std::void_t<...>} & Zamienia poprawne podstawienia typów na \code{void}; przydatne w starszych technikach wykrywania. \\
\code{if constexpr} & W szablonie wybiera gałąź podczas kompilacji i odrzuca niepotrzebną gałąź po instancjacji. \\
\bottomrule
\end{tabularx}
\end{center}
\label{sec:cpp9-decltype}

\texttt{decltype(expr)} pozwala odczytać typ wyrażenia bez wykonywania
samego wyrażenia. Subtelność stanowi różnica między nieujętym w nawiasy
identyfikatorem a ogólnym wyrażeniem:

\begin{lstlisting}
int x = 5;
static_assert(std::is_same_v<decltype(x), int>);
static_assert(std::is_same_v<decltype((x)), int&>);
\end{lstlisting}

Wyrażenie \texttt{(x)} jest \textit{lvalue}, więc \texttt{decltype((x))}
zwraca \texttt{int\&}. Reguła ta ma bezpośredni związek z kategoriami
wartości (rozdział~6).

\texttt{std::declval<T>()} pozwala użyć argumentu typu \texttt{T}
w \emph{kontekście niewartościowanym}, nawet jeśli nie możemy utworzyć
jego obiektu. Nie wolno rzeczywiście wykonywać wywołania
\texttt{declval}.

\begin{lstlisting}
#include <utility>
#include <type_traits>

template<class T>
using dereference_t = decltype(*std::declval<T&>());

static_assert(std::is_same_v<dereference_t<int*>, int&>);

struct NonDefault {
    NonDefault() = delete;
    int method() const;
};

using Result = decltype(std::declval<const NonDefault&>().method());
static_assert(std::is_same_v<Result, int>);
\end{lstlisting}

Typ wyniku wywołania można również wyznaczyć za pomocą gotowej metafunkcji:
\begin{center}
\texttt{std::invoke\_result\_t<F, Args...>}.
\end{center}

\section{SFINAE: błędne podstawienie to nie zawsze błąd kompilacji}
\label{sec:cpp9-sfinae}

Ta część pomaga czytać biblioteki napisane przed C++20.
W nowym kodzie wymagania wobec typów zwykle prościej wyrazić
przez \texttt{requires}, opisane w rozdziale~13.

\textbf{SFINAE} (\textit{Substitution Failure Is Not An Error}) mówi,
że niektóre błędy powstałe podczas podstawiania argumentów do szablonu
\emph{nie kończą kompilacji}: dana specjalizacja zostaje wyłączona
z zestawu kandydatów lub odpowiednia specjalizacja częściowa nie pasuje.
Dotyczy to określonych kontekstów bezpośredniego podstawiania;
\textbf{nie} jest mechanizmem przechwytującym dowolne błędy w ciele funkcji
ani skutki uboczne instancjowania innych typów.

\subsection{\texttt{std::enable\_if\_t}}

\begin{lstlisting}
#include <type_traits>

template<class T>
constexpr std::enable_if_t<std::is_integral_v<T>, T>
doubled(T x) {
    return x + x;
}

static_assert(doubled(3) == 6);
// doubled(2.5); // no matching function
\end{lstlisting}

Dla typu \texttt{double} nie istnieje \texttt{type} w
\texttt{std::enable\_if<false,double>}; funkcja zostaje wyłączona
z odpowiedniego zbioru kandydatów. Do konstruktorów, operatorów i
różnych wariantów przeciążeń istnieją odmienne miejsca umieszczania
\texttt{enable\_if}; łatwo przy tym uzyskać przypadkowe powtórne
deklaracje lub nieczytelne błędy.

\subsection{\texttt{void\_t} — detection idiom}

Od C++17 wygodnie wykrywać istnienie nazw i poprawność wyrażeń przy
użyciu \texttt{std::void\_t}. Jeżeli wszystkie podane argumenty
są prawidłowymi typami, rezultatem jest po prostu \texttt{void}:

\begin{lstlisting}
#include <type_traits>
#include <vector>

template<class, class = void>
struct has_value_type : std::false_type {};

template<class T>
struct has_value_type<T, std::void_t<typename T::value_type>>
    : std::true_type {};

static_assert(has_value_type<std::vector<int>>::value);
static_assert(!has_value_type<int>::value);
\end{lstlisting}

W podstawieniu \texttt{T=int} wyrażenie \texttt{int::value\_type}
nie ma sensu; specjalizacja częściowa nie jest więc wybierana.
Analogicznie można badać poprawność wyrażenia z \texttt{decltype}:

\begin{lstlisting}
#include <type_traits>
#include <utility>

template<class, class = void>
struct has_preincrement : std::false_type {};

template<class T>
struct has_preincrement<T,
    std::void_t<decltype(++std::declval<T&>())>>
    : std::true_type {};

static_assert(has_preincrement<int>::value);
\end{lstlisting}

\subsection{SFINAE kontra \texttt{requires} (C++20)}

Konstrukcję wykrywającą operację możemy wyrazić bezpośrednio za pomocą
koncepcji z rozdziału~13:

\begin{lstlisting}
#include <concepts>

template<class T>
concept Preincrementable = requires(T& x) {
    { ++x } -> std::same_as<T&>;
};

template<Preincrementable T>
void advance_one(T& x) { ++x; }
\end{lstlisting}

Metafunkcja \texttt{has\_preincrement<T>} sprawdza jedynie poprawność
operacji \texttt{++x}. Z kolei koncepcja \texttt{Preincrementable}
dodatkowo wymaga, aby jej \emph{typ wyniku} odpowiadał dokładnie \texttt{T\&}.
Różnica ta jest zamierzona: ograniczenia należy definiować zgodnie z
kontraktem, a nie tylko według dostępności pewnego wyrażenia.

\section{\texttt{if constexpr}: wybór implementacji na podstawie typu}
\label{sec:cpp9-ifconstexpr}

Od C++17 często zamiast rekurencyjnych specjalizacji wybieramy gałąź
wewnątrz pojedynczego szablonu. Niebrana gałąź zależna od argumentu
szablonu nie jest instancjowana dla danego zestawu argumentów:

\begin{lstlisting}
#include <type_traits>

template<class T>
constexpr auto number_or_value(const T& x) {
    if constexpr (std::is_arithmetic_v<T>)
        return x;
    else
        return x.value();
}

struct Box {
    constexpr int value() const { return 7; }
};
static_assert(number_or_value(3) == 3);
static_assert(number_or_value(Box{}) == 7);
\end{lstlisting}

\paragraph{Dlaczego nie zwykłe \texttt{if}?}
W zwykłym \texttt{if} kompilator musi przeanalizować obie gałęzie
pod kątem poprawności dla instancjowanego kodu. Dla \texttt{int}
wyrażenie \texttt{x.value()} nie jest poprawne, a sam warunek
\texttt{std::is\_arithmetic\_v<int>} tego nie zmieni.

\paragraph{Ważna subtelność.}
\texttt{if constexpr(false)} nie jest zamiennikiem
\texttt{\#if 0}: poza kontekstem szablonowym odrzucona gałąź nadal
podlega standardowym kontrolom poprawności. Ponadto mechanizm nie
usprawiedliwia konstrukcji zawsze niepoprawnych niezależnie od
argumentów szablonu.

\section{Pakiety typów, sekwencje indeksów i fold expressions}
\label{sec:cpp9-type-lists}

Metaprogramowanie dobrze współpracuje z pakietami parametrów
(rozdział~11). Zamiast tworzyć klasyczne rekurencyjne listy typów,
często możemy używać rozwinięcia pakietu i wyrażenia zwijającego.

\begin{lstlisting}
#include <type_traits>

template<class... Ts>
struct type_list {};

template<class T, class... Ts>
inline constexpr bool contains_v =
    (std::is_same_v<T, Ts> || ...); // empty pack -> false

static_assert(contains_v<int, char, int, double>);
static_assert(!contains_v<long, char, int, double>);
static_assert(!contains_v<int>);
\end{lstlisting}

Zamiast tworzyć specjalizacje dla każdego elementu pakietu, korzystamy
z jednego \texttt{fold expression}. Gdy potrzebna jest reprezentacja
indeksów \texttt{0,1,...,N-1}, biblioteka udostępnia
\texttt{std::make\_index\_sequence<N>}:

\begin{lstlisting}
#include <tuple>
#include <utility>
#include <cstddef>

template<class Tuple, std::size_t... I>
constexpr auto reversed_impl(const Tuple& t,
                             std::index_sequence<I...>) {
    return std::make_tuple(std::get<sizeof...(I)-1-I>(t)...);
}

template<class... Ts>
constexpr auto reversed(const std::tuple<Ts...>& t) {
    return reversed_impl(t, std::index_sequence_for<Ts...>{});
}

static_assert(reversed(std::tuple{1, 2, 3}) ==
              std::tuple{3, 2, 1});
\end{lstlisting}

\paragraph{Czy zawsze potrzebujemy indeksów?}
Nie. W przypadku prostego wykonania funkcji na elementach krotki
w C++17 często wystarczy \texttt{std::apply}; sekwencja indeksów
jest potrzebna zwłaszcza wtedy, gdy chcemy \emph{przekształcić pozycje}
lub mieć indywidualny dostęp do każdego indeksu na etapie kompilacji.

\section{\texttt{constexpr}, \texttt{consteval} i obliczenia na danych}
\label{sec:cpp9-constant-eval}

Szablony nie są jedynym sposobem wykonania pracy w trakcie kompilacji.
Od C++11 rozwijano możliwości \texttt{constexpr}, a w C++20 mamy
również \texttt{consteval}. Różnica jest istotna:
\begin{itemize}
\item funkcja \texttt{constexpr} \emph{może} uczestniczyć w obliczeniach
stałych, a w dozwolonych wywołaniach działa też podczas wykonania;
\item wywołanie funkcji \texttt{consteval} musi spełniać reguły
właściwe obliczaniu w czasie kompilacji (\textit{immediate function}).
\end{itemize}

\begin{lstlisting}
constexpr int square(int x) { return x * x; }
consteval int cube_ct(int x) { return x * x * x; }

static_assert(square(4) == 16);
static_assert(cube_ct(3) == 27);

int n = 4;
int a = square(n); // runtime evaluation possible
// int b = cube_ct(n); // error: n is not constant
\end{lstlisting}

Przykładem praktycznym jest haszowanie \texttt{std::string\_view}
w stałym wyrażeniu. Poniższy wariant FNV-1a nadaje się do zbudowania
stałej identyfikującej np. nazwę polecenia. Nie jest to hash
kryptograficzny i nie wyklucza kolizji.

\begin{lstlisting}
#include <cstdint>
#include <string_view>

consteval std::uint64_t fnv1a(std::string_view text) {
    std::uint64_t h = 14695981039346656037ull;
    for (unsigned char c : text) {
        h ^= c;
        h *= 1099511628211ull;
    }
    return h;
}

static_assert(fnv1a("abc") == 0xe71fa2190541574bull);
\end{lstlisting}

\paragraph{\texttt{constinit} i \texttt{std::is\_constant\_evaluated}.}
\texttt{constinit} zapewnia statyczną inicjalizację odpowiedniej
zmiennej, lecz sam nie czyni jej stałą. Funkcja
\texttt{std::is\_constant\_evaluated()} (C++20) pozwala w
pewnych przypadkach rozpoznać, czy bieżąca ewaluacja ma charakter
stały; nadużywanie jej może jednak zaciemniać semantykę funkcji.

\section{CRTP: polimorfizm statyczny}
\label{sec:cpp9-crtp}

\textbf{CRTP} (\textit{Curiously Recurring Template Pattern}) polega na
tym, że klasa dziedziczy po specjalizacji szablonu, której argumentem
jest ona sama:

\begin{lstlisting}
template<class Derived>
struct Printable {
    void print() const {
        static_cast<const Derived&>(*this).print_impl();
    }
};

struct Number : Printable<Number> {
    int value = 0;
    void print_impl() const;
};
\end{lstlisting}

Mechanizm \texttt{static\_cast} jest tutaj poprawny pod założeniem,
że obiekt bazowy faktycznie stanowi podobiekt \texttt{Derived}.
CRTP nie zapewnia dynamicznego sprawdzania typu ani wirtualnego
wywołania — poprawność wynika z konstrukcji hierarchii. Wywołanie
\texttt{print\_impl()} może być statycznie związane i często
łatwiej je wbudować w kod wynikowy (\textit{inline}).

Bardziej praktycznym przykładem jest domieszka (\textit{mixin})
dostarczająca wspólną operację, lecz działająca na danych
konkretnej klasy pochodnej:

\begin{lstlisting}
template<class Derived>
struct TranslationMixin {
    constexpr void translate(int dx, int dy) {
        auto& self = static_cast<Derived&>(*this);
        self.x += dx;
        self.y += dy;
    }
};

struct Point : TranslationMixin<Point> {
    int x{}, y{};
    constexpr Point(int a, int b) : x(a), y(b) {}
};

constexpr bool test_translation() {
    Point p{1, 2};
    p.translate(3, 4);
    return p.x == 4 && p.y == 6;
}
static_assert(test_translation());
\end{lstlisting}

Baza \texttt{TranslationMixin<Point>} jest instancjowana z typem
\texttt{Point}, chociaż sam typ \texttt{Point} nie jest jeszcze
kompletny w momencie deklarowania dziedziczenia. Odwołania do składowych
\texttt{x} i \texttt{y} zależą jednak od \texttt{Derived} i są
sprawdzane dopiero podczas instancjowania odpowiedniej funkcji.
W rzeczywistym API warto opisać wymagane składowe koncepcją lub
przenieść operację do jawnie wskazanego interfejsu.

\paragraph{CRTP a wirtualne funkcje.}
Klasyczny polimorfizm dynamiczny pozwala wybierać implementację
na podstawie rzeczywistego typu obiektu przez interfejs klasy bazowej,
także gdy typ jest nieznany podczas kompilacji kodu wywołującego.
CRTP wymaga znajomości konkretnego typu podczas instancjowania.
Nie jest ogólnym zamiennikiem klas abstrakcyjnych. W C++20 prostsze
zastosowania domieszek i interfejsów statycznych można często oprzeć
na \texttt{concepts} i funkcjach wolnostojących.

\section{Policy-based design: wybór zachowania przez typ}
\label{sec:cpp9-policy}

Szablon może przyjąć \emph{politykę}, której metody definiują wariant
algorytmu. Umożliwia to wybór zachowania bez wirtualnych wywołań,
a kompilator może całkowicie usunąć niepotrzebne operacje:

\begin{lstlisting}
#include <limits>
#include <stdexcept>

struct Checked {
    static int divide(int a, int b) {
        if (b == 0) throw std::invalid_argument("zero");
        if (a == std::numeric_limits<int>::min() && b == -1)
            throw std::overflow_error("result outside int");
        return a / b;
    }
};

struct Fast {
    static int divide(int a, int b) { return a / b; }
};

template<class Policy>
struct Calculator {
    int divide(int a, int b) const {
        return Policy::divide(a, b);
    }
};

// Calculator<Checked> checks both invalid integer cases;
// Calculator<Fast> requires a nonzero divisor and no overflow.
\end{lstlisting}

Polityka \texttt{Fast} nie zabezpiecza dzielenia przez zero ani
przepełnienia \texttt{INT\_MIN / -1}; oba przypadki mają
niezdefiniowane zachowanie dla \texttt{int}. \texttt{Checked}
odrzuca je przed dzieleniem. Dobór polityki jest częścią kontraktu API.

W realnym kodzie szablon polityki można ograniczyć koncepcją, a jeśli
liczba wariantów jest niewielka, zwykłe \texttt{if constexpr} może być
łatwiejsze w utrzymaniu niż rozbudowana hierarchia polityk.

\section{Kompletny przykład: dobór typu akumulatora}
\label{sec:cpp9-worked}

Załóżmy, że chcemy obliczać sumę liczb, lecz wybierać typ akumulatora
w zależności od argumentu: dla typów całkowitych \texttt{long long},
a dla zmiennoprzecinkowych \texttt{long double}. Rozwiązanie łączy
\texttt{concepts}, \texttt{type\_traits} i \texttt{conditional\_t}:

\begin{lstlisting}
#include <concepts>
#include <type_traits>
#include <vector>

template<class T>
concept Numeric = std::integral<T> || std::floating_point<T>;

template<Numeric T>
using sum_type_t = std::conditional_t<
    std::integral<T>, long long, long double>;

template<Numeric T>
sum_type_t<T> accumulate_values(const std::vector<T>& xs) {
    sum_type_t<T> result{};
    for (const T& x : xs) {
        result += static_cast<sum_type_t<T>>(x);
    }
    return result;
}

static_assert(std::is_same_v<sum_type_t<int>, long long>);
static_assert(std::is_same_v<sum_type_t<double>, long double>);
\end{lstlisting}

\paragraph{Ograniczenia.}
Przykład wyjaśnia selekcję typu, lecz \emph{nie daje gwarancji
braku przepełnienia}. Nawet \texttt{long long} ma ograniczony zakres,
a suma zmiennoprzecinkowa może akumulować błędy zaokrągleń.
Najpierw należy ustalić wymagania numeryczne, a dopiero potem
projektować odpowiednią metafunkcję.

\section{Czas kompilacji, diagnostyka i praktyczne ograniczenia}

Metaprogramowanie przenosi koszt części pracy z czasu wykonania na
czas kompilacji, ale samo w sobie nie oznacza wydajniejszego programu.
Każda odmienna specjalizacja szablonu może wymagać osobnego
instancjowania i wygenerowania kodu. Powstają wtedy m.in.:
\begin{itemize}
\item dłuższe kompilacje i większe zużycie pamięci kompilatora;
\item \emph{code bloat}, czyli powiększenie plików wynikowych przez wiele
podobnych instancji;
\item skomplikowane komunikaty błędów przy zagnieżdżonych metafunkcjach;
\item problemy z zależnościami między plikami nagłówkowymi;
\item ograniczenia głębokości instancjowania i kosztu ewaluacji stałej.
\end{itemize}

\paragraph{Pliki nagłówkowe i jawne instancjowanie.}
Definicje szablonów zazwyczaj umieszczamy w nagłówkach, ponieważ
kompilator musi je znać przy instancjowaniu z konkretnymi argumentami.
Możliwe jest również \emph{jawne instancjowanie} dla ustalonych typów,
co pozwala kontrolować zestaw generowanych specjalizacji:

\begin{lstlisting}
// In a header:
template<class T> T triple(T x);
extern template int triple<int>(int);

// In one .cpp file (definition must be visible here):
template<class T> T triple(T x) { return x + x + x; }
template int triple<int>(int);
\end{lstlisting}

Ten przykład przedstawia ideę rozdzielenia deklaracji i jawnego
instancjowania; obie oznaczone części należy umieścić w odpowiednich
plikach, a dla wywołań innych typów zapewnić dostęp do definicji.
Warto pamiętać o regule ODR z rozdziału~4.

\section{Jaką technikę wybrać?}

\begin{center}
\small
\begin{tabular}{|p{5.7cm}|p{6.8cm}|}
\hline
\textbf{Potrzeba} & \textbf{Najczęściej odpowiedni mechanizm} \\
\hline
Sprawdzenie wymaganych operacji w nowym API & \texttt{concepts}, \texttt{requires} (C++20) \\
\hline
Obliczenie wartości dla stałych danych & Funkcja \texttt{constexpr} \\
\hline
Wymuszenie ewaluacji w czasie kompilacji & \texttt{consteval} (C++20) \\
\hline
Warunkowy kod zależny od typu & \texttt{if constexpr} \\
\hline
Ustalenie własności typu & \texttt{std::is\_...\_v}, inne \texttt{type\_traits} \\
\hline
Zmiana lub wybór typu wynikowego & \texttt{...\_t}, specjalizacja pomocniczej klasy \\
\hline
Wsparcie biblioteki sprzed C++20 & SFINAE, \texttt{enable\_if}, \texttt{void\_t} \\
\hline
Wieloargumentowe cechy typów & Pakiety typów, \texttt{fold expressions} \\
\hline
Dostęp do pozycji elementów krotki & \texttt{index\_sequence}, \texttt{std::apply} \\
\hline
Wielokrotne użycie zachowania klasowego & CRTP / mixin, rozważ też koncepcje \\
\hline
Wybór wariantu strategii w typie & Szablon polityki (policy-based design) \\
\hline
\end{tabular}
\end{center}

\section{Typowe błędy i zasady projektowe}

\begin{enumerate}
\item \textbf{Mylenie typu z wartością.} \texttt{typename T::value\_type}
      to typ, a \texttt{Trait<T>::value} to wartość; nie są wymienne.
\item \textbf{Brak warunku stopu.} Rekurencyjny szablon klasy może
      prowadzić do nieograniczonego instancjowania i błędów głębokości.
\item \textbf{Założenie, że SFINAE przechwytuje wszystkie błędy.}
      Działa tylko w określonych kontekstach podstawiania.
\item \textbf{Nadużywanie \texttt{conditional\_t}.} Oba argumenty
      typowe trzeba poprawnie sformować, zanim zostanie dokonany wybór.
\item \textbf{Próba wywołania \texttt{declval}.} Można go używać wyłącznie
      w odpowiednich kontekstach niewartościowanych.
\item \textbf{Mylenie \texttt{decltype(x)} z \texttt{decltype((x))}.}
      Dodatkowe nawiasy mogą zmienić wynik na typ referencyjny.
\item \textbf{Mylenie \texttt{constexpr} z obowiązkiem obliczenia
      podczas kompilacji.} Nie każde wywołanie funkcji \texttt{constexpr}
      jest ewaluowane w czasie kompilacji.
\item \textbf{Przenoszenie semantyki do traits.} Typ spełniający
      syntaktyczne wymaganie nie musi spełniać własności biznesowej,
      np. braku utraty precyzji czy bezpiecznej serializacji.
\item \textbf{Niepotrzebny CRTP.} Gdy wystarczy wolna funkcja
      z koncepcją, hierarchia klas często niczego nie upraszcza.
\item \textbf{Przedwczesne metaprogramowanie.} Najprostsze czytelne
      rozwiązanie jest zazwyczaj lepsze od dodatkowej warstwy abstrakcji.
\end{enumerate}

\paragraph{Pytania kontrolne.}
\begin{enumerate}
\item Czym różni się instancjowanie szablonu od wykonania funkcji?
\item Dlaczego \texttt{Factorial<0>} wymaga osobnej specjalizacji?
\item Kiedy \texttt{std::decay\_t<T>} i \texttt{std::remove\_cvref\_t<T>}
      dają różne wyniki?
\item Dlaczego \texttt{std::conditional\_t} nie rozwiązuje wszystkich
      problemów warunkowego tworzenia typów?
\item Co dokładnie sprawdza \texttt{void\_t<decltype(++declval<T\&>())>}?
\item Kiedy niepowodzenie podstawienia jest SFINAE, a kiedy zwykłym
      błędem kompilacji?
\item Dlaczego \texttt{if constexpr} nie zastępuje preprocesora?
\item Czym różnią się \texttt{constexpr} i \texttt{consteval}?
\item W jakich sytuacjach CRTP nie zastąpi funkcji wirtualnych?
\item Kiedy lepiej zastosować \texttt{concepts} niż \texttt{enable\_if}?
\end{enumerate}

\paragraph{Dalsza lektura.}
Warto sięgnąć do dokumentacji \emph{cppreference} dotyczącej
\emph{Templates}, \emph{Type traits} i \emph{SFINAE}.
Pomocne są również hasła \texttt{std::void\_t}, \texttt{std::enable\_if},
\texttt{if constexpr}, \texttt{consteval} oraz
\texttt{std::integer\_sequence}. Do diagnozowania złożonych
błędów szablonów przydają się opcje \texttt{-Wall -Wextra -Wpedantic},
\texttt{-ftemplate-backtrace-limit=0} (GCC) oraz ograniczanie zasięgu
instancjowanego kodu do minimalnego przykładu.

% ROZDZIAL 5 -- Concepts i requires w C++20
% Fragment do wlaczenia do glownego dokumentu po rozdziale 4.
% Wymagane w preambule: pakiet listings i ustawienie language=C++.
% Uzycie: \input{rozdzial_5_concepts_cpp.tex}

\chapter[Concepts i requires w C++]{Koncepcje i ograniczenia szablonów (\textit{concepts} i \texttt{requires})}
\label{chap:concepts-cpp}

\section{Po co wprowadzono koncepcje?}

Szablony C++ pozwalają pisać kod niezależny od konkretnego typu.
Tradycyjny szablon deklaruje jednak zwykle \emph{co kompilator ma zrobić},
nie wyjaśniając wprost, \emph{jakich własności} wymaga od argumentów.
Jeśli operacja jest niepoprawna, błąd może zostać wykryty dopiero w ciele
szablonu, a komunikat diagnostyczny bywa długi.

\begin{lstlisting}
template<class T>
T twice(const T& x) {
    return x + x;           // requires operator+
}

// twice(7);                // OK
// twice(SomeType{});       // may fail inside function body
\end{lstlisting}

Od C++20 można umieścić wymagania na poziomie interfejsu:

\begin{lstlisting}
#include <concepts>

template<class T>
concept Addable = requires(const T& a, const T& b) {
    { a + b } -> std::convertible_to<T>;
};

template<Addable T>
T twice(const T& x) {
    return x + x;
}
\end{lstlisting}

Koncepcja (\textit{concept}) to nazwany predykat sprawdzany podczas
kompilacji, opisujący wymagania wobec argumentów szablonu.
Może dotyczyć istnienia operatorów, typów składowych, własności typów,
kategorii iteratorów i wielu innych cech. \emph{Constraints}
(ograniczenia) wpływają zarówno na dopuszczalność konkretyzacji,
jak i na wybór przeciążenia funkcji.

\paragraph{Wersja języka.}
W całym rozdziale zakładamy \textbf{C++20} oraz kompilację
np. poleceniem \texttt{g++ -std=c++20 -Wall -Wextra -Wpedantic main.cpp}.
Definicje koncepcji standardowych znajdują się przede wszystkim
w nagłówkach \texttt{<concepts>}, \texttt{<iterator>} i
\texttt{<ranges>}.

\section{Pierwsza koncepcja: składnia i zastosowanie}
\label{sec:concepts-pierwsza}

Definicja nazwanej koncepcji ma postać
\texttt{template<...> concept Nazwa = warunek;}. Warunek powinien
być wyrażeniem stałym typu \texttt{bool}. Definicja koncepcji
występuje na poziomie przestrzeni nazw; nie jest zwykłą funkcją.

\begin{lstlisting}
#include <concepts>
#include <type_traits>

template<class T>
concept Integer = std::is_integral_v<T>;

static_assert(Integer<int>);
static_assert(Integer<bool>);    // bool is integral
static_assert(!Integer<double>);
\end{lstlisting}

Koncepcję można sprawdzić wyrażeniem \texttt{Integer<int>} lub
\texttt{static\_assert}, ale nie należy jej mylić ze zwykłą
zmienną. W szczególności \texttt{concept} nie daje się jawnie
specjalizować w sposób analogiczny do szablonu klasy.

\subsection{Cztery podstawowe sposoby ograniczania funkcji}

Poniższe deklaracje opisują ten sam warunek: funkcja jest przeznaczona
dla typów całkowitych w sensie \texttt{std::integral}.

\begin{lstlisting}
#include <concepts>

template<class T>
requires std::integral<T>
T square1(T x) { return x * x; }

template<class T>
T square2(T x) requires std::integral<T> {
    return x * x;
}

template<std::integral T>
T square3(T x) { return x * x; }

std::integral auto square4(std::integral auto x) {
    return x * x;
}
\end{lstlisting}

W ostatnim przykładzie ograniczony jest nie tylko typ argumentu,
ale też \emph{typ wyniku}. Sama dedukcja typu wyniku nadal odbywa
się na podstawie wyrażenia \texttt{x * x}; koncept jest sprawdzany
po jej przeprowadzeniu. Zapis \texttt{std::integral auto} w parametrze
jest skrótem szablonu funkcji, tzw. \emph{abbreviated function template}.

\paragraph{Kiedy którą postać wybrać?}
Dla jednego prostego ograniczenia zapis z ograniczonym
parametrem szablonu, np. \texttt{template<\allowbreak std::integral T>},
jest krótki i czytelny. Jeśli zależność dotyczy kilku parametrów,
najwygodniejsza bywa klauzula \texttt{requires}.
Skrócona postać z \texttt{auto} przydaje się w niewielkich funkcjach
i lambdach.

\section{Dwie role słowa \texttt{requires}}

\begin{center}\small
\begin{tabularx}{\textwidth}{@{}p{0.34\textwidth}X@{}}
\toprule
Zapis & Za co odpowiada \\
\midrule
\code{concept C = ...} & Nadaje nazwę logicznemu ograniczeniu parametrów szablonu. \\
\code{requires C<T>} & Ogranicza udział szablonu w przeciążaniu do typów spełniających \code{C}. \\
\code{requires(T x) { x.f(); }} & Sprawdza poprawność wskazanego wyrażenia dla podstawionego typu. \\
\code{requires requires(...) { ... }} & Pierwsze \code{requires} wprowadza ograniczenie, drugie wyrażenie sprawdzające. \\
\code{static_assert(C<T>)} & Wymusza spełnienie warunku w danym miejscu i zgłasza błąd kompilacji, gdy jest fałszywy. \\
\bottomrule
\end{tabularx}
\end{center}
\label{sec:concepts-dwa-requires}

Identyczny zapis słowny \texttt{requires} występuje w dwóch
różnych konstrukcjach:
\begin{enumerate}
  \item \emph{requires-clause} -- klauzula ograniczająca deklarację;
  \item \emph{requires-expression} -- wyrażenie logiczne, które bada
  poprawność podanych wymagań.
\end{enumerate}

\begin{lstlisting}
#include <concepts>

template<class T>
concept Incrementable = requires(T value) { // requires-expression
    ++value;
};

template<class T>
requires Incrementable<T>                // requires-clause
void increment(T& x) {
    ++x;
}
\end{lstlisting}

Można je połączyć bez nadawania koncepcji osobnej nazwy:

\begin{lstlisting}
template<class T>
requires requires(T x) { ++x; }
void increment_again(T& x) {
    ++x;
}
\end{lstlisting}

Pierwsze \texttt{requires} w ostatnim przykładzie zaczyna klauzulę;
drugie wprowadza wyrażenie sprawdzające wymagania. To poprawna
składnia, choć przy wielokrotnym użyciu lepiej utworzyć nazwaną
koncepcję.

\section{Cztery rodzaje wymagań w \texttt{requires-expression}}

\begin{center}\small
\begin{tabularx}{\textwidth}{@{}p{0.31\textwidth}X@{}}
\toprule
Rodzaj wymagania & Co sprawdza \\
\midrule
Proste: \code{a + b;} & Czy wyrażenie jest poprawne składniowo dla danych typów. \\
Typowe: \code{typename T::value_type;} & Czy wskazany typ zależny istnieje. \\
Złożone: \code{{ f(x) } -> C;} & Czy wyrażenie jest poprawne oraz czy jego typ spełnia ograniczenie \code{C}. \\
Zagnieżdżone: \code{requires C<T>;} & Czy dodatkowy warunek logiczny jest spełniony. \\
\bottomrule
\end{tabularx}
\end{center}
\label{sec:concepts-cztery-wymagania}

Wyrażenie \texttt{requires} umożliwia opisanie czterech rodzajów
warunków. Parametry wpisane w nawiasach \texttt{requires(T x)}
są jedynie \emph{lokalnymi parametrami do sprawdzania wyrażeń};
nie tworzy się dla nich obiektów podczas działania programu.

\subsection{Wymaganie proste (\textit{simple requirement})}

\begin{lstlisting}
#include <ostream>

template<class T>
concept Printable = requires(std::ostream& os, const T& x) {
    os << x;                    // must be well-formed
};
\end{lstlisting}

Wyrażenie \texttt{os << x}
nie jest tutaj wykonywane; kompilator sprawdza jego poprawność.
Proste wymaganie nie nakłada warunku na typ zwracany.

\subsection{Wymaganie typu (\textit{type requirement})}

\begin{lstlisting}
#include <vector>

template<class T>
concept HasValueType = requires {
    typename T::value_type;    // nested type must exist
};

static_assert(HasValueType<std::vector<int>>);
static_assert(!HasValueType<int>);
\end{lstlisting}

Słowo \texttt{typename} jest tu składnikiem składni wymagania,
a nie deklaracją nowego typu.

\subsection{Wymaganie złożone (\textit{compound requirement})}

\begin{lstlisting}
#include <concepts>

template<class T>
concept AddableNoThrow = requires(const T& a, const T& b) {
    { a + b } noexcept -> std::same_as<T>;
};
\end{lstlisting}

Zapis \texttt{\{ wyrazenie \} noexcept -> Koncepcja;}
może sprawdzać jednocześnie: poprawność wyrażenia,
to, że jest ono \texttt{noexcept}, i własność typu
\texttt{decltype((wyrazenie))}. Obie części, \texttt{noexcept}
oraz \texttt{-> Koncepcja}, są opcjonalne.

\begin{lstlisting}
template<class T>
concept PlusReturnsT = requires(const T& a, const T& b) {
    { a + b } -> std::same_as<T>;
};

static_assert(PlusReturnsT<int>);
// A type with a proxy return value may not satisfy this.
\end{lstlisting}

\paragraph{Uwaga na typ i konwersje.}
Koncepcja \texttt{std::same\_as<T>} żąda \emph{dokładnie}
typu \texttt{T}, natomiast koncepcja
\texttt{std::convertible\_\allowbreak to<T>} dopuszcza konwersję.
Dla \texttt{int} operator \texttt{+} daje \texttt{int}, ale np.
dla małych typów całkowitych może zajść promocja do \texttt{int}.

\subsection{Wymaganie zagnieżdżone (\textit{nested requirement})}

\begin{lstlisting}
#include <type_traits>

template<class T>
concept SmallTrivial = requires {
    requires sizeof(T) <= 16;
    requires std::is_trivially_copyable_v<T>;
};

static_assert(SmallTrivial<int>);
\end{lstlisting}

Wymaganie zagnieżdżone zaczyna się kolejnym \texttt{requires}
i musi zawierać warunek stały. Jest wygodne, gdy wymaganie
dotyczy własności typu, a nie określonego wyrażenia.

\subsection{Przykład: połączenie różnych wymagań}

\begin{lstlisting}
#include <concepts>
#include <cstddef>

template<class T>
concept Indexable = requires(T& obj, std::size_t i) {
    typename T::value_type;
    { obj.size() } -> std::convertible_to<std::size_t>;
    obj[i];
};
\end{lstlisting}

To przykład \emph{minimalnego kontraktu składniowego}.
Koncepcja nie dowodzi, że indeks \texttt{i} jest w zakresie,
że operacja jest szybka ani że wynik indeksowania zachowuje się
jak zwykła referencja. Takie własności wymagają dodatkowych założeń.

\section{Standardowe koncepcje z biblioteki C++20}
\label{sec:concepts-standardowe}

Przed napisaniem własnej koncepcji warto sprawdzić, czy istnieje
gotowa w standardowej bibliotece.

\subsection{Typy, konwersje i dziedziczenie}

\begin{lstlisting}
#include <concepts>

static_assert(std::same_as<int, int>);
static_assert(!std::same_as<int, long>);
static_assert(std::integral<unsigned>);
static_assert(std::signed_integral<int>);
static_assert(std::unsigned_integral<unsigned>);
static_assert(std::floating_point<double>);
static_assert(std::convertible_to<int, double>);

struct Base {};
struct Derived : Base {};
static_assert(std::derived_from<Derived, Base>);
\end{lstlisting}

\texttt{std::derived\_from<D,B>} obejmuje publiczne i jednoznaczne
dziedziczenie oraz przypadek tego samego typu. Różni się pod tym
względem od samego \texttt{std::is\_base\_of\_v}.

\subsection{Konstruowalność, kopiowanie, przenoszenie i porównywanie}

Przydatne są m.in. koncepcje wyrażające inicjalizowalność,
konstruowalność, możliwość przenoszenia i kopiowania oraz porównywania:
\begin{lstlisting}
default_initializable    constructible_from
move_constructible       copy_constructible
movable                  copyable
semiregular              regular
equality_comparable      totally_ordered
\end{lstlisting}

\begin{lstlisting}
#include <concepts>
#include <memory>
#include <string>

static_assert(std::regular<int>);
static_assert(std::copyable<std::string>);
static_assert(std::movable<std::unique_ptr<int>>);
static_assert(!std::copyable<std::unique_ptr<int>>);
static_assert(std::totally_ordered<int>);
\end{lstlisting}

Koncepcje te opisują interfejs i zakładają także określone
\emph{prawa semantyczne} (np. sensowne porównywanie). Sam kompilator
nie potrafi udowodnić dowolnej własności semantycznej implementacji.

\subsection{Funkcje wywoływalne i predykaty}

\begin{lstlisting}
#include <concepts>

struct IsPositive {
    bool operator()(int x) const { return x > 0; }
};

static_assert(std::invocable<IsPositive, int>);
static_assert(std::predicate<IsPositive, int>);
static_assert(!std::predicate<IsPositive, void*>);
\end{lstlisting}

\texttt{std::invocable<F,Args...>} sprawdza możliwość wywołania,
natomiast \texttt{std::predicate} wymaga dodatkowo wyniku,
którego można użyć jako warunku logicznego, i wiąże się
z określonymi wymaganiami semantycznymi.

\subsection{Iteratory i zakresy}

Koncepcje iteratorów są udostępniane przez \texttt{<iterator>},
a koncepcje zakresów przez \texttt{<ranges>}.

\begin{lstlisting}
#include <ranges>
#include <vector>
#include <list>

static_assert(std::ranges::random_access_range<
              std::vector<int>>);
static_assert(std::ranges::bidirectional_range<
              std::list<int>>);
static_assert(!std::ranges::random_access_range<
              std::list<int>>);
\end{lstlisting}

W szczególności przydatne są koncepcje \emph{input},
\emph{forward}, \emph{bidirectional}, \emph{random access}
i \emph{contiguous range}. Pozwalają opisywać wymagania
algorytmów bez uzależniania ich od konkretnego kontenera.

\section{Łączenie koncepcji i definiowanie własnych kontraktów}

Wymagania można łączyć operatorami logicznymi
\texttt{\&\&} i \texttt{||}. Często lepiej budować niewielkie,
nazwane koncepcje niż powtarzać jedno długie wyrażenie.

\begin{lstlisting}
#include <concepts>

template<class T>
concept Arithmetic = std::integral<T> || std::floating_point<T>;

template<class T>
concept NonBooleanInteger =
    std::integral<T> && !std::same_as<T, bool>;

template<Arithmetic T>
T cube(T x) { return x * x * x; }

static_assert(Arithmetic<double>);
static_assert(!NonBooleanInteger<bool>);
\end{lstlisting}

Dla połączenia \texttt{A \&\& B} wymaganie \texttt{B} jest
sprawdzane dopiero, gdy \texttt{A} jest spełnione; analogicznie
dla \texttt{A || B} drugi składnik jest potrzebny dopiero,
gdy pierwszy nie jest spełniony. Ta reguła chroni często przed
podstawianiem wyrażeń, które dla danego typu nie mają sensu.

\begin{lstlisting}
#include <concepts>
#include <type_traits>

template<class T>
concept HasIntegralValueType =
    requires { typename T::value_type; } &&
    std::integral<typename T::value_type>;
\end{lstlisting}

W powyższym przykładzie odwołanie do \texttt{T::value\_type}
po prawej stronie jest istotne dopiero po potwierdzeniu,
że typ składowy istnieje. Należy jednak pamiętać, że złożone
reguły podstawiania i normalizacji koncepcji mają swoje szczegóły;
nie każdą źle sformułowaną deklarację naprawi samo zwarcie logiczne.

\section{Koncepcje z wieloma parametrami}

Koncepcja nie musi opisywać własności jednego typu.
Może także opisywać \emph{relację między typami}.

\begin{lstlisting}
#include <concepts>

template<class From, class To>
concept SameType = std::same_as<From, To>;

template<class Derived, class Base>
concept PublicChildOf = std::derived_from<Derived, Base>;

template<class T, class U>
concept AddableTogether = requires(const T& a, const U& b) {
    a + b;
};

static_assert(AddableTogether<int, double>);
\end{lstlisting}

Wymienione koncepcje \texttt{SameType} i
\texttt{PublicChildOf} służą tu jedynie ilustracji składni.
W praktyce używaj bezpośrednio standardowych
\texttt{std::same\_as} i \texttt{std::derived\_from},
jeżeli nie dodajesz do nich istotnego znaczenia domenowego.

\paragraph{Wstawianie argumentu przy ograniczonym parametrze.}

\begin{lstlisting}
struct Animal {};

template<std::derived_from<Animal> T>
void inspect(const T&) {}

// Equivalent constraint: std::derived_from<T, Animal>
\end{lstlisting}

W zapisie powyżej typ \texttt{T} jest podstawiany
\emph{na pierwszej pozycji} koncepcji, zatem
chodzi o warunek z poniższego przykładu:
\begin{lstlisting}
std::derived_from<T, Animal>
\end{lstlisting}

\section{Przeciążanie funkcji i wybór bardziej szczegółowego ograniczenia}
\label{sec:concepts-przeciazanie}

Ograniczenia uczestniczą w wyborze przeciążenia. Jeśli dwie funkcje
są dopasowane do argumentów, kompilator może preferować tę,
której ograniczenie jest bardziej szczegółowe (ang.
\emph{subsumption}).

\begin{lstlisting}
#include <concepts>
#include <string_view>

template<class T>
requires std::integral<T>
constexpr std::string_view category(T) {
    return "integral";
}

template<class T>
requires std::signed_integral<T>
constexpr std::string_view category(T) {
    return "signed integral";
}

static_assert(category(12) == "signed integral");
static_assert(category(12u) == "integral");
\end{lstlisting}

\texttt{std::signed\_integral} jest zbudowane na
\texttt{std::integral} i dodaje dodatkowe wymaganie. Dzięki temu
bardziej szczegółowe przeciążenie zostaje wybrane dla \texttt{int}.

\subsection{Ważna pułapka: podobny zapis to nie to samo ograniczenie}

Porządek ograniczeń nie jest ogólnym dowodzeniem implikacji
matematycznych. Kompilator porównuje znormalizowane ograniczenia
atomowe i ich \emph{tożsamość}, a nie tylko sens pozornie podobnych
wyrażeń.

\begin{lstlisting}
#include <type_traits>

template<class T>
concept HasSize = requires(const T& x) { x.size(); };

template<class T>
concept HasSizeAndCopy =
    HasSize<T> && std::is_copy_constructible_v<T>;

template<class T> requires HasSize<T>
int classify(const T&) { return 1; }

template<class T> requires HasSizeAndCopy<T>
int classify(const T&) { return 2; }
\end{lstlisting}

Druga koncepcja \emph{wykorzystuje tę samą} koncepcję
\texttt{HasSize}, dlatego kompilator rozpoznaje wspólny
składnik. Jeśli powtórzymy ręcznie nawet identycznie wyglądające
wymaganie opisujące metodę \texttt{size()},
nie mamy automatycznej gwarancji identyczności ograniczeń
atomowych.

\paragraph{Rekomendacja.}
Twórz proste nazwane koncepcje i buduj bardziej szczegółowe
z ich wykorzystaniem. Ułatwia to konserwację kodu i
przewidywanie wyboru przeciążeń.

\section{Koncepcje dla szablonów klas i ich specjalizacji}

Ograniczać można nie tylko funkcje, lecz również szablony klas.

\begin{lstlisting}
#include <concepts>

template<std::integral T>
class Counter {
    T value_{};
public:
    void add(T delta) { value_ += delta; }
    T value() const { return value_; }
};

// Counter<int> a;       // OK
// Counter<double> b;    // ERROR: constraint not satisfied
\end{lstlisting}

Ograniczenia mogą też sterować wyborem częściowej specjalizacji.

\begin{lstlisting}
#include <concepts>
#include <string_view>

template<class T>
struct TypeLabel {
    static constexpr std::string_view value = "other";
};

template<std::integral T>
struct TypeLabel<T> {
    static constexpr std::string_view value = "integer";
};

template<std::floating_point T>
struct TypeLabel<T> {
    static constexpr std::string_view value = "floating";
};

static_assert(TypeLabel<int>::value == "integer");
static_assert(TypeLabel<double>::value == "floating");
\end{lstlisting}

Nie używaj jednak złożonego systemu specjalizacji, jeżeli
prosta funkcja przeciążona lub \texttt{if constexpr}
lepiej wyraża zamysł.

\section{Koncepcje i pakiety parametrów: kontynuacja rozdziału 4}
\label{sec:concepts-variadic}

Koncepcje bardzo dobrze współpracują z \emph{variadic templates}.
Możemy nakładać warunek na każdy element pakietu typów
za pomocą \emph{fold expression}.

\begin{lstlisting}
#include <concepts>

template<class... Ts>
requires (std::integral<Ts> && ...)
constexpr auto sum_integers(Ts... xs) {
    return (0 + ... + xs);
}

static_assert(sum_integers() == 0);
static_assert(sum_integers(1, 2, 3) == 6);
// sum_integers(1, 2.5);    // ERROR: double not integral
\end{lstlisting}

Zapis \texttt{(P<Ts> \&\& ...)} reprezentuje
koniunkcję wymagań dla wszystkich elementów pakietu.
Dla pustego pakietu taki fold logiczny daje \texttt{true};
nie oznacza to jednak, że każde ciało funkcji poradzi sobie
z zerową liczbą argumentów. W przykładzie użyliśmy
\texttt{0} jako elementu neutralnego sumowania.

\subsection{Ograniczony pakiet typów i uniwersalne przekazywanie}

\begin{lstlisting}
#include <concepts>
#include <utility>

template<std::integral... Ts>
void accept_integers(Ts&&... xs) {
    ((void)std::forward<Ts>(xs), ...); // illustrative use
}
\end{lstlisting}

Zapis \texttt{std::integral... Ts} również ogranicza wszystkie
elementy pakietu. Uwaga: przy dedukcji z przekazywanymi referencjami
\texttt{Ts\&\&}, typ \texttt{Ts} może być \emph{referencją},
gdy argument jest lvalue. W takim przypadku
\texttt{std::integral<Ts>} może być fałszywe.
Jeżeli chcemy akceptować obie kategorie wartości, często
właściwsze jest użycie \texttt{std::remove\_cvref\_t}:

\begin{lstlisting}
#include <concepts>
#include <type_traits>
#include <utility>

template<class... Ts>
requires (std::integral<std::remove_cvref_t<Ts>> && ...)
void accept_any_int(Ts&&... xs) {
    ((void)xs, ...); // demonstration; no ownership transfer
}
\end{lstlisting}

Ta różnica jest szczególnie ważna przy projektowaniu generycznych
wrapperów, które mają przyjmować także obiekty \texttt{const}
i referencje.

\section{Concepts a \texttt{if constexpr}, \texttt{static\_assert} i SFINAE}
\label{sec:concepts-porownanie}

Są to narzędzia pokrewne, ale mają różne cele.

\subsection{\texttt{if constexpr}: wybór implementacji}

\begin{lstlisting}
#include <concepts>
#include <string>
#include <string_view>

template<class T>
std::string describe(const T& value) {
    if constexpr (std::integral<T>) {
        return "an integer";
    } else if constexpr (std::floating_point<T>) {
        return "a floating-point value";
    } else {
        return "another type";
    }
}
\end{lstlisting}

\texttt{if constexpr} umożliwia wybranie gałęzi kodu
\emph{wewnątrz} szablonu. Nie tworzy samo przez się
ograniczonego interfejsu. Z kolei \texttt{requires}
może odrzucić niewłaściwe wywołanie jeszcze na etapie
dobierania przeciążeń.

\subsection{\texttt{static\_assert}: własny komunikat}

\begin{lstlisting}
#include <concepts>

template<class T>
void save_number(const T& value) {
    static_assert(std::integral<T>,
                  "save_number expects an integral type");
    // ...
}
\end{lstlisting}

\texttt{static\_assert} jest użyteczne, gdy potrzebujemy
jednoznacznego komunikatu albo sprawdzamy niezmiennik wewnątrz
implementacji. Nie zastępuje jednak ograniczenia uczestniczącego
w rozstrzyganiu przeciążeń.

\subsection{Przed C++20: \texttt{std::enable\_if} i idiom wykrywania}

\begin{lstlisting}
#include <type_traits>

// C++11/14/17 technique:
template<class T,
         std::enable_if_t<std::is_integral_v<T>, int> = 0>
T old_square(T x) {
    return x * x;
}

// C++20:
#include <concepts>
template<std::integral T>
T new_square(T x) {
    return x * x;
}
\end{lstlisting}

Dla nowego kodu C++20 najczęściej korzystniej wybrać
koncepcję: intencja staje się czytelna i spójna z mechanizmem
porządkowania przeciążeń. \texttt{std::enable\_if}
nadal ma znaczenie w bibliotekach wspierających starsze
wersje standardu.

\section{Praktyczny przykład: średnia z zakresu liczb}
\label{sec:concepts-praktyka}

Zamiast pisać funkcję tylko dla \texttt{std::vector<double>},
możemy określić dokładne wymagania wobec \emph{zakresu}.

\begin{lstlisting}
#include <concepts>
#include <cstddef>
#include <optional>
#include <ranges>
#include <vector>
#include <array>

template<class R>
concept NumericInputRange =
    std::ranges::input_range<R> &&
    std::convertible_to<std::ranges::range_reference_t<R>, double>;

template<NumericInputRange R>
std::optional<double> average(R&& range) {
    double sum = 0.0;
    std::size_t count = 0;
    for (auto&& item : range) {
        sum += static_cast<double>(item);
        ++count;
    }
    if (count == 0) return std::nullopt;
    return sum / static_cast<double>(count);
}

int main() {
    std::vector<int> v{2, 4, 6};
    std::array<double, 2> a{3.0, 7.0};
    auto x = average(v); // 4.0
    auto y = average(a); // 5.0
    auto z = average(std::vector<int>{}); // empty optional
    return !(x && y && !z && *x == 4.0 && *y == 5.0);
}
\end{lstlisting}

Funkcja działa dla różnych kontenerów i zakresów dostarczających
iterację jednokierunkową. Warunek \texttt{input\_range}
wyraża strukturę danych, a \texttt{convertible\_to}
-- oczekiwany rodzaj elementów. Zastosowaliśmy
\texttt{std::optional<double>}, aby pusty zakres nie powodował
dzielenia przez zero. Dla zastosowań numerycznych należy
dodatkowo rozważyć stabilność sumowania, przepełnienia i
znaczenie konwersji do \texttt{double}.

\paragraph{Kiedy \texttt{ranges}?}
Koncepcje standardowych zakresów są szczególnie użyteczne,
gdy funkcja operuje na sekwencji, ale nie wymaga konkretnie
\texttt{std::vector} czy \texttt{std::list}. Zamiast wymuszać
silniejsze wymaganie \texttt{random\_access\_range},
wybieraj najsłabszą koncepcję, która wystarcza do działania
algorytmu.

\section{Koncepcje dla lambd i algorytmów}

Funkcje anonimowe również mogą mieć ograniczone parametry.

\begin{lstlisting}
#include <concepts>

int main() {
    auto twice = [](std::integral auto x) { return x + x; };
    return twice(5) != 10;
}
\end{lstlisting}

Można też tworzyć generyczne algorytmy przyjmujące funkcje
spełniające \texttt{std::predicate}.

\begin{lstlisting}
#include <concepts>
#include <cstddef>
#include <vector>

template<class T, class Pred>
requires std::predicate<Pred&, const T&>
std::size_t count_if_custom(const std::vector<T>& values,
                            Pred pred) {
    std::size_t count = 0;
    for (const T& x : values)
        if (pred(x)) ++count;
    return count;
}
\end{lstlisting}

Zwróć uwagę na argumenty typu w
\texttt{std::predicate<Pred\&, const T\&>}:
sprawdzają możliwość wywołania \emph{konkretnego rodzaju}
predykatu z \emph{konkretnym rodzajem} argumentu.
Kategoria wartości ma znaczenie i nie należy bezrefleksyjnie
wstawiać tu samych nazw typów.

\section{Granice gwarancji: najczęstsze pułapki}
\label{sec:concepts-pulapki}

\subsection{Concept nie sprawdza dowolnej własności semantycznej}

\begin{lstlisting}
#include <concepts>

struct Strange {
    int value{};
    bool operator==(const Strange&) const { return false; }
};

static_assert(std::equality_comparable<Strange>);
// Syntactically valid, but equality lacks reflexivity.
\end{lstlisting}

Kompilator akceptuje składnię porównania, lecz
\texttt{a == a} jest zawsze fałszywe. Kod nie spełnia
matematycznych oczekiwań związanych z relacją równości.
To obowiązek autora typu, niezależny od samego mechanizmu
\texttt{concepts}.

\subsection{Koncepcje nie zamieniają błędów w poprawny kod}

Wyrażenia \texttt{requires} pozwalają w kontekście zależnym
od parametru szablonu traktować niepoprawne podstawienia
jako niespełnione wymagania. Nie oznacza to jednak, że
\emph{każdy} błędny fragment C++ staje się automatycznie
\texttt{false}. Niezależne od parametru, trwale niepoprawne
konstrukcje mogą prowadzić do błędu kompilacji.

\subsection{Różnica między \texttt{T}, \texttt{T\&} i \texttt{const T\&}}

\begin{lstlisting}
#include <concepts>
#include <type_traits>

template<class T>
concept PlainInteger =
    std::integral<std::remove_cvref_t<T>>;

static_assert(!std::integral<int&>);
static_assert(PlainInteger<int&>);
static_assert(PlainInteger<const int&>);
\end{lstlisting}

Przy dedukcji typu generycznego parametru
\texttt{T\&\&} ta różnica często decyduje o tym,
czy z pozoru poprawne wywołanie spełni koncepcję.

\subsection{\texttt{std::same\_as} czy \texttt{std::convertible\_to}?}

\begin{lstlisting}
#include <concepts>

template<class T>
concept StrictAdd = requires(T a, T b) {
    { a + b } -> std::same_as<T>;
};

template<class T>
concept FlexibleAdd = requires(T a, T b) {
    { a + b } -> std::convertible_to<T>;
};

static_assert(!StrictAdd<short>);  // short + short -> int
static_assert(FlexibleAdd<short>); // int -> short allowed
\end{lstlisting}

Koncepcja \texttt{FlexibleAdd} nie oznacza \emph{bezstratnej}
konwersji; \texttt{int} można niejawnie przekształcić do
\texttt{short}, tracąc informację. Nazwa koncepcji powinna
odzwierciedlać rzeczywistą, a nie postulowaną gwarancję.

\subsection{Funkcja ograniczona nie zawsze jest optymalniejsza}

Koncepcje opisują poprawność i wybór instancji szablonów.
Nie wymuszają \emph{inline}, wektoryzacji, eliminacji kopii
ani lepszego kodu maszynowego. Można je stosować bez obietnicy
przyspieszenia wykonywania programu.

\section{Zasady projektowania: kiedy używać, a kiedy nie?}

\begin{center}
\small
\begin{tabular}{|p{0.34\textwidth}|p{0.55\textwidth}|}
\hline
\textbf{Sytuacja} & \textbf{Zalecane rozwiązanie} \\
\hline
Funkcja generyczna wymaga konkretnej kategorii typów & Nazwany \texttt{concept}, np. \texttt{std::integral} \\
\hline
Wymagania dotyczą istnienia operacji lub typu składowego & Własny \texttt{requires-expression} \\
\hline
Ważne są różne wersje algorytmu dla różnych typów & Ograniczone przeciążenia lub specjalizacje \\
\hline
Jedna implementacja ma kilka gałęzi zależnych od typu & \texttt{if constexpr} \\
\hline
Sprawdzasz założenie wewnątrz ciała funkcji & \texttt{static\_assert} z komunikatem \\
\hline
Weryfikujesz jednakowy warunek dla całego pakietu & Fold \texttt{(Concept<Ts> \&\& ...)} \\
\hline
Potrzebujesz iteracji po sekwencji & \texttt{std::ranges::input\_range} lub silniejszy koncept \\
\hline
Funkcja działa dla jednego, konkretnego typu & Zwykła funkcja często jest prostsza niż szablon \\
\hline
Projekt musi działać także w C++17 & \texttt{type\_traits} i (w razie potrzeby) SFINAE \\
\hline
Wymagasz warunku wykonywania programu, np. \texttt{n>0} & Zwykłe sprawdzenie w czasie działania; koncept go nie zastępuje \\
\hline
\end{tabular}
\end{center}

\section{Podsumowanie i pytania kontrolne}

\begin{enumerate}
\item \texttt{concept} nazywa predykat nad parametrami szablonu,
sprawdzany w czasie kompilacji.
\item \texttt{requires-clause} ogranicza deklarację;
\texttt{requires-expression} jest wyrażeniem badającym wymagania.
\item Istnieją cztery formy wymagań: proste, typu, złożone
i zagnieżdżone.
\item \texttt{std::same\_as}, \texttt{std::convertible\_to},
\texttt{std::integral}, \texttt{std::predicate} i koncepcje zakresów
powinny być zwykle pierwszym wyborem.
\item Ograniczenia mogą wpływać na wybór przeciążenia, ale mechanizm
porządkowania nie jest dowolnym dowodzeniem implikacji logicznej.
\item Nazwane koncepcje warto wykorzystywać ponownie przy
budowaniu koncepcji bardziej szczegółowych.
\item \texttt{std::remove\_cvref\_t<T>} ma często znaczenie
przy koncepcjach typów dedukowanych z referencji przekazujących.
\item Koncepcje opisują przede wszystkim wymagania składniowe,
a dodatkowe prawa semantyczne muszą respektować autorzy typów.
\end{enumerate}

\paragraph{Pytania do samodzielnego sprawdzenia.}
Zastanów się, co kompilator zaakceptuje i dlaczego:

\begin{lstlisting}
#include <concepts>
#include <type_traits>

static_assert(std::integral<bool>);
static_assert(!std::integral<int&>);
static_assert(std::integral<std::remove_cvref_t<const int&>>);
static_assert(!std::same_as<int, long>);

// 1. Why does StrictAdd<short> fail?
// 2. Why can a concept be true for a logically broken ==?
// 3. How do requires-clause and requires-expression differ?
// 4. Which overload of category(5) is chosen?
// 5. Why is NumericInputRange more general than vector<double>?
\end{lstlisting}

\paragraph{Źródła i dalsza lektura.}
Oficjalne reguły języka określa standard ISO C++, a przystępne
opisy i przykłady znajdują się w serwisie \emph{cppreference}.
Przydatne adresy:
\begin{lstlisting}
en.cppreference.com/cpp/language/constraints
en.cppreference.com/cpp/language/requires
en.cppreference.com/cpp/concepts
\end{lstlisting}

% ROZDZIAL 8 — Iteratory, ranges i views (C++20)
% Fragment do włączenia po rozdziale 7.
% Wymagane pakiety: listings, array, amsmath (opcjonalnie), hyperref.
% Użycie: \input{rozdzial_8_iteratory_ranges_views_cpp.tex}

\chapter[Iteratory, Ranges i Views]{Iteratory, zakresy i widoki: \textit{Ranges} oraz \textit{Views} w C++20}
\label{chap:ranges-views-cpp}

\section{Dlaczego wprowadzono Ranges?}

Klasyczne algorytmy STL operują na \emph{parze iteratorów}: pierwszym
wskazującym początek danych i drugim wyznaczającym koniec. Dzięki temu
\texttt{std::sort} lub \texttt{std::find} mogą współpracować z wieloma
kontenerami. Wymaga to jednak przekazywania obu granic i dodatkowej
ostrożności, aby pochodziły z tego samego zakresu.

\begin{lstlisting}
#include <algorithm>
#include <vector>

std::vector<int> xs{4, 1, 3, 2};
std::sort(xs.begin(), xs.end()); // STL before C++20
std::ranges::sort(xs);          // C++20
\end{lstlisting}

Nowy model pozwala przekazać \emph{zakres} (\textit{range}) jako jeden
obiekt oraz formułować wymagania wobec niego poprzez \emph{koncepcje}
(\textit{concepts}, rozdział~5). Oddzielną innowacją są \emph{widoki}
(\textit{views}): lekkie reprezentacje danych, często przekształcane
\textbf{leniwe}, to znaczy dopiero podczas ich przeglądania.

\begin{lstlisting}
#include <ranges>

std::vector<int> xs{1, 2, 3, 4, 5, 6};
auto pipeline = xs
    | std::views::filter([](int x) { return x % 2 == 0; })
    | std::views::transform([](int x) { return x * x; });

// No vector of squares has been created.
// Iteration produces 4, 16, 36.
for (int x : pipeline) { /* use x */ }
\end{lstlisting}

\paragraph{Założenia.}
W całym rozdziale przyjmujemy standard \textbf{C++20} oraz nagłówki
\texttt{<algorithm>}, \texttt{<iterator>}, \texttt{<ranges>},
\texttt{<concepts>}, \texttt{<vector>} i pozostałe nagłówki wymienione
w przykładach. W dalszych sekcjach posłużymy się semantyką przenoszenia
(rozdział~6) i lambdami (rozdział~7).

\section{Iteratory: uogólnienie wskaźników}
\label{sec:cpp8-iterators}

\subsection{Operacje \texttt{begin}, \texttt{end}, \texttt{*} i \texttt{++}}

Iterator jest obiektem umożliwiającym poruszanie się po elementach
zakresu. Wskaźnik \texttt{int*} jest przykładem iteratora, ale iterator
może także przechowywać indeks, wskaźnik do węzła listy lub stan
wewnętrzny generatora. Podstawowe operacje przedstawia program:

\begin{lstlisting}
#include <vector>

std::vector<int> v{10, 20, 30};
auto it = v.begin();     // points to 10
int a = *it;            // 10
++it;                   // points to 20
int b = *it;            // 20
bool ok = it != v.end();
\end{lstlisting}

Iterator \texttt{end()} zwykle oznacza pozycję \emph{za ostatnim}
elementem; nie wolno go dereferencjonować. Dzięki temu pusty zakres
spełnia \texttt{begin()==end()}, a przechodzenie po elementach ma
jednolity schemat:

\begin{lstlisting}
for (auto it = v.begin(); it != v.end(); ++it) {
    // *it is a valid element
}
\end{lstlisting}

\paragraph{Pętla zakresowa.}
Zapis \texttt{for (auto\& x : v)} jest wygodniejszą postacią iteracji,
ale w dalszym ciągu wykorzystuje operacje początku, końca,
porównania, zwiększania i dereferencji. Aby \emph{modyfikować} elementy,
stosujemy \texttt{auto\&}; aby uniknąć kopiowania bez modyfikowania,
\texttt{const auto\&}; dla niewielkich wartości liczbowych często
wystarczy \texttt{auto}.

\subsection{Iteratory stałe i nieważność iteratorów}

\begin{lstlisting}
std::vector<int> v{1, 2, 3};
auto m = v.begin();   // can modify *m
auto c = v.cbegin();  // cannot modify *c
*m = 10;
// *c = 20;         // compilation error
\end{lstlisting}

\texttt{const\_iterator} chroni elementy przed modyfikacją poprzez
ten iterator; nie oznacza, że cały kontener jest niezmienny. Równie
ważna jest \emph{nieważność iteratorów} (\textit{iterator invalidation}).
Na przykład ponowna alokacja pamięci przez \texttt{std::vector::push\_back}
może unieważnić \textbf{wszystkie} iteratory, wskaźniki i referencje do
jego elementów.

\begin{lstlisting}
std::vector<int> v{1, 2, 3};
auto it = v.begin();
v.push_back(4);       // may reallocate storage
// *it;              // unsafe if reallocation happened
it = v.begin();      // obtain a fresh iterator
\end{lstlisting}

Reguły zależą od kontenera i operacji. W \texttt{std::list} wstawienie
elementu nie unieważnia iteratorów do pozostałych elementów, ale nie
oznacza to, że można dalej używać iteratora do elementu usuniętego.

\section{Kategorie iteratorów i ich wymagania}
\label{sec:cpp8-iterator-categories}

Zestaw koncepcji C++20 opisuje operacje i własności iteratorów, a nie
jedynie nazwy klas. Najważniejsze kategorie zestawia tabela.

\begin{center}
\small
\begin{tabular}{|p{5.2cm}|p{7.3cm}|}
\hline
\textbf{Koncepcja} & \textbf{Możliwości i przykłady} \\
\hline
\texttt{input\_iterator} & Odczytywanie kolejnych wartości; przeglądanie może być jednokrotne (np. strumień wejściowy). \\
\hline
\texttt{output\_iterator} & Zapisywanie kolejnych wartości do celu (np. \texttt{back\_insert\_iterator}). \\
\hline
\texttt{forward\_iterator} & Wielokrotne przejścia, porównywanie iteratorów, krok do przodu. \\
\hline
\texttt{bidirectional\_iterator} & Dodatkowo krok wstecz, czyli \texttt{--it}; np. \texttt{list}. \\
\hline
\texttt{random\_access\_iterator} & Indeksowanie i skoki \texttt{it+n} w stałym czasie; np. \texttt{deque}. \\
\hline
\texttt{contiguous\_iterator} & Elementy leżą kolejno w pamięci; np. iterator \texttt{vector<int>} i surowy wskaźnik. \\
\hline
\end{tabular}
\end{center}

Hierarchia wymagań dla iteratorów odczytujących jest w przybliżeniu
następująca (każda kolejna koncepcja jest silniejsza):

\begin{center}
\texttt{input $\Rightarrow$ forward $\Rightarrow$ bidirectional}
\\[2pt]
\texttt{$\Rightarrow$ random\_access $\Rightarrow$ contiguous}.
\end{center}

Koncepcja \texttt{output\_iterator} dotyczy osobnej rodziny wymagań.
Nowe koncepcje nie są w każdym szczególe identyczne ze starymi
kategoriami \emph{LegacyIterator}. Na przykład model
\texttt{forward\_iterator} nie wymaga zwracania referencji
przy dereferencji.

\begin{lstlisting}
#include <concepts>
#include <iterator>
#include <list>
#include <vector>

static_assert(std::contiguous_iterator<int*>);
static_assert(std::random_access_iterator<
    std::vector<int>::iterator>);
static_assert(std::bidirectional_iterator<
    std::list<int>::iterator>);
static_assert(!std::random_access_iterator<
    std::list<int>::iterator>);
\end{lstlisting}

\paragraph{Dlaczego to ważne?}
\texttt{std::ranges::sort} wymaga iteratorów o dostępie swobodnym,
więc można użyć go do \texttt{std::vector}, ale \textbf{nie} do
\texttt{std::list}. Listę sortuje się metodą \texttt{list::sort()}.
Wybór algorytmu jest zatem ograniczony kategorią iteratora.

\section{Własny iterator i własny zakres}
\label{sec:cpp8-own-iterator}

Zbudujemy iterator generujący kolejne liczby całkowite. Nie będzie on
przechowywać ani wskaźnika, ani kontenera: całą potrzebną informacją
jest bieżąca wartość. Zdefiniujemy \emph{iterator do przodu}
(\texttt{forward\_iterator}).

\begin{lstlisting}
#include <concepts>
#include <cstddef>
#include <iterator>
#include <ranges>

struct CountingIterator {
    using value_type = int;
    using difference_type = std::ptrdiff_t;
    using iterator_concept = std::forward_iterator_tag;

    int current = 0;

    int operator*() const { return current; }

    CountingIterator& operator++() {
        ++current;
        return *this;
    }
    CountingIterator operator++(int) {
        auto old = *this;
        ++*this;
        return old;
    }
    bool operator==(const CountingIterator&) const = default;
};

struct CountingRange {
    int first = 0;
    int last = 0;

    CountingIterator begin() const { return {first}; }
    CountingIterator end() const { return {last}; }
};

static_assert(std::forward_iterator<CountingIterator>);
static_assert(std::ranges::forward_range<CountingRange>);
\end{lstlisting}

Teraz kod \texttt{for (int i : CountingRange\{3, 7\})} wypisze
\texttt{3, 4, 5, 6}. Funkcje \texttt{begin()} i \texttt{end()}
wyznaczają \emph{półotwarty} zakres \([3,7)\). Prawidłowe użycie
zakłada \texttt{first <= last}; dla \texttt{first > last} nasz
minimalny iterator nie zbiegnie do wartości końcowej.

\paragraph{Co trzeba zagwarantować?}
Samo przejście \texttt{static\_assert} nie dowodzi wszystkich
semantycznych własności iteratora. Dla \texttt{forward\_iterator}
ważna jest m.in. własność \emph{multi-pass}: kopiowanie iteratora
i przejście jedną kopią nie powinno zmieniać wartości zwracanych
przez drugą. Nasz iterator ją spełnia, ponieważ każda kopia ma własne
pole \texttt{current}.

\section{Zwykłe \textit{range} i \textit{sentinel}}
\label{sec:cpp8-range-sentinel}

\emph{Zakres} jest obiektem, dla którego można otrzymać początek
\texttt{std::ranges::begin(r)} i koniec \texttt{std::ranges::end(r)}.
Typ wyniku obu wywołań \textbf{nie musi być identyczny}. Obiekt
wyznaczający koniec to \emph{sentinel}; musi dać się porównać z
iteratorem, ale sam nie musi obsługiwać jego pozostałych operacji.

\begin{lstlisting}
#include <ranges>
#include <vector>

std::vector<int> xs{1, 2, 3};
auto first = std::ranges::begin(xs);
auto last = std::ranges::end(xs);
while (first != last) { ++first; }
\end{lstlisting}

W tym przypadku obie granice mają ten sam typ: zakres spełnia
koncepcję \texttt{common\_range}. W innych zakresach koniec może
wyznaczać osobny obiekt sentinela. Odpowiednia koncepcja
\texttt{std::sentinel\_for} określa warunki porównywania granic:

\begin{lstlisting}
static_assert(std::ranges::common_range<decltype(xs)>);
static_assert(std::sentinel_for<
    decltype(last), decltype(first)>);
\end{lstlisting}

Przykładem osobnego typu końca jest \texttt{std::default\_sentinel\_t}.

\paragraph{Rozmiar a liczba kroków.}
Operacja \texttt{std::ranges::size(r)} jest dostępna dla zakresów
spełniających \texttt{sized\_range}. Nie każdy zakres ma tanią operację
wyznaczenia rozmiaru. Funkcja \texttt{distance(r)} z przestrzeni
\texttt{std::ranges} może wymagać
\emph{przejścia przez wszystkie elementy}. Dla niektórych
jednokrotnych źródeł danych przejście oznacza również ich zużycie.

\section{Koncepcje zakresów w praktyce}

Własności zakresów sprawdzamy za pomocą koncepcji z rozdziału~5.
Najważniejsze z nich to:
\begin{itemize}
\item \texttt{input\_range} i \texttt{forward\_range},
\item \texttt{bidirectional\_range} i \texttt{random\_access\_range},
\item \texttt{contiguous\_range}, \texttt{common\_range} i \texttt{sized\_range}.
\end{itemize}

\begin{lstlisting}
#include <ranges>
#include <vector>

template <std::ranges::input_range R>
long long sum_integers(R&& range) {
    long long result = 0;
    for (auto value : range) result += value;
    return result;
}

std::vector<int> v{1, 2, 3};
long long s = sum_integers(v); // 6
\end{lstlisting}

Parametr \texttt{R\&\&} jest tutaj \emph{referencją przekazującą}
(\textit{forwarding reference}), a nie żądaniem przejęcia własności.
Funkcja obsługuje także źródła jednokrotne, gdyż przechodzi zakres tylko
raz. Dla naprawdę uniwersalnej funkcji trzeba osobno wyrazić warunki
na typ elementu i wykonywaną na nim operację dodawania.

\section{Algorytmy \texttt{std::ranges} i projekcje}
\label{sec:cpp8-algorithms}

Funkcje z rodziny \texttt{std::ranges}, np. \texttt{sort},
\texttt{find}, \texttt{copy} i \texttt{transform}, przyjmują
bezpośrednio zakres albo parę iterator--sentinel.
Ich wykonanie jest na ogół \emph{natychmiastowe} (w odróżnieniu od
samych widoków \texttt{views::transform}).

\begin{lstlisting}
#include <algorithm>
#include <iterator>
#include <vector>

std::vector<int> v{4, 2, 3, 1};
std::ranges::sort(v);  // modifies v immediately

auto pos = std::ranges::find(v, 3);
if (pos != v.end()) { *pos = 30; }

std::vector<int> out;
std::ranges::copy(v, std::back_inserter(out));
\end{lstlisting}

\subsection{Projekcja: sortowanie po wskazanym polu}

Algorytmy Ranges często posiadają argument \emph{projekcji}: funkcję
wywoływaną na elemencie przed porównaniem albo wyszukiwaniem. Dzięki
temu zamiast pisać własny komparator, można wskazać pole struktury.

\begin{lstlisting}
#include <string>
#include <vector>

struct Person {
    std::string name;
    int age;
};

std::vector<Person> people{
    {"Ada", 40}, {"Ewa", 25}, {"Jan", 32}
};

std::ranges::sort(people, {}, &Person::age);
// ascending ages: 25, 32, 40

auto it = std::ranges::find(people, 32, &Person::age);
// it->name == "Jan"
\end{lstlisting}

W wywołaniu \texttt{sort(people, \{\}, \&Person::age)} puste
nawiasy oznaczają domyślny komparator \texttt{std::ranges::less},
a wskaźnik do składowej to projekcja. Standard używa mechanizmu
\texttt{std::invoke}, omówionego w rozdziale~7.

\section{Widoki: \textit{view} nie jest synonimem kontenera}
\label{sec:cpp8-views}

\emph{Widok} jest zakresem przystosowanym do taniego tworzenia i
kopiowania lub przenoszenia \emph{samego obiektu widoku}. Nie musi mieć
własnych danych: może przechowywać referencję do kontenera, mały zestaw
iteratorów albo parametry generowania wartości.

\begin{lstlisting}
#include <ranges>
#include <vector>

std::vector<int> xs{10, 20, 30, 40};
auto first_two = xs | std::views::take(2);
// first_two refers to the original xs
\end{lstlisting}

Stworzenie widoku nie jest równoznaczne ze skopiowaniem elementów do
nowego kontenera. Odwołanie do danych źródłowych wymaga, by pozostawały
one ważne przez cały czas korzystania z widoku. Szczegółowe reguły
własności omówimy dalej.

\paragraph{Dwa znaczenia nazwy \texttt{views}.}
\texttt{std::views} jest aliasem przestrzeni nazw z obiektami
adaptującymi zakresy. \texttt{std::ranges::view} to natomiast
\emph{koncepcja} spełniana przez dany typ. Nie należy mylić ich
z pojęciem \texttt{std::ranges::range}.

\section{Leniwe potoki \texttt{filter}, \texttt{transform}, \texttt{take}}
\label{sec:cpp8-pipelines}

Najczęściej stosowane widoki łączymy za pomocą operatora \texttt{|}.
Adaptory przekazują sobie zakresy bez potrzeby budowania pośrednich
wektorów:

\begin{lstlisting}
std::vector<int> xs{1,2,3,4,5,6,7,8,9,10};

auto pipeline = xs
    | std::views::filter([](int x) {
          return x % 2 == 0;
      })
    | std::views::transform([](int x) {
          return x * x;
      })
    | std::views::take(3);

// Iteration produces 4, 16, 36.
\end{lstlisting}

\texttt{filter} pomija elementy niespełniające predykatu,
\texttt{transform} wyznacza nową wartość przy dereferencji,
a \texttt{take(3)} ogranicza liczbę wyników. Algorytm
\texttt{std::ranges::transform} działa od razu; natomiast
utworzenie widoku \texttt{views::transform} nie przelicza
wszystkich wyników.

\subsection{Kiedy następują obliczenia?}

\begin{lstlisting}
int calls = 0;
auto squares = xs | std::views::transform([&](int x) {
    ++calls;
    return x * x;
});
// calls == 0 here
int first = *squares.begin();
// calls has increased (transform was evaluated)
\end{lstlisting}

W tym przykładzie pokazujemy moment obliczenia, a nie zalecany styl
projektowania: \textbf{predykaty i transformacje powinny zasadniczo
unikać efektów ubocznych}. Ponowne przejście widoku, ponowna
dereferencja lub implementacyjnie dopuszczone buforowanie może
powodować inne liczby wywołań, niż sugerowałaby pojedyncza pętla.
Nie należy projektować logiki programu na podstawie rzekomej liczby
wywołań funkcji przez widok.

\paragraph{Ważny szczegół \texttt{filter\_view}.}
Dla pewnych kategorii zakresu \texttt{filter\_view} może zapamiętywać
pierwszy znaleziony element. Ponadto modyfikowanie elementu tak,
aby przestał spełniać predykat filtrowania, narusza wymagania
\texttt{filter\_view} i może prowadzić do niezdefiniowanego
zachowania. Predykat powinien być stabilny wobec przebiegu iteracji.

\section{Pozostałe przydatne adaptory C++20}
\label{sec:cpp8-adaptors}

\subsection{\texttt{iota}: wartości wyznaczane bez tablicy}

\begin{lstlisting}
auto ids = std::views::iota(1, 5);
// Iteration produces 1, 2, 3, 4.

// This is an unbounded range, limited by take(5):
auto first_five = std::views::iota(1)
                | std::views::take(5);
// Produces 1, 2, 3, 4, 5.
\end{lstlisting}

\texttt{iota} przydaje się do indeksowania lub generowania sekwencji.
Dla nieograniczonego widoku trzeba pamiętać o warunku zakończenia;
ponadto operacje arytmetyczne nie mogą przepełniać typu całkowitego
ze znakiem.

\subsection{\texttt{drop}, \texttt{take} i \texttt{reverse}}

\begin{lstlisting}
std::vector<int> data{1, 2, 3, 4, 5};
auto middle = data | std::views::drop(1)
                   | std::views::take(3);
// 2, 3, 4

auto backwards = data | std::views::reverse;
// 5, 4, 3, 2, 1
\end{lstlisting}

\texttt{reverse} wymaga co najmniej zakresu dwukierunkowego;
\texttt{take} i \texttt{drop} są przydatne także dla sekwencji
generowanych leniwie. Dla niektórych widoków koszt \texttt{begin()}
lub dojścia do pierwszego wyniku może być większy niż stały.

\subsection{\texttt{keys}, \texttt{values} i \texttt{split}}

\begin{lstlisting}
std::vector<std::pair<int, std::string>> pairs{
    {1, "one"}, {2, "two"}
};
for (int key : pairs | std::views::keys) { /* 1, 2 */ }
for (const auto& value : pairs | std::views::values) {
    /* "one", "two" */
}

std::string csv = "aa,b,cc";
for (auto part : csv | std::views::split(',')) {
    std::string piece(part.begin(), part.end());
    // Copy each part into an owning string if needed.
}
\end{lstlisting}

\texttt{split} nie tworzy automatycznie \texttt{std::string} dla
każdego fragmentu; otrzymujemy zakresy wskazujące fragmenty źródła.
Zachowując wynik na później, trzeba zapewnić odpowiedni czas życia
źródła. Nie należy też zakładać, że splitowanie tekstu CSV obsługuje
cudzysłowy i escapowanie pól zgodnie z pełnym formatem CSV.

\paragraph{Uwaga na wersję standardu.}
\texttt{std::views::zip} i \texttt{std::views::enumerate} są
konstrukcjami \textbf{C++23}, podobnie jak \texttt{std::ranges::to}.
Nie są dostępne w czystym C++20. W standardzie C++20 wynik można zebrać do kontenera np.
algorytmem \texttt{ranges::copy} z iteratorem wstawiającym
\texttt{std::back\_inserter}.

\section{Materializacja wyniku i koszt przetwarzania}
\label{sec:cpp8-materialize}

Jeśli wynik ma być niezależny od źródła i przechowywany na później,
możemy zamienić potok w kontener. W C++20 nie ma standardowego
\texttt{std::ranges::to<std::vector>()}.

\begin{lstlisting}
#include <algorithm>
#include <iterator>
#include <vector>

std::vector<int> data{1,2,3,4,5,6};
auto even = data | std::views::filter([](int x) {
    return x % 2 == 0;
});

std::vector<int> result;
std::ranges::copy(even, std::back_inserter(result));
// result == {2, 4, 6}
\end{lstlisting}

Taka materializacja wymaga pamięci proporcjonalnej do wielkości wyniku.
Zaletą potoku jest unikanie pośrednich alokacji i często możliwość
\emph{łączenia} transformacji w jednym przebiegu. Nie jest to jednak
gwarancja, że widoki zawsze będą szybsze: koszty iteracji, predykatów,
przepływu sterowania i optymalizacji zależą od konkretnego kodu i
kompilatora. W zadaniach krytycznych wykonuj pomiary.

\section{Własność danych, czas życia i \texttt{borrowed\_range}}
\label{sec:cpp8-lifetime}

To najważniejsza część praktycznego korzystania z Ranges. Widoki
\textbf{nie zawsze posiadają} bazowe dane. W przypadku
\texttt{std::views::all} przeprowadzamy istotne rozróżnienie:

\begin{lstlisting}
std::vector<int> source{2, 4, 6};
auto borrowed = std::views::all(source);
// Refers to source: source must stay alive.

auto owned = std::views::all(std::vector<int>{3, 5, 7});
// In C++20, this can be an owning_view.
// The temporary vector is moved into the view.
\end{lstlisting}

Nie oznacza to, że \emph{każdy} widok utworzony z wartości tymczasowej
przejmuje własność, ani że wszystkie widoki zachowują się tak samo.
Różne adaptory przechowują różne rodzaje obiektów bazowych.

\subsection{Wskaźniki wiszące wewnątrz potoku}

\begin{lstlisting}
auto make_bad_view() {
    std::vector<int> local{1, 2, 3};
    return local | std::views::filter([](int x) {
        return x > 1;
    });
}
// Using the returned view would access destroyed local.
\end{lstlisting}

Ponieważ \texttt{local} jest \emph{lvalue}, adapter przechowuje
nieposiadający danych widok do lokalnego kontenera. Po wyjściu z
funkcji \texttt{local} zostaje zniszczony, a zwrócony widok zawiera
nieaktualne odwołanie. Bezpieczne wyjścia to utworzenie widoku z
przekazanego na własność kontenera (np. obiektu tymczasowego lub
\texttt{std::move(local)}) albo zwrócenie zwykłego kontenera wynikowego.

\begin{lstlisting}
auto make_owning_view() {
    std::vector<int> local{1, 2, 3};
    return std::views::all(std::move(local))
         | std::views::filter([](int x) { return x > 1; });
}
// The returned view owns its vector through owning_view.
\end{lstlisting}

\subsection{Co naprawdę mówi \texttt{borrowed\_range}?}

Koncepcja \texttt{borrowed\_range<R>} oznacza w uproszczeniu, że
iteratory otrzymane z zakresu typu \texttt{R} mogą pozostać ważne
po zniszczeniu \emph{samego obiektu zakresu} (o ile dane, do których
odsyłają, nadal istnieją). Nie jest to synonim ``widok nie posiada
danych'' ani ``wszystkie wskaźniki są zawsze bezpieczne''.

\begin{lstlisting}
#include <concepts>
#include <ranges>
#include <vector>

static_assert(std::ranges::borrowed_range<
    std::ranges::ref_view<std::vector<int>>>);
static_assert(!std::ranges::borrowed_range<
    std::vector<int>>);
\end{lstlisting}

\subsection{\texttt{std::ranges::dangling} jako zabezpieczenie}

Algorytm zwracający iterator do \emph{tymczasowego} kontenera nie może
bezpiecznie zwrócić rzeczywistego iteratora do zniszczonych danych.
Dlatego pewne algorytmy Ranges zwracają specjalny typ
\texttt{std::ranges::dangling}.

\begin{lstlisting}
#include <type_traits>

using Result = decltype(std::ranges::find(
    std::vector<int>{1, 2, 3}, 2));
static_assert(std::same_as<Result, std::ranges::dangling>);
\end{lstlisting}

Nie należy jednak oczekiwać, że mechanizm ten wykryje wszystkie
problemy z czasem życia. Jeżeli zapisujemy nieposiadający danych
widok do lokalnej zmiennej, kompilator może nie zauważyć, że docelowy
kontener umrze wcześniej niż widok.

\section{\texttt{std::span}, \texttt{std::string\_view} i \texttt{subrange}}
\label{sec:cpp8-span}

\texttt{std::span<T>} opisuje ciągły fragment pamięci jako wskaźnik
i rozmiar (lub rozmiar znany w czasie kompilacji). Jest to lekka,
\textbf{nieposiadająca danych} alternatywa dla pary
\texttt{(T*, size)}.

\begin{lstlisting}
#include <span>

void normalize(std::span<double> xs) {
    for (double& x : xs) x /= 100.0;
}

std::vector<double> v{100, 200, 300};
normalize(v);  // v becomes {1, 2, 3}
\end{lstlisting}

\texttt{std::span<const T>} umożliwia odczyt bez modyfikowania
elementów, a \texttt{std::span<T>} pozwala modyfikować elementy, mimo
że sam widok nie posiada pamięci.

\begin{lstlisting}
int a[]{1, 2, 3, 4};
std::span<int> all{a};
auto head = all.first(2);
head[1] = 20; // a[1] == 20
\end{lstlisting}

\texttt{std::string\_view} spełnia podobną rolę dla tekstu; nie
alokuje pamięci na znaki (oraz nie musi kończyć się bajtem zero).
Trzeba uważać, aby nie zapisać go po zniszczeniu źródłowego
\texttt{std::string}.

\begin{lstlisting}
#include <string_view>

std::string text = "abcdef";
std::string_view view{text};
std::string_view prefix = view.substr(0, 3); // "abc"
// Do not use view after text is destroyed or invalidated.
\end{lstlisting}

\texttt{std::ranges::subrange<I, S>} grupuje iterator \texttt{I}
i sentinel \texttt{S} w zakres bez kopiowania elementów. Także tutaj
obowiązuje odpowiedzialność za czas życia wskazywanych danych.

\begin{lstlisting}
auto sub = std::ranges::subrange(v.begin(), v.end());
// sub refers to elements of v; it does not own v.
\end{lstlisting}

\paragraph{Ostrożnie z realokacją.}
\texttt{std::span} do elementów wektora, \texttt{string\_view} do
znaków łańcucha oraz \texttt{subrange} z iteratorów do kontenera mogą
stać się niepoprawne po operacji zmieniającej miejsce przechowywania
danych. Przenoszenie lub realokacja nie są automatycznie śledzone.

\section{Własny widok za pomocą \texttt{view\_interface}}
\label{sec:cpp8-custom-view}

Na końcu zbudujemy \emph{własny widok} generujący kwadraty liczb
całkowitych od zera do \(n-1\). Użyjemy
\texttt{std::ranges::view\_interface}, które dostarcza wygodnych
operacji opartych na \texttt{begin()} i \texttt{end()}. W klasie
pozostają wyłącznie parametry generatora, więc nie ma alokacji pamięci
na elementy sekwencji.

\begin{lstlisting}
#include <concepts>
#include <cstddef>
#include <iterator>
#include <ranges>

struct SquaresView : std::ranges::view_interface<SquaresView> {
    struct iterator {
        using value_type = int;
        using difference_type = std::ptrdiff_t;
        using iterator_concept = std::forward_iterator_tag;

        int index = 0;
        int operator*() const { return index * index; }
        iterator& operator++() {
            ++index;
            return *this;
        }
        iterator operator++(int) {
            auto old = *this;
            ++*this;
            return old;
        }
        bool operator==(const iterator&) const = default;
    };

    int n = 0;
    SquaresView() = default;
    explicit SquaresView(int count) : n(count) {}

    iterator begin() const { return {0}; }
    iterator end() const { return {n}; }
};

static_assert(std::ranges::view<SquaresView>);
static_assert(std::ranges::forward_range<SquaresView>);
\end{lstlisting}

\begin{lstlisting}
SquaresView squares{6};
for (int x : squares) { /* 0, 1, 4, 9, 16, 25 */ }

// Our view works with standard adapters:
auto even_squares = squares | std::views::filter(
    [](int x) { return x % 2 == 0; });
// Iteration yields 0, 4, 16.
\end{lstlisting}

Konstruktor z argumentem \texttt{count} powinien otrzymać wartość
nieujemną. Ponieważ liczymy \texttt{index * index} w typie
\texttt{int}, należy ponadto ograniczyć rozmiar zakresu, aby uniknąć
przepełnienia. W rozwiązaniu produkcyjnym należałoby te warunki
weryfikować oraz dopracować obsługę szerszych typów całkowitych.

\paragraph{Widok a adapter.}
\texttt{SquaresView} jest własnym \emph{widokiem}, który można
przekazać do standardowych adapterów. Nie jest jeszcze własnym
\emph{range adaptor closure}, czyli obiektem pozwalającym zapisać
\texttt{some\_range | squares\_adaptor}. Implementacja takiego
adaptera stanowi osobny poziom abstrakcji i nie jest konieczna do
budowania własnych widoków.

\section{Różnica między \texttt{transform} i \texttt{views::transform}}

Poniższe operacje mogą prowadzić do takich samych liczb, ale różnią
się \emph{czasem wykonania} i \emph{sposobem przechowywania wyniku}:

\begin{lstlisting}
std::vector<int> source{1, 2, 3};
std::vector<int> eager;
std::ranges::transform(source,
    std::back_inserter(eager),
    [](int x) { return x * x; });
// eager owns {1, 4, 9} right now

auto lazy = source | std::views::transform(
    [](int x) { return x * x; });
// lazy does not own a vector of squared values
\end{lstlisting}

Wariant pierwszy wymaga pamięci na trzy elementy i kończy obliczenia
przed kolejną instrukcją. Wariant drugi może obliczać wartości dopiero
przy dereferencji iteratorów. Po zmianie \texttt{source} kolejne
odczyty widoku \texttt{lazy} mogą odzwierciedlać nowe dane.

\section{Projektowanie API z wykorzystaniem Ranges}
\label{sec:cpp8-api}

Jeżeli funkcja tylko przegląda elementy dowolnego źródła, preferuj
parametr ograniczony koncepcją \texttt{input\_range}. Dla wielokrotnych
przejść korzystaj z \texttt{forward\_range}. Jeśli potrzebujesz
ciągłego bloku pamięci i znasz typ elementów, wygodnym wyborem jest
\texttt{std::span<T>}.

\begin{lstlisting}
template <std::ranges::input_range R>
void print_each(R&& r) {
    for (auto&& x : r) {
        // Print or consume x.
    }
}

void process_buffer(std::span<const double> values) {
    for (double value : values) {
        // Read numeric values.
    }
}
\end{lstlisting}

Zależnie od celu można zwracać widok albo kontener. Widok jest
wygodny, gdy chcemy dalej budować potok i mamy kontrolę nad czasem
życia źródła. Kontener jest często bezpieczniejszy, gdy zwracany
wynik ma istnieć niezależnie od wszystkich argumentów.

\section{Najczęstsze błędy i praktyczne reguły}
\label{sec:cpp8-pitfalls}

\begin{enumerate}
\item \textbf{Dereferencja \texttt{end()}.} Iterator końcowy jest
granicą, a nie elementem. Przed \texttt{*it} trzeba mieć
\texttt{it != end}.

\item \textbf{Mylenie kategorii iteratorów.} Lista nie obsługuje
skoków w stałym czasie. Nie można użyć dla niej algorytmu
\texttt{ranges::sort}; należy wywołać metodę \texttt{list::sort}.

\item \textbf{Traktowanie widoku jak kontenera.} \texttt{views::filter}
nie tworzy kopii wybranych elementów. Zmiany i zniszczenie źródła
mogą wpływać na wynik.

\item \textbf{Wiszące referencje w lambdach.} Zwrócony widok
z predykatem przechwytującym lokalny obiekt przez referencję może
być niebezpieczny nawet wtedy, gdy sam kontener ma odpowiednio długi
czas życia.

\item \textbf{Efekty uboczne w predykacie.} Liczba i kolejność
wywołań w leniwym potoku nie są odpowiednią podstawą działania
programu; predykaty filtrowania powinny zachować stabilną semantykę.

\item \textbf{Zakładanie stałego kosztu operacji.} Obliczenie
\texttt{distance} lub pierwszego pasującego elementu może wymagać
przejrzenia dużej części źródła.

\item \textbf{Mieszanie standardów.} \texttt{ranges::to},
\texttt{views::zip} i \texttt{views::enumerate} nie należą do C++20.

\item \textbf{Używanie nieważnych iteratorów.} Realokacja wektora
lub usuwanie elementu może unieważnić iteratory także wtedy, gdy
znajdują się wewnątrz wcześniej utworzonego widoku.
\end{enumerate}

\section{Której konstrukcji użyć?}

\begin{center}
\small
\begin{tabular}{|p{6.0cm}|p{6.5cm}|}
\hline
\textbf{Potrzeba} & \textbf{Zalecany wybór} \\
\hline
Przejście po zwykłym kontenerze & Pętla \texttt{for (auto\& x : v)} \\
\hline
Jednorazowe sortowanie albo wyszukiwanie & \texttt{std::ranges::sort}, \texttt{std::ranges::find} \\
\hline
Leniwe filtrowanie i przekształcanie & \texttt{std::views::filter}, \texttt{std::views::transform} \\
\hline
Generowanie sekwencji indeksów & \texttt{std::views::iota} \\
\hline
Odczyt lub modyfikacja ciągłego bufora & \texttt{std::span<T>} (bez własności) \\
\hline
Lekki odczyt fragmentu tekstu & \texttt{std::string\_view} (bez własności) \\
\hline
Zakres wyznaczony iteratorami & \texttt{std::ranges::subrange} \\
\hline
Wynik niezależny od źródła & Kontener, np. \texttt{std::vector} \\
\hline
Własny generator wartości & Typ modelujący \texttt{std::ranges::view} \\
\hline
Wydajny ogólny algorytm generyczny & Szablon z \texttt{std::ranges::input\_range} lub silniejszą koncepcją \\
\hline
\end{tabular}
\end{center}

\paragraph{Pytania kontrolne.}
\begin{enumerate}
\item Dlaczego iterator \texttt{end()} nie jest elementem kolekcji?
\item Co odróżnia \texttt{forward\_iterator} od \texttt{input\_iterator}?
\item Czy \texttt{std::ranges::sort} działa dla \texttt{std::list}? Dlaczego?
\item Po co \texttt{std::ranges} dopuszcza sentinel innego typu niż iterator?
\item W którym momencie \texttt{views::transform} wykonuje funkcję?
\item Dlaczego \texttt{views::filter} nie powinien zmieniać elementów, od których zależy jego predykat?
\item Czym różni się \texttt{views::all(lvalue)} od \texttt{views::all(rvalue)} dla \texttt{std::vector}?
\item Co gwarantuje, a czego nie gwarantuje \texttt{borrowed\_range}?
\item Dlaczego \texttt{ranges::find(std::vector\{\dots\}, x)} może zwrócić \texttt{dangling}?
\item Jak zebrać elementy \texttt{filter\_view} do wektora w czystym C++20?
\item Czym różni się \texttt{std::span<const T>} od \texttt{std::vector<T>}?
\item Co należy zaimplementować, aby własny obiekt był \texttt{forward\_range}?
\end{enumerate}

\paragraph{Dalsza lektura.}
Dobrym punktem odniesienia jest dokumentacja \emph{cppreference}:
\emph{Ranges library}, \emph{Iterator library}, \emph{Range adaptors},
\emph{std::ranges::borrowed\_range}, \emph{std::ranges::dangling}
oraz \emph{std::ranges::view\_interface}. Warto testować kod z opcjami kompilatora
\texttt{-std=c++20} i \texttt{-Wall -Wextra -Wpedantic},
a także korzystać z AddressSanitizer przy badaniu czasu życia danych.

\chapter{Współbieżność i asynchroniczność}
\label{chap:concurrency}

\section{Motywacja i podstawowe pojęcia}

W programie jednowątkowym instrukcje wykonywane są w jednym \emph{wątku sterowania}.
Nie znaczy to, że procesor nie wykonuje w tle innych czynności; chodzi o model
samego programu. W programie współbieżnym istnieje kilka niezależnych ciągów
instrukcji, których wykonanie może się przeplatać. Zamiast jednego zadania
rozpatrujemy np. odczyt danych, obliczenia i zapisywanie wyników.

\begin{description}[style=nextline]
\item[Współbieżność (\emph{concurrency})]
Wiele zadań może robić postępy w nakładających się okresach czasu.
Nie wymaga to, aby fizycznie wykonywały się w tej samej chwili.
\item[Równoległość (\emph{parallelism})]
Co najmniej dwa zadania wykonują się fizycznie jednocześnie, na przykład
na dwóch rdzeniach procesora.
\item[Asynchroniczność (\emph{asynchrony})]
Rozpoczęcie operacji nie wymusza natychmiastowego oczekiwania na jej wynik.
Wywołujący może kontynuować pracę, a wynik odebrać później.
\end{description}

Pojęcia te nie są synonimami. Można mieć współbieżność bez równoległości
(jeden rdzeń i przełączanie zadań), a także interfejs asynchroniczny bez nowego
wątku (np. opóźnione wykonanie funkcji).

\nota{Zegar i kolejność komunikatów w przykładach współbieżnych nie stanowią
gwarantowanego wyniku programu. Planista systemu może przeplatać wątki inaczej
przy każdym uruchomieniu.}

\subsection{Proces, wątek, zadanie i współdzielone dane}

Proces zwykle posiada własną przestrzeń adresową. Wątki tego samego procesu
współdzielą między innymi pamięć dynamiczną i zmienne globalne, ale mają własne
stosy i własne wykonanie instrukcji. \emph{Zadanie} to jednostka pracy,
którą można przypisać wątkowi; nie każde zadanie musi oznaczać nowy wątek.

Dostęp do wspólnego obiektu wymaga szczególnej ostrożności, jeżeli co najmniej
jeden wątek go modyfikuje. Sam fakt istnienia wielu wątków nie jest błędem;
problem pojawia się przy niezsynchronizowanym dostępie do pamięci.

\section{Uruchamianie wątków: \texttt{std::thread}}
\label{sec:thread}

Nagłówek \code{<thread>} udostępnia typ \code{std::thread} (C++11).
Utworzenie obiektu z funkcją uruchamia nowy wątek. Wywołanie \code{join()}
czeka na jego zakończenie; \code{joinable()} informuje, czy obiekt nadal
reprezentuje wątek wymagający obsłużenia.

\begin{lstlisting}[caption={Pierwszy program z dwoma wątkami (C++11).}]
#include <iostream>
#include <thread>

void work(int id) {
    std::cout << "Worker " << id << "\n";
}

int main() {
    std::thread a(work, 1);
    std::thread b(work, 2);

    a.join();
    b.join();
}
\end{lstlisting}

Kolejność wypisywania nie jest określona, a fragmenty wypisywane przez
różne wątki mogą się przeplatać. Jeżeli wymagana jest spójna linia wyjścia,
można chronić wypisywanie mutexem albo użyć \code{std::osyncstream} (C++20).

\subsection{Przekazywanie argumentów i czasu życia}

Argumenty przekazywane do konstruktora \code{std::thread} są odpowiednio
przechowywane na potrzeby nowego wątku; domyślnie nie oznacza to przekazania
modyfikowalnego argumentu przez referencję. Referencję można przekazać
za pomocą \code{std::ref} albo przechwycenia referencji w lambdzie.

\begin{lstlisting}[caption={Bezpieczne przekazanie referencji z zachowaniem jej czasu życia.}]
#include <functional>
#include <iostream>
#include <thread>

void increment(int& value) { ++value; }

int main() {
    int counter = 0;
    std::thread worker(increment, std::ref(counter));
    worker.join(); // counter nadal istnieje
    std::cout << counter << '\n'; // 1
}
\end{lstlisting}

Destruktor obiektu \code{std::thread}, który nadal jest
\code{joinable()}, wywołuje \code{std::terminate()}.
Dlatego przy wielu ścieżkach wyjścia lepiej stosować RAII,
np. \code{std::jthread}. Operacja \code{detach()} wymaga z kolei
gwarancji, że wątek nie będzie używał już zniszczonych obiektów.

\section{Wyścig danych, warunek wyścigu i relacja \emph{happens-before}}
\label{sec:race}

\term{Wyścig danych (\emph{data race})} występuje, gdy co najmniej dwa
potencjalnie współbieżne dostępy do tego samego miejsca pamięci są konfliktowe
(co najmniej jeden modyfikuje), nie są odpowiednio uporządkowane przez
synchronizację i nie są to dostępy atomowe w wymaganym sensie. W standardowym
C++ prowadzi to do \textbf{niezdefiniowanego zachowania}.

\begin{lstlisting}[caption={Błędny przykład: niezdefiniowane zachowanie.}]
// Fragment celowo niepoprawny; NIE uruchamiac jako poprawnego przykladu.
int counter = 0;
// Watki rownoczesnie wykonuja ++counter bez synchronizacji.
\end{lstlisting}

\term{Warunek wyścigu (\emph{race condition})} jest pojęciem szerszym:
wynik programu zależy od względnej kolejności zdarzeń. Warunek wyścigu nie
zawsze oznacza niezdefiniowane zachowanie; np. dwa poprawnie zsynchronizowane
wątki mogą rywalizować o to, który pierwszy wykona zadanie.

\term{Relacja \emph{happens-before}} formalizuje, jakie skutki wcześniejszych
operacji muszą być widoczne dla późniejszych. Tworzy się m.in. przez porządek
operacji w jednym wątku i przez operacje synchronizujące (np. odpowiednio
sparowane \code{release}/\code{acquire}, odblokowanie/zablokowanie mutexa).
To centralny element modelu pamięci C++.

\section{Ochrona danych: \texttt{mutex} i blokady RAII}
\label{sec:mutex}

\code{std::mutex} pozwala maksymalnie jednemu wątkowi utrzymywać blokadę
w danej chwili. Jeśli kilka operacji ma tworzyć niepodzielną \emph{sekcję
krytyczną}, umieszczamy je pod tą samą blokadą.

\begin{lstlisting}[caption={Bezpieczny licznik przy użyciu mutexa (C++11).}]
#include <iostream>
#include <mutex>
#include <thread>
#include <vector>

int main() {
    int counter = 0;
    std::mutex mtx;
    std::vector<std::thread> workers;

    for (int i = 0; i < 4; ++i) {
        workers.emplace_back([&] {
            for (int j = 0; j < 10000; ++j) {
                std::lock_guard<std::mutex> lock(mtx);
                ++counter;
            }
        });
    }
    for (auto& t : workers) t.join();
    std::cout << counter << '\n'; // 40000
}
\end{lstlisting}

\code{lock_guard} automatycznie przejmuje mutex w konstruktorze i zwalnia
w destruktorze. Dzięki temu zwolnienie blokady następuje także podczas
rozwijania stosu po wyjątku. \code{unique_lock} jest bardziej elastyczny:
pozwala jawnie odblokować i ponownie zablokować mutex oraz współpracuje
z \code{condition_variable}.

\subsection{Kilka mutexów i zakleszczenie}

Jeżeli dwa wątki blokują mutexy w przeciwnej kolejności, może wystąpić
\emph{deadlock}: pierwszy czeka na drugi mutex, drugi czeka na pierwszy.
W C++17 można bezpiecznie przejmować kilka mutexów przez \code{scoped_lock}.

\begin{lstlisting}[caption={Łączne przejęcie dwóch mutexów (C++17).}]
#include <mutex>

int main() {
    std::mutex a, b;
    {
        std::scoped_lock lock(a, b);
        // Operacje chronione jednoczesnie przez a i b.
    }
}
\end{lstlisting}

\nota{Mutexy nie chronią ,,samej zmiennej'' automatycznie: poprawność wymaga,
aby \emph{wszystkie} konfliktowe dostępy do chronionego stanu stosowały
uzgodniony mechanizm synchronizacji.}

\section{Oczekiwanie na stan: \texttt{condition\_variable}}
\label{sec:cv}

Mutex zapewnia wyłączność, ale nie sygnalizuje gotowości danych.
\code{std::condition_variable} pozwala czekać bez aktywnego odpytywania.
Do oczekiwania z predykatem używamy blokady
\begin{center}
\code{std::unique_lock<std::mutex>}.
\end{center}
Predykat opisuje stan, na który czeka konsument.

\begin{lstlisting}[caption={Producent i konsument z warunkiem gotowości (C++11).}]
#include <condition_variable>
#include <iostream>
#include <mutex>
#include <thread>

int main() {
    std::mutex mtx;
    std::condition_variable cv;
    bool ready = false;
    int data = 0;

    std::thread consumer([&] {
        std::unique_lock<std::mutex> lock(mtx);
        cv.wait(lock, [&] { return ready; });
        std::cout << data << '\n'; // 42
    });

    std::thread producer([&] {
        {
            std::lock_guard<std::mutex> lock(mtx);
            data = 42;
            ready = true;
        }
        cv.notify_one();
    });

    producer.join();
    consumer.join();
}
\end{lstlisting}

Wywołanie \code{wait(lock, pred)} ma następujący sens:
\begin{enumerate}
\item podczas trzymania mutexa sprawdzany jest predykat;
\item jeśli jest fałszywy, wątek atomowo zwalnia mutex i przechodzi w oczekiwanie;
\item po powiadomieniu (lub pozornym wybudzeniu) odzyskuje mutex i znów sprawdza predykat;
\item kończy czekanie dopiero wtedy, gdy predykat jest prawdziwy.
\end{enumerate}

Dopuszczalne są \emph{spurious wakeups}, czyli wybudzenia bez zmiany interesującego
nas stanu. Dlatego samo \code{cv.wait(lock)} zwykle nie wystarcza.
Powiadomienie \code{notify_one()} budzi co najwyżej jeden oczekujący wątek,
a \code{notify_all()} umożliwia wybudzenie wszystkich. Powiadomienia nie są
kolejkowane jak komunikaty; stan \code{ready}, a nie samo powiadomienie,
stanowi źródło prawdy.

\subsection{Kolejka producent--konsument}

W ogólniejszym przypadku producent umieszcza wiele elementów w kolejce,
a konsument czeka na warunek ,,kolejka niepusta lub zakończono produkcję''.
Flaga zakończenia jest konieczna, aby konsument nie oczekiwał bez końca
po przetworzeniu ostatniego elementu.

\begin{lstlisting}[caption={Współdzielona kolejka i sygnał zakończenia (C++11).}]
#include <condition_variable>
#include <iostream>
#include <mutex>
#include <queue>
#include <thread>

int main() {
    std::mutex m;
    std::condition_variable cv;
    std::queue<int> q;
    bool done = false;

    std::thread consumer([&] {
        for (;;) {
            int item;
            {
                std::unique_lock<std::mutex> lock(m);
                cv.wait(lock, [&] { return done || !q.empty(); });
                if (q.empty() && done) break;
                item = q.front();
                q.pop();
            }
            std::cout << "Consumed " << item << '\n';
        }
    });

    std::thread producer([&] {
        for (int x = 1; x <= 3; ++x) {
            {
                std::lock_guard<std::mutex> lock(m);
                q.push(x);
            }
            cv.notify_one();
        }
        {
            std::lock_guard<std::mutex> lock(m);
            done = true;
        }
        cv.notify_all();
    });

    producer.join();
    consumer.join();
}
\end{lstlisting}

\section{Operacje atomowe: \texttt{std::atomic}}
\label{sec:atomic}

\code{std::atomic<T>} udostępnia atomowe operacje na obiekcie \code{T}.
Nie oznacza to ogólnie ,,braku jakichkolwiek blokad'' w implementacji:
nie każda instancja atomika jest \emph{lock-free}. Można to sprawdzić
metodą \code{is_lock_free()}.

\begin{lstlisting}[caption={Licznik atomowy z porządkiem relaxed (C++11).}]
#include <atomic>
#include <iostream>
#include <thread>
#include <vector>

int main() {
    std::atomic<int> counter{0};
    std::vector<std::thread> threads;
    for (int i = 0; i < 4; ++i) {
        threads.emplace_back([&] {
            for (int j = 0; j < 10000; ++j)
                counter.fetch_add(1, std::memory_order_relaxed);
        });
    }
    for (auto& t : threads) t.join();
    std::cout << counter.load() << '\n'; // 40000
}
\end{lstlisting}

W przypadku samego licznika potrzebujemy niepodzielności inkrementacji,
a nie synchronizacji dodatkowych danych. Stąd wystarcza
\code{memory_order_relaxed}. Domyślnie operacje atomowe stosują
\code{memory_order_seq_cst}, który narzuca silniejsze gwarancje porządkowania.

\subsection{Operacja porównaj-i-zamień}

\code{compare_exchange_weak(expected, desired)} atomowo sprawdza, czy
obecna wartość jest równa \code{expected}. Przy powodzeniu zapisuje
\code{desired}; przy niepowodzeniu aktualizuje \code{expected} do wartości
odczytanej. Wersja \code{weak} może zgłosić niepowodzenie także pozornie,
dlatego często umieszcza się ją w pętli.

\begin{lstlisting}[caption={Atomowe podnoszenie wartości do zadanego minimum (C++11).}]
#include <atomic>
#include <iostream>

int main() {
    std::atomic<int> maximum{3};
    int desired = 10;
    int expected = maximum.load(std::memory_order_relaxed);

    while (expected < desired &&
           !maximum.compare_exchange_weak(
               expected, desired,
               std::memory_order_relaxed)) {
        // expected zawiera najnowsza odczytana wartosc
    }
    std::cout << maximum.load() << '\n'; // 10
}
\end{lstlisting}

\subsection{Porządek pamięci: \texttt{release} i \texttt{acquire}}

Atomowość jednej zmiennej nie wystarcza, gdy jednym sygnałem chcemy
bezpiecznie opublikować inne dane. Stosujemy wówczas odpowiednio sparowane
operacje \code{release} i \code{acquire}.

\begin{lstlisting}[caption={Bezpieczna publikacja zwykłej zmiennej (C++11).}]
#include <atomic>
#include <iostream>
#include <thread>

int main() {
    int data = 0;
    std::atomic<bool> ready{false};

    std::thread producer([&] {
        data = 42;
        ready.store(true, std::memory_order_release);
    });
    std::thread consumer([&] {
        while (!ready.load(std::memory_order_acquire)) {
            std::this_thread::yield();
        }
        std::cout << data << '\n'; // 42
    });

    producer.join();
    consumer.join();
}
\end{lstlisting}

Gdy odczyt \code{acquire} zaobserwuje właściwy zapis
\code{release}, operacje \emph{synchronizes-with} gwarantują
widoczność wcześniejszego zapisu \code{data = 42}.
Nie wolno tu zamieniać obu porządków na \code{relaxed}.

\begin{center}
\begin{tabularx}{\textwidth}{@{}lX@{}}
\toprule
\textbf{Porządek} & \textbf{Zasadnicza gwarancja} \\
\midrule
\code{relaxed} & Atomowość bez synchronizacji pozostałych danych. \\
\code{release} & Publikowanie wcześniejszych efektów do odpowiedniego odczytu. \\
\code{acquire} & Odczyt odbierający efekty opublikowane przez \code{release}. \\
\code{acq_rel} & Połączenie obu kierunków dla operacji odczyt--modyfikacja--zapis. \\
\code{seq_cst} & Dodatkowy wspólny porządek operacji sekwencyjnie spójnych. \\
\bottomrule
\end{tabularx}
\end{center}

Istnieje również \code{memory_order_consume}, ale jego historia i obsługa
przez kompilatory są skomplikowane; nie wykorzystujemy go w przykładach.
Programiście początkującemu zalecamy domyślny \code{seq_cst}, a słabsze
porządki dopiero po uzasadnieniu poprawności.

\subsection{Czekanie na atomik bez aktywnej pętli: C++20}

Od C++20 obiekt atomowy udostępnia \code{wait(old)} oraz
\code{notify_one()/notify_all()}. Oczekiwanie kończy się, kiedy
zaobserwowana wartość różni się od \code{old}; wątek nie musi stale
wykonywać aktywnego odpytywania.

\begin{lstlisting}[caption={Atomowe czekanie i wybudzanie (C++20).}]
#include <atomic>
#include <iostream>
#include <thread>

int main() {
    std::atomic<bool> ready{false};
    int payload = 0;

    std::jthread worker([&] {
        payload = 42;
        ready.store(true, std::memory_order_release);
        ready.notify_one();
    });

    ready.wait(false, std::memory_order_acquire);
    std::cout << payload << '\n'; // 42
}
\end{lstlisting}

\nota{Atomowy znacznik nie jest zamiennikiem mutexa chroniącego wiele
powiązanych pól struktury. Algorytmy \emph{lock-free} bywają trudniejsze
od rozwiązania opartego na blokadzie; dodatkowo należy rozumieć m.in.
problem ABA i czas życia współdzielonych obiektów.}

% === Rozszerzenie: zaawansowany model pamięci C++ ===
\section{Model pamięci C++: formalne podstawy synchronizacji}
\label{sec:advanced-memory-model}

Sekcja~\ref{sec:race} wprowadziła pojęcie wyścigu danych. Teraz sprecyzujemy,
\emph{dlaczego} poprawny program wielowątkowy nie może polegać na fizycznym
porządku wykonywania instrukcji. Model pamięci C++ określa dopuszczalne
obserwacje programu również wtedy, gdy kompilator zmienia kolejność instrukcji,
procesor używa buforów zapisów, a rdzenie mają prywatne pamięci podręczne.
Nie jest to model jednej konkretnej architektury (np. x86 albo ARM).

\subsection{Trzy relacje: \emph{sequenced-before}, \emph{synchronizes-with} i \emph{happens-before}}

Niech $A$ i $B$ oznaczają dwa zdarzenia (ewaluacje) programu.
\begin{description}[style=nextline]
\item[\emph{Sequenced-before}]
Relacja porządkująca ewaluacje wewnątrz jednego wątku, wynikająca z zasad
sekwencjonowania w C++. Nie wolno jej utożsamiać z samą kolejnością
występowania tokenów w źródle: nie każda para podwyrażeń jest sekwencjonowana.
\item[\emph{Synchronizes-with}]
Relacja między odpowiednimi zdarzeniami synchronizacyjnymi, np. skutecznym
odblokowaniem mutexa i późniejszym jego zablokowaniem albo zapisem
\code{release} i odczytem \code{acquire}, który obserwuje zapis
(lub odpowiedni element jego sekwencji \emph{release}).
\item[\emph{Happens-before}]
Relacja logicznego porządku zdarzeń. W przykładach bez osobliwego
porządku \emph{consume} można ją konstruować z dwóch poprzednich
relacji i przechodniości. Jeżeli $A$ \emph{happens-before} $B$,
skutki $A$ ograniczają dopuszczalne obserwacje zdarzenia $B$.
\end{description}

Dla trzech zdarzeń reguła przechodniości ma postać
\[
 A\mathrel{\mathrm{hb}}B\ \land\ B\mathrel{\mathrm{hb}}C
 \quad\Longrightarrow\quad A\mathrel{\mathrm{hb}}C.
\]

Przykładowy łańcuch synchronizacji można zapisać jako
\[
\begin{aligned}
A:\text{ zapis }x
  &\xrightarrow{\mathrm{sequenced\mbox{-}before}}
    B:\text{ store(release)},\\
B&\xrightarrow{\mathrm{synchronizes\mbox{-}with}}
    C:\text{ load(acquire)},\\
C&\xrightarrow{\mathrm{sequenced\mbox{-}before}}
    D:\text{ odczyt }x.
\end{aligned}
\]
Z przechodniości wynika $A\mathrel{\mathrm{hb}}D$. Istotne jest jednak,
że operacja \code{load(acquire)} musi rzeczywiście odczytać wartość pochodzącą
z odpowiedniego zapisu \code{release} (lub jego sekwencji); samo użycie
obu etykiet w dwóch wątkach nie wystarcza.

\subsection{Wyścig danych a wyścig logiczny: przykład}

\begin{lstlisting}[caption={Dwa atomiki eliminują data race, ale nie ustalają zwycięzcy.}]
#include <atomic>
#include <thread>

int main() {
    std::atomic<int> winner{0};
    std::jthread a([&] {
        int expected = 0;
        winner.compare_exchange_strong(expected, 1);
    });
    std::jthread b([&] {
        int expected = 0;
        winner.compare_exchange_strong(expected, 2);
    });
    a.join();
    b.join();
    // winner == 1 albo winner == 2.
}
\end{lstlisting}

Wynik zależy od harmonogramu, ale program nie ma wyścigu danych.
\emph{Race condition} jest błędem dopiero wtedy, gdy program wymaga
konkretnego zwycięzcy albo narusza swój kontrakt w części harmonogramów.

\subsection{Porządek modyfikacji pojedynczego atomika}

Każdy obiekt atomowy ma własny, spójny dla wszystkich wątków
\emph{modification order}: całkowity porządek jego atomowych modyfikacji.
Nie oznacza to wspólnego globalnego porządku wszystkich atomików.
Dla modyfikacji \code{x} oraz \code{y} mogą więc istnieć obserwacje,
które intuicyjnie zaskakują, mimo że program nie zawiera \emph{data race}.

W poniższym przykładzie obie wartości startują od zera:
\begin{lstlisting}[caption={Litmus test: dwa zapisy, dwa odczyty (fragment dwóch wątków).}]
std::atomic<int> x{0}, y{0};
int r1 = -1, r2 = -1;

// Watek A:
x.store(1, std::memory_order_relaxed);
r1 = y.load(std::memory_order_relaxed);

// Watek B:
y.store(1, std::memory_order_relaxed);
r2 = x.load(std::memory_order_relaxed);
\end{lstlisting}

Przy porządku \code{relaxed} \emph{dopuszczalne} jest $r_1=r_2=0$.
Zastąpienie \emph{wszystkich czterech} operacji przez \code{seq_cst}
wyklucza ten wynik: jeden wspólny porządek musiałby wtedy spełniać
cykliczne sprzeczne ograniczenia. Test uruchomiony tylko na jednym
procesorze niczego tu nie dowodzi: brak obserwacji nie oznacza zakazu
w standardzie.

\subsection{Porządki atomowe w praktyce}

\begin{center}
\small
\begin{tabularx}{\textwidth}{@{}p{0.25\textwidth}X@{}}
\toprule
\textbf{Porządek} & \textbf{Praktyczny sens i typowe zastosowanie} \\
\midrule
\code{relaxed} & Niepodzielność i porządek modyfikacji danego atomika; liczniki statystyk, gdy nie synchronizują innego stanu. \\
\code{release} & Publikuje efekty poprzedzających operacji; stosowany do zapisów lub odpowiednich operacji RMW. \\
\code{acquire} & Odbiera publikację po odczytaniu właściwej wartości; stosowany do odczytów lub odpowiednich RMW. \\
\code{acq_rel} & Oba kierunki synchronizacji podczas operacji typu read--modify--write. \\
\code{seq_cst} & Domyślny, mocniejszy porządek; dodatkowo wspólne uporządkowanie wszystkich operacji \code{seq_cst}. \\
\bottomrule
\end{tabularx}
\end{center}

\code{memory_order_consume} jest historycznie problematyczny; nie stanowi
dobrego punktu wyjścia do pisania przenośnego kodu. Wybór słabszego
porządku należy poprzedzić argumentem poprawności, a nie samym wynikiem
benchmarku. Nie wszystkie porządki są dozwolone dla każdej operacji,
np. zwykły \code{store} nie może użyć \code{memory_order_acquire}.

\subsection{Publikacja danych i sekwencja \emph{release}}

Rozpatrzmy trzy zdarzenia: producent wypełnia dane, drugi wątek
zwiększa atomowy licznik, a trzeci odbiera wynik. W C++20 (i wcześniej)
sekwencja \emph{release} po zapisie \code{release} może być kontynuowana
przez kolejne atomowe operacje \emph{read--modify--write} tego samego
atomika, także wykonywane przez inne wątki. Dlatego odczyt \code{acquire}
wartości zapisanej przez taką operację RMW może synchronizować się
z początkowym zapisem \code{release}.

\begin{lstlisting}[caption={Sekwencja release: fragment przy ustalonej kolejności operacji.}]
int payload = 0;
std::atomic<int> state{0};

// Producent (A):
payload = 42;
state.store(1, std::memory_order_release);

// Posrednik (B), DOPIERO po wykonaniu powyzszego store:
state.fetch_add(1, std::memory_order_relaxed); // 1 -> 2

// Konsument (C): odczytuje 2 z RMW posrednika:
if (state.load(std::memory_order_acquire) == 2) {
    // Wolno odczytac payload == 42.
}
\end{lstlisting}

Warunek o kolejności jest częścią przykładu: bez niego pośrednik mógłby
zwiększyć \code{state} jeszcze przed publikacją. Zwykły zapis wykonany
przez inny wątek przerywa tę sekwencję; nie należy jej mylić z dowolną
serią zapisów do tego samego atomika.

\subsection{Operacje CAS: sukces i porażka}

\code{compare_exchange_weak} i \code{compare_exchange_strong} mają
przeciążenia z osobnymi porządkami pamięci dla sukcesu i porażki.
Operacja nieudana jest \emph{odczytem}, a nie zapisem. Dlatego jej porządek
nie może być \code{release} ani \code{acq_rel} oraz nie może być silniejszy
od porządku sukcesu w sensie wymagań standardu.

\begin{lstlisting}[caption={Pętla CAS: zwiększanie licznika, przykład praktyczny.}]
#include <atomic>

void add_one(std::atomic<int>& x) {
    int expected = x.load(std::memory_order_relaxed);
    while (!x.compare_exchange_weak(
        expected, expected + 1,
        std::memory_order_acq_rel,
        std::memory_order_acquire)) {
        // expected zostalo zaktualizowane po niepowodzeniu.
    }
}
\end{lstlisting}

W tym przypadku \code{fetch_add(1)} jest prostsze i zwykle lepsze.
CAS ma sens zwłaszcza przy warunkowej zmianie stanu lub złożonych
strukturach atomowych. \code{weak} może nie powieść się pozornie,
więc jest wygodny w pętli; \code{strong} lepiej nadaje się do
pojedynczych prób. Użycie \code{acq_rel} w powyższym liczniku służy
zilustrowaniu API; gdy licznik niczego nie publikuje, wystarczy
\code{memory_order_relaxed}.

\subsection{Bariery pamięci: \texttt{atomic\_thread\_fence}}

\code{std::atomic_thread_fence(order)} nakłada ograniczenia porządkowania
pamięci, ale \emph{sama} bariera nie synchronizuje magicznie dwóch wątków.
Potrzebna jest odpowiednia operacja atomowa, która połączy strony
synchronizacji. Przykład celowo przenosi \code{release}/\code{acquire}
z atomika na bariery:

\begin{lstlisting}[caption={Publikacja przez dwie bariery i atomowy znacznik (C++11).}]
#include <atomic>
#include <thread>
#include <cassert>

int main() {
    int payload = 0;
    std::atomic<bool> ready{false};

    std::thread producer([&] {
        payload = 42;
        std::atomic_thread_fence(std::memory_order_release);
        ready.store(true, std::memory_order_relaxed);
    });
    std::thread consumer([&] {
        while (!ready.load(std::memory_order_relaxed)) {
            std::this_thread::yield();
        }
        std::atomic_thread_fence(std::memory_order_acquire);
        assert(payload == 42);
    });
    producer.join();
    consumer.join();
}
\end{lstlisting}

Dla tego układu bariery łączą się poprzez odczyt \code{ready}, który
obserwuje zapis producenta. Zwykłe \code{ready.store(release)} oraz
\code{ready.load(acquire)} byłoby prostsze. \code{atomic_signal_fence}
ma inny cel: ogranicza określone przestawienia przez kompilator,
a nie pełni roli zwykłej synchronizacji między wątkami.

\subsection{Atomowość, blokady i gwarancje postępu}

\begin{description}[style=nextline]
\item[\emph{Blocking}] Wątek może oczekiwać na zwolnienie blokady przez
innego uczestnika, który np. został wywłaszczony przez system.
\item[\emph{Lock-free}] W standardowym ujęciu algorytmowym system jako
całość może robić postęp mimo tego, że pewne wątki wielokrotnie ponawiają
operację. Pojedynczy wątek może ulegać zagłodzeniu.
\item[\emph{Wait-free}] Każdy uczestnik kończy operację w ograniczonej
liczbie własnych kroków, niezależnie od harmonogramu innych uczestników.
\end{description}

\code{atomic<T>::is_lock_free()} bada własność implementacji operacji
na tym atomiku; nie jest dowodem, że dowolny algorytm zbudowany z
atomików jest \emph{lock-free}. Nie wolno też wnioskować, że
\emph{lock-free} zawsze oznacza szybszy kod. Rywalizacja o jedną
linię cache może dominować koszt obliczeń.

\subsection{Problem ABA i odzyskiwanie pamięci}

Niech wierzchołek współbieżnego stosu zawiera wskaźnik \code{A}.
Wątek~1 odczytuje \code{A} i zapamiętuje wskaźnik do następnego
węzła, po czym zostaje wstrzymany. Wątek~2 usuwa \code{A}, usuwa
kolejny element \code{B} i ponownie umieszcza \code{A} na szczycie.
Wątek~1 widzi znowu \code{A}: porównanie wskaźnika przez CAS może się
powieść, choć pomiędzy odczytem i zapisem struktura zmieniła się
istotnie. Co gorsza, wcześniej zapamiętany węzeł mógł zostać już
zwolniony. To właśnie \emph{problem ABA}.

\[
\boxed{\text{A}\to\text{B}\to\text{C}}
\quad\longrightarrow\quad
\boxed{\text{B}\to\text{C}}
\quad\longrightarrow\quad
\boxed{\text{A}\to\text{C}}.
\]

Do typowych narzędzi należą: oznaczanie wskaźnika licznikiem wersji
(\emph{tagged pointer}), \emph{hazard pointers}, \emph{epoch-based
reclamation} i \emph{reference counting}. Zabezpieczenie samego
porównania wskaźników nie załatwia problemu \emph{czasu życia obiektu}.
Dla zwykłej kolejki aplikacyjnej \code{std::mutex} bywa
najprostszy i najbardziej wiarygodny.

\subsection{Zagłodzenie, \emph{livelock} i błędne podwójne sprawdzenie}

\emph{Deadlock} oznacza brak postępu wskutek cyklicznego oczekiwania.
\emph{Livelock} występuje, gdy wątki aktywnie reagują na siebie,
ale nie kończą pracy. \emph{Starvation} to brak postępu konkretnego
wątku mimo że inne wątki pracują. Te trzy zjawiska należy odróżniać.

W starszym kodzie spotyka się tzw. \emph{double-checked locking}.
Bez poprawnych atomików albo mutexa nie można go naprawić przez samo
\code{volatile}. W nowym kodzie znacznie lepiej używać bezpiecznej
inicjalizacji statycznej zmiennej lokalnej (C++11) lub \code{std::call_once}.
\code{volatile} nie zapewnia atomowości ani synchronizacji wątków.



\section{Wynik przyszłego obliczenia: \texttt{future} i \texttt{promise}}
\label{sec:future-promise}

\code{std::future<T>} reprezentuje możliwość odebrania \emph{jednego}
wyniku typu \code{T}, który może jeszcze nie być dostępny.
\code{std::promise<T>} umożliwia producentowi ten wynik ustawić.
Te obiekty odnoszą się do wspólnego \emph{stanu asynchronicznego}
(ang. \emph{shared state}). Sam \code{promise} nie uruchamia wątku.

\begin{center}
\fbox{\small\strut\code{promise.set_value(x)}\quad
$\longrightarrow$\quad stan współdzielony\quad $\longrightarrow$\quad
\code{future.get()}}
\end{center}

\begin{lstlisting}[caption={Ręczne przekazanie wyniku z wątku (C++11).}]
#include <future>
#include <iostream>
#include <thread>
#include <utility>

void compute(std::promise<int> producer) {
    producer.set_value(6 * 7);
}

int main() {
    std::promise<int> promise;
    std::future<int> result = promise.get_future();

    std::thread worker(compute, std::move(promise));
    std::cout << result.get() << '\n'; // 42
    worker.join();
}
\end{lstlisting}

\code{future.get()} czeka, jeśli wynik nie jest jeszcze gotowy, po czym
odbiera go (lub ponownie zgłasza wyjątek). Zwykłe \code{future}
jest jednokrotnego odbioru: po \code{get()} obiekt przestaje być
ważnym uchwytem do stanu i ponowne \code{get()} nie jest dozwolone.
\code{wait()} tylko czeka, natomiast \code{wait_for()} i \code{wait_until()}
pozwalają ograniczyć czas oczekiwania.

\subsection{Przekazywanie wyjątku i \emph{broken promise}}

Producent może przekazać wyjątek za pomocą \code{set_exception()}.
Jeśli \code{promise} zostanie zniszczony bez przekazania wyniku,
stan otrzymuje błąd \code{broken_promise}, który zostanie zgłoszony
podczas odbioru przez \code{future.get()}.

\begin{lstlisting}[caption={Przekazanie wyjątku przez promise (C++14).}]
#include <exception>
#include <future>
#include <iostream>
#include <stdexcept>
#include <thread>
#include <utility>

int main() {
    std::promise<int> producer;
    auto f = producer.get_future();

    std::thread worker([p = std::move(producer)]() mutable {
        try {
            throw std::runtime_error("Calculation failed");
        } catch (...) {
            p.set_exception(std::current_exception());
        }
    });

    try {
        std::cout << f.get() << '\n';
    } catch (const std::exception& e) {
        std::cout << e.what() << '\n';
    }
    worker.join();
}
\end{lstlisting}

\subsection{Wiele odbiorców: \texttt{shared\_future}}

\code{std::shared_future<T>} można kopiować i wielokrotnie odczytywać.
Każdy odbiorca uzyskuje dostęp do tego samego wyniku. Uzyskujemy go
np. przez \code{future.share()}.

\begin{lstlisting}[caption={Dwa wątki odczytują ten sam wynik (C++11).}]
#include <future>
#include <iostream>
#include <thread>

int main() {
    std::promise<int> p;
    std::shared_future<int> f = p.get_future().share();
    int a = 0, b = 0;

    std::thread first([&] { a = f.get(); });
    std::thread second([&] { b = f.get(); });

    p.set_value(42);
    first.join();
    second.join();
    std::cout << a << ' ' << b << '\n'; // 42 42
}
\end{lstlisting}

\section{Uruchamianie zadań: \texttt{std::async}}
\label{sec:async}

\code{std::async} wiąże wywołanie funkcji ze stanem asynchronicznym
i zwraca \code{std::future<R>} dla typu wyniku \code{R}.
W odróżnieniu od \code{promise} nie trzeba samodzielnie obsługiwać
przekazywania wartości ani wyjątków. Wariant
\code{std::launch::async} uruchamia funkcję w innym wątku.

\begin{lstlisting}[caption={Dwa niezależne obliczenia przez async (C++11).}]
#include <future>
#include <iostream>

int square(int x) { return x * x; }

int main() {
    auto f1 = std::async(std::launch::async, square, 6);
    auto f2 = std::async(std::launch::async, square, 7);

    // Oba zadania wystartowaly przed pobraniem wynikow.
    std::cout << f1.get() + f2.get() << '\n'; // 85
}
\end{lstlisting}

\subsection{Polityki uruchamiania}

\begin{center}
\begin{tabularx}{\textwidth}{@{}lX@{}}
\toprule
\textbf{Polityka} & \textbf{Znaczenie} \\
\midrule
\code{std::launch::async} & Funkcja zostaje uruchomiona w osobnym wątku; wywołujący może kontynuować pracę. \\
\code{std::launch::deferred} & Funkcja zostaje wykonana dopiero przy pierwszym nieczasowym oczekiwaniu (np. \code{get()} lub \code{wait()}), w wątku, który tę operację wykonuje. \\
Brak jawnej polityki & Implementacja może wybrać \code{async}, \code{deferred} (lub dopuszczoną politykę implementacyjną). \\
\bottomrule
\end{tabularx}
\end{center}

Nie wolno utożsamiać samego wywołania \code{std::async(f)}
z gwarancją równoległości. Jeśli potrzebujesz niezależnego wątku,
wybierz jawnie \code{std::launch::async}.

\begin{lstlisting}[caption={Opóźnione wykonanie: deferred (C++11).}]
#include <future>
#include <iostream>

int main() {
    auto f = std::async(std::launch::deferred, [] {
        std::cout << "Executing the task\n";
        return 42;
    });
    std::cout << "Before get\n";
    std::cout << f.get() << '\n';
}
\end{lstlisting}

W tym przypadku komunikat \code{Before get} zawsze poprzedza wykonanie
funkcji odroczonej, ponieważ została ona uruchomiona dopiero przez \code{get()}.

\subsection{Destruktor tymczasowego \texttt{future} może blokować}

Istotna pułapka polega na tym, że zniszczenie \emph{ostatniego} uchwytu
do niezakończonego stanu utworzonego przez \code{std::async}
z polityką \code{async} może czekać na ukończenie zadania.
Zapis poniżej nie uruchamia efektywnie obu prac obok siebie:

\begin{lstlisting}[caption={Pułapka: tymczasowe future; fragment ilustracyjny.}]
// Kazdy tymczasowy future jest niszczony na koncu instrukcji.
std::async(std::launch::async, expensive_a);
std::async(std::launch::async, expensive_b);
\end{lstlisting}

Poprawnie przechowujemy wyniki obu wywołań, zanim zaczniemy czekać:

\begin{lstlisting}[caption={Zadania mogą pracować równolegle; fragment.}]
auto a = std::async(std::launch::async, expensive_a);
auto b = std::async(std::launch::async, expensive_b);
a.wait();
b.wait();
\end{lstlisting}

\subsection{Oczekiwanie z limitem czasu i polityka deferred}

\code{wait_for()} zwraca jeden z trzech statusów:
\code{ready}, \code{timeout} lub \code{deferred}.
Przy zadaniu odroczonym samo \code{wait_for()} nie uruchamia pracy.
Niepoprawna pętla sprawdzająca wyłącznie \code{ready}
może więc kręcić się bez końca.

\begin{lstlisting}[caption={Rozróżnianie statusów future (C++11).}]
#include <chrono>
#include <future>
#include <iostream>

int main() {
    auto f = std::async(std::launch::deferred, [] { return 42; });
    auto status = f.wait_for(std::chrono::milliseconds(1));

    if (status == std::future_status::deferred) {
        std::cout << "Deferred: get starts the computation\n";
    } else if (status == std::future_status::timeout) {
        std::cout << "Not ready yet\n";
    } else {
        std::cout << "Already ready\n";
    }
    std::cout << f.get() << '\n';
}
\end{lstlisting}

\subsection{Wyjątki z zadań i brak automatycznego anulowania}

Jeśli funkcja uruchomiona przez \code{std::async} zgłosi wyjątek,
wyjątek jest przechowywany w stanie asynchronicznym, a \code{get()}
zgłasza go w wątku odbierającym wynik. \code{wait_for()}
\emph{nie anuluje} pracy przy przekroczeniu czasu: tylko informuje,
że wynik jeszcze nie jest gotowy.

\begin{lstlisting}[caption={Wyjątek z async zostaje przekazany przez get() (C++11).}]
#include <future>
#include <iostream>
#include <stdexcept>

int main() {
    auto f = std::async(std::launch::async, []() -> int {
        throw std::runtime_error("Task failed");
    });
    try {
        std::cout << f.get() << '\n';
    } catch (const std::runtime_error& e) {
        std::cout << e.what() << '\n';
    }
}
\end{lstlisting}

\section{\texttt{std::packaged\_task}: funkcja z wynikiem future}
\label{sec:packaged}

\code{std::packaged_task<R(Args...)>} opakowuje funkcję w zadanie,
którego wynik lub wyjątek trafia do powiązanego \code{future}.
W odróżnieniu od \code{std::async} samo utworzenie
\code{packaged_task} nie uruchamia żadnego wątku.

\begin{lstlisting}[caption={Packaged task i ręcznie uruchomiony wątek (C++11).}]
#include <future>
#include <iostream>
#include <thread>
#include <utility>

int main() {
    std::packaged_task<int(int)> task([](int x) {
        return x * x;
    });
    auto f = task.get_future();
    std::thread worker(std::move(task), 7);

    std::cout << f.get() << '\n'; // 49
    worker.join();
}
\end{lstlisting}

To jeden z elementów potrzebnych do zbudowania puli wątków:
\emph{task} trafia do kolejki, wątek roboczy go pobiera,
a klient otrzymuje \code{future}. Pula wątków jest szczególnie
przydatna dla bardzo wielu małych zadań, ponieważ tworzenie
odrębnego wątku dla każdego zadania może być kosztowne.

\section{\texttt{std::jthread} i kooperacyjne zatrzymywanie (C++20)}
\label{sec:jthread}

\code{std::jthread} jest obiektem zarządzającym wątkiem zgodnie z RAII.
Jeżeli jest dołączalny, destruktor żąda zatrzymania przez
\code{request_stop()}, a następnie wywołuje \code{join()}.
Żądanie zatrzymania nie wymusza przerwania funkcji: wątek musi
sam sprawdzać odpowiedni \code{stop_token}.

\begin{lstlisting}[caption={Kooperacyjne zatrzymywanie wątku (C++20).}]
#include <chrono>
#include <iostream>
#include <thread>

int main() {
    std::jthread worker([](std::stop_token stop) {
        while (!stop.stop_requested()) {
            // Przykladowa porcja pracy.
            std::this_thread::sleep_for(std::chrono::milliseconds(5));
        }
        std::cout << "Stopped\n";
    });

    std::this_thread::sleep_for(std::chrono::milliseconds(25));
    worker.request_stop();
    // Destruktor worker dolaczy watek.
}
\end{lstlisting}

Można również przekazywać \code{stop_token} do zadań, które
reagują na anulowanie w większej liczbie punktów. Dla złożonych
kolejek C++20 oferuje \code{condition_variable_any} z przeciążeniem
obsługującym żądanie zatrzymania; zwykłe \code{condition_variable}
nie ma analogicznego przeciążenia ze \code{stop_token}.

\section{Inne przydatne mechanizmy biblioteki standardowej}

\subsection{\texttt{shared\_mutex}: wielu czytelników, jeden pisarz}

\code{std::shared_mutex} (C++17) pozwala wielu czytelnikom jednocześnie
uzyskać blokadę współdzieloną, a pisarzowi --- blokadę wyłączną.
Czytanie chronimy przez \code{std::shared_lock}, zapisy przez
\code{std::unique_lock}. Nie gwarantuje to braku zagłodzenia wątków;
polityka przydziału blokad może zależeć od implementacji.

\begin{lstlisting}[caption={Blokada współdzielona i wyłączna (C++17).}]
#include <iostream>
#include <mutex>
#include <shared_mutex>

int main() {
    std::shared_mutex m;
    int value = 10;
    {
        std::shared_lock<std::shared_mutex> read_lock(m);
        std::cout << value << '\n';
    }
    {
        std::unique_lock<std::shared_mutex> write_lock(m);
        value = 20;
    }
}
\end{lstlisting}

\subsection{\texttt{call\_once} i inicjalizacja jednokrotna}

\code{std::call_once} oraz \code{std::once_flag} pozwalają
wykonać funkcję inicjalizującą co najwyżej raz pomyślnie.
Jeżeli inicjalizacja rzuci wyjątek, późniejsze wywołanie może
ponowić próbę. Lokalna statyczna zmienna ma w C++11
bezpieczną współbieżnie inicjalizację, co często jest prostszą alternatywą.

\begin{lstlisting}[caption={Inicjalizacja wykonywana raz (C++11).}]
#include <iostream>
#include <mutex>
#include <thread>

std::once_flag init_flag;
int shared_value = 0;

void initialize() { shared_value = 42; }
void work() {
    std::call_once(init_flag, initialize);
    // Kazdy watek po call_once widzi zakonczona inicjalizacje.
}

int main() {
    std::thread a(work), b(work);
    a.join();
    b.join();
    std::cout << shared_value << '\n';
}
\end{lstlisting}

\subsection{\texttt{thread\_local}: osobny egzemplarz dla wątku}

Zmienne zadeklarowane \code{thread_local} mają osobny egzemplarz
w każdym wątku. Nie rozwiązują automatycznie problemu komunikacji
między wątkami, ale przydają się np. do lokalnych buforów roboczych.

\begin{lstlisting}[caption={Pamięć lokalna dla każdego wątku (C++11).}]
#include <iostream>
#include <thread>

thread_local int local_counter = 0;
void count() { ++local_counter; }

int main() {
    std::thread a(count), b(count);
    a.join();
    b.join();
    std::cout << local_counter << '\n'; // 0 w watku main
}
\end{lstlisting}

\subsection{Semafory, zapadki i bariery w C++20}

\begin{description}[style=nextline]
\item[\code{std::counting_semaphore<N>}]
Ogranicza liczbę równoczesnych wejść. Operacja \code{acquire()}
pobiera zezwolenie, a \code{release()} je zwraca.
\item[\code{std::latch}]
Jednorazowy punkt synchronizacji: po odpowiedniej liczbie zgłoszeń
\code{count_down()} czekające wątki mogą kontynuować.
\item[\code{std::barrier}]
Punkt synchronizacji wielokrotnego użytku dla kolejnych faz pracy.
\end{description}

\begin{lstlisting}[caption={Semafor dopuszcza maksymalnie dwa zadania naraz (C++20).}]
#include <chrono>
#include <iostream>
#include <semaphore>
#include <thread>
#include <vector>

int main() {
    std::counting_semaphore<2> slots(2);
    std::vector<std::jthread> workers;
    for (int i = 0; i < 4; ++i) {
        workers.emplace_back([&, i] {
            slots.acquire();
            // Jednoczesnie w sekcji moze byc najwyzej 2 watki.
            std::this_thread::sleep_for(std::chrono::milliseconds(10));
            slots.release();
        });
    }
    // Destruktory jthread dolaczaja watki przed destruktorem slots.
}
\end{lstlisting}

\nota{W przykładzie z semaforem zmienne zostały zadeklarowane w kolejności,
dzięki której wektor \code{workers} zostaje zniszczony przed \code{slots}.
Ma to znaczenie dla czasu życia obiektów współdzielonych przez wątki.}

% === Rozszerzenie: zadania, pula wątków, anulowanie ===
\section{Zadania, pule wątków i planowanie pracy}
\label{sec:thread-pool}

Wątek jest zasobem systemowym, a zadanie jest jednostką pracy. Wiele
krótkich zadań nie powinno wymagać utworzenia osobnego wątku na każde
wywołanie. \emph{Thread pool} utrzymuje ograniczoną liczbę pracowników,
którzy pobierają zadania ze współdzielonej kolejki.

\subsection{Minimalna pula wątków ze zwracaniem \texttt{future}}

Poniższa implementacja jest przykładem dydaktycznym w C++20.
Łączy kilka poznanych mechanizmów:
\begin{itemize}
\item \code{std::jthread} i \code{std::condition_variable} do obsługi pracowników;
\item \code{std::packaged_task} i \code{std::future} do przekazywania wyników;
\item \code{std::function} do przechowywania zadań w jednej kolejce.
\end{itemize}
Metoda \code{submit} zwraca \code{future} z wynikiem.
Wyjątek zgłoszony przez zadanie zostaje ponowiony podczas \code{get()}.

\begin{lstlisting}[caption={Prosta pula: kompletna implementacja C++20.}]
#include <condition_variable>
#include <cstddef>
#include <functional>
#include <future>
#include <memory>
#include <mutex>
#include <queue>
#include <stdexcept>
#include <thread>
#include <type_traits>
#include <utility>
#include <vector>

class ThreadPool {
    std::mutex m_;
    std::condition_variable cv_;
    std::queue<std::function<void()>> tasks_;
    bool closing_ = false;
    // workers_ jest ostatni: jego destruktory wywolaja join
    // PRZED zniszczeniem m_, cv_ i kolejki.
    std::vector<std::jthread> workers_;

public:
    explicit ThreadPool(std::size_t count) {
        if (count == 0) throw std::invalid_argument("count=0");
        workers_.reserve(count);
        for (std::size_t i = 0; i < count; ++i) {
            workers_.emplace_back([this] {
                for (;;) {
                    std::function<void()> work;
                    {
                        std::unique_lock lock(m_);
                        cv_.wait(lock, [&] {
                            return closing_ || !tasks_.empty();
                        });
                        if (closing_ && tasks_.empty()) return;
                        work = std::move(tasks_.front());
                        tasks_.pop();
                    }
                    work(); // poza sekcja krytyczna
                }
            });
        }
    }

    ThreadPool(const ThreadPool&) = delete;
    ThreadPool& operator=(const ThreadPool&) = delete;

    template<class F>
    auto submit(F&& f)
        -> std::future<std::invoke_result_t<std::decay_t<F>&>> {
        using Result = std::invoke_result_t<std::decay_t<F>&>;
        auto task = std::make_shared<std::packaged_task<Result()>>(
            std::forward<F>(f));
        auto result = task->get_future();
        {
            std::lock_guard lock(m_);
            if (closing_)
                throw std::runtime_error("Pool is closing");
            tasks_.emplace([task] { (*task)(); });
        }
        cv_.notify_one();
        return result;
    }

    ~ThreadPool() {
        {
            std::lock_guard lock(m_);
            closing_ = true;
        }
        cv_.notify_all();
        // Destruktory jthread czekaja na pracownikow.
    }
};
\end{lstlisting}

\begin{lstlisting}[caption={Użycie puli i odbiór wyników; dołącz do programu z klasą ThreadPool.}]
#include <iostream>

int main() {
    ThreadPool pool(3);
    auto a = pool.submit([] { return 40 + 2; });
    auto b = pool.submit([] { return 6 * 7; });
    std::cout << a.get() << " " << b.get() << '\n';
}
\end{lstlisting}

Przykład nie implementuje wszystkich funkcji puli produkcyjnej:
nie ma limitu długości kolejki, priorytetów, \emph{work stealing},
możliwości natychmiastowego porzucenia zadania ani kooperacyjnego
anulowania pojedynczego zadania. Destruktor wykonuje model
\emph{graceful shutdown}: przestaje przyjmować zadania, pozwala
dokończyć prace oczekujące w kolejce i czeka na zakończenie pracowników.
Wywołujący musi zadbać, żeby inne wątki nie używały obiektu puli podczas
jego niszczenia. Nie należy też niszczyć puli z jej własnego pracownika.

\nota{Gdy pula ma jednego pracownika, zadanie oczekujące synchronicznie
na \code{future} innego zadania umieszczonego w \emph{tej samej} puli
może wywołać zakleszczenie. Równie ważna jest kolejność niszczenia
składowych wykorzystywanych przez pracowników.}

\subsection{Anulowanie zadania: prośba, nie wymuszone przerwanie}

\code{std::stop_source} dostarcza możliwość wysłania żądania, a
\code{std::stop_token} pozwala je obserwować. Mogą istnieć liczni
odbiorcy jednego tokenu, a sygnał pozostaje ustawiony po wywołaniu
\code{request_stop()}. Mechanizm sam nie przerywa instrukcji,
blokującego wywołania systemowego ani obliczenia, które ignoruje token.

\begin{lstlisting}[caption={Jedna prośba o zatrzymanie i kilka punktów kontroli.}]
#include <chrono>
#include <iostream>
#include <stop_token>
#include <thread>

void work(std::stop_token token) {
    for (int i = 0; i < 100; ++i) {
        if (token.stop_requested()) {
            std::cout << "Cancelled after " << i << " steps\n";
            return;
        }
        std::this_thread::sleep_for(std::chrono::milliseconds(1));
    }
}

int main() {
    std::stop_source source;
    std::jthread worker(work, source.get_token());
    std::this_thread::sleep_for(std::chrono::milliseconds(5));
    source.request_stop();
}
\end{lstlisting}

Możliwa jest także reakcja w callbacku rejestrowanym przez
\code{std::stop_callback}. Trzeba uważać, by callback nie blokował
mutexa, który jest już trzymany przez wątek wywołujący
\code{request_stop()}, ponieważ callback może zostać wykonany
\emph{synchronicznie} w tym wywołaniu.

\subsection{Czekanie z możliwością zatrzymania}

\code{std::condition_variable_any} ma w C++20 przeciążenie
\code{wait(lock, token, pred)}. Dzięki niemu kolejka zadań może
wybudzać się również po prośbie o zatrzymanie.

\begin{lstlisting}[caption={Oczekiwanie na stan lub anulowanie w C++20.}]
#include <condition_variable>
#include <mutex>
#include <stop_token>

std::mutex m;
std::condition_variable_any cv;
bool ready = false;

void wait_for_work(std::stop_token token) {
    std::unique_lock lock(m);
    bool completed = cv.wait(lock, token, [&] { return ready; });
    if (!completed) {
        // Stop requested zanim predykat stal sie prawdziwy.
        return;
    }
    // W tym miejscu trzymamy m, a ready jest prawdziwe.
}
\end{lstlisting}

Sam token nie zastępuje sygnalizacji \code{cv.notify_one()} po zmianie
predykatu. Wykorzystanie tokenów nie zwalnia także z wymagania, by
stan zadania i jego zasoby pozostawały ważne aż do zakończenia pracy.

\subsection{Wiele zadań i wyjątki: kiedy użyć \texttt{async}, a kiedy puli?}

\code{std::async(std::launch::async, f)} jest dobry, gdy liczba zadań
jest mała, a każde zadanie względnie kosztowne. Pula sprawdza się,
gdy zadań jest wiele i nie chcemy tworzyć osobnego wątku dla każdego.
Z kolei \code{std::jthread} jest wygodny, gdy potrzebujemy długowiecznego
pracownika wykonującego pętlę. W praktyce można też użyć gotowego
\emph{executora} lub biblioteki wyspecjalizowanej w planowaniu zadań.

Gdy funkcja w wątku rzuca wyjątek i nie przechwytuje go lokalnie,
program może zakończyć się przez \code{std::terminate}. Wymianę
wyjątków między wątkami organizujemy przez \code{future},
\code{promise::set_exception} lub \code{std::exception_ptr}.

\section{Skalowanie i koszt współbieżności}

Przyspieszenie równoległego algorytmu jest ograniczone nie tylko przez
liczbę rdzeni, ale także przez odsetek części sekwencyjnej i koszty
synchronizacji. Jeśli część $s\in[0,1]$ pracy musi pozostać sekwencyjna,
to idealizowane prawo Amdahla daje
\[
S(p)\le \frac{1}{s+(1-s)/p},
\]
gdzie $p$ to liczba procesorów/rdzeni roboczych. Jest to uproszczony
model pomijający np. pamięć, uruchamianie wątków i nierównomierność pracy.

W praktyce należy rozważyć: narzut utworzenia i obudzenia wątku,
przeciążenie wspólnego mutexa, współdzielenie linii cache
(\emph{false sharing}), ograniczoną przepustowość pamięci i liczbę
rzeczywiście fizycznych rdzeni. \code{std::thread::hardware_concurrency()}
może być użyteczną wskazówką, ale nie gwarantuje liczby dostępnych
rdzeni i może zwrócić zero.

Równoległość nie jest synonimem wydajności: drobne zadania,
trudne do przewidzenia skoki oraz intensywna konkurencja o pamięć
mogą działać wolniej od wariantu jednowątkowego. Pomiary należy
wykonywać na reprezentatywnych danych w konfiguracji Release,
a poprawność sprawdzać niezależnie od benchmarków.



\section{Przykład numeryczny: równoległe sumowanie}
\label{sec:numerical}

Załóżmy, że mamy duży wektor wartości i chcemy policzyć sumę. Możemy
podzielić zakres na dwie rozłączne części, jedną zsumować w zadaniu
\code{std::async}, a drugą w bieżącym wątku.

\begin{lstlisting}[caption={Dwa niezależne fragmenty obliczenia sumy (C++11).}]
#include <future>
#include <iostream>
#include <numeric>
#include <vector>

int main() {
    std::vector<double> values(100000, 1.0);
    auto middle = values.begin() + values.size() / 2;

    auto right = std::async(std::launch::async, [&] {
        return std::accumulate(middle, values.end(), 0.0);
    });

    double left = std::accumulate(values.begin(), middle, 0.0);
    double total = left + right.get();
    std::cout << total << '\n'; // 100000
}
\end{lstlisting}

Odczytywane fragmenty \code{values} pozostają niezmienne aż do zakończenia
obliczeń, więc nie ma wyścigu danych. Trzeba jednak pamiętać, że
\emph{koszt uruchomienia zadania} może przewyższać korzyść, jeżeli
obliczenia są zbyt krótkie. Przy sumach zmiennoprzecinkowych inna
kolejność dodawania może ponadto zmienić wynik z powodu zaokrągleń.

% === Rozszerzenie: OpenMP ===
\section{OpenMP: programowanie równoległe przez dyrektywy}
\label{sec:openmp}

\emph{OpenMP} (\emph{Open Multi-Processing}) to standard API do
programowania równoległego na pamięci współdzielonej, stosowany
w C, C++ i Fortranie. Nie jest pojedynczą klasą ani tradycyjną
biblioteką szablonową C++: łączy \emph{dyrektywy kompilatora},
\emph{bibliotekę funkcji wykonawczych} (\code{omp.h}) i \emph{zmienne
środowiskowe}. Posiada również konstrukcje do pracy na akceleratorach,
lecz tutaj skupimy się na wielordzeniowym CPU.

Zaletą OpenMP jest możliwość stopniowego równoleglenia istniejących pętli.
Kod źródłowy pozostaje zazwyczaj podobny do wersji sekwencyjnej.
Ceną jest konieczność rozumienia prywatności zmiennych, zależności
pomiędzy iteracjami i synchronizacji.

\nota{Specyfikacja OpenMP ma własne wersje (m.in. 5.2 i 6.0),
niezależne od C++20. W przykładach poniżej używamy głównie
klasycznych dyrektyw dostępnych w wielu kompilatorach. Samo istnienie
mechanizmu w specyfikacji nie gwarantuje, że konkretna wersja GCC,
Clang, MSVC lub Intel oneAPI obsługuje go w pełni.}

\subsection{Kompilacja i konfiguracja OpenMP}

Dla GCC i typowych konfiguracji Clang można użyć flagi
\code{-fopenmp}. W Clangu konieczne może być dodatkowo
zainstalowanie odpowiedniego środowiska wykonawczego OpenMP
(np. \code{libomp}). W MSVC wsparcie i zestaw flag zależą od
wersji oraz wybranego trybu OpenMP.

\begin{lstlisting}[language=bash,caption={Kompilacja programu OpenMP w GCC.}]
g++ -std=c++20 -O2 -Wall -Wextra -fopenmp example.cpp -o example
./example

# Przykladowe uruchomienie z 4 watkami:
OMP_NUM_THREADS=4 ./example
\end{lstlisting}

W CMake można uniknąć ręcznego wpisywania flag:
\begin{lstlisting}[language=bash,caption={Minimalny CMakeLists.txt (tekst pliku).}]
cmake_minimum_required(VERSION 3.16)
project(OpenMPDemo LANGUAGES CXX)
set(CMAKE_CXX_STANDARD 20)
find_package(OpenMP REQUIRED COMPONENTS CXX)
add_executable(example example.cpp)
target_link_libraries(example PRIVATE OpenMP::OpenMP_CXX)
\end{lstlisting}

Biblioteka wykonawcza udostępnia funkcje do odczytu stanu środowiska:
\begin{itemize}
\item \code{omp_get_thread_num()} --- identyfikator wątku;
\item \code{omp_get_num_threads()} --- liczba uczestników aktualnego zespołu;
\item \code{omp_get_max_threads()} --- górna granica następnego zespołu;
\item \code{omp_set_num_threads()} --- ustawianie żądanej liczby wątków;
\item \code{omp_get_wtime()} --- czas ścienny w sekundach.
\end{itemize}
Poza regionem równoległym \code{omp_get_num_threads()} zwraca~1.
Wartość \code{omp_get_max_threads()} nie jest liczbą fizycznych rdzeni.

\subsection{Pierwszy region: \texttt{parallel}}

\code{\#pragma omp parallel} tworzy zespół wątków, którego uczestnicy
wykonują blok. Wątek napotykający dyrektywę staje się wątkiem
głównym tego zespołu (numer~0). Zakończenie bloku obejmuje domyślną
synchronizację uczestników.

\begin{lstlisting}[caption={Region równoległy OpenMP.}]
#include <iostream>
#include <omp.h>

int main() {
    #pragma omp parallel num_threads(4)
    {
        int id = omp_get_thread_num();
        int count = omp_get_num_threads();
        #pragma omp critical
        {
            std::cout << "Thread " << id
                      << " of " << count << '\n';
        }
    }
    std::cout << "Region finished\n";
}
\end{lstlisting}

\code{num_threads(4)} jest żądaniem liczby wątków, a faktyczna liczba
może zależeć od ograniczeń środowiska wykonawczego. \code{critical}
chroni tu przed przeplataniem fragmentów wypisywanego tekstu,
nie narzuca kolejności identyfikatorów wątków.

\subsection{Równoległa pętla \texttt{parallel for}}

Najczęstsze zastosowanie OpenMP polega na podziale \emph{niezależnych}
iteracji jednej pętli między uczestników zespołu.

\begin{lstlisting}[caption={Wypełnianie wektora: niezależne iteracje.}]
#include <cstddef>
#include <vector>

int main() {
    constexpr long long n = 100000;
    std::vector<double> squares(static_cast<std::size_t>(n));

    #pragma omp parallel for default(none) shared(squares, n) \
        schedule(static)
    for (long long i = 0; i < n; ++i) {
        squares[static_cast<std::size_t>(i)] =
            static_cast<double>(i) * i;
    }
}
\end{lstlisting}

\code{default(none)} wymaga jawnego określenia udziału zmiennych,
które nie mają z góry ustalonych atrybutów współdzielenia. Jest
wartościowym narzędziem kontroli poprawności. Zmienna sterująca
pętlą \code{i} jest prywatna na mocy reguł dyrektywy \code{for}.
Każdy wątek zapisuje inny element wektora; nie ma konfliktu,
zakładając że nikt nie zmienia równolegle rozmiaru ani struktury wektora.

\subsection{\texttt{shared}, \texttt{private}, \texttt{firstprivate}, \texttt{lastprivate}}

\begin{description}[style=nextline]
\item[\code{shared(x)}] Wątki odnoszą się do tego samego obiektu.
Dostępy konfliktowe wymagają ochrony lub podziału danych.
\item[\code{private(x)}] Każdy uczestnik ma własny egzemplarz
(z zachowaniem reguł inicjalizacji odpowiednich typów);
nie należy zakładać skopiowania początkowej wartości zmiennej pierwotnej.
\item[\code{firstprivate(x)}] Nowy egzemplarz zostaje zainicjalizowany
wartością odpowiadającej mu zmiennej sprzed regionu.
\item[\code{lastprivate(x)}] Po regionie wynik z logicznie ostatniej
iteracji (niekoniecznie wykonanej na końcu w czasie rzeczywistym)
zostaje przypisany do zmiennej pierwotnej.
\end{description}

\begin{lstlisting}[caption={Przykład firstprivate i lastprivate w OpenMP.}]
#include <iostream>

int main() {
    int seed = 10;
    int last = -1;

    #pragma omp parallel for firstprivate(seed) lastprivate(last)
    for (int i = 0; i < 8; ++i) {
        int local = seed + i;
        last = local;
    }
    std::cout << last << '\n'; // 17
}
\end{lstlisting}

Zmienna \code{seed} jest kopiowana do prywatnych egzemplarzy.
Wartość \code{last} po zakończeniu odpowiada iteracji $i=7$,
nie zaś temu wątkowi, który jako ostatni fizycznie zakończył pracę.

\subsection{Wyścig w redukcji i klauzula \texttt{reduction}}

Poniższa konstrukcja jest \textbf{błędna}:
\begin{lstlisting}[caption={NIEPOPRAWNIE: konfliktowe modyfikacje sumy.}]
// Fragment celowo bledny:
// double sum = 0;
// #pragma omp parallel for
// for (int i = 0; i < n; ++i) sum += values[i];
\end{lstlisting}

Każdy wątek modyfikuje wspólny \code{sum}. Możemy chronić dodawanie
blokadą, lecz w klasycznym sumowaniu wygodniejsze są prywatne sumy,
łączone pod koniec obliczeń:

\begin{lstlisting}[caption={Redukcja równoległa w OpenMP.}]
#include <iostream>
#include <vector>

int main() {
    std::vector<double> values(1000000, 1.0);
    const long long n = static_cast<long long>(values.size());
    double sum = 0.0;

    #pragma omp parallel for reduction(+:sum) schedule(static)
    for (long long i = 0; i < n; ++i) {
        sum += values[static_cast<std::size_t>(i)];
    }
    std::cout << sum << '\n'; // 1000000
}
\end{lstlisting}

OpenMP tworzy prywatne wartości redukcji i łączy je operatorem
\code{+}. Dostępne są również inne redukcje, np. \code{*},
\code{min} i \code{max} (zależnie od typu oraz konstrukcji).
Ważne: dla liczb zmiennoprzecinkowych kolejność działań może być
inna przy różnych zespołach i harmonogramach. Wyniki mogą zatem
różnić się na ostatnich cyfrach, bo
\[
 (a+b)+c\ne a+(b+c)
\]
może zachodzić w arytmetyce IEEE~754. Dla algorytmów numerycznych
mających wymagania reprodukowalności nie wystarczy zakładać
identycznych bitowo wyników każdej redukcji.

\subsection{Harmonogramowanie iteracji: \texttt{schedule}}

\begin{center}
\small
\begin{tabularx}{\textwidth}{@{}p{0.29\textwidth}X@{}}
\toprule
\textbf{Tryb} & \textbf{Znaczenie} \\
\midrule
\code{static} & Przydział iteracji ustalany na początku; mały narzut, dobry dla równych kosztów. \\
\code{static,k} & Przydział kolejnych porcji po $k$ iteracji w sposób cykliczny. \\
\code{dynamic,k} & Porcje po $k$ iteracji przydzielane pracownikom w trakcie pracy; dobre dla nierównych obciążeń, większy narzut. \\
\code{guided,k} & Rozmiary przydzielanych porcji zasadniczo maleją, do narzuconego minimum; kompromis między narzutem a równoważeniem. \\
\code{runtime} & Ustawienia pobierane z środowiska, m.in. \code{OMP_SCHEDULE}. \\
\bottomrule
\end{tabularx}
\end{center}

\begin{lstlisting}[caption={Równoważenie kosztów obliczeń w pętli.}]
#pragma omp parallel for schedule(dynamic, 8)
for (int i = 0; i < n; ++i) {
    perform_variable_cost_work(i);
}
\end{lstlisting}

Przy równomiernych obliczeniach \code{static} bywa tańszy od
\code{dynamic}. Dla zadań nieregularnych dynamiczny harmonogram
zmniejsza bezczynność pracowników, ale może pogorszyć lokalność pamięci.
Doboru dokonujemy na podstawie pomiarów, nie wyłącznie intuicji.

\subsection{Synchronizacja: \texttt{critical}, \texttt{atomic}, \texttt{barrier}, \texttt{nowait}}

\begin{description}[style=nextline]
\item[\code{critical}] Chroni region wykonywany jednocześnie najwyżej
przez jeden wątek w danej grupie nazwanych regionów krytycznych.
\item[\code{atomic}] Zapewnia atomową obsługę wspieranego wyrażenia
(np. inkrementacji); nie obejmuje dowolnego bloku instrukcji.
\item[\code{barrier}] Wymaga, żeby uczestnicy zespołu dotarli do
punktu synchronizacji, zanim którykolwiek przejdzie dalej.
\item[\code{nowait}] Usuwa domyślną barierę tam, gdzie klauzula jest
dopuszczalna; wymaga osobnego rozumowania o zależnościach danych.
\end{description}

\begin{lstlisting}[caption={Dwie alternatywy zwiększania licznika.}]
int counter = 0;

#pragma omp parallel for
for (int i = 0; i < 10000; ++i) {
    #pragma omp atomic update
    ++counter;
}
// counter == 10000
\end{lstlisting}

Dla tej operacji \code{reduction(+:counter)} zwykle ogranicza
rywalizację lepiej niż atomowa inkrementacja przy każdej iteracji.
Z kolei \code{critical} jest przydatny, gdy atomowo musi wykonać się
cały fragment kodu, a nie tylko jedno proste wyrażenie.

\begin{lstlisting}[caption={Bariera wewnątrz istniejącego zespołu.}]
#pragma omp parallel
{
    prepare_local_data();
    #pragma omp barrier
    consume_results_after_everyone_prepared();
}
\end{lstlisting}

Nie można wstawiać bariery w taki sposób, aby część wątków zespołu
ją pomijała, a pozostałe na niej czekały. Podobnie \code{nowait}
nie może być dodawane mechanicznie, jeśli następna pętla potrzebuje
wyników pierwszej.

\subsection{\texttt{sections} i \texttt{single}}

Gdy mamy kilka niezależnych, różnych funkcji, wygodne bywa
\code{sections}; gdy chcemy pojedynczego producenta zadań,
wybieramy \code{single}:

\begin{lstlisting}[caption={Kilka różnych zadań i jedna sekcja inicjalizacyjna.}]
#pragma omp parallel sections
{
    #pragma omp section
    load_input();

    #pragma omp section
    preprocess();
}

#pragma omp parallel
{
    #pragma omp single
    initialize_once();
    // Wszystkie watki mijaja koniec single po barierze.
}
\end{lstlisting}

Powyższy przykład \code{sections} ma sens tylko wtedy, gdy
\code{load_input()} i \code{preprocess()} są rzeczywiście niezależne.
W przeciwnym razie potrzebna jest synchronizacja albo zadania
z opisanymi zależnościami.

\subsection{Zadania OpenMP: \texttt{task}, \texttt{taskwait}, \texttt{taskgroup}}

Podział pętli \code{for} pasuje do regularnych obliczeń. Dla pracy
nieregularnej, np. rekursywnego przetwarzania drzewa, używa się
\code{task}. Dyrektywa tworzy jednostkę pracy, którą środowisko
wykonawcze może przydzielić pracownikowi zespołu. Nie jest to
równoznaczne z utworzeniem nowego wątku systemowego.

\begin{lstlisting}[caption={Tworzenie i oczekiwanie na zadania OpenMP.}]
#include <iostream>

int main() {
    int results[4]{};
    #pragma omp parallel
    {
        #pragma omp single
        {
            for (int i = 0; i < 4; ++i) {
                #pragma omp task firstprivate(i) shared(results)
                {
                    results[i] = i * i;
                }
            }
            #pragma omp taskwait
        }
    }
    for (int x : results) std::cout << x << ' ';
}
\end{lstlisting}

\code{taskwait} oczekuje na odpowiednie bezpośrednie zadania potomne
bieżącego zadania. Gdy chcemy objąć również ich potomków,
praktyczną konstrukcją jest \code{taskgroup}. Koniec regionu
\code{parallel} również zapewnia zakończenie prac należących
do tej struktury zadaniowej.

\subsection{Zależności pomiędzy zadaniami: \texttt{depend}}

\begin{lstlisting}[caption={Zadanie drugie korzysta z wyniku pierwszego.}]
#include <cassert>

int main() {
    int a = 0, b = 0;
    #pragma omp parallel
    {
        #pragma omp single
        {
            #pragma omp task shared(a) depend(out:a)
            {
                a = 21;
            }
            #pragma omp task shared(a,b) depend(in:a)
            {
                b = a * 2;
            }
            #pragma omp taskwait
        }
    }
    assert(b == 42);
}
\end{lstlisting}

Klauzule \code{depend(in:...)}, \code{depend(out:...)} i
\code{depend(inout:...)} opisują zależności danych m.in.
pomiędzy zadaniami rodzeństwa. Nie zastępują synchronizacji
przy dowolnych dostępach do innych, nieopisanych zasobów.

\subsection{\texttt{collapse} i \texttt{simd}: pętle wielowymiarowe i wektoryzacja}

\code{collapse(2)} łączy przestrzeń iteracji dwóch odpowiednio
zagnieżdżonych pętli, zwiększając liczbę możliwych porcji pracy.
\code{simd} natomiast przekazuje kompilatorowi informację,
że iteracje dają się wykonać w modelu SIMD. Wektoryzacja i
równoległość wielowątkowa są różnymi poziomami przyspieszania.

\begin{lstlisting}[caption={Tablica 2D: parallel for z collapse(2).}]
#include <vector>

int main() {
    const int rows = 400, cols = 400;
    std::vector<double> a(rows * cols);
    #pragma omp parallel for collapse(2) schedule(static)
    for (int i = 0; i < rows; ++i)
        for (int j = 0; j < cols; ++j)
            a[i * cols + j] = i + 0.5 * j;
}
\end{lstlisting}

\begin{lstlisting}[caption={Wektoryzacja bez obowiązkowego tworzenia wielu wątków.}]
void scale(double* dest, const double* src, int n, double alpha) {
    #pragma omp simd
    for (int i = 0; i < n; ++i)
        dest[i] = alpha * src[i];
}
\end{lstlisting}

Prawidłowość \code{simd} wymaga faktycznej niezależności iteracji,
w tym braku niedozwolonego nakładania się zapisów i odczytów.
Nie należy wprowadzać dyrektywy w kodzie o nieustalonym aliasingu
bez wcześniejszej analizy zależności.

\subsection{Zmienne środowiskowe OpenMP i diagnostyka}

\begin{center}
\small
\begin{tabularx}{\textwidth}{@{}p{0.3\textwidth}X@{}}
\toprule
\textbf{Zmienna} & \textbf{Rola} \\
\midrule
\code{OMP_NUM_THREADS} & Żądana liczba wątków w zespołach. \\
\code{OMP_DYNAMIC} & Pozwala środowisku dynamicznie dostosowywać liczbę wątków. \\
\code{OMP_SCHEDULE} & Harmonogram stosowany przy \code{schedule(runtime)}. \\
\code{OMP_PROC_BIND} & Strategia wiązania wątków z zasobami obliczeniowymi. \\
\code{OMP_PLACES} & Definiuje miejsca wykonania dla mechanizmu wiązania. \\
\code{OMP_MAX_ACTIVE_LEVELS} & Ogranicza liczbę aktywnych poziomów zagnieżdżonych regionów. \\
\bottomrule
\end{tabularx}
\end{center}

Przykładowa sesja uruchomieniowa:
\begin{lstlisting}[language=bash,caption={Sterowanie środowiskiem bez zmiany źródeł.}]
OMP_NUM_THREADS=8 OMP_SCHEDULE=dynamic,32 ./example
OMP_NUM_THREADS=4 OMP_PROC_BIND=close ./example
\end{lstlisting}

Rzeczywista dostępność i interpretacja opcji dotyczących wiązania
wątków zależą od środowiska wykonawczego i maszyny. Zagnieżdżone
regiony OpenMP oraz łączenie OpenMP z \code{std::thread}
może spowodować \emph{oversubscription}: więcej aktywnych
wątków roboczych niż zasobów obliczeniowych.

\subsection{OpenMP a model pamięci C++}

OpenMP określa własny model pamięci, operację \emph{flush} oraz
semantykę synchronizacji dyrektyw, takich jak \code{barrier},
\code{critical} i \code{atomic}. \emph{Flush} nie jest funkcją
\code{std::atomic_thread_fence}, a samo słowo ,,flush'' nie znaczy
,,wyczyść wszystkie cache procesora''.

W praktyce warto trzymać się jednej, konsekwentnej reguły:
\begin{itemize}
\item dane dzielone między iteracje OpenMP synchronizujemy przede
wszystkim przez odpowiednie konstrukcje OpenMP;
\item dane wymieniane pomiędzy własnymi wątkami C++ synchronizujemy
przez \code{std::mutex}, \code{std::atomic} lub inne narzędzia C++;
\item przy \emph{mieszaniu} obu modeli wymagane jest szczególnie
staranne rozumowanie o zgodności mechanizmów, zamiast zakładania,
że przypadkowe \code{flush} naprawi wyścig danych.
\end{itemize}

Wyjątków C++ nie należy wypuszczać poza regiony OpenMP w sposób
naruszający zasady specyfikacji. Jeżeli zadanie może zgłosić wyjątek,
trzeba obsłużyć go wewnątrz odpowiedniego regionu i przekazać
informację o błędzie w kontrolowany sposób po zakończeniu pracy.

\subsection{OpenMP czy \texttt{std::thread}, \texttt{async}, a może Ranges?}

\begin{center}
\small
\begin{tabularx}{\textwidth}{@{}p{0.28\textwidth}X@{}}
\toprule
\textbf{Mechanizm} & \textbf{Typowy przypadek użycia} \\
\midrule
\code{std::thread}/\code{jthread} & Długowieczny wątek, własny protokół komunikacji, kontrola czasu życia. \\
\code{std::async} & Niewiele niezależnych obliczeń z wynikiem w \code{future}. \\
Pula wątków & Wiele krótkich lub średnich zadań, ograniczona liczba pracowników. \\
OpenMP \code{parallel for} & Obliczenia naukowe i numeryczne z niezależnymi iteracjami. \\
OpenMP \code{task} & Dynamiczne drzewo prac i zależności. \\
\code{std::execution::par} & Algorytmy standardowe C++17+ z odpowiednią polityką wykonania i dostępnym backendem. \\
\bottomrule
\end{tabularx}
\end{center}

Polityka \code{std::execution::par} nie stanowi gwarancji faktycznego
wielowątkowego wykonania na każdym kompilatorze i bibliotece
standardowej. \emph{Ranges views} z C++20 są zazwyczaj przede
wszystkim leniwymi adaptatorami zakresów, nie automatycznym
mechanizmem równoległego planowania.

\subsection{Studium numeryczne: porównanie trzech sposobów sumowania}

Rozważmy trzy warianty: zwykły \code{std::accumulate}, równoległe
sumowanie przez \code{std::async} oraz OpenMP \code{reduction}.
W przykładzie\footnote{Przy dzieleniu dużego wektora na wiele części
należy uwzględnić balans obciążenia, rozmiary porcji i koszt
agregacji wyników.} obliczamy sumę kwadratów wejściowych liczb
przez OpenMP, a następnie porównujemy z referencyjną pętlą:

\begin{lstlisting}[caption={Praktyczny przykład: suma kwadratów OpenMP z pomiarem.}]
#include <cmath>
#include <iostream>
#include <omp.h>
#include <vector>

int main() {
    constexpr int n = 1000000;
    std::vector<double> a(n);
    for (int i = 0; i < n; ++i) a[i] = i * 1e-5;

    double sequential = 0.0;
    double t0 = omp_get_wtime();
    for (double x : a) sequential += x * x;
    double t1 = omp_get_wtime();

    double parallel = 0.0;
    #pragma omp parallel for reduction(+:parallel) schedule(static)
    for (int i = 0; i < n; ++i)
        parallel += a[i] * a[i];
    double t2 = omp_get_wtime();

    std::cout << "sequential: " << sequential
              << ", time: " << t1 - t0 << " s\n";
    std::cout << "OpenMP: " << parallel
              << ", time: " << t2 - t1 << " s\n";
    std::cout << "relative error: "
              << std::abs(sequential - parallel) /
                 (1.0 + std::abs(sequential)) << '\n';
}
\end{lstlisting}

\nota{Ten program jest demonstracją interfejsu, a nie metodologicznie
pełnym benchmarkiem. Należy wykonać rozgrzewkę, wiele powtórzeń,
uwzględnić rozrzut wyników i sprawdzić, czy kompilator nie usunął
obliczeń. Koszt dostępu do pamięci może ograniczać przyspieszenie;
przy małym \code{n} równoległa wersja może być wolniejsza.}



\section{Porównanie i zasady projektowania}
\label{sec:choose}

Po zapoznaniu się z modelem pamięci i OpenMP potrafimy rozróżnić
trzy poziomy pracy: \emph{własność wątku} (\code{thread}/\code{jthread}),
\emph{własność zadania i jego wyniku} (\code{future}, pule) oraz
\emph{deklaratywny podział iteracji/zadań} (OpenMP). Wybór nie
powinien zależeć jedynie od liczby linijek kodu.


\begin{center}
\small
\begin{tabularx}{\textwidth}{@{}lX@{}}
\toprule
\textbf{Narzędzie} & \textbf{Kiedy używać} \\
\midrule
\code{thread} & Potrzebna jest bezpośrednia kontrola nad wątkiem i jego czasem życia. \\
\code{jthread} & W C++20, gdy przydaje się RAII i kooperacyjny sygnał zatrzymania. \\
\code{mutex} & Kilka operacji musi spójnie chronić współdzielony stan. \\
\code{condition_variable} & Wątek powinien zasnąć do spełnienia warunku na stanie. \\
\code{atomic} & Niezależne, atomowe operacje na zmiennej lub starannie zaprojektowana synchronizacja. \\
\code{promise} & Producent sam decyduje, kiedy dostarczyć wynik albo wyjątek. \\
\code{future} & Odbiór jednego wyniku i ewentualne oczekiwanie. \\
\code{shared_future} & Kilku odbiorców potrzebuje tej samej wartości. \\
\code{async} & Chcemy uruchomić funkcję i odebrać jej wynik bez zarządzania wątkiem. \\
\code{packaged_task} & Trzeba oddzielić zadanie od mechanizmu jego uruchamiania. \\
Pula wątków & Dużo zadań i ograniczona liczba pracowników; kontrolowane zamykanie. \\
OpenMP & Duże obliczenia z niezależnymi iteracjami albo grafem zadań. \\
\bottomrule
\end{tabularx}
\end{center}

\subsection{Najczęstsze błędy}

\begin{enumerate}
\item Użycie zwykłej zmiennej współdzielonej bez synchronizacji
(\emph{data race}, niezdefiniowane zachowanie).
\item Założenie, że \code{++x} na zwykłym \code{int} jest atomowe.
\item Zniszczenie aktywnego \code{std::thread} bez \code{join()} lub \code{detach()}.
\item Przekazanie referencji do obiektu, którego czas życia kończy się
przed zakończeniem pracy wątku.
\item Oczekiwanie przez \code{condition_variable} bez predykatu.
\item Mylenie powiadomienia z trwałym zapisem informacji o stanie.
\item Założenie, że domyślne \code{std::async(f)} zawsze uruchamia nowy wątek.
\item Natychmiastowe niszczenie tymczasowego \code{future} z \code{std::async}.
\item Traktowanie \code{wait_for()} jako anulowania zadania.
\item Próba wielokrotnego \code{get()} na zwykłym \code{future}.
\item Nadmierne używanie \code{memory_order_relaxed} bez dowodu poprawności.
\item Zakładanie, że dla każdego drobnego zadania opłaca się uruchomić
osobny wątek.
\item Przenoszenie \code{memory_order_relaxed} do algorytmu publikującego dane
bez analizy zależności między wątkami.
\item Stosowanie \code{\#pragma omp parallel for} do pętli z zależnościami
między iteracjami albo wspólnym niechronionym akumulatorem.
\item Mieszanie własnej puli wątków z OpenMP bez kontroli liczby aktywnych
pracowników; pojawia się ryzyko \emph{oversubscription}.
\end{enumerate}

\subsection{Standard języka a dostępność API}

\begin{center}
\begin{tabularx}{\textwidth}{@{}lX@{}}
\toprule
\textbf{Standard} & \textbf{Wybrane mechanizmy} \\
\midrule
C++11 & \code{thread}, \code{mutex}, \code{condition_variable}, \code{atomic}, \code{promise}, \code{future}, \code{async}, \code{packaged_task}, \code{call_once}, \code{thread_local}. \\
C++14 & M.in. ułatwienia związane z lambdami, np. inicjalizujące przechwycenie użyte w przykładzie \code{promise}. \\
C++17 & \code{scoped_lock}, \code{shared_mutex}, \code{shared_lock} (ta ostatnia od C++14). \\
C++20 & \code{jthread}, \code{stop_token}, \code{atomic::wait}, \code{semaphore}, \code{latch}, \code{barrier}, \code{osyncstream}. \\
\bottomrule
\end{tabularx}
\end{center}

Przykład z inicjalizującym przechwyceniem:
\begin{center}
\code{[p = std::move(producer)]}
\end{center}
wymaga C++14, mimo że sam \code{std::promise} istnieje od C++11.

\section{Kompilacja, sprawdzanie i diagnostyka}

Każdy kompletny przykład można zapisać w osobnym pliku z rozszerzeniem
\code{.cpp}. Dla kodu C++20 typowe wywołanie kompilatora GCC lub Clang to:

\begin{lstlisting}[language=bash,caption={Kompilacja przykładu i uruchomienie programu.}]
g++ -std=c++20 -O2 -Wall -Wextra -pthread example.cpp -o example
./example
\end{lstlisting}

Podczas nauki szczególnie pomocny jest \emph{ThreadSanitizer}, który
może wykrywać wiele wyścigów danych (dostępność zależy od platformy
oraz kompilatora):

\begin{lstlisting}[language=bash,caption={Przykładowe uruchomienie z ThreadSanitizer.}]
g++ -std=c++20 -O1 -g -pthread -fsanitize=thread example.cpp -o example_tsan
./example_tsan
\end{lstlisting}

Nawet jeśli test przejdzie tysiąc razy, nie dowodzi to braku błędów
współbieżności. Poprawność wymaga analizy synchronizacji i czasu życia;
narzędzia diagnostyczne stanowią jej uzupełnienie.

\section*{Podsumowanie rozdziału}
\addcontentsline{toc}{section}{Podsumowanie rozdziału}

Z \code{std::thread} i \code{std::jthread} sterujemy \emph{wątkami},
a z \code{std::async} i \code{std::packaged_task} opisujemy raczej
\emph{zadania}. \code{future}/\code{promise} tworzą interfejs dostarczania
i odbioru jednego wyniku. \code{mutex} i \code{condition_variable}
pozwalają chronić współdzielony stan i czekać na jego zmianę.
\code{atomic} udostępnia operacje atomowe oraz mechanizmy porządkowania
pamięci. Algorytmy \emph{lock-free} wymagają ponadto zarządzania czasem
życia węzłów i analizy problemu ABA. Pula wątków separuje zadania
od fizycznych pracowników, a OpenMP udostępnia deklaratywny podział
pracy na iteracje, redukcje i zadania. Żadne z tych narzędzi nie
zastępuje pozostałych we wszystkich zastosowaniach.

Dobra architektura zaczyna się od pytania, \emph{które dane są współdzielone,
kto jest właścicielem zadania, kiedy wolno zakończyć jego pracę oraz
jak przekazać wynik i wyjątek}. Dopiero potem należy wybrać konkretne
narzędzia biblioteki standardowej.

\section*{Literatura i dokumentacja}
\addcontentsline{toc}{section}{Literatura i dokumentacja}

\begin{itemize}
\item \url{https://cppreference.com/cpp/thread/async} --- \code{std::async} i polityki uruchamiania.
\item \url{https://en.cppreference.com/w/cpp/thread/future} --- \code{std::future}.
\item \url{https://cppreference.com/cpp/thread/condition_variable/wait} --- oczekiwanie z predykatem.
\item \url{https://en.cppreference.com/w/cpp/atomic/atomic/wait} --- \code{atomic::wait} w C++20.
\item \url{https://en.cppreference.com/w/cpp/atomic/memory_order} --- porządek pamięci.
\item \url{https://en.cppreference.com/w/cpp/atomic/atomic/compare_exchange} --- CAS.
\item \url{https://www.openmp.org/specifications/} --- oficjalne wersje i dokumentacja OpenMP (w tym 5.2 i 6.0).
\item \url{https://www.openmp.org/spec-html/5.2/openmp.html} --- treść specyfikacji OpenMP 5.2.
\item \url{https://www.openmp.org/spec-html/5.2/openmpse4.html} --- model pamięci OpenMP.
\item Anthony Williams, \emph{C++ Concurrency in Action}, wyd. 2, Manning, 2019.
\end{itemize}

% ROZDZIAL 15 — Korutyny, zadania i porownanie ze std::async (C++20).
% Fragment do wspolnego skryptu: bez preambuly i bez \begin{document}.
% Wymagane pakiety: listings, amsmath, amssymb, booktabs, tabularx,
% array, hyperref (jak w rozdzialach poprzednich).
% Dolaczenie: \input{rozdzial_15_korutyny_vs_async_cpp.tex}
% Listing oznaczony jako „kompletny program” mozna skompilowac oddzielnie.

\chapter{Korutyny, zadania asynchroniczne i \texttt{std::async}}
\label{chap:cpp15-coroutines}

\section{Po co C++ potrzebuje korutyn?}

Zwykła funkcja ma zasadniczo jeden punkt wejścia. Może wywoływać inne
funkcje, ale po powrocie oddaje sterowanie wywołującemu i nie posiada
mechanizmu późniejszego wznowienia dokładnie w miejscu, w którym
zakończyła wykonywanie. \emph{Korutyna} (ang. \emph{coroutine}) jest
funkcją, której wykonanie można \emph{zawiesić}, a następnie
\emph{wznowić}, zachowując wybrane zmienne lokalne i miejsce kontynuacji.

Korutyny C++20 wspierają między innymi:
\begin{itemize}
\item \emph{generatory} --- wytwarzanie kolejnych wartości na żądanie,
      bez budowania całej kolekcji;
\item \emph{zadania asynchroniczne} --- opis operacji, których kolejne
      etapy uruchamiane są po zakończeniu wcześniejszych;
\item \emph{kooperacyjne planowanie} --- oddawanie sterowania
      harmonogramowi zadań;
\item \emph{maszyny stanów} --- zachowywanie stanu pomiędzy wznowieniami
      bez ręcznego kodowania każdej możliwej kontynuacji.
\end{itemize}

Najważniejsza zasada brzmi:
\begin{center}
\emph{Korutyna nie jest wątkiem, a \texttt{co\_await} nie oznacza
automatycznie ,,uruchom w nowym wątku''.}
\end{center}

Korutyna określa \emph{jak zatrzymać i kontynuować obliczenie};
\texttt{std::async}, \texttt{std::thread} i pula wątków dotyczą również
\emph{gdzie i kiedy} wykonywać pracę. W szczególności kod korutyny
może działać od początku do końca w jednym wątku, nawet jeśli
wielokrotnie się zawiesza. Odwrotnie, korutyna może być wznawiana
przez różne wątki, jeżeli tak zdecyduje dostarczony mechanizm planowania.

\subsection{Trzy słowa kluczowe}

\begin{center}\small
\begin{tabularx}{\textwidth}{@{}p{0.32\textwidth}X@{}}
\toprule
Konstrukcja & Działanie w korutynie \\
\midrule
\code{co_await expr} & Przekazuje sterowanie obiektowi oczekującemu; może zawiesić korutynę bez blokowania wątku. \\
\code{co_yield value} & Udostępnia kolejną wartość i zwykle zawiesza generator do wznowienia. \\
\code{co_return value} & Kończy korutynę i przekazuje wynik zgodnie z protokołem jej typu zwracanego. \\
\code{await_ready()} & Określa, czy wynik jest już gotowy i można pominąć zawieszenie. \\
\code{await_suspend(h)} & Obsługuje zawieszenie i uchwyt \code{h}; może zaplanować późniejsze wznowienie. \\
\code{await_resume()} & Zwraca rezultat oczekiwania albo zgłasza zapisany wyjątek po wznowieniu. \\
\bottomrule
\end{tabularx}
\end{center}

C++20 wprowadził następujące konstrukcje składniowe:
\begin{description}
\item[\texttt{co\_await expr}] Zastosuj protokół oczekiwania do wyrażenia;
      zależnie od niego możesz zawiesić korutynę i później ją wznowić.
\item[\texttt{co\_yield expr}] Udostępnij kolejną wartość generatora;
      składniowo jest to szczególna postać oczekiwania, wykorzystująca
      metodę \texttt{yield\_value} obiektu \texttt{promise\_type}.
\item[\texttt{co\_return expr}] Zakończ korutynę i przekaż wynik
      przez odpowiednie \texttt{return\_value} albo \texttt{return\_void}.
\end{description}

Funkcja zawierająca jeden z tych operatorów (w odpowiednim kontekście)
rozpoznawana jest jako korutyna. Typ jej zwracanej wartości musi
udostępniać właściwy \emph{typ obietnicy} \texttt{promise\_type}
(lub odpowiednią specjalizację \texttt{std::coroutine\_traits}).
To wymóg mechanizmu korutyn, a nie znanego już \texttt{std::promise}.

\begin{lstlisting}[caption={Niepoprawna intuicja: std::future nie jest automatycznie typem korutyny (C++20); fragment celowo błędny.}]
#include <future>

// Standardowe std::future<int> nie zapewnia promise_type:
// std::future<int> calculate() { co_return 42; } // BLAD!
\end{lstlisting}

W C++20 biblioteka standardowa dostarcza \emph{mechanizm} korutyn
(nagłówek \texttt{<coroutine>}), ale nie dostarcza gotowego, ogólnego
\texttt{Task<T>} ani standardowego executora. Takie typy definiuje
biblioteka lub programista. \texttt{std::generator} jest rozwiązaniem
późniejszym, wprowadzonym dopiero w C++23, a jego wsparcie zależy
od użytej biblioteki standardowej.

\section{Różnica między zawieszeniem a zablokowaniem wątku}
\label{sec:cpp15-suspend-block}

Rozważmy dwa sposoby wstrzymania działania programu:
\begin{itemize}
\item \texttt{std::this\_thread::sleep\_for(100ms)} \emph{blokuje
      bieżący wątek} na pewien czas; sam wątek pozostaje zajętym
      zasobem systemowym, choć może nie zużywać wówczas czasu CPU;
\item dobrze zaprojektowane \texttt{co\_await} \emph{zawiesza
      korutynę}, oddając wątek do wykonywania innych zadań;
      wznowienie może nastąpić po nadejściu wyniku operacji I/O,
      komunikatu lub upływie czasu.
\end{itemize}

Nie każde \texttt{co\_await} daje zawieszenie: awaiter może zgłosić,
że wynik jest już gotowy. Również samo użycie \texttt{co\_await}
nie gwarantuje nieblokującego I/O --- zależy to od jego implementacji.

\begin{center}
\small
\begin{tabularx}{\textwidth}{@{}lXX@{}}
\toprule
\textbf{Właściwość} & \textbf{Blokada wątku} & \textbf{Zawieszenie korutyny} \\
\midrule
Zachowanie stanu & Na stosie wywołań wątku & Zwykle w ramce korutyny \\
Wątek potrzebny podczas oczekiwania & Tak, istnieje i czeka & Nie jest konieczny \emph{dla samego oczekiwania} \\
Wznowienie & Planista systemu/odblokowanie & Mechanizm awaitera i właściciel kontynuacji \\
Koszt & M.in. wątek i jego stos & M.in. ramka i obsługa zawieszenia \\
\bottomrule
\end{tabularx}
\end{center}

\section{Ramka korutyny i jej cykl życia}
\label{sec:cpp15-frame}

Kompilator przekształca korutynę w obiekt reprezentujący jej stan.
Ten stan nazywamy \emph{ramką korutyny} (ang. \emph{coroutine frame}).
Może ona obejmować obiekt \texttt{promise\_type}, parametry i stan
potrzebny do kontynuowania obliczenia po zawieszeniu. To, które
zmienne lokalne trzeba zachować, wynika z przebiegu programu.

W ogólnym modelu wywołanie korutyny przebiega w następujących etapach:
\begin{enumerate}
\item Alokowana jest ramka, chyba że implementacja legalnie
      wyeliminuje tę alokację (np. przez zagnieżdżenie ramki w innym obiekcie).
\item Kopiowane lub przenoszone są parametry odpowiednio do ich rodzajów;
      \emph{parametry referencyjne pozostają referencjami}.
\item Konstruowany jest obiekt obietnicy \texttt{promise\_type}.
\item Wywoływana jest metoda \texttt{get\_return\_object()}.
\item Wykonywany jest \texttt{co\_await promise.initial\_suspend()}.
\item Uruchamiane jest ciało korutyny (od razu lub po wznowieniu).
\item \texttt{co\_return} uruchamia \texttt{return\_value} lub
      \texttt{return\_void}; nieobsłużony wyjątek trafia do
      \texttt{unhandled\_exception()}.
\item Niszczone są obiekty lokalne, a korutyna wykonuje końcowy
      \texttt{co\_await} na wyniku \texttt{promise.final\_suspend()}.
\item Ramka zostaje ostatecznie zniszczona we właściwym momencie.
\end{enumerate}

Kolejność i czas \emph{zwolnienia ramki} zależą od przyjętej polityki
własności. Dla \texttt{final\_suspend = suspend\_always} ramkę zwykle
niszczy właściciel uchwytu, nawet po zakończeniu ciała korutyny.
Przy \texttt{final\_suspend = suspend\_never} mechanizm korutyny
może zakończyć się bez pozostawiania ramki zawieszonej do późniejszego
zniszczenia przez właściciela. Tych strategii nie wolno mieszać.

\subsection{Początek i koniec: \texttt{initial\_suspend} i \texttt{final\_suspend}}

Dwa najczęściej używane typy z nagłówka \texttt{<coroutine>} to:
\begin{description}
\item[\texttt{std::suspend\_never}] Oczekiwanie jest natychmiast
      gotowe i nie powoduje zawieszenia.
\item[\texttt{std::suspend\_always}] Oczekiwanie zawsze prowadzi
      do zawieszenia, dopóki ktoś nie wznowi korutyny.
\end{description}

Jeżeli \texttt{initial\_suspend()} zwraca \texttt{suspend\_always},
korutyna jest \emph{lazy}: nie rozpoczyna wykonywania ciała przy samym
wywołaniu funkcji. Natomiast \texttt{suspend\_never} oznacza
\emph{eager start}: ciało może rozpocząć się natychmiast w wątku
wywołującym i pracować aż do pierwszego rzeczywistego zawieszenia.

\textbf{Ważne:} nie wolno wywoływać \texttt{resume()} na korutynie,
która jest zawieszona w swoim \texttt{final\_suspend}. Jest to
niezdefiniowane zachowanie. Uchwyt do takiej ramki może nadal służyć
do \texttt{destroy()}, jeżeli ramka nie została już zniszczona.

\section{Uchwyt \texttt{std::coroutine\_handle} i własność ramki}

\texttt{std::coroutine\_handle<Promise>} jest lekkim uchwytem
umożliwiającym m.in. wznowienie korutyny i dostęp do jej obietnicy.
Sam \texttt{coroutine\_handle} \emph{nie zarządza automatycznie
czasem życia ramki}. Skopiowanie uchwytu nie jest kopią korutyny,
a zniszczenie uchwytu nie wywołuje automatycznie \texttt{destroy()}.
Dlatego wygodna klasa będąca właścicielem ramki powinna stosować RAII,
uniemożliwiać kopiowanie oraz poprawnie obsługiwać przenoszenie.

\begin{lstlisting}[caption={Operacje na uchwycie; fragment ilustracyjny.}]
#include <coroutine>

using Handle = std::coroutine_handle<>;

// Zakladamy, ze h wskazuje na istniejaca, zawieszona korutyne.
// if (h && !h.done()) h.resume();
// h.destroy(); // dokladnie raz, przez wlasciciela ramki
// Po destroy() wszystkie kopie h sa wiszacymi uchwytami.
\end{lstlisting}

Uchwyt może zostać przekazany do innego wątku lub executora, ale
wówczas trzeba zagwarantować, że ramka nie zostanie zniszczona,
zanim planista faktycznie wznowi zawieszoną korutynę. Jest to jedna
z głównych trudności implementacyjnych asynchronicznego API C++20.

\section{Własny generator: \texttt{co\_yield} od podstaw}
\label{sec:cpp15-generator}

Zbudujemy teraz typ \texttt{Generator}, który produkuje liczby całkowite
kolejno, \emph{na żądanie}. Poniższy przykład jest samowystarczalny.
Prywatny uchwyt do ramki stanowi jedynego właściciela zasobu;
\texttt{next()} wznawia korutynę tylko wtedy, gdy nie zakończyła
wykonywania. Własną obsługą wyjątku uzupełniamy przykład o
poprawną propagację błędów.

\begin{lstlisting}[caption={Generator liczb całkowitych: co\_yield i RAII (C++20); kompletny program.}]
#include <cassert>
#include <coroutine>
#include <exception>
#include <limits>
#include <stdexcept>
#include <utility>

class Generator {
public:
    struct promise_type {
        int current = 0;
        std::exception_ptr error;

        Generator get_return_object() {
            return Generator{Handle::from_promise(*this)};
        }
        std::suspend_always initial_suspend() noexcept { return {}; }
        std::suspend_always final_suspend() noexcept { return {}; }
        std::suspend_always yield_value(int x) noexcept {
            current = x;
            return {};
        }
        void return_void() noexcept {}
        void unhandled_exception() noexcept {
            error = std::current_exception();
        }
    };

    using Handle = std::coroutine_handle<promise_type>;

    explicit Generator(Handle h) noexcept : handle_(h) {}
    Generator(const Generator&) = delete;
    Generator& operator=(const Generator&) = delete;

    Generator(Generator&& other) noexcept
        : handle_(std::exchange(other.handle_, {})),
          ready_(std::exchange(other.ready_, false)) {}

    Generator& operator=(Generator&& other) noexcept {
        if (this != &other) {
            if (handle_) handle_.destroy();
            handle_ = std::exchange(other.handle_, {});
            ready_ = std::exchange(other.ready_, false);
        }
        return *this;
    }

    ~Generator() {
        if (handle_) handle_.destroy();
    }

    bool next() {
        ready_ = false;
        if (!handle_ || handle_.done()) return false;
        handle_.resume();
        if (handle_.promise().error)
            std::rethrow_exception(handle_.promise().error);
        ready_ = !handle_.done();
        return ready_;
    }

    int value() const {
        if (!ready_)
            throw std::logic_error("No current value");
        return handle_.promise().current;
    }

private:
    Handle handle_{};
    bool ready_ = false;
};

Generator squares(int n) {
    for (int i = 1; i <= n; ++i) {
        if (i > std::numeric_limits<int>::max() / i)
            throw std::overflow_error("square outside int");
        co_yield i * i;
    }
}

int main() {
    auto gen = squares(4); // Cialo jeszcze sie nie wykonalo.
    int sum = 0;
    while (gen.next()) sum += gen.value();
    assert(sum == 1 + 4 + 9 + 16);
    assert(!gen.next());
}
\end{lstlisting}

\subsection{Co dzieje się podczas \texttt{co\_yield}?}

Instrukcja \texttt{co\_yield i * i} powoduje wywołanie
\texttt{promise.yield\_value(i * i)} i zastosowanie mechanizmu
\texttt{co\_await} do otrzymanego awaitera. W naszym przypadku
\texttt{yield\_value} zapisuje wartość w obietnicy i zwraca
\texttt{suspend\_always}. Funkcja \texttt{next()} otrzymuje wtedy
sterowanie z powrotem i udostępnia zapisany element. Następne
\texttt{next()} kontynuuje pętlę od poprzedniego miejsca zawieszenia.

Generator posiada tylko jeden bieżący wynik. Nie przechowuje
całego ciągu jak \texttt{std::vector<int>}. To użyteczne przy dużych
lub potencjalnie nieskończonych sekwencjach. \emph{Nie oznacza to}
jednak automatycznie mniejszego kosztu obliczeń --- dochodzą koszty
ramki i kolejnych wznowień.

\nota{Przedstawiony \texttt{Generator} jest celowo minimalny: nie
udostępnia interfejsu iteratorów \texttt{begin()/end()}, nie jest
bezpieczny do jednoczesnego wznawiania z kilku wątków i nie powinien
być stosowany jako gotowa biblioteka produkcyjna.}

\section{Protokół \texttt{co\_await}: \texttt{await\_ready},
\texttt{await\_suspend}, \texttt{await\_resume}}
\label{sec:cpp15-awaiter}

Operator \texttt{co\_await} nie jest synonimem oczekiwania na
\texttt{future}. Kompilator uzyskuje obiekt \emph{awaiter}, którego
interfejs ma trzy fundamentalne operacje:
\begin{description}
\item[\texttt{await\_ready()}] Zwraca wartość logiczną: jeśli
      \texttt{true}, zawieszenie jest pomijane.
\item[\texttt{await\_suspend(handle)}] Jeżeli \texttt{await\_ready()}
      było fałszywe, metoda podejmuje decyzję o kontynuacji.
      Dopuszczalne postacie zwracają \texttt{void}, \texttt{bool}
      albo inny \texttt{std::coroutine\_handle}.
\item[\texttt{await\_resume()}] Wyznacza wynik całego wyrażenia
      \texttt{co\_await}; może też zgłosić wyjątek.
\end{description}

Gdy \texttt{await\_suspend} zwraca \texttt{bool}, wartość
\texttt{false} oznacza kontynuację bez pozostawania w zawieszeniu,
a \texttt{true} --- pozostawienie korutyny w stanie zawieszonym.
Przy zwrocie \texttt{void} korutyna pozostaje zawieszona do późniejszego
wznowienia. Zwrócenie innego uchwytu powoduje \emph{symmetric transfer}:
sterowanie przechodzi do wskazanej korutyny.

Nie każdy typ bezpośrednio dostarcza trzy wymienione metody: język
obsługuje też przeciążenie \texttt{operator co\_await} oraz
\texttt{await\_transform} zdefiniowane w \texttt{promise\_type}.
Sposób uzyskania awaitera zależy od tych reguł.

\subsection{Czysto lokalne zawieszenie bez nowego wątku}

Poniższy typ liczy wznowienia. Wykorzystamy go z prostą korutyną,
której właściciel ręcznie wywołuje \texttt{resume()}.

\begin{lstlisting}[caption={co\_await z własnym awaiterem: żadnego nowego wątku (C++20); kompletny program.}]
#include <cassert>
#include <coroutine>
#include <exception>
#include <utility>

struct Stepper {
    struct promise_type {
        Stepper get_return_object() {
            return Stepper{handle::from_promise(*this)};
        }
        std::suspend_always initial_suspend() noexcept { return {}; }
        std::suspend_always final_suspend() noexcept { return {}; }
        void return_void() noexcept {}
        void unhandled_exception() { std::terminate(); }
    };

    using handle = std::coroutine_handle<promise_type>;
    handle h;
    explicit Stepper(handle x) : h(x) {}
    Stepper(const Stepper&) = delete;
    Stepper(Stepper&& other) noexcept
        : h(std::exchange(other.h, {})) {}
    ~Stepper() { if (h) h.destroy(); }
    bool resume() {
        if (!h || h.done()) return false;
        h.resume();
        return !h.done();
    }
};

struct Pause {
    int& count;
    bool await_ready() const noexcept { return false; }
    void await_suspend(std::coroutine_handle<>) const noexcept {
        ++count; // Zawiesz, az wlasciciel wywola resume().
    }
    void await_resume() const noexcept {}
};

Stepper example(int& count) {
    ++count;                 // wykona sie przy pierwszym resume
    co_await Pause{count};  // pierwsze zawieszenie
    ++count;                 // wykona sie przy drugim resume
}

int main() {
    int count = 0;
    auto step = example(count);
    assert(count == 0);
    assert(step.resume() && count == 2);
    assert(!step.resume() && count == 3);
}
\end{lstlisting}

Kod niczego nie oblicza równolegle. Wszystkie instrukcje wykonał
wątek \texttt{main}. Korutyna jedynie oddała sterowanie, zachowując
miejsce wznowienia.

\section{Własne \texttt{Task<T>}: leniwe obliczenie a asynchroniczność}
\label{sec:cpp15-task}

Zaprojektujemy najmniejszy rozsądny typ \texttt{Task<T>} zwracający
wynik i obsługujący wyjątek. Celowo nie implementujemy jeszcze
harmonogramu: pozwala to zrozumieć, że \emph{sam typ korutyny nie jest
schedulerem}. \texttt{get()} w tym przykładzie działa wyłącznie
wtedy, gdy ciało może zakończyć się po pojedynczym wznowieniu.

\begin{lstlisting}[caption={Lazy Task<T> bez schedulera (C++20); kompletny program.}]
#include <cassert>
#include <coroutine>
#include <exception>
#include <optional>
#include <stdexcept>
#include <utility>

template<class T>
class Task {
public:
    struct promise_type {
        std::optional<T> result;
        std::exception_ptr error;
        Task get_return_object() {
            return Task{Handle::from_promise(*this)};
        }
        std::suspend_always initial_suspend() noexcept { return {}; }
        std::suspend_always final_suspend() noexcept { return {}; }
        void return_value(T x) { result.emplace(std::move(x)); }
        void unhandled_exception() noexcept {
            error = std::current_exception();
        }
    };
    using Handle = std::coroutine_handle<promise_type>;

    explicit Task(Handle h) noexcept : h_(h) {}
    Task(const Task&) = delete;
    Task& operator=(const Task&) = delete;
    Task(Task&& x) noexcept
        : h_(std::exchange(x.h_, {})),
          taken_(std::exchange(x.taken_, false)) {}
    ~Task() { if (h_) h_.destroy(); }

    T get() {
        if (!h_) throw std::logic_error("Empty task");
        if (taken_) throw std::logic_error("Result already consumed");
        if (!h_.done()) h_.resume();
        if (!h_.done())
            throw std::logic_error("Task needs an external scheduler");
        taken_ = true;
        if (h_.promise().error)
            std::rethrow_exception(h_.promise().error);
        return std::move(*h_.promise().result);
    }

private:
    Handle h_{};
    bool taken_ = false;
};

Task<int> calculate(int x) {
    co_return x * x;
}

int main() {
    auto task = calculate(7); // lazy, obliczenie jeszcze nie ruszylo
    assert(task.get() == 49);
}
\end{lstlisting}

Ten \texttt{Task<T>} jest pedagogiczny: nie ma przeciążenia
\texttt{operator co\_await}, nie wspiera wielokrotnego pobierania
wyniku (drugie \texttt{get()} zgłasza błąd) i nie potrafi samodzielnie
czekać na zdarzenia zewnętrzne.
W szczególności \texttt{get()} nie może po prostu wielokrotnie
wywoływać \texttt{resume()} na rzeczywistym zadaniu I/O: wznowienie
przed ukończeniem operacji mogłoby naruszyć protokół awaitera.
Produkcyjne \texttt{Task<T>} wymaga m.in. protokołu kontynuacji,
jasnej semantyki własności i poprawnej synchronizacji z executorem.

\section{Korutyna a \texttt{std::async}: fundamentalne porównanie}
\label{sec:cpp15-async-comparison}

\texttt{std::async} (C++11, nagłówek \texttt{<future>}) jest gotową
funkcją uruchamiającą zadanie i zwracającą \texttt{std::future<T>}.
Nie wymaga od użytkownika implementowania \texttt{promise\_type}.
Przy polityce \texttt{std::launch::async} funkcja pracuje w odrębnym
wątku; polityka \texttt{std::launch::deferred} odracza jej uruchomienie
do pierwszego nieczasowego oczekiwania na wynik.

Z kolei korutyna C++20 definiuje sposób zapisywania \emph{kontynuacji}.
Nawet jeśli wynik ma charakter asynchroniczny, \emph{faktyczne
uruchamianie pracy} zależy od awaitera, event loop, schedulera lub
innego mechanizmu zewnętrznego. Nie istnieje standardowa ogólna
reguła ,,jedna korutyna = jeden wątek''.


\paragraph{Czy można napisać \texttt{co\_await future}?}
Nie wprost: w C++20 standardowy \texttt{std::future<T>} nie udostępnia
protokołu awaitera wymaganego przez \texttt{co\_await}.
\texttt{future.get()} jest operacją potencjalnie blokującą
bieżący wątek. Istnieją biblioteki dostarczające awaitery dla
przyszłych wyników, ale sam adapter nie zamienia blokującej
operacji w nieblokującą. Naiwne uruchamianie osobnego wątku
tylko po to, by wykonał \texttt{future.wait()} i wznowił korutynę,
może zniweczyć zalety oszczędnego używania wątków.
Dla rzeczywiście nieblokującego oczekiwania na I/O potrzeba
integracji z właściwym mechanizmem asynchronicznym.

\begin{center}
\small
\begin{tabularx}{\textwidth}{@{}lXX@{}}
\toprule
\textbf{Cecha} & \textbf{\texttt{std::async}} & \textbf{Korutyny C++20} \\
\midrule
Co opisuje? & Zadanie ze stanem wyniku & Wykonanie z punktami zawieszenia \\
Tworzy nowy wątek? & \texttt{launch::async}: osobny wątek & Samo \texttt{co\_await}: nie \\
Kiedy zaczyna pracę? & Zależnie od polityki & Zależnie od \texttt{initial\_suspend} i schedulera \\
Odbieranie wyniku & \texttt{future::get()} (może blokować) & Zależy od typu \texttt{Task} i awaitera \\
Wyjątki & Przechowywane w stanie \texttt{future} & Obsługa przez \texttt{unhandled\_exception} i typ wyniku \\
Anulowanie & Brak automatycznego anulowania przez \texttt{future} & Możliwe kooperacyjnie; wymaga protokołu \\
Pamięć/zasoby & Stan asynchroniczny, przy \texttt{async} również wątek & Ramka korutyny, ewentualnie executor \\
Gotowe narzędzie w C++20? & Tak & Składnia tak; ogólny \texttt{Task<T>} nie \\
\bottomrule
\end{tabularx}
\end{center}

\subsection{Ten sam rachunek zapisany na dwa sposoby}

Porównanie dwóch krótkich funkcji pokazuje różnicę semantyczną.

\begin{lstlisting}[caption={std::async uruchamia zadanie w odrębnym wątku (C++11); kompletny program.}]
#include <cassert>
#include <future>

int square(int x) { return x * x; }

int main() {
    auto f = std::async(std::launch::async, square, 7);
    assert(f.get() == 49);
}
\end{lstlisting}

\begin{lstlisting}[caption={Korutyna bez awaitera zewnętrznego nie tworzy wątku; fragment.}]
// Przyjmujemy implementacje Task<T> z poprzedniej sekcji.
Task<int> square_coroutine(int x) {
    co_return x * x;
}

// Task<int> t = square_coroutine(7);
// int result = t.get(); // 49, w watku wywolujacym get()
\end{lstlisting}

W obu przypadkach można otrzymać tę samą wartość, ale nie jest to
to samo zachowanie wykonawcze. \texttt{std::async} z \texttt{launch::async}
zarządza wykonaniem w innym wątku, a pokazany \texttt{Task}
jedynie odkłada wywołanie do momentu ręcznego wznowienia.

\subsection{Pułapka destruktora \texttt{future}}

Ważna różnica dotyczy czasu życia. Gdy ostatni uchwyt do stanu
asynchronicznego utworzonego przez \texttt{std::async} z polityką
\texttt{launch::async} jest niszczony i praca wciąż trwa, destruktor
czeka na zakończenie powiązanego wątku.
Pierwsze z poniższych wywołań zakończy się, zanim rozpocznie się drugie:

\begin{lstlisting}[caption={Zignorowany future wymusza oczekiwanie; fragment.}]
// Funkcje slow_a i slow_b wykonuja dlugie obliczenia.
// std::async(std::launch::async, slow_a);
// std::async(std::launch::async, slow_b);
\end{lstlisting}

Jeżeli chcemy umożliwić nakładanie się zadań, trzeba przechować ich
uchwyty przed oczekiwaniem:
\begin{lstlisting}[caption={Dwa przyszłe wyniki powinny żyć jednocześnie; fragment.}]
// auto fa = std::async(std::launch::async, slow_a);
// auto fb = std::async(std::launch::async, slow_b);
// auto a = fa.get();
// auto b = fb.get();
\end{lstlisting}

Analogicznie w systemie korutyn należy znać politykę niszczenia
\texttt{Task}: usunięcie właściciela zawieszonej korutyny bez
koordynacji z executorem może zostawić wiszący uchwyt. Różne
biblioteki oferują odmienne kontrakty anulowania i oczekiwania.

\subsection{Asynchroniczne I/O a zadania CPU-bound}

Przy \emph{operacji I/O} większość czasu może upływać na oczekiwaniu
na dysk, sieć lub timer. Duża liczba korutyn może opisywać wiele
rozpoczętych operacji, podczas gdy mała pula wątków obsługuje
zdarzenia gotowości i wznowienia. Korzyść nie bierze się z tego,
że \texttt{co\_await} ,,jest szybsze'', lecz z możliwości
nieblokowania wątków podczas oczekiwania.

Przy \emph{operacjach CPU-bound} (obliczenia numeryczne, kompresja,
przetwarzanie obrazów) samo zawieszenie korutyny nie zwiększa
liczby rdzeni wykonujących kod. Tu potrzebne jest rzeczywiste
rozłożenie pracy na wiele wątków (np. \texttt{std::async}, pula
wątków lub OpenMP). Korutyna może być wygodnym \emph{interfejsem}
do takiego executora, ale nie zastępuje go.

\section{Rzeczywista asynchroniczność: korutyna i własny executor}
\label{sec:cpp15-executor}

Żeby zestawienie z \texttt{std::async} nie pozostało czysto
teoretyczne, utworzymy prosty executor z jednym wątkiem roboczym.
\texttt{co\_await executor.schedule()} odłoży kontynuację do
kolejki. Wznowienie wykona \emph{wątek executora}, nie wywołujący.
Wartością zwracaną będzie nasza klasa \texttt{FutureTask<T>},
wykorzystująca wewnętrznie \texttt{std::promise<T>} i
\texttt{std::future<T>}.

W przeciwieństwie do poprzedniego przykładu ramka ma strategię
\emph{samozniszczenia} przez \texttt{final\_suspend = suspend\_never}.
Nie przechowujemy już zewnętrznego uchwytu właścicielskiego:
po zakończeniu korutyna sama domyka swój cykl życia, a klient posiada
wyłącznie \texttt{future}. W tej uproszczonej konstrukcji executor
musi żyć aż do zakończenia zadań, a pojedyncza korutyna jest
wznawiana w kolejce co najwyżej raz.

\begin{lstlisting}[caption={Korutyna przechodząca do innego wątku poprzez executor (C++20); kompletny program.}]
#include <cassert>
#include <condition_variable>
#include <coroutine>
#include <deque>
#include <exception>
#include <future>
#include <mutex>
#include <stdexcept>
#include <thread>
#include <utility>

class Executor {
    std::mutex m_;
    std::condition_variable cv_;
    std::deque<std::coroutine_handle<>> queue_;
    bool closed_ = false;
    std::thread worker_;

    void run() {
        for (;;) {
            std::coroutine_handle<> h;
            {
                std::unique_lock lock(m_);
                cv_.wait(lock, [&] { return closed_ || !queue_.empty(); });
                if (queue_.empty() && closed_) return;
                h = queue_.front();
                queue_.pop_front();
            }
            h.resume(); // H moze sie tu samozniszczyc.
            // Nie korzystamy juz z h po resume().
        }
    }

public:
    Executor() : worker_([this] { run(); }) {}
    Executor(const Executor&) = delete;
    Executor& operator=(const Executor&) = delete;

    ~Executor() {
        {
            std::lock_guard lock(m_);
            closed_ = true;
        }
        cv_.notify_one();
        worker_.join();
    }

    void post(std::coroutine_handle<> h) {
        {
            std::lock_guard lock(m_);
            if (closed_) throw std::runtime_error("Executor closed");
            queue_.push_back(h);
        }
        cv_.notify_one();
    }

    struct Schedule {
        Executor& ex;
        bool await_ready() const noexcept { return false; }
        void await_suspend(std::coroutine_handle<> h) const {
            ex.post(h); // Po przekazaniu h nie uzywamy juz awaitera.
        }
        void await_resume() const noexcept {}
    };

    Schedule schedule() noexcept { return Schedule{*this}; }
};

template<class T>
class FutureTask {
    std::future<T> future_;
public:
    explicit FutureTask(std::future<T> f) : future_(std::move(f)) {}
    FutureTask(FutureTask&&) = default;
    FutureTask(const FutureTask&) = delete;

    T get() { return future_.get(); }

    struct promise_type {
        std::promise<T> result;
        FutureTask get_return_object() {
            return FutureTask{result.get_future()};
        }
        std::suspend_never initial_suspend() noexcept { return {}; }
        std::suspend_never final_suspend() noexcept { return {}; }
        void return_value(T x) { result.set_value(std::move(x)); }
        void unhandled_exception() noexcept {
            try { result.set_exception(std::current_exception()); }
            catch (...) { std::terminate(); }
        }
    };
};

FutureTask<int> square_on(Executor& ex, int x) {
    co_await ex.schedule(); // Tu oddajemy sterowanie.
    co_return x * x;        // Wykona sie w watku roboczym.
}

int main() {
    Executor ex;
    auto a = square_on(ex, 6);
    auto b = square_on(ex, 7);
    assert(a.get() + b.get() == 85);
    // Przed zniszczeniem ex oba zadania sie zakonczyly.
}
\end{lstlisting}

\paragraph{Co to porównanie rzeczywiście dowodzi?}
Kod z \texttt{square\_on} jest asynchroniczny względem wywołującego,
ponieważ po \texttt{co\_await} kontynuację przejmuje inny wątek.
Nie oznacza to jeszcze \emph{równoległego wykonywania} obu obliczeń:
executor ma tylko jednego pracownika, więc zadania są przez niego
wykonywane kolejno. Do równoległego przetwarzania dwóch obliczeń
należałoby dodać drugiego pracownika albo wykorzystać istniejącą
pulę wątków.

Pokazany executor jest \emph{szkicem edukacyjnym}. Nie obsługuje
ogólnego anulowania ani ponownego planowania tej samej korutyny,
nie rozwiązuje wszystkich problemów czasu życia i nie stanowi
gotowej biblioteki. W szczególności klient nie ma uchwytu pozwalającego
anulować i bezpiecznie zniszczyć zawieszone zadanie. W systemie
produkcyjnym własność korutyny, zamknięcie executora i czas życia
powiązanych zasobów muszą stanowić jeden, spójny kontrakt.

\section{Obsługa wyjątków, anulowanie i czas życia danych}
\label{sec:cpp15-errors}

\subsection{Wyjątki a \texttt{unhandled\_exception}}

Nieobsłużony wyjątek z ciała korutyny nie może zostać pozostawiony
bez decyzji. Kompilator wywołuje \texttt{promise.unhandled\_exception()}.
W generatorze zapisaliśmy \texttt{std::exception\_ptr}; przy kolejnym
\texttt{next()} błąd zostanie zgłoszony w kodzie klienta. W naszym
\texttt{FutureTask<T>} wyjątek trafia do stanu \texttt{std::future}
i ujawnia się podczas \texttt{get()}. Jest to analogiczne do
propagacji błędów w \texttt{std::async}, ale sam mechanizm wymagał
jawnej implementacji w \texttt{promise\_type}.

Należy też uwzględnić wyjątki zgłaszane przez operacje awaitera
oraz błędy alokacji ramki. Awaria podczas tworzenia ramki może
wystąpić jeszcze przed wykonaniem ciała korutyny. Typ obietnicy
może opcjonalnie udostępniać \texttt{get\_return\_object\_on\_allocation\_failure()}
razem z odpowiednią, niezgłaszającą wyjątku alokacją, ale nie jest to
mechanizm automatycznie dostępny każdemu typowi korutyny.

\subsection{Referencje przechowywane w ramce}

Argumenty referencyjne nie wydłużają czasu życia obiektów.
Jeżeli korutyna przechowuje referencję do zmiennej lokalnej, która
zniknie podczas zawieszenia, jej późniejsze wznowienie może prowadzić
do użycia nieistniejącego już obiektu.

\begin{lstlisting}[caption={Niebezpieczny czas życia referencji w korutynie; fragment koncepcyjny.}]
// Zakladamy lazy Task<int> z poprzedniej sekcji.
Task<int> read_later(const int& ref) {
    co_return ref;
}

// NIEPOPRAWNE:
// auto task = read_later(123); // 123 znika na koncu tej instrukcji
// int value = task.get();     // odczyt przez wiszaca referencje
\end{lstlisting}

Z reguły bezpieczniej przechowywać wewnątrz ramki wartości
przeznaczone do późniejszego użycia (np. obiekt przekazany przez
wartość), a nie referencje do tymczasowych obiektów.
Trzeba też zwracać uwagę na czas życia domknięć lambda:
korutyna będąca lambdą i wykorzystująca przechwycenia może
po zawieszeniu odwoływać się do już zniszczonego obiektu domknięcia.
Sam fakt przeniesienia wynikowej korutyny nie musi temu zapobiegać.

\subsection{Anulowanie jest kooperacyjne}

\texttt{std::stop\_token} (C++20) umożliwia \emph{żądanie}, a nie
wymuszenie zakończenia zadania. Korutyna może sprawdzać token
przed rozpoczęciem kolejnych faz albo korzystać z awaitera,
który potrafi odpowiednio reagować na anulowanie. Nie wynika
z tego jednak, że każde \texttt{co\_await} automatycznie obsługuje
\texttt{stop\_token}.

\begin{lstlisting}[caption={Kooperacyjne anulowanie korutyny: sprawdzenie żądania w punktach kontrolnych; fragment.}]
// Zakladamy Task<int> z poprzedniej sekcji.
#include <stop_token>

Task<int> calculate(std::stop_token token) {
    if (token.stop_requested()) co_return -1;
    // Pierwszy krok obliczen...
    if (token.stop_requested()) co_return -1;
    // Kolejny krok obliczen...
    co_return 42;
}
\end{lstlisting}

W powyższym przykładzie token jest przekazywany przez wartość.
Sam \texttt{stop\_token} nie steruje czasem życia ramki i nie
wznawia korutyny, która czeka na zdarzenie. Do pełnego anulowania
operacji I/O potrzeba jeszcze protokołu z usługą, która wykonuje
rzeczywistą operację, oraz reguł określających \emph{kto} zniszczy
korutynę i \emph{kiedy} wolno to zrobić.

\section{Pułapki synchronizacji i planowania}

\subsection{Równoczesne wznowienie tej samej korutyny}

Nie można traktować \texttt{coroutine\_handle::resume()} jak
operacji atomowej nadającej się do równoczesnego wywoływania przez
wiele wątków. Trzeba zapewnić, że dana zawieszona korutyna ma
jednego logicznego właściciela uprawnienia do wznowienia.
Ponowne wznowienie korutyny już aktywnej lub zakończonej może
prowadzić do niezdefiniowanego zachowania.

\subsection{Utracone powiadomienie i wyścig publikacji uchwytu}

Awaiter może przekazać uchwyt do innego wątku wewnątrz
\texttt{await\_suspend}. Wówczas druga strona może wznowić korutynę
\emph{zanim} pierwotne wywołanie \texttt{await\_suspend} powróci.
Dlatego po opublikowaniu uchwytu kod awaitera nie powinien korzystać
z obiektów, których czas życia mógł się już zakończyć.

W bardziej złożonych awaiterach trzeba również uniknąć zgubienia
zdarzenia kończącego operację między sprawdzeniem gotowości a
zarejestrowaniem oczekiwania. Poprawne rozwiązania stosują np.
mutex, operacje atomowe albo protokół jednorazowego przejęcia
kontynuacji. Samo \texttt{await\_ready()} nie zastępuje synchronizacji.

\subsection{Brak własności \texttt{coroutine\_handle}}

Gdy planista ma w kolejce jedynie \texttt{coroutine\_handle<>},
nie posiada informacji, kto może zniszczyć ramkę. Usunięcie
obiektu \texttt{Task} podczas oczekiwania może prowadzić do
późniejszego wznowienia już zwolnionej pamięci.
Najbezpieczniejszy projekt definiuje własność od początku:
\emph{Task posiada ramkę}, \emph{executor posiada ramkę} albo
\emph{ramka posiada samą siebie}, z dokładnym protokołem przekazania
własności i zakończenia.

\section{Koszty, optymalizacje i pomiary}
\label{sec:cpp15-performance}

Wydajność należy rozpatrywać w całym systemie, a nie tylko przez
porównanie czasu jednego \texttt{co\_await} i jednego \texttt{std::async}.
\begin{itemize}
\item \textbf{Utworzenie wątku.} \texttt{std::async(launch::async, ...)}
      uruchamia pracę w odrębnym wątku; bardzo duża liczba krótkich
      zadań może generować istotny narzut zarządzania wątkami.
\item \textbf{Ramka.} Korutyna może wymagać alokacji pamięci,
      zachowywania stanu i wznawiania. Kompilator może czasem
      wyeliminować alokację ramki, ale nie jest to gwarancja.
\item \textbf{Executor.} Kolejka, mutexy, przełączanie między wątkami
      oraz współdzielenie pamięci również kosztują.
\item \textbf{Lokalność pamięci.} Duża liczba korutyn z rozproszonymi
      ramkami może prowadzić do dodatkowych nietrafień cache.
\item \textbf{Skalowalność.} Dla I/O ważniejsza może być liczba
      oczekujących operacji niż przepustowość pojedynczego zadania.
      Dla CPU-bound kluczowa jest liczba rdzeni, wielkość porcji pracy
      i ograniczenie dostępu do współdzielonych zasobów.
\end{itemize}

Rzetelny benchmark powinien mierzyć co najmniej: całkowitą
przepustowość, opóźnienia (również wartości odstające), liczbę wątków,
liczbę alokacji oraz maksymalne zużycie pamięci. Porównywanie
\texttt{std::async} z \emph{nieuruchamiającą nowych wątków} korutyną
wykonującą natychmiast prostą arytmetykę nie jest miarodajnym
testem wydajności asynchronicznej.

\subsection{Praktyczna decyzja: kiedy używać którego narzędzia?}

\begin{center}
\small
\begin{tabularx}{\textwidth}{@{}lX@{}}
\toprule
\textbf{Problem} & \textbf{Rozsądny pierwszy wybór} \\
\midrule
Jedno niezależne zadanie CPU z wynikiem & \texttt{std::async(std::launch::async, ...)} \\
Wiele krótkich zadań CPU & Pula wątków, executor lub OpenMP \\
Wiele połączeń I/O z długim oczekiwaniem & Korutyny + asynchroniczny runtime I/O \\
Leniwe generowanie sekwencji & Generator korutynowy \\
Wspólny licznik lub krótka sekcja krytyczna & \texttt{atomic} albo mutex, nie korutyna \\
Anulowalne zadanie działające długo & Przemyślany protokół \texttt{stop\_token} \\
Interfejs biblioteki o złożonych zależnościach & \texttt{Task<T>} nad zdefiniowanym executorem \\
\bottomrule
\end{tabularx}
\end{center}

\section{Związek z OpenMP, \texttt{thread} i \texttt{future}}

Rozszerzony rozdział o współbieżności omawiał \texttt{std::thread},
\texttt{std::jthread}, atomiki, blokady, \texttt{std::async} oraz OpenMP.
Korutyny nie unieważniają tych narzędzi; stanowią kolejną warstwę.
\begin{itemize}
\item \texttt{std::thread} opisuje konkretny wątek wykonania,
      a \texttt{std::jthread} dodatkowo pomaga zarządzać jego
      czasem życia i kooperacyjnym zatrzymywaniem.
\item \texttt{std::async} tworzy stan przyszłego wyniku i
      uruchamia wywołanie zgodnie z przyjętą polityką.
\item OpenMP rozdziela pętle i zadania między wątki obsługiwane
      przez środowisko wykonawcze OpenMP; szczególnie przydaje się
      w obliczeniach CPU-bound i kodzie numerycznym.
\item Korutyny pozwalają wygodnie opisać wieloetapową logikę
      oczekiwania i kontynuacji. Ich wydajność i stopień równoległości
      zależą od dostarczonego środowiska wykonawczego.
\end{itemize}

Możliwe jest użycie \emph{korutyn nad pulą wątków}: \texttt{co\_await}
przekazuje kontynuację do kolejki, a pracownik puli później ją wznawia.
Nie należy jednak od razu łączyć wszystkich mechanizmów.
Najpierw określamy, czy program czeka głównie na I/O, czy wykonuje
obliczenia, a następnie wybieramy model wykonania i własności zasobów.

\section{Najczęstsze błędy projektowe}

\begin{enumerate}
\item Zakładanie, że \texttt{co\_await} uruchamia nowy wątek.
\item Próba zwrócenia bezpośrednio \texttt{std::future<T>}
      z własnej korutyny bez odpowiedniego \texttt{promise\_type}.
\item Dwukrotne wywołanie \texttt{destroy()} dla tej samej ramki.
\item Wywołanie \texttt{resume()} po osiągnięciu
      \texttt{final\_suspend}.
\item Pozostawienie referencji lub przechwycenia lambdy do obiektu,
      który znika podczas zawieszenia.
\item Zniszczenie \texttt{Task} mimo tego, że executor posiada
      w kolejce uchwyt do jego ramki.
\item Mylenie automatycznej propagacji wyjątków w \texttt{std::async}
      z obowiązkami implementatora \texttt{promise\_type}.
\item Traktowanie \texttt{std::stop\_token} jako mechanizmu
      przymusowego przerwania dowolnego \texttt{co\_await}.
\item Równoczesne wznawianie tej samej korutyny z kilku wątków.
\item Uznanie, że zawieszenie i wznowienie zawsze kosztuje mniej
      niż zwykłe wywołanie funkcji.
\item Uruchamianie \texttt{std::async} dla tysięcy mikro-zadań
      bez rozważenia puli wątków.
\item Zwracanie z \texttt{await\_suspend} po opublikowaniu uchwytu
      i dalsze używanie potencjalnie zniszczonego awaitera.
\end{enumerate}

\section*{Podsumowanie rozdziału}
\addcontentsline{toc}{section}{Podsumowanie rozdziału}

Korutyny C++20 są mechanizmem \emph{zawieszenia i kontynuacji},
a nie samodzielnym środowiskiem asynchronicznym. Generator używa
\texttt{co\_yield}, aby leniwie udostępniać wartości. Typ zadania
używa \texttt{co\_return}, obiektu \texttt{promise\_type} i awaiterów,
aby opisać wykonanie z wynikiem i wyjątkami. \texttt{co\_await}
przekazuje kontrolę zgodnie z protokołem awaitera, a dopiero executor
lub inne środowisko może wznowić kod w osobnym wątku.

\texttt{std::async} jest prostym narzędziem do uruchomienia funkcji
i odebrania przyszłego wyniku. Korutyny są lepszym budulcem dla
złożonych przepływów wieloetapowych, szczególnie I/O, ale wymagają
większej odpowiedzialności za pamięć, synchronizację, anulowanie
i czas życia ramki. W obu przypadkach należy rozdzielać trzy pytania:
\emph{kto wykonuje kod, kto posiada stan i kto decyduje o zakończeniu}.

\section*{Literatura i dokumentacja}
\addcontentsline{toc}{section}{Literatura i dokumentacja}
\begin{itemize}
\item \url{https://en.cppreference.com/w/cpp/language/coroutines} --- model korutyn C++20.
\item \url{https://en.cppreference.com/w/cpp/coroutine/coroutine_handle} --- uchwyty i zarządzanie ramką.
\item \url{https://en.cppreference.com/w/cpp/coroutine/suspend_always} --- podstawowe awaitery.
\item \url{https://en.cppreference.com/w/cpp/thread/async} --- \texttt{std::async} i polityki uruchamiania.
\item \url{https://en.cppreference.com/w/cpp/thread/future} --- \texttt{std::future}.
\item \url{https://en.cppreference.com/w/cpp/coroutine/generator} --- \texttt{std::generator} (C++23).
\item Lewis Baker, \emph{C++ Coroutines: Understanding the Compiler Transformations}
      --- seria publikacji o modelu wykonania korutyn.
\item David Mazières, \emph{My Tutorial and Take on C++20 Coroutines}, 2021.
\end{itemize}

% Rozdzial 16. Fragment do skryptu Nowoczesny C++ (bez preambuly).
% C++20. Wymagane pakiety obecne w poprzednich rozdzialach:
% listings, amsmath, amssymb, booktabs, tabularx, array, hyperref, enumitem.
% Wymagane makra: \code, \term, \nota (jak w pliku glownym skryptu).
% Dolaczenie: \input{rozdzial_16_zaawansowana_wspolbieznosc_model_pamieci_cpp.tex}
% Kazdy listing opisany jako kompletny program jest samodzielny.

\chapter{Zaawansowana współbieżność, model pamięci i algorytmy bez blokad}
\label{chap:cpp16-advanced-concurrency}

\section{Cel rozdziału i powiązanie z podstawami}

W rozdziale~\ref{chap:concurrency} poznaliśmy wątki, mutexy, zmienne warunkowe, operacje
atomowe, \code{std::async}, pule wątków i~OpenMP. Sama znajomość tych
interfejsów nie wystarcza jednak do samodzielnego projektowania
algorytmów współbieżnych. Trzeba odpowiedzieć na trudniejsze pytania:
\begin{itemize}
\item Kiedy zapis wykonany przez jeden wątek \emph{musi} być widoczny dla drugiego?
\item Które wyniki są dozwolone przez model pamięci, nawet gdy wydają się
      sprzeczne z intuicją opartą na sekwencyjnym wykonywaniu?
\item Czym różni się atomowość pojedynczej operacji od niepodzielności
      całego algorytmu?
\item Jak zaprojektować \emph{postęp} algorytmu i~bezpiecznie zwalniać pamięć,
      gdy inne wątki mogą jeszcze posiadać wskaźniki do obiektów?
\item Kiedy rozwiązanie bez mutexa rzeczywiście poprawia wydajność?
\end{itemize}

Skupimy się na \emph{uzasadnianiu poprawności}, a~nie tylko na składni.
W~szczególności przedstawimy proste testy modelu pamięci (\emph{litmus tests}),
algorytm kolejki SPSC, publikację niezmiennych struktur danych, problem ABA,
techniki odzyskiwania pamięci oraz praktyczne zasady pomiaru i~diagnostyki.
Nie zakładamy, że każdy \code{std::atomic<T>} jest realizowany bez blokad.

\nota{Wszystkie gwarancje sformułowane w~rozdziale dotyczą abstrakcyjnego
modelu języka C++. Rzeczywisty procesor może wykonywać instrukcje
w~innej kolejności i~stosować bufory zapisów, ale program musi być
poprawny na gruncie standardu, a~nie tylko na jednym komputerze.}

\section{Formalny porządek zdarzeń w~wielowątkowym programie}
\label{sec:cpp16-formal}

\subsection{Zdarzenia i~relacje \emph{sequenced-before}, \emph{synchronizes-with}}

Rozpatrzmy pewne operacje na pamięci jako \emph{zdarzenia}. Wśród relacji
używanych do opisu ich uporządkowania najważniejsze są:
\begin{description}[style=nextline]
\item[\emph{Sequenced-before} (\(\mathrm{sb}\))]
  Porządek określony regułami języka pomiędzy wybranymi operacjami
  w~obrębie jednego wątku. Nie należy go mylić z~dowolną kolejnością
  pojawiania się instrukcji w~kodzie źródłowym: język zawiera także
  wyrażenia z~nieokreśloną kolejnością ewaluacji.
\item[\emph{Synchronizes-with} (\(\mathrm{sw}\))]
  Synchronizacja pomiędzy odpowiednimi operacjami różnych wątków:
  na przykład zapis \code{release} i~odczyt \code{acquire}, który czyta
  wartość z~odpowiedniego zapisu lub jego sekwencji \emph{release}.
  Podobną synchronizację zapewniają właściwe operacje na mutexach.
\item[\emph{Happens-before} (\(\mathrm{hb}\))]
  Relacja porządkująca zdarzenia, budowana w~typowych przykładach
  przez \(\mathrm{sb}\), \(\mathrm{sw}\) i~ich przechodnie złożenia.
  Jeśli \(A\;\mathrm{hb}\;B\), skutki \(A\) są uporządkowane przed \(B\)
  w~sensie formalnego modelu; nie jest to stwierdzenie o~czasie zegarowym.
\end{description}

Dla przejrzystości będziemy używać następującego schematu:
\[
\underbrace{A\;\xrightarrow{\mathrm{sb}}\;R}_{\text{producent}}
\quad\xrightarrow{\mathrm{sw}}\quad
\underbrace{Q\;\xrightarrow{\mathrm{sb}}\;B}_{\text{konsument}}
\quad\Longrightarrow\quad A\;\xrightarrow{\mathrm{hb}}\;B.
\]
Operacja \(R\) jest publikacją, \(Q\) odpowiadającym jej odbiorem;
\(A\) i~\(B\) mogą dotyczyć zwykłej, nieatomowej zmiennej.

\subsection{Atomowość nie zastępuje synchronizacji}

W~przeciwieństwie do zwykłego odczytu lub zapisu operacje na
\code{std::atomic<T>} nie powodują wyścigu danych pomiędzy sobą.
Nie oznacza to jednak, że każdy odczyt atomowy widzi ostatni zapis
według zegara ściennego. Taki globalny \emph{ostatni zapis} zwykle
w~ogóle nie istnieje jako zdefiniowane pojęcie modelu.

Z~kolei jeśli dwa wątki odczytują i~modyfikują to samo miejsce pamięci
nieatomowo, a~konfliktowe operacje nie są uporządkowane relacją
\emph{happens-before}, program ma \emph{data race} i~niezdefiniowane zachowanie.
Samo oznaczenie obiektu przez \code{volatile} nie rozwiązuje tego problemu.

\subsection{Porządek modyfikacji (\emph{modification order}) i~spójność}

Dla \emph{każdego pojedynczego} obiektu atomowego istnieje całkowity
porządek jego modyfikacji (\emph{modification order}). Nie wynika z~tego
jeden wspólny porządek modyfikacji \emph{wszystkich} atomików.
Ważne konsekwencje to między innymi:
\begin{itemize}
\item zapis do \code{x} uporządkowany przed innym zapisem do \code{x}
      musi być odpowiednio umieszczony przed nim w~porządku modyfikacji;
\item jeśli odczyt \code{x} jest uporządkowany przed kolejnym odczytem
      \code{x} w~tym samym wątku, drugi odczyt nie może cofnąć się
      do wcześniejszej modyfikacji w~porządku tego atomika;
\item ograniczenia spójności dla jednego atomika nie zapewniają
      automatycznie spójnego obrazu \emph{dwóch różnych} atomików.
\end{itemize}

\section{Testy modelu pamięci: wyniki, które zaskakują}
\label{sec:cpp16-litmus}

\subsection{\emph{Store buffering}: dwa odczyty zera}

Załóżmy, że \code{x} i~\code{y} początkowo są równe zero, a~wątki
wykonują następujące operacje:
\[
\begin{array}{c|c}
\text{Wątek 1} & \text{Wątek 2}\\ \hline
x\gets 1 & y\gets 1\\
r_1\gets y & r_2\gets x
\end{array}
\]

\begin{lstlisting}[caption={Store buffering: kompletna, wolna od wyścigu danych próba (C++20).}]
#include <atomic>
#include <iostream>
#include <thread>

int main() {
    std::atomic<int> x{0}, y{0};
    int r1 = -1, r2 = -1;

    std::jthread a([&] {
        x.store(1, std::memory_order_relaxed);
        r1 = y.load(std::memory_order_relaxed);
    });
    std::jthread b([&] {
        y.store(1, std::memory_order_relaxed);
        r2 = x.load(std::memory_order_relaxed);
    });
    a.join();
    b.join();
    std::cout << r1 << ' ' << r2 << '\n';
}
\end{lstlisting}

W~tym programie \(r_1=r_2=0\) jest \textbf{dozwolonym} wynikiem
modelu C++ dla \code{relaxed}; nie ma gwarancji, że zostanie
zaobserwowany w~praktycznym uruchomieniu. Jeśli \emph{wszystkie cztery}
operacje atomowe zamienimy na \code{memory_order_seq_cst}, wynik
\((0,0)\) stanie się niemożliwy: wspólny porządek operacji
sekwencyjnie spójnych musiałby przeczyć kolejności w~każdym z~wątków.

\nota{Wynik pojedynczego eksperymentu, nawet powtórzony milion razy,
nie dowodzi niedopuszczalności innego zachowania. Test tego rodzaju
to materiał do rozumowania nad modelem, a~nie pełny dowód poprawności.}

\subsection{\emph{Message passing}: publikacja zwykłych danych}

\begin{lstlisting}[caption={Publikacja danych przy użyciu release/acquire (C++11).}]
#include <atomic>
#include <cassert>
#include <thread>

int main() {
    int payload = 0;
    std::atomic<bool> published{false};

    std::thread producer([&] {
        payload = 42;
        published.store(true, std::memory_order_release);
    });
    std::thread consumer([&] {
        while (!published.load(std::memory_order_acquire)) {
            std::this_thread::yield();
        }
        assert(payload == 42);
    });

    producer.join();
    consumer.join();
}
\end{lstlisting}

Odczyt flagi, który widzi \code{true}, synchronizuje się z~zapisem
\code{release}. Zapis \code{payload = 42} następuje \emph{happens-before}
odczytu \code{payload} przez konsumenta. Flaga \emph{publikuje}
wcześniejszy stan; nie należy po publikacji modyfikować
\code{payload} bez dalszej synchronizacji.

\section{Porządki atomowe i~ich dokładne zastosowania}
\label{sec:cpp16-orders}

\begin{center}
\small
\begin{tabularx}{\textwidth}{@{}lX@{}}
\toprule
\textbf{Porządek} & \textbf{Typowe zastosowanie i~ograniczenie} \\
\midrule
\code{relaxed} & Liczniki, statystyki, pola kontrolne, gdy
sama atomowość wystarcza, a~relacja synchronizacji wynika skądinąd. \\
\code{release} & Zapis lub operacja RMW publikująca wcześniejsze efekty
wątku. \\
\code{acquire} & Odczyt lub operacja RMW odbierająca stan od publikującego wątku. \\
\code{acq_rel} & Operacja odczyt--modyfikacja--zapis wymagająca obu kierunków. \\
\code{seq_cst} & Operacje z~dodatkową gwarancją pojedynczego wspólnego
porządku dla operacji sekwencyjnie spójnych. \\
\bottomrule
\end{tabularx}
\end{center}

\code{memory_order_consume} ma długą i~skomplikowaną historię
standaryzacji. W~praktycznym kodzie C++20 należy preferować
\code{acquire} zamiast opierania przenośnej poprawności na
\code{consume}. Uwaga: porządek \code{release} jest właściwy dla
zapisu, \code{acquire} dla odczytu, a~\code{acq_rel} dla operacji
odczyt--modyfikacja--zapis; nie wszystkie porządki są poprawne dla
każdego przeciążenia funkcji atomowej.

\subsection{Operacje \emph{read--modify--write}}

Metody \code{fetch_add}, \code{fetch_sub}, \code{exchange} i~udany
\code{compare_exchange_*} wykonują operację RMW. Są one niepodzielne
w~odniesieniu do danego atomika, ale nie tworzą automatycznie
transakcji obejmującej inne pola.

\begin{lstlisting}[caption={Przydzielanie kolejnych numerów: wystarcza relaxed (C++11).}]
#include <atomic>
#include <cassert>
#include <thread>

int main() {
    std::atomic<int> next_id{0};
    int first = -1, second = -1;
    std::thread a([&] {
        first = next_id.fetch_add(1, std::memory_order_relaxed);
    });
    std::thread b([&] {
        second = next_id.fetch_add(1, std::memory_order_relaxed);
    });
    a.join();
    b.join();
    assert(first != second);
    assert(first + second == 1);
}
\end{lstlisting}

Uzyskujemy unikatowe wartości, bo RMW mają jednoznaczny porządek
modyfikacji \code{next_id}. Nie uzyskujemy przez to żadnej dodatkowej
synchronizacji niezwiązanych danych.

\subsection{\emph{Release sequence}: łańcuch operacji RMW}

W~C++20 zapis z~porządkiem \code{release} może rozpoczynać sekwencję
\emph{release} składającą się z~następujących po nim operacji RMW
na \emph{tym samym} atomiku, wykonywanych także przez inne wątki.
Odczyt \code{acquire} wartości pochodzącej z~tej sekwencji może
synchronizować się z~początkowym zapisem \code{release}.

\begin{lstlisting}[caption={Przekazanie publikacji przez relaxed RMW (C++20).}]
#include <atomic>
#include <cassert>
#include <thread>

int main() {
    int payload = 7;
    std::atomic<int> state{0};

    std::thread producer([&] {
        payload = 42;
        state.store(1, std::memory_order_release);
    });
    std::thread intermediary([&] {
        int expected = 0;
        while ((expected = state.load(std::memory_order_acquire))
               == 0) {
            std::this_thread::yield();
        }
        state.fetch_add(1, std::memory_order_relaxed);
    });
    std::thread consumer([&] {
        while (state.load(std::memory_order_acquire) != 2) {
            std::this_thread::yield();
        }
        assert(payload == 42);
    });

    producer.join();
    intermediary.join();
    consumer.join();
}
\end{lstlisting}

Wątek pośredni czeka, aż stan osiągnie 1, i~wykonuje RMW, który
zapisuje 2. Odczyt wartości 2 z~\code{acquire} synchronizuje się
z~pierwotnym \code{release} przez sekwencję \emph{release}.
Nie można w~ogólności zastąpić \code{fetch_add} zwykłym
\code{store(2, relaxed)} i~utrzymać tej samej argumentacji.

\section{CAS: porównaj i~zamień, pętle ponawiania i~porządek niepowodzenia}

\begin{center}\small
\begin{tabularx}{\textwidth}{@{}p{0.38\textwidth}X@{}}
\toprule
Operacja atomowa & Działanie \\
\midrule
\code{compare_exchange_strong} & Jednorazowo porównuje i próbuje podmienić wartość; nie zgłasza pozornego niepowodzenia. \\
\code{compare_exchange_weak} & Może zgłosić pozorne niepowodzenie; nadaje się do pętli ponawiania. \\
\code{expected} przy CAS & Przy niepowodzeniu otrzymuje rzeczywistą wartość atomika, którą trzeba uwzględnić przy kolejnej próbie. \\
\code{memory_order_acq_rel} & Przy udanej operacji odczytująco-zapisującej łączy własności acquire i release. \\
\code{memory_order_relaxed} & Zapewnia atomowość bez dodatkowej synchronizacji z innymi danymi. \\
\bottomrule
\end{tabularx}
\end{center}
\label{sec:cpp16-cas}

\code{compare_exchange_weak(expected, desired, success, failure)}
atomowo porównuje bieżącą wartość z~\code{expected}. Gdy porównanie
się powiedzie, zapisuje \code{desired}; w~przeciwnym razie aktualizuje
\code{expected} do wartości zastanej w~atomiku. Wersja \code{weak}
może ponadto nie powieść się pozornie (\emph{spurious failure}).
W~pętli jest zwykle dobrym wyborem.

\begin{lstlisting}[caption={Atomowe maksimum z~pętlą CAS (C++11).}]
#include <atomic>
#include <cassert>

void atomic_max(std::atomic<int>& target, int candidate) {
    int old = target.load(std::memory_order_relaxed);
    while (old < candidate &&
           !target.compare_exchange_weak(
               old, candidate,
               std::memory_order_relaxed,
               std::memory_order_relaxed)) {
        // Przy niepowodzeniu old jest aktualizowane.
    }
}

int main() {
    std::atomic<int> maximum{10};
    atomic_max(maximum, 8);
    atomic_max(maximum, 17);
    assert(maximum.load() == 17);
}
\end{lstlisting}

Parametr \code{failure} opisuje porządek operacji \emph{odczytu}
w~przypadku niepowodzenia. Nie może to być \code{release} ani
\code{acq_rel}; musi też spełniać ograniczenia względem porządku
sukcesu określone standardem. Typowy wzorzec publikacji nowej
wartości to sukces z~\code{acq_rel} i~niepowodzenie z~\code{acquire}
(albo \code{relaxed}, jeśli tak wynika z~dowodu).

\subsection{\code{weak} kontra \code{strong}}

\code{compare_exchange_strong} nie zgłasza pozornego niepowodzenia
przy poprawnie porównanych wartościach. Może być użyteczne dla
pojedynczej próby; \code{weak} bywa wygodniejsze w~pętlach CAS.
Nie ma uniwersalnej gwarancji, że jedna z~tych wersji zawsze działa
szybciej na danej platformie.

\section{Bariery pamięci i~\texttt{std::atomic\_thread\_fence}}
\label{sec:cpp16-fences}

Bariera \code{std::atomic_thread_fence(order)} jest operacją
porządkowania pamięci, która sama nie odczytuje ani nie zapisuje
konkretnego atomika. Synchronizacja z~innym wątkiem wciąż wymaga
odpowiednich operacji atomowych łączących obie strony. Samo
wstawienie \code{fence} obok zwykłej zmiennej nie usuwa wyścigu danych.

\begin{lstlisting}[caption={Publikacja ze zwolnieniową i~przejmującą barierą pamięci (C++11).}]
#include <atomic>
#include <cassert>
#include <thread>

int main() {
    int payload = 0;
    std::atomic<bool> flag{false};

    std::thread producer([&] {
        payload = 42;
        std::atomic_thread_fence(std::memory_order_release);
        flag.store(true, std::memory_order_relaxed);
    });
    std::thread consumer([&] {
        while (!flag.load(std::memory_order_relaxed)) {
            std::this_thread::yield();
        }
        std::atomic_thread_fence(std::memory_order_acquire);
        assert(payload == 42);
    });

    producer.join();
    consumer.join();
}
\end{lstlisting}

To przykład synchronizacji \emph{fence--fence}: odczyt atomika
widzi zapis następujący po barierze \code{release}, a~bariera
\code{acquire} jest wykonywana po tym odczycie. W~nowym kodzie
zwykle prostsze i~czytelniejsze jest zastosowanie pary operacji
\code{release}/\code{acquire} bez jawnych barier. Bariery nie są
substytutem przemyślanego protokołu komunikacji.

\section{Gwarancje postępu: \emph{blocking}, \emph{lock-free} i~\emph{wait-free}}
\label{sec:cpp16-progress}

Wielowątkowy algorytm można analizować pod względem \emph{bezpieczeństwa}
(\emph{safety}: nic niedozwolonego się nie dzieje) i~\emph{żywotności}
(\emph{liveness}: coś pożądanego w~końcu się wydarzy). Są to różne kwestie.

\begin{description}[style=nextline]
\item[\emph{Blocking}]
  Wątek może oczekiwać na odblokowanie mutexa przez konkretny inny wątek.
  Jeśli właściciel przestanie robić postępy, pozostali również mogą utknąć.
\item[\emph{Lock-free}]
  W~standardowym ujęciu postępu algorytmowego, mimo rywalizacji wątków,
  system jako całość może kontynuować wykonywanie operacji:
  niepowodzenia jednej próby CAS zwykle oznaczają, że inny uczestnik
  wykonał postęp. Nie gwarantuje to postępu \emph{każdego} wątku.
\item[\emph{Wait-free}]
  Każda operacja danego wątku ma ograniczoną liczbę własnych kroków
  niezależnie od działań innych wątków (przy przyjętym modelu).
  To silniejszy warunek niż \emph{lock-free}.
\item[\emph{Obstruction-free}]
  Operacja postępuje, jeśli wątek może działać dostatecznie długo
  bez konfliktów ze strony innych wątków. Jest to słabsza gwarancja.
\end{description}

\emph{Lock-free} nie znaczy ,,zawsze szybciej''. Nie oznacza również
braku aktywnego oczekiwania ani braku zagłodzenia konkretnego wątku.
Odróżniajmy też własność implementacji pojedynczego atomika
(\code{is_lock_free()}) od postępu kompletnego algorytmu.

\subsection{Prosta blokada wirująca (\emph{spinlock})}

\begin{lstlisting}[caption={Poprawna blokada wirująca z flagą atomową; ilustracja (C++20).}]
#include <atomic>
#include <cassert>
#include <mutex>
#include <thread>

class SpinLock {
    std::atomic_flag flag = ATOMIC_FLAG_INIT;
public:
    void lock() noexcept {
        while (flag.test_and_set(std::memory_order_acquire)) {
            // Zamiast stale rywalizowac o cache line, czekamy.
            flag.wait(true, std::memory_order_relaxed);
        }
    }
    void unlock() noexcept {
        flag.clear(std::memory_order_release);
        flag.notify_one();
    }
};

int main() {
    SpinLock lock;
    int counter = 0;
    auto work = [&] {
        for (int i = 0; i < 10000; ++i) {
            std::lock_guard<SpinLock> guard(lock);
            ++counter;
        }
    };
    std::jthread a(work), b(work);
    a.join();
    b.join();
    assert(counter == 20000);
}
\end{lstlisting}

\code{test_and_set} ma semantykę RMW; odblokowanie jest operacją
\code{release}. Metoda \code{wait(true)} może blokować, a~zatem
powyższy \code{SpinLock} \emph{nie} jest algorytmem \emph{lock-free}.
Ponadto nie gwarantuje sprawiedliwości. W~realnym oprogramowaniu
\code{std::mutex} zwykle pozostaje lepszym wyborem, zwłaszcza gdy
sekcja krytyczna może być długa lub obciążenie jest wysokie.

\section{Problem ABA: poprawne CAS nie wystarcza}
\label{sec:cpp16-aba}

Rozważmy wierzchołek stosu reprezentowany wskaźnikiem atomowym.
Wątek \(A\) czyta wskaźnik \(p\) i~zamierza wykonać
\code{compare_exchange} z~wartości \(p\) na \code{p->next}.
W~tym czasie wątek \(B\) może:
\begin{enumerate}
\item zdjąć element \(p\) ze stosu;
\item zdjąć jeszcze jeden element;
\item ponownie umieścić \(p\) na wierzchołku stosu.
\end{enumerate}

Wartość wskaźnika uległa zmianie \(A\to B\to A\) (stąd \emph{ABA}),
ale CAS wątku \(A\) może zobaczyć ponownie tę samą wartość \(p\).
W~międzyczasie struktura stosu mogła się całkowicie zmienić.

\begin{center}
\(\boxed{\text{head}=p}\;\longrightarrow\;
  \boxed{\text{head}=q}\;\longrightarrow\;
  \boxed{\text{head}=p}\)
\end{center}

Dwa problemy trzeba odróżnić:
\begin{enumerate}
\item \emph{ABA logiczne}: porównanie tej samej wartości nie wykrywa
      zmian stanu struktury.
\item \emph{Bezpieczeństwo pamięci}: wątek \(A\) może próbować odczytać
      \code{p->next}, gdy inny wątek zniszczył już obiekt \code{*p}.
      Taki odczyt nie staje się poprawny tylko dlatego, że \code{head}
      jest atomowy.
\end{enumerate}

\subsection{Wskaźnik z~licznikiem generacji}

Jedna z~technik polega na porównywaniu pary
\((\text{wskaźnik},\text{licznik wersji})\) zamiast samego wskaźnika.
Wtedy ponowne pojawienie się tego samego adresu nie oznacza równości
całego stanu. Licznik również ma ograniczony zakres i~może się
przepełnić; trzeba uzasadnić brak niebezpiecznego \emph{wrap-around}.

\nota{Nie zakładaj, że \code{std::atomic<SomeStruct>} będzie bez blokad:
implementacja może wykorzystać wewnętrzny mutex lub bibliotekę
atomowych operacji. Możliwość zapisania stanu złożonego w~\code{atomic}
nie rozwiązuje też samodzielnie problemu czasu życia obiektów.}

\subsection{Hazard pointers, epoch-based reclamation i~RCU}

\term{Hazard pointers} pozwalają wątkom publikować informację:
,,zamierzam korzystać z~tego obiektu''. Wątek usuwający obiekt
odkłada go na listę \emph{retired} i~odracza zniszczenie do chwili,
gdy żaden aktywny hazard pointer go nie wskazuje. Poprawna procedura
zwykle musi ponownie zweryfikować odczyt wskaźnika po publikacji
hazardu; samo opublikowanie dowolnego, potencjalnie już nieważnego
adresu nie wystarcza.

\term{Epoch-based reclamation} grupuje operacje w~epoki. Usunięte
obiekty mogą zostać odzyskane dopiero wtedy, gdy wszystkie aktywne
wątki zakończyły okresy korzystania z~danych obejmujące ich usunięcie.
Dłużej zatrzymany wątek może opóźniać zwolnienie pamięci.

\term{Read-copy-update (RCU)} korzysta z~odczytów starych wersji
danych i~publikacji nowych wersji; usuwanie poprzednich odbywa się
po odpowiednim okresie ochronnym (\emph{grace period}). Szczegóły
różnią się pomiędzy implementacjami.

\begin{center}
\small
\begin{tabularx}{\textwidth}{@{}lX@{}}
\toprule
\textbf{Metoda} & \textbf{Kiedy rozważyć} \\
\midrule
Hazard pointers & Wiele odczytów wskaźników, częste usuwanie
węzłów; możliwa kosztowna administracja hazardami. \\
Epoki & Duży ruch, kontrolowane sekcje dostępu; zatrzymany czytelnik
może opóźnić odzyskanie pamięci. \\
RCU & Wzorce ,,wielu czytelników, niewielu pisarzy'', publikacja
niezmiennych wersji danych. \\
\code{shared_ptr} & Gdy łatwość dowodu poprawności jest ważniejsza
od kosztu liczenia referencji; nie obiecuje \emph{lock-free}. \\
\bottomrule
\end{tabularx}
\end{center}

W~C++20 standardowa biblioteka nie udostępnia jeszcze kompletnej,
przenośnej abstrakcji hazard pointerów. Bezpieczna implementacja
wymaga precyzyjnego protokołu pamięciowego, dlatego zamiast podawać
uproszczony, ale błędny stos ze zwalnianiem węzłów, przedstawimy
algorytm bez dynamicznego odzyskiwania pamięci: kolejkę SPSC.

\section{Kolejka SPSC: praktyczny algorytm bez mutexa}
\label{sec:cpp16-spsc}

\emph{Single Producer, Single Consumer} (SPSC) oznacza dokładnie
jednego producenta i~jednego konsumenta. Bufor cykliczny może wtedy
używać dwóch liczników:
\begin{itemize}
\item \code{head} --- numer kolejnego miejsca zapisu, modyfikowany
      wyłącznie przez producenta;
\item \code{tail} --- numer kolejnego miejsca odczytu, modyfikowany
      wyłącznie przez konsumenta.
\end{itemize}

Producent zapisuje komórkę, a~następnie publikuje nowy \code{head}
operacją \code{release}. Konsument odczytuje \code{head} przez
\code{acquire}, a~więc może bezpiecznie zobaczyć zawartość wcześniej
zapisanej komórki. Analogicznie konsument zwalnia komórkę, publikując
nowy \code{tail}, a~producent odbiera tę informację operacją
\code{acquire}. Warunek pełności bufora o~pojemności \(C\) to
\(\mathrm{head}-\mathrm{tail}=C\), a~pustości:
\(\mathrm{head}=\mathrm{tail}\).

\begin{lstlisting}[caption={Kompletna kolejka SPSC dla int, ze statyczną pamięcią (C++20).}]
#include <array>
#include <atomic>
#include <cassert>
#include <cstddef>
#include <thread>

// Tylko jeden producent i dokladnie jeden konsument.
template<std::size_t Capacity>
class SpscQueue {
    static_assert(Capacity > 0);
    static_assert((Capacity & (Capacity - 1)) == 0);
    std::array<int, Capacity> buffer_{};
    std::atomic<std::size_t> head_{0};
    std::atomic<std::size_t> tail_{0};
public:
    bool try_push(int value) {
        const auto head = head_.load(std::memory_order_relaxed);
        const auto tail = tail_.load(std::memory_order_acquire);
        if (head - tail == Capacity) return false;
        buffer_[head & (Capacity - 1)] = value;
        head_.store(head + 1, std::memory_order_release);
        return true;
    }

    bool try_pop(int& result) {
        const auto tail = tail_.load(std::memory_order_relaxed);
        const auto head = head_.load(std::memory_order_acquire);
        if (tail == head) return false;
        result = buffer_[tail & (Capacity - 1)];
        tail_.store(tail + 1, std::memory_order_release);
        return true;
    }
};

int main() {
    constexpr int Count = 100000;
    SpscQueue<1024> queue;
    long long result = 0;

    std::jthread producer([&] {
        for (int x = 1; x <= Count; ++x) {
            while (!queue.try_push(x)) std::this_thread::yield();
        }
    });
    std::jthread consumer([&] {
        int item = 0;
        for (int n = 0; n < Count; ++n) {
            while (!queue.try_pop(item)) std::this_thread::yield();
            result += item;
        }
    });
    producer.join();
    consumer.join();
    assert(result == 1LL * Count * (Count + 1) / 2);
}
\end{lstlisting}

\subsection{Dlaczego kod jest poprawny?}

Dla danych w~\code{buffer_} nie stosujemy operacji atomowych.
Nie ma jednak konfliktowego dostępu bez synchronizacji, ponieważ:
\begin{enumerate}
\item przed odczytem komórki konsument otrzymuje informację o~zapisie
      przez \code{head_.load(acquire)};
\item przed ponownym nadpisaniem tej komórki producent otrzymuje
      informację o~jej zwolnieniu przez \code{tail_.load(acquire)};
\item każdy z~liczników ma jednego piszącego, więc nie potrzeba CAS
      do aktualizowania \code{head_} ani \code{tail_}.
\end{enumerate}

Przykład zakłada, że operacji \code{try_push} nie wywołuje kilku
producentów, a~\code{try_pop} kilku konsumentów, nawet jeśli
same metody są dostępne publicznie. Złamane założenie SPSC
powoduje utratę poprawności. Ograniczenia bufora, przepełnienie
liczników i~kwestie typów niebanalnych wymagają osobnej analizy.
Użyliśmy potęgi dwójki, żeby indeksowanie przez maskę pozostało
poprawne także przy modularnym zawijaniu licznika bez znaku.

\nota{Ten przykład jest edukacyjną, \emph{nieblokującą} konstrukcją
w~sensie, że nie przejmuje mutexa, ale aktywne pętle producenta
i~konsumenta mogą czekać dowolnie długo. Metody \code{try_push}
i~\code{try_pop} są krótkie i~nie wykonują pętli CAS. Bez uzasadnienia
własności użytych atomików na danej platformie nie deklarujemy
silniejszej gwarancji wydajnościowej.}

\section{Atomowa publikacja niezmiennych wersji danych}
\label{sec:cpp16-snapshot}

Wiele rzeczywistych aplikacji (konfiguracja serwera, słowniki,
parametry modeli) posiada licznych czytelników i~rzadkich pisarzy.
Dobrą architekturą bywa utworzenie nowej, niezmiennej wersji
struktury i~jednorazowe opublikowanie wskaźnika do niej.
W~C++20 można do tego wykorzystać specjalizację
\code{std::atomic<std::shared_ptr<T>>}.

\begin{lstlisting}[caption={Bezpieczny snapshot konfiguracji (C++20).}]
#include <atomic>
#include <cassert>
#include <memory>
#include <thread>

struct Config {
    int threshold;
    double scale;
};

int main() {
    using Ptr = std::shared_ptr<const Config>;
    std::atomic<Ptr> current{
        std::make_shared<const Config>(Config{10, 1.0})
    };

    std::jthread writer([&] {
        Ptr next = std::make_shared<const Config>(Config{20, 2.0});
        current.store(std::move(next), std::memory_order_release);
    });
    std::jthread reader([&] {
        Ptr snap = current.load(std::memory_order_acquire);
        assert((snap->threshold == 10 && snap->scale == 1.0) ||
               (snap->threshold == 20 && snap->scale == 2.0));
    });
    writer.join();
    reader.join();
}
\end{lstlisting}

Czytelnik otrzymuje \emph{spójny} snapshot obu pól, a~stara wersja
jest automatycznie utrzymywana przy życiu, dopóki istnieją do niej
uchwyty \code{shared_ptr}. Nie oznacza to, że przechowywanie
\code{shared_ptr} w~atomiku jest \emph{lock-free}. Referencyjne
liczenie i~alokacje bywają istotnym kosztem przy bardzo częstych
odczytach. Dla obciążenia z~dużą przewagą czytelników należy
porównać to podejście z~\code{shared_mutex}, kopiowaniem małych
struktur i~technikami RCU.

\subsection{\code{std::atomic_ref}: synchronizacja istniejącego obiektu}

\code{std::atomic_ref<T>} (C++20) pozwala wykonywać operacje atomowe
na istniejącym, odpowiednio wyrównanym obiekcie typu \code{T}.
Przydatne bywa przy interoperacyjności ze strukturami danych,
które nie przechowują \code{std::atomic<T>} jako pola.

\begin{lstlisting}[caption={Operacje atomowe na istniejącej zmiennej (C++20).}]
#include <atomic>
#include <cassert>
#include <thread>

int main() {
    alignas(std::atomic_ref<int>::required_alignment) int counter = 0;
    {
        std::atomic_ref<int> shared_counter(counter);
        auto add = [&] {
            for (int i = 0; i < 10000; ++i)
                shared_counter.fetch_add(1, std::memory_order_relaxed);
        };
        std::jthread a(add), b(add);
        a.join();
        b.join();
    } // koniec zycia atomic_ref
    assert(counter == 20000);
}
\end{lstlisting}

Dopóki istnieje \code{atomic_ref} odnoszący się do obiektu, dostęp
do tego obiektu należy realizować przez instancje \code{atomic_ref}
--- nie tylko w~chwilach równoczesnego dostępu. Dlatego zwykłe
odczytanie \code{counter} wykonujemy \emph{po} zakończeniu czasu
życia referencji atomowej. Sam obiekt musi żyć dostatecznie długo,
a~jego wyrównanie musi spełniać wymagania \code{required_alignment}.

\section{Planowanie zadań: pule wątków, \emph{work stealing} i~\emph{backpressure}}
\label{sec:cpp16-scheduling}

Poza poprawną synchronizacją potrzebujemy rozsądnej strategii
przypisywania prac do wątków. Wątek systemowy ma koszt utworzenia,
przełączenia kontekstu i~utrzymania stosu. Z~tego powodu tysiące
krótkich zadań nie powinny automatycznie oznaczać tysiąca wątków.

\begin{description}[style=nextline]
\item[\emph{Thread pool}]
  Ustalona liczba wątków pobiera zadania z~kolejki. Ogranicza koszt
  częstego tworzenia wątków i~może ograniczać równoległe obciążenie.
\item[\emph{Work stealing}]
  Każdy wątek ma własną kolejkę, a~bezczynny pracownik próbuje
  przejąć część pracy od innych. Dobrze sprawdza się przy nierównych
  czasach trwania podzadań. Wymaga ostrożnego protokołu synchronizacji
  kolejek, szczególnie gdy korzystają z~nich różni producenci.
\item[\emph{Backpressure}]
  Mechanizm spowalniania producenta, gdy konsumenci nie nadążają.
  Bez niego kolejka bez ograniczenia pojemności może zużyć całą pamięć.
\item[\emph{Structured concurrency}]
  Sposób organizacji zadań, w~którym ich czas życia, anulowanie i~błędy
  pozostają związane z~zakresem nadrzędnej operacji. Sam \code{std::async}
  nie implementuje kompletnego frameworka takiego stylu.
\end{description}

\subsection{Pula wątków: bezpieczne zamknięcie i~propagacja wyjątków}

Następny przykład rozwija pulę z~rozdziału~\ref{chap:concurrency}. Używa mutexa
oraz zmiennej warunkowej. Łatwy dowód poprawności i~przewidywalne
zachowanie przy zamykaniu są często ważniejsze niż
niepotwierdzony \emph{lock-free}.

\begin{lstlisting}[caption={Prosta pula wątków odprowadzająca kolejkę przed zamknięciem (C++20).}]
#include <condition_variable>
#include <functional>
#include <future>
#include <iostream>
#include <memory>
#include <mutex>
#include <queue>
#include <stdexcept>
#include <thread>
#include <type_traits>
#include <utility>
#include <vector>

class ThreadPool {
    std::mutex mutex_;
    std::condition_variable cv_;
    std::queue<std::function<void()>> jobs_;
    bool closed_ = false;
    std::vector<std::jthread> workers_;

    void worker_loop() {
        for (;;) {
            std::function<void()> job;
            {
                std::unique_lock lock(mutex_);
                cv_.wait(lock, [&] {
                    return closed_ || !jobs_.empty();
                });
                if (closed_ && jobs_.empty()) return;
                job = std::move(jobs_.front());
                jobs_.pop();
            }
            job();  // packaged_task przechwyci wyjatek
        }
    }

public:
    explicit ThreadPool(unsigned count) {
        if (count == 0) throw std::invalid_argument("count=0");
        try {
            for (unsigned i = 0; i < count; ++i)
                workers_.emplace_back([this] { worker_loop(); });
        } catch (...) {
            close();  // obudz pracownikow; ich destruktory wykonaja join
            throw;
        }
    }
    ThreadPool(const ThreadPool&) = delete;
    ThreadPool& operator=(const ThreadPool&) = delete;
    ~ThreadPool() { close(); }

    void close() {
        {
            std::lock_guard lock(mutex_);
            closed_ = true;
        }
        cv_.notify_all();
    }

    template<class F>
    auto submit(F&& callable)
        -> std::future<std::invoke_result_t<std::decay_t<F>&>> {
        using R = std::invoke_result_t<std::decay_t<F>&>;
        auto task = std::make_shared<std::packaged_task<R()>>(
            std::forward<F>(callable)
        );
        auto result = task->get_future();
        {
            std::lock_guard lock(mutex_);
            if (closed_) throw std::runtime_error("pool closed");
            jobs_.emplace([task] { (*task)(); });
        }
        cv_.notify_one();
        return result;
    }
};

int main() {
    ThreadPool pool(2);
    auto a = pool.submit([] { return 40; });
    auto b = pool.submit([] { return 2; });
    std::cout << a.get() + b.get() << '\n'; // 42
    pool.close(); // od tej chwili bez nowych zadan
}
\end{lstlisting}

Ważne własności pokazanej implementacji:
\begin{itemize}
\item \code{close()} jest bezpieczne wobec równoległego \code{submit()}:
      oba działania są serializowane przez ten sam mutex;
\item po \code{close()} pracownicy \emph{kończą kolejkę}, a~dopiero
      następnie opuszczają pętlę; destruktory \code{jthread} dołączają
      pracowników, zanim zniszczone zostaną mutex i~kolejka;
\item \code{packaged_task} zapisuje wartość lub wyjątek w~stanie
      asynchronicznym, więc zwykłe wyjątki z~zadań nie wychodzą
      bezpośrednio przez \code{worker_loop};
\item ta minimalna pula nie ma ograniczenia długości kolejki,
      priorytetów, terminów ani anulowania uruchomionych prac.
\end{itemize}

\nota{Nie wolno niszczyć takiej puli z~zadania wykonywanego w~tej
samej puli: dołączenie własnego wątku prowadzi do błędu.
Podobnie zadanie czekające na inne zadanie z~tej samej puli może
zakleszczyć program, jeśli wszyscy pracownicy pozostają zajęci
oczekiwaniem. Produkcyjny executor wymaga rozwiązań dla takich sytuacji.}

\subsection{Kiedy potrzebny jest \emph{work stealing}?}

Załóżmy, że przetwarzamy drzewo lub dzielimy dużą macierz na bloki.
Część bloków kończy się po mikrosekundach, inne zajmują milisekundy.
Statyczny przydział zadań może zostawić jeden wątek przeciążony,
gdy inne nie mają już pracy. Przechwytywanie pracy pomiędzy lokalnymi
kolejkami poprawia \emph{równoważenie obciążenia}.

Typowy projekt stosuje:
\begin{enumerate}
\item właścicielski dostęp pracownika do jednego końca jego kolejki;
\item kradzież elementów przez inne wątki z~drugiego końca;
\item atomowe indeksy i~protokół obsługujący konflikt właściciela
      ze złodziejem o~ostatni element;
\item kontrolę czasu życia tablicy zadań oraz jej ewentualnej realokacji.
\end{enumerate}

Takie struktury są znacznie bardziej subtelne niż \code{std::deque}
z~mutexem. W~praktyce warto użyć sprawdzonego executora/biblioteki,
a~własną implementację traktować jako osobny, formalnie weryfikowany projekt.

\section{Anulowanie, sygnały stopu i~punkty oczekiwania}
\label{sec:cpp16-cancel}

Kooperacyjne zatrzymywanie w~C++20 jest oparte na
\code{std::stop_source}, \code{std::stop_token} i~\code{std::stop_callback}.
Żądanie zatrzymania nie zabija wątku; umożliwia mu dobrowolne
zakończenie po zweryfikowaniu stanu i~odpowiednim uporządkowaniu zasobów.

\begin{lstlisting}[caption={Wybudzanie i anulowanie czekania (C++20).}]
#include <condition_variable>
#include <iostream>
#include <mutex>
#include <thread>

int main() {
    std::mutex mutex;
    std::condition_variable_any cv;
    bool ready = false;

    std::jthread worker([&](std::stop_token stop) {
        std::unique_lock lock(mutex);
        bool completed = cv.wait(lock, stop, [&] { return ready; });
        std::cout << (completed ? "ready\n" : "cancelled\n");
    });

    worker.request_stop();
    worker.join();
}
\end{lstlisting}

W~tym przeciążeniu \code{wait} zwraca wartość predykatu po zakończeniu
czekania. Jeśli żądanie zatrzymania nastąpi przed spełnieniem
predykatu, przykład wypisze \code{cancelled}. Jeśli stan zmieniłby
się w~międzyczasie, wynik mógłby być inny. Zwykła
\code{std::condition_variable} nie posiada w~C++20 odpowiednika
tego przeciążenia z~\code{stop_token}.

\subsection{Anulowanie i~bezpieczeństwo wyjątkowe}

W~dobrym API warto określić, co dzieje się z~wynikiem po anulowaniu:
\begin{itemize}
\item czy zwracany jest częściowy wynik;
\item czy do \code{future} trafia własny wyjątek oznaczający anulowanie;
\item czy potrzebne jest odwrócenie skutków wcześniej wykonanych operacji;
\item jak zapewnić, że wszystkie zależne zadania zostaną zakończone
      przed zniszczeniem współdzielonych danych.
\end{itemize}

Żadne z~tych pytań nie ma jednej odpowiedzi dla wszystkich bibliotek.
Sama metoda \code{future.wait_for()} nie anuluje zadań, a~korutyna
z~rozdziału~\ref{chap:cpp15-coroutines} wymaga mechanizmu propagacji żądania anulowania
przez awaitery i~executor.

\section{Fałszywe współdzielenie (\emph{false sharing}) i~cache coherence}
\label{sec:cpp16-false-sharing}

Dwa niezależne atomiki mogą być zapisane blisko siebie w~pamięci,
tak że mieszczą się w~tej samej jednostce przenoszenia danych
pomiędzy pamięciami podręcznymi. Wtedy zapis przez różne rdzenie
może wymagać kosztownych transferów linii cache, choć z~punktu
widzenia C++ nie ma wyścigu danych.

\begin{lstlisting}[caption={Rozdzielenie intensywnie zapisywanych liczników (C++20; stała wyrównania od C++17).}]
#include <atomic>
#include <cstdint>
#include <iostream>
#include <new>
#include <thread>

struct alignas(std::hardware_destructive_interference_size) Counter {
    std::atomic<std::uint64_t> value{0};
};

int main() {
    Counter a, b;
    auto increment = [](Counter& c) {
        for (int i = 0; i < 100000; ++i)
            c.value.fetch_add(1, std::memory_order_relaxed);
    };
    std::jthread first([&] { increment(a); });
    std::jthread second([&] { increment(b); });
    first.join();
    second.join();
    std::cout << a.value << ' ' << b.value << '\n';
}
\end{lstlisting}

\code{std::hardware_destructive_interference_size} jest wskazówką
implementacji dotyczącą rozdzielania obiektów często modyfikowanych
przez różne wątki; w~praktyce jej obsługa zależy od kompilatora
i~biblioteki standardowej. \code{alignas} może zwiększać zużycie
pamięci, zatem nie należy bez pomiarów wyrównywać w~ten sposób
każdego pola. Rozdział~\ref{chap:cpp-memory-performance} omawia szerzej cache locality i~organizację
danych AoS/SoA.

\subsection{Co mierzyć przy ocenie skalowalności?}

Pojedynczy czas wykonania jest mało informacyjny. Przy porównywaniu
wersji jednowątkowej i~wielowątkowej warto mierzyć:
\begin{itemize}
\item przepustowość (operacje/sekundę), a~nie tylko opóźnienie
      pojedynczego zadania;
\item \emph{speedup} \(S(p)=T(1)/T(p)\), dla \(p\) wątków;
\item rozkład opóźnień, np. medianę i~wysokie percentyle;
\item liczbę alokacji, przełączeń kontekstu, braków w~cache i~konfliktów CAS;
\item oddzielnie okres rozgrzewki oraz właściwy okres pomiarowy.
\end{itemize}

Według prawa Amdahla, jeśli część \(s\) programu musi pozostać
sekwencyjna, to w~idealizacji
\[
S(p)\leq\frac{1}{s+\frac{1-s}{p}}.
\]
Wzór ten nie obejmuje kosztów synchronizacji, a~więc rzeczywiste
przyspieszenie może być jeszcze niższe. Jednocześnie program
z~dominującym kosztem dostępu do pamięci może nie skalować się
nawet przy bardzo małej części sekwencyjnej.

\section{Linearizowalność i~niezmienniki współdzielonych struktur}
\label{sec:cpp16-linearizability}

Dla współbieżnych kolekcji często potrzebujemy mocniejszej własności
niż brak wyścigów danych: \emph{linearizowalności}.
Operacja jest linearizowalna, jeśli można przypisać jej jeden
umowny moment pomiędzy wywołaniem i~zakończeniem, w~którym
wywołanie ma taki sam skutek jak operacja wykonana sekwencyjnie.

Dla prostej atomowej inkrementacji punktem linearyzacji jest
sam RMW. Dla kolejki algorytm jest bardziej subtelny:
\begin{itemize}
\item producent najpierw zapisuje komórkę, a~następnie publikuje
      nowy \code{head};
\item konsument najpierw odczytuje i~przetwarza komórkę, a~następnie
      publikuje nowy \code{tail};
\item sama komórka nie musi być atomowa, o~ile zachowana jest
      właściwa synchronizacja i~protokół własności miejsca w~buforze.
\end{itemize}

W~algorytmach złożonych, np. z~dwoma atomikami \code{size} i~\code{head},
każde z~pól może być wolne od wyścigu danych, ale ich jednoczesny
odczyt nie musi stanowić spójnego snapshotu. Osobny \emph{niezmiennik}
struktury wymaga własnego uzasadnienia.

\subsection{Dlaczego prosty \emph{seqlock} nie wystarczy dla zwykłych danych?}

Popularny pomysł polega na tym, że pisarz zwiększa licznik wersji
przed i~po aktualizacji obiektu, a~czytelnik ponawia odczyt, jeżeli
wersja zmieniła się podczas odczytu. Jeżeli pola obiektu są zwykłymi
zmiennymi C++, czytelnik nadal może czytać je jednocześnie z~zapisem,
co oznacza \emph{data race}. Ponowienie odczytu \emph{po fakcie}
nie naprawia niezdefiniowanego zachowania, które już nastąpiło.

Poprawny wariant dla przenośnego C++ musi zastosować odpowiednią
synchronizację lub atomowe pola danych, a~nie tylko atomowy licznik
wersji. Alternatywą jest pokazana wcześniej publikacja niezmiennych
snapshotów.

\section{Współbieżność a~parallel STL, Ranges i~OpenMP}
\label{sec:cpp16-interop}

Nie każda równoległa praca wymaga własnej kolejki czy atomowych
liczników. Gdy zadanie polega na niezależnym przetwarzaniu elementów,
lepszą abstrakcją może być algorytm równoległy lub dyrektywa OpenMP.

\subsection{Polityki wykonania standardowych algorytmów}

Nagłówek \code{<execution>} udostępnia od C++17 między innymi:
\begin{description}[style=nextline]
\item[\code{std::execution::seq}] Wykonanie sekwencyjne.
\item[\code{std::execution::par}] Implementacja może wykonywać pracę
      równolegle; nie daje gwarancji, że rzeczywiście użyje wielu wątków.
\item[\code{std::execution::par_unseq}] Zezwala dodatkowo na
      niesekwencyjne wykonywanie, w~tym wektoryzację; nakłada ważne
      ograniczenia na wykonywane funkcje i~synchronizację.
\end{description}

\begin{lstlisting}[caption={Równoległa transformacja bez współdzielonych zapisów (C++17).}]
#include <algorithm>
#include <execution>
#include <numeric>
#include <vector>

int main() {
    std::vector<double> a(10000), b(a.size());
    std::iota(a.begin(), a.end(), 0.0);
    std::transform(std::execution::par,
                   a.begin(), a.end(), b.begin(),
                   [](double x) { return x * x; });
}
\end{lstlisting}

Kod nie zakłada kolejności przetwarzania elementów. Wyjściowy zakres
\code{b} nie nakłada się z~\code{a}, a~każde wywołanie funkcji
zapisuje inny element. Rzeczywiste zrównoleglenie zależy od
implementacji biblioteki standardowej i~konfiguracji środowiska;
niektóre kompilacje wymagają dodatkowej biblioteki wykonawczej.

\subsection{C++ i~OpenMP: dwa poziomy abstrakcji}

OpenMP, omawiany szczegółowo w~rozdziale~\ref{chap:concurrency}, umożliwia między innymi
równoległe pętle, redukcje i~zależności pomiędzy zadaniami.
W~obliczeniach numerycznych typowe rozwiązanie wygląda tak:

\begin{lstlisting}[caption={Wektorowa redukcja z OpenMP; wymagane -fopenmp.}]
#include <iostream>
#include <vector>

int main() {
    const std::vector<double> x(100000, 1.0);
    double sum = 0.0;

    #pragma omp parallel for reduction(+:sum) schedule(static)
    for (int i = 0; i < static_cast<int>(x.size()); ++i) {
        sum += x[static_cast<std::size_t>(i)];
    }
    std::cout << sum << '\n';
}
\end{lstlisting}

Kodowanie redukcji przez \code{reduction} pozwala uniknąć współdzielenia
jednego licznika zapisywanego przez wszystkie wątki. Kolejność
sumowania liczb zmiennoprzecinkowych może być odmienna niż przy
obliczeniu sekwencyjnym, a~zatem wynik może różnić się przez
zaokrąglenia. Liczbę wątków można ustawić m.in. przez
\code{OMP_NUM_THREADS}; \code{OMP_SCHEDULE} wpływa na harmonogram
tylko przy \code{schedule(runtime)}. Powyższy przykład jawnie
wybiera \code{schedule(static)}.

\nota{Mieszanie pul wątków, OpenMP, równoległych algorytmów STL
oraz zadań asynchronicznych może prowadzić do \emph{oversubscription}:
liczba aktywnych wątków znacząco przewyższa dostępne zasoby CPU.
Przy nakładaniu kilku warstw równoległości trzeba świadomie
kontrolować budżet wątków i~mierzyć efekt końcowy.}

\section{Diagnostyka: kompilatory, sanitizatory, testy przeciążeniowe}
\label{sec:cpp16-testing}

\subsection{Kompilacja i~uruchamianie przykładów}

Dla przykładów wykorzystujących \code{std::jthread}, \code{atomic::wait}
i~\code{atomic<shared_ptr<T>>} wybieramy tryb C++20:

\begin{lstlisting}[language=bash,caption={Kompilacja wielowątkowego programu C++20.}]
g++ -std=c++20 -O2 -Wall -Wextra -Wpedantic -pthread file.cpp -o app
./app
\end{lstlisting}

Dla przykładów OpenMP w~GCC lub odpowiednio skonfigurowanym Clangu:
\begin{lstlisting}[language=bash,caption={Włączenie OpenMP.}]
g++ -std=c++20 -O2 -Wall -Wextra -fopenmp file.cpp -o app
./app
\end{lstlisting}

\subsection{ThreadSanitizer i~inne narzędzia}

\emph{ThreadSanitizer} (TSan) potrafi wykrywać wiele wyścigów danych
i~wybrane błędy synchronizacji w~uruchamianym programie:

\begin{lstlisting}[language=bash,caption={Budowanie i uruchamianie z ThreadSanitizer.}]
g++ -std=c++20 -O1 -g -pthread -fsanitize=thread file.cpp -o app_tsan
./app_tsan
\end{lstlisting}

Nie wszystkie środowiska poprawnie wspierają TSan; szczególnie
przy nietypowych bibliotekach wykonawczych, OpenMP i~własnych
mechanizmach synchronizacji należy weryfikować zakres instrumentacji.
AddressSanitizer i~UndefinedBehaviorSanitizer mogą dodatkowo wykrywać
część błędów czasu życia oraz naruszeń zasad języka, lecz nie są
zamiennikami TSan ani formalnego dowodu synchronizacji.

\subsection{Jak testować algorytm współbieżny?}

Dla poważniejszej struktury danych zaleca się następującą procedurę:
\begin{enumerate}
\item \textbf{Zdefiniuj kontrakt}: ilu producentów i~konsumentów,
      dozwolone operacje oraz zasady zamknięcia i~niszczenia struktury.
\item \textbf{Zidentyfikuj własność danych}: kto może zapisywać dane
      w~każdej chwili; kiedy przechodzi własność konkretnej komórki.
\item \textbf{Narysuj relacje}: wskaż atomik, który publikuje zapis
      i~odczyt, który tę publikację odbiera.
\item \textbf{Wskaż punkt linearyzacji}: jeśli API ma być linearizowalne,
      określ, kiedy operacja staje się widoczna logicznie.
\item \textbf{Przeanalizuj zakończenie}: co dzieje się przy anulowaniu,
      wyjątku, opóźnieniu konsumenta i~niszczeniu obiektu.
\item \textbf{Testuj wielokrotnie}: zmieniaj liczbę wątków, rozmiary
      danych i~częstotliwość przełączeń; uruchamiaj TSan tam,
      gdzie jest dostępny.
\item \textbf{Benchmarkuj osobno}: włącz optymalizacje i~mierz
      przy realistycznym obciążeniu, ale nie myl szybkiego testu
      z~dowodem poprawności.
\end{enumerate}

\section{Typowe błędy i~zalecenia projektowe}
\label{sec:cpp16-guidelines}

\begin{center}
\small
\begin{tabularx}{\textwidth}{@{}lX@{}}
\toprule
\textbf{Błąd} & \textbf{Dlaczego jest groźny / co wybrać} \\
\midrule
\code{volatile} jako synchronizacja & Nie eliminuje \emph{data race};
użyj mutexa lub atomików z~właściwym porządkiem. \\
\code{relaxed} do publikacji zwykłych danych & Nie wytwarza
relacji synchronizacji; zazwyczaj potrzebne \code{release/acquire}. \\
Bezrefleksyjne \code{seq_cst} & Zwykle poprawniejsze niż źle
dobrane słabsze porządki, ale nie zastąpi logicznego kontraktu obiektu. \\
Kasowanie węzła po udanym CAS & Inny wątek może nadal używać
wskaźnika; zaprojektuj odzyskiwanie pamięci. \\
Dwa atomiki jako ,,jedna transakcja'' & Brak automatycznego
atomowego snapshotu pary; użyj blokady lub niezmiennej wersji. \\
\emph{Spinlock} dla długiej pracy & Prowadzi do rywalizacji
o~CPU; zwykle użyj \code{std::mutex}. \\
Nieograniczona kolejka zadań & Może zużyć pamięć; dodaj limit
i~mechanizm \emph{backpressure}. \\
Nadmierna liczba wątków & Konkurencja o~rdzenie i~cache;
kontroluj równoległość i~profiluj. \\
Mieszanie OpenMP i~własnej puli bez planu & Może prowadzić
do \emph{oversubscription} i~spadku wydajności. \\
\bottomrule
\end{tabularx}
\end{center}

\subsection{Praktyczne reguły wyboru}

\begin{itemize}
\item \textbf{Mutex} wybieraj jako punkt wyjścia, jeśli współdzielisz
      kilka powiązanych pól i~potrzebujesz łatwo sprawdzalnych niezmienników.
\item \textbf{Atomik \code{relaxed}} jest właściwy dla prostego licznika
      niezależnego od publikacji innych danych.
\item \textbf{\code{release/acquire}} stosuj, gdy istnieje wyraźny
      punkt przekazania danych lub własności pomiędzy wątkami.
\item \textbf{Lock-free} rozważaj dopiero wtedy, gdy pomiary pokażą
      realny koszt blokad, a~projekt ma przejrzysty dowód poprawności
      oraz plan odzyskiwania pamięci.
\item \textbf{OpenMP / parallel STL} wybieraj do operacji na niezależnych
      elementach, redukcji i~równoległości danych.
\item \textbf{Thread pool} wybieraj do dużej liczby krótkich zadań,
      gdy potrzebujesz ograniczać liczbę wątków, priorytety lub
      przepustowość; osobno zdefiniuj zasady zamykania.
\end{itemize}

\section*{Podsumowanie rozdziału}
\addcontentsline{toc}{section}{Podsumowanie rozdziału}

Zaawansowana współbieżność nie polega na mechanicznym zastępowaniu
mutexów przez \code{std::atomic}. Najpierw trzeba określić
\emph{własność danych i~kontrakt algorytmu}, następnie uzasadnić
relacje \emph{happens-before} i~wyznaczyć punkty widoczności
poszczególnych operacji. Dopiero wtedy warto wybierać porządki
pamięciowe, typy kolejek, planistę i~sposób odzyskiwania pamięci.

W~szczególności brak wyścigu danych nie oznacza jeszcze poprawności
logicznej, a~nieblokujący dostęp do atomika nie gwarantuje
nieblokującego ani szybkiego zachowania całego programu. Bezpieczna
współbieżność to połączenie \emph{modelu formalnego}, \emph{kontroli
czasu życia} i~\emph{rzetelnego pomiaru}.

\section*{Literatura i dokumentacja}
\addcontentsline{toc}{section}{Literatura i dokumentacja}
\begin{itemize}
\item Anthony Williams, \emph{C++ Concurrency in Action}, wyd.~2,
      Manning, 2019.
\item ISO/IEC 14882:2020, rozdziały o~modelu pamięci, wątkach,
      atomikach i~synchronizacji.
\item \url{https://en.cppreference.com/w/cpp/atomic/memory_order}
      --- porządki pamięci i~relacje synchronizacji.
\item \url{https://en.cppreference.com/w/cpp/atomic/atomic}
      --- typy atomowe i~operacje RMW.
\item \url{https://en.cppreference.com/w/cpp/atomic/atomic_ref}
      --- atomowe odwołania do istniejących obiektów.
\item \url{https://en.cppreference.com/w/cpp/thread/jthread}
      --- \code{jthread}, czas życia i~kooperacyjne zatrzymywanie.
\item OpenMP Architecture Review Board,
      \emph{OpenMP Application Programming Interface};
      \url{https://www.openmp.org/specifications/}.
\end{itemize}

% Rozdzial 17 do skryptu "Nowoczesny C++" (C++20).
% Plik zawiera wylacznie tresc rozdzialu: bez \documentclass i \begin{document}.
% Wymaga pakietow: listings, amsmath, amssymb, booktabs, tabularx,
% array, enumitem, hyperref; makr \code, \term, \nota z dotychczasowego skryptu.
% Uzycie: \input{rozdzial_17_kompilacja_linkowanie_abi_moduly_cpp.tex}
% Komendy powloki sa przeznaczone glownie dla GCC/Clang na Linux/ELF;
% roznice Windows/MSVC i macOS zaznaczono w tekscie.

\chapter{Kompilacja, linkowanie, ABI i moduły C++20}
\label{chap:cpp17-build-abi-modules}

\section{Dlaczego trzeba rozumieć proces budowania?}

W~poprzednich rozdziałach projektowaliśmy klasy, szablony, mechanizmy
zarządzania pamięcią i~współbieżność. Dopóki cały program mieści się
w~jednym pliku, wystarcza zwykle wywołanie \code{g++ main.cpp -o app}.
W~realnym projekcie kod dzieli się jednak na dziesiątki lub tysiące
\emph{jednostek translacji}, biblioteki zewnętrzne i~moduły. To wtedy
pojawiają się komunikaty \code{undefined reference}, \code{multiple definition},
konflikty wersji bibliotek, długi czas budowania oraz błędy ujawniające
się dopiero po wymianie biblioteki współdzielonej.

Musimy rozdzielić trzy różne zagadnienia:
\begin{description}[style=nextline]
\item[Poprawność programu w~języku C++]
Czy deklaracje, definicje, typy i~wywołania spełniają reguły C++,
w~szczególności \emph{One Definition Rule} (ODR)?
\item[Model budowania i~linkowania]
Jakie pliki są kompilowane osobno, gdzie kompilator zapisuje symbole,
w~jakiej kolejności linker przeszukuje biblioteki oraz jak ładowane są
biblioteki dynamiczne?
\item[Zgodność binarna (ABI)]
Czy dwa osobno zbudowane komponenty zgadzają się co do reprezentacji
typów, sposobu wywoływania funkcji, nazw symboli i~pozostałych reguł
interfejsu binarnego?
\end{description}

\nota{Norma ISO C++ opisuje semantykę programu, ale nie narzuca jednego
formatu pliku obiektowego, konwencji \emph{name mangling}, organizacji
bibliotek \code{.so}/\code{.dll} ani poleceń linkera. Poniżej
odróżniamy reguły języka od typowych rozwiązań narzędziowych.}

\section{Od pliku źródłowego do programu}

\begin{center}\small
\begin{tabularx}{\textwidth}{@{}p{0.31\textwidth}X@{}}
\toprule
Etap & Za co odpowiada \\
\midrule
Preprocesor & Rozwija dyrektywy \code{#include}, makra i warunkową kompilację. \\
Kompilator & Sprawdza składnię i typy oraz wytwarza kod dla jednostki translacji. \\
Asembler & Zamienia kod asemblera na plik obiektowy z symbolami i relokacjami. \\
Linker & Łączy pliki obiektowe i biblioteki, rozwiązując odwołania do symboli. \\
Loader systemowy & Ładuje program i wymagane biblioteki współdzielone przed uruchomieniem. \\
\bottomrule
\end{tabularx}
\end{center}
\label{sec:cpp17-pipeline}

\subsection{Preprocesor, kompilator, asembler, linker i~loader}

W~uproszczeniu klasyczny projekt z~nagłówkami przechodzi przez etapy:
\[
\begin{aligned}
\text{źródła i nagłówki}&\ \longrightarrow\ \text{preprocessing}
\ \longrightarrow\ \text{kompilacja}\\
&\phantom{\ \longrightarrow\ }\longrightarrow\ \text{plik obiektowy}
\ \longrightarrow\ \text{linkowanie}
\ \longrightarrow\ \text{program}.
\end{aligned}
\]

Na niektórych platformach kompilator wytwarza tekst asemblerowy
(\code{.s}) jako etap pośredni; w~innych przejście do kodu maszynowego
jest zintegrowane. Linkowanie może odbywać się przy budowaniu programu,
a~dodatkowe rozwiązywanie symboli może następować w~czasie ładowania
lub później.

Prześledźmy etapy na samodzielnym programie.

\begin{lstlisting}[caption={Plik hello.cpp: prosty program do obserwowania kompilacji.}]
#include <iostream>

int twice(int x) { return 2 * x; }

int main() {
    std::cout << twice(21) << '\n';
}
\end{lstlisting}

\begin{lstlisting}[language=bash,caption={Wybrane etapy budowania przy uzyciu GCC/Clang.}]
g++ -std=c++20 -E hello.cpp -o hello.ii  # Preprocesor
# Produkt poprzedniego etapu moze byc duzy.
g++ -std=c++20 -S hello.cpp -o hello.s   # Asembler

g++ -std=c++20 -c hello.cpp -o hello.o   # Kod obiektowy
g++ hello.o -o hello                      # Linkowanie
./hello                                   # 42
\end{lstlisting}

Opcja \code{-c} nie tworzy wykonywalnego programu; produkuje plik
obiektowy, który może zawierać nierozwiązane odwołania do symboli
pochodzących z~innych jednostek translacji. Polecenie \code{-S}
kończy budowanie na asemblerze, a~\code{-E} na wstępnym przetwarzaniu.
Warto odróżnić program sterujący kompilacją (np.~\code{g++}) od samego
frontendu C++, asemblera i~linkera.

\subsection{Pliki obiektowe, symbole i~relokacje}

Plik obiektowy zawiera zwykle kod maszynowy, dane, tablice symboli
oraz \emph{relokacje}, czyli informacje, które adresy muszą zostać
skorygowane po ustaleniu ostatecznego układu programu. W~środowisku
Linux powszechny jest format ELF; Windows używa PE/COFF, a~macOS
Mach-O. Szczegóły tych formatów nie są częścią języka C++.

\begin{lstlisting}[language=bash,caption={Ogladanie symboli i relokacji (Linux/ELF).}]
nm -C hello.o                 # Symbole; -C demangluje nazwy C++
readelf -Ws hello.o          # Tablica symboli ELF
readelf -r hello.o           # Relokacje
objdump -drC hello.o         # Asembler i relokacje
\end{lstlisting}

Przykładowo litera \code{T} w~typowym wyjściu \code{nm} informuje
o~symbolu definiowanym w~sekcji kodu, a~\code{U} o~symbolu
niezdefiniowanym w~danym obiekcie, który trzeba rozwiązać. Nie
należy traktować oznaczeń \code{nm} jako przenośnej własności C++.

\subsection{Jednostka translacji i~kompilacja oddzielna}

W~tradycyjnym modelu \emph{jednostka translacji} (\emph{translation unit},
TU) powstaje z~pliku źródłowego po uwzględnieniu dyrektyw preprocesora,
w~szczególności zawartości dołączanych nagłówków. Każda TU jest
analizowana niezależnie, a~linker łączy otrzymane obiekty.

Rozważmy projekt:
\begin{lstlisting}[caption={Trzy pliki klasycznego projektu C++: API, implementacja i klient.}]
// Plik: math.hpp
#ifndef CPP17_MATH_HPP
#define CPP17_MATH_HPP
int square(int x);
#endif

// Plik: math.cpp
#include "math.hpp"
int square(int x) { return x * x; }

// Plik: main.cpp
#include "math.hpp"
#include <iostream>
int main() { std::cout << square(7) << '\n'; }
\end{lstlisting}

\begin{lstlisting}[language=bash,caption={Niezalezne budowanie TU i koncowe linkowanie.}]
g++ -std=c++20 -Wall -Wextra -c math.cpp -o math.o
g++ -std=c++20 -Wall -Wextra -c main.cpp -o main.o
g++ main.o math.o -o app
./app                        # 49
\end{lstlisting}

\code{math.hpp} dostarcza deklarację klientowi i~implementacji,
ale nie jest zwykle oddzielnym obiektem linkowanym. Zmiana samego
\code{math.cpp} nie musi wymuszać ponownej kompilacji \code{main.cpp},
o~ile nie zmienił się kontrakt w~nagłówku lub pozostałe zależności.

\subsection{Przykładowe komendy w~MSVC}

W~\emph{Developer Command Prompt for Visual Studio} można zbudować
analogiczny przykład za pomocą \code{cl.exe} i~\code{link.exe}:

\begin{lstlisting}[language=bash,caption={Kompilacja dwóch plików w środowisku MSVC/Windows.}]
cl /std:c++20 /EHsc /c math.cpp
cl /std:c++20 /EHsc /c main.cpp
link main.obj math.obj /OUT:app.exe
app.exe
\end{lstlisting}

Przy budowaniu programów C++ w~Visual Studio warto również
kontrolować zgodność wersji MSVC, ustawienia biblioteki uruchomieniowej
(np.~\code{/MD} wobec \code{/MT}) i~konfigurację Debug/Release.
Takie flagi nie mają bezpośrednich odpowiedników narzuconych przez
standard języka; należą do konkretnego środowiska.

\section{Nagłówki, makra i~zasada jednej definicji}
\label{sec:cpp17-headers-odr}

\subsection{Include guards i~\texttt{\#pragma once}}

Zawartość \code{\#include} jest w~uproszczeniu wstawiana do źródła
przez preprocesor. Wielokrotne dołączenie tego samego nagłówka
w~jednej TU mogłoby więc prowadzić do powtórzenia definicji klasy.
Chronimy się przez \emph{include guard}:

\begin{lstlisting}[caption={Klasyczna os?ona nag?owka.}]
#ifndef LIB_WIDGET_HPP
#define LIB_WIDGET_HPP

struct Widget {
    int value{};
};

#endif // LIB_WIDGET_HPP
\end{lstlisting}

W~praktyce popularna jest też dyrektywa \code{\#pragma once},
obsługiwana przez najważniejsze kompilatory, lecz niestandaryzowana
przez ISO C++. Oba rozwiązania chronią \emph{jedną TU} przed
powtórnym dołączeniem pliku. Nie rozwiązują automatycznie problemu
\emph{wielu definicji tego samego nie-inline symbolu w~różnych TU}.

\subsection{ODR: deklaracja kontra definicja}

Deklaracja wprowadza nazwę i~informację o~typie, natomiast definicja
zwykle dostarcza byt danej funkcji lub obiektu. Deklaracji pewnej
funkcji może być wiele, ale \emph{nie-inline} funkcja z~zewnętrzną
łącznością, używana w~programie w~sposób wymagający definicji,
powinna mieć jedną definicję w~programie. Reguły ODR zawierają
wyjątki pozwalające na identyczne definicje w~różnych TU, np. dla
odpowiednich funkcji \code{inline}, klas i~szablonów.

\begin{lstlisting}[caption={Typowa pu?apka: definicja zwyk?ej funkcji w nag?owku.}]
// BAD: wrong.hpp, dolaczony do dwoch plikow .cpp
int twice(int x) { return 2 * x; }
// Dwie zewnetrzne definicje -> naruszenie ODR.
\end{lstlisting}

\begin{lstlisting}[caption={Definicja inline przeznaczona do nag?owka.}]
// Plik: twice.hpp
#pragma once
inline int twice(int x) { return 2 * x; }
\end{lstlisting}

\code{inline} nie nakazuje kompilatorowi wstawienia instrukcji funkcji
bezpośrednio w~miejscu wywołania. Jego znaczenie dla organizacji
programu polega przede wszystkim na dopuszczeniu definicji w~wielu TU
zgodnie z~ODR (z~odpowiednimi warunkami). Wszystkie takie definicje
muszą być zgodne i~spełniać reguły ODR; samo dopisanie \code{inline}
nie naprawia różniących się implementacji w~dwóch plikach.

\subsection{Zmienne \texttt{extern}, \texttt{inline}, \texttt{const} i~\texttt{static}}

\begin{lstlisting}[caption={Jedna definicja zmiennej globalnej, wiele deklaracji.}]
// Plik: settings.hpp
#pragma once
extern int request_limit;       // Tylko deklaracja

// Plik: settings.cpp
#include "settings.hpp"
int request_limit = 100;        // Definicja

// Plik: main.cpp
#include "settings.hpp"
int main() { return request_limit == 100 ? 0 : 1; }
\end{lstlisting}

Od C++17 w~nagłówkach można również definiować \emph{zmienne inline}:
\begin{lstlisting}[caption={Zmienna inline w nag?owku (C++17).}]
#pragma once
inline constexpr int default_timeout_ms = 250;
inline int global_counter = 0;
\end{lstlisting}

W~tym przypadku odpowiednie definicje we wszystkich TU mogą
wskazywać na ten sam obiekt o~zewnętrznej łączności. W~kontraście
zmienna \code{const} zdefiniowana w~przestrzeni nazw bez \code{extern}
ma w~typowym przypadku łączność wewnętrzną; oddzielne TU mogą mieć
wówczas odrębne egzemplarze. \code{static} na poziomie przestrzeni
nazw także zapewnia łączność wewnętrzną, lecz na współczesnym C++
zwykle preferujemy anonimową przestrzeń nazw.

\begin{lstlisting}[caption={Nazwy pomocnicze ograniczone do jednej TU.}]
// Plik: helper.cpp
namespace {
    int internal_counter = 0;
    int twice_local(int x) { return x * 2; }
}

int perform(int x) {
    ++internal_counter;
    return twice_local(x);
}
\end{lstlisting}

\subsection{Szablony: dlaczego definicje zwykle znajdują się w~nagłówkach?}

Kompilator instancjuje szablon dla określonych argumentów typowych
w~miejscu, w~którym wymagana jest jego definicja. Samo udostępnienie
deklaracji z~pliku \code{.hpp} i~ukrycie ogólnej definicji szablonu
w~oddzielnym \code{.cpp} może zatem prowadzić do błędu linkera.

\begin{lstlisting}[caption={Najprostszy, tradycyjny sposob udostepniania szablonu.}]
// Plik: multiply.hpp
#pragma once

template<class T>
T multiply(T a, T b) {
    return a * b;
}
\end{lstlisting}

Można także jawnie instancjować \emph{konkretne} typy w~pliku
implementacji. Oddzielne jednostki translacji nie muszą wtedy
instancjować tego samego kodu niezależnie:

\begin{lstlisting}[caption={Jawna instancjacja okreslonego wariantu szablonu.}]
// Plik: operation.hpp
#pragma once
template<class T>
T add(T a, T b);
extern template int add<int>(int, int);

// Plik: operation.cpp
#include "operation.hpp"
template<class T>
T add(T a, T b) { return a + b; }
template int add<int>(int, int);

// Plik: main.cpp
#include "operation.hpp"
int main() { return add(10, 32) == 42 ? 0 : 1; }
\end{lstlisting}

W~tym przykładzie klient może korzystać z~\code{add<int>}, ale
instancjacja innych typów nie zostaje mu udostępniona automatycznie.
Definicje szablonów umieszczane w~interfejsach modułów podlegają
nieco innemu modelowi widoczności, który omówimy dalej.

\subsection{Naruszenie ODR bez łatwego komunikatu}

Nie wszystkie naruszenia ODR skutkują natychmiastowym błędem
linkera. Szczególnie groźne jest warunkowe definiowanie tej samej
klasy w~dwóch TU w~inny sposób, np. przy użyciu sprzecznych
makr kompilacji:

\begin{lstlisting}[caption={Przyk?ad niezgodnych definicji typu (celowo b?edny).}]
// Plik: config.hpp
struct Config {
#ifdef EXTRA_FIELD
    int extra;
#endif
    int value;
};

// Jedna TU z -DEXTRA_FIELD, inna bez tej flagi:
// rozne definicje klasy -> naruszenie ODR.
\end{lstlisting}

Taki program może się nawet zlinkować, lecz pozostaje niepoprawny.
Dlatego opcje zmieniające publiczne definicje klas i~struktur muszą
być spójne we wszystkich TU danego programu.

\section{Symbole, łączność i~typowe błędy linkera}
\label{sec:cpp17-symbols}

\subsection{Łączność a~widoczność symbolu}

\emph{Linkage} w~języku C++ mówi, czy deklaracje w~różnych zakresach
i~TU mogą nazywać tę samą encję. W~praktyce spotykamy między innymi
łączność zewnętrzną, wewnętrzną, a~w~jednostkach modułów także
łączność modułową. Nie należy jednak utożsamiać tej własności
języka z~\emph{visibility} na poziomie ELF ani z~eksportem symbolu
z~biblioteki dynamicznej.

\begin{center}
\begin{tabularx}{\textwidth}{@{}lX@{}}
\toprule
Mechanizm & Co określa?\\
\midrule
\code{static} w~przestrzeni nazw & Łączność wewnętrzna nazwy w~TU.\\
\code{extern} & M.in. deklarowanie obiektu definiowanego w~innym miejscu.\\
Anonimowa przestrzeń nazw & Wewnętrzna łączność elementów pomocniczych.\\
\code{inline} & Warunki wielokrotnych definicji zgodnych z~ODR.\\
\code{-fvisibility=hidden} & Domyślna widoczność symboli w~typowych bibliotekach ELF.\\
\code{__declspec(dllexport)} & Udostępnianie symboli biblioteki DLL w~MSVC.\\
\bottomrule
\end{tabularx}
\end{center}

\subsection{Name mangling i~\texttt{extern~\textquotedbl{}C\textquotedbl{}}}

C++ dopuszcza przeciążanie funkcji, przestrzenie nazw oraz rozbudowane
typy szablonowe. Aby móc odróżniać symbole po linkowaniu,
implementacje zwykle kodują część tych informacji w~nazwach
binarnych (\emph{name mangling}). Dokładny sposób kodowania nie jest
ustalony przez ISO C++.

\begin{lstlisting}[caption={Przeciazanie i punkt wejscia C ABI.}]
namespace calc {
    int f(int x) { return x + 1; }
    double f(double x) { return x + 1.0; }
}

extern "C" int api_add(int a, int b) {
    return a + b;
}
\end{lstlisting}

Na typowej platformie ELF/Itanium ABI funkcje przeciążone otrzymują
różne symbole zakodowane przez kompilator, podczas gdy \code{api_add}
ma nazwę zgodną z~łącznością języka C. Nie jest to gwarancja pełnej
zgodności ABI z~dowolną biblioteką C: znaczenie mają także deklaracje
typów i~konwencja wywołania ustalona dla danej platformy.

\begin{lstlisting}[language=bash,caption={Demanglowanie i inspekcja symboli (przyk?ady narzedziowe).}]
g++ -std=c++20 -c overload.cpp -o overload.o
nm overload.o
nm -C overload.o
c++filt _Z3fooi      # Ilustracja popularnego schematu manglingu
\end{lstlisting}

\subsection{\texttt{undefined reference}, \texttt{multiple definition}}

Błąd \code{undefined reference} oznacza zwykle, że linker znalazł
odwołanie do symbolu, lecz nie znalazł odpowiedniej definicji.
Przykładowe przyczyny to pominięcie \code{math.o}, pominięcie
biblioteki, niezgodny podpis funkcji albo brak jawnej instancjacji
szablonu. Z~kolei \code{multiple definition} sugeruje m.in. zwykłą
funkcję lub zmienną zdefiniowaną w~nagłówku dołączanym do wielu TU.

Przydatna lista kontrolna:
\begin{enumerate}
\item Sprawdź, czy wszystkie wymagane \code{.cpp} są kompilowane i~linkowane.
\item Porównaj podpis funkcji: przestrzeń nazw, typy, kwalifikatory i~ABI.
\item Upewnij się, że deklaracja i~definicja pochodzą z~tego samego interfejsu.
\item Sprawdź kolejność bibliotek statycznych i~ich zależności.
\item Dla klas polimorficznych sprawdź brakujące definicje funkcji
      wirtualnych, w~tym destruktora wirtualnego, który może wymagać
      definicji nawet jako funkcja \code{= 0}.
\item Dla szablonów sprawdź dostępność definicji w~punkcie instancjacji
      lub odpowiednią instancjację jawną.
\end{enumerate}

\section{Biblioteki statyczne}
\label{sec:cpp17-static-libraries}

\emph{Biblioteka statyczna} jest typowo archiwum plików obiektowych.
W~Linux ma zwykle rozszerzenie \code{.a}, w~MSVC \code{.lib}.
Biblioteka sama nie jest osobnym procesem ani automatycznie tworzonym
wątkiem. Linker wybiera z~niej obiekty potrzebne do rozwiązania
zależności i~tworzy wynikowy program lub inną bibliotekę.

\begin{lstlisting}[language=bash,caption={Budowanie i uzycie biblioteki statycznej (Linux/GCC).}]
g++ -std=c++20 -O2 -c math.cpp -o math.o
ar rcs libmathstuff.a math.o

g++ -std=c++20 -O2 -c main.cpp -o main.o
g++ main.o -L. -lmathstuff -o app
./app
\end{lstlisting}

W~typowym poleceniu GNU linkera biblioteki archiwalne warto umieszczać
\emph{po} obiektach, które odwołują się do ich symboli. Przykładowo
\code{g++ main.o -lmathstuff} jest zazwyczaj właściwszym porządkiem
niż \code{g++ -lmathstuff main.o}. Linkowanie zależności cyklicznych
bibliotek statycznych może wymagać ponownego wskazania biblioteki
lub użycia grupy linkera (np.~\code{--start-group}/\code{--end-group}
w~GNU ld); lepiej jednak unikać takich zależności architektonicznie.

\nota{Statyczna biblioteka nie oznacza automatycznie w~pełni
statycznego programu. Aplikacja może dalej korzystać z~bibliotek
współdzielonych systemu i~standardowej biblioteki C++.}

\section{Biblioteki współdzielone i~dynamiczne ładowanie}
\label{sec:cpp17-shared-libraries}

\subsection{Budowanie biblioteki współdzielonej}

Na systemach ELF biblioteki współdzielone są zazwyczaj plikami
\code{.so}. Obiekty wykorzystywane w~takich bibliotekach zwykle
buduje się jako kod niezależny od położenia w~pamięci (\emph{PIC}).

\begin{lstlisting}[language=bash,caption={Biblioteka .so oraz program klienta (Linux/GCC).}]
g++ -std=c++20 -O2 -fPIC -c math.cpp -o math_pic.o
g++ -shared math_pic.o -o libmathstuff.so

g++ -std=c++20 -O2 main.cpp -L. -lmathstuff \
    '-Wl,-rpath,$ORIGIN' -o app
./app
\end{lstlisting}

Wyrażenie \code{\$ORIGIN} odnosi się w~typowym loaderze ELF do
katalogu, w~którym znajduje się ładowany obiekt. Apostrofy w~poleceniu
powłoki zapobiegają wcześniejszej ekspansji \code{\$ORIGIN} przez
powłokę. W~projekcie produkcyjnym należy przemyśleć ścieżki
instalacyjne, pakowanie bibliotek i~bezpieczeństwo ładowania.

Biblioteki współdzielone pozwalają wielu procesom mapować ich kod
z~tego samego pliku oraz wymieniać implementację bez konieczności
pełnego linkowania programu w~czasie jego budowania. Nie gwarantuje
to jednak, że \emph{dowolna} nowa wersja zachowa zgodność ABI.

\subsection{Wyszukiwanie bibliotek: RPATH/RUNPATH, SONAME}

Na Linux/ELF linker i~loader korzystają z~mechanizmów takich jak
\emph{SONAME}, \emph{RPATH} i~\emph{RUNPATH}. Reguły wyszukiwania
zależą od systemu i~konfiguracji; środowisko wykonania może również
uwzględniać ścieżki systemowe lub zmienne środowiskowe.

\begin{lstlisting}[language=bash,caption={Diagnostyka zaleznosci dynamicznych (Linux).}]
readelf -d ./app | grep -E 'NEEDED|RPATH|RUNPATH'
ldd ./app                  # Tylko dla zaufanych plikow
readelf -Ws libmathstuff.so
\end{lstlisting}

Programów nieznanego pochodzenia nie należy lekkomyślnie przekazywać
do narzędzi, które mogą uruchamiać kod lub loader; do statycznej
inspekcji lepiej nadaje się \code{readelf}.

\subsection{Eksportowanie symboli i~kontrola API}

Domyślne eksportowanie zbyt wielu symboli może powodować kolizje
nazw i~powiększać powierzchnię ABI. W~ELF można budować bibliotekę
z~\code{-fvisibility=hidden}, a~publiczne symbole jawnie oznaczać:

\begin{lstlisting}[caption={Przenosne w granicach GCC/Clang/MSVC makro eksportu.}]
// Plik: api.hpp
#pragma once

#if defined(_WIN32)
  #if defined(MYLIB_BUILD_SHARED)
    #define MYLIB_API __declspec(dllexport)
  #elif defined(MYLIB_USE_SHARED)
    #define MYLIB_API __declspec(dllimport)
  #else
    #define MYLIB_API
  #endif
#elif defined(__GNUC__) || defined(__clang__)
  #define MYLIB_API __attribute__((visibility("default")))
#else
  #define MYLIB_API
#endif

extern "C" MYLIB_API int mylib_add(int, int);
\end{lstlisting}

W~MSVC zwykle stosuje się projekt DLL oraz odpowiednią bibliotekę
importową; na macOS biblioteki mają najczęściej rozszerzenie
\code{.dylib}. Flagi \code{-fPIC}, \code{-shared} i~\code{-Wl,...}
w~powyższej postaci dotyczą konkretnego toolchainu, nie ISO C++.

\subsection{\texttt{dlopen}, \texttt{dlsym} i~granica interfejsu wtyczki}

Bibliotekę można także załadować podczas działania procesu, bez
linkowania jej wprost z~aplikacją. Na POSIX często używa się
\code{dlopen}/\code{dlsym}, natomiast na Windows odpowiedników
\code{LoadLibrary}/\code{GetProcAddress}.

\begin{lstlisting}[caption={Prosta wtyczka C ABI: biblioteka i klient POSIX (dwa pliki).}]
// Plik: plugin.cpp
extern "C" int plugin_square(int x) {
    return x * x;
}

// Plik: loader.cpp
#include <dlfcn.h>
#include <iostream>
#include <stdexcept>

int main() {
    void* lib = dlopen("./libplugin.so", RTLD_NOW);
    if (!lib) throw std::runtime_error(dlerror());

    using Fn = int (*)(int);
    void* address = dlsym(lib, "plugin_square");
    if (!address) {
        dlclose(lib);
        throw std::runtime_error("Missing plugin_square");
    }

    // POSIX dopuszcza konwersje dlsym -> wskaznik do funkcji.
    Fn square = reinterpret_cast<Fn>(address);
    std::cout << square(7) << '\n'; // 49
    dlclose(lib);
}
\end{lstlisting}

\begin{lstlisting}[language=bash,caption={Budowanie przyk?adu wtyczki (Linux/ELF).}]
g++ -std=c++20 -fPIC -shared plugin.cpp -o libplugin.so
g++ -std=c++20 loader.cpp -ldl -o loader
./loader
\end{lstlisting}

\nota{To przykład dla POSIX, a~nie w~pełni przenośna konstrukcja
ISO C++. Również rzutowanie wyniku \code{dlsym} na wskaźnik funkcji
opiera się na wymaganiach POSIX, a~nie na ogólnej gwarancji C++.
Dodatkowo nie wolno wyładowywać biblioteki, dopóki jakakolwiek funkcja
lub obiekt z~tej biblioteki nadal może być użyty. Przy bardziej
złożonych wtyczkach potrzebny jest kontrakt wersji ABI i~zarządzania
czasem życia.}

\section{ABI --- Application Binary Interface}
\label{sec:cpp17-abi}

\subsection{API a~ABI: dwa różne kontrakty}

\emph{API} opisuje przede wszystkim dostępne nazwy, typy, funkcje
i~semantykę na poziomie źródła. \emph{ABI} określa, jak te elementy
reprezentowane są w~kodzie maszynowym. ABI obejmuje między innymi:
\begin{itemize}
\item konwencje wywołania: przekazywanie parametrów, wartość zwracana,
      rejestry i~ułożenie stosu;
\item reprezentację i~wyrównanie obiektów, w~tym pól klas;
\item nazwy symboli C++, RTTI, mechanizmy wyjątków i~wirtualnych wywołań;
\item reguły udostępniania i~ładowania symboli bibliotek;
\item ABI standardowej biblioteki oraz ustawienia kompilatora.
\end{itemize}

Dwie biblioteki mogą być zgodne źródłowo i~jednocześnie binarnie
niezgodne. W~szczególności identyczny nagłówek nie oznacza, że plik
\code{.so} zbudowany innym toolchainem lub z~innym zestawem opcji
na pewno współpracuje z~aplikacją.

\subsection{Zmiana układu obiektu łamie zgodność binarną}

Rozważmy klasę publiczną:
\begin{lstlisting}[caption={Zmiana prywatnych pol nadal moze zmienic ABI.}]
// Wersja biblioteki 1
class Record {
    int id_{};
public:
    int id() const;
};

// Wersja biblioteki 2
class Record {
    int id_{};
    double score_{}; // Nawet prywatne pole zmienia uklad.
public:
    int id() const;
};
\end{lstlisting}

Klient kompiluje np.~\code{sizeof(Record)} według nagłówka,
który widział podczas kompilacji. Gdy zamienimy jedynie bibliotekę
na wersję z~innym układem \code{Record}, może dojść do naruszenia
kontraktu pamięci, nawet jeśli \code{id()} nadal ma taki sam podpis.
Podobne ryzyko dotyczy dodawania wirtualnych funkcji i~zmian
hierarchii dziedziczenia.

\subsection{Wzorzec PImpl}

Aby zmniejszyć zależności kompilacyjne i~ukryć układ implementacji,
można przechowywać wskaźnik do niekompletnego typu. Dzięki temu
zmiana \emph{pól implementacyjnych} często nie zmienia rozmiaru
publicznego obiektu.

\begin{lstlisting}[caption={PImpl: plik widget.hpp.}]
#pragma once
#include <memory>

class Widget {
public:
    Widget();
    ~Widget();

    Widget(Widget&&) noexcept;
    Widget& operator=(Widget&&) noexcept;

    Widget(const Widget&) = delete;
    Widget& operator=(const Widget&) = delete;

    int value() const;
    void set_value(int value);

private:
    struct Impl;
    std::unique_ptr<Impl> impl_;
};
\end{lstlisting}

\begin{lstlisting}[caption={PImpl: plik widget.cpp.}]
#include "widget.hpp"
#include <utility>

struct Widget::Impl {
    int value = 0;
    // Tu moga pojawiac sie dalsze pola.
};

Widget::Widget() : impl_(std::make_unique<Impl>()) {}
Widget::~Widget() = default;
Widget::Widget(Widget&&) noexcept = default;
Widget& Widget::operator=(Widget&&) noexcept = default;

int Widget::value() const { return impl_->value; }
void Widget::set_value(int value) { impl_->value = value; }
\end{lstlisting}

Destruktor jest zdefiniowany w~\code{widget.cpp}, gdzie \code{Impl}
jest kompletny; to istotne przy \code{std::unique_ptr<Impl>}.
PImpl nie zapewnia automatycznie pełnej stabilności ABI: nadal
liczy się zgodność kompilatora, standardowej biblioteki, eksportu
symboli, konwencji wywołania i~pozostałych publicznych deklaracji.
Dodatkowy wskaźnik i~alokacja na stercie mogą zwiększać koszt dostępu.

\subsection{C ABI jako ostrożna granica interoperacyjności}

Dla wtyczek i~bibliotek wykorzystywanych z~różnych języków często
projektuje się wąski interfejs oparty na funkcjach \code{extern "C"}
oraz nieprzezroczystych uchwytach (\emph{opaque handles}). Zamiast
przekazywać \code{std::string} lub klasy z~funkcjami wirtualnymi,
przekazuje się proste typy i~wskaźniki.

\begin{lstlisting}[caption={Przyk?adowy szkic stabilniejszego interfejsu binarnego.}]
#ifdef __cplusplus
extern "C" {
#endif

struct engine_handle;

engine_handle* engine_create(void);
void engine_destroy(engine_handle* p);
int engine_eval(engine_handle* p, double x, double* result);

#ifdef __cplusplus
}
#endif
\end{lstlisting}

Podpis funkcji nadal wymaga uzgodnienia ABI danej platformy,
konwencji wywołania i~kontraktu błędów. Przekraczanie takiej granicy
wyjątkami C++, pamięcią zwalnianą przez niezgodny alokator lub obiektami
standardowej biblioteki jest niebezpieczne. Bezpieczniej zwracać
kody błędu, a~alokację i~zwalnianie łączyć w~pary po tej samej
stronie interfejsu.

\subsection{Wersje kompilatora i~biblioteki standardowej}

Istnieją różne rodziny ABI C++ (np. rodzina Itanium ABI stosowana na
wielu platformach uniksowych oraz ABI MSVC w~Windows). Nawet w~ramach
jednej rodziny znaczenie mają flagi kompilatora, ustawienia standardowej
biblioteki, opcje wyjątków, RTTI i~ABI konkretnych typów.
Przykładowo GCC/libstdc++ używa mechanizmu tzw. \emph{dual ABI}
dla niektórych typów standardowych (m.in. \code{std::string}); mieszanie
niezgodnych ustawień może skutkować błędami linkera lub działania.
Nie należy zakładać, że biblioteki zbudowane dowolnymi wersjami
GCC, Clanga i~MSVC są wymienne binarnie.

\section{System budowania: CMake i~projekt wieloplikowy}
\label{sec:cpp17-cmake}

\subsection{Dlaczego targety są ważniejsze niż globalne flagi?}

Ręczne wywoływanie kompilatora dobrze tłumaczy mechanikę, lecz dla
większych projektów potrzebujemy grafu zależności, wariantów
Debug/Release, testów i~kontroli opcji kompilatora. Współczesny
CMake modeluje budowanie przez \emph{targety}: bibliotekę, program,
interfejsową bibliotekę nagłówkową i~ich właściwości.

Przykładowa struktura:
\begin{lstlisting}[caption={Uk?ad plikow klasycznej biblioteki CMake.}]
project/
  CMakeLists.txt
  include/
    mathstuff/math.hpp
  src/
    math.cpp
    main.cpp
\end{lstlisting}

W~tej konkretnej strukturze pliki źródłowe mogą zaczynać się od
\code{\#include "mathstuff/math.hpp"}, a~\code{main.cpp} również
korzysta z~tej samej ścieżki nagłówka. To dzięki
\code{target_include_directories} kompilator znajduje katalog
\code{include} bez wpisywania ręcznych flag \code{-I}.

\begin{lstlisting}[language={},caption={CMakeLists.txt dla klasycznej biblioteki.}]
cmake_minimum_required(VERSION 3.20)
project(MathStuff LANGUAGES CXX)

add_library(mathstuff STATIC src/math.cpp)
target_compile_features(mathstuff PUBLIC cxx_std_20)
target_include_directories(mathstuff PUBLIC
    $<BUILD_INTERFACE:${CMAKE_CURRENT_SOURCE_DIR}/include>
    $<INSTALL_INTERFACE:include>
)

add_executable(app src/main.cpp)
target_link_libraries(app PRIVATE mathstuff)
\end{lstlisting}

\begin{lstlisting}[language=bash,caption={Konfigurowanie i budowanie projektu z CMake.}]
cmake -S . -B build -DCMAKE_BUILD_TYPE=Debug
cmake --build build -j
./build/app
\end{lstlisting}

\code{PRIVATE}, \code{PUBLIC} i~\code{INTERFACE} określają
\emph{propagację wymagań użycia} (usage requirements) między
targetami:
\begin{center}
\begin{tabularx}{\textwidth}{@{}lX@{}}
\toprule
Słowo & Znaczenie\\
\midrule
\code{PRIVATE} & Wymaganie tylko na potrzeby budowania danego targetu.\\
\code{INTERFACE} & Wymaganie wyłącznie dla jego konsumentów.\\
\code{PUBLIC} & Wymaganie dla targetu i~jego konsumentów.\\
\bottomrule
\end{tabularx}
\end{center}

Gdy publiczny nagłówek biblioteki wymaga pewnego katalogu
nagłówków, ten katalog powinien zwykle propagować się do klientów.
Gdy jednak biblioteka korzysta z~zależności jedynie wewnątrz
\code{.cpp}, często wystarcza \code{PRIVATE}. Właściwy zakres
zależy też od tego, czy target jest biblioteką statyczną lub
współdzieloną i~jakie zależności są konieczne przy ostatecznym
linkowaniu.

\subsection{Biblioteka nagłówkowa i~\texttt{INTERFACE}}

\begin{lstlisting}[language={},caption={Header-only target i wymagania dla klientow.}]
add_library(my_headers INTERFACE)
target_include_directories(my_headers INTERFACE
    ${CMAKE_CURRENT_SOURCE_DIR}/include
)
target_compile_features(my_headers INTERFACE cxx_std_20)

add_executable(tool src/tool.cpp)
target_link_libraries(tool PRIVATE my_headers)
\end{lstlisting}

Taki target nie tworzy archiwum \code{.a}. Reprezentuje zestaw
wymagań i~publiczne nagłówki używane przez konsumentów.

\subsection{Diagnostyka, sanitizatory, LTO i~build types}

\begin{lstlisting}[language=bash,caption={Przyk?adowe flagi diagnostyczne GCC/Clang.}]
g++ -std=c++20 -O0 -g -Wall -Wextra -Wpedantic \
    -fsanitize=address,undefined main.cpp -o app_debug

g++ -std=c++20 -O2 -flto main.cpp math.cpp -o app_lto
\end{lstlisting}

\emph{Link-Time Optimization} (LTO) może umożliwiać optymalizacje
pomiędzy TU, np. dodatkowe wstawianie funkcji do miejsc wywołania.
Jest to jednak kosztowna funkcja toolchainu; jej efekty zależą od
programu i~opcji. Profilowanie i~powtarzalne pomiary są ważniejsze
niż sam fakt włączenia \code{-O3} lub \code{-flto}.

\section{Moduły C++20: nowy mechanizm organizacji kodu}
\label{sec:cpp17-modules}

\subsection{Motywacja: problem tekstowego dołączania nagłówków}

W~klasycznym modelu ten sam popularny nagłówek może być parsowany
oddzielnie w~wielu TU. Makra i~kolejność \code{\#include} wpływają
na wynik preprocessingu. Moduły C++20 umożliwiają jawne eksportowanie
deklaracji i~importowanie ich bez tekstowego wklejania całych
interfejsów do każdego klienta.

Zaletami modułów mogą być izolacja makr, bardziej jednoznaczne
zależności i~krótsze czasy kompilacji. Nie ma jednak gwarancji,
że \emph{każdy} program po migracji przyspieszy, a~wdrożenie
wymaga wsparcia kompilatora i~systemu budowania.

\subsection{Pierwszy moduł: \texttt{export module} i~\texttt{import}}

Rozważmy dwa pliki:
\begin{lstlisting}[caption={Interfejs modu?u geometry: plik geometry.cppm.}]
export module geometry;

export int square(int x) {
    return x * x;
}

int hidden_helper(int x) {
    return x + 1;
}
\end{lstlisting}

\begin{lstlisting}[caption={Klient modu?u: plik main.cpp.}]
#include <iostream>
import geometry;

int main() {
    std::cout << square(7) << '\n'; // 49
    // hidden_helper(7); // Nie jest eksportowana.
}
\end{lstlisting}

\code{export module geometry;} deklaruje \emph{główną jednostkę
interfejsu modułu} (\emph{primary module interface unit}).
\code{export int square(...)} udostępnia deklarację klientowi.
Samo \code{import geometry;} nie oznacza linkowania kodu funkcji
z~powietrza: jeśli moduł generuje kod obiektowy, trzeba go również
uwzględnić podczas linkowania.

\subsection{BMI/CMI: moduł nie jest biblioteką binarną}

Kompilator wytwarza dla interfejsu modułu własną reprezentację,
często nazywaną \emph{BMI} (\emph{Binary Module Interface}) lub
\emph{CMI} (\emph{Compiled Module Interface}). W~Clangu spotyka się
pliki \code{.pcm}, w~GCC \code{.gcm}, a~w~MSVC pliki \code{.ifc}.
Nie są to uniwersalne, przenośne między kompilatorami artefakty
podobne do biblioteki \code{.a}. Importer musi znać interfejs
przetworzony przez zgodną konfigurację kompilatora.

Zależności muszą być budowane w~poprawnym porządku:
\[
\text{źródło interfejsu modułu}
\longrightarrow \text{BMI/CMI}
\longrightarrow \text{kompilacja importera}
\longrightarrow \text{linkowanie obiektów}.
\]

\subsection{Rzeczywiste komendy: Clang i~GCC}

Dla prostego modułu \code{geometry.cppm} w~Clangu obsługującym
standardowe moduły można zastosować:
\begin{lstlisting}[language=bash,caption={Budowanie przyk?adowego modu?u w Clangu (tryb BMI).}]
clang++ -std=c++20 geometry.cppm --precompile -o geometry.pcm
clang++ -std=c++20 -c geometry.pcm -o geometry.o
clang++ -std=c++20 -fmodule-file=geometry=geometry.pcm \
    -c main.cpp -o main.o
clang++ geometry.o main.o -o app
./app
\end{lstlisting}

Dla GCC~14 (wsparcie modułów C++20 było w~tej wersji nadal rozwijane)
trzeba dodatkowo włączyć eksperymentalny mechanizm
\code{-fmodules-ts}. W~przykładach z~GCC~14 przydatne bywa jawne
wskazanie języka przy rozszerzeniu \code{.cppm}:
\begin{lstlisting}[language=bash,caption={Budowanie tego samego modu?u z GCC 14.}]
g++ -std=c++20 -fmodules-ts -x c++ -c geometry.cppm \
    -o geometry.o
# GCC tworzy potrzebne dane interfejsu w gcm.cache/.
g++ -std=c++20 -fmodules-ts -c main.cpp -o main.o
g++ geometry.o main.o -o app
./app
\end{lstlisting}

Dokładne opcje i~formaty modułów mogą zmieniać się między wersjami
narzędzi. W~dużym projekcie lepiej powierzyć zależności systemowi
budowania obsługującemu skanowanie importów, zamiast ręcznie
zarządzać plikami BMI.

\subsection{Jednostka interfejsu i~jednostka implementacji}

Moduł może rozdzielać interfejs od implementacji bez tradycyjnego
nagłówka. Jednostka implementacji deklaruje \code{module name;},
a~nie \code{export module name;}:

\begin{lstlisting}[caption={Dwa pliki jednego modu?u: interface.cppm i implementation.cpp.}]
// Plik: engine.cppm
export module engine;
export int run(int value);

// Plik: engine.cpp
module engine;
int run(int value) { return 2 * value; }
\end{lstlisting}

Importer widzi tylko deklaracje eksportowane przez interfejs.
Jednostkę implementacji również trzeba skompilować i~zlinkować.
Konkretny system budowania musi uwzględnić fakt, że jednostka
implementacji importuje interfejs swojego modułu niejawnie.

\subsection{Global module fragment i~nagłówki klasyczne}

Tradycyjny nagłówek może wymagać makr konfigurujących jego treść.
W~takim przypadku można skorzystać z~\emph{global module fragment}:

\begin{lstlisting}[caption={Global module fragment: tradycyjny include przed modu?em.}]
module;
#include <string>
#include <utility>

export module greeting;

export std::string greet(std::string name) {
    return "Hello " + std::move(name);
}
\end{lstlisting}

Dyrektywa \code{module;} rozpoczyna fragment globalny, a~następna
deklaracja \code{export module greeting;} rozpoczyna właściwą
jednostkę modułu. Importowanie nagłówków jednostkowych przez
\code{import <string>;} jest innym mechanizmem, tzw. \emph{header units},
których dostępność zależy od toolchainu. Nie utożsamiajmy go z~nazwanym
modułem standardowej biblioteki \code{import std;}.

\subsection{Partycje i~ponowny eksport importu}

Duży moduł można podzielić na partycje. Nazwa po dwukropku oznacza
partycję wewnątrz jednego modułu, a~nie odrębny nazwany moduł.

\begin{lstlisting}[caption={Dwie jednostki interfejsu modu?u z partycja.}]
// Plik: math_detail.cppm
export module mathlib:detail;
export int triple(int x) { return 3 * x; }

// Plik: mathlib.cppm
export module mathlib;
export import :detail;
export int times_six(int x) { return 2 * triple(x); }

// Plik: main.cpp
import mathlib;
int main() { return times_six(7) == 42 ? 0 : 1; }
\end{lstlisting}

Klient importuje \code{mathlib}, a~nie \code{mathlib:detail}.
W~tym przykładzie \code{export import :detail;} re-eksportuje
eksporty partycji. Wymaga to prawidłowej kolejności budowania
\emph{najpierw partycji, potem głównego interfejsu, a~na końcu klienta}.

\subsection{Łączność modułowa i~ODR}

C++20 wprowadza reguły własności deklaracji (\emph{module ownership}).
Deklaracje przypisane do nazwanego modułu należą do niego, a~pewne
nieeksportowane nazwy uzyskują \emph{module linkage}. Różni się to od
mechanicznego zastąpienia \code{\#include} przez \code{import}.

Przykładowo dodanie \code{export} nie sprawia automatycznie, że
interfejs C++ będzie zgodny ABI między kompilatorami. Ponadto
źródłowy kod szablonów nadal musi być semantycznie dostępny
odpowiedniemu mechanizmowi instancjacji, choć nie musi być już
tekstowo wklejany przez \code{\#include}.

\subsection{CMake i~\texttt{FILE\_SET CXX\_MODULES}}

Od CMake~3.28 dostępne są mechanizmy skanowania modułów i~specjalny
typ zbioru plików \code{CXX_MODULES}. Przykład projektu:

\begin{lstlisting}[language={},caption={CMakeLists.txt projektu korzystajacego z modu?ow C++20.}]
cmake_minimum_required(VERSION 3.28)
project(GeometryModules LANGUAGES CXX)

add_library(geometry)
target_compile_features(geometry PUBLIC cxx_std_20)
target_sources(geometry PUBLIC
    FILE_SET CXX_MODULES FILES geometry.cppm
)

add_executable(app main.cpp)
target_link_libraries(app PRIVATE geometry)
\end{lstlisting}

\begin{lstlisting}[language=bash,caption={Konfiguracja modu?ow przy uzyciu generatora Ninja.}]
cmake -S . -B build -G Ninja
cmake --build build
./build/app
\end{lstlisting}

Ten przykład pokazuje \emph{składnię} CMake, ale jego zbudowanie
wymaga zgodnego zestawu wersji CMake, kompilatora, generatora i~skanera
zależności. W~szczególności kompilator akceptujący składnię modułów
niekoniecznie współpracuje z~automatycznym skanowaniem w~danej
konfiguracji CMake. Na współczesnych systemach używa się m.in.
Ninja (co najmniej 1.11 w~typowych konfiguracjach) i~odpowiednich
wersji Clanga lub MSVC. Obsługa modułów standardowej biblioteki oraz
\emph{header units} ma odrębne ograniczenia.

\subsection{\texttt{import std;} nie jest tym samym co \texttt{\#include}}

Sam standard C++20 wprowadza językowy mechanizm modułów, ale
\emph{nie zobowiązuje} każdej biblioteki standardowej do dostarczenia
nazwanego modułu \code{std} w~tym standardzie. \code{import std;}
jest standaryzowane w~C++23; obsługa zależy od kompilatora, dostawcy
biblioteki standardowej i~wersji systemu budowania.

\begin{lstlisting}[caption={Nazwany modu? standardowej biblioteki (C++23, zalezny od wsparcia).}]
import std;

int main() {
    std::vector<int> xs{1, 2, 3};
    std::println("size = {}", xs.size());
}
\end{lstlisting}

Powyższy przykład wykorzystuje także \code{std::println} z~C++23.
Nie należy automatycznie kompilować go z~\code{-std=c++20} ani
zakładać, że \code{import std;} działa w~każdym środowisku C++23.

\section{Optymalizacja procesu budowania}
\label{sec:cpp17-build-perf}

\subsection{Incremental build i~zależności nagłówkowe}

Zmiana publicznego nagłówka może wymusić ponowną kompilację wielu TU,
podczas gdy zmiana prywatnego \code{.cpp} zwykle dotyka mniejszej
części projektu. Dlatego ograniczanie niepotrzebnych dołączeń
nagłówków, stosowanie deklaracji wyprzedzających tam, gdzie są
prawidłowe, i~PImpl mogą poprawiać czas budowania.

Moduły zmieniają strukturę zależności, ale nie usuwają konieczności
ich śledzenia; importer musi otrzymać aktualną reprezentację
interfejsu modułu.

\subsection{Precompiled headers, unity builds, ccache}

Popularne techniki przyspieszania kompilacji to:
\begin{description}[style=nextline]
\item[Precompiled headers (PCH)]
Przechowywanie wyniku parsowania stabilnych, obszernych nagłówków.
Pomagają w~projektach z~dużą liczbą TU, ale wymagają spójności
opcji kompilacji i~zarządzania unieważnianiem pamięci podręcznej.
\item[Unity builds]
Łączenie kilku plików źródłowych we wspólną TU. Może ograniczać
narzut kompilacji, ale też ujawniać kolizje anonimowych przestrzeni
nazw, makr i~definicji o~łączności wewnętrznej.
\item[Compilation cache]
Narzędzia takie jak \code{ccache} unikają ponownego wykonywania
kompilacji, jeżeli odpowiednio zdefiniowane wejścia pozostają takie same.
\item[Parallel build]
Kompilowanie niezależnych TU równolegle (np. \code{cmake --build build -j})
wykorzystuje wiele rdzeni, ale zwiększa chwilowe użycie pamięci.
\end{description}

\subsection{Powtarzalność buildów i~zależności transytywne}

Dla bibliotek publicznych warto utrzymywać przejrzystą strukturę
zależności. Przykładowo użytkownik nie powinien polegać na tym,
że \code{\#include <vector>} ,,przypadkiem'' pojawiło się wskutek
dołączenia innego nagłówka. Każdy nagłówek powinien zawierać
potrzebne deklaracje w~sposób niezależny od kolejności dołączeń.

Warto również zapisywać wersje kompilatora i~bibliotek, flagi
wpływające na ABI, wariant budowania oraz sposób instalowania
artefaktów. \emph{Reproducible builds} wymagają dodatkowo kontroli
źródeł niedeterministyczności, np. znaczników czasu i~ścieżek
zapisywanych w~plikach wynikowych.

\section{Praktyczny przewodnik diagnostyczny}
\label{sec:cpp17-troubleshooting}

\begin{center}
\small
\begin{tabularx}{\textwidth}{@{}p{0.34\textwidth}X@{}}
\toprule
Objaw & Najbardziej prawdopodobne źródło problemu\\
\midrule
\code{undefined reference} & Brakujący plik obiektowy, biblioteka,
definicja albo instancjacja szablonu.\\
\code{multiple definition} & Niepoprawna definicja zwykłej
funkcji lub zmiennej w~nagłówku albo podwójnie linkowany obiekt.\\
\code{undefined reference to vtable} & Niezdefiniowana wymagana
funkcja wirtualna (często destruktor) lub brak pliku implementacji.\\
\code{cannot find -lfoo} & Brak biblioteki lub nieprawidłowe
katalogi przeszukiwania linkera.\\
\code{cannot open shared object file} & Problem ścieżki lub
zależności ładowania, niekoniecznie problem kompilacji.\\
Kompilacja klienta działa, program kończy się błędem & Możliwa
niezgodność ABI, naruszenie ODR, zły czas życia obiektu lub błąd programu.\\
Importer nie znajduje modułu & Nie zbudowano BMI/CMI, nie wskazano
ścieżki albo konfiguracja skanera nie obsługuje importu.\\
Import działa, ale linker zgłasza brak funkcji & Brak pliku
obiektowego modułu/implementacji lub potrzebnej biblioteki.\\
Zmiana nagłówka przebudowuje prawie wszystko & Silne zależności
nagłówkowe; rozważ redukcję dołączeń, PImpl, moduły lub PCH.\\
\bottomrule
\end{tabularx}
\end{center}

\subsection{Kiedy wybrać nagłówki, bibliotekę czy moduł?}

\begin{itemize}
\item \textbf{Header-only} sprawdza się dla małych narzędzi, szablonów
      i~funkcji \code{inline}; ceną mogą być zależności kompilacyjne.
\item \textbf{Biblioteka statyczna} upraszcza dystrybucję kodu obiektowego
      wewnątrz aplikacji, ale nie usuwa wymagań ABI w~czasie linkowania.
\item \textbf{Biblioteka współdzielona} pozwala współdzielić kod
      i~wymieniać implementację, lecz wymaga jasnej polityki ABI,
      widoczności symboli i~wersjonowania.
\item \textbf{Moduł C++20} służy do semantycznego podziału interfejsu
      źródłowego i~kontroli importów; nie jest alternatywą wobec
      procesu linkowania ani automatycznym formatem dystrybucji ABI.
\item \textbf{PImpl i~C API} rozważaj na granicach bibliotek o~długim
      cyklu życia, szczególnie gdy stabilność ABI ma duże znaczenie.
\end{itemize}

\section*{Podsumowanie rozdziału}
\addcontentsline{toc}{section}{Podsumowanie rozdziału}

Kompletny program C++ jest wynikiem wspólnej pracy preprocessora,
kompilatora, linkera i~loadera. Organizując projekt, trzeba osobno
rozumieć \emph{deklaracje i~definicje}, \emph{ODR i~łączność},
\emph{symbole i~biblioteki} oraz \emph{ABI}. Nagłówki i~moduły
pomagają udostępniać interfejsy źródłowe; biblioteki statyczne
i~współdzielone dostarczają kod do połączenia lub załadowania.

Dobra organizacja projektu polega na jasnym rozdziale API i~implementacji,
ograniczaniu zależności, korzystaniu z~targetów CMake oraz świadomym
wersjonowaniu interfejsu binarnego. Nowoczesne moduły upraszczają część
problemów nagłówków, ale wprowadzają wymagania dotyczące narzędzi,
BMI/CMI i~kolejności kompilacji.

\section*{Literatura i~dokumentacja}
\addcontentsline{toc}{section}{Literatura i dokumentacja}
\begin{itemize}
\item ISO/IEC 14882:2020, zagadnienia \emph{translation units},
      \emph{linkage}, \emph{One Definition Rule} i~\emph{modules}.
\item \url{https://en.cppreference.com/w/cpp/language/definition}
      --- ODR, deklaracje i~definicje.
\item \url{https://en.cppreference.com/w/cpp/language/modules}
      --- składnia i~semantyka modułów.
\item \url{https://gcc.gnu.org/onlinedocs/gcc/C_002b_002b-Modules.html}
      --- implementacja modułów w~GCC.
\item \url{https://clang.llvm.org/docs/StandardCPlusPlusModules.html}
      --- moduły w~Clangu, BMI i~komendy budowania.
\item \url{https://cmake.org/cmake/help/latest/manual/cmake-cxxmodules.7.html}
      --- CMake i~zależności modułów.
\item \url{https://cmake.org/cmake/help/latest/command/target_sources.html}
      --- \code{FILE_SET CXX_MODULES}.
\item \url{https://gcc.gnu.org/onlinedocs/gcc/Code-Gen-Options.html}
      --- widoczność symboli w~ELF.
\item John Lakos, \emph{Large-Scale C++ Software Design}, Addison-Wesley, 1996.
\end{itemize}

% Rozdzial 18: Testowanie, debugowanie i profilowanie w C++20.
% Wylacznie tresc rozdzialu; bez \documentclass i \begin{document}.
% Zgodny ze skryptem: listings, amsmath, amssymb, booktabs, tabularx,
% enumitem, hyperref oraz makra \code, \term, \nota.
% Przyklady terminalowe: zasadniczo Linux/macOS z GCC/Clang;
% roznice na Windows/MSVC zaznaczono osobno.
% Uzycie: \input{rozdzial_18_testowanie_debugowanie_profilowanie_cpp.tex}

\chapter{Testowanie, debugowanie i profilowanie C++}
\label{chap:cpp18-test-debug-profile}

\section{Trzy różne pytania o program}

Po napisaniu i skompilowaniu programu nie wiemy jeszcze, czy jest on
poprawny, bezpieczny i wydajny. Każdą z tych własności sprawdza się
innymi metodami:
\begin{description}[style=nextline]
\item[Testowanie (\emph{testing})]
Czy program daje oczekiwane wyniki na wybranych danych oraz czy
spełnia wymagane kontrakty? Testy wykrywają błędy, lecz ich pozytywny
wynik nie jest ogólnym dowodem poprawności.
\item[Debugowanie (\emph{debugging})]
Gdzie i dlaczego program odbiega od oczekiwanego zachowania?
Debugger pozwala obserwować stos wywołań, zmienne, wątki i kolejne
instrukcje.
\item[Profilowanie (\emph{profiling})]
Gdzie program zużywa czas procesora, pamięć i inne zasoby?
Profilowanie wskazuje rzeczywiste wąskie gardła; samo zgadywanie
na podstawie długości kodu często prowadzi do błędnych optymalizacji.
\end{description}

Istotne jest także rozróżnienie \emph{testu poprawności} od
\emph{benchmarku}: pierwszy sprawdza warunek logiczny, drugi mierzy
wielkość fizyczną (np. czas, przepustowość, alokacje). Szybki, lecz
błędny algorytm nie jest dobrym rozwiązaniem.

\nota{W~tym rozdziale podstawowym standardem jest C++20. GoogleTest,
Google Benchmark, GDB, LLDB, Valgrind i~narzędzia \code{perf} nie są
częścią standardowej biblioteki C++; wymagają odpowiednich instalacji.
Polecenia zależą od systemu operacyjnego i~kompilatora.}

\section{Strategia testowania i testowalny projekt}
\label{sec:cpp18-strategy}

\subsection{Poziomy testów}

\begin{center}
\small
\begin{tabularx}{\textwidth}{@{}p{0.29\textwidth}X@{}}
\toprule
Rodzaj testu & Co sprawdza i kiedy jest przydatny?\\
\midrule
Jednostkowy (\emph{unit}) & Mała, izolowana funkcja lub klasa; szybka informacja zwrotna.\\
Integracyjny & Współdziałanie komponentów, np. parsera i~bazy danych.\\
Systemowy / end-to-end & Cały scenariusz użytkownika, często z~rzeczywistym wejściem/wyjściem.\\
Regresyjny & Przypadek, który wcześniej ujawniał błąd; nie powinien powrócić.\\
Własnościowy / losowy & Niezmienniki dla wielu zestawów danych, nie tylko pojedyncze wyniki.\\
Wydajnościowy & Czas, przepustowość, opóźnienie, pamięć; wymaga kontrolowanego środowiska.\\
\bottomrule
\end{tabularx}
\end{center}

Dobry test powinien być \emph{powtarzalny}, \emph{niezależny} od innych
testów i~możliwie mały. W~praktyce wydzielamy obliczenia od efektów
ubocznych (sieci, zegara, plików, losowości), dzięki czemu zależności
możemy zastąpić kontrolowanymi implementacjami.

\subsection{Kontrakty, przypadki brzegowe i błędy numeryczne}

Rozważmy funkcję obliczającą średnią. Jej kontrakt powinien ustalać
zachowanie dla pustego zakresu; poniżej zgłaszamy wyjątek.

\begin{lstlisting}[caption={Prosta funkcja z jawnym kontraktem dla pustego zakresu (C++20).}]
#include <numeric>
#include <span>
#include <stdexcept>

double average(std::span<const double> xs) {
    if (xs.empty())
        throw std::invalid_argument("empty range");
    return std::accumulate(xs.begin(), xs.end(), 0.0)
         / static_cast<double>(xs.size());
}
\end{lstlisting}

Testujemy nie tylko wartości typowe, lecz także rozmiary $0$ i~$1$,
wartości ujemne, powtórzenia, skrajne zakresy oraz dane naruszające
warunki wejściowe. Dla liczb zmiennoprzecinkowych zwykła równość
\code{a == b} nie zawsze jest właściwa: wynik może zależeć od porządku
dodawania i~optymalizacji. Stosujemy tolerancję dostosowaną do skali
i~wymagań zadania, a~nie arbitralne \code{1e-6} dla każdego algorytmu.

\section{Testy bez zewnętrznego frameworka}

\subsection{\texttt{assert} i \texttt{static\_assert}}

\code{assert(expr)} z~nagłówka \code{<cassert>} sprawdza warunek
w~czasie działania, lecz może zostać usunięte przy zdefiniowanym
\code{NDEBUG}. \code{static_assert} sprawdza warunek w~czasie
kompilacji. W~szczególności zwykłe \code{assert} nie nadaje się do
walidacji niepoprawnych danych wejściowych w~kodzie produkcyjnym.

\begin{lstlisting}[caption={Asercja kompilacyjna i kontrola zachowania funkcji.}]
#include <cassert>
#include <concepts>
#include <type_traits>

constexpr int square(int x) { return x * x; }
static_assert(square(6) == 36);
static_assert(std::integral<int>);
static_assert(!std::is_copy_constructible_v<void>);

int main() {
    assert(square(7) == 49);
}
\end{lstlisting}

\nota{Przy \code{-DNDEBUG} argument \code{assert} może w~ogóle nie być
wykonywany. Nie wolno umieszczać w~nim niezbędnych efektów ubocznych,
np.~\code{assert(++counter == 1)}. \code{static_assert} również nie
zastępuje testów dla danych znanych dopiero podczas wykonania.}

\subsection{Niezależne testy danych wejściowych}

Poniższą funkcję \code{main} umieść pod definicją \code{average}
z~poprzedniego listingu. Test nie wymaga zewnętrznej biblioteki.
Sprawdza wynik, wyjątek i~przypadek jednoelementowy. Kompiluj go
bez \code{-DNDEBUG}, aby wyrażenia \code{assert} były wykonywane.

\begin{lstlisting}[caption={Test funkcji average bez frameworka; dołącz pod jej definicją (C++20).}]
#include <cassert>
#include <cmath>
#include <stdexcept>
#include <vector>

// Zakladamy definicje average z poprzedniego przykladu.
int main() {
    const std::vector<double> xs{1.0, 2.0, 3.0, 4.0};
    assert(std::abs(average(xs) - 2.5) < 1e-12);
    const std::vector<double> one{7.0};
    assert(average(one) == 7.0);

    bool caught = false;
    const std::vector<double> empty;
    try {
        average(empty);
    } catch (const std::invalid_argument&) {
        caught = true;
    }
    assert(caught);
}
\end{lstlisting}

W~realnym projekcie lepiej użyć frameworka, który rejestruje testy,
raportuje niepowodzenia i~pozwala uruchamiać wybrane przypadki.

\section{GoogleTest: testy jednostkowe w praktyce}
\label{sec:cpp18-gtest}

Biblioteka \emph{GoogleTest} (często \code{gtest}) udostępnia makra
\code{TEST}, \code{TEST_F}, \code{EXPECT_*} i~\code{ASSERT_*}.
Assercje \code{EXPECT_*} zapisują niepowodzenie i~pozwalają wykonać
resztę funkcji testowej, podczas gdy \code{ASSERT_*} przerywają
bieżącą funkcję testową. To nie są mechanizmy obsługi wyjątków C++.

\begin{lstlisting}[caption={GoogleTest: wyniki, wyjatki i przyblizona rownosc.}]
// Plik: average_test.cpp
#include <gtest/gtest.h>
#include <stdexcept>
#include <vector>

#include "average.hpp"
// Funkcja average jest zdefiniowana w pliku average.cpp.
TEST(AverageTest, ComputesMean) {
    const std::vector<double> xs{2.0, 4.0, 6.0};
    EXPECT_NEAR(average(xs), 4.0, 1e-12);
}

TEST(AverageTest, RejectsEmptyInput) {
    const std::vector<double> empty;
    EXPECT_THROW(average(empty), std::invalid_argument);
}

TEST(AverageTest, HandlesSingleton) {
    const std::vector<double> xs{7.0};
    ASSERT_EQ(xs.size(), 1u);
    EXPECT_DOUBLE_EQ(average(xs), 7.0);
}
\end{lstlisting}

\code{EXPECT_NEAR(a,b,tol)} sprawdza bezwzględną różnicę, dlatego
niekiedy lepiej użyć własnego testu tolerancji względnej. Makro
\code{EXPECT_DOUBLE_EQ} porównuje liczby z~tolerancją opartą na ULP;
nie należy utożsamiać go z~matematyczną równością liczb rzeczywistych.

\subsection{Fixture i testy parametryzowane}

\emph{Fixture} przygotowuje stan wspólny dla testów. Każdy test
\code{TEST_F} otrzymuje osobny egzemplarz fixture, dzięki czemu zmiany
jednego testu nie powinny wpływać na inny.

\begin{lstlisting}[caption={Fixture z przygotowanym wektorem.}]
#include <gtest/gtest.h>
#include <algorithm>
#include <vector>

class VectorTest : public ::testing::Test {
protected:
    std::vector<int> values{3, 1, 4, 1, 5};
};

TEST_F(VectorTest, SortsAscending) {
    std::sort(values.begin(), values.end());
    EXPECT_TRUE(std::is_sorted(values.begin(), values.end()));
}

TEST_F(VectorTest, StartsFromOriginalState) {
    EXPECT_EQ(values.front(), 3);
    EXPECT_EQ(values.size(), 5u);
}
\end{lstlisting}

Dla wielu par wejście--wynik przydatne są testy parametryzowane:

\begin{lstlisting}[caption={GoogleTest: test parametryzowany z zestawem wartosci.}]
#include <gtest/gtest.h>
#include <tuple>

constexpr int square(int x) { return x * x; }

class SquareTest
    : public ::testing::TestWithParam<std::tuple<int, int>> {};

TEST_P(SquareTest, MatchesExpected) {
    auto [x, expected] = GetParam();
    EXPECT_EQ(square(x), expected);
}

INSTANTIATE_TEST_SUITE_P(
    Values, SquareTest,
    ::testing::Values(
        std::tuple{0, 0},
        std::tuple{2, 4},
        std::tuple{-3, 9}));
\end{lstlisting}

\subsection{Mocki i testowanie współpracy komponentów}

\emph{GoogleMock} pozwala tworzyć kontrolowane implementacje
interfejsów i~sprawdzać wywołania między obiektami. Używaj go tam,
gdzie zachowanie naprawdę zależy od interakcji; nadmierne mockowanie
utrwala przypadkowe szczegóły implementacji.

\begin{lstlisting}[caption={GoogleMock: sprawdzenie pojedynczego wywolania.}]
#include <gmock/gmock.h>
#include <string_view>

struct Sender {
    virtual ~Sender() = default;
    virtual bool send(std::string_view) = 0;
};

struct MockSender : Sender {
    MOCK_METHOD(bool, send, (std::string_view), (override));
};

void notify(Sender& s) { (void)s.send("ready"); }

TEST(NotificationTest, SendsOneMessage) {
    MockSender sender;
    EXPECT_CALL(sender, send("ready"))
        .Times(1).WillOnce(::testing::Return(true));
    notify(sender);
}
\end{lstlisting}

Własne wyjątki, klasy niewspierające kopiowania i~asynchroniczne
wywołania wymagają odpowiedniego projektu fixture/mocków.
Dla operacji wykonywanych w~innym wątku szczególnie istotne jest
kontrolowane oczekiwanie na koniec pracy przed zniszczeniem mocka.

\subsection{CMake i CTest: mały projekt wieloplikowy}
\label{sec:cpp18-cmake-test}

Poniższa organizacja zakłada, że \code{average} ma deklarację w~pliku
\code{average.hpp}, definicję w~\code{average.cpp}, a~testy w~pliku
\code{average_test.cpp}. W~przykładzie \code{average_test.cpp}
należy dodać \verb|#include "average.hpp"| zamiast deklarować
funkcję ręcznie. Plik nagłówkowy może wyglądać następująco:

\begin{lstlisting}[caption={Plik average.hpp: deklaracja API testowanej funkcji.}]
#pragma once
#include <span>
double average(std::span<const double> xs);
\end{lstlisting}

Plik \code{average.cpp} zawiera kod pierwszego przykładu
z~tej części rozdziału (z~dołączonym \code{average.hpp});
w~pliku testów umieszczamy testy z~sekcji~\ref{sec:cpp18-gtest}.

\begin{lstlisting}[language=bash,caption={Struktura projektu z GoogleTest.}]
project/
  CMakeLists.txt
  average.hpp
  average.cpp
  average_test.cpp
\end{lstlisting}

\begin{lstlisting}[language=bash,caption={CMakeLists.txt: biblioteka, testy i automatyczna rejestracja w CTest.}]
cmake_minimum_required(VERSION 3.20)
project(cpp18_tests LANGUAGES CXX)

add_library(core average.cpp)
target_compile_features(core PUBLIC cxx_std_20)
target_include_directories(core PUBLIC ${CMAKE_CURRENT_SOURCE_DIR})

# GoogleTest musi byc zainstalowany lub dostarczony przez menedzer pakietow.
find_package(GTest CONFIG REQUIRED)

add_executable(core_tests average_test.cpp)
target_link_libraries(core_tests PRIVATE core GTest::gtest_main)

enable_testing()
include(GoogleTest)
gtest_discover_tests(core_tests)
\end{lstlisting}

\begin{lstlisting}[language=bash,caption={Budowanie i uruchamianie testow przez CMake/CTest.}]
cmake -S . -B build -DCMAKE_BUILD_TYPE=Debug
cmake --build build -j
ctest --test-dir build --output-on-failure
./build/core_tests --gtest_filter='AverageTest.*'
\end{lstlisting}

\nota{W~systemach z~generatorem wielokonfiguracyjnym (np.~Visual
Studio) wariant budowania wybiera się przez \code{--config Debug}.
Przykład wymaga zainstalowanego GoogleTest; alternatywą jest
\code{FetchContent} z~\emph{przypiętą} wersją zależności, a~nie
pobieranie nieokreślonej wersji podczas każdego budowania.}

\section{Testy własnościowe, generatywne i fuzzing}
\label{sec:cpp18-property}

Sprawdzanie setek wartości jest bardziej użyteczne, jeśli testuje
\emph{własność}, a~nie jedynie listę ręcznie wpisanych wyników.
Na przykład po posortowaniu ciągu musi on być niemalejący i~mieć
tę samą multizbiorowość wartości co przed sortowaniem.

\begin{lstlisting}[caption={Powtarzalny test wlasnosciowy przy uzyciu generatora losowego.}]
#include <algorithm>
#include <cassert>
#include <random>
#include <vector>

int main() {
    std::mt19937 gen(123456u);  // Ustalony seed.
    std::uniform_int_distribution<int> count(0, 300);
    std::uniform_int_distribution<int> value(-1000, 1000);

    for (int trial = 0; trial < 1000; ++trial) {
        std::vector<int> x(count(gen));
        for (int& v : x) v = value(gen);
        auto original = x;
        std::sort(x.begin(), x.end());
        assert(std::is_sorted(x.begin(), x.end()));
        std::sort(original.begin(), original.end());
        assert(x == original);
    }
}
\end{lstlisting}

W~testach losowych przechowuj \emph{seed} i~minimalny przykład,
który ujawnił błąd. Warto również stosować \emph{shrinking}:
automatyczne zmniejszanie wejścia do możliwie prostego kontrprzykładu.

\emph{Fuzzing} automatycznie generuje wejścia, kierując się m.in.
pokryciem kodu. W~Clangu narzędzie \emph{libFuzzer} może wywoływać
poniższą funkcję z~wieloma binarnymi danymi wejściowymi.

\begin{lstlisting}[caption={Minimalny punkt wejsciowy dla libFuzzer (Clang).}]
#include <cstddef>
#include <cstdint>
#include <vector>

extern "C" int LLVMFuzzerTestOneInput(
    const std::uint8_t* data, std::size_t size) {
    // Podmien na testowany parser lub dekoder.
    std::vector<std::uint8_t> bytes(data, data + size);
    if (bytes.size() > 2 && bytes[0] == 0x42) {
        // Przypadek zalezy od danych; nie wolno zakladac ich poprawnosci.
        (void)bytes[1];
    }
    return 0;
}
\end{lstlisting}

\begin{lstlisting}[language=bash,caption={Budowanie harnessu fuzzingu w Clangu.}]
clang++ -std=c++20 -O1 -g \
  -fsanitize=fuzzer,address,undefined fuzz.cpp -o fuzz
./fuzz -max_total_time=30
\end{lstlisting}

Po zgłoszeniu awarii zapisujemy minimalny przypadek jako test
regresyjny. Fuzzer nie dowodzi braku błędów, a~jego skuteczność zależy
od jakości punktów wejściowych i~stanu pokrycia.

\section{Debugowanie: od komunikatu do przyczyny}
\label{sec:cpp18-debug}

\subsection{Informacje debugowe i rodzaje kompilacji}

Użyteczne opcje GCC/Clang to \code{-g} (symbole debugowe),
\code{-O0} (łatwe śledzenie krokowe) oraz \code{-Og} (optymalizacje
starające się zachować dobrą obserwowalność w~debuggerze). Program
kompilowany z~\code{-O2} może nie mieć osobnej zmiennej ani instrukcji
odpowiadającej każdej linii kodu źródłowego.

\begin{lstlisting}[language=bash,caption={Budowanie programu do debugowania.}]
g++ -std=c++20 -g -O0 -Wall -Wextra -Wpedantic \
    app.cpp -o app_debug

gdb ./app_debug
\end{lstlisting}

Nie wszystkie błędy są odtwarzalne przy \code{-O0}; niektóre pojawiają
się dopiero przy \code{-O2} lub \code{-O3}. Do ich analizy warto
pozostawić \code{-g} także w~wersji zoptymalizowanej i~zachować
niestripowane symbole debugowe.

\subsection{GDB: breakpoint, stos i obserwacja zmiennych}

\begin{lstlisting}[caption={Program do cwiczenia debugowania.}]
#include <stdexcept>
#include <vector>

int element_at(const std::vector<int>& v, std::size_t index) {
    return v.at(index);  // Moze zglosic std::out_of_range.
}

int main() {
    const std::vector<int> v{3, 5, 8};
    return element_at(v, 5);
}
\end{lstlisting}

\begin{lstlisting}[language=bash,caption={Typowa sesja w GDB.}]
(gdb) break element_at        # Punkt zatrzymania.
(gdb) run                     # Uruchom program.
(gdb) info args               # Argumenty funkcji.
(gdb) print index             # Wartosc zmiennej.
(gdb) backtrace               # Stos wywolan.
(gdb) next                    # Nastepna linia bez wchodzenia.
(gdb) step                    # Wejdz do wywolywanej funkcji.
(gdb) catch throw             # Zatrzymaj przy zglaszaniu wyjatku.
(gdb) continue                # Kontynuuj.
\end{lstlisting}

\code{break}, \code{watch}, \code{next}, \code{step},
\code{finish}, \code{frame}, \code{print} i~\code{bt} to podstawowe
narzędzia. \emph{Watchpoint} śledzi zmianę pamięci, ale możliwość
sprzętowego śledzenia zależy od architektury i~debuggera.

\subsection{Wątki i deadlock w debuggerze}

W~aplikacji wielowątkowej sama analiza wątku \code{main} nie
wystarczy. W~GDB:

\begin{lstlisting}[language=bash,caption={Diagnostyka wielowatkowosci w GDB.}]
(gdb) info threads
(gdb) thread 3
(gdb) bt
(gdb) thread apply all bt
(gdb) set print pretty on
\end{lstlisting}

Jeśli wątki zatrzymały się w~\code{mutex::lock()} lub
\code{condition_variable::wait()}, trzeba odtworzyć graf
zależności: który wątek posiada dany mutex i~na co czeka.
Program zakleszczony nie musi zawierać wyścigu danych, więc
ThreadSanitizer nie zastąpi analizy zakleszczeń.

\subsection{LLDB: odpowiednik na macOS i innych platformach}

LLDB oferuje analogiczne operacje:

\begin{lstlisting}[language=bash,caption={Podstawowe polecenia LLDB.}]
(lldb) breakpoint set --name element_at
(lldb) run
(lldb) frame variable
(lldb) thread backtrace
(lldb) thread list
(lldb) thread backtrace all
(lldb) next
(lldb) step
(lldb) continue
\end{lstlisting}

W~VS Code można podłączyć GDB/LLDB poprzez odpowiednie rozszerzenie,
a~następnie korzystać z~breakpointów, podglądu zmiennych i~stosów.
Dokładny adapter oraz konfiguracja \code{launch.json} zależą od
systemu, rozszerzenia i~narzędzi zainstalowanych lokalnie.

\section{Sanitizatory i narzędzia wykrywania błędów}
\label{sec:cpp18-sanitizers}

Sanitizatory dodają do programu kod kontrolujący określone klasy
błędów w~czasie wykonywania. W~odróżnieniu od debuggera mogą
wykryć nieprawidłową operację \emph{w~chwili jej wykonania}, zanim
zauważymy błędny rezultat.

\begin{center}
\small
\begin{tabularx}{\textwidth}{@{}p{0.28\textwidth}X@{}}
\toprule
Narzędzie & Główne zastosowanie\\
\midrule
AddressSanitizer (ASan) & Wyjście poza granice, use-after-free, część przypadków wycieków.\\
UndefinedBehaviorSanitizer (UBSan) & Np. przepełnienie \code{signed int}, błędne przesunięcia, null dereference, alignment.\\
ThreadSanitizer (TSan) & Wiele wyścigów danych i~część problemów synchronizacji.\\
MemorySanitizer (MSan) & Odczyty niezainicjalizowanej pamięci; wymaga specjalnych warunków instrumentacji.\\
Valgrind Memcheck & Wykrywanie wybranych błędów pamięci i~wycieków (na wspieranych systemach).\\
\bottomrule
\end{tabularx}
\end{center}

\subsection{AddressSanitizer i UndefinedBehaviorSanitizer}

Poniższy program jest \emph{celowo niepoprawny}; dostęp do elementu
po zwolnieniu pamięci jest niezdefiniowanym zachowaniem. Nie należy
interpretować jego wyniku bez instrumentacji.

\begin{lstlisting}[caption={Celowo bledny program: use-after-free.}]
int main() {
    int* p = new int(42);
    delete p;
    return *p;  // UB: odczyt z nieistniejacego obiektu.
}
\end{lstlisting}

\begin{lstlisting}[language=bash,caption={Instrumentacja ASan i UBSan; GCC/Clang.}]
g++ -std=c++20 -O1 -g -fno-omit-frame-pointer \
    -fsanitize=address,undefined -fno-sanitize-recover=all \
    bad.cpp -o bad_san
./bad_san
\end{lstlisting}

ASan i~UBSan można często połączyć. ASan/TSan/MSan nie są natomiast
zwykle łączone ze sobą w~jednym procesie: używamy osobnych
konfiguracji i~uruchomień. Raporty sanitizatorów zależą od wersji
kompilatora, wsparcia runtime i~platformy.

\subsection{ThreadSanitizer i wyścig danych}

Przykładowy wyścig występuje przy równoczesnym modyfikowaniu
zwykłego licznika. To również \emph{celowo błędny} kod:

\begin{lstlisting}[caption={Celowo bledne wspoldzielenie licznika.}]
#include <thread>

int main() {
    int x = 0;
    std::thread t1([&] { ++x; });
    std::thread t2([&] { ++x; });
    t1.join();
    t2.join();
}
\end{lstlisting}

\begin{lstlisting}[language=bash,caption={Oddzielna konfiguracja z ThreadSanitizer.}]
g++ -std=c++20 -O1 -g -pthread \
    -fsanitize=thread race.cpp -o race_tsan
./race_tsan
\end{lstlisting}

\nota{Brak zgłoszenia w~danym uruchomieniu \emph{nie} dowodzi braku
wyścigów danych. TSan ma ograniczenia platformowe, może zgłaszać
wyniki fałszywie dodatnie dla niestandardowej synchronizacji i~nie
gwarantuje wykrywania deadlocków.}

\subsection{Valgrind i testy wycieków}

Na wspieranych systemach Linux przydatna jest kontrola programów
bez kompilacyjnej instrumentacji:

\begin{lstlisting}[language=bash,caption={Przykladowe badanie pamieci przez Valgrind Memcheck.}]
valgrind --tool=memcheck --leak-check=full \
         --show-leak-kinds=all ./app_debug
\end{lstlisting}

Różne narzędzia mają inne modele wykrywania i~inne koszty
wykonania. Wybór zależy od konkretnej klasy błędów; wyniki dobrze
weryfikować na małym, odtwarzalnym przypadku testowym.

\section{Analiza statyczna, ostrzeżenia i pokrycie kodu}
\label{sec:cpp18-static}

Kompilator potrafi zgłosić wiele usterek jeszcze przed wykonaniem
programu. Warto utrzymywać ostrzeżenia w~CI i~unikać wyciszania ich
bez zrozumienia przyczyny.

\begin{lstlisting}[language=bash,caption={Przykladowe ostrzezenia GCC/Clang.}]
g++ -std=c++20 -Wall -Wextra -Wpedantic -Wconversion \
    -Wshadow -Werror -c core.cpp
\end{lstlisting}

Opcje zależą od kompilatora; \code{-Werror} bywa pożyteczne przy
stabilnej bazie kodu, lecz aktualizacja kompilatora może wprowadzić
nowe diagnostyki. \emph{clang-tidy} wykonuje zestaw analiz
statycznych, np. dotyczących nowoczesnych idiomów C++, błędów
potencjalnych i~stylu.

\begin{lstlisting}[language=bash,caption={Analiza projektu ze zbiorem komend kompilacji.}]
cmake -S . -B build -DCMAKE_EXPORT_COMPILE_COMMANDS=ON
clang-tidy -p build src/core.cpp
\end{lstlisting}

\subsection{Pokrycie testami: co liczby naprawdę oznaczają?}

\emph{Line coverage} mierzy wykonane linie, \emph{branch coverage}
wykonane gałęzie, a~\emph{function coverage} wywołane funkcje.
Wysoki procent nie oznacza sprawdzenia wszystkich kombinacji wejść
ani dowodu poprawności. Na przykład test może wykonać gałąź, nie
sprawdzając rezultatu.

\begin{lstlisting}[language=bash,caption={Pokrycie testow z GCC i gcov (przykladowy projekt).}]
g++ -std=c++20 --coverage -O0 -g core.cpp tests.cpp -o tests
./tests
gcov core.cpp
\end{lstlisting}

Dla Clanga dostępny jest osobny zestaw instrumentacji LLVM:
\begin{lstlisting}[language=bash,caption={Clang: instrumentacja i raport pokrycia.}]
clang++ -std=c++20 -O0 -g -fprofile-instr-generate \
    -fcoverage-mapping core.cpp tests.cpp -o tests
LLVM_PROFILE_FILE=tests.profraw ./tests
llvm-profdata merge -sparse tests.profraw -o tests.profdata
llvm-cov report ./tests -instr-profile=tests.profdata
\end{lstlisting}

W~szczególności należy testować \emph{warunki brzegowe}, a~nie
maksymalizować jedynie procent pokrycia.

\section{Benchmarking: jak mierzyć szybkość bez samooszukiwania}
\label{sec:cpp18-benchmark}

\subsection{Pułapki mikrobenchmarków}

Nowoczesny kompilator usuwa obliczenia o~nieobserwowalnych efektach,
wykonuje \emph{constant folding} i~wektoryzuje pętle. Dlatego mierzenie
pustej pętli lub stałego obliczenia może nie mieć związku z~realną
wydajnością. Wynik zależy ponadto od rozgrzania cache, alokacji,
obciążenia systemu, zegarów CPU, temperatury i~danych testowych.

\begin{lstlisting}[caption={Pomiar algorytmu z rzeczywistym wynikiem (bez biblioteki benchmark).}]
#include <chrono>
#include <iostream>
#include <numeric>
#include <vector>

int main() {
    std::vector<int> data(1'000'000, 1);
    using clock = std::chrono::steady_clock;
    long long checksum = 0;

    for (int repetition = 0; repetition < 10; ++repetition) {
        auto start = clock::now();
        const auto result = std::accumulate(
            data.begin(), data.end(), 0LL);
        auto finish = clock::now();
        checksum += result;  // Wynik jest wykorzystany po pomiarze.
        auto us = std::chrono::duration_cast<
            std::chrono::microseconds>(finish - start).count();
        std::cout << "Trial " << repetition << ": " << us << " us\n";
    }
    std::cout << "checksum=" << checksum << '\n';
}
\end{lstlisting}

To ilustracja, nie wzorcowy mikrobenchmark: optymalizator może
ułatwić sobie obliczenia na regularnych danych; dodatkowo
\code{std::cout} poza odcinkiem pomiarowym zmienia warunki
między próbami. Użyteczne benchmarki wymagają wielu powtórzeń,
przygotowanych danych i~analizy rozkładu czasów.

\subsection{Google Benchmark}

\emph{Google Benchmark} automatyzuje wielokrotne uruchamianie i~raportuje
miary czasowe. W~krytycznym fragmencie wykorzystujemy
\code{benchmark::DoNotOptimize}, aby wynik obliczenia nie został
trywialnie usunięty; nie wyłącza to jednak wszystkich optymalizacji
na badanym wyrażeniu.

\begin{lstlisting}[caption={Samodzielny benchmark sumowania wielu elementow.}]
// Plik: benchmark_sum.cpp
#include <benchmark/benchmark.h>
#include <numeric>
#include <vector>

static void BM_Accumulate(benchmark::State& state) {
    const std::vector<int> data(state.range(0), 1);
    for (auto _ : state) {
        const auto sum = std::accumulate(
            data.begin(), data.end(), 0LL);
        benchmark::DoNotOptimize(sum);
    }
    state.SetItemsProcessed(state.iterations() * state.range(0));
}

BENCHMARK(BM_Accumulate)->Arg(1024)->Arg(1 << 20);
BENCHMARK_MAIN();
\end{lstlisting}

\begin{lstlisting}[language={},caption={CMakeLists.txt dla przykładu Google Benchmark.}]
cmake_minimum_required(VERSION 3.20)
project(BenchmarkSum LANGUAGES CXX)
find_package(benchmark CONFIG REQUIRED)
add_executable(bench benchmark_sum.cpp)
target_compile_features(bench PRIVATE cxx_std_20)
target_link_libraries(bench PRIVATE benchmark::benchmark)
\end{lstlisting}

\begin{lstlisting}[language=bash,caption={Budowanie i powtarzanie pomiaru.}]
cmake -S . -B build -DCMAKE_BUILD_TYPE=Release
cmake --build build
./build/bench --benchmark_repetitions=10
./build/bench --benchmark_out=results.json \
              --benchmark_out_format=json
\end{lstlisting}

W~porównaniach równoległych trzeba ponadto ustalić liczbę wątków,
rozmiary bloków, warm-up, powtarzalność danych oraz ewentualne
\emph{false sharing}. Sam najkrótszy pojedynczy czas nie opisuje
stabilnej wydajności.

\section{Profilowanie CPU i pamięci}
\label{sec:cpp18-profiling}

\subsection{Profilowanie próbkowe: Linux \texttt{perf}}

Zamiast mierzyć wyłącznie całkowity czas, profiler pomaga wykrywać
\emph{hot spots}: funkcje zajmujące znaczącą część wykonania.
Przykładowe polecenia systemu Linux:

\begin{lstlisting}[language=bash,caption={Profilowanie CPU przez perf (Linux).}]
g++ -std=c++20 -O2 -g -fno-omit-frame-pointer \
    app.cpp -o app
perf stat ./app
perf record -g ./app
perf report
\end{lstlisting}

\code{perf stat} podaje liczniki zdarzeń udostępniane przez jądro
i~sprzęt, np. cykle i~instrukcje. \code{perf record} zbiera próbki
umożliwiające oglądanie gorących funkcji i~stosu. Dostępne zdarzenia
i~uprawnienia zależą od komputera i~ustawień systemu.

\subsection{Flame graph}

\emph{Flame graph} to agregat stosów wywołań: szerokość prostokąta
odpowiada udziałowi próbek, a~wysokość zagnieżdżeniu wywołań.
Nie jest to dokładna oś czasu. Przykładowy przepływ z~zewnętrznymi
narzędziami FlameGraph:

\begin{lstlisting}[language=bash,caption={Tworzenie flame graph z probek perf (Linux).}]
perf record -F 99 -g -- ./app
perf script | stackcollapse-perf.pl | flamegraph.pl > flame.svg
\end{lstlisting}

Ostatnie dwa skrypty pochodzą z~osobnego projektu FlameGraph i~muszą
być zainstalowane. W~projektach o~dużej liczbie wątków pomagają też
profile opóźnień, I/O i~blokad, a~nie tylko wykorzystania CPU.

\subsection{Alokacje i zużycie pamięci}

Wydajność pamięci to nie tylko liczba bajtów. Ważne są rozkład
alokacji, fragmentacja, lokalność cache i~czas życia obiektów.
Na wspieranych systemach można użyć m.in.:

\begin{lstlisting}[language=bash,caption={Przykladowe pomiary alokacji i pamieci.}]
valgrind --tool=massif ./app
ms_print massif.out.*
/usr/bin/time -v ./app
\end{lstlisting}

\code{massif} służy do profilowania sterty (zgodnie z~możliwościami
narzędzia); \code{/usr/bin/time -v} raportuje miary procesu, np.
\emph{maximum resident set size}, ale nie mówi, które alokacje były
odpowiedzialne. Instrumentacja pamięci również wpływa na czas pracy.

\section{Asembler i związek kodu źródłowego z wydajnością}
\label{sec:cpp18-asm}

Dwa podobne zapisy C++ mogą dać różny kod maszynowy, szczególnie po
włączeniu \code{-O2} lub \code{-O3}. Przy analizie trzeba jednak
uważać, aby nie wyciągać wniosków na podstawie sztucznych przykładów,
w których kompilator może wyeliminować całą pracę.

\begin{lstlisting}[caption={Niewielka funkcja do analizy wygenerowanego kodu.}]
#include <cstddef>

double dot(const double* a, const double* b, std::size_t n) {
    double sum = 0.0;
    for (std::size_t i = 0; i < n; ++i)
        sum += a[i] * b[i];
    return sum;
}
\end{lstlisting}

\begin{lstlisting}[language=bash,caption={Asembler i raport wektoryzacji w GCC (Linux).}]
g++ -std=c++20 -O3 -S -masm=intel dot.cpp -o dot.s

g++ -std=c++20 -O3 -fopt-info-vec-optimized \
    -c dot.cpp -o dot.o
objdump -drC dot.o
\end{lstlisting}

Dla wektoryzacji ważne są: zależności między iteracjami, wyrównanie,
\emph{aliasing} oraz semantyka arytmetyki zmiennoprzecinkowej.
Brak instrukcji SIMD nie musi oznaczać błędu kompilatora; dla
dokładnie zachowanego porządku sumowania kompilator może nie móc
zmienić kolejności operacji. Flaga \code{-ffast-math} pozwala na
istotnie inne przekształcenia numeryczne i~nie powinna być
włączana automatycznie w~programach wymagających ścisłej semantyki.

\subsection{Profile-guided optimization (PGO)}

Kompilator może wykorzystać profile zebrane na reprezentatywnych
danych, aby lepiej rozłożyć kod i~podjąć decyzje optymalizacyjne.
Schemat dla GCC:

\begin{lstlisting}[language=bash,caption={Dwa przebiegi budowania z PGO (GCC).}]
mkdir -p pgo-data
g++ -std=c++20 -O2 -fprofile-generate=pgo-data \
    app.cpp -o app
./app  # Wykonaj reprezentatywny workload.
g++ -std=c++20 -O2 -fprofile-use=pgo-data \
    app.cpp -o app
\end{lstlisting}

PGO nie zastępuje pomiarów: profil z~innego rozkładu danych może
nawet pogorszyć wyniki. Należy też zachować zgodność źródeł, flag
i~miejsc profilu pomiędzy dwoma etapami budowania.

\section{Testowanie kodu wielowątkowego}
\label{sec:cpp18-concurrency-tests}

W~systemach współbieżnych testy wynikowe muszą badać także
\emph{zakończenie pracy}, własność danych i~synchronizację.
Powtórzenie testu wiele razy zwiększa szansę wykrycia błędu, ale
nigdy nie zastępuje rozumowania o~relacjach \emph{happens-before}
(rozdziały~\ref{chap:concurrency} i~\ref{chap:cpp16-advanced-concurrency}).

\begin{lstlisting}[caption={Poprawny test licznika atomowego z jthread (C++20).}]
#include <atomic>
#include <cassert>
#include <thread>
#include <vector>

int main() {
    std::atomic<int> count{0};
    {
        std::vector<std::jthread> workers;
        for (int i = 0; i < 4; ++i) {
            workers.emplace_back([&] {
                for (int k = 0; k < 10000; ++k)
                    count.fetch_add(1, std::memory_order_relaxed);
            });
        }
    } // Destruktory jthread czekaja na wszystkie watki.
    assert(count.load(std::memory_order_relaxed) == 40000);
}
\end{lstlisting}

Do koordynowania momentu rozpoczęcia testu można wykorzystać
\code{std::latch} lub \code{std::barrier} (C++20). W~testach
\emph{timeout} zapobiega nieskończonemu oczekiwaniu w~infrastrukturze
testowej, ale nie może być uznawany za mechanizm anulowania
wątku. Zwracaj uwagę na wyścigi związane z~destrukcją fixture
oraz na kolejność zamykania wątków i~obiektów synchronizacyjnych.

\section{Praktyczny workflow: od błędu do poprawki}

Prześledźmy zalecaną kolejność działania:
\begin{enumerate}
\item \textbf{Odtwórz błąd}: zanotuj wejście, wersję programu, flagi
      kompilacji, środowisko i~opcjonalnie seed.
\item \textbf{Zminimalizuj przypadek}: mały, powtarzalny przykład jest
      łatwiejszy do analizy niż cały system.
\item \textbf{Włącz diagnostykę}: ostrzeżenia, ASan/UBSan, a~dla
      kodu współbieżnego oddzielne uruchomienie TSan.
\item \textbf{Postaw hipotezę}: zidentyfikuj zmienną lub niezmiennik,
      który naruszono; sprawdź stos i~stan programu w~debuggerze.
\item \textbf{Dodaj test regresyjny}: powinien zawodzić przed poprawką
      i~przechodzić po niej.
\item \textbf{Popraw przyczynę}: nie tylko usuń widoczny symptom.
\item \textbf{Zmierz koszt}: jeżeli poprawka ma znaczenie wydajnościowe,
      porównaj benchmarki i~profile przed oraz po zmianie.
\item \textbf{Uruchom pełen zestaw testów}: szczególnie z~różnymi
      konfiguracjami optymalizacji i~kompilatorami.
\end{enumerate}

\section{Automatyzacja w CI i macierz konfiguracji}

W~projektach utrzymywanych długo testy powinny wykonywać się
automatycznie po każdej istotnej zmianie. Praktyczna macierz obejmuje
zwykle konfigurację debugową z~sanitizatorami oraz konfigurację
optymalizowaną z~testami wynikowymi. TSan uruchamiamy w~osobnej
konfiguracji i~na wspieranym środowisku.

\begin{center}
\small
\begin{tabularx}{\textwidth}{@{}p{0.28\textwidth}X@{}}
\toprule
Konfiguracja & Główne zadanie\\
\midrule
Debug, \code{-O0 -g} & Debugowanie źródłowe i~lokalizacja błędów.\\
Debug/san, \code{-O1 -g}, ASan+UBSan & Dostęp do niepoprawnej pamięci i~wybrane UB.\\
Oddzielny build TSan & Wyścigi danych przy konkurencyjnym wykonaniu.\\
Release, \code{-O2} lub \code{-O3} & Testy produkcyjnej konfiguracji i~benchmarki.\\
Clang i GCC & Różne diagnostyki i~różne decyzje optymalizatora.\\
\bottomrule
\end{tabularx}
\end{center}

\nota{Benchmarki są wrażliwe na obciążenie współdzielonych runnerów
CI. W~zwykłym pipeline CI warto przede wszystkim weryfikować
poprawność i~duże regresje; wiarygodne mikrobenchmarki powinny
wykonywać się na kontrolowanym sprzęcie.}

\section{Najczęstsze błędy i wybór narzędzia}

\begin{center}
\small
\begin{tabularx}{\textwidth}{@{}p{0.40\textwidth}X@{}}
\toprule
Problem & Zalecane podejście\\
\midrule
Funkcja daje złe wyniki & Test jednostkowy, debugger, test regresyjny.\\
Awaria tylko dla niektórych danych & Fuzzing, test własnościowy, ASan/UBSan.\\
Podejrzenie wiszącego wskaźnika & ASan, kontrola czasu życia, debugger.\\
Podejrzenie data race & TSan, analiza synchronizacji i~\emph{happens-before}.\\
Program zawiesza się przy wielu wątkach & Stosy wszystkich wątków, graf blokad, timeout testu.\\
Wysokie zużycie pamięci & Massif, profiler alokacji, pomiar RSS.\\
Program jest wolny & \code{perf}/profiler, później benchmark gorącej funkcji.\\
Spadek szybkości po zmianie & Reprodukowalny benchmark, kontrola warunków, profil porównawczy.\\
Wątpliwości co do wektoryzacji & Raport optymalizacji i~asembler, weryfikacja numeryczna.\\
\bottomrule
\end{tabularx}
\end{center}

\section*{Podsumowanie rozdziału}
\addcontentsline{toc}{section}{Podsumowanie rozdziału}

Testy określają, czy otrzymujemy oczekiwane zachowanie; debugger
pozwala prześledzić przyczyny problemu; sanitizatory wykrywają
określone naruszenia reguł wykonywania; profiler i~benchmarki
pomagają podejmować decyzje dotyczące wydajności. Nie istnieje
jedno narzędzie, które zastępuje wszystkie pozostałe.

Praktyczna sekwencja brzmi: \emph{najpierw poprawność, potem
bezpieczeństwo i~reprodukcyjność, następnie pomiary, wreszcie
optymalizacja}. Każda wykryta usterka powinna prowadzić do testu
regresyjnego, a~każda deklarowana poprawa wydajności --- do
porównywalnego pomiaru. Dla systemów wielowątkowych konieczne jest
ponadto świadome rozumienie synchronizacji, czasu życia obiektów
oraz gwarancji modelu pamięci C++.

\section*{Literatura i dokumentacja}
\addcontentsline{toc}{section}{Literatura i dokumentacja}
\begingroup
\sloppy
\begin{itemize}
\item \url{https://google.github.io/googletest/primer.html} ---
      GoogleTest: asercje, fixture i organizacja testów.
\item \url{https://github.com/google/googletest/blob/main/docs/quickstart-cmake.md}
      --- integracja GoogleTest i CMake.
\item \url{https://cmake.org/cmake/help/latest/module/GoogleTest.html}
      --- \code{gtest_discover_tests} i CTest.
\item \url{https://github.com/google/benchmark/blob/main/docs/user_guide.md}
      --- Google Benchmark i stabilne pomiary.
\item \url{https://clang.llvm.org/docs/AddressSanitizer.html} --- ASan.
\item \url{https://clang.llvm.org/docs/UndefinedBehaviorSanitizer.html}
      --- UBSan.
\item \url{https://clang.llvm.org/docs/ThreadSanitizer.html} --- TSan.
\item \url{https://clang.llvm.org/docs/LibFuzzer.html} --- fuzzing.
\item \url{https://sourceware.org/gdb/current/onlinedocs/gdb.html} --- GDB.
\item \url{https://lldb.llvm.org/use/tutorial.html} --- LLDB.
\item \url{https://valgrind.org/docs/manual/manual.html} --- Memcheck i Massif.
\item \url{https://perf.wiki.kernel.org/} --- profiler \code{perf}.
\item Brendan Gregg, \emph{Systems Performance}, wyd. 2,
      Addison-Wesley, 2020.
\end{itemize}
\endgroup

% Rozdzial 20. Fragment do \input; bez preambuly.
% Wymaga listings, enumitem, tabularx, booktabs, hyperref i makr
% \code, \term, \nota. Bazowo C++20.
\chapter{Projektowanie bibliotek i nowoczesnych interfejsów API}
\label{chap:cpp20-api}

\section{Czym jest dobry interfejs biblioteki?}

Biblioteka to nie tylko poprawny algorytm, lecz przede wszystkim
\emph{umowa} między jej autorem a użytkownikiem. Umowa obejmuje typy,
sygnatury, ograniczenia, własność zasobów, czas życia, obsługę błędów,
bezpieczeństwo wyjątkowe, własności współbieżne i kompatybilność wersji.
Każda decyzja ujawniona w publicznym nagłówku może stać się
kosztownym zobowiązaniem na lata. Celem jest interfejs, w którym
\emph{poprawne użycie jest naturalne, a błędne możliwie trudne}.

\begin{center}\small
\begin{tabularx}{\textwidth}{@{}p{0.30\textwidth}X@{}}
\toprule
Własność API & Pytanie kontrolne \\
\midrule
Kontrakt & Co jest dozwolonym wejściem, a co wynikiem? \\
Własność & Kto posiada obiekt, kto go zwalnia i kiedy? \\
Bezpieczeństwo & Czy błędne wejście może spowodować UB albo wyciek? \\
Wyjątki & Jakie gwarancje pozostają po niepowodzeniu? \\
Współbieżność & Co wolno wywoływać równocześnie i co jest synchronizowane? \\
Wydajność & Jakie kopiowania, alokacje i synchronizacje są narzucane? \\
Kompatybilność & Które zmiany są zgodne źródłowo i binarnie? \\
\bottomrule
\end{tabularx}
\end{center}

\nota{Rozdziały o smart pointerach, concepts, type erasure, wyjątkach,
linkowaniu i ABI dostarczają narzędzi. Tu rozpatrujemy je razem, jako
wybory projektowe podejmowane przy budowaniu stabilnych bibliotek.}

\section{Granica interfejsu: dane, widoki i własność}
\label{sec:cpp20-ownership}

Najpierw wybieramy, czy argument jest \emph{wejściem}, \emph{wyjściem},
\emph{stanem współdzielonym} czy \emph{zasobem przekazywanym na własność}.
Dopiero potem wybieramy konkretny typ C++.

\begin{center}\small
\begin{tabularx}{\textwidth}{@{}p{0.37\textwidth}X@{}}
\toprule
Parametr & Interpretacja w typowym interfejsie \\
\midrule
\code{T} & Wartość własna; przekazanie lub przejęcie kopii/przeniesienia. \\
\texttt{const T\&} & Pożyczony, niezmieniany obiekt, który musi przeżyć wywołanie. \\
\texttt{T\&} & Obowiązkowy, modyfikowalny obiekt pożyczony. \\
\code{T*} & Zwykle opcjonalny, niewłasnościowy dostęp; dopuszcza \code{nullptr}. \\
\code{unique_ptr<T>} & Przeniesienie lub przechowywanie wyłącznej własności. \\
\code{shared_ptr<T>} & Współdzielona własność, nie sam odczyt obiektu. \\
\code{span<T>} & Pożyczony widok ciągłego bufora, bez informacji o własności. \\
\code{string_view} & Pożyczony widok tekstu, możliwy problem czasu życia. \\
\bottomrule
\end{tabularx}
\end{center}

\subsection{Parametry i brak zbędnej alokacji}

\begin{lstlisting}[caption={Interfejs operujacy na dowolnym ciaglym buforze (C++20).}]
#include <cstddef>
#include <span>
#include <stdexcept>

// Nie przechowuje span po powrocie z funkcji.
double weighted_sum(std::span<const double> a,
                    std::span<const double> b) {
    if (a.size() != b.size())
        throw std::invalid_argument("size mismatch");
    double sum = 0.0;
    for (std::size_t i = 0; i < a.size(); ++i)
        sum += a[i] * b[i];
    return sum;
}
\end{lstlisting}

Stosowanie \code{span} pozwala przyjmować
\code{vector}, \code{array} i tablice bez narzucania własności.
Jednocześnie nie wolno zachować widoku dłużej niż trwa jego bufor.

\subsection{Wartość czy referencja?}

Dla małych i tanio kopiowanych typów (np. \code{int}) przekazanie przez
wartość jest często najlepsze. Dla tekstu, który ma być zapisany
wewnątrz obiektu, idiom \code{std::string} przez wartość i
\code{std::move} do pola może uprościć przeciążenia konstruktorów.
Nie jest jednak bezwarunkowo najszybszy dla każdego rozkładu wejść:
dla samych \emph{lvalue} może wykonać dodatkowe kopiowanie.

\begin{lstlisting}[caption={Jeden konstruktor obslugujacy lvalue i rvalue (C++11).}]
#include <string>
#include <utility>

class Person {
    std::string name_;
public:
    explicit Person(std::string name)
        : name_(std::move(name)) {}
    const std::string& name() const noexcept { return name_; }
};
\end{lstlisting}

\section{Wyrażanie kontraktów w systemie typów}
\label{sec:cpp20-contracts}

Silny typ eliminuje klasę błędów polegających na pomyleniu wartości o
tej samej reprezentacji. \code{explicit} blokuje niepożądaną konwersję
niejawną. \code{[[nodiscard]]} pomaga wykryć ignorowanie wyniku;
nie gwarantuje jednak obowiązkowego ostrzeżenia we wszystkich trybach.

\begin{lstlisting}[caption={Strong types rozrozniajace dwa identyfikatory (C++20).}]
#include <cstdint>

struct UserId {
    std::uint64_t value;
    explicit constexpr UserId(std::uint64_t v) : value(v) {}
    friend constexpr bool operator==(UserId, UserId) = default;
};

struct OrderId {
    std::uint64_t value;
    explicit constexpr OrderId(std::uint64_t v) : value(v) {}
    friend constexpr bool operator==(OrderId, OrderId) = default;
};

[[nodiscard]] constexpr bool valid(UserId id) noexcept {
    return id.value != 0;
}

static_assert(valid(UserId{1}));
// valid(OrderId{1}); // blad kompilacji
\end{lstlisting}

\subsection{Concepts: wymóg składniowy nie zawsze jest wymogiem semantycznym}

\code{requires} pozwala sprawdzić, czy wyrażenie ma poprawną składnię,
ale nie dowodzi prawa matematycznego, np. przechodniości porządku,
asocjatywności czy stabilności haszowania. Dokumentacja musi dopowiadać
wymagania semantyczne. Przykładowo komparator dla sortowania musi
realizować \emph{strict weak ordering}; przypadkowy predykat typu
\code{bool(const T&,const T&)} nie wystarcza.

\begin{lstlisting}[caption={Ograniczenie generycznego sumowania (C++20).}]
#include <concepts>
#include <ranges>
#include <type_traits>

template<class R>
concept SummableRange = std::ranges::input_range<R> &&
    std::default_initializable<std::ranges::range_value_t<R>> &&
    requires(std::ranges::range_value_t<R>& result,
             std::ranges::range_reference_t<R> value) {
        result = result + value;
    };

template<SummableRange R>
auto sum_values(R&& range) {
    using T = std::ranges::range_value_t<R>;
    T result{};
    for (auto&& value : range) result = result + value;
    return result;
}
\end{lstlisting}

Wymóg \code{default_initializable} odpowiada instrukcji \code{T result\{\}}.
Wyrażenie w \code{requires} sprawdza dokładnie przypisanie wykonywane
w pętli, także gdy element zakresu jest typem pośrednim, a nie \code{T&}.
To nadal nie dowodzi sensu matematycznego sumowania. Jeśli wartość
domyślna nie jest elementem neutralnym dla dodawania, funkcja powinna
przyjmować wartość początkową od wywołującego.

\section{Obsługa błędów: wyjątki, typ wyniku i status}

\begin{center}\small
\begin{tabularx}{\textwidth}{@{}p{0.33\textwidth}X@{}}
\toprule
Kontrakt API & Kiedy i jak go użyć \\
\midrule
Wyjątek & Przerywa zwykły tok działania i przenosi błąd do obsługi; zasoby powinny być chronione przez RAII. \\
\code{std::optional<T>} & Oznacza obecność lub brak wyniku, gdy przyczyna braku nie jest potrzebna. \\
\code{std::expected<T,E>} & Zwraca wynik albo jawny błąd typu \code{E}; standard od C++23. \\
\code{std::error_code} & Przenosi kod błędu bez wyjątku, często obok wariantu zgłaszającego wyjątki. \\
\code{noexcept} & Jest częścią kontraktu: wyjątek opuszczający taką funkcję kończy program. \\
\bottomrule
\end{tabularx}
\end{center}

W C++20 podstawowymi wyborami są wyjątki oraz jawne typy wyniku,
np. \code{optional<T>} dla naturalnego braku wartości. Od C++23
\code{expected<T,E>} reprezentuje wartość lub opis błędu. Żadna
strategia nie jest uniwersalna: trzeba ustalić politykę dla całej
warstwy biblioteki, a nie mieszać bez powodu sposoby sygnalizowania.

\begin{lstlisting}[caption={Jawny blad jako typ wyniku w C++20.}]
#include <optional>
#include <string_view>

std::optional<unsigned> parse_digit(std::string_view x) {
    if (x.size() != 1 || x[0] < '0' || x[0] > '9')
        return std::nullopt;
    return static_cast<unsigned>(x[0] - '0');
}
\end{lstlisting}

\begin{lstlisting}[caption={Rozbudowany blad jako expected (C++23; fragment).}]
#include <cstddef>
#include <expected>
#include <string>
#include <string_view>

std::expected<std::size_t, std::string>
nonempty_length(std::string_view text) {
    if (text.empty())
        return std::unexpected<std::string>{"empty text"};
    return text.size();
}
\end{lstlisting}

\subsection{Gwarancje wyjątkowe i \texttt{noexcept}}

Funkcja modyfikująca stan może gwarantować: \emph{brak skutków po
niepowodzeniu} (silna gwarancja), zachowanie niezmienników mimo częściowej
zmiany (podstawowa gwarancja) albo brak wyjątków (\code{noexcept}).
\code{noexcept} nie znaczy, że błąd jest niemożliwy: nieobsłużony
wyjątek kończy wtedy program przez \code{std::terminate}.

\begin{lstlisting}[caption={Silna gwarancja przez przygotowanie kopii.}]
#include <string>
#include <utility>

class Settings {
    std::string data_;
public:
    void update(std::string next) {
        // Wszystkie kosztowne operacje wykonano przy
        // tworzeniu argumentu 'next'. Swap string ma noexcept.
        data_.swap(next);
    }
    const std::string& get() const noexcept { return data_; }
};
\end{lstlisting}

\section{Krycie implementacji: PImpl i granica ABI}
\label{sec:cpp20-pimpl}

Nagłówki udostępniające wszystkie pola klas narzucają zależności
kompilacyjne, a zmiana ich układu może naruszać ABI bibliotek
współdzielonych. Idiom PImpl (\emph{pointer to implementation})
przenosi szczegóły do osobnego typu. Cena: alokacja, pośredni dostęp
i potrzeba poprawnego zarządzania funkcjami specjalnymi.

\begin{lstlisting}[caption={widget.hpp -- publiczny naglowek.}]
#pragma once
#include <memory>
#include <string>

class Widget {
public:
    explicit Widget(std::string name);
    ~Widget();
    Widget(Widget&&) noexcept;
    Widget& operator=(Widget&&) noexcept;
    Widget(const Widget&) = delete;
    Widget& operator=(const Widget&) = delete;
    std::string name() const;
private:
    struct Impl;
    std::unique_ptr<Impl> impl_;
};
\end{lstlisting}

\begin{lstlisting}[caption={widget.cpp -- pelna definicja ukrytego typu.}]
#include "widget.hpp"
#include <utility>

struct Widget::Impl {
    std::string name;
    explicit Impl(std::string n) : name(std::move(n)) {}
};

Widget::Widget(std::string name)
    : impl_(std::make_unique<Impl>(std::move(name))) {}
Widget::~Widget() = default;
Widget::Widget(Widget&&) noexcept = default;
Widget& Widget::operator=(Widget&&) noexcept = default;
std::string Widget::name() const { return impl_->name; }
\end{lstlisting}

\begin{lstlisting}[caption={main.cpp -- uzycie niezalezne od szczegolow Impl.}]
#include "widget.hpp"
#include <iostream>

int main() {
    Widget w("sample");
    std::cout << w.name() << '\n';
}
\end{lstlisting}

Destruktor \code{Widget} jest definiowany \emph{po} pełnej definicji
\code{Impl}. Jest to istotne dla destrukcji \code{unique_ptr<Impl>}.
Wersja z domyślnym przenoszeniem dopuszcza obiekt \code{Widget}
w stanie \emph{moved-from} z pustym \code{impl_}; użytkownik nie może
wtedy wywoływać \code{name()} bez dodatkowo określonego kontraktu.
Jeśli wymagane jest, by każdy \code{Widget} zawsze miał implementację,
należy inaczej zaprojektować przenoszenie.

\subsection{Stabilność ABI nie jest automatyczna}

PImpl stabilizuje \emph{układ prywatnego stanu}, ale samo w sobie nie
zapewnia zgodności między kompilatorami, standardowymi bibliotekami,
wersjami runtime ani opcjami kompilacji. Publiczne typy
\code{std::string} lub \code{std::vector} mogą również wpływać na ABI.
Na granicy modułów współdzielonych wymagających zgodności wielu
środowisk warto rozważyć wąski interfejs C: nieprzezroczysty uchwyt,
stałe formaty danych i jawne funkcje tworzenia/niszczenia.

\section{Polimorfizm: dynamiczny, statyczny i wymazywanie typów}

Do rozszerzalnego API można zastosować wirtualną klasę bazową,
typ generyczny oparty na concepts, \code{variant} dla zamkniętego
zbioru alternatyw lub \emph{type erasure}. Mechanizmy różnią się
kosztami, możliwością rozbudowy i wymaganiami własności.

\begin{lstlisting}[caption={Interfejs dynamiczny i bezpieczne niszczenie polimorficzne.}]
#include <memory>
#include <string>

struct Renderer {
    virtual ~Renderer() = default;
    virtual std::string render(int value) const = 0;
};

struct SimpleRenderer final : Renderer {
    std::string render(int value) const override {
        return std::to_string(value);
    }
};

std::unique_ptr<Renderer> make_renderer() {
    return std::make_unique<SimpleRenderer>();
}
\end{lstlisting}

\begin{center}\small
\begin{tabularx}{\textwidth}{@{}p{0.25\textwidth}X@{}}
\toprule
Mechanizm & Atuty i koszty \\
\midrule
\code{virtual} & Otwarta hierarchia, wywołania pośrednie, ABI wrażliwe na zmianę interfejsu. \\
\code{template}/concept & Weryfikacja i specjalizacja przy kompilacji, więcej instancji kodu. \\
\code{variant} & Zamknięta lista typów, jawna obsługa przypadków, możliwy wzrost rozmiaru obiektu. \\
Type erasure & Jednolity typ opakowania, ukryta implementacja, zwykle pośredni dispatch i ewentualna alokacja. \\
\bottomrule
\end{tabularx}
\end{center}

\section{Punkty dostosowania: ADL, hidden friends i CPO}
\label{sec:cpp20-adl}

\emph{Argument-dependent lookup} (ADL) wyszukuje wolne funkcje
w przestrzeniach nazw związanych z typami argumentów. Pozwala
implementować operacje jak \code{swap}, \code{begin} czy
\code{operator==} bez modyfikowania biblioteki głównej.
\emph{Hidden friend} jest funkcją zaprzyjaźnioną zdefiniowaną
wewnątrz klasy; można ją odnaleźć przez ADL, ale nie przez zwykłe
wyszukiwanie nazwy bez odpowiedniego argumentu.

\begin{lstlisting}[caption={Hidden friend i idiom wywolywania swap przez ADL.}]
#include <utility>

namespace geom {
struct Point {
    int x{};
    int y{};
    friend bool operator==(const Point&, const Point&) = default;
    friend void swap(Point& a, Point& b) noexcept {
        using std::swap;
        swap(a.x, b.x);
        swap(a.y, b.y);
    }
};
}

template<class T>
void exchange(T& a, T& b) {
    using std::swap;
    swap(a, b); // moze znalezc geom::swap przez ADL
}
\end{lstlisting}

Nowoczesna biblioteka Ranges stosuje \emph{customization point objects}
(CPO), np. \code{std::ranges::begin}. Nie należy utożsamiać ich
z dowolnymi wolnymi funkcjami; obiekt CPO ma własne reguły wywołania,
weryfikacji wymagań i wyszukiwania implementacji.

\subsection{Wzorzec \texttt{tag\_invoke} -- nie standardowe API}

\code{tag_invoke} jest popularną konwencją proponowaną dla rozszerzalnych
CPO, ale \emph{nie jest standardową funkcją C++20 ani C++23}.
W bibliotekach, które ją przyjmują, użytkownik implementuje
\code{tag_invoke(tag, arguments...)} w przestrzeni nazw swego typu,
a biblioteka znajduje ją za pośrednictwem ADL. Bez własnej implementacji
mechanizmu nie wolno pisać \code{std::tag_invoke}.

\begin{lstlisting}[caption={Minimalna ilustracja konwencji tag\_invoke (fragment).}]
struct serialize_tag {};

namespace model {
struct Record { int id; };

// Odnajdywane przez ADL, gdy argumentem jest Record.
int tag_invoke(serialize_tag, const Record& r) {
    return r.id;
}
}

int serialize(const model::Record& r) {
    return tag_invoke(serialize_tag{}, r);
}
\end{lstlisting}

\section{Policy-based design i wstrzykiwanie zależności}

W \emph{policy-based design} strategia jest parametrem typu.
Wybrana implementacja może zostać wbudowana przez kompilator,
ale zwiększa liczbę instancji szablonów i zależności od nagłówków.
W \emph{dependency injection} obiekt otrzymuje zależności w konstruktorze
lub jako argument, zamiast samodzielnie tworzyć globalne usługi.

\begin{lstlisting}[caption={Polityka zapisu z kontrola w czasie kompilacji (C++20).}]
#include <concepts>
#include <string_view>

struct NullSink {
    void write(std::string_view) const noexcept {}
};

struct CounterSink {
    int* counter;
    void write(std::string_view) const noexcept { ++*counter; }
};

template<class T>
concept Sink = requires(T s, std::string_view text) {
    { s.write(text) } -> std::same_as<void>;
};

template<Sink S>
class Logger {
    S sink_;
public:
    explicit Logger(S sink) : sink_(sink) {}
    void log(std::string_view text) { sink_.write(text); }
};
\end{lstlisting}

Ograniczenie kompilacyjne nie zastępuje dokumentacji: każda polityka
\code{Sink} musi sama zapewniać odpowiednie własności współbieżne,
a \code{Logger} nie powinien obiecywać większej odporności na błędy
niż implementacja przekazanej polityki.
W pokazanym \code{CounterSink} wskaźnik musi pozostać ważny przez czas
używania loggera; dwa wątki nie mogą jednocześnie zwiększać tego samego
zwykłego licznika bez synchronizacji.

\section{Bezpieczeństwo wątkowe i reentrancy}

\emph{Thread-safe} nie jest jednoznacznym stwierdzeniem bez określenia
zakresu. Dokumentacja powinna mówić, czy dwa wątki mogą jednocześnie
wywoływać metody na tym samym obiekcie, czy tylko na różnych obiektach,
oraz czy zezwala się na jednoczesne czytanie i modyfikowanie.
\emph{Reentrant} oznacza, że wywołanie może bezpiecznie nakładać się
z kolejnym wywołaniem tej samej funkcji w określonym modelu; nie jest
automatycznym synonimem pełnej \emph{thread safety}.

Warto jawnie określić, czy callback dostarczony przez użytkownika
jest uruchamiany pod blokadą. Wywołanie obcego kodu przy zablokowanym
mutexie może spowodować deadlock, rekurencję lub trudny do przewidzenia
czas trzymania blokady.

\section{Publiczny nagłówek, moduły i konfiguracja CMake}

Czytelna biblioteka oddziela pliki \code{include/}, \code{src/},
\code{tests/} i \code{examples/}. W publicznym nagłówku umieszczamy
minimalny zbiór deklaracji potrzebnych odbiorcy, a nie \code{using
namespace std;} ani zależności od szczegółów implementacji.

\begin{lstlisting}[caption={Przykladowe CMakeLists.txt dla biblioteki publicznej.},language={}]
cmake_minimum_required(VERSION 3.20)
project(geometry VERSION 1.0 LANGUAGES CXX)

add_library(geometry src/geometry.cpp)
target_compile_features(geometry PUBLIC cxx_std_20)
target_include_directories(geometry PUBLIC
    $<BUILD_INTERFACE:${CMAKE_CURRENT_SOURCE_DIR}/include>
    $<INSTALL_INTERFACE:include>)

add_executable(example examples/main.cpp)
target_link_libraries(example PRIVATE geometry)
\end{lstlisting}

\nota{Deklaracje modułów C++20 mogą zmniejszyć zależność od klasycznego
preprocesora, ale same w sobie nie zapewniają stabilności ABI.
Konfiguracja modułów i dostępność \code{import std;} zależą od
wersji narzędzi; szczegóły omówiono w rozdziale~\ref{chap:cpp17-build-abi-modules}.}

\section{Wersjonowanie i testowanie kontraktu}

Zmiana kodu wewnętrznego nie jest zmianą API, dopóki zachowuje
zewnętrznie obiecane zachowanie. Zmiana przeciążenia, domyślnego
argumentu, wyjątków, czasu życia zwracanych referencji lub układu
publicznej klasy może mieć znacznie większe konsekwencje niż sugeruje
rozmiar poprawki. Warto odróżniać zgodność źródłową, ABI oraz zgodność
semantyczną.

\begin{enumerate}
\item Testuj kontrakt na granicy publicznych interfejsów, a nie wyłącznie
      metody prywatne; dodawaj testy przypadków błędnych.
\item Kompiluj testowego klienta w osobnym projekcie, korzystając tylko
      z publicznych nagłówków lub eksportowanych modułów.
\item Sprawdzaj zachowanie dla lvalue, rvalue, typów const oraz wyjątków
      we wszystkich operacjach z własnością zasobów.
\item Jeśli biblioteka obiecuje ABI, testuj zmianę wersji z programem
      zlinkowanym względem poprzedniego wydania.
\item Mierz koszt użycia biblioteki: narzut alokacji, wielkość binarium,
      opóźnienia i zachowanie w wielu wątkach.
\end{enumerate}

\section{Studium przypadku: API biblioteki numerycznej}

Załóżmy, że biblioteka oblicza iloczyn skalarny i przechowuje operator
liniowy. Dobry interfejs rozdziela widok danych od stanu własnego,
kontroluje zgodność wymiarów i określa zachowanie dla wyjątków.

\begin{lstlisting}[caption={Niewlasnosciowe wejscie do funkcji numerycznej (C++20).}]
#include <span>
#include <stdexcept>

[[nodiscard]] double dot(std::span<const double> lhs,
                         std::span<const double> rhs) {
    if (lhs.size() != rhs.size())
        throw std::invalid_argument("dimension mismatch");
    double acc = 0;
    for (std::size_t i = 0; i != lhs.size(); ++i)
        acc += lhs[i] * rhs[i];
    return acc;
}
\end{lstlisting}

Przy zastosowaniach wymagających powtarzalności numerycznej trzeba
udokumentować metodę sumowania i ewentualne przestawianie operacji.
Implementacja równoległa nie musi dać bitowo identycznego wyniku jak
sekwencyjna. Jeżeli biblioteka ma wspierać akceleratory i macierze
wielowymiarowe, należy oddzielić interfejs algorytmu od szczegółów
alokacji i układu pamięci.

\section{C API jako granica bibliotek współdzielonych}

Kiedy biblioteka ma być używana przez różne języki lub kompilatory,
wąski interfejs C bywa stabilniejszą granicą niż publiczna klasa C++.
Użytkownik nie musi znać rozmiaru obiektu, a biblioteka zachowuje
pełną kontrolę nad alokacją i zniszczeniem. Dwa podstawowe warunki to
\emph{nieprzekraczanie granicy przez wyjątki C++} oraz jawny kontrakt
zwalniania każdego zasobu.

\begin{lstlisting}[caption={calc\_api.h -- przyklad nieprzezroczystego uchwytu.}]
#pragma once
#ifdef __cplusplus
extern "C" {
#endif

typedef struct calc_handle calc_handle;
calc_handle* calc_create(int initial_value);
int calc_get(const calc_handle* handle, int* output);
void calc_destroy(calc_handle* handle);

#ifdef __cplusplus
}
#endif
\end{lstlisting}

\begin{lstlisting}[caption={calc\_api.cpp -- implementacja z obsluga bledu alokacji.}]
#include "calc_api.h"
#include <new>

struct calc_handle { int value; };

extern "C" calc_handle* calc_create(int value) {
    return new (std::nothrow) calc_handle{value};
}
extern "C" int calc_get(const calc_handle* h, int* out) {
    if (!h || !out) return 1;
    *out = h->value;
    return 0;
}
extern "C" void calc_destroy(calc_handle* h) {
    delete h;
}
\end{lstlisting}

Interfejs jest ilustracyjny. W realnym systemie warto wprowadzić
kody błędów, wersję ABI, zasady bezpieczeństwa wątkowego i politykę
alokacji. Nie należy zakładać, że wyjątki C++ mogą bezpiecznie
przechodzić przez granicę C lub modułu skompilowanego innym
kompilatorem. Szczególnie ważne jest, aby obiekt utworzony przez
\code{calc_create} zawsze niszczyć przez \code{calc_destroy},
a nie przez \code{free} ani \code{delete} po stronie klienta.

\section{Projektowanie wydajności: alokatory, PMR i batching}

Interfejs może narzucać niepotrzebne alokacje przez zwracanie
osobnych kontenerów dla każdej małej operacji. Często warto
rozważyć wariant zapisujący do bufora dostarczonego przez
wywołującego albo przetwarzanie wsadowe (\emph{batching}). Nie znaczy
to jednak, że każde API powinno przyjmować ręcznie zarządzane wskaźniki:
lepsza wydajność musi być uzasadniona pomiarem.

Biblioteka \code{std::pmr} pozwala oddzielić politykę alokacji od
typu danych. Przyjmowanie \code{std::pmr::memory_resource*} oznacza
jednak dodatkowe zobowiązanie: zasób pamięci musi istnieć co najmniej
tak długo, jak kontenery wykorzystujące go do zwalniania pamięci.

\begin{lstlisting}[caption={Kontener PMR z buforem przydzielonym przez klienta (C++17).}]
#include <array>
#include <cstddef>
#include <memory_resource>
#include <vector>

int main() {
    std::array<std::byte, 2048> storage{};
    std::pmr::monotonic_buffer_resource pool(
        storage.data(), storage.size());
    std::pmr::vector<int> values(&pool);
    for (int i = 0; i < 100; ++i) values.push_back(i);
}
\end{lstlisting}

\code{monotonic_buffer_resource} najpierw korzysta z tablicy
\code{storage}, ale po jej zapełnieniu może pobrać dodatkową pamięć
z zasobu nadrzędnego. Sama tablica nie narzuca więc twardego limitu.

\section{Dokumentacja API na przykładzie kontraktu operacji}

Dla każdej publicznej funkcji warto zapisać krótki, ale precyzyjny
kontrakt. W przypadku funkcji \code{dot(span,span)} może on brzmieć:

\begin{description}[style=nextline]
\item[Warunki wstępne] Obydwa widoki wskazują istniejące, czytelne
bufory, których czas życia obejmuje całe wywołanie.
\item[Wynik] Iloczyn skalarny w arytmetyce \code{double}.
\item[Błędy] Dla różnych długości zgłaszany jest
\code{std::invalid_argument}.
\item[Wyjątki] Przy błędzie argumentów żaden bufor nie jest
modyfikowany; sama funkcja nie przejmuje własności.
\item[Współbieżność] Równoczesne wywołania są dozwolone,
jeśli inni uczestnicy nie modyfikują tych samych elementów
w sposób powodujący wyścig danych.
\item[Wydajność] Koszt liniowy względem liczby elementów,
bez konieczności dodatkowej alokacji.
\end{description}

Umowa nie powinna obiecywać braku zależności od kolejności dodawania
dla liczb zmiennoprzecinkowych. Jeżeli kod ma gwarantować określony
błąd numeryczny, taki warunek trzeba jawnie sformułować i uzasadnić.

\section{Ewolucja API i polityka kompatybilności}

Podczas wydawania kolejnej wersji warto rozważyć odrębnie:
\emph{source compatibility}, \emph{binary compatibility} i
\emph{behavioral compatibility}. Na przykład dodanie publicznego
przeciążenia może prowadzić do niejednoznaczności wcześniej
poprawnego wywołania, mimo że żadna stara sygnatura nie została
usunięta. Zmiana wartości domyślnego argumentu może z kolei
prowadzić do innego zachowania po kompilacji klienta, choć ABI
pozostaje bez zmian.

W projektach opartych na bibliotekach współdzielonych stosuje się
niekiedy \emph{symbol versioning} lub oddzielne numery wersji ABI.
Nie ma jednak jednego mechanizmu przenośnego między wszystkimi
platformami. Dla bibliotek publicznych szczególnie przydatne są
przykłady minimalnych klientów kompilowane w macierzy wersji oraz
testy odtwarzające zachowanie poprzedniej wersji.

\section{Lista kontrolna przeglądu biblioteki}

\begin{enumerate}
\item Czy publiczne nagłówki kompilują się samodzielnie bez przypadkowych
zależności pośrednich?
\item Czy własność, możliwość \code{nullptr}, czas życia oraz
odpowiedzialność za zwalnianie zasobów są jednoznaczne?
\item Czy koncepty odpowiadają wszystkim wymaganym operacjom kodu?
\item Czy zachowanie przy błędzie jest testowane, łącznie z wyjątkami
z konstruktorów i callbacków użytkownika?
\item Czy funkcje \code{noexcept} naprawdę nie przepuszczają wyjątków?
\item Czy zbędne alokacje i niejawne kopiowanie są zmierzone, a nie
tylko zakładane?
\item Czy API można przetestować z implementacją testową zależności,
bez uruchamiania sieci ani bazy danych?
\item Czy publiczne typy nie ujawniają szczegółów, które trzeba
będzie zmienić w przyszłości?
\item Czy dokumentacja dokładnie określa dopuszczalne wywołania
współbieżne?
\end{enumerate}


\section*{Podsumowanie rozdziału}
\addcontentsline{toc}{section}{Podsumowanie rozdziału}

Projektowanie bibliotek C++ polega przede wszystkim na precyzyjnym
wyrażeniu zobowiązań. Typy takie jak \code{span}, \code{unique_ptr},
\code{optional}, concepts i interfejsy wirtualne są narzędziami do
modelowania tych zobowiązań. PImpl pomaga ukryć implementację; ADL
i CPO pomagają projektować rozszerzalność; testy, wersjonowanie
oraz dokumentacja wyznaczają granice rzeczywistej niezawodności.
Najważniejsza zasada brzmi: nie obiecuj silniejszych gwarancji własności,
współbieżności, bezpieczeństwa wyjątkowego ani ABI, niż możesz
zapewnić i przetestować.

\section*{Literatura i dokumentacja}
\addcontentsline{toc}{section}{Literatura i dokumentacja}
\begingroup\sloppy
\begin{itemize}
\item ISO C++ Core Guidelines, \url{https://isocpp.github.io/CppCoreGuidelines/CppCoreGuidelines}.
\item \url{https://en.cppreference.com/w/cpp/language/adl} --- ADL i hidden friends.
\item \url{https://en.cppreference.com/w/cpp/ranges} --- customization point objects Ranges.
\item \url{https://en.cppreference.com/w/cpp/language/pimpl} --- idiom PImpl.
\item \url{https://cmake.org/cmake/help/latest/guide/tutorial/} --- CMake.
\item Klaus Iglberger, \emph{C++ Software Design}, O'Reilly, 2022.
\end{itemize}
\endgroup

% Rozdzial 21. Samodzielna tresc rozdzialu do \input (bez preambuly).
% Wymaga listings, enumitem, tabularx, booktabs, hyperref, makr
% \code, \term, \nota. Przyklady: standard oznaczony przy kazdym kodzie.
% Stan omowienia: pazdziernik 2026. Nalezy sprawdzac stan wdrozenia API.
\chapter{Nowości C++23 i C++26: język, biblioteka i przenośność}
\label{chap:cpp21-new}

\section{Jak czytać informacje o nowym standardzie?}

Zatwierdzenie funkcjonalności przez komitet C++ nie oznacza, że
wszystkie kompilatory i biblioteki standardowe już ją dostarczają.
W szczególności trzeba odróżniać: (1) zgodność kompilatora z regułami
\emph{języka}, (2) implementację odpowiedniego nagłówka w \emph{bibliotece
standardowej} oraz (3) gotowość systemu budowania i pakietowania.
Opisy w tym rozdziale dotyczą C++23 i pakietu funkcji C++26
według stanu na październik 2026; dostępność należy sprawdzić dla
konkretnego środowiska GCC/libstdc++, Clang/libc++ lub MSVC/STL.

\begin{center}\small
\begin{tabularx}{\textwidth}{@{}p{0.30\textwidth}X@{}}
\toprule
Warstwa & Praktyczna weryfikacja \\
\midrule
Standard języka & \code{-std=c++23}, \code{-std=c++2c} (lub wariant kompilatora), \code{__cplusplus} i makra \code{__cpp_*}. \\
Biblioteka standardowa & \code{<version>}, makra \code{__cpp_lib_*}, dostępność nagłówków i linkowania. \\
Narzędzia budowania & Wsparcie CMake, modułów, diagnostyki i zależności BMI. \\
Platforma & Implementacja wielowątkowości, wektoryzacji, I/O, architektury procesora. \\
\bottomrule
\end{tabularx}
\end{center}

\subsection{Feature-test macros zamiast testowania numeru kompilatora}

Nie należy zakładać, że \code{__cplusplus >= 202302L} oznacza obecność
każdego nowego nagłówka. Wykrywamy \emph{konkretną funkcję}.
\code{__has_include} sprawdza wyłącznie obecność pliku nagłówkowego;
nie dowodzi, że zawartość została w pełni zaimplementowana.

\begin{lstlisting}[caption={Bezpieczny wybor std::expected albo wariantu zgodnego z C++20.}]
#include <version>
#include <iostream>

#if defined(__cpp_lib_expected) && __cpp_lib_expected >= 202202L
    #include <expected>
    #define HAS_EXPECTED 1
#else
    #include <optional>
    #define HAS_EXPECTED 0
#endif

int main() {
#if HAS_EXPECTED
    std::expected<int, int> result = 42;
    std::cout << result.value() << '\n';
#else
    std::optional<int> result = 42;
    std::cout << *result << '\n';
#endif
}
\end{lstlisting}

Wartości makr zmieniają się, kiedy standard rozszerza dane API.
W bibliotekach produkcyjnych zaleca się test konfiguracji
\emph{compile-and-link}, a nie tylko sprawdzenie preprocesora.

\section{C++23: \texttt{std::expected<T,E>} i jawne błędy}

\begin{center}\small
\begin{tabularx}{\textwidth}{@{}p{0.35\textwidth}X@{}}
\toprule
Operacja & Znaczenie \\
\midrule
\code{std::expected<T,E>} & Przechowuje wartość \code{T} albo błąd \code{E}. \\
\code{std::unexpected(error)} & Tworzy wariant błędu przekazywany do \code{expected}. \\
\code{has_value()} & Sprawdza, czy dostępna jest wartość, a nie błąd. \\
\code{value()} & Zwraca wartość; gdy przechowywany jest błąd, zgłasza \code{std::bad_expected_access}. \\
\code{error()} & Udostępnia błąd; należy używać, gdy obiekt przechowuje błąd. \\
\code{value_or(fallback)} & Zwraca wartość albo podany zamiennik. \\
\bottomrule
\end{tabularx}
\end{center}
\label{sec:cpp21-expected}

\code{std::expected<T,E>} reprezentuje wynik typu \code{T} albo błąd
typu \code{E}. W przeciwieństwie do \code{std::optional<T>} przenosi
informację, \emph{dlaczego} operacja się nie powiodła. Nie jest
zamiennikiem wyjątków w każdej sytuacji: użytkownik musi obsłużyć
stan błędu, a \code{value()} może zgłosić \code{bad_expected_access}.

\begin{lstlisting}[caption={Parser z jawnym typem bledu (kompletny program C++23).}]
#include <charconv>
#include <expected>
#include <iostream>
#include <string>
#include <string_view>

std::expected<int, std::string> parse_integer(std::string_view text) {
    int value{};
    auto [ptr, error] = std::from_chars(text.data(),
                                        text.data() + text.size(),
                                        value);
    if (error != std::errc{} || ptr != text.data() + text.size())
        return std::unexpected("invalid integer");
    return value;
}

int main() {
    auto result = parse_integer("123");
    if (result) std::cout << *result << '\n';
    else std::cerr << result.error() << '\n';
}
\end{lstlisting}

Nowsze operacje monadyczne \code{and_then}, \code{transform},
\code{or_else} i \code{transform_error} można wykorzystywać po
sprawdzeniu właściwego poziomu makra \code{__cpp_lib_expected}.
\code{expected} daje się również stosować z \code{void} jako typem
sukcesu, co jest przydatne w operacjach zapisujących dane.

\section{C++23: \texttt{std::print}, \texttt{std::format} i formatowanie}

\code{std::format} jest częścią C++20, natomiast \code{std::print}
i \code{std::println} wprowadzono w C++23. Ich wsparcie w praktyce
zależy głównie od wersji biblioteki standardowej. Formatowanie
stabilnych protokołów danych nadal wymaga jawnego kontraktu, a nie
przypadkowego wydruku do terminala.

\begin{lstlisting}[caption={Formatowany wydruk C++23 (wymaga naglowka <print>).}]
#include <print>

int main() {
    std::println("value = {:.2f}", 3.14159);
    std::print("without newline: {}", 42);
}
\end{lstlisting}

\code{<format>} oferuje bogatszy język formatowania niż ręczne
sklejanie napisów, ale jego użycie dla danych nieufnych wymaga
kontroli formatu i argumentów. Należy odróżniać tekst formatu
znany przy kompilacji od wariantów wykorzystujących format dynamiczny.

\section{C++23: \texttt{std::mdspan} i macierze wielowymiarowe}
\label{sec:cpp21-mdspan}

\code{std::mdspan} jest niewłasnościowym widokiem na pamięć
wielowymiarową. Podobnie jak \code{std::span} nie zarządza czasem życia
bufora. Pozwala opisać ekstensje, odwzorowanie indeksów na adresy
i sposób dostępu. Układy \code{layout_right} i \code{layout_left}
mają inny sposób wyliczania przesunięcia w pamięci; nie należy
przypisywać im automatycznie takiej samej wydajności.

\begin{lstlisting}[caption={Macierz 2 x 3 przechowywana w vector (C++23).}]
#include <mdspan>
#include <vector>
#include <iostream>

int main() {
    std::vector<double> a{1, 2, 3, 4, 5, 6};
    std::mdspan matrix(a.data(), 2, 3);
    matrix[1, 2] = 10.0;  // operator indeksowania C++23
    std::cout << a[5] << '\n';
}
\end{lstlisting}

\nota{Dokładna składnia operatora indeksowania i dostępność nagłówka
\code{<mdspan>} mogą być zależne od implementacji i późniejszych
poprawek. Nie wolno zwracać \code{mdspan} do lokalnego wektora,
który zostanie zniszczony po wyjściu z funkcji.}

\section{C++23: generatory, wymazywanie typu i inne klasy użytkowe}

\subsection{\texttt{std::generator}: od korutyny do zakresu}

\code{std::generator<T>} z nagłówka \code{<generator>} jest standardowym
narzędziem C++23 do leniwego przekazywania wartości przez \code{co_yield}.
Nie każdy kompilator udostępnia tę bibliotekę nawet wtedy, gdy
obsługuje składnię korutyn C++20. Generator nie tworzy automatycznie
oddzielnego wątku ani mechanizmu planowania asynchronicznego.

\begin{lstlisting}[caption={Leniwy generator liczb calkowitych (C++23; wymaga <generator>).}]
#include <generator>
#include <iostream>

std::generator<int> numbers(int limit) {
    for (int i = 0; i < limit; ++i)
        co_yield i;
}

int main() {
    for (int v : numbers(4))
        std::cout << v << ' ';
}
\end{lstlisting}

\subsection{\texttt{move\_only\_function} i koszt własności callbacku}

\code{std::move_only_function<R(Args...)>} jest odpowiednikiem
opisu wywoływalnego obiektu, który \emph{nie wymaga kopiowalności}
przechwyconego stanu. Ma zastosowanie m.in. w kolejkach zadań
przechowujących lambdy zawierające \code{unique_ptr}.

\begin{lstlisting}[caption={Callback z obiektem posiadajacym wylaczna wlasnosc (C++23).}]
#include <functional>
#include <iostream>
#include <memory>
#include <utility>

int main() {
    auto payload = std::make_unique<int>(42);
    std::move_only_function<int()> get =
        [p = std::move(payload)] { return *p; };
    std::cout << get() << '\n';
}
\end{lstlisting}

Ponadto C++23 dodaje m.in. \code{std::byteswap}, \code{std::to_underlying},
\code{std::forward_like}, \code{std::bind_back}, \code{std::flat_map},
\code{std::flat_set}, \code{std::stacktrace} i \code{<spanstream>}.
Każda z tych konstrukcji ma własny zakres użycia i niezależne
wymagania implementacyjne.

\section{C++23: Ranges i materializacja widoków}
\label{sec:cpp21-ranges}

C++23 rozbudowuje bibliotekę Ranges m.in. o \code{views::zip},
\code{views::chunk}, \code{views::slide}, \code{views::stride},
\code{views::enumerate} i \code{ranges::to}. Część algorytmów
czyta wyniki leniwie; \code{ranges::to} materializuje widok w
kontenerze. To duża zmiana ergonomii, ale reguły czasu życia
widoków i koszt materializacji nadal obowiązują.

\begin{lstlisting}[caption={Przekazanie wynikow widoku do wektora (C++23).}]
#include <iostream>
#include <ranges>
#include <vector>

int main() {
    auto values = std::views::iota(0, 10)
                | std::views::filter([](int x) { return x % 2 == 0; });
    auto result = std::ranges::to<std::vector<int>>(values);
    std::cout << result.size() << '\n'; // 5
}
\end{lstlisting}

Dla funkcji, których implementacja nie wspiera jeszcze \code{ranges::to},
można zastosować ręczną iterację i \code{push_back} albo
\code{ranges::copy} wraz z odpowiednim inserterem.

\section{C++23: składnia języka --- deducing this i \texttt{if consteval}}

\emph{Explicit object parameter} (często nazywany \emph{deducing this})
pozwala jawnie opisać obiekt, na którym wywoływana jest funkcja
składowa. Może zastąpić rodzinę przeciążeń kwalifikowanych
\code{&}, \code{const&} oraz \code{&&}, ale nadmierna ogólność
może uczynić kontrakt funkcji mniej czytelnym.

\begin{lstlisting}[caption={Jawny parametr obiektu (C++23).}]
#include <iostream>

struct Counter {
    int value = 0;
    void add(this Counter& self, int step) {
        self.value += step;
    }
};

int main() {
    Counter x;
    x.add(3);
    std::cout << x.value << '\n';
}
\end{lstlisting}

\code{if consteval} pozwala wykonać inną gałąź podczas ewaluacji
stałej, bez traktowania zwykłej funkcji czasu wykonywania jako
\code{consteval}. Trzeba jednak nadal pamiętać o ograniczeniach
konstruowania stałych wyrażeń.

\begin{lstlisting}[caption={if consteval (C++23, kompletny program).}]
#include <iostream>

constexpr int double_value(int n) {
    if consteval {
        return n * 2;
    } else {
        return n + n;
    }
}
static_assert(double_value(21) == 42);
int main() { std::cout << double_value(3) << '\n'; }
\end{lstlisting}

W rzeczywistym projekcie rozdzielenie ścieżek powinno być
uzasadnione innym dostępnym zestawem mechanizmów na etapie
kompilacji i wykonania, a nie sztucznym powielaniem identycznego kodu.

\section{C++26: stały rozmiar przechowywania bez alokacji}
\label{sec:cpp21-inplace}

\code{std::inplace_vector<T,N>} jest kontenerem o zmiennym rozmiarze,
którego \emph{maksymalna pojemność} wynosi \code{N}, a pamięć dla
elementów jest umieszczona w samym obiekcie. Różni się więc od
\code{std::array<T,N>}, która zawsze zawiera \code{N} elementów.
Nie jest zwykłym \code{vector}: przy przekroczeniu pojemności
nie może zwiększyć jej przez alokację.

\begin{lstlisting}[caption={Kontener inplace\_vector (C++26; wymagane wsparcie biblioteki).}]
#include <inplace_vector>
#include <iostream>

int main() {
    std::inplace_vector<int, 8> values;
    values.push_back(10);
    values.push_back(20);
    std::cout << values.size() << '\n';
}
\end{lstlisting}

\code{std::hive} to kolejna nowa struktura kontenerowa C++26,
ukierunkowana na efektywne wstawianie i usuwanie elementów oraz
określone gwarancje stabilności istniejących elementów.
Nie stanowi zamiennika \code{vector} dla obliczeń wymagających
ciągłego bloku pamięci i lokalności odczytów.

\section{C++26: SIMD, algebra liniowa i wielowątkowość}
\label{sec:cpp21-simd}

Nagłówek \code{<simd>} i typy \code{std::simd} służą do jawnego
wyrażania obliczeń na paczkach wartości, potencjalnie mapowanych
na instrukcje wektorowe procesora. Ich wykorzystanie nie zwalnia
z pomiarów; wydajność zależy od architektury, wyrównania,
przepustowości pamięci i zachowania kompilatora.

\begin{lstlisting}[caption={Schemat obliczenia na SIMD (C++26, fragment pogladowy).}]
// #include <simd>
// Wymaga biblioteki z implementacja C++26.
// std::simd<float> a(1.0f), b(2.0f);
// auto c = a + b;
// Typ i sposob ladowania zalezy od przyjetego ABI/implementacji.
\end{lstlisting}

C++26 rozwija także narzędzia algebry liniowej w \code{<linalg>}
i mechanizmy asynchroniczne oparte na \emph{senders/receivers}
(\code{<execution>}). Nowy model wykonania opisuje kompozycję zadań,
obsługę błędów, anulowanie i relację z harmonogramem pracy.
Sama obecność nagłówka nie dowodzi, że dany executor lub scheduler
działa w pożądany sposób; szczegółowe semantyki zależą od wybranych
operacji i implementacji.

\subsection{Bezpieczne odzyskiwanie pamięci}

Nagłówki \code{<hazard_pointer>} i \code{<rcu>} w C++26 mają pomóc
w zarządzaniu czasem życia obiektów w strukturach współbieżnych.
Nie usuwają one automatycznie problemu ABA ani nie sprawiają,
że dowolna struktura z \code{std::atomic<T*>} jest bezpieczna.
Należy osobno uzasadnić porządek pamięci, gwarancje postępu
i moment fizycznego zniszczenia obiektów (por. rozdział 16).

\section{C++26: język -- pack indexing, refleksja i kontrakty}

C++26 wprowadza mechanizm \emph{pack indexing}, czyli wybór
określonego elementu pakietu parametrów bez rekurencyjnego
rozpakowywania. Jego wsparcie zależy od kompilatora.

\begin{lstlisting}[caption={Indeksowanie pakietu typow (C++26, tylko kompilator ze wsparciem).}]
template<class... Ts>
using first_type = Ts...[0];

static_assert(sizeof(first_type<int, double>) == sizeof(int));
\end{lstlisting}

W C++26 rozwijane są również: \emph{static reflection} (nagłówek
\code{<meta>} i odpowiednia składnia refleksji) oraz
\emph{kontrakty} (np. asercje kontraktowe i wsparcie z
\code{<contracts>}). Są to znaczące nowe możliwości,
ale składnia i zakres praktycznego użycia wymagają konkretnego
kompilatora z odpowiednim wsparciem. Nie podajemy tu pozornie
przenośnego przykładu refleksji lub kontraktów, który mógłby
nie kompilować się w typowej instalacji C++26.

\subsection{Co dzieje się z \texttt{constexpr}?}

W kolejnych rewizjach standardu zwiększa się zakres dozwolonych
obliczeń stałych. Nie wynika z tego, że wszystkie operacje I/O,
zasoby systemowe i dynamiczne struktury mogą automatycznie istnieć
w stałym wyrażeniu. Warto oddzielać możliwość \emph{wykonania obliczenia}
przy kompilacji od zdolności utworzenia \emph{trwałego obiektu}
z zasobami zarządzanymi dynamicznie.

\section{Jak wdrażać C++23/26 do istniejącej biblioteki?}
\label{sec:cpp21-migration}

Włączenie nowego standardu dla całego projektu z dnia na dzień jest
często mniej korzystne niż migracja etapami. Proponujemy strategię:

\begin{enumerate}
\item Spisać minimalne wersje obsługiwanych GCC, Clang, MSVC oraz
      implementacji bibliotek standardowych.
\item Dla każdego nowego API dodać test \emph{compile-and-link},
      a nie tylko sprawdzenie \code{__cplusplus}.
\item Rozróżnić funkcje opcjonalne (np. wygodniejszy wydruk) i wymagane
      dla poprawności (np. określony model zarządzania pamięcią).
\item Zachować wspierany wariant alternatywny, jeśli API ma nadal działać
      z C++20 lub starszym środowiskiem.
\item Wprowadzić osobne zadania CI z różnymi kompilatorami,
      architekturami i wariantami standardowej biblioteki.
\item Przed zmianą interfejsu publicznego przeanalizować kompatybilność
      źródłową, ABI, wyjątki, alokację i obietnice dotyczące czasu życia.
\end{enumerate}

\begin{lstlisting}[caption={Weryfikacja obecnosci nowych funkcji biblioteki (C++23).}]
#include <iostream>
#include <version>

int main() {
#ifdef __cpp_lib_expected
    std::cout << "expected: " << __cpp_lib_expected << '\n';
#else
    std::cout << "expected unavailable\n";
#endif
#ifdef __cpp_lib_mdspan
    std::cout << "mdspan: " << __cpp_lib_mdspan << '\n';
#else
    std::cout << "mdspan unavailable\n";
#endif
#ifdef __cpp_lib_print
    std::cout << "print: " << __cpp_lib_print << '\n';
#else
    std::cout << "print unavailable\n";
#endif
}
\end{lstlisting}

\section{Zestawienie funkcjonalności i rekomendacje}
\begin{center}\footnotesize
\begin{tabularx}{\textwidth}{@{}p{0.33\textwidth}p{0.11\textwidth}X@{}}
\toprule
Funkcjonalność & Standard & Główne zastosowanie \\
\midrule
\code{expected}, \code{print} & C++23 & Jawny wynik/błąd; formatowane I/O. \\
\code{mdspan}, \code{generator} & C++23 & Widoki wielowymiarowe; leniwe zakresy korutyn. \\
\code{move_only_function} & C++23 & Wywoływalne obiekty tylko przenoszalne. \\
\code{ranges::to}, \code{views::zip} & C++23 & Składanie i materializacja zakresów. \\
Deducing this, \code{if consteval} & C++23 & Ergonomia szablonów i ewaluacja stała. \\
\code{inplace_vector}, \code{hive} & C++26 & Kontenery o szczególnych gwarancjach pamięci i stabilności. \\
SIMD, linalg, senders/receivers & C++26 & Obliczenia i asynchroniczne wykonanie. \\
Refleksja, kontrakty & C++26 & Metaprogramowanie i opis warunków poprawności. \\
\bottomrule
\end{tabularx}
\end{center}

\nota{Powyższe przyporządkowanie do wersji standardu nie jest
obietnicą dostępności w dowolnej instalacji C++23/C++26. Najlepszym
źródłem danych o narzędziu jest wynik własnego testu kompilacji,
zgodny z dokumentacją danej standardowej biblioteki.}

\section{C++23: kontenery uporządkowane o zwartej reprezentacji}

\code{std::flat_map} i \code{std::flat_set} to adaptery kontenerowe,
których klucze i wartości są przechowywane w sekwencyjnych
strukturach danych, a nie w węzłach drzewa takich jak \code{std::map}.
Pozwala to uzyskać dobrą lokalność danych i efektywne wyszukiwanie,
zwłaszcza gdy zbiór jest stosunkowo statyczny. Z drugiej strony
wstawianie lub usuwanie w środku może wymagać przesuwania elementów.

\begin{lstlisting}[caption={Mapa o sekwencyjnym przechowywaniu (C++23).}]
#include <flat_map>
#include <iostream>

int main() {
    std::flat_map<int, int> lookup;
    lookup.emplace(2, 20);
    lookup.emplace(1, 10);
    std::cout << lookup.at(2) << '\n';
}
\end{lstlisting}

Dla zastosowań numerycznych często opłaca się samemu porównać
\code{flat_map}, \code{unordered_map} i posortowany \code{vector}
z wyszukiwaniem \code{lower_bound}. Wynik zależy m.in. od liczby
rekordów, częstotliwości modyfikacji i wielkości elementów.

\section{C++23: \texttt{spanstream}, \texttt{stacktrace} i funkcje pomocnicze}

\code{std::spanstream} umożliwia operacje strumieniowe na zewnętrznym,
ciągłym buforze. Nie przejmuje jego własności, dlatego po zakończeniu
czasu życia bufora korzystanie ze strumienia jest błędem.
\code{std::stacktrace} ułatwia uzyskanie śladu wywołań; jakość
symbolizacji i obecność informacji o nazwach funkcji są zależne
od platformy, opcji kompilacji oraz sposobu linkowania.

Mniejsze, lecz praktyczne ułatwienia obejmują:
\begin{description}[style=nextline]
\item[\code{std::byteswap}] Odwrócenie kolejności bajtów dla
odpowiedniego typu całkowitego, przydatne w kodowaniu protokołów.
\item[\code{std::to_underlying}] Jawna konwersja wartości typu wyliczeniowego
na jego typ bazowy.
\item[\code{std::forward_like}] Przenoszenie wzorca kwalifikatorów
\code{const} i kategorii wartości z jednego typu na inne wyrażenie,
przydatne przy pisaniu generycznych adapterów.
\item[\code{std::bind_back}] Wygodne częściowe wiązanie argumentów
od prawej strony listy parametrów.
\end{description}

\begin{lstlisting}[caption={Konwersja enum i zamiana kolejnosci bajtow (C++23).}]
#include <bit>
#include <cstdint>
#include <iostream>
#include <utility>

enum class Status : std::uint32_t { ready = 0x12345678u };

int main() {
    auto raw = std::to_underlying(Status::ready);
    auto swapped = std::byteswap(raw);
    std::cout << std::hex << raw << ' ' << swapped << '\n';
}
\end{lstlisting}

\section{C++26: \texttt{function\_ref} i funkcje bez własności}

\code{std::function_ref} reprezentuje niewłasnościowy widok na
obiekt wywoływalny. Przydaje się, gdy funkcja przyjmuje callback
na czas jednego wywołania i nie zamierza go zachować.
W odróżnieniu od \code{std::function}, nie przejmuje własności
lambdy ani jej przechwyconych danych. Podstawowym zagrożeniem
jest zatem pozostawienie widoku po zakończeniu życia callbacku.

\begin{lstlisting}[caption={Przyklad zastosowania function\_ref (C++26; fragment).}]
#include <functional>

void repeat(int n, std::function_ref<void(int)> callback) {
    for (int i = 0; i < n; ++i)
        callback(i);
}
\end{lstlisting}

Jest to typ zgodny z celem opisanym w rozdziale 20: odróżniamy
przekazanie callbacku na czas wywołania od przechowywania go we
współdzielonym stanie zadań.

\section{C++26: interpretacja nowego modelu wykonania}

Model \emph{senders/receivers} wprowadza pojęcia opisujące wynik
operacji asynchronicznej oraz jej kontynuację. W odróżnieniu od
prostego \code{std::async}, wyrażenie złożone z wielu etapów może
przenosić informacje o błędzie i anulowaniu, a wykonanie może być
zależne od schedulera dostarczonego przez bibliotekę.

Nie należy jednak utożsamiać wyrażenia opisującego zadanie z natychmiast
uruchomioną pracą. Ważne są trzy oddzielne momenty:
\begin{enumerate}
\item utworzenie opisu obliczenia (\emph{sender});
\item połączenie z odbiorcą wyniku (\emph{receiver});
\item rozpoczęcie operacji w wybranym środowisku wykonawczym.
\end{enumerate}

Gdy aplikacja ma już własny executor i korutyny, migracja powinna
zacząć się od próbnego adaptera i testów anulowania, a nie od
mechanicznego zastąpienia wszystkich \code{future}. Biblioteki
systemowe mogą udostępniać własne executory, często o innym modelu
kolejek i innych gwarancjach zamknięcia.

\section{Eksperymentalne wdrażanie nowego API bez utraty przenośności}

Niektóre biblioteki \code{<experimental/...>} dostarczały wybrane
mechanizmy przed przyjęciem ich do standardu. Nie należy zakładać,
że ich typy, przestrzenie nazw czy ABI są identyczne ze standardowymi.
Dla kodu produkcyjnego możliwe są trzy strategie:

\begin{enumerate}
\item Pozostać przy C++20, używając bibliotek zewnętrznych dla brakujących
funkcji i oddzielając te zależności od interfejsu publicznego.
\item Wymagać C++23 i wskazanych minimalnych wersji standardowej
biblioteki, walidowanych w systemie CI.
\item Udostępniać implementację warunkową dla różnych standardów,
przy czym wszystkie warianty muszą zachowywać ten sam kontrakt i
przechodzić wspólny zestaw testów.
\end{enumerate}

Dla funkcji odwołujących się do kompilatora lub sprzętu warto
budować niezależne zestawy benchmarków. Mechanizm szybszy na
maszynie deweloperskiej może okazać się mniej korzystny na systemie
o innym cache, szerokości SIMD albo liczbie rdzeni.


\section*{Podsumowanie rozdziału}
\addcontentsline{toc}{section}{Podsumowanie rozdziału}

C++23 wyraźnie ulepsza projektowanie interfejsów (\code{expected}),
pracę na danych (\code{mdspan}), korutyny (\code{generator}) oraz
programowanie zakresowe. C++26 rozwija kontenery, wektoryzację,
wykonanie asynchroniczne i mechanizmy metaprogramowania.
Najważniejszą umiejętnością nie jest jednak mechaniczne używanie
najświeższej składni, lecz świadome odróżnianie gwarancji
standardu, stanu wdrożenia i wymagań aplikacji.

\section*{Literatura i dokumentacja}
\addcontentsline{toc}{section}{Literatura i dokumentacja}
\begingroup\sloppy
\begin{itemize}
\item \url{https://en.cppreference.com/w/cpp/23} --- zestawienie C++23.
\item \url{https://en.cppreference.com/w/cpp/26} --- zestawienie C++26.
\item \url{https://en.cppreference.com/w/cpp/compiler_support} --- wsparcie kompilatorów i bibliotek.
\item \url{https://en.cppreference.com/w/cpp/feature_test} --- makra do testowania funkcjonalności.
\item \url{https://en.cppreference.com/w/cpp/utility/expected} --- \code{std::expected}.
\item \url{https://en.cppreference.com/w/cpp/container/mdspan} --- \code{std::mdspan}.
\item \url{https://wg21.link} --- oficjalne dokumenty WG21 i propozycje standardu.
\end{itemize}
\endgroup

\end{document}
